\documentclass[11pt]{article}

\usepackage[T1]{fontenc}
\usepackage{lmodern}
\usepackage[letterpaper,margin=1in]{geometry}
\usepackage{microtype}
\usepackage{mathtools}
\usepackage{amssymb}
\usepackage{mathrsfs}
\usepackage{amsthm}
\usepackage{graphicx}
\usepackage[font=small,labelfont=bf,skip=6pt]{caption}
\usepackage{needspace}
\usepackage{booktabs}
\usepackage{tabularx}
\usepackage{array}
\usepackage{enumitem}
\usepackage{tikz}
\usetikzlibrary{arrows.meta,positioning,shapes.geometric,calc}
\IfFileExists{fontawesome5.sty}{%
  \usepackage{fontawesome5}%
  \newcommand{\GitHubIcon}{\mbox{{\small\faGithub}}\nobreakspace}%
}{%
  \newcommand{\GitHubIcon}{}%
}
\usepackage{aliascnt}
\usepackage[hidelinks]{hyperref}
\hypersetup{
  pdftitle={A Complexity Dichotomy for Complex-Valued Boolean Holant Problems with Binary Disequality},
  pdfauthor={Chenghua Liu and Boning Meng}
}
\usepackage[capitalise,nameinlink,noabbrev]{cleveref}
\usepackage{authblk}

\definecolor{repositorylink}{HTML}{1F5A94}

\renewcommand\Authfont{\normalsize}
\renewcommand\Affilfont{\small}
\setlength{\affilsep}{0.5em}
\setlength{\emergencystretch}{2em}
\setlist[itemize]{leftmargin=1.6em,itemsep=0.25em,topsep=0.3em}
\setlist[enumerate]{leftmargin=1.8em,itemsep=0.3em,topsep=0.3em}
\Urlmuskip=0mu plus 1mu\relax
\allowdisplaybreaks[2]

\theoremstyle{plain}
\newtheorem{theorem}{Theorem}[section]

\newaliascnt{lemma}{theorem}
\newtheorem{lemma}[lemma]{Lemma}
\aliascntresetthe{lemma}

\newaliascnt{proposition}{theorem}
\newtheorem{proposition}[proposition]{Proposition}
\aliascntresetthe{proposition}

\newaliascnt{corollary}{theorem}
\newtheorem{corollary}[corollary]{Corollary}
\aliascntresetthe{corollary}

\newaliascnt{claim}{theorem}
\newtheorem{claim}[claim]{Claim}
\aliascntresetthe{claim}

\theoremstyle{definition}
\newaliascnt{definition}{theorem}
\newtheorem{definition}[definition]{Definition}
\aliascntresetthe{definition}


\theoremstyle{remark}
\newaliascnt{remark}{theorem}
\newtheorem{remark}[remark]{Remark}
\aliascntresetthe{remark}

\newenvironment{proofstrategy}
  {\par\noindent\textit{Proof strategy.}\ }
  {\par}

\crefname{theorem}{Theorem}{Theorems}
\Crefname{theorem}{Theorem}{Theorems}
\crefname{lemma}{Lemma}{Lemmas}
\Crefname{lemma}{Lemma}{Lemmas}
\crefname{proposition}{Proposition}{Propositions}
\Crefname{proposition}{Proposition}{Propositions}
\crefname{corollary}{Corollary}{Corollaries}
\Crefname{corollary}{Corollary}{Corollaries}
\crefname{claim}{Claim}{Claims}
\Crefname{claim}{Claim}{Claims}
\crefname{definition}{Definition}{Definitions}
\Crefname{definition}{Definition}{Definitions}
\crefname{remark}{Remark}{Remarks}
\Crefname{remark}{Remark}{Remarks}

\newcommand{\AlgNums}{\overline{\mathbb Q}}
\newcommand{\Ftwo}{\mathbb F_2}
\newcommand{\cF}{\mathcal F}
\newcommand{\cG}{\mathcal G}
\newcommand{\cS}{\mathcal S}
\newcommand{\cA}{\mathscr A}
\newcommand{\cP}{\mathscr P}
\newcommand{\cL}{\mathscr L}
\newcommand{\cM}{\mathscr M}
\newcommand{\cT}{\mathscr T}
\newcommand{\cC}{\mathscr C}
\newcommand{\cE}{\mathscr E}
\newcommand{\angles}[1]{\left\langle #1\right\rangle}
\newcommand{\Holant}{\operatorname{Holant}}
\newcommand{\KHolant}{\operatorname{K\text{-}Holant}}
\newcommand{\EOTract}{\mathsf{EOTract}}
\newcommand{\Tract}{\mathsf{Tract}}
\newcommand{\TractX}{\mathsf{Tract}_{X}}
\newcommand{\PresK}{\mathsf{Pres}_{K}}
\newcommand{\Stable}{\mathsf{Stable}}
\newcommand{\Even}{\mathsf{Even}}
\newcommand{\Close}{\mathsf{Close}}

\newcommand{\FLag}{\mathsf{FLag}}
\newcommand{\FP}{\mathsf{FP}}
\newcommand{\SharpP}{\#\mathsf P}
\newcommand{\PGL}{\operatorname{PGL}}
\newcommand{\GL}{\operatorname{GL}}
\newcommand{\Mat}{\operatorname{Mat}}
\newcommand{\diag}{\operatorname{diag}}
\newcommand{\Span}{\operatorname{span}}
\newcommand{\rank}{\operatorname{rank}}
\newcommand{\wt}{\operatorname{wt}}
\newcommand{\supp}{\operatorname{supp}}
\newcommand{\arity}{\operatorname{arity}}
\newcommand{\EO}{\mathrm{EO}}
\newcommand{\Eq}{\mathrm{EQ}}
\newcommand{\Neq}{\mathrm{NEQ}}
\newcommand{\eps}{\varepsilon}
\newcommand{\one}{\mathbf 1}
\newcommand{\zero}{\mathbf 0}
\newcommand{\DeltaZero}{\Delta_0}
\newcommand{\DeltaOne}{\Delta_1}
\newcommand{\RMcore}{\mathrm{RM}_8}
\newcommand{\Hcore}{\mathrm{H}_6}
\newcommand{\leT}{\leq_{\mathrm T}}
\newcommand{\equivT}{\equiv_{\mathrm T}}
\DeclarePairedDelimiter{\abs}{\lvert}{\rvert}
\DeclarePairedDelimiter{\norm}{\lVert}{\rVert}

\newcolumntype{L}[1]{>{\raggedright\arraybackslash}p{#1}}

\begin{document}

\title{A Dichotomy for Complex Boolean Holant\\[0.15em]
with Binary Disequality}

\author[1]{Chenghua Liu}
\author[2]{Boning Meng}

\affil[1]{Institute of Software, Chinese Academy of Sciences, Beijing, China}
\affil[2]{University of Regensburg, Regensburg, Germany}
\affil[ ]{\texttt{liuch.russell@gmail.com}, \texttt{mengboning2013@gmail.com}}
\date{}

\maketitle

\begin{abstract}
We classify Boolean Holant problems with a supplied binary disequality
signature and an arbitrary fixed finite set of algebraic complex-valued
signatures.  Each problem is polynomial-time computable or
\(\SharpP\)-hard, and the tractability criterion is decidable from exact
algebraic descriptions.  After a fixed change of coordinates, the criterion
has four alternatives: tensor products of signatures of arity at most two,
or one common affine, product, or local-affine presentation.
The proof follows the smallest nonbinary tensor-prime signature together
with the finite projective group of realizable binary transfers.
Pair contractions connect this local information to a factorization,
a smaller prime, or a common tractable presentation.
The resulting cyclic, dihedral, Klein-four, and Platonic cases combine
structural arguments with specified exact finite classifications.
We also prove a signature-set dichotomy when an endpoint-nondegenerate,
tensor-prime eight-vertex signature is available; this result supplies
the four-port boundary of the main proof.
\end{abstract}

\section*{Acknowledgments}
The authors used OpenAI's ChatGPT in preparing
this manuscript, including for language editing and \LaTeX{} preparation.
ChatGPT also contributed to the exploratory development of most of the
arguments.  All companion verification code was written by OpenAI Codex and
subsequently reviewed by the authors.  All claims, proofs, and computational
results were independently checked and finalized by the authors, who take
full responsibility for the content.

\section*{Code Availability}
The accompanying Lean formalization contains the public mathematical
specifications in \path{Holant/Statements.lean}, their proofs in
\path{Holant/Proofs.lean}, and exact-input recognition results in
\path{Holant/Decidability.lean}.  Supplementary lemmas and exact finite
certificates are included through \path{Holant/Library.lean}.
The statement-only reference \path{Verification/Challenge.lean} and
\path{Verification/comparator.json} support independent comparison using
Lean's Comparator.  Build and verification instructions are given in the
repository README.

The original Python companion repository is
\href{https://github.com/liuchliuch/Holant-X-Dichotomy}%
{\textcolor{repositorylink}{\GitHubIcon\textit{Holant-X-Dichotomy}}}.
References below to \path{certificates/} and its Python scripts refer to
that companion repository.  The Lean release supplies its own exact
certificates under \path{Holant/Certificates/}; the Python companion is
not required to build or check the Lean proofs.

\clearpage
\tableofcontents
\clearpage
\section{Introduction}
\label{sec:p1-introduction}

A Boolean signature of arity \(n\) is a function
\(f:\{0,1\}^n\to\mathbb C\).  In a Holant instance, a graph vertex
of degree \(n\) carries such a function, with its incident edges ordered.
For a fixed finite set \(\cF\) of signatures, the problem
\(\Holant(\cF)\) evaluates
\[
  \sum_{\sigma:E\to\{0,1\}}
  \prod_{v\in V}f_v\bigl(\sigma|_{E(v)}\bigr),
  \qquad f_v\in\cF.
\]
Here \(E(v)\) is the ordered list of edge incidences at \(v\);
\cref{subsec:holant-definitions} gives the conventions for loops and
multiple edges.  The graph and its vertex labels form the input; the set
\(\cF\) is fixed.  This framework originated in holographic algorithms
\cite{Valiant2008,CaiLuXia2009}.  A complexity dichotomy specifies which
fixed signature sets admit polynomial-time evaluation and proves
\(\SharpP\)-hardness for all remaining sets.

We study arbitrary algebraic complex-valued signatures with the binary
disequality relation supplied:
\[
 X=\Neq_2=\begin{pmatrix}0&1\\1&0\end{pmatrix},
 \qquad I=\Eq_2=\begin{pmatrix}1&0\\0&1\end{pmatrix}.
\]
Supplying \(X\) means that any number of degree-two vertices may carry
this signature.  It enforces opposite values on its two incident edges.
Our main result is a decidable dichotomy for \(\Holant(\cF,X)\), where
\(\Holant(\cF,X)\) abbreviates \(\Holant(\cF\cup\{X\})\).

\paragraph{Background and scope.}
Several earlier classifications prescribe auxiliary signatures.
The model \(\Holant^*\) supplies every unary signature
\cite{CaiLuXia2011HolantStar}, whereas \(\Holant^c\) supplies the two
pinnings \(\DeltaZero=[1,0]^{\mathsf T}\) and
\(\DeltaOne=[0,1]^{\mathsf T}\).  The latter model was classified for
symmetric complex signatures by Cai, Huang, and Lu
\cite{CaiHuangLu2012Holantc}, for arbitrary real signatures by Cai, Lu,
and Xia \cite{CaiLuXia2018RealHolantc}, and for algebraic complex
signatures by Backens \cite{Backens2021,Backens2025}.
The real-valued classification introduced the local-affine tractable
family \(\cL\), which also occurs here.  Binary disequality supplies
a relation between two bits, while a pinning fixes one bit.  Accordingly,
our reductions account explicitly for any unary information extracted
from the supplied binary relation.

Without a prescribed auxiliary signature, dichotomies are known for
complex symmetric signatures \cite{CaiGuoWilliams2016}, nonnegative
signatures \cite{LinWang2018}, and arbitrary real Boolean signatures
\cite{ShaoCai2020}.  Meng, Wang, Xia, and Zheng classified the complex
case when a nonzero odd-arity signature is present and established a
tensor-factor extraction theorem \cite{MengWangXiaZheng2025}.
Their two support-defined tractable alternatives have polynomial-time
algorithms by Guan, Shao, and Shi \cite{GuanShaoShi2026}.
The main remaining work in our proof concerns sets whose nonbinary
tensor-prime signatures all have even arity.

There is also a precedent for prescribing disequality in counting CSP.
In \(\#\mathrm{CSP}_d\), each variable occurs a number of times divisible
by \(d\).  Cai, Fu, and Shao classified
\(\#\mathrm{CSP}_d(\Neq_2,\cF)\) for complex signatures
\cite[full version, Section~5.1 and Theorem~5.3]{CaiFuShao2020Entanglement};
Shao and Cai used its hardness conclusion in their real Holant dichotomy
\cite[full version, Theorem~2.24]{ShaoCai2020}.
Here the instance is an arbitrary Holant graph: the supplied resource
is \(\Neq_2\), with no degree-divisible family of equality constraints.

\paragraph{Two presentations and the main theorems.}
It is useful to place \(X\) on every edge and give each incident
vertex its own bit.  We denote this presentation by \(\KHolant\).
Joining two bits \(a,b\) then contributes \(X_{ab}\); in ordinary
Holant the corresponding edge contributes \(I_{ab}\).
Let \(\equivT\) denote polynomial-time Turing equivalence and fix
\[
 K=\frac1{\sqrt2}\begin{pmatrix}1&1\\i&-i\end{pmatrix},
 \qquad K^{\mathsf T}K=X,
 \qquad K(\cS)=\{K^{\otimes\arity(f)}f:f\in\cS\}.
\]
The basis change acts on every vertex and on the edge pairing, giving
\(\KHolant(\cS)\equivT\Holant(K(\cS))\).
A separate wire substitution gives
\begin{equation}
 \Holant(\cF,X)\equivT\KHolant(\cF,I).
 \label{eq:intro-wire-bridge}
\end{equation}
In this second equivalence the signatures in \(\cF\) retain their
entries: paths of degree-two \(X\)-vertices are replaced according to
their parity, and closed components contribute known scalars.
\Cref{lem:p1-equality-wire-equivalence} proves both substitutions.

The predicate \(\Tract\) in \cref{subsec:tract-predicate} records the
seven tractable alternatives in the odd-arity classification, with the
polynomial-time improvements just cited.  On a set containing
\(Z=K^{\otimes2}I=\diag(1,-1)\), it reduces to the four-alternative
predicate \(\TractX\) by \cref{lem:tract-X-anchor}.
These alternatives are tensor products of signatures of arity at most
two, or transformability into the affine, product, or local-affine class.
For a transformable set, one basis must work for the edge and every
signature; the precise definitions and decision procedure are in
\cref{subsec:standard-classes,subsec:predicate-algorithms}.

\begin{theorem}[Binary-disequality Holant dichotomy]
\label{thm:p1-deq2-dichotomy}
For every finite set of algebraic complex-valued Boolean signatures \(\cF\),
the problem \(\Holant(\cF,X)\) is polynomial-time computable when
\(\TractX(K(\cF\cup\{I\}))\) holds and is \(\SharpP\)-hard otherwise.
The tractability criterion is decidable from exact algebraic input.
\end{theorem}

By \cref{eq:intro-wire-bridge}, the ordinary-Holant statement is
equivalent to the following formulation with native disequality edges.

\begin{theorem}[Literal-equality \(\KHolant\) dichotomy]
\label{thm:p1-eq2-kholant-dichotomy}
For every finite set \(\cF\) of algebraic complex-valued Boolean signatures
with \(I=\Eq_2\in\cF\), the problem \(\KHolant(\cF)\) is polynomial-time
computable when \(\TractX(K(\cF))\) holds and is \(\SharpP\)-hard otherwise.
The tractability criterion is decidable from exact algebraic input.
\end{theorem}

The proof also needs a four-port classification without a supplied
\(I\)-vertex.  A quaternary signature \(q\) is
\emph{endpoint-nondegenerate} when
\[
 q(x)=0\quad\text{for odd }\wt(x),\qquad q(0000)q(1111)\ne0.
\]
Here \(\wt(x)\) is the number of ones in the Boolean word \(x\).
A nonzero signature is \emph{tensor-prime} if it has no factorization
\(a\otimes b\) on two nonempty blocks of ports, allowing a port
permutation and a nonzero scalar.

\begin{theorem}[Quaternary \(\KHolant\) dichotomy]
\label{thm:p1-quaternary-kholant-dichotomy}
Let \(\cF\) be a finite set of algebraic complex-valued Boolean signatures.
Suppose that \(\cF\) contains an endpoint-nondegenerate tensor-prime
quaternary signature.  Then \(\KHolant(\cF)\) is polynomial-time computable
if \(\Tract(K(\cF))\) holds and is \(\SharpP\)-hard otherwise.  The
tractability criterion is decidable from exact algebraic input.
\end{theorem}

\Cref{sec:p1-eight-vertex} proves this signature-set theorem using the
single-signature eight-vertex boundary of Cai and Fu
\cite[Theorem~3.1]{CaiFu2023}.  Tensor primality removes its
factorization outcome.  The full predicate \(\Tract\) is needed here
because the signature set need not contain the binary anchor \(I\).
For the other two theorems, \cref{proofsec:matching-deck-synthesis}
proves the structural classification used in the hardness argument,
and \cref{sec:p1-synthesis} then gives a separate proof of each stated
dichotomy.  Its literal-equality proof uses the given set \(\cF\);
its binary-disequality proof applies that result to
\(\cF\cup\{I\}\) through \cref{eq:intro-wire-bridge}.
The decision assertion in all three statements is supplied by
\cref{thm:predicate-algorithms}.

\paragraph{The proof's main obstacle.}
A contraction may reveal a useful small tensor while giving little
information about the other signatures.  A tractable conclusion,
however, requires one presentation of the entire signature set.
We therefore retain the signatures already available whenever a gadget
or a tensor factor is added.  A realized gadget gives a direct reduction;
the tensor-factor theorem gives a Turing reduction to the retained
problem.  \Cref{sec:retained-normalization} records these two operations
and the common tractability witness they must preserve.

The first stage handles nonzero odd signatures, singular binaries,
infinite-order binary transfers, and suitable equalities using
\cref{thm:external-odd,thm:external-factor,lem:binary-interpolation,%
lem:equality-anchor,lem:pure-ge}.  The weighted Eulerian-orientation
and eight-vertex classifications provide further reductions
\cite{CaiFuXia2018SixVertex,CaiFu2023,CaiFuShao2020ARS,
MengWangXia2025EO,GuanShaoShi2026}.
Weighted Eulerian orientations use \(X\)-edges and signatures supported
on their central Hamming layer, so their statements can be applied in
the same coordinates.  At the four-port boundary the resulting
signature is checked within the whole retained set.  A common
local-affine presentation is kept as such unless the hypotheses of
\cref{lem:p1-eight-vertex-local-affine-collapse} give an affine one.

For the remaining cases, let \(\Lambda\) be the retained signature set
and let \(f\in\Lambda\) have the smallest arity among its nonbinary
tensor-prime signatures.  An ordered binary with matrix \(B\) has
transfer matrix \(T_B=BX\), so chaining binaries multiplies their
transfers.  After the preceding reductions, the directly realizable
projective transfer classes form a finite group
\[
 G\le\PGL_2(\AlgNums),\qquad
 x=[X]\in G,\qquad
 j=[Z]\in N_{\PGL_2(\AlgNums)}(G).
\]
Square brackets identify matrices differing by a nonzero scalar;
\(N(G)\) consists of projective matrices whose conjugation preserves
\(G\).  The first mark comes from the supplied \(I\)-vertex, and the
second describes reversal of the two ordered ports.
This uses Xia's finite-group approach
\cite[Section~4 and Table~2]{Xia2026FrameworkV1}, together with the
arity descent used in real Holant \cite{ShaoCai2020}.
The relative position of \(x,j\), as well as the abstract group,
determines which matrices the subsequent constructions may use.

We examine the \emph{cards} of \(f\): tensors obtained by contracting
two selected ports against available binary kernels, with allowed
binary transformations on the remaining ports.  Closing all such
contractions and their factor decompositions either resolves the
problem, lowers the chosen prime's arity, enlarges the complete group
\(G\), or leaves every card zero or a tensor product of group binaries
on a matching of its ports.  In the last case the pair \((f,G)\) is
stable in the precise sense of \cref{sec:weighted-equality}.
The coefficient-field bound on finite groups and the finite factor
worklist prove termination; \cref{lem:token-coverage} accounts for
the additional constructions used in the individual cases.

\paragraph{Recovering a tensor from its cards.}
At a stable pair, every proper two-port contraction has a binary
matching decomposition, although \(f\) itself is tensor-prime.
This discrepancy supplies the rigidity argument.  If the matchings
agree, their common factors constrain the original tensor; if they
disagree, a small alternating cycle localizes the discrepancy to four
or six ports.  The four marked group families use this information
in different ways.

For dihedral groups, \cref{sec:normalized-dihedral} combines the
geometry of root-of-unity pencils with Boolean support constraints.
A pencil is a family \(u+t v\) generated by two fixed tensors.
The permitted values of \(t\) constrain both its support and its
coefficient ratios.  A six-port classification isolates the possible
small configurations; alternating-cycle localization and Laurent
interpolation then recover binary factors at larger arities.

For proper cyclic groups, \cref{sec:proper-deck} first constructs a
four-port output or a binary reflection outside the current group.
At higher arity, the contraction supports reconstruct the first-order
Reed--Muller code \(RM(1,3)\), whose words are evaluations of affine
Boolean functions on \(\mathbb F_2^3\).  The subsequent network
construction realizes sums over specified affine fibres.  These
sums test whether every retained signature has the same tractable
presentation; \cref{lem:rm-radon-accessible} supplies their actual
Holant realization.

For the two Klein-four forms, \cref{sec:v4-deck} uses different
contraction bases.  Standard Pauli contractions determine a tensor
line above arity four, allowing a non-real tensor to be reduced to
a non-real four-port tensor.  If all retained signatures are real
up to individual scalars, the real Holant dichotomy applies to the
whole set.  In the other form, the six-port classification and
explicit crossed contractions exclude a stable nonbinary prime.
The four-copy construction from \(\Hcore\) to \(\RMcore\) connects
this analysis to the tetrahedral case.

For Platonic groups, \cref{sec:platonic-decks} represents binaries by
four quaternion coordinates.  Incidences of the resulting finite
root configurations control the linear maps that send a kernel to
its contracted card.  A common matching gives a tensor factor;
three different matchings on a projective line identify a four-port
block and reduce the parent to six ports.  Octahedral and icosahedral
outputs then reach a four-port boundary or produce a new binary.
For the tetrahedral group, one permutation class reaches that
boundary; the other uses stabilizer codes to shorten the tensor
while preserving failure of affine membership.  An eight-port
classification completes the induction to a common affine
presentation of \(\Lambda\).

The low-arity classifications that still require exact finite
enumeration are stated in
Appendix~\ref{sec:exact-finite-certificates}, with their mathematical
inputs, outputs, and completeness arguments.  The general arity
arguments and the reductions linking them are proved in the text.

\paragraph{Organization.}
\Cref{sec:preliminaries} defines the signatures, coordinate changes,
tractable classes, and mathematical tools used throughout.
\Cref{sec:p1-shared-terminals,sec:p1-eight-vertex} establish the unary,
equality, and four-port reductions.
\Cref{sec:weighted-equality} sets up the finite-group and contraction
analysis; \cref{sec:normalized-dihedral,sec:proper-deck,sec:v4-deck,%
sec:platonic-decks} prove its four branches.
\Cref{sec:p1-synthesis} completes the main proofs.
\section{Preliminaries}
\label{sec:preliminaries}

We fix the tensor and reduction conventions first, then state the tractable
classes and imported dichotomies.  The later subsections introduce the
specialized geometric and stabilizer tools at the level used in the proofs.

\subsection{Foundational algebraic conventions}
\label{subsec:foundational-algebra}

\subsubsection{Sets, maps, and exact fields}
\label{subsubsec:set-field-exact}

We write \([n]=\{1,\ldots,n\}\), \([0]=\varnothing\), and
\(\one[\mathsf P]\in\{0,1\}\) for the indicator of a proposition.
For a finite set \(S\), \(|S|\) is its cardinality and \(S^c\) its
complement in the stated ambient set.  The symbols
\(\dot\cup\) and \(\bigsqcup\) denote disjoint unions.
A restriction is written \(f|_S\), and
\(\pi^{-1}(y)=\{x:\pi(x)=y\}\) is the fibre of a map \(\pi\).
Tuples and Boolean words use lexicographic order.

For a field \(L\), write
\[
 L^\times:=L\setminus\{0\},\qquad
 \Mat_{m\times n}(L):=\{m\text{-by-}n\text{ matrices over }L\},\qquad
 \Mat_n(L):=\Mat_{n\times n}(L).
\]
We use \(\mathbb Z_m=\mathbb Z/m\mathbb Z\) and the binary field
\(\Ftwo=\mathbb Z_2\).  The function \(\varphi(m)\) counts the residue
classes modulo \(m\) coprime to \(m\).

Fix the field \(\AlgNums=\overline{\mathbb Q}\subseteq\mathbb C\) of
algebraic numbers.  An algebraic closure \(\overline L\) of a field \(L\)
contains only elements algebraic over \(L\), and every nonconstant
polynomial over \(\overline L\) has a root there.  We write
\([L':L]=\dim_L L'\) for a finite field-extension degree and \(L(\alpha)\)
for the field generated by \(L\) and \(\alpha\).  A number field is a
finite extension of \(\mathbb Q\); the compositum of finitely many
subfields of \(\mathbb C\) is the field they generate.  An
\(L\)-conjugate of \(\alpha\) is its image under an embedding
\(L(\alpha)\hookrightarrow\overline L\) fixing \(L\), whereas
\(\bar z\) always denotes complex conjugation.
An effective number field supports exact arithmetic and equality tests.
One representation is \(\mathbb Q(\theta)\), with the monic minimal
polynomial of \(\theta\) and a rational isolating region selecting its
complex root \cite{Cohen1993}.

\subsubsection{Linear, polynomial, and tensor conventions}
\label{subsubsec:linear-polynomial-tensor}

All vector spaces are finite-dimensional unless stated otherwise.
The notation \(\Span_L S\) specifies the field of linear combinations;
\(W^\vee=\operatorname{Hom}_L(W,L)\) is the dual space.  For a linear map
\(T:U\to V\), we write
\[
\begin{aligned}
 \ker T&:=\{u:T(u)=0\},&
 \operatorname{im}T&:=\{T(u):u\in U\},\\
 \rank T&:=\dim\operatorname{im}T,&
 \operatorname{Hom}_L(U,V)&:=\{T:U\to V:T\text{ is \(L\)-linear}\}.
\end{aligned}
\]
For subspaces \(U,V\) of \(W\), the notation \(W=U\oplus V\) means that each
\(w\in W\) has a unique expression \(u+v\) with \(u\in U\) and \(v\in V\).
For \(w\in W\), put \(w+U:=\{w+u:u\in U\}\).  The quotient is
\(W/U:=\{w+U:w\in W\}\), and, when \(W\) is
finite-dimensional, \(\operatorname{codim}_W U:=\dim W-\dim U\).  An
\emph{affine subspace}
is a coset \(a+U\), its translation space is \(U\), and an affine hyperplane
has codimension one.

For a square matrix \(A\), nonsingular means \(\det A\ne0\).  Define
\(\GL_n(L):=\{A\in\Mat_n(L):\det A\ne0\}\) and, for an \(L\)-space \(V\),
\(\GL(V)\) as the group of invertible \(L\)-linear maps \(V\to V\).
Matrix rank is the dimension of the column space.  Also \(A^{\mathsf T}\) is the transpose,
\(\operatorname{tr}A:=\sum_i a_{ii}\),
\(I_n=(\one[i=j])_{i,j=1}^n\), and a \emph{minor} is the determinant of a
square submatrix.  A \emph{pivot} is a selected nonzero entry used in
elimination or normalization.  An eigenpair satisfies \(Av=\lambda v\)
with \(v\ne0\); an eigenbasis is a basis of eigenvectors, and a matrix is
diagonalizable when it has one.  A matrix \(N\) is nilpotent when \(N^k=0\)
for some \(k\ge1\), and
\(J_k(\lambda)=(\lambda\one[i=j]+\one[j=i+1])_{i,j=1}^k\) is a Jordan
block.  An \(L\)-matrix algebra is an \(L\)-subspace of \(\Mat_n(L)\)
closed under multiplication; it is unital when it contains \(I_n\), and
transpose-stable when it is closed under transpose.
For a matrix whose rows and columns have an explicitly ordered index set
\(J\), the matrix unit \(E_{ij}\), \(i,j\in J\), has its sole \(1\) in
position \((i,j)\).  Thus Boolean-indexed matrices use
\(E_{00},E_{01},E_{10},E_{11}\), while ordinary one-based indexing uses
\(E_{11},E_{12},\ldots\); the stated index set distinguishes the convention.

A monomial is \(x^u=\prod_jx_j^{u_j}\) for
\(u\in\mathbb Z_{\ge0}^n\), with total degree \(\sum_ju_j\).  The polynomial
ring \(L[x_1,\ldots,x_n]\) consists of the finite sums
\(\sum_uc_ux^u\) with \(c_u\in L\).  A polynomial is \emph{homogeneous}
when all its nonzero
monomials have one total degree.  For \(p(t)=\sum_ja_jt^j\), its formal
derivative is \(p'(t)=\sum_{j\ge1}ja_jt^{j-1}\); a repeated root is a common
root of \(p\) and \(p'\), and \(p\) is square-free when it has none.

The word \emph{multilinear} means linear in each argument with the others
fixed.  The tensor product \(V_1\otimes\cdots\otimes V_r\) comes with a
multilinear map
\(\tau(v_1,\ldots,v_r)=v_1\otimes\cdots\otimes v_r\) such that every
multilinear map \(m:V_1\times\cdots\times V_r\to W\) factors uniquely as
\(m=\widetilde m\circ\tau\) through a linear map \(\widetilde m\).  The
displayed tensor \(v_1\otimes\cdots\otimes v_r\) is a \emph{simple tensor}.  The tensor power \(A^{\otimes r}\) applies \(A\)
to every factor.  A bipartite \emph{flattening} groups selected tensor
coordinates into a row index and the complementary coordinates into a
column index; its flattening rank is the resulting matrix rank.

\subsubsection{Groups and projective spaces}
\label{subsubsec:group-projective}

A homomorphism preserves the displayed group multiplication, and an
isomorphism is a bijective homomorphism.  A subgroup is written \(H\le G\),
and \(H\lhd G\) means \(gHg^{-1}=H\) for all \(g\in G\).  A left coset is
\(gH\), its index \([G:H]\) is the number of left cosets, and, for
\(N\lhd G\), the quotient group \(G/N=\{gN:g\in G\}\) has multiplication
\((gN)(hN)=ghN\).  The subgroup generated by \(S\) is \(\langle S\rangle\);
\(\langle S:\mathcal R\rangle\) denotes a presentation by generators and
relations.

The order of \(g\in G\) is the least positive \(m\) with \(g^m=1\), or is
infinite; an involution has order two.  The order of a finite group is its
cardinality, and its exponent is the least common multiple of its element
orders.  The center is \(Z(G):=\{z:zg=gz\text{ for all }g\in G\}\).  A
conjugacy class is \(\{gxg^{-1}:g\in G\}\), and a nontrivial group is simple
when its only normal subgroups are \(1\) and \(G\).  An automorphism is an
isomorphism \(G\to G\); the automorphisms form \(\operatorname{Aut}(G)\).
An automorphism is inner when it has the form \(g\mapsto hgh^{-1}\); the
inner automorphisms form the normal subgroup \(\operatorname{Inn}(G)\) of
\(\operatorname{Aut}(G)\), and
\(\operatorname{Out}(G):=\operatorname{Aut}(G)/\operatorname{Inn}(G)\).

A group action of \(G\) on \(X\) is a map \((g,x)\mapsto gx\) satisfying
\(1x=x\) and \((gh)x=g(hx)\).  Its orbit and stabilizer are
\(Gx:=\{gx:g\in G\}\) and \(G_x:=\{g:gx=x\}\).  The action is transitive
when it has one orbit and two-transitive when it is transitive on ordered
pairs of distinct points.  A representation is a group action by linear maps,
equivalently a homomorphism \(G\to\GL(V)\), and is faithful when that
homomorphism is injective.  Given an action of \(H\) on \(N\) for which every
map \(n\mapsto h\cdot n\) is an automorphism, the semidirect product
\(N\rtimes H\) has multiplication
\((n,h)(n',h')=(n(h\cdot n'),hh')\).  An element is torsion when it has
finite order; a group has bounded torsion when there is an \(m\ge1\) such
that \(g^m=1\) for every element \(g\).  For \(H\le G\), its normalizer and centralizer are
\[
 N_G(H):=\{g:gHg^{-1}=H\},\qquad
 C_G(H):=\{g:gh=hg\text{ for every }h\in H\}.
\]
For groups \(G\) and \(H\), the direct product \(G\times H\) has
componentwise multiplication \((g,h)(g',h')=(gg',hh')\).  The symbols
\(S_m,A_m,C_m,D_{2m}\), and \(V_4\) denote the symmetric group,
its subgroup of permutations with an even number of inversions, the cyclic group
\(\langle r:r^m=1\rangle\), the dihedral group
\[
 D_{2m}=\langle r,s:r^m=s^2=1,\ srs=r^{-1}\rangle,
\]
and \(C_2\times C_2\), respectively.  In \(D_{2m}\), elements of
\(\langle r\rangle\) are rotations and those in its other coset are
reflections.
For any finite set \(V\), \(\operatorname{Sym}(V)\) is the group of all
permutations of \(V\) under composition.  A \emph{transposition} swaps two
points and fixes the others; a \emph{double transposition} is the product of
two transpositions with disjoint moved points.

For a finite-dimensional \(L\)-space \(V\), its projective space is
\(\mathbb P_L(V):=(V\setminus\{0\})/L^\times\); when \(L=\mathbb C\) or is
clear we write \(\mathbb P(V)\).  The class of \(v\ne0\) is \([v]\), and
the line through distinct \([u],[v]\) is
\(\mathbb P(\Span_L\{u,v\})\).  A projective subspace \(\mathbb P(U)\) has
dimension \(\dim U-1\).  An invertible linear map induces a projective
automorphism, and scalar multiples induce the same one.  We put
\[
 \PGL(V):=\GL(V)/(L^\times I_V),\qquad
 \PGL_n(L):=\GL_n(L)/(L^\times I_n).
\]
Here \(I_V\) is the identity linear map on \(V\).
The diagonal torus in \(\PGL_n(L)\) is the subgroup of projective classes of
nonsingular diagonal matrices.  Thus brackets denote a projective class when
their argument is a nonzero vector or matrix; group-index brackets have two
group arguments.

\subsection{Complexity conventions}
\label{subsec:complexity-conventions}

We use polynomial-time Turing reductions \(\leT\), allowing adaptive exact
oracle queries; \(\equivT\) means reduction in both directions.  The
classes \(\FP\), \(\SharpP\), and \(\FP^{\mathsf{NP}}\) have their usual
meanings.  For algebraic-valued problems, \(\SharpP\)-hardness means that
every integer-valued \(\SharpP\) function reduces to the exact algebraic
oracle, including polynomial-time postprocessing of its answers.

Signature sets, their coefficient fields, and fixed gadgets are fixed
independently of the input network.  Thus polynomial time measures the
network's encoding length; exact answers use the chosen effective number
field.  Recognition of a tractable signature set is a separate task whose
input consists of finite exact signature tables.  We prove decidability
for that task without an encoding-size complexity bound.

\subsection{Signatures, tensors, and Holant instances}
\label{subsec:holant-definitions}

We use the fields \(\AlgNums\) and \(\Ftwo\) fixed above.  For \(r\ge0\), an
arity-\(r\) Boolean signature is a function
\(f:\{0,1\}^r\to\AlgNums\), identified with its lexicographically ordered
column tensor; its \(r\) argument positions are its \emph{ports}, and their
order is part of the data.  Arity zero, one, two, three, and four are called
nullary, unary, binary, ternary, and quaternary.  We write \(\arity(f)=r\) and
\(\supp(f)=\{x:f(x)\ne0\}\).  For a Boolean
word \(x\in\{0,1\}^r\), we also write
\(\supp(x):=\{j\in[r]:x_j=1\}\); the argument type distinguishes word
support from signature support.
All binary-vector arithmetic is over \(\Ftwo\), with \(\oplus\) used when
bitwise addition should be explicit.  The symbols \(\zero^r\) and
\(\one^r\) denote the all-zero and all-one words, respectively; at \(r=0\),
both denote the unique empty word.
For scalars \(a,b\), define the diagonal matrix
\[
  \diag(a,b)=\begin{pmatrix}a&0\\0&b\end{pmatrix}.
\]

For \(\tau\in S_r\), define the port reordering
\((\tau f)(x)=f(x_{\tau^{-1}(1)},\ldots,x_{\tau^{-1}(r)})\).
With this convention, \(\tau(\sigma f)=(\sigma\circ\tau)f\).
On disjoint ordered ports, \((f\otimes g)(x,y)=f(x)g(y)\).
For \(T\subseteq[r]\), \(\one_T\in\{0,1\}^r\) denotes the incidence
vector of \(T\); in particular \(\one_\varnothing=\zero^r\).  We write
\(e_j=\one_{\{j\}}\) for the \(j\)-th standard basis vector whenever the
ambient arity is clear.
A nonzero positive-arity signature is \emph{tensor-prime} if, up to port permutation and nonzero scalar, it cannot
be written as \(a\otimes b\) with both \(a\) and \(b\) of positive arity.
Prime factorization is unique up to these operations and factor order
\cite[full version, Lemma~2.13]{CaiFuShao2020ARS}.  For any class \(\cC\) of
signatures, its tensor closure \(\angles{\cC}\) consists of tensor products
of members of \(\cC\), up to port permutation and nonzero scalar.  By
convention, it also contains zero and the arity-zero empty product.
Thus \(\angles{\cC}\) denotes tensor closure, whereas \(\langle S\rangle\)
denotes the subgroup generated by \(S\).
A \emph{tensor factorization} partitions the ordered ports into blocks and
writes the tensor, up to a recorded nonzero scalar and port permutation, as
a tensor product on those blocks.  For example, a \(2+2\) factorization of a quaternary has two binary
factors.

For \(a\in\{0,1\}^r\), define \(\delta_a:\{0,1\}^r\to\AlgNums\) by
\[
  \delta_a(x)=\one[x=a].
\]
For \(r\ge1\), a \emph{weighted generalized equality} is a signature
\begin{equation*}
  c_0\delta_a+c_1\delta_{\bar a},
  \qquad c_0,c_1\in\AlgNums,\quad c_0c_1\ne0,
\end{equation*}
where \(\bar a\) is the bitwise complement.  It is \emph{pure} when
\(a=\zero^r\).  Replacing each argument bit \(x_j\) by \(x_j\oplus a_j\)
makes it pure, but these portwise bit flips are used only when
actually available or inside a stated global transform.

The binary equality and disequality tensors are
\[
  I=\Eq_2=
  \begin{pmatrix}1&0\\0&1\end{pmatrix},
  \qquad
  X=\Neq_2=
  \begin{pmatrix}0&1\\1&0\end{pmatrix}.
\]

Graphs are finite multigraphs with ordered incidences at each vertex.
Loops and parallel edges are allowed; a loop contributes two incidences
to the degree.  Cycles include loops and two-cycles formed by parallel
edges, and are identified under reversal and cyclic changes of starting
point.  A unicyclic component has exactly one cycle.  Subdivision inserts
degree-two vertices on an edge; suppression reverses this operation.
Splicing joins two exposed incidences and removes their boundary
designations.  A factor graph records signature vertices and their
incidences with edge kernels.  A matching is a family of disjoint
two-element subsets, and is perfect when it covers its ground set.

We use a half-edge semantics throughout.  A (possibly open) signature
network \(\Omega\) consists of a finite set \(H(\Omega)\) of port half-edges,
a finite set of ordinary tensor vertices, and, for each such vertex \(v\), an
assigned signature \(f_v\) and a subset \(H_v\subseteq H(\Omega)\) equipped
with a linear order.  These subsets are pairwise disjoint, satisfy
\(\lvert H_v\rvert=\arity(f_v)\), and obey
\(H(\Omega)=\bigl(\bigsqcup_vH_v\bigr)\sqcup H_{\rm free}(\Omega)\), where
the free stubs \(H_{\rm free}(\Omega)\) are half-edges not attached to a
tensor vertex.  Every free stub is a paired boundary endpoint; no isolated
unpaired free stub is allowed.  A
symbol \(\mathsf E_W\) is a zero-cost
structural edge box for the model's default edge kernel \(W\) (\(W=X\) in
the native model and \(W=I\) in the ordinary model): it replaces one wire
pair of half-edges (whose endpoints may be ordinary or free) and contributes
the single factor \(W(a,b)\) on their endpoint bits \(a,b\).  Thus \(\mathsf E_W\)
is not an additional
tensor vertex, has no separate \(H_v\), and is never multiplied again as a
vertex factor.  Any displayed \(I\) or \(X\) supplied in the signature set
(including the tensor equal to the default kernel) is instead an ordinary
tensor vertex and is transformed and counted like the other input signatures.
An explicitly ordered tuple \(\partial\Omega\) of distinct exposed ports is
drawn from \(H(\Omega)\).  The set \(E_{\rm int}(\Omega)\) is a matching of unordered
pairs of half-edges.  An ordinary vertex port is unpaired exactly when it is
in the boundary; every other ordinary port belongs to exactly one pair.
Every free stub belongs to the boundary and to exactly one pair.  Thus a
boundary port is either an unpaired ordinary port or a paired free endpoint
of an edge box, and no half-edge belongs to two pairs.  We call these wire pairs
\(E_{\rm int}(\Omega)\) (for a closed
network they are exactly the usual internal edges).  A pair may join two
ports of the same vertex (a loop), and parallel pairs are allowed.  For each
wire pair \(e=\{h,h'\}\) choose a kernel
\(W_e\in\{I,X\}\), with \(I\) and \(X\) as fixed above, and put
\(\mathbf W=(W_e)_{e\in E_{\rm int}(\Omega)}\).  For a boundary word
\(z_{\partial\Omega}\in\{0,1\}^{|\partial\Omega|}\), interpreted in the
fixed order of \(\partial\Omega\), first form the raw table
\[
 R_\Omega(z_{\partial\Omega})
 =\sum_{\substack{\tau:H(\Omega)\to\{0,1\}\\\tau|_{\partial\Omega}=z_{\partial\Omega}}}
    \prod_v f_v(\tau|_{H_v})
    \prod_{\{h,h'\}\in E_{\rm int}(\Omega)}W_{\{h,h'\}}(\tau_h,\tau_{h'}).
\]
For the \(j\)-th boundary port, put \(A_j=I\) if it is an ordinary port,
and \(A_j=W_e\) if it is a free endpoint of the wire pair \(e\).
The boundary tensor, in the coordinates used to replace a tensor vertex,
and the closed partition function are
\begin{align}
 T_\Omega&=\left(\bigotimes_{j=1}^{|\partial\Omega|}A_j\right)R_\Omega,
       \label{eq:half-edge-boundary-tensor}\\
 Z_{\mathbf W}(\Omega)&:=T_\Omega(\varnothing)
   \qquad(\partial\Omega=\varnothing).
 \label{eq:half-edge-partition-function}
\end{align}
An open signature network, equipped with this boundary order and viewed
through its boundary tensor \(T_\Omega\), is an \emph{open gadget}.
If every boundary port is ordinary, then \(T_\Omega=R_\Omega\).
For a wire with one free endpoint, the boundary action cancels that wire
because \(W_e^2=I\): deleting the free endpoint and its wire and exposing its
ordinary mate gives the same \(T_\Omega\), in the inherited boundary order.
A wire with two free endpoints is an isolated \(\mathsf E_W\), whose
boundary tensor is \(WWW=W\).  These operations reduce every open network
to ordinary exposed ports and isolated edge boxes without adding resources.
When every wire kernel is the same matrix \(W\), we abbreviate
\(Z_{\mathbf W}\) by \(Z_W\).
In the product, \(f_v\) always means the assigned signature at an ordinary
vertex; global signature transformations act only on ordinary signature
vertices.
The edge-box notation is shorthand for one occurrence of the default edge
kernel: when an \(\mathsf E_W\) is spliced into a network, its two ports are
connected by the displayed \(W\) exactly once (equivalently, one wire pair's
kernel is replaced by \(W\)); the same \(W\) is not multiplied again.
The edge boxes just described are the formal realization of the native
two-port primitive (and of the ordinary equality primitive); the empty
network is the nullary unit \(1\), not a nontrivial two-port tensor.  The
order of every \(H_v\) is part of the input, and the order of the boundary
ports is part of an open gadget.  In particular, a native \(X\)-loop has two
independent endpoint bits; it is not a single shared edge variable.  For
example, an \(\mathsf E_X\) with two exposed endpoints \(h,h'\) has boundary
tensor \(T_{\mathsf E_X}(a,b)=X(a,b)\).  Closing its two endpoints by the
native contraction gives the known scalar \(\sum_{a,b}X(a,b)^2=2\).

The notation \(Z_{\mathbf W}\) is a bookkeeping device for the specified
edge labelling; an instance is allowed to use only the
homogeneous ordinary choice \(Z_I\) or the homogeneous native choice \(Z_X\),
unless a displayed \(I\) or \(X\) is supplied as an explicit degree-two
signature/gadget.  The ordinary problem \(\Holant(\cF)\) is \(Z_I\) on
closed networks over \(\cF\), and the native-edge problem is
\(\KHolant(\cF):=Z_X\); problem instances have no free stubs (free stubs
occur only in open gadgets).  When all wire kernels are \(I\), summing the
half-edge variables forces the two endpoints of each edge to agree and
recovers the usual one-bit-per-edge multigraph definition (including loops).
For signature sets \(\cE\) and \(\cF\), the bipartite problem
\(\Holant(\cE\mid\cF)\) is the ordinary (\(I\)-edge) half-edge problem on a
bipartite network: vertices in the left part are assigned signatures from
\(\cE\), and vertices in the right part are assigned signatures from
\(\cF\).  For \(X\)-edges no one-bit compression is made.  Subdividing every native
\(X\)-edge by a degree-two \(X\)-vertex, with ordinary equality wires on its
two sides, gives the equivalent bipartite presentation
\begin{equation*}
  \KHolant(\cF)\equivT\Holant(X\mid\cF).
\end{equation*}
The empty network is the nullary unit.  Throughout,
comma-separated signature or signature-set arguments mean adjoining their
union; for example, \(\Holant(\cF,a,b)\) abbreviates
\(\Holant(\cF\cup\{a,b\})\).

We use the unary pinnings
\[
  \DeltaZero=[1,0]^{\mathsf T},
  \qquad
  \DeltaOne=[0,1]^{\mathsf T},
\]
which are point masses forcing the incident bit to \(0\) and \(1\),
respectively.  
We also fix the sign and skew matrices
\[
  Z=\begin{pmatrix}1&0\\0&-1\end{pmatrix},
  \qquad
  Y=ZX=\begin{pmatrix}0&1\\-1&0\end{pmatrix}.
\]
A \emph{weighted equality} is a binary table \(\diag(p,q)\), with
\(p,q\in\AlgNums\); it is nonsingular when \(pq\ne0\), and literal equality
is \(I=\diag(1,1)\).  A weighted equality is never a free auxiliary
signature in this paper.  In \cref{sec:p1-eight-vertex} it may instead be
produced by an actual gadget and retained with its reduction provenance; after
that interface returns, the equality-accessible analysis uses the supplied
literal \(I\) in the original coordinates.
Define the nullary equality/unit by \(\Eq_0(\varnothing):=1\).  More generally,
for \(r\ge1\), define the arity-\(r\) equality signature by
\[
  \Eq_r(x_1,\ldots,x_r)=\one[x_1=\cdots=x_r].
\]
We write
\[
  \Holant(\cF,X)=\Holant(\cF\cup\{X\})
\]
for the same ordinary equality-wire problem with the additional binary
signature \(X\) available as a vertex.
Subdividing every ordinary edge by an \(\Eq_2\)-vertex gives
\[
  \Holant(\cF)=\Holant(\Eq_2\mid\cF).
\]

\subsection{Holographic transformations and the
  \texorpdfstring{\(K\)}{K}-basis}
\label{subsec:holographic}

Let \(i\in\mathbb C\) satisfy \(i^2=-1\), and fix the \(K\)-coordinate matrix
\begin{equation}
  K=\frac1{\sqrt2}
  \begin{pmatrix}1&1\\ i&-i\end{pmatrix},
  \qquad
  K^{-1}=\frac1{\sqrt2}
  \begin{pmatrix}1&-i\\1&i\end{pmatrix}.
  \label{eq:K-matrix}
\end{equation}

The transformations below rewrite every tensor and its incident edge
consistently, giving exact identities between partition functions.  Their
fixed algebraic coefficients lie in a finite extension of the input
coefficient field.  Known nonzero rescalings are restored by exact
postprocessing, as in \cref{subsec:complexity-conventions}.

For tensors or matrices \(A\) and \(B\), write
\[
  A\doteq B
\]
when
\[
  A=\lambda B
  \qquad\text{for some }\lambda\in\AlgNums^\times.
\]
In every gadget or reduction identity, this scalar is known, tracked, and
restored by exact algebraic postprocessing.

For \(M\in\GL_2(\AlgNums)\) and an arity-\(r\) right-side signature \(f\),
write
\[
  Mf=M^{\otimes r}f,
  \qquad
  M\cF=\{Mf:f\in\cF\}.
\]
For any signature set \(\cS\), \(M(\cS)\) means the same tensorwise image
\(M\cS\).  In particular, the expressions \(K(\cF)\) and
\(K(\cF\cup\{I\})\) in the two main-theorem predicates are not matrix
products; they are these set images.  In the native theorem the hypothesis
\(I\in\cF\) supplies the equality signature, whereas the ordinary-Holant
formulation adjoins \(I\) through the equality-wire equivalence.  For a signature
in ordinary coordinates, a hat denotes \(\widehat f=K^{-1}f\), and
\(\widehat{\cS}=K^{-1}\cS\).
For an arity-\(r\) left-side signature \(g\), juxtaposition denotes the dual
action \(gM=(M^{\mathsf T})^{\otimes r}g\).  In particular,
\((\Eq_2)M\) has matrix \(M^{\mathsf T}M\), and the holographic identity is
\begin{equation*}
  \Holant(\Eq_2\mid\cF)
  =
  \Holant((\Eq_2)M\mid M^{-1}\cF),
  \qquad
  (\Eq_2)M=M^{\mathsf T}M.
\end{equation*}
This exact identity makes the transformed edge a genuine weighted binary,
not a componentwise scalar.

We use \(Z_X(\Omega)\) for the homogeneous native-\(X\) partition function
defined in \cref{eq:half-edge-partition-function}.  We also write
\begin{equation*}
  GO(X)
  =
  \{D\in\GL_2(\AlgNums):D^{\mathsf T}XD=\lambda X
    \text{ for some }\lambda\in\AlgNums^\times\}
\end{equation*}
for legal global transformations preserving the native disequality edge
\(X\) projectively.  For a closed network, if \(D\Omega\) applies \(D\) at
every ordinary vertex and
\(D^{\mathsf T}XD=\lambda X\), then
\begin{equation}
  Z_X(D\Omega)=\lambda^{|E_{\rm int}(\Omega)|}Z_X(\Omega).
  \label{eq:GOX-global-scalar}
\end{equation}
Thus the two problems differ only by the known edge-count scalar.  The
matrix acts globally and is not thereby a local gadget.  For \(m\ge1\), define
the group of \(m\)-th roots of unity by
\[
  \mu_m=\{z\in\AlgNums:z^m=1\}.
\]
An element \(\omega\in\mu_m\) is a \emph{primitive} \(m\)-th root of
unity when \(m\) is the least positive integer satisfying \(\omega^m=1\).
For a nonsingular \(2\)-by-\(2\) matrix \(T\), its class
\([T]\in\PGL_2(\AlgNums)\) uses the projective convention of
\cref{subsubsec:group-projective}.  The \emph{projective order} of \([T]\)
is the least \(m\ge1\) for which
\(T^m=\lambda I\) for some \(\lambda\in\AlgNums^\times\), and is infinite
if no such \(m\) exists.
We use the quaternion
matrix basis
\begin{equation}
  Q_0=\begin{pmatrix}1&0\\0&1\end{pmatrix},\quad
  Q_1=\begin{pmatrix}0&i\\i&0\end{pmatrix},\quad
  Q_2=\begin{pmatrix}0&1\\-1&0\end{pmatrix},\quad
  Q_3=\begin{pmatrix}i&0\\0&-i\end{pmatrix}.
  \label{eq:quaternion-basis}
\end{equation}
It is a basis of \(\Mat_2(\mathbb C)\); moreover, \(Q_0=I\),
\(Q_j^2=-I\) for \(1\le j\le3\), and
\(Q_jQ_k=-Q_kQ_j\) for distinct \(j,k\in\{1,2,3\}\).  These relations are
what the adjective \emph{quaternion} records.

\begin{proposition}[Finite subgroups of \(\PGL_2(\mathbb C)\)]
\label{prop:finite-pgl2-classification}
Every finite subgroup of \(\PGL_2(\mathbb C)\) is conjugate to exactly one
of the following types:
\[
  C_m,\qquad D_{2m},\qquad A_4,\qquad S_4,\qquad A_5,
\]
where \(m\ge1\) for \(C_m\), \(m\ge2\) for \(D_{2m}\),
\(|C_m|=m\), and \(|D_{2m}|=2m\).  If \(\omega\) is a primitive
\(m\)-th root of unity, explicit cyclic and dihedral models are
\[
 C_m^{\rm std}
 =
 \left\{
   \left[\begin{pmatrix}1&0\\0&\omega^k\end{pmatrix}\right]:
   0\le k<m
 \right\},
\]
and
\[
 D_{2m}^{\rm std}
 =
 C_m^{\rm std}
 \cup
 \left\{
   \left[
     X\begin{pmatrix}1&0\\0&\omega^k\end{pmatrix}
   \right]:
   0\le k<m
 \right\}.
\]
The diagonal classes are the rotations, and the anti-diagonal coset
consists of the reflections.

In the basis \cref{eq:quaternion-basis}, take the tetrahedral model
\begin{equation}
\begin{aligned}
 \mathcal T
 ={}&
 \{[Q_0],[Q_1],[Q_2],[Q_3]\}\\
 &{}\cup
 \{[Q_0+\epsilon_1Q_1+\epsilon_2Q_2+\epsilon_3Q_3]:
       \epsilon_j\in\{1,-1\}\},
\end{aligned}
\label{eq:tetrahedral-domain}
\end{equation}
and the octahedral model
\begin{equation}
 \mathcal O
 =
 \mathcal T
 \cup
 \{[Q_\mu+\epsilon Q_\nu]:
       0\le\mu<\nu\le3,\ \epsilon\in\{1,-1\}\}.
 \label{eq:octahedral-domain}
\end{equation}
Finally, define the golden-ratio scalar and its reciprocal by
\[
  \phi:=\frac{1+\sqrt5}{2},
  \qquad
  \tau:=\phi^{-1},
\]
and define the icosahedral model by
\begin{equation}
\begin{aligned}
 \mathcal I
 =\mathcal T\ \cup\
 \bigl\{[v_0Q_0+v_1Q_1+v_2Q_2+v_3Q_3]:{}&
 (v_0,v_1,v_2,v_3)\\
 &{}=\pi(0,\epsilon_1,\epsilon_2\phi,\epsilon_3\tau),\
 \ \pi\in A_4,\ \epsilon_j\in\{1,-1\}\bigr\},
\end{aligned}
\label{eq:icosahedral-domain}
\end{equation}
where \(A_4\) acts by even coordinate permutations and projective
duplicates are removed.  The additional part of \(\mathcal I\) has
\(48\) projective classes.  Thus
\[
  \mathcal T\cong A_4,\qquad
  \mathcal O\cong S_4,\qquad
  \mathcal I\cong A_5,
\]
of orders \(12,24,60\), respectively.
\end{proposition}

\begin{proof}
The conjugacy classification is classical
\cite[Section~5.1]{Beardon1983}.  The first two displays are the standard
rotation and reflection models.  Direct matrix multiplication verifies
closure, cardinalities, and element-order distributions of the three
remaining displayed sets, identifying them respectively with
\(A_4,S_4,A_5\).
\end{proof}

\begin{lemma}[Finite projective-group facts]
\label{lem:finite-pgl2-facts}
The following properties will be used below.
\begin{enumerate}
\item A cyclic group has at most one involution, and it has one exactly when
      its order is even.
\item For
      \[
        D_{2m}=\langle r,s:r^m=s^2=1,\ srs=r^{-1}\rangle,
      \]
      the rotations \(\langle r\rangle\cong C_m\) form a normal index-two
      subgroup, and every element outside it is a reflection.  If \(m\) is
      odd, the rotation group is the unique index-two subgroup and the
      center is trivial.  If \(m\) is even and \(m\ge4\), the center is
      \(\{1,r^{m/2}\}\); for \(m=2\), \(D_4\cong V_4\) and all three
      nonidentity elements are central.  Every subgroup of a cyclic or
      dihedral group is cyclic or dihedral.
\item The three involutions of \(A_4\), together with the identity, form
      its unique normal \(V_4\).  The group has no index-two subgroup, and
      the centralizer of an involution is \(V_4\).
\item The unique index-two subgroup of \(S_4\) is \(A_4\).  Its
      involutions are six transpositions and three double transpositions.
      The centralizer of a transposition has order \(4\), whereas the
      centralizer of a double transposition is \(D_8\).  The three double
      transpositions are the nonidentity elements of the unique normal
      \(V_4\).  Conjugation is transitive both on ordered pairs of distinct
      double transpositions and on ordered pairs of disjoint transpositions.
\item The group \(A_5\) is simple and has no index-two subgroup.  Its
      fifteen involutions form one conjugacy class, and the centralizer of
      each is \(V_4\).
\item The element-order distributions are
      \[
      \begin{array}{c|rrrr}
          &1&2&3&\text{order }4\text{ or }5\\ \hline
        A_4&1&3&8&0\\
        S_4&1&9&8&6\text{ of order }4\\
        A_5&1&15&20&24\text{ of order }5.
      \end{array}
      \]
\item For the displayed models,
      \[
        N_{\PGL_2(\mathbb C)}(\mathcal T)=\mathcal O,
        \qquad
        N_{\PGL_2(\mathbb C)}(\mathcal O)=\mathcal O,
        \qquad
        N_{\PGL_2(\mathbb C)}(\mathcal I)=\mathcal I.
      \]
\item The projective centralizers of the distinguished involutions are
      \[
      C_{\PGL_2(\mathbb C)}([Z])
      =
      \left\{
      \left[\begin{pmatrix}a&0\\0&b\end{pmatrix}\right]:ab\ne0
      \right\}
      \sqcup
      \left\{
      \left[\begin{pmatrix}0&a\\b&0\end{pmatrix}\right]:ab\ne0
      \right\}.
      \]
      Thus \([Z]\) is the only nonidentity diagonal involution, while
      every nonsingular anti-diagonal class is an involution.  Moreover,
      \[
        C_{\PGL_2(\mathbb C)}\bigl(\langle[X],[Z]\rangle\bigr)
        =\langle[X],[Z]\rangle\cong V_4.
      \]
      Every finite subgroup of the displayed diagonal torus is cyclic.
\end{enumerate}
\end{lemma}

\begin{proof}
Items~1--6 follow directly from the cyclic and dihedral presentations and
the usual permutation models of \(A_4,S_4,A_5\)
\cite{DixonMortimer1996}.  For item~7, conjugation
embeds each normalizer modulo its projective centralizer into the
corresponding automorphism group.  The projective centralizers are trivial;
\(\operatorname{Aut}(A_4)\cong S_4\), the displayed \(\mathcal O\)
normalizes \(\mathcal T\), and every automorphism of \(S_4\) is inner.
Since \(\operatorname{Aut}(A_5)\cong S_5\), an additional outer normalizer
of \(\mathcal I\) would give a finite projective group \(H\) of order \(120\)
containing \(A_5\).  The classification forces \(H\cong C_{120}\) or
\(D_{120}\); every subgroup of a cyclic or dihedral group is cyclic or
dihedral, so \(H\) cannot contain \(A_5\).  Item~8 follows by solving
\(UZ\doteq ZU\), and then imposing the same relation with \(X\).  The ratio
map \([\diag(a,b)]\mapsto a/b\) embeds the diagonal torus into
\(\mathbb C^\times\); every finite subgroup of \(\mathbb C^\times\) is a cyclic group
of roots of unity.
\end{proof}

The conjugacy classification above is group-theoretic only.  Its
conjugating matrix need not lie in \(GO(X)\), so it does not by itself
define a legal global \(\KHolant\) normalization or provide an actual binary
gadget.  Later applications retain the distinguished classes \([X]\) and
\([Z]\).

Since \(K^{\mathsf T}K=X\), putting
\(\widehat{\cF}=K^{-1}\cF\) gives
\begin{equation}
  \Holant(\cF)
  \equivT
  \Holant(X\mid\widehat{\cF})
  \equivT
  \KHolant(\widehat{\cF}).
  \label{eq:K-equivalence}
\end{equation}
For original-coordinate signatures and signature sets, the hats in
\cref{eq:K-equivalence} record passage to \(K\)-coordinates.  Once we work
with a signature set already in \(K\)-coordinates, we suppress
these hats and denote its member signatures simply by \(f\).  Unless stated
otherwise, signatures and signature sets appearing in a \(\KHolant\) problem
are understood to be in \(K\)-coordinates, with native edge \(X\).  The Walsh
coordinates used below are analytic bookkeeping, not another legal global
transformation.  Define the Walsh matrix \(P\) and, for a signature \(f\)
already in \(K\)-coordinates, its Walsh-coordinate tensor \(\widetilde f\) by
\begin{equation}
  P:=\begin{pmatrix}1&1\\1&-1\end{pmatrix},
  \qquad
  \widetilde f:=(P^{-1})^{\otimes\arity(f)}f.
  \label{eq:Walsh-coordinates}
\end{equation}
Thus \(\widetilde f\) denotes the analytic Walsh representation of a
signature \(f\).  For
\(s\in\AlgNums\), the identities
\begin{equation*}
  P^{-1}\!\left[\frac12
  \begin{pmatrix}1+s&1-s\\1-s&1+s\end{pmatrix}\right]P
  =\diag(1,s),
  \qquad
  P^{\mathsf T}XP=2Z
\end{equation*}
explain its use in the odd cyclic and dihedral calculations.

\subsection{Gadgets, interpolation, and factor reductions}
\label{subsec:reduction-conventions}

For a signature \(f\), a \emph{card} means an output
\(\partial_p^C f\) of one of the two-port contractions defined below.

A \emph{direct gadget} over \(\cG\) is a finite native-\(X\) network of
\(\cG\)-vertices, possibly using the primitive \(\mathsf E_X\), with an
ordered exposed boundary.  If its boundary tensor
in the sense of \cref{eq:half-edge-boundary-tensor} is
\(T_\Omega=\rho h\), where \(\rho\in\AlgNums^\times\) is a known nonzero
scalar, then
\begin{equation}
  \KHolant(\cG,h)
  \leT
  \KHolant(\cG).
  \label{eq:direct-gadget-reduction}
\end{equation}
A fixed gadget has constant size.  A gadget realizing \(\rho h\),
\(\rho\ne0\), implements the same Turing reduction after an instance with
\(N\) occurrences is rescaled by \(\rho^{-N}\).
We call such an \(h\) an \emph{actual signature} over \(\cG\).  More
generally, a finite set \(S\) of signatures is \emph{Turing-exposed} over
\(\cG\) if
\[
  \KHolant(\cG,S)\leT\KHolant(\cG).
\]
In either homogeneous model, gluing two boundary ports contracts their
boundary-tensor coordinates against the default kernel \(W\); the remaining
boundary order is inherited.  Two unpaired ordinary ports receive a new
\(W\)-wire.  For one paired free endpoint and one unpaired ordinary port,
splice the latter into the former's incidence, retaining its existing wire.
For two free endpoints on distinct wires, delete those endpoints and their
two wires and join their mates by one \(W\)-wire.  If both endpoints belong
to the same isolated \(\mathsf E_W\), remove that component and record the
scalar \(\sum_{a,b}W(a,b)^2=2\).  These rules follow from \(W^2=I\) and
\(W^3=W\).  In particular, vertex replacement first cancels the one-free-end
wires as above and then inserts the resulting ordinary-boundary gadget;
each remaining isolated \(\mathsf E_W\) splices into a single wire.
Closed components containing only such edge boxes contribute known factors
of \(2\).  This proves \cref{eq:direct-gadget-reduction}, including loops,
without leaving free stubs in the oracle instance.

When transporting an open native gadget by \(D\in GO(X)\), first cancel
its one-free-end wires in this way.  Every remaining internal wire between
ordinary ports contributes the known scalar from \(D^{\mathsf T}XD=\lambda X\).
An isolated primitive remains available because \(DXD^{\mathsf T}=\lambda X\).
Thus its boundary tensor is transported by \(D\) up to a recorded nonzero
scalar; the closed-network edge count in \cref{eq:GOX-global-scalar} is used
only for closed instances.
Thus every actual signature is Turing-exposed, but a factor obtained from
the tensor-factor theorem or a signature obtained by interpolation need not
be an actual gadget.  This provenance is persistent: a gadget built over an
augmentation containing a merely Turing-exposed resource is, relative to the
predecessor, only Turing-exposed.  Slices, Walsh expansions, linear inverses,
similarities, projective normalizations, tomography, and decoder coordinates
are analytic arguments; they may prove a tensor identity or select an actual
contraction, but they do not themselves add a signature resource.
A \emph{slice} is a partial evaluation after selected tensor coordinates
have been fixed; a matrix similarity has the form \(B=SAS^{-1}\) with
\(S\) invertible; and \emph{tomography} is reconstruction from a separating
family of linear measurements, meaning that their joint value map is
injective on the tensors under consideration.  An interpolation recovers finitely many
unknown coefficients from oracle evaluations of a parameterized gadget
family by solving an invertible linear system.  Its output is Turing-exposed
unless a separate direct gadget realizes it.

Let \(f\) have arity \(r\).  For an ordered pair of distinct ports
\(p=(i,j)\), a binary kernel
\(C:\Ftwo^2\to\AlgNums\), and a residual word \(z\) indexed by
\([r]\setminus\{i,j\}\) in increasing port order, let
\(z^{i\leftarrow a,j\leftarrow b}\) be the arity-\(r\) word obtained by
inserting \(a,b\) at ports \(i,j\).  Define the two-port contraction by
\begin{equation*}
  (\partial_p^C f)(z)
  =
  \sum_{a,b\in\Ftwo}C(a,b)
  f\bigl(z^{i\leftarrow a,j\leftarrow b}\bigr).
\end{equation*}
In the notation of Shao--Cai
\cite[full version, Section~2.4]{ShaoCai2020}, their
\(\widehat\partial_{ij}\widehat f\) and \(\partial_{ij}f\) are,
respectively, our \(\partial_{(i,j)}^X\widehat f\) and
\(\partial_{(i,j)}^I f\).  Port order fixes the ordered pair \((i,j)\).
For binary matrices
\(U,V\), the Frobenius pairing is bilinear:
\begin{equation*}
  \langle U,V\rangle_F
  =\operatorname{tr}(U^{\mathsf T}V)
  =\sum_{a,b\in\Ftwo}U(a,b)V(a,b),
\end{equation*}
without conjugation; \(\operatorname{tr}\) is the matrix trace fixed in
\cref{subsubsec:linear-polynomial-tensor}.  A \emph{binary chain} is a serial
composition of ordered binary gadgets joined by native \(X\)-edges.  A
\emph{dressed card} \(\partial_p^Cf\) uses a kernel realized by actual binary
chains over the current retained set and one native edge.  Suppose the closing
subnetwork has ordered boundary tensor \(B\).  Its two attachment
\(X\)-edges complement both boundary bits, so
\begin{equation*}
  C(a,b)
  =\sum_{u,v\in\Ftwo}X(a,u)B(u,v)X(v,b)
  =(XBX)(a,b)
  =B(1-a,1-b).
\end{equation*}
Exchanging the two ordered boundary inputs replaces \(B\) by
\(B^{\mathsf T}\), and the same attachment therefore gives
\(XB^{\mathsf T}X\).  Likewise, two inward-ordered binary chains \(L\) and
\(R\) joined by one native \(X\)-edge have external tensor
\(LXR^{\mathsf T}\).  These formulas explain where transposes arise in
physical gadget constructions and fix the ordered-port convention used
below.  Every dressed card so constructed is an actual gadget; linear
combinations of cards are only analytic comparisons.

\subsubsection{Two-port slices and contractions}
\label{subsubsec:two-port-fibres}

For an arity-\(r\) tensor \(f\), ordered ports \(p=(i,j)\), and a residual
word \(z\), the \emph{two-port slice} is the binary
matrix
\[
 F^f_{p,z}(a,b):=
 f\bigl(z^{i\leftarrow a,j\leftarrow b}\bigr),
 \qquad a,b\in\Ftwo.
\]
If \(\mathscr K\le\Mat_2(\mathbb C)\) is a space of kernel directions, its
Frobenius-orthogonal space is
\[
 \mathscr K^{\perp_F}:=
 \{F\in\Mat_2(\mathbb C):
   \langle K,F\rangle_F=0\text{ for every }K\in\mathscr K\}.
\]
Thus contractions against \(\mathscr K\) determine the slice modulo
\(\mathscr K^{\perp_F}\).

Fix the preprocessed signature set \(\Gamma_\star\) at the beginning of a
structural branch.  A \emph{retained signature set} is a finite set reached
from \(\Gamma_\star\) by recorded direct gadget augmentation, retained factor
exposure, or, inside the eight-vertex interface, a legal global \(GO(X)\)
change of coordinates.  Every such step carries its reduction direction and
provenance.  Generic interpolation is terminal only, and an analytic tensor
identity does not by itself adjoin a retained signature.

We record the current coordinates and signature set as \((A,\Lambda)\), with
\[
 A\in GO(X),\qquad A\Gamma_\star\subseteq\Lambda,
 \qquad \KHolant(\Lambda)\equivT\KHolant(\Gamma_\star).
\]
A literal retained augmentation is the special case \(A=I\).  After a
global normalization the current signature set need not literally contain
\(\Gamma_\star\); the invariant is the transported inclusion
\(A\Gamma_\star\subseteq\Lambda\).  The normalization data in
\cref{lem:normalization-ledger} records the additional coefficient-field and
group data.  This transported form is used for the self-contained
eight-vertex interface in \cref{sec:p1-eight-vertex}.  Before that interface
hands a continuing branch to the matching-deck analysis, it pulls the state
back to the caller's coordinates.  From \cref{sec:p1-binary-phase} onward,
every continuing retained set contains the distinguished literal \(I\), and
no realized weighted equality replaces it as the anchor.  Any continuing
coordinate change must preserve both native \(X\) and literal \(I\)
projectively; otherwise it is used only inside a terminal reduction and is
not a state update.

A card is \emph{proper} if its arity is strictly smaller than
its parent's; \emph{actual} always means directly realized, not analytically
combined.

The zero-cost native edge box \(\mathsf E_X\) is a legal two-port gadget
with tensor \(X\); the empty network is the nullary unit.  Thus, when a
transfer group is introduced below, we take the identity-class binary gadget
\(B_1:=\mathsf E_X=X\) and define its transfer matrix
\(R_1:=B_1X=I\).  This is an
actual gadget identity rather than an implicit free signature.

For every ordered binary signature \(B\), define its transfer matrix and
transfer class by
\[
 T_B:=BX,
 \qquad [T_B]=[BX].
\]
For a finite realized projective group
\(G\le\PGL_2(\AlgNums)\), choose one
nonsingular ordered gadget \(B_g\) for each class \(g\).  Define its transfer
matrix \(R_g:=T_{B_g}=B_gX\), require \([R_g]=g\), and normalize \(R_1=I\), where
\(1=[I]\).  The matrix of
every actual ordered binary \(B\) having transfer class \(g\) is
\(B=\rho B_g\) for a unique \(\rho\in\AlgNums^\times\); we likewise record
the multiplication and reversal scalars.  Consequently all
displayed projective identities below stand for actual gadget identities
with their known scalar restored.

Whenever the proof says that an external port is X-translated, it means
attaching the retained two-port chain whose transfer class is \([X]\) and
recording its nonzero scalar in the port data.  This is a direct binary
chain, not a free unary bit flip.

For each signature, record its realizing gadget or reduction, ordered ports,
discarded nonzero scalars, and coordinate changes.  To express it in earlier
coordinates, apply the inverses of the recorded changes.

Interpolation records parameter count, gadget size, and oracle calls.
Linear combinations are not gadgets.  The tensor-factor extraction result
used later is stated as \cref{thm:external-factor}; its common-one-sided
exception is always kept explicit.

\subsection{Zero and nullary preprocessing}
\label{subsec:zero-nullary}

For a closed network \(\Omega\), define its model-dependent partition-function
value by \(Z(\Omega):=Z_I(\Omega)\) in the ordinary model and
\(Z(\Omega):=Z_X(\Omega)\) in the native model under discussion.  Let
\(\Omega^\circ\) be the network obtained by deleting all nullary vertices
(and their scalar contributions); a zero-signature vertex instead makes the
value zero.  Let the reduced signature set \(\cF^\circ\) be obtained from
\(\cF\) by deleting all zero signatures
and all arity-zero signatures.  If an instance \(\Omega\) contains a
zero-signature vertex, then \(Z(\Omega)=0\).  Otherwise, writing
\(m_c(\Omega)\) for the number of occurrences of a nullary \(c\),
\[
 Z(\Omega)=
 \left(\prod_{\substack{c\in\cF\\\arity(c)=0}}
 c^{m_c(\Omega)}\right) Z(\Omega^\circ),
\]
where the partition function of the empty residual instance is the empty
product \(1\).  Thus structural classification applies only to the nonzero
positive-arity signature set \(\cF^\circ\), while the root evaluator restores
all original nullary scalars exactly.  The tractable predicates ignore the
deleted signatures.
We call this deletion-and-restoration step \emph{preprocessing}; a
\emph{preprocessed signature set} contains only nonzero positive-arity
signatures, with all deleted nullaries handled by the root evaluator.

Algebraic numbers are represented exactly, for example by a minimal
polynomial together with an isolating rectangle selecting the intended
complex root.  All zero tests, complex conjugations
\(z\mapsto\bar z\), field operations, and oracle
answers are performed in a finite effective compositum and returned in the
chosen source embedding.  For a fixed finite signature set, the number and
size of oracle instances and the exact field arithmetic are polynomial in
the Holant instance size.  Effective recognition is asserted to be
decidable for finite exact algebraic input; no polynomial bound in the
encoding length of the fixed signature set is claimed.

\subsection{Charge and Hamming layers}
\label{subsec:charge}

All charge terminology in this subsection is in \(K\)-coordinates.
For \(x\in\{0,1\}^r\), define its Hamming weight by
\(\wt(x):=|\{j:x_j=1\}|\) and its doubled charge by
\[
  q_r(x)=2\wt(x)-r.
\]
A \emph{Hamming layer} is the set of words of one fixed Hamming weight.  A
word has even or odd parity according as its weight is even or odd, and a
signature has \emph{even support} when it vanishes on every odd-parity word.
We distinguish \(\supp f\subseteq M\) from \(\supp f=M\); coefficients
in a prescribed support set may vanish unless equality is required.
For even \(r\), we also use the integer half-charge
\[
  \chi_r(x)=\frac{q_r(x)}2=\wt(x)-\frac r2.
\]
In the all-even analysis, gcd, divisibility, height, and interval statements
use \(\chi_r\); \emph{doubled charge} explicitly means \(q_r\).
Define the Hamming half-spaces
\[
  HW_r^{\ge}=\{x:q_r(x)\ge0\},
  \quad
  HW_r^{\le}=\{x:q_r(x)\le0\},
  \quad
  HW_r^{=}=\{x:q_r(x)=0\},
\]
and define \(HW_r^{>},HW_r^{<}\) analogously.  A signature is
\(HW^{\ge}\) (respectively, \(HW^{\le}\)) when its support lies in
\(HW_r^{\ge}\) (respectively, \(HW_r^{\le}\)).  A signature set is \emph{common
one-sided} when all members are \(HW^{\ge}\) or all are \(HW^{\le}\);
otherwise it is \emph{two-sided}, equivalently both charge signs occur among
its support words.  An arity-\(2d\) \emph{EO signature} is supported in
\(HW_{2d}^=\).

For a homogeneous native-\(X\) network (all structural wire pairs have
kernel \(X\); any explicitly supplied \(I\) or \(X\) is an ordinary tensor
vertex), let
\(U(\Omega)=H(\Omega)\setminus\bigcup_{e\in E_{\rm int}(\Omega)}e\) be the
unpaired boundary half-edges (by definition,
\(U(\Omega)\subseteq\partial\Omega\)).  Call an assignment
\emph{edge-supported} when every structural edge-kernel factor is nonzero on
the two bits assigned to its endpoints.  Every edge-supported assignment has
one \(0\)-bit and one \(1\)-bit on each
\(X\)-pair; assignments violating an \(X\)-edge have zero weight and are not
contributing assignments.  Hence, pointwise for every edge-supported
assignment \(\tau\),
\begin{equation*}
  \sum_v q_{\arity(f_v)}(\tau|_{H_v})
  +\sum_{h\in H_{\rm free}(\Omega)}(2\tau_h-1)
  =\sum_{h\in U(\Omega)}(2\tau_h-1).
\end{equation*}
For a contribution to the boundary tensor \(T_\Omega(z)\), an ordinary
boundary bit equals \(\tau_h\), while a free boundary bit equals
\(1-\tau_h\) by \cref{eq:half-edge-boundary-tensor}.  Moving the free-stub
sum to the right therefore gives
\(q_{|\partial\Omega|}(z)=\sum_v q_{\arity(f_v)}(\tau|_{H_v})\).
In the standard closed instances below there are
no free stubs and \(U(\Omega)=\varnothing\), so
\(\sum_vq_{\arity(f_v)}(\tau|_{H_v})=0\).  Thus native \(X\)-contraction preserves
the total charge in every closed contributing assignment.

\begin{remark}[Notation relative to prior work]
\label{rem:mwxz-notation}
We follow \cite[Sections~2.1--2.3]{MengWangXiaZheng2025} in writing
\(\widehat{\cF}=K^{-1}\cF\), \(HW^{\ge}\), \(HW^{\le}\), \(HW^=\),
and \emph{EO signature}.  Their \(\neq_2\) and \(=_r\) are our \(X\) and
\(\Eq_r\), respectively; subscripts in \(HW_r^\bullet\) record arity.  Thus
\(q_r(x)=|\{j:x_j=1\}|-|\{j:x_j=0\}|\), while the height arguments use its
integer half \(\chi_r=q_r/2\).

Our \(\Span\) is linear span; the affine span denoted
\(\operatorname{Span}(f)\) in \cite{MengWangXiaZheng2025} is written
\(\operatorname{aff}(\supp(f))\).
\end{remark}

\subsection{Binary codes and affine geometry}
\label{subsec:binary-affine-conventions}

In \cref{sec:v4-deck}, signature supports become binary affine spaces:
contraction imposes affine equations on the support, while coefficient
ratios record the remaining phase.  The code and bilinear-form conventions
below describe these two parts separately.

\subsubsection{Boolean and affine tools}
\label{subsubsec:boolean-affine-tools}

For a finite coordinate set \(S\), put
\(\Ftwo^S:=\{x:S\to\Ftwo\}\).  A \emph{binary code} of length \(n\) is a
subset of \(\Ftwo^n\); its elements are \emph{codewords}.  It is linear when
it is a subspace and affine when it is a coset of a linear code.  The truth
table of a Boolean function lists its values in the fixed lexicographic input
order.  For a commutative ring \(R\), a Boolean polynomial over \(R\) is an
element of the quotient ring of \(R[x_1,\ldots,x_n]\) by the ideal generated
by the relations \(x_j^2-x_j\), \(j\in[n]\), and hence has a unique
multilinear representative.  Its degree is the largest number of
distinct variables in a nonzero monomial, and a mixed term contains at least
two distinct variables.  An affine-linear form is a constant plus a linear
form.  A \emph{quadratic phase} is \(i^{Q(x)}\), where
\(Q:\Ftwo^n\to\mathbb Z_4\) has multilinear degree at most two and even
coefficients on mixed terms.  A character of an additive group is a group
homomorphism into the multiplicative group of the coefficient field.  These
Boolean normal forms are instances of Boolean-lattice Möbius inversion
\cite[Chapters~1--2]{Carlet2021}.

For an affine subspace \(a+U\), \(U\) is its \emph{translation space}; the
affine hull \(\operatorname{aff}(D)\) is the smallest affine subspace
containing \(D\).  An affine map \(a+U\to b+V\) has the form
\(a+u\mapsto b+T(u)\) for a linear map \(T:U\to V\).  If an affine surjection
\(\pi:a+U\twoheadrightarrow D\) onto an affine space \(D\) is fixed, an
\emph{affine section} is an
affine map \(\sigma:D\to a+U\) with
\(\pi\circ\sigma=\mathrm{id}_D\).  For \(S\subseteq[n]\), in increasing
coordinate order, define coordinate projection and zero extension by
\[
 \pi_S:\Ftwo^n\to\Ftwo^S,\quad \pi_S(x)=(x_i)_{i\in S},
 \qquad
 \iota_S:\Ftwo^S\to\Ftwo^n,\quad
 (\iota_Sy)_i=\begin{cases}y_i,&i\in S,\\0,&i\notin S.\end{cases}
\]

For \(d\ge1\),
\[
 \operatorname{AGL}(d,2):=\Ftwo^d\rtimes\GL_d(\Ftwo)
\]
acts by \(x\mapsto Ax+b\).  An \emph{affine relabelling} of ports indexed by
\(\Ftwo^d\) is the corresponding permutation of those ports.  Indexing the
eight coordinates by \(\Ftwo^3\) in lexicographic order, define the
first-order Reed--Muller code by
\[
 RM(1,3):=
 \{(\ell(t))_{t\in\Ftwo^3}:
   \ell(t)=a+u_1t_1+u_2t_2+u_3t_3,\ 
   a,u_1,u_2,u_3\in\Ftwo\}.
\]
Thus \(RM(1,3)\) consists of the sixteen truth tables of affine-linear
functions on \(\Ftwo^3\).  The natural action satisfies
\(\operatorname{AGL}(2,2)\cong S_4\);
\(\operatorname{AGL}(3,2)\) is two-transitive and preserves \(RM(1,3)\),
because affine precomposition preserves affine truth tables
\cite[Section~4.1]{Carlet2021}.

A \emph{spread} in a binary vector space is a family of nonzero subspaces
whose nonzero vectors partition the nonzero vectors of the ambient space; a
spread coset is a coset of one of its member subspaces.  For a function
\(F:\Ftwo^d\to R\) into an abelian group, its additive derivative is
\[
 (\Delta_aF)(x):=F(x+a)-F(x).
\]
The identities \(\Delta_a\Delta_bF=0\) for all \(a,b\) hold exactly when
\(F(x)=F(0)+\ell(x)\) for a group homomorphism
\(\ell:(\Ftwo^d,+)\to(R,+)\): the identities make every
\(\Delta_aF\) constant, and the resulting
\(\ell(x):=F(x)-F(0)\) is additive.

For \(Q:\Ftwo^d\to R\), with \(R\) a commutative ring, its
\emph{Boolean Möbius expansion} is the unique multilinear expression
\[
 Q(x)=\sum_{S\subseteq[d]}\widehat Q_S\prod_{j\in S}x_j,
 \qquad
 \widehat Q_S
 =\sum_{T\subseteq S}(-1)^{|S|-|T|}Q(\one_T).
\]
The coefficient table \((\widehat Q_S)_S\) is the Boolean Möbius transform;
the formula is interpreted in the stated coefficient ring, including
\(\Ftwo\) and \(\mathbb Z_4\).  Direct substitution is the Möbius-inversion
proof of uniqueness.  This additive transform is distinct from the
multiplicative Möbius inversion used later.

\subsubsection{Symmetric binary spaces}
\label{subsubsec:symmetric-binary-spaces}

A bilinear form \(b:W\times W\to\Ftwo\) is linear in each argument and is
symmetric when \(b(x,y)=b(y,x)\).  A \emph{symmetric binary space} is a
finite-dimensional \(\Ftwo\)-space with a specified symmetric form, written
\(x\cdot y\).  On \(\Ftwo^m\), the standard dot product is
\(x\cdot y:=\sum_jx_jy_j\).  For \(U\le W\), define
\[
 U^{\perp_W}:=\{y\in W:x\cdot y=0\text{ for every }x\in U\},
 \qquad
 \operatorname{rad}_W(U):=U\cap U^{\perp_W}.
\]
We omit \(W\) when it is clear.  The restricted form on \(U\) is
\emph{nondegenerate} when \(\operatorname{rad}_W(U)=\{0\}\).
The notation \(W=U\perp P\) means \(W=U\oplus P\) and \(U\cdot P=0\);
then \(\pi_U(u+p)=u\) and \(\pi_P(u+p)=p\) are the orthogonal projections.

If \(R\le\operatorname{rad}_W(U)\), the quotient \(U/R\) has the induced
form \((x+R)\cdot(y+R):=x\cdot y\).  When \(R\) acts by translations on an affine space with
translation space \(U\), its orbit space has translation space \(U/R\).
A linear map \(\phi:U\to U'\) is an \emph{isometry} when it is bijective and
preserves the form; the corresponding affine isometry is
\(a+u\mapsto a'+\phi(u)\).

A vector \(x\) is \emph{isotropic} when \(x\cdot x=0\) and anisotropic
otherwise.  A subspace is \emph{totally isotropic} when
\(U\subseteq U^{\perp_W}\), and is \emph{self-dual in \(W\)} when
\(U=U^{\perp_W}\).  In a nondegenerate symmetric binary space, the map
\(x\mapsto x\cdot x\) is linear, so nondegeneracy identifies it with a
unique vector \(c\) satisfying \(x\cdot x=c\cdot x\); this is the
\emph{characteristic vector}.  The space is alternating when \(c=0\), and
nonalternating otherwise.

For linearly independent vectors \(r_1,\ldots,r_s\), a \emph{dual family}
is a family \(r_1^*,\ldots,r_s^*\) with
\(r_i\cdot r_j^*=\one[i=j]\).  Pairing the radical means adjoining such a
family, chosen orthogonal to a previously fixed nondegenerate complement.
Indeed, when the ambient \(W\) is nondegenerate and
\(R=\Span_{\Ftwo}\{r_1,\ldots,r_s\}\), the pairing map
\[
 W\longrightarrow R^\vee,
 \qquad w\longmapsto(r\mapsto r\cdot w),
\]
is surjective: otherwise a nonzero element of \(R\) would pair to zero with
all of \(W\).  Choosing preimages gives a dual family.  If
\(U=P\perp R\) with \(P\) nondegenerate and
\(R=\operatorname{rad}_W(U)\), the same argument inside \(P^{\perp_W}\)
makes the family orthogonal to \(P\).  The new vectors are independent
modulo \(U\), as pairing a linear combination with each \(r_i\) shows.
On \(R\) and its dual family the Gram matrix is
\(\left(\begin{smallmatrix}0&I\\I&A\end{smallmatrix}\right)\), which is
nonsingular.  Adjoining the family therefore gives a nondegenerate extension
of dimension \(\dim U+\dim R\).
We use the following facts about symmetric binary spaces:
nondegenerate subspaces split off orthogonally; every
even-dimensional nondegenerate symmetric binary space has a self-dual
subspace; every nondegenerate nonalternating \(m\)-space is isometric to
the standard dot-product space \(\Ftwo^m\); and every odd-dimensional
nondegenerate symmetric binary space is nonalternating and has an
anisotropic vector.  Here are short proofs.  For a nondegenerate subspace
\(P\), its pairing map identifies \(P\) with \(P^\vee\), so each ambient
vector differs from a unique vector of \(P\) by an element of \(P^\perp\).
This gives the orthogonal splitting.  A nonzero nondegenerate alternating
space contains vectors \(u,v\) with \(u\cdot v=1\); their Gram matrix is
\(H=\left(\begin{smallmatrix}0&1\\1&0\end{smallmatrix}\right)\).
Splitting off these planes shows that an alternating space has even
dimension, and one vector from each plane spans a self-dual subspace.
In a nonalternating space, split off norm-one vectors until the remainder
is alternating.  A norm-one vector \(e\) orthogonal to such a plane can
absorb it: \(e+u,e+v,e+u+v\) are an orthogonal norm-one basis of their
three-dimensional span.  Repeating this gives the standard dot-product
basis.  In even dimension \(2k\), the vectors
\(e_1+e_2,\ldots,e_{2k-1}+e_{2k}\) span a self-dual subspace of that
space (the empty span applies in dimension zero).  Finally, odd dimension
excludes the alternating case and hence supplies a norm-one vector.

\subsection{Bell signatures and the two affine cores}
\label{subsec:pauli-notation}

The four Bell tensors form a basis for binary contractions.  Their action
on the two signatures below supplies the local data reconstructed in
\cref{sec:v4-deck}; the dictionary also fixes the signs inherited from
Shao--Cai's conventions.

For \(d,\ell,u,v\in\Ftwo\), define the ordered Bell signature
\begin{equation*}
  B_{d,\ell}(u,v)
  =
  \one[u\oplus v=d](-1)^{\ell u}.
\end{equation*}
Write \(B_\gamma=B_{d,\ell}\) for \(\gamma=(d,\ell)\).  Port order matters
for \(B_{1,1}\), which is negated under reversal.

With truth-table order \(00,01,10,11\), these four signatures are
\begin{equation}
 \begin{aligned}
  B_{0,0}&=(1,0,0,1)=I,&
  B_{0,1}&=(1,0,0,-1)=Z,\\
  B_{1,0}&=(0,1,1,0)=X,&
  B_{1,1}&=(0,1,-1,0)=Y.
 \end{aligned}
 \label{eq:bell-matrix-dictionary}
\end{equation}

\paragraph{Matrix and marked-pair dictionary.}
Combining \cref{eq:quaternion-basis,eq:bell-matrix-dictionary} gives the
actual identities
\begin{equation}
 I=Q_0=B_{0,0},\qquad
 iX=Q_1=iB_{1,0},\qquad
 Y=Q_2=B_{1,1},\qquad
 iZ=Q_3=iB_{0,1}.
 \label{eq:global-matrix-aliases}
\end{equation}
Consequently the marked projective notation used throughout is
\begin{equation}
 \begin{aligned}
  1&=[I]=[Q_0]=[B_{0,0}],&
  x&=[X]=[Q_1]=[B_{1,0}],\\
  j&=[Z]=[Q_3]=[B_{0,1}],&
  xj&=[Q_2]=[B_{1,1}].
 \end{aligned}
 \label{eq:global-mark-aliases}
\end{equation}
Here the last identity is projective: \(XZ=-Y\), whereas \(ZX=Y\).  For
the transfer convention \(T_B=BX\), the four Bell binaries give
\begin{equation*}
\begin{array}{c|c|c|c}
 B&\text{Bell name}&[B]&[T_B]=[BX]\\ \hline
 I&B_{0,0}&1&x\\
 X&B_{1,0}&x&1\\
 Z&B_{0,1}&j&xj\\
 Y&B_{1,1}&xj&j
\end{array}
\end{equation*}
The subscripted \(B_{d,\ell}\) is a fixed ordered Bell signature, whereas
\(B_g\) later denotes a chosen actual binary representative of a general
transfer class \(g\).  Reversal fixes \(I,X,Z\) and sends
\(Y=B_{1,1}\) to \(-Y\), without changing its projective class.
Throughout, a \emph{Pauli binary} means a nonzero scalar multiple of one of
these four ordered Bell tensors \(B_\gamma\), with its displayed input order
retained; the adjective \emph{actual} singles out a chosen representative
rather than only its projective class.

\begin{remark}[The Bell-signature dictionary of Shao--Cai]
The tuple in \cref{eq:bell-matrix-dictionary} is Shao--Cai's
\((=^+_2,=^-_2,\neq^+_2,\neq^-_2)\) \cite[Section~2]{ShaoCai2020};
\(Y^{\mathsf T}=-Y\).  Their holographic matrix named \(Z\) is our \(K\),
whereas our \(Z\) is only the diagonal sign matrix.
\end{remark}

Use the lexicographic indexing of the eight ports by \(\Ftwo^3\) fixed in
\cref{subsubsec:boolean-affine-tools}.
Define the eight-ary Reed--Muller indicator signature
\(\RMcore:\Ftwo^8\to\{0,1\}\) by
\begin{equation}
  \RMcore(x):=\one[x\in RM(1,3)].
  \label{eq:RM-core}
\end{equation}

Define the matrix
\[
  \Theta=
  \begin{pmatrix}0&1\\1&1\end{pmatrix}
  \in\GL_2(\Ftwo),
  \qquad \Theta^3=I.
\]
In a tensor \(A\otimes B\otimes C\) of three ordered binary signatures,
\(A\), \(B\), and \(C\) act on port pairs \((1,2)\), \((3,4)\), and
\((5,6)\), respectively.  With this convention, the six-port Pauli core is
\begin{equation}
  \begin{aligned}
    \Hcore
    &=\sum_{\gamma\in\Ftwo^2}
      B_\gamma\otimes B_{\Theta\gamma}\otimes B_{\Theta^2\gamma}\\
    &=I\otimes I\otimes I
      +Z\otimes Y\otimes X
      +X\otimes Z\otimes Y
      +Y\otimes X\otimes Z.
  \end{aligned}
  \label{eq:H6}
\end{equation}

For an ordered binary tensor \(A\) and an ordered pair of ports \((p,q)\),
the notation \(A^{(p,q)}\) places the first and second inputs of \(A\) on
ports \(p\) and \(q\), respectively.  For an even finite port set \(R\), a
\emph{Bell product} on \(R\) is a nonzero scalar multiple of
\[
 \bigotimes_{e=(p_e,q_e)\in M}B_{\gamma_e}^{(p_e,q_e)},
 \qquad \gamma_e\in\Ftwo^2,
\]
where \(M\) is a perfect matching of \(R\) and each pair of \(M\) has been
given the displayed orientation.  Reversing a pair carrying \(B_{1,1}\)
only changes the overall scalar sign, so this class is independent of the
chosen orientations.  For \(R=\varnothing\), a Bell product is a nonzero
nullary signature.  A Bell product is a \emph{same-\(\gamma\) Bell product}
when every factor label \(\gamma_e\) equals one fixed \(\gamma\).

An even-arity tensor-prime signature \(f\) has the \emph{Bell property}
when, for every ordered pair of distinct ports \((i,j)\) and every
\(\gamma\in\Ftwo^2\), the contraction
\(\partial_{(i,j)}^{B_\gamma}f\) is a Bell product on the remaining ports.
It has the \emph{strong Bell property} when each such contraction is, more
specifically, a same-\(\gamma\) Bell product.  These are algebraic contraction
properties and do not by themselves assert actual gadget provenance.  This
is the ordered-port, projective form of the standard definitions
\cite[full version, Definition~6.12]{ShaoCai2020} and
\cite[full version, Definition~1.10; conference version, Definition~12]%
{CaiFuShao2020Entanglement}.

\begin{remark}[The Shao--Cai \(f_8\) and \(\widehat f_6\) cores]
\label{rem:sc-f8-core}
\label{rem:sc-f6-core}
With the preceding port order, \(\RMcore\) is
literally the signature \(f_8=\widehat f_8\) of
\cite[full version, Section~8]{ShaoCai2020}; it has the strong Bell property:
contracting any ordered port pair by \(B_\gamma\) gives, up to a nonzero
scalar, a product of three copies of \(B_\gamma\) on the remaining ports
\cite[Theorems~13--14]{CaiFuShao2020Entanglement}.
Let \(\widehat f_6\) denote the six-ary signature of Shao--Cai given in
\cite[full version, Eq.~(6.1)]{ShaoCai2020}.  Direct comparison with
\cref{eq:H6} gives the literal port-permutation identity
\[
 \Hcore(x_1,\ldots,x_6)
 =
 \widehat f_6(x_1,x_3,x_5,x_4,x_6,x_2).
\]
The same results show that \(\Hcore\) has the Bell property: every ordered
Bell contraction is, up to a nonzero scalar, a product of two ordered Bell
binaries
\cite[Theorems~13--14]{CaiFuShao2020Entanglement}; see also
\cite[full version, Definition~6.12 and Lemma~6.13]{ShaoCai2020}.
We retain \(\RMcore,\Hcore\) as names because their code and Pauli geometry
drive the relative argument.  The inverse deck classifications, the actual
\(\Hcore\)-to-\(\RMcore\) circuit, and the signature-set-wide lifts are not
supplied by these forward Bell-property results and are proved here.
\end{remark}

The Bell tables satisfy the
exact bilinear orthogonality relation
\[
 \langle B_\gamma,B_{\gamma'}\rangle_F
 =2\,\one[\gamma=\gamma'].
\]

\subsection{Tractable signature sets and effective recognition}
\label{sec:predicate}

\subsubsection{The standard signature classes}
\label{subsec:standard-classes}

We include the zero signature in every class below.  The affine class
\(\cA\) consists of signatures
\begin{equation}
  f(x)=\lambda\,\one[Ax=b]\,i^{Q(x)},
  \label{eq:affine-class}
\end{equation}
where \(\lambda\in\AlgNums\setminus\{0\}\) and, for
\(r=\arity(f)\), there is an \(m\ge0\) with
\(A\in\Ftwo^{m\times r}\) and \(b\in\Ftwo^m\).  The function
\(Q:\{0,1\}^r\to\mathbb Z_4\), where
\(\mathbb Z_4=\mathbb Z/4\mathbb Z\), has a multilinear representative
\[
 Q(x)=q_0+\sum_{j=1}^r q_jx_j
      +\sum_{1\le j<\ell\le r}q_{j\ell}x_jx_\ell\pmod4,
 \qquad q_{j\ell}\in\{0,2\},
\]
so it has degree at most two and even mixed coefficients.
Equivalently, the phase is a product of fourth-root phases of affine linear
forms interpreted as \(\{0,1\}\)-valued integers.
Membership in \(\cA\) always refers to the displayed computational
coordinates.

For a nowhere-zero function \(f:\Ftwo^k\to\AlgNums\setminus\{0\}\), define
its multiplicative third-difference ratio
\(\mathcal T_f:(\Ftwo^k)^4\to\AlgNums^\times\) by
\begin{equation*}
 \mathcal T_f(z,u,v,w)=
 \frac{f(z+u+v+w)f(z+u)f(z+v)f(z+w)}
 {f(z+u+v)f(z+u+w)f(z+v+w)f(z)}.
\end{equation*}
Here \(z,u,v,w\in\Ftwo^k\), and repeated directions are permitted.
For fixed \(u,v\), let \(D_{u,v}\) be the multiplicative
second-difference operator acting on the entire function \(f\):
\begin{equation*}
 \bigl(D_{u,v}f\bigr)(z)
 :=
 \frac{f(z+u+v)f(z)}
      {f(z+u)f(z+v)}.
\end{equation*}

\begin{lemma}[Affine phase criterion]
\label{lem:affine-phase-criterion}
Let \(f:\Ftwo^k\to\AlgNums\setminus\{0\}\).  Then \(f\) belongs to
\(\cA\) if and only if
\(\mathcal T_f(z,u,v,w)=1\) for all \(z,u,v,w\in\Ftwo^k\).
\end{lemma}

\begin{proof}
Define the scalar \(\lambda=f(\zero^k)\), the scalars
\(A_j=f(e_j)/\lambda\) for \(j\in[k]\), and, for
\(1\le j<\ell\le k\), the scalars
\[
 B_{j\ell}=\frac{\lambda f(e_j+e_\ell)}{f(e_j)f(e_\ell)}.
\]
Substitution of \(z=\zero^k\) and \(u=v=w=e_j\) gives \(A_j^4=1\); taking
\(u=v=e_j\) and \(w=e_\ell\) gives \(B_{j\ell}^2=1\).  Moreover, the identity
\(\mathcal T_f(z,u,v,w)=1\) says exactly that
\[
 \bigl(D_{u,v}f\bigr)(z+w)=\bigl(D_{u,v}f\bigr)(z).
\]
Because \(w\) is arbitrary, \(D_{u,v}f\) is constant on \(\Ftwo^k\).
In particular,
\[
 \bigl(D_{e_j,e_\ell}f\bigr)(z)=
 \bigl(D_{e_j,e_\ell}f\bigr)(\zero^k)=B_{j\ell}
 \qquad(1\le j<\ell\le k).
\]

We now give the Hamming-weight induction explicitly.  Extend the notation
symmetrically by \(B_{\ell j}=B_{j\ell}\).  If \(z_j=0\), start with
\(f(e_j)/f(\zero^k)=A_j\) and add the basis vectors in \(\supp(z)\) one at a
time.  The identity \(D_{e_j,e_\ell}f=B_{j\ell}\), rearranged after the
previously added vector \(y\), is
\[
 \frac{f(y+e_\ell+e_j)}{f(y+e_\ell)}
 =
 B_{j\ell}\frac{f(y+e_j)}{f(y)}.
\]
It follows by induction on \(\wt(z)\) that
\begin{equation}
 \frac{f(z+e_j)}{f(z)}
 =A_j\prod_{\ell\ne j}B_{j\ell}^{z_\ell}
 \qquad(z_j=0).
 \label{eq:affine-phase-ratio}
\end{equation}
We next prove the desired formula for \(f(z)\), again by induction on
\(\wt(z)\).  It is immediate for \(z=\zero^k\).  For
\(z\ne\zero^k\), choose
\(j\in\supp(z)\) and define the predecessor word \(y:=z+e_j\), so \(y_j=0\) and
\(\wt(y)=\wt(z)-1\).  Applying the induction hypothesis to \(f(y)\) and
\cref{eq:affine-phase-ratio} to \(f(y+e_j)/f(y)\) yields
\[
 f(z)=\lambda\prod_jA_j^{z_j}\prod_{j<\ell}B_{j\ell}^{z_jz_\ell},
 \qquad A_j\in\mu_4,\quad B_{j\ell}\in\{1,-1\},
\]
because the new factors are precisely \(A_j\) and
\(B_{j\ell}\) for \(\ell\in\supp(y)\).  This is the full-support normal form in
\cref{eq:affine-class}.  Conversely, each linear factor and each
even-quadratic factor cancels in the multiplicative third difference, so
\(\mathcal T_f=1\).
\end{proof}

\begin{lemma}[Two-translate closure of the affine class]
\label{lem:affine-two-translate-closure}
If \(h\in\cA\) and \(w\) is a Boolean translation of its input
coordinates, then \(h+T_wh\in\cA\) (with the zero signature allowed), where
the translation operator \(T_w\) is defined by \((T_wh)(x):=h(x+w)\).
\end{lemma}

\begin{proof}
The zero case is immediate.  Let \(n\) denote the arity of \(h\), so
\(n:=\arity(h)\).  For \(h\ne0\), write
\[
 h(x)=c\,\one[x\in a+L]i^{Q(x)},
 \qquad c\in\AlgNums^\times,\quad a\in\Ftwo^n,\quad
 L\le\Ftwo^n,
\]
where \(Q:\Ftwo^n\to\mathbb Z_4\) is a quadratic phase with even mixed
coefficients.  If \(w\notin L\), every word in the union of the two supports
has a unique form \(a+y+tw\), with \(y\in L\) and \(t\in\Ftwo\).  On this
affine space,
\[
 (h+T_wh)(a+y+tw)=c\,i^{Q(a+y)}.
\]
The inverse coordinates \((y,t)\) are affine, so the displayed exponent
is a quadratic phase with even mixed coefficients.  Thus the sum is affine.
For \(w\in L\), expand
\(Q(x+w)-Q(x)=c_w+2\ell_w(x)\pmod4\), where
\(c_w\in\mathbb Z_4\) is constant and
\(\ell_w:\Ftwo^n\to\Ftwo\) is linear.  This gives
\[
 h(x)+h(x+w)=c\,i^{Q(x)}\bigl(1+i^{c_w}(-1)^{\ell_w(x)}\bigr)
 \quad(x\in a+L).
\]
For even \(c_w\) the bracket is a scalar times an affine hyperplane
indicator; for odd \(c_w\) it is a nonzero scalar times a fourth-root
linear phase.  Restriction of \(Q\) and adjoining a linear phase preserve
\(\cA\), proving the claim.  The displayed difference formula follows by
expanding the multilinear representative of \(Q\); all nonconstant
coefficients of that difference are even.
\end{proof}

Let \(L\) be a linear subspace of \(\Ftwo^{2k}\).  Its orthogonal complement
under the standard dot product
\(x\mathbin{\cdot}y:=\sum_{j=1}^{2k}x_jy_j\in\Ftwo\) is
\[
 L^\perp
 =
 \{y\in\Ftwo^{2k}:x\mathbin{\cdot}y=0
     \text{ for every }x\in L\}.
\]
Call \(L\) \emph{Lagrangian} when \(L=L^\perp\), and an arity-\(2k\)
nonzero signature \emph{flat-Lagrangian} when it has the form
\begin{equation*}
  f(x)=c\,\one[x\in a+L],\qquad
  0\ne c\in\AlgNums,\quad a\in\Ftwo^{2k},\quad L=L^\perp.
\end{equation*}
Let \(\FLag\) denote the zero signature together with all flat-Lagrangian
signatures of every even arity.
We call an arity-\(2k\) signature \emph{character-flat} when a linear sign
character is also allowed:
\begin{equation}
 \mathsf{CFlat}_{2k}
 =\left\{
 c(-1)^{\ell(x)}\one[x\in a+L]:
 c\in\AlgNums^\times,\ a\in\Ftwo^{2k},\ \ell\in(\Ftwo^{2k})^\vee,
 \ L=L^\perp
 \right\}\cup\{0\}.
 \label{eq:character-flat-class}
\end{equation}
Define the signature class
\(\mathsf{CFlat}=\bigcup_{k\ge0}\mathsf{CFlat}_{2k}\).  This is a
subclass of \(\cA\); it is the closure needed when a Reed--Muller core is
used together with all four Pauli signs rather than only \(I/X\).

The product class \(\cP\) consists of zero and all signatures of the form
\(\lambda\prod_jh_j\), where \(0\ne\lambda\in\AlgNums\) and every factor is
either a nonzero unary signature \(u(x_i)\), a binary equality
\(I(x_i,x_j)\), or a binary disequality \(X(x_i,x_j)\) on designated
signature variables.  Empty products and repeated variables are allowed;
an inconsistent product denotes the zero signature.

Define the primitive eighth root and the phase matrix by
\[
 \zeta_8:=e^{\pi i/4}=\frac{1+i}{\sqrt2},
 \qquad T_2:=\diag(1,\zeta_8).
\]
The local-affine class \(\cL\) consists of the signatures \(f\) such that
\begin{equation*}
 \left(\bigotimes_{j=1}^{\arity(f)}T_2^{\alpha_j}\right)f\in\cA
 \qquad\text{for every }\alpha\in\supp(f).
\end{equation*}

Let \(\cT\) be the set of all unary and binary signatures, and define the
signature class \(\cM\) by
\begin{equation*}
  \cM=\{f:f(x)=0\text{ whenever }\wt(x)>1\}.
\end{equation*}

\begin{definition}[Common transformability and \(K\)-presentation]
\label{def:common-transformability}
\label{def:common-K-presentation}
Let \(\cC\) be a signature class closed under nonzero scalars and port permutations.  A
finite signature set \(\cF\) is \(\cC\)-transformable if one
\(M\in\GL_2(\AlgNums)\) satisfies
\begin{equation*}
 M^{\mathsf T}M\in\cC,\qquad M^{-1}\cF\subseteq\cC.
\end{equation*}
Equivalently, for \(\widehat{\cF}=K^{-1}\cF\), one matrix
\(U=K^{-1}M\) gives the common \(K\)-presentation
\begin{equation*}
 U^{\mathsf T}XU\in\cC,\qquad U^{-1}\widehat{\cF}\subseteq\cC.
\end{equation*}
We abbreviate this signature-set-wide condition by
\begin{equation}
 \PresK(\cF;\cC)
 \quad\Longleftrightarrow\quad
 \exists U\in\GL_2(\AlgNums):
 \quad U^{\mathsf T}XU\in\cC,
 \qquad U^{-1}(K^{-1}\cF)\subseteq\cC.
 \label{eq:common-K-presentation-predicate}
\end{equation}
If \(\Gamma\) is already in \(K\)-coordinates, its common
\(\cC\) \(K\)-presentation is precisely
\[
\begin{aligned}
 \PresK(K\Gamma;\cC)\quad\Longleftrightarrow{}&
 \exists U\in\GL_2(\AlgNums):\\
 &U^{\mathsf T}XU\in\cC,\qquad U^{-1}\Gamma\subseteq\cC.
\end{aligned}
\]
These conditions include the edge and require one basis
for the entire signature set, not one basis per signature.  We use this definition
for \(\cC\in\{\cA,\cP,\cL\}\), and for the flat-Lagrangian and
\(\mathsf{CFlat}\) subclasses.
\end{definition}

A common flat-Lagrangian \(K\)-presentation requires the transformed edge
and every nonzero transformed signature to be flat-Lagrangian; it strengthens,
rather than enlarges, the affine family.
A common character-flat presentation is defined analogously and likewise
implies one common affine presentation.

\subsubsection{The Eulerian-orientation predicate}
\label{subsec:eo-predicate}

For an arity-\(r\) signature \(g\), define its central restriction
\begin{equation*}
  g|_{\EO}(x)=g(x)\one[q_r(x)=0].
\end{equation*}
It is zero for odd \(r\), and zero restrictions are deleted.  Define the
resulting central-restriction signature set by
\begin{equation*}
  \cG|_{\EO}=\{g|_{\EO}:g\in\cG,\ g|_{\EO}\not\equiv0\}.
\end{equation*}

Let \(g\ne0\) be a signature of arity \(k\) supported on the single Hamming
layer of weight \(d\).  Its \emph{EO completion} is the ordered tensor
\begin{equation}
  g_{\to\EO}
  =
  \begin{cases}
    g\otimes\DeltaZero^{\otimes(2d-k)},
      & 2d\ge k,\\
    g\otimes\DeltaOne^{\otimes(k-2d)},
      & 2d<k.
  \end{cases}
  \label{eq:eo-completion}
\end{equation}
The original ports precede the added unaries; a zeroth power is the nullary
unit.  The completed signature has even arity and central support, hence is
EO\@.  For a signature set of nonzero single-layer signatures, define its
EO-completion signature set by
\begin{equation*}
  \cG_{\to\EO}=\{g_{\to\EO}:g\in\cG\}.
\end{equation*}

For an arity-\(2d\) EO signature \(g\), an oriented perfect pairing
\(\pi=((a_1,b_1),\ldots,(a_d,b_d))\) has unordered pairs partitioning
\([2d]\).  For \(z\in\Ftwo^d\), let \(y\in\Ftwo^{2d}\) be the unique word
with \(y_{a_j}=z_j\) and \(y_{b_j}=1-z_j\); the pairing induces
\begin{equation*}
  g^\pi(z_1,\ldots,z_d)=g(y).
\end{equation*}

\begin{definition}[The predicate \(\EOTract\)]
\label{def:eo-tract}
For a finite EO signature set \(\cG\), \(\EOTract(\cG)\) holds if there exist a
single sign \(\varepsilon\in\{+1,-1\}\) and a single class
\(\cC\in\{\cA,\cP\}\), both shared by the entire signature set, such that:
\begin{enumerate}
\item for every \(g\in\cG\) and all
      \(\alpha,\beta,\gamma\in\supp(g)\), if
      \(\delta=\alpha\oplus\beta\oplus\gamma\), then
      \(\delta\in\supp(g)\) or
      \(\varepsilon q_{\arity(g)}(\delta)>0\);
\item for every \(g\in\cG\) and every oriented perfect pairing \(\pi\)
      of its ports, \(g^\pi\in\cC\).
\end{enumerate}
We adopt the vacuous convention \(\EOTract(\varnothing)=\mathrm{true}\).
\end{definition}

This is the compressed form of the opposite-pair restriction used in
\cite{MengWangXiaZheng2025,GuanShaoShi2026}: substitute
\(y_{b_j}=1-y_{a_j}\).  To make the interface literal, define the full
zero-extended restriction by
\[
 R_\pi g(x)=g(x)\prod_{j=1}^{d}X(x_{a_j},x_{b_j}).
\]
After permuting ports and writing the two incidences of each pair as
\((z_j,w_j)\), this becomes
\[
 R_\pi g(z,w)=g^\pi(z)\prod_{j=1}^{d}X(z_j,w_j).
\]
Thus \(g^\pi\in\cA\) (respectively \(\cP\)) implies
\(R_\pi g\in\cA\) (respectively \(\cP\)) by adjoining binary affine
(respectively product) factors, while contracting each \(w_j\) with the
unary vector \([1,1]\) recovers \(g^\pi\).  This is a closure statement
about the normal forms, not an assertion that \([1,1]\) is a free input
gadget.  In the native half-edge model the factors in \(R_\pi g\) are
realized by the structural \(X\)-edges, so the converse lift has exactly
the source's value-preserving meaning.  Conversely, \(\cA\) is closed under
this affine-variable summation (Gaussian elimination).  The product class
has the analogous elementary star elimination: after multiplying all unary
factors at an eliminated variable, choose one incident \(I/X\) edge as a
representative, replace every other incident edge by the corresponding
\(I/X\) edge between its neighbor and the representative, and absorb the
remaining unary weight into that representative (an inconsistent component
gives zero; an isolated variable gives a scalar).  In the displayed bridge
each \(w_j\) is a leaf, and explicitly
\[
  \sum_{w_j\in\{0,1\}}X(z_j,w_j)[1,1](w_j)=1.
\]
Thus the contraction simply deletes that factor and preserves the product
form (up to the displayed scalar).  Hence \(R_\pi g\in\cC\) if and only if
\(g^\pi\in\cC\), for \(\cC\in\{\cA,\cP\}\).

The source papers use unoriented pairs.  Reversing one pair applies the
unary coordinate permutation \(z_j\mapsto1-z_j\), and reordering pairs is
a port permutation; both preserve \(\cA\) and \(\cP\).  Since the signature
set is fixed, enumerating orientations and pairings is constant-size.
Consequently the oriented predicate above and the source predicate are
equivalent, with the same exact polynomial-time tests.

\subsubsection{The tractability predicates}
\label{subsec:tract-predicate}

The tractability predicate used here combines the seven alternatives
of Meng, Wang, Xia, and Zheng with the polynomial-time algorithms
of Guan, Shao, and Shi for the two support-defined alternatives
\cite{MengWangXiaZheng2025,GuanShaoShi2026}.  We package this boundary as
\(\Tract\).  It is the general tractability criterion underlying
\cref{thm:p1-deq2-dichotomy,thm:p1-eq2-kholant-dichotomy}: after the stated
preprocessing, its truth specifies the polynomial-time side of the main
dichotomies, and its failure specifies the \(\SharpP\)-hard side.  We retain
the full seven-clause predicate because it is also applied to intermediate
signature sets without a distinguished binary signature.

When either predicate below is applied to an arbitrary finite input set, its
argument is first reduced to the nonzero positive-arity part according to
\cref{subsec:zero-nullary}.  The definitions may therefore be stated for
preprocessed signature sets without carrying a separate superscript in the
public criteria.

For a preprocessed signature set \(\cF\), define its \(K\)-coordinate image
\(\widehat{\cF}=K^{-1}\cF\).  Define \(\Tract(\cF)\) to be the disjunction
of the following seven alternatives.

\begin{enumerate}[label=\textnormal{(\arabic*)}]
\item \label{case:one-sided}
      The signature set \(\widehat{\cF}\) is common one-sided, in the explicit
      sense that
      \[
       [\,\forall\widehat f\in\widehat{\cF},\
       \supp(\widehat f)\subseteq HW_{\arity(\widehat f)}^{\ge}\,]
       \quad\text{or}\quad
       [\,\forall\widehat f\in\widehat{\cF},\
       \supp(\widehat f)\subseteq HW_{\arity(\widehat f)}^{\le}\,],
      \]
      and \(\EOTract(\widehat{\cF}|_{\EO})\) holds.

\item \label{case:single-weight}
      Every \(\widehat f\in\widehat{\cF}\) is supported on one Hamming
      weight \(d_f\).  Across the signature set, some support word has negative
      charge and some (possibly other) support word has positive charge, and
      \(\EOTract(\widehat{\cF}_{\to\EO})\) holds, with the completion
      operation defined in \cref{eq:eo-completion}.

\item \label{case:arity-two}
      \(\cF\subseteq\angles{\cT}\).

\item \label{case:matching}
      \(\cF\subseteq\angles{K\cM}\), or
      \(\cF\subseteq\angles{KX\cM}\).

\item \label{case:affine}
      \(\cF\) is \(\cA\)-transformable.

\item \label{case:product}
      \(\cF\) is \(\cP\)-transformable.

\item \label{case:local-affine}
      \(\cF\) is \(\cL\)-transformable.

\end{enumerate}

\begin{remark}
The matrix symbols in \cref{case:matching} act tensorwise at the arity of
each signature.  In particular,
\[
\begin{aligned}
 f\in\angles{K\cM}
 &\Longleftrightarrow
 \text{ every tensor-prime factor of }K^{-1}f\text{ lies in }\cM,\\
 f\in\angles{KX\cM}
 &\Longleftrightarrow
 \text{ every tensor-prime factor of }XK^{-1}f\text{ lies in }\cM.
\end{aligned}
\]
\end{remark}

For the binary-disequality problems studied here, however, the tractable
boundary collapses to four alternatives.  A literal \(X\) in the signature
set---or, in the coordinate image tested by the main theorems, the equivalent
anchor \(Z=K^{\otimes2}I\)---rules out the one-sided, single-weight, and
matching alternatives.  This motivates the following reduced predicate.

\begin{definition}[The reduced predicate \(\TractX\)]
\label{def:tract-X}
For a preprocessed signature set \(\cF\), define
\begin{equation}
\begin{split}
 \TractX(\cF)\quad\Longleftrightarrow\quad
 &\cF\subseteq\angles{\cT},
 \\[-0.2em]
 &\text{or \(\cF\) is \(\cA\)-transformable,}
 \\[-0.2em]
 &\text{or \(\cF\) is \(\cP\)-transformable,}
 \\[-0.2em]
 &\text{or \(\cF\) is \(\cL\)-transformable.}
\end{split}
\label{eq:tract-X}
\end{equation}
Thus \(\TractX\) is exactly the disjunction of
\cref{case:arity-two,case:affine,case:product,case:local-affine}.
\end{definition}

\begin{lemma}[The binary anchor removes the other three alternatives]
\label{lem:tract-X-anchor}
Let \(\cF\) be a preprocessed signature set.  If \(X\in\cF\), then
\[
 \Tract(\cF)\quad\Longleftrightarrow\quad\TractX(\cF).
\]
The same conclusion holds if \(Z=\diag(1,-1)\in\cF\).
\end{lemma}

\begin{proof}
If \(X\in\cF\), then
\[
 -iZ=K^{-1}X\in\widehat{\cF};
\]
if \(Z\in\cF\), then instead \(I=K^{-1}Z\in\widehat{\cF}\).
In either case \(\widehat{\cF}\) contains a nonsingular binary signature
supported precisely on \(\{00,11\}\), with both endpoint values nonzero.
The words \(00\) and \(11\) have opposite charges, so
\cref{case:one-sided} is impossible, and they have different Hamming
weights, so \cref{case:single-weight} is impossible.

The retained binary has matrix rank two and is therefore tensor-prime.  It
does not belong to \(\cM\), since it is nonzero at \(11\); applying
\(X^{\otimes2}\) only exchanges its two endpoint values, so the resulting
tensor is again tensor-prime and not in \(\cM\).  The two characterizations
in the preceding remark therefore exclude both orientations in
\cref{case:matching}.  Hence only
\cref{case:arity-two,case:affine,case:product,case:local-affine} can hold,
which proves the forward implication.  The reverse implication follows
immediately because these four alternatives are clauses of \(\Tract\).
\end{proof}

Both coordinate images tested in the main theorems contain
\(Z=K^{\otimes2}I\).  Thus the second part of \cref{lem:tract-X-anchor}
identifies their reduced \(\TractX\) tests with the full \(\Tract\) criterion.

\subsubsection{Effective algebraic geometry}
\label{subsubsec:effective-algebraic-geometry}

The transformability test in \cref{lem:algebraic-transform-witness} asks
for one matrix satisfying polynomial conditions for the edge and every
signature.  The tools below turn that existence question into ideal
membership and recover an exact algebraic matrix when one exists.

Let \(R=L[x_1,\ldots,x_n]\) be a polynomial ring over the stated exact
field.  An ideal \(I\lhd R\) is an additive subgroup with \(rf\in I\) for
all \(r\in R,f\in I\).  It is proper when \(I\ne R\), prime when it is
proper and \(fg\in I\) implies \(f\in I\) or \(g\in I\).  For
\(f_1,\ldots,f_s\in R\), their generated ideal is
\[
 (f_1,\ldots,f_s):=
 \left\{\sum_{j=1}^s r_jf_j:r_1,\ldots,r_s\in R\right\};
\]
an ideal is finitely generated when it equals one of this form.  We use
\emph{algebraically closed} in the sense of
\cref{subsubsec:set-field-exact}.  For an algebraically
closed \(\Omega\supseteq L\), the affine zero locus is
\[
 V_\Omega(I):=\{a\in\Omega^n:f(a)=0\text{ for every }f\in I\}.
\]
A projective zero locus is defined by homogeneous equations in the specified
projective space.  A Zariski-closed set is a zero locus, a Zariski-open set
is its complement, a locally closed set is an intersection of one of each,
and a constructible set is a finite union of locally closed sets.

A monomial order is a total order \(\prec\) in which every nonempty monomial
set has a least member, \(1\preceq m\) for every monomial \(m\), and
\(m\prec n\) implies \(mp\prec np\) for every monomial \(p\).  For
\(f\ne0\), \(\operatorname{LM}(f)\) is its largest monomial, and
\(\operatorname{LM}(I)\) is the ideal generated by all
\(\operatorname{LM}(f)\) with \(0\ne f\in I\).
A finite \(\mathcal G\subseteq I\) is a Gr\"obner basis when its leading
monomials generate \(\operatorname{LM}(I)\).  It is reduced when every
element is monic---its leading coefficient is \(1\)---and no monomial of one
element is divisible by the leading monomial of another.  A \emph{monomial
ideal} is an ideal generated by monomials.  For a fixed order, the reduced basis is unique;
division by a Gr\"obner basis decides ideal membership, and Buchberger's
algorithm constructs one.  It terminates by the ascending-chain condition:
there is no infinite strictly increasing chain of monomial ideals
\cite[Chapters~2--3]{CoxLittleOShea2015}.

For \(\Delta\in R\), abbreviate
\(V=V_{\overline L}\) and \(V(\Delta):=V((\Delta))\).  The weak
Nullstellensatz and the geometric effect of saturation are
\[
 V_{\overline L}(I)=\varnothing\Longleftrightarrow1\in I,
 \qquad
 I:\Delta^\infty:=
 \{f:\Delta^mf\in I\text{ for some }m\ge0\},
\]
\[
 V(I:\Delta^\infty)
 =\overline{\,V(I)\setminus V(\Delta)\,}^{\,\mathrm{Zar}},
\]
where Zariski closure means the smallest Zariski-closed set containing the
displayed set,
so saturation retains precisely the irreducible components not contained in
\(V(\Delta)\) when \(I\) is radical
\cite[Chapters~4 and~8]{CoxLittleOShea2015}.  Here an ideal is
\emph{radical} when \(f^m\in I\) implies \(f\in I\), and a prime
decomposition expresses a radical ideal as a finite intersection of prime
ideals.  This ideal-theoretic radical is unrelated to the radical of a
bilinear form.

A nonempty closed set is irreducible when it is not the union of two proper
closed subsets, and an irreducible component is a maximal irreducible closed
subset.  For a nonempty algebraic set \(X\), its dimension \(\dim X\) is the
largest \(d\) for which there is a chain
\(Z_0\subsetneq\cdots\subsetneq Z_d\subseteq X\) of nonempty irreducible
closed subsets.  For an ordered variable split, the elimination
ideal is \(I\cap L[x_{k+1},\ldots,x_n]\).  In a triangular set, each
successive equation has a new largest variable, called its main variable,
and its initial is its leading coefficient in that variable.  An element
\(a\) is a non-zero-divisor modulo \(J\) when
\(ab\in J\Rightarrow b\in J\).  A regular chain is a triangular set in
which every successive initial is a non-zero-divisor modulo the ideal
obtained by saturating the preceding equations by their initials.  Its
constructible branch requires all initials to be nonzero.  A triangular or
regular-chain decomposition splits a solution set into finitely many such
constructible branches.  A \emph{generic algebraic specialization} assigns
algebraic values to the remaining free coordinates outside the finite union
of exceptional zero loci on which a required initial or denominator
vanishes.  Regular-chain decomposition, generic specialization, and minimal
polynomials with isolating data produce exact algebraic sample points
\cite{AubryLazardMorenoMaza1999,Cohen1993}.

For a field extension \(L'/L\), base change sends \(I\) to
\(I_{L'}:=I\,L'[x_1,\ldots,x_n]\) and extends the associated solution
spaces.  When \(L\subseteq\mathbb C\), \emph{complexification} means this
base change with \(L'=\mathbb C\).  In finite certificates, RREF means
reduced row-echelon form: every
nonzero row has a leading \(1\), pivot columns have no other nonzero entry,
pivots move strictly rightward down the rows, and zero rows are last.

\begin{lemma}[Effective algebraic transformability witnesses]
\label{lem:algebraic-transform-witness}
Let \(E\in\Mat_2(\AlgNums)\), let \(\cF\) be a finite algebraic
signature set, and let \(\cC\in\{\cA,\cP,\cL\}\).  Suppose there is a matrix
\(U\in\GL_2(\mathbb C)\) such that
\[
  U^{\mathsf T}EU\in\cC,
  \qquad
  U^{-1}\cF\subseteq\cC,
\]
where membership in \(\cC\) allows complex values for the free scalar and
unary parameters in its normal forms.  Then such a witness exists in
\(\GL_2(\AlgNums)\).  Moreover, for exact algebraic input, the existence of a
witness is decidable and an algebraic witness can be produced.
\end{lemma}

\begin{proof}
Let \(L\) be a number field containing the entries of \(E\) and \(\cF\),
as well as \(i\) and \(\zeta_8=(1+i)/\sqrt2\) (equivalently \(i\) and
\(\sqrt2\)).  There are only finitely many combinatorial normal-form
types compatible with the arities: affine supports and phase coefficients,
product partitions with equality or disequality signs and unary weights,
and, for \(\cL\), the finitely many prescribed
\(T_2\)-twisted affine conditions.  Enumerate one shared type choice for the transformed edge and
all transformed signatures.

For \(\cP\), repeated factors cause no infinitude: at fixed arity we
canonically enumerate the support set, tensor partition, factor
multiplicities, equality/disequality signs, unary phase type, and port
permutation.  On a fixed pair of variables, repeated copies of the same
\(I\)- or \(X\)-factor are idempotent and are retained as a single factor;
the simultaneous presence of an \(I\)- and an \(X\)-factor makes the whole
product zero.  A loop \(I(x,x)\) is the scalar one and a loop \(X(x,x)\)
is zero.  Unary factors on one variable are multiplied and combined into
one unary (or zero).  This gives a finite canonical enumeration without
using a mod--two rule for product factors.
For \(\cL\), enumerate the prospective transformed support set and affine phase data.  Vanishing outside that mask, nonvanishing on it,
and the corresponding twist conditions then become finitely many polynomial
equalities and inequations.  Thus every candidate, including one whose
transformed support is not known in advance, is one of finitely many
constructible systems.

Fix one such branch, and let \(\mathbf x\) be the tuple of all its scalar
variables: the four entries of \(U\), the scalar and unary parameters in the
chosen normal forms, and the auxiliary variables introduced below.  Introduce
an auxiliary scalar variable \(z\) and impose
\[
 z\det U=1.
\]
Here \(\operatorname{adj}(U)\) denotes the classical adjugate, characterized
by \(U\operatorname{adj}(U)=(\det U)I\); hence on this solution set
\(U^{-1}=z\operatorname{adj}(U)\).  After this substitution, equality of
every transformed tensor entry with its proposed normal form is a polynomial
equation over \(L\).  If the branch also has the nonvanishing polynomial
constraints \(p_1,\ldots,p_s\ne0\), introduce one more auxiliary scalar
variable \(t\) and replace all these inequations by the single equation
\[
 t\prod_{\ell=1}^s p_\ell=1.
\]
Thus the branch is exactly the zero set of a finitely generated ideal
\(I\subseteq L[\mathbf x]\).

Compute a Gr\"obner basis of \(I\) over the exact number field \(L\), using
the terminating ideal-membership algorithm described above.  Let
\(\overline L\subseteq\mathbb C\) be the algebraic closure of the embedded
number field \(L\) inside \(\mathbb C\).  Writing
\(V_{\overline L}(I)\) for the affine zero locus of \(I\) over
\(\overline L\), the standard ideal-membership and weak-Nullstellensatz
theorems give
\[
 V_{\overline L}(I)=\varnothing
 \quad\Longleftrightarrow\quad
 1\in I,
\]
and the latter is equivalent to the reduced Gr\"obner basis being
\(\{1\}\).  Hence emptiness of each branch is decidable.  Moreover, a
complex solution makes \(I\) proper, so the weak Nullstellensatz supplies
a point over \(\overline L=\AlgNums\); conversely every such point is a
complex solution.  In a zero-dimensional branch, standard elimination
followed by triangular decomposition produces exact minimal polynomials and
isolating data for one point.  A nonempty positive-dimensional branch is
also effective: enumerate algebraic tuples fairly, first by total exact
encoding height and then lexicographically, and test every displayed
equation exactly.  The weak Nullstellensatz guarantees that this search
eventually reaches a point of the branch.  Equivalently one may use a
regular-chain decomposition with a generic algebraic specialization of its
free coordinates.  Reading the \(U\)-coordinates gives the required
algebraic witness.  These are standard effective consequences of
Gr\"obner-basis theory \cite{Cohen1993,CoxLittleOShea2015}; no
encoding-size or running-time bound is asserted here.
\end{proof}

\subsubsection{Imported dichotomies and decomposition}
\label{subsec:imported-results}

Unless an ordinary \(\Holant\) symbol or an explicit source presentation
\(\Holant(X\mid\cdot)\) is displayed, the imported \(\KHolant\) statements
below use native \(K\)-coordinates (so the edge tensor is \(X\)); applying
the fixed \(K\)-change gives the corresponding ordinary statement.  All
signature sets in these statements are fixed and preprocessed.
In particular, every later use of \cref{thm:external-odd} on a native set
\(\Gamma\) means the typed bridge
\[
 \KHolant(\Gamma)\equivT\Holant(K\Gamma),
 \qquad
 K\Gamma=\{K^{\otimes\arity(g)}g:g\in\Gamma\},
\]
followed by the ordinary theorem on \(K\Gamma\) and the inverse bridge.
The known nonzero edge and gadget scalars are restored by the record.

\begin{theorem}[Weighted Eulerian orientations
{\cite{MengWangXia2025EO}; \cite[Theorems~1.1 and~2.17]{GuanShaoShi2026}}]
\label{thm:external-eo}
For each fixed finite preprocessed algebraic EO signature set \(\cG\) (zero
signatures and nullary scalar factors are deleted and their known scalar
contributions are restored separately), if \(\EOTract(\cG)\) holds,
then \(\KHolant(\cG)\in\FP\); otherwise
\(\KHolant(\cG)\) is \(\SharpP\)-hard.
\end{theorem}

\begin{theorem}[Odd-arity Holant
{\cite[full version, Theorem~33; conference version, Theorem~30]{MengWangXiaZheng2025}}]
\label{thm:external-odd}
Let \(\cF\) be a finite preprocessed algebraic signature set containing a
nonzero odd-arity signature.  Then \(\Holant(\cF)\in\FP\) if
\(\Tract(\cF)\) holds, and is \(\SharpP\)-hard otherwise.
\end{theorem}

Its first two alternatives were originally in \(\FP^{\mathsf{NP}}\); the
stated \(\FP\) conclusion follows from
\cite[Theorem~1.1 and Corollary~1.2]{GuanShaoShi2026}.

\begin{theorem}[Tensor-factor extraction
{\cite[full version, Theorem~32 and Remark~39; conference version,
Theorem~29]{MengWangXiaZheng2025}}]
\label{thm:external-factor}
Let \(\Lambda\) be a finite preprocessed algebraic signature set in native
\(K\)-coordinates, and let \(h=a\otimes b\) be a nonzero tensor product with
algebraic factors \(a,b\) of positive arity (equivalently, apply the source theorem
to the ordinary set \(K(\Lambda\cup\{h\})\)).  Any nullary scalar factor is
absorbed during preprocessing.  Unless
\(\Lambda\cup\{h\}\) is common one-sided,
\[
  \KHolant(\Lambda,a,b)\equivT\KHolant(\Lambda,h).
\]
In the common-one-sided exception, \(\KHolant(\Lambda,h)\) is in \(\FP\)
when \(\EOTract((\Lambda\cup\{h\})|_{\EO})\) holds and is
\(\SharpP\)-hard otherwise, by the one-sided clause of the cited theorem
and \cref{thm:external-eo}.
\end{theorem}

\begin{theorem}[Quaternary eight-vertex boundary
{\cite[Theorem~3.1]{CaiFu2023}}]
\label{thm:external-eight-vertex}
Let \(M\) be an algebraic quaternary signature satisfying
\[
 \supp(M)\subseteq\{x\in\{0,1\}^4:\wt(x)=2\}\cup\{0000,1111\},\qquad
 M(0000)M(1111)\ne0.
\]
Then \(\Holant(X\mid M)\) is \(\SharpP\)-hard unless there exist
\(\cC\in\{\cA,\cP,\cL\}\) and \(U\in\GL_2(\AlgNums)\) such that
\[
  U^{\mathsf T}XU\in\cC,\qquad U^{-1}M\in\cC.
\]
Every exception is in \(\FP\).  The source permits \(U\) over \(\mathbb C\);
\cref{lem:algebraic-transform-witness} makes it algebraic.
\end{theorem}

\begin{lemma}[One-hot matching evaluation]
\label{lem:one-hot-matching-evaluation}
Fix a finite algebraic signature set, and let \(\Gamma\) be a closed
factor-expanded finite native-\(X\) network over it (each
occurrence has first been replaced by its fixed tensor-prime factors, with
the induced port wiring and recorded scalar).  If a vertex signature is
zero, return zero; otherwise, for every vertex \(w\), let \(h_w\) denote its
nonzero signature, define its arity by \(d_w:=\arity(h_w)\),
and suppose
\(\supp(h_w)\subseteq\{\zero^{d_w},e_1,\ldots,e_{d_w}\}\).  Define its
zero-word weight \(a_w:=h_w(\zero^{d_w})\) and weight-one port coefficients
\(b_{w,p}:=h_w(e_p)\).  Then \(Z_X(\Gamma)\) is computable using
\(O(|V(\Gamma)|+|E_{\rm int}(\Gamma)|)\) arithmetic operations in the fixed
coefficient field, and hence in polynomial bit time; \(V(\Gamma)\) is the
vertex set.
\end{lemma}

\begin{proof}
For an \(X\)-edge exactly one of its two half-edges receives bit \(1\).
Consequently assignments are in bijection with choices of one endpoint of
each edge, subject to at most one chosen endpoint at each vertex.  For such a
feasible endpoint-choice assignment \(\omega\), write
\(\operatorname{occ}_\omega(w)\in\{0,1\}\) for the number of selected
endpoints at \(w\), and let \(s_w(\omega)\) be the selected port when it
exists.  The weight of \(\omega\) is
\[
  \prod_{\operatorname{occ}_\omega(w)=0}a_w
  \prod_{\operatorname{occ}_\omega(w)=1}b_{w,s_w(\omega)}.
\]
The product of the fixed nonzero scalar factors recorded during the
factor expansion is multiplied into the final component product (and a
zero expansion scalar was already handled by returning zero), so it does
not change the recurrence below.
For a connected component, let \(e\) and \(v\) denote its edge and vertex
counts.  A component with \(e>v\) has value zero.  Every surviving component
has \(e=v-1\) (a tree) or \(e=v\) (unicyclic, with loops and parallel edges
included).  On a rooted tree, let \(p_w\) be the parent-edge port at \(w\),
and let \(p_{wc}\) be the port at \(w\) of the child edge \(wc\).  If
\(F_w(1)\) (resp. \(F_w(0)\)) denotes the subtree sum when the parent edge
selects \(w\) (resp.\ its parent), then
\[
\begin{aligned}
 F_w(1)&=b_{w,p_w}\prod_cF_c(1),\\
 F_w(0)&=a_w\prod_cF_c(1)+
 \sum_c b_{w,p_{wc}}F_c(0)\prod_{d\ne c}F_d(1).
\end{aligned}
\]
At the root use the second line without a parent term.  Order the children as
\(c_1,\ldots,c_m\).  Products and sums are evaluated without division by
the prefix products
\(P_j:=\prod_{\ell<j}F_{c_\ell}(1)\) and suffix products
\(S_j:=\prod_{\ell>j}F_{c_\ell}(1)\), for which
\(P_jS_j=\prod_{\ell\ne j}F_{c_\ell}(1)\) even
when some factors vanish.  For a unicyclic component delete one cycle
edge \(e\), with endpoint ports \(p_u\) at \(u\) and \(p_v\) at \(v\).
Enumerate the two legal choices \(p\in\{p_u,p_v\}\), and let \(v(p)\)
be the vertex incident to the chosen port.  The ports are distinct even
when \(e\) is a loop.  For each choice, root the resulting tree at
\(v(p)\) and impose that \(p\) is already occupied through \(e\):
at that root use
\[
 b_{v(p),p}\prod_cF_c(1).
\]
This forces every remaining incident edge at that root to select its child.
Use the ordinary
tree recurrence on all proper subtrees and perform this same forced-root
calculation separately for each \(p\); sum the two values.  The same
boundary-condition description handles a loop (the two ports are distinct
incidences at one vertex) and parallel cycle edges.
The prefix/suffix evaluations use linearly many field operations.  Every
intermediate entry is a sum of at most exponentially many products of
\(O(|V|+|E_{\rm int}|)\) fixed coefficients.  In a fixed rational basis of
their number field, a common denominator and the numerator sizes therefore
have \(O(|V|+|E_{\rm int}|)\) bits, giving polynomial bit time.
\end{proof}

\subsubsection{Decidability and polynomial-time evaluation}
\label{subsec:predicate-algorithms}

\begin{theorem}[Recognition and evaluation]
\label{thm:predicate-algorithms}
For every finite algebraic signature set \(\cF\), the predicate \(\Tract(\cF)\),
with the preprocessing convention of \cref{subsec:zero-nullary}, is
decidable.  If it holds, then \(\Holant(\cF)\in\FP\).
\end{theorem}

\begin{proof}
The two EO alternatives are decided by finitely many exact tests of support,
layers, affine subspaces, and phases; enumerate all oriented pairings and the
four shared \((\varepsilon,\cC)\)-choices.  A signature's tensor
decomposition can be found by enumerating all nontrivial bipartitions
\(S\mid S^c\), where
\(\varnothing\ne S\subsetneq[\arity(f)]\) and
\(S^c=[\arity(f)]\setminus S\), and testing whether the corresponding
matrix flattening has rank one; recursively applying this test to the
factors decides \cref{case:arity-two,case:matching}.  Meanwhile,
\cref{lem:algebraic-transform-witness} with \(E=I\) decides
\cref{case:affine,case:product,case:local-affine} and returns the common
witness.

For evaluation, charge conservation reduces \cref{case:one-sided} to
\cref{thm:external-eo}: in a closed contributing assignment all local
charges have the common sign and sum to zero, so every local charge is zero
and only the central restrictions remain.  Thus the off-central entries may
be deleted in this evaluator (this is an identity on closed instances, not a
claim that the restrictions are free gadgets).  In \cref{case:single-weight}, an
occurrence of arity \(k_f\) and weight \(d_f\) has charge
\(c_f=2d_f-k_f\); each connected component \(C\) must satisfy
\begin{equation*}
  \sum_{v\in V(C)}c_{f_v}=0.
\end{equation*}
Failure gives zero.  Otherwise construct an auxiliary EO instance solely for
the evaluator: replace each occurrence \(f_v\) by its fixed completion
\((f_v)_{\to\EO}\), and in each component pair the newly appended
\(\DeltaZero\)-ports (from \(c_{f_v}>0\)) with the appended
\(\DeltaOne\)-ports (from \(c_{f_v}<0\)) by native \(X\)-edges.  The charge
balance makes the two numbers equal, and every original assignment has a
unique extension to this auxiliary instance (and conversely), so its
partition function is unchanged.  These pins are explicit entries of the
derived tensors in an algorithmic reduction; no pinning signature is assumed
to be a free gadget of the input problem.  The finite completed signature set
is EO, so \cref{thm:external-eo} supplies the polynomial-time evaluator.
After tensor-prime expansion, the arity-two branch is a graph of maximum
degree two and is evaluated by transfer matrices on paths and cycles.
For the matching alternatives, first apply the common \(K\)-coordinate
equivalence, the possible global \(X\)-toggle, and tensor-prime expansion;
the toggle is a legal native change because \(X^{\mathsf T}XX=X\).
After this change, every counted prime factor \(h\) satisfies
\(h(x)=0\) for \(\wt(x)>1\).  Only structural equality kernels, or degree-two equality
representations introduced solely to display the source interface, are
suppressed before the factor graph is built; an explicit arity-two signature
occurrence remains a factor vertex unless it has first been eliminated by the
path/cycle preprocessing of \cref{lem:p1-equality-wire-equivalence}, with its
known scalar recorded.  Thus
\cref{lem:one-hot-matching-evaluation} gives a polynomial-time evaluator
using linearly many field operations; the
expansion has constant size per occurrence because the signature set is
fixed.  Finally, a common witness \(M\) gives
\[
  \Holant((\Eq_2)M\mid M^{-1}\cF),
  \qquad(\Eq_2)M=M^{\mathsf T}M,
\]
with both sides in the same class
\(\cC\in\{\cA,\cP,\cL\}\).  Forgetting the
bipartition yields an ordinary \(\Holant(\cC)\) instance, computable in
\(\FP\): affine instances are quadratic Gauss sums, product instances split
into unary/equality/disequality components, and local-affine instances are a
restriction of the corrected complex-valued \(\Holant^c\) evaluator of
\cite{Backens2021,Backens2025} (the real-valued presentation
\cite[Corollary~3.1]{CaiLuXia2018RealHolantc} gives the same local routine).
\end{proof}

\subsection{Projective tensor and finite-incidence geometry}
\label{subsec:projective-incidence-conventions}

In \cref{sec:normalized-dihedral,sec:platonic-decks}, a matching product
is a simple tensor after its ports are grouped into pairs.  Segre geometry
therefore turns tensor factorization into rank equations.  Incidences among
the available binary kernels then constrain which pairings those equations
can describe.

\subsubsection{Projective tensor geometry}
\label{subsubsec:projective-tensor-geometry}

We use the projective-space conventions of
\cref{subsubsec:group-projective}.  Whenever a projective object occurs in a
literal tensor equation, a displayed or previously fixed nonzero
representative is understood; changing it changes only the separately
recorded scalar.  We set \(\mathbb P(0):=\varnothing\).

A nonzero tensor \(v_1\otimes\cdots\otimes v_s\) is \emph{simple} or
\emph{decomposable}.  For finite-dimensional spaces \(V_1,\ldots,V_s\), the
Segre map is
\[
 \mathbb P(V_1)\times\cdots\times\mathbb P(V_s)
 \longrightarrow \mathbb P(V_1\otimes\cdots\otimes V_s),
 \qquad ([v_1],\ldots,[v_s])\longmapsto[v_1\otimes\cdots\otimes v_s].
\]
Its image is the \emph{Segre variety}; its affine cone consists of zero and
the simple tensors.  The Segre ideal is
generated by the \(2\)-by-\(2\) minors of the bipartite flattenings.  A
\emph{one-factor Segre ruling} is a projective linear family in which all
factors but one are fixed, and every projective line contained in a Segre
variety belongs to such a ruling \cite[Chapter~2]{Landsberg2012}.

\subsubsection{Möbius geometry}
\label{subsubsec:mobius-geometry}

The Riemann sphere is
\(\widehat{\mathbb C}:=\mathbb C\cup\{\infty\}\cong\mathbb P(\mathbb C^2)\),
with affine chart \(z=[z:1]\) and \(\infty=[1:0]\).  A Möbius transformation
is a projective automorphism of this sphere; in the affine chart it is
\[
 z\longmapsto\frac{az+b}{cz+d},
 \qquad ad-bc\ne0,
\]
where \(cz+d=0\) is sent to \(\infty\) and \(\infty\) is interpreted
projectively.  A \emph{generalized circle} is a Euclidean circle
\(\{z:|z-z_0|=r\}\), \(r>0\), or a real affine line together with
\(\infty\); such a line has the form
\(\{z:\operatorname{Re}(\overline\alpha z)=\beta\}\) for
\(\alpha\in\mathbb C^\times\) and \(\beta\in\mathbb R\).
Möbius transformations carry generalized circles to
generalized circles \cite[Chapter~3]{Beardon1983}.  The unit circle is
\(S^1:=\{z\in\mathbb C:|z|=1\}\); its counterclockwise cyclic order is the
order of arguments modulo \(2\pi\).

\subsubsection{Root configurations and incidence}
\label{subsubsec:root-incidence}

Let \(E\) be a real vector space with a positive-definite symmetric bilinear
form \(\langle\ ,\ \rangle\); this is a real inner-product space.  A
\emph{reduced finite root system} in \(E\) is a finite spanning set \(\Phi\)
of nonzero vectors such that
each reflection
\[
 v\longmapsto
 v-2\frac{\langle v,\alpha\rangle}{\langle\alpha,\alpha\rangle}\alpha
 \qquad(\alpha\in\Phi)
\]
permutes \(\Phi\), and
\(\Phi\cap\mathbb R\alpha=\{\alpha,-\alpha\}\).  Its root-line
configuration is
\[
 \mathcal R(\Phi):=\{[\alpha]:\alpha\in\Phi\}
 \subseteq\mathbb P(E\otimes_{\mathbb R}\mathbb C).
\]
In root-system notation, \(\Phi(D_4),\Phi(F_4),\Phi(H_4)\) denote the
standard systems of those types, and
\[
 \mathcal R(D_4):=\mathcal R(\Phi(D_4)),\quad
 \mathcal R(F_4):=\mathcal R(\Phi(F_4)),\quad
 \mathcal R(H_4):=\mathcal R(\Phi(H_4)).
\]
These are root-system types \cite{Humphreys1990}; in particular, the type
\(D_4\) differs from the order-four dihedral group \(D_4\cong V_4\).

A finite projective point configuration is a finite subset \(X\) of a
projective space.  Points are collinear when they lie on one projective
line.  A secant line contains at least two points of \(X\), a \(k\)-rich
line contains at least \(k\), and a plane section is \(X\cap\Pi\) for a
projective plane \(\Pi\).  A point--line incidence structure is a pair
\((P,\mathcal L)\), where \(\mathcal L\) is a family of subsets of \(P\);
an automorphism is a permutation of \(P\) preserving \(\mathcal L\).  More
generally, a finite incidence hypergraph has the named points as vertices
and the named incident subsets as hyperedges, and its automorphisms preserve
every distinguished hyperedge family.

For four distinct points \(p_i=[v_i]\) on a projective line
\(\ell=\mathbb P(U)\), choose an ordered basis of the two-dimensional vector
space \(U\).  Let \(\det_U(v,w)\) be the \(2\)-by-\(2\) determinant of the
coordinate columns of \(v,w\) in that basis, and define
\[
 [p_1,p_2;p_3,p_4]
 :=
 \frac{\det_U(v_1,v_3)\det_U(v_2,v_4)}
      {\det_U(v_1,v_4)\det_U(v_2,v_3)}.
\]
Changing a representative scales one numerator and one denominator factor,
and changing the ordered basis scales every determinant by the same nonzero
factor.  Hence the scalar is well defined.  The ordered quadruple is
\emph{harmonic} when this value
is \(-1\).  A harmonic projectivity between two four-point lines is a
projective isomorphism carrying one set to the other and preserving exactly
the ordered quadruples of cross-ratio \(-1\).  This projective cross-ratio
is distinct from any locally defined table cross-ratio.

\subsection{Laurent gauges and product weights}
\label{subsec:laurent-gauge-conventions}

The diagonal parameters in \cref{sec:normalized-dihedral} attach
independent weights to ports.  Laurent monomials distinguish their
contributions, so coefficient identities can recover restrictions on the
original tensor without choosing a new basis for the Holant instance.

\subsubsection{Laurent polynomials}
\label{subsubsec:laurent-polynomials}

For a commutative ring \(R\), a Laurent polynomial over \(R\) is a finite sum
\[
 P(\lambda)=\sum_{u\in\mathbb Z^r}c_u\lambda^u,
 \qquad
 \lambda^u:=\prod_{j=1}^r\lambda_j^{u_j},
\]
with \(c_u\in R\) and all but finitely many \(c_u\) zero; equivalently,
\(P\in R[\lambda_1^{\pm1},\ldots,\lambda_r^{\pm1}]\).  Coefficient
extraction is \([\lambda^u]P:=c_u\).  When a
Boolean vector is used as an exponent, its entries are lifted from
\(\{0,1\}\) to \(\mathbb Z\), so exponent arithmetic is never modulo two.
For nonzero \(P\), put
\[
 m_j:=\max\!\bigl(0,-\min\{u_j:c_u\ne0\}\bigr),
 \qquad m=(m_1,\ldots,m_r).
\]
Then \(m\in\mathbb Z_{\ge0}^r\) is coordinatewise minimal with
\(\lambda^mP\in R[\lambda_1,\ldots,\lambda_r]\); thus multiplication by
\(\lambda^m\) clears all negative exponents.  The coordinate degree in
\(\lambda_j\) of a nonzero \(P\) is
\[
 \deg_{\lambda_j}P
 :=\max\{u_j:c_u\ne0\}-\min\{u_j:c_u\ne0\}.
\]
By convention, \(\deg_{\lambda_j}0:=-\infty\); consequently, every finite
upper degree bound used below includes the zero Laurent polynomial.

\subsubsection{Gauges and product weights}
\label{subsubsec:gauge-product-weights}

A \emph{portwise gauge} is an invertible diagonal rescaling on specified
tensor ports; on a binary flattening it gives invertible diagonal row or
column scaling.  It is an analytic coordinate operation, not an available
gadget unless separately realized.  A genuine character of
\((\Ftwo^r,+)\) has the form
\(z\mapsto\prod_j\epsilon_j^{z_j}\) with
\(\epsilon_j\in\{1,-1\}\).  More generally,
\[
 z\longmapsto\prod_j\gamma_j^{z_j},
 \qquad \gamma_j\in\mathbb C^\times,
\]
is a \emph{product weight}; it is a
character exactly when every \(\gamma_j^2=1\).

\subsection{Hilbert-space, Pauli, and stabilizer conventions}
\label{subsec:stabilizer-conventions}

The stabilizer arguments in \cref{sec:v4-deck,sec:platonic-decks} view
signature coefficients as amplitudes and Pauli contractions as linear
measurements.  Commuting Pauli labels describe affine support and phase
constraints; an encoder normal form recovers tensor factors from those
constraints.  These are analytic descriptions of signatures.  Each gadget
realization is established separately with the native bilinear contraction.

\subsubsection{Hilbert-space conventions}
\label{subsubsec:hilbert-conventions}

A finite-dimensional complex Hilbert space is a complex vector space with a
Hermitian inner product, conjugate-linear in its first argument, linear in
its second, satisfying
\(\langle v,u\rangle=\overline{\langle u,v\rangle}\), and positive definite
in the sense that \(\langle u,u\rangle>0\) for \(u\ne0\).  For a matrix \(A\), \(A^\dagger\)
is its conjugate transpose and
\(\langle u,v\rangle:=u^\dagger v\).  Put
\(\mathcal H_n:=(\mathbb C^2)^{\otimes n}\), and write
\[
 |x\rangle:=|x_1\rangle\otimes\cdots\otimes|x_n\rangle,
 \qquad
 \langle x|:=|x\rangle^\dagger
 \quad(x\in\Ftwo^n).
\]
A unitary on \(\mathcal H_n\) satisfies
\(U^\dagger U=UU^\dagger=I_{\mathcal H_n}\), while an isometry
\(V:\mathcal H_k\to\mathcal H_n\) satisfies
\(V^\dagger V=I_{\mathcal H_k}\).
These Hermitian conventions are distinct from the bilinear Frobenius pairing
used for Holant contractions and from isometries of symmetric binary forms.

A density operator is a positive-semidefinite operator of trace one, where
positive semidefinite means \(u^\dagger\rho u\ge0\) for all \(u\).  A nonzero
vector \(|\psi\rangle\) has normalized density operator
\[
 \rho_\psi=\frac{|\psi\rangle\langle\psi|}
                  {\langle\psi,\psi\rangle}.
\]
Thus a nonzero ket, modulo multiplication by a nonzero scalar, determines a
\emph{pure state} through \(\rho_\psi\).
For a density operator \(\rho\) on
\(\mathcal H_A\otimes\mathcal H_B\), fix an orthonormal basis
\(\{|b_j\rangle\}_j\) of \(\mathcal H_B\).  The partial trace is specified,
for every \(|a\rangle,|a'\rangle\in\mathcal H_A\), by
\[
 \langle a|\operatorname{Tr}_B(\rho)|a'\rangle
 =\sum_j\langle a,b_j|\rho|a',b_j\rangle.
\]
The result is independent of the chosen orthonormal basis of
\(\mathcal H_B\).  The operator
\(\operatorname{Tr}_B(\rho)\) is the \emph{marginal} on \(\mathcal H_A\); orthonormal means
\(\langle u_i,u_j\rangle=\one[i=j]\).  A one-qubit marginal is maximally
 mixed when it is \(I_{\mathbb C^2}/2\).  Every bipartite pure state has a Schmidt
decomposition
\[
 |\psi\rangle=\sum_{j=1}^s
 \sigma_j|u_j\rangle\otimes|v_j\rangle,\qquad \sigma_j>0,
\]
with both displayed families orthonormal; \(s\) is the Schmidt rank
\cite[Section~2.5]{NielsenChuang2010}.  For \(n\ge1\), \(r\in[n]\), and a
one-qubit linear operator \(R:\mathbb C^2\to\mathbb C^2\), define
\[
 R_r:=I_{\mathbb C^2}^{\otimes(r-1)}\otimes R\otimes
      I_{\mathbb C^2}^{\otimes(n-r)}:\mathcal H_n\to\mathcal H_n.
\]

\subsubsection{Pauli and stabilizer conventions}
\label{subsubsec:pauli-stabilizer}

The actual \(n\)-qubit Pauli group is
\[
 \mathcal P_n^{(4)}
 :=
 \{i^c\textstyle\bigotimes_{j=1}^nX^{a_j}Z^{b_j}:
   c\in\mathbb Z_4,\ a,b\in\Ftwo^n\}.
\]
Its projectivization by the scalar subgroup
\(\{\pm I_{\mathcal H_n},\pm iI_{\mathcal H_n}\}\) is identified with
\(\mathsf P([n])=\Ftwo^n\oplus\Ftwo^n\); the label
\(a\mid b\) is the concatenation of the two binary words.  This label space
has the symplectic form
\[
 \omega((a\mid b),(a'\mid b'))=a\cdot b'+b\cdot a'.
\]
Here symplectic means alternating
\(\omega(v,v)=0\) and nondegenerate
\(\{v:\omega(v,w)=0\text{ for all }w\}=\{0\}\).  For a subspace \(S\),
put \(S^{\perp_\omega}:=\{v:\omega(v,s)=0\text{ for every }s\in S\}\);
\(S\) is \emph{Lagrangian} when \(S=S^{\perp_\omega}\).  If \(S\) is
Lagrangian, a \emph{Lagrangian complement} of \(S\) is a Lagrangian subspace
\(T\) satisfying \(\mathsf P([n])=S\oplus T\).

A \emph{projective Clifford} is the phase-class, modulo multiplication by a
scalar in \(S^1\), of a unitary \(U\) satisfying
\(U\mathcal P_n^{(4)}U^\dagger=\mathcal P_n^{(4)}\).  Clifford conjugation induces a symplectic linear map
on Pauli labels \cite[Chapter~10]{NielsenChuang2010}.  A stabilizer code is
the common eigenspace
\[
 \{v:P v=\chi(P)v\text{ for every }P\in\widetilde S\}
\]
of a commuting Pauli subgroup \(\widetilde S\), where
\(\chi:\widetilde S\to\mu_4\) is an eigenvalue character agreeing with the
scalar by which every scalar Pauli acts.  A
one-dimensional stabilizer code is a stabilizer state, and a nonzero tensor
spanning it is a stabilizer tensor.

A stabilizer isometry is an isometry whose image is a stabilizer code and
whose logical Paulis have physical Pauli representatives; equivalently it is
a Clifford encoder with fixed stabilizer-state ancillas.  Here an
\emph{encoder} is the code isometry, a \emph{physical operator} acts on its
ambient output space, the corresponding \emph{logical operator} is the
induced operator on the encoded input space, and an \emph{ancilla} is a fixed
initialized tensor factor.  If \(S\) is the
projective stabilizer label space, the logical Pauli label space is
\(S^{\perp_\omega}/S\); for a \(2^k\)-dimensional code it has binary
dimension \(2k\).  A Clifford decoder is a Clifford unitary \(D\) satisfying
\[
 DV|\psi\rangle=|\psi\rangle\otimes|\eta\rangle
 \qquad(|\psi\rangle\in\mathcal H_k)
\]
for a fixed stabilizer ancilla \(|\eta\rangle\), namely an additional tensor
factor initialized in that state.  These standard
encoder--decoder and logical-Pauli facts follow from the stabilizer normal
form \cite[Chapters~3--4]{Gottesman1997}.

A stabilizer tableau is a binary symplectic generator matrix whose rows label
Pauli generators, together with their phases; multiplying generators is an
elementary row operation, while a Clifford coordinate change acts by a
symplectic change of columns \cite[Chapter~4]{Gottesman1997}.  A hyperbolic
pair is a pair \(p,q\) with
\(\omega(p,q)=1\); symplectic Gaussian elimination is row reduction by
changes of symplectic basis.  A nowhere-zero stabilizer quadratic phase is
\[
 \zeta(z)=\zeta_0i^{Q(z)},\qquad \zeta_0\in\mathbb C^\times,
\]
where \(Q:\Ftwo^k\to\mathbb Z_4\) is multilinear of degree at most two with
even mixed coefficients.  After fixed physical and logical Clifford
coordinate changes, a stabilizer isometry has amplitude normal form
\[
 V|z\rangle=\zeta(z)|c+M^{\mathsf T}z\rangle,
 \qquad
 c\in\Ftwo^n,\quad M\in\Mat_{k\times n}(\Ftwo),\quad\rank M=k,
\]
a standard consequence of stabilizer-tableau reduction
\cite[Chapter~4]{Gottesman1997}.

\subsubsection{Stabilizer error correction}
\label{subsubsec:stabilizer-error-correction}

Erasure of physical output \(r\) of an isometry \(V\) is correctable exactly
when the single-site Knill--Laflamme condition holds:
\[
 V^\dagger R_rV\in\mathbb C I_{\mathcal H_k}
 \qquad\text{for every one-qubit operator }R.
\]
This is the single-subsystem specialization of the Knill--Laflamme
criterion \cite{KnillLaflamme1997}.
In the stabilizer setting, an encoder tableau is the tableau of a Clifford
encoder; its symplectic row reduction records physical representatives of
logical Paulis and the residual stabilizer generators.

\subsection{Finite induction and termination}
\label{subsec:worklist-termination}

The reductions in \cref{subsec:safe-matching-cards} retain tensor factors
and sometimes enlarge the finite group of binary transfers.  The orders
below distinguish finite processing of a fixed family from a restart with
a smaller arity or larger group.

\subsubsection{Orders and termination}
\label{subsubsec:worklist-termination}

A strict order is irreflexive and transitive.  It is \emph{well founded}
when it has no infinite descending chain.  A finite multiset over a set
\(E\) records elements of \(E\) with multiplicity; in this paragraph
\(\dot\cup\) adds multiplicities, rather than denoting a disjoint union of
sets.  Given a strict order \(<\) on \(E\) and finite multisets \(A,B\) over
\(E\), its strict multiset extension is defined by \(A<_{\rm mul}B\) when
\[
 B=Z\mathbin{\dot\cup}X,\qquad
 A=Z\mathbin{\dot\cup}Y,\qquad X\ne\varnothing,
\]
for finite multisets \(X,Y,Z\) over \(E\), and every \(y\in Y\) satisfies
\(y<x\) for some \(x\in X\).  The multiset
extension of a well-founded strict order is well founded
\cite{DershowitzManna1979}.  Lexicographic order compares the first
coordinate at which two tuples differ; finite lexicographic products of
well-founded orders are well founded by induction on the number of
coordinates.  Well-founded induction proves a property at \(x\) assuming it
at every strict predecessor of \(x\).

We index unfinished signature occurrences by a finite multiset.  Each
occurrence retains its realizing gadget or reduction and its ordered ports.
To identify duplicate signatures, normalize by the first nonzero coefficient
in the fixed port order, recording the removed scalar separately.  Repeated
occurrences are still counted with multiplicity in the induction measures.
The reductions test their listed cases in order; the first applicable case
determines the next step.

\subsection{Retained reductions and normalization}
\label{sec:retained-normalization}

\subsubsection{Retained factor saturation}
\label{subsec:factor-saturation}

Record each recursive tensor factorization as a tree, retaining the ordered
port blocks, port permutation, and nonzero scalar at every split.  The leaves
are the tensor-prime factors.

For a preprocessed finite signature set \(\Lambda\), choose an exact
tensor-prime factorization of every member and let \(\Pi(\Lambda)\) be the
joint signature set of all prime factors, with ordered port blocks retained.
We call \(\Lambda\) \emph{factor-saturated} when
\(\Pi(\Lambda)\subseteq\Lambda\).

\begin{lemma}[Retained factor saturation]
\label{lem:factor-saturation}
If \(\Lambda\) is not common one-sided, then
\begin{equation}
  \KHolant(\Lambda)
  \equivT
  \KHolant(\Lambda,\Pi(\Lambda)).
  \label{eq:factor-saturation}
\end{equation}
The recursive state is the literal augmentation
\(\Lambda\cup\Pi(\Lambda)\), retaining in particular the original
two-sided charge witnesses.
Moreover, every invocation of factor saturation is a finite worklist
computation.  Start with one root occurrence for each member of the current
finite retained set and expose a factor only through the decomposition
theorem.  Let \(h_u\) denote the signature carried by a nonzero node \(u\).
Such a node, carrying its ordered port block, is either
certified tensor-prime, or is replaced by one split
\[
 h_u\doteq a_u\otimes b_u,
 \qquad 1\leq\arity(a_u),\arity(b_u)<\arity(h_u),
\]
with the nonzero scalar and the port permutation recorded in the provenance.
Nullary scalars are recorded and discarded, and unary nodes are certified
immediately and never enter the unresolved queue.  Fix the representative of a node by dividing by its first
nonzero entry and retaining its ordered port block; identical canonical keys
may share one retained signature, while their occurrences may still be
kept in the queue.  Let \(\mathcal W\) be the multiset of unresolved open
occurrences, all of arity at least two.  The potential
\begin{equation}
  \Phi(\mathcal W)
  =
  \sum_{u\in\mathcal W}(\arity(h_u)-1)
  \label{eq:factor-potential}
\end{equation}
is a nonnegative integer and strictly decreases at every worklist step.  Thus
\(\Pi(\Lambda)\) is the finite set of canonical prime leaves.  Its
independence of processing order follows from the tensor-prime factorization
uniqueness stated in \cref{subsec:holant-definitions}.
\end{lemma}

\begin{proof}
Apply the decomposition theorem at each queued split
\(h_u=a_u\otimes b_u\) to the current retained signature set
\cite[full version, Theorem~32 and Remark~39]{MengWangXiaZheng2025}.
The original non-one-sided signature set remains literally present, so the
exceptional case never reappears.  Iteration exposes every factor, while
literal inclusion gives the reverse reduction; these two directions give
\cref{eq:factor-saturation}.  Replacing a nonunary work item of arity
\(r+s\) removes \(r+s-1\) and can add at most
\((r-1)+(s-1)=r+s-2\); already queued or certified canonical children only
lower this amount further.  Certifying a prime removes its summand.  Hence the
factorization tree and the retained augmentation are finite, and all reduction
scalars and ordered blocks are available for later gadgets.
\end{proof}

For a factor-saturated signature set \(\Lambda\) containing a tensor-prime
signature of arity at least three, define its minimum nonbinary prime arity by
\begin{equation*}
 \nu(\Lambda)
 =\min\{\arity(h):h\in\Lambda
          \text{ is tensor-prime and }\arity(h)\ge3\}.
\end{equation*}
This parameter is used only after the branch in which every prime factor has
arity at most two has been removed.

\subsubsection{Downward inheritance and global normalization}

\begin{lemma}[Downward inheritance]
\label{lem:downward-inheritance}
If \(\cG\subseteq\cG'\) are preprocessed signature sets and
\(\Tract(K\cG')\) holds, then \(\Tract(K\cG)\) holds.
\end{lemma}

\begin{proof}
All alternatives except the single-weight one are hereditary by definition.
In the latter case a signature subset either retains both charge signs and the
same EO-completion witness, or becomes one-sided.  In the latter event, its
central members are unchanged by both restriction and EO completion, whereas
all off-central members restrict to zero.  Hence its nonzero central
restrictions form a subset of the original completion signature set, and the
same shared EO witness applies.
\end{proof}

\begin{lemma}[\(GO(X)\)-invariance of transported tractability]
\label{lem:GOX-tract-invariance}
Let \(\cG\) be a finite preprocessed signature set and let
\(D\in GO(X)\).  Then
\begin{equation*}
  \Tract(K\cG)
  \quad\Longleftrightarrow\quad
  \Tract\bigl(K(D\cG)\bigr).
\end{equation*}
\end{lemma}

\begin{proof}
Write the matrix \(D=\begin{psmallmatrix}a&b\\ c&d\end{psmallmatrix}\), and
let \(\lambda\in\AlgNums^\times\) be its conformal multiplier, defined by
\(D^{\mathsf T}XD=\lambda X\).  Comparing entries gives
\(ac=bd=0\) and \(ad+bc=\lambda\ne0\).  Thus \(D\) is diagonal or
anti-diagonal.  The former preserves every support word; the latter
complements it; on a fixed layer its multiplier is constant.  Hence the two
EO alternatives are preserved (with the support side and shared sign
possibly exchanged), tensor-prime factor arities are unchanged, and
\(\angles{\cM}\) and \(\angles{X\cM}\) are preserved or exchanged.

Define the conjugate matrix \(A=KDK^{-1}\), so
\(A^{\mathsf T}A=\lambda I\).  If the matrix \(W\) is a common
\(\cC\)-transformability witness, define the transported witness matrix
\(W':=AW\).  It satisfies
\[
  (W')^{-1}A(K\cG)=W^{-1}(K\cG)\subseteq\cC,
  \qquad
  (W')^{\mathsf T}W'
  =\lambda W^{\mathsf T}W\in\cC
\]
for \(\cC\in\{\cA,\cP,\cL\}\).  Applying the same argument to \(D^{-1}\)
proves the converse.
\end{proof}

Every reduction, basis change, group enlargement, hardness conclusion, and
tractable presentation obtained inside a coordinate-transport branch can be
transported correctly back to the original signature set.  The following
lemma records the data needed for that transport.  Its continuing-state use
is confined to the self-contained eight-vertex interface; later sections may
invoke it to interpret a call to that interface, but do not use a realized
weighted equality to replace the supplied literal \(I\).

\begin{lemma}[Normalization data]
\label{lem:normalization-ledger}
Let \(\cG_0\) be the signature set at the beginning of a structural
branch in which coordinate transport is permitted.  Let \(L_0\) be a fixed
number field containing all entries of
\(\cG_0\), their complex conjugates, and all algebraic constants already fixed by
the coordinate conventions.  Thus \(L_0\) is the frozen initial coefficient field, not a
field enlarged after each normalization.  Every continuing normalization
inside that branch is an explicitly recorded \(GO(X)\) change of coordinates, and tensor
splits are performed after pulling back by the current \(A\), using a nonzero
pivot and row--column ratios.  Thus the pulled-back data (rather than the
displayed current coordinates) stay over \(L_0\), and the pulled-back group
remains inside \(\PGL_2(L_0)\).  The
structural reductions maintain a current retained signature set \(\Lambda\),
a matrix \(A\in GO(X)\), and its nonzero conformal factor
\(\mu_A\in\AlgNums^\times\), satisfying
\begin{equation}
  A\cG_0\subseteq\Lambda,\qquad
  A^{\mathsf T}XA=\mu_A X,\qquad
  \KHolant(\Lambda)\equivT\KHolant(\cG_0).
  \label{eq:normalization-ledger}
\end{equation}
Every positive-arity \(h\in\Lambda\), up to a nonzero scalar, also has the
field-compatible form
\begin{equation*}
 h=A^{\otimes\arity(h)}h_0
 \quad\text{for a tensor }h_0\text{ over }L_0.
\end{equation*}
At a residual node where every directly realizable nonsingular transfer has
finite projective order, these transfers form a group; denote by
\(G(\Lambda)\) the group of their classes \([BX]\), where \(B\) ranges over
directly realizable nonsingular ordered binaries over \(\Lambda\).  Then the
pulled-back group
\begin{equation*}
 \overline G(A,\Lambda)
 :=[A]^{-1}G(\Lambda)[A]
 \le \PGL_2(L_0)
\end{equation*}
is unchanged by a global normalization and is monotone under retained
augmentation.
Direct gadget augmentation and retained factor saturation preserve \(A\).
If \(D\in GO(X)\) is applied globally to the current signature set, the record
updates by
\begin{equation}
  (A,\Lambda)\longmapsto(DA,D\Lambda).
  \label{eq:normalization-ledger-update}
\end{equation}
The new state again satisfies \cref{eq:normalization-ledger}, up to the
known edge scalar from \cref{eq:GOX-global-scalar}.  Before a continuing
branch leaves \cref{sec:p1-eight-vertex}, apply the inverse of the accumulated
matrix \(A\).  The returned state therefore has \(A=I\) in the caller's fixed
coordinates and retains any supplied literal equality literally.
Generic interpolation is used only as a terminal transition and is not a
record-preserving update.  If an interpolation output were ever allowed to
re-enter a continuing finite-group state, the effective coefficient field
and the pulled-back group would first have to be enlarged explicitly.

If \(\Tract(K\Lambda)\) holds, then
\(\Tract(K\cG_0)\) holds.  Separately, if
\(V\) is a common \(\cC\) \(K\)-presentation of \(\Lambda\), where
\(\cC\) is closed under nonzero scalar multiplication, then the pullback
matrix
\begin{equation*}
  W:=A^{-1}V
\end{equation*}
is a common \(\cC\) \(K\)-presentation of \(\cG_0\).  In particular, this
applies to \(\cA,\cP,\cL\) and to the flat-Lagrangian subclass.
\end{lemma}

\begin{proof}
The assertions hold initially for \((I,\cG_0)\).  Direct augmentation and
factor saturation preserve them by
\cref{eq:direct-gadget-reduction,eq:factor-saturation}.  A global
\(D\in GO(X)\), with nonzero conformal factor
\(\mu_D\in\AlgNums^\times\) defined by \(D^{\mathsf T}XD=\mu_DX\), replaces \(A\) by
\(DA\); then
\((DA)^{\mathsf T}X(DA)=\mu_D\mu_AX\), and
\cref{eq:GOX-global-scalar} supplies the stated equivalence.  Conversely,
the transported inclusion \(A\cG_0\subseteq\Lambda\), followed by the
inverse global transformation, reduces the initial problem to the current
one; this is why the record records equivalence rather than only one
direction.

Inductively pull every construction through \(A\).  Gadget contraction
stays over \(L_0\); at a tensor split, normalizing one nonzero entry and
taking row--column ratios chooses the factors over \(L_0\).  Moreover
\(A^{\mathsf T}X\doteq XA^{-1}\).  Thus, if a current binary \(B\) has the
record form \(B=A^{\otimes2}B_0\) for a binary tensor \(B_0\) over \(L_0\),
its projective transfer pulls back by
conjugation:
\[
 [BX]=[A(B_0X)A^{-1}].
\]
The identity transfer is realized by the native edge, and chaining realizes
products.  At such a finite-projective-order residual node, the inverse of each transfer is a
positive power of that transfer, so the realizable classes form a group.
The displayed conjugation shows that global normalization leaves
\(\overline G(A,\Lambda)\) unchanged, while retained augmentation can only
add realizable transfers.  This proves the field and group assertions.
Finally, if \(\Tract(K\Lambda)\) holds, downward inheritance first gives
\(\Tract(K(A\cG_0))\), and \cref{lem:GOX-tract-invariance} then gives
\(\Tract(K\cG_0)\).  If \(V\) is an explicit common presentation, then
\(W=A^{-1}V\) satisfies
\[
 W^{-1}\cG_0\subseteq V^{-1}\Lambda\subseteq\cC,
 \qquad
 W^{\mathsf T}XW=\mu_A^{-1}V^{\mathsf T}XV\in\cC,
\]
which proves the last assertion.
\end{proof}
\section{Reductions Using Unaries and Equalities}
\label{sec:p1-shared-terminals}

The imported odd-arity dichotomy and factor theorem
(\cref{subsec:imported-results}) resolve a branch once a nonzero unary is
accessible.  This section supplies the reductions that obtain such a unary
from a singular or infinite-order binary, a single-occurrence carrier, or
a higher equality.  We use the distinction between direct realization and
Turing access fixed in \cref{subsec:reduction-conventions}.

\subsection{Binary Interpolation}
\label{proofsubsec:central-descent}

\begin{lemma}[Binary interpolation]
\label{lem:binary-interpolation}
Let \(\Gamma\) be a finite algebraic, retained, two-sided signature set,
let \(B\ne0\) be an ordered binary directly realizable over \(\Gamma\),
and define its transfer matrix \(A=BX\).
\begin{enumerate}
\item If \(B\) has the rank-one factorization \(B=u\otimes v\), where
      \(u\) and \(v\) are nonzero algebraic unary signatures, then
      \[
        \KHolant(\Gamma,u,v)\leT\KHolant(\Gamma).
      \]
\item If \(A\) is nonsingular and has infinite projective order, there is
      a nonzero rank-one binary \(H\) such that
      \[
        \KHolant(\Gamma,H)\leT\KHolant(\Gamma).
      \]
      For \(m\ge1\) designated occurrences of \(H\), the interpolation uses
      \(m+1\) oracle queries, \(O(m^2)\) added vertices per query, and hence
      \(O(m^3)\) added vertices over all queries.
\end{enumerate}
\end{lemma}

\begin{proof}
For the singular case, choose a nonzero algebraic pivot of the rank-one table
and take the corresponding algebraic row and column factors (absorbing the
pivot into one factor); then \(B=u\otimes v\) with \(u,v\) algebraic, and
\cref{thm:external-factor} applies.  Two-sidedness excludes its one-sided
exception.  For the other case, a length-\(t\) chain realizes the ordered
binary signature \(B_t:=A^tX\).  If \(A\) is diagonalizable, choose an eigenbasis matrix
\(V\in\GL_2(\AlgNums)\) and nonzero eigenvalues
\(\lambda_0,\lambda_1\in\AlgNums^\times\) such that
\(A=V\diag(\lambda_0,\lambda_1)V^{-1}\).  For \(j\in\{0,1\}\), let
\(E_{jj}\) be the \(2\times2\) matrix with \((a,b)\)-entry
\(\one[a=b=j]\), for \(a,b\in\{0,1\}\), and define the rank-one binary
matrix \(H_j=VE_{jj}V^{-1}X\).  Then
\[
 B_t=\lambda_0^tH_0+\lambda_1^tH_1,
\]
with \(H_0,H_1\) rank one.  Replacing the \(m\) designated \(H_0\)-vertices by \(B_t\) and
dividing by \(\lambda_1^{tm}\) yields evaluations of a polynomial \(P(z)\)
of degree at most \(m\) at the distinct points
\(z=(\lambda_0/\lambda_1)^t\), \(1\le t\le m+1\); the coefficient of
\(z^m\) is the all-\(H_0\) value.  These points are distinct because
\([A]\) has infinite projective order, so \(\lambda_0/\lambda_1\) is not a
root of unity.

In the remaining infinite-order Jordan case, write \(A=\lambda I+N\), where
\(\lambda\in\AlgNums^\times\) is the eigenvalue and \(N\) is a nonzero
nilpotent \(2\times2\) matrix satisfying \(N^2=0\).  Then
\(B_t=\lambda^t(X+t\lambda^{-1}NX)\).  After normalization the oracle value
is a polynomial in \(t\) of degree at most \(m\), whose coefficient of
\(t^m\) is the value with the rank-one binary \(\lambda^{-1}NX\) at every
designated vertex.
Indeed, \(N\ne0\) and \(N^2=0\) in dimension two imply
\(\rank(N)=1\), and right multiplication by \(X\) preserves
rank.
Both cases use \(m+1\) queries; \(t\le m+1\) gives \(O(m^2)\) replacement
size for each query and \(O(m^3)\) in total.  The eigenvalues and basis
matrices are fixed algebraic numbers.  Their powers and the interpolation
systems have polynomial encoding length, so exact interpolation is
polynomial-time postprocessing.
\end{proof}

\subsection{Single-occurrence carriers and unary terminals}
\label{subsec:single-occurrence-carriers}

The next reductions initially permit only one occurrence of an auxiliary
signature.  A context that returns this signature together with a new
tensor can be nested to make arbitrarily many copies of the latter.
The reserved carrier ports keep their order under nesting, with no extra
port identifications.
Define the ordered carrier signature
\[
  \Delta=\DeltaZero\otimes\DeltaOne,
\]
with the displayed order of its two ports.  For a signature \(g\), the
superscripts \(g^{=1}\) and \(g^{\leq1}\) mean that an oracle instance contains,
respectively, exactly one or at most one \(g\)-vertex.  They do not denote
tensor powers.

\begin{lemma}[Universal one-pair reduction]
\label{lem:universal-directed-pair}
For every fixed finite algebraic signature set \(\Gamma\),
\[
  \KHolant(\Gamma,\Delta^{\leq1})
  \leT
  \KHolant(\Gamma).
\]
No EO, parity, charge, or tractability hypothesis is required.
\end{lemma}

\begin{proof}
The zero-occurrence case is already an instance over \(\Gamma\).  Consider an
instance with one \(\Delta\)-vertex.  Delete that vertex, retaining its two
incident native \(X\)-attachments as wires to ordered free boundary ports.
Write the raw table, before the boundary-coordinate correction in
\cref{eq:half-edge-boundary-tensor}, as
\[
  H=\begin{pmatrix}h_{00}&h_{01}\\h_{10}&h_{11}\end{pmatrix},
  \qquad\text{the desired one-\(\Delta\) value is }h_{01}.
\]
Closing its two exposed ports by a native \(X\)-edge gives an oracle value
\[
  S=h_{01}+h_{10}.
\]
Let \(\mathcal B_\Gamma\) be the set of ordered binary tensors realized by
finite \(\Gamma\)-networks.  If every \(B\in\mathcal B_\Gamma\) satisfies
  \(B_{01}=B_{10}\), then \(XHX\in\mathcal B_\Gamma\) is symmetric, so
  \(H\) is symmetric and the desired value is
  \(S/2\).

Otherwise fix once and for all a realized ordered binary
\[
  B=\begin{pmatrix}a&b\\c&d\end{pmatrix},
  \qquad b\ne c.
\]
Insert its fixed network at the deleted location in the two boundary
orientations.  The two oracle values are
\[
  T_+=ah_{00}+bh_{01}+ch_{10}+dh_{11},
  \qquad
  T_-=ah_{00}+ch_{01}+bh_{10}+dh_{11}.
\]
Thus
\[
  h_{01}
  =\frac12\left(S+\frac{T_+-T_-}{b-c}\right).
\]
The ordered ports are essential in the subtraction.  The witness network for
\(B\) is fixed with \(\Gamma\), hence has constant size, and the reduction
uses at most three oracle queries.
\end{proof}

The next elementary lemma isolates the only additional ingredient needed to
turn one-occurrence access into unrestricted access to a new signature.  It
applies verbatim to ordinary Holant and to \(\KHolant\); write \(\mathsf H\) for
either setting, with its corresponding edge semantics, and write
\(Z_{\mathsf H}\) for the corresponding closed-instance value.

\begin{lemma}[Single-occurrence replication]
\label{lem:single-occurrence-replication}
Let \(\Gamma\) be a fixed signature set and let \(g,h\) be nonzero signatures.  Suppose that
\begin{enumerate}
\item \(\mathsf H(\Gamma,g^{=1})\leT\mathsf H(\Gamma)\);
\item a fixed carrier-preserving context \(C\) over \(\Gamma\) has a
      distinguished arity-\(\arity(g)\) unassigned vertex and an ordered output
      boundary consisting of two disjoint ordered blocks \(P\) followed by
      \(H\), where \(|P|=\arity(g)\) and \(|H|=\arity(h)\).  The block \(P\)
      is the \(g\)-carrier, and filling the hole by \(g\) realizes
      \(\rho\,g(P)\otimes h(H)\) for a known nonzero scalar
      \(\rho\in\AlgNums^\times\).  Substituting a
      copy of \(C\) into the same distinguished hole (plugging its carrier
      into the next unassigned vertex, with no other port identifications) is a
      legal composition; hence, for every \(k\ge1\), the \(k\)-fold nested
      context has exactly one \(g\)-occurrence and realizes
      \(\rho^k g(P)\otimes\bigotimes_{\ell=1}^k h(H_\ell)\), with disjoint
      output blocks \(H_1,\ldots,H_k\); and
\item some fixed \(\mathsf H(\Gamma,g^{=1})\)-instance \(J\) has
      \(Z_{\mathsf H}(J)\ne0\).
\end{enumerate}
Then
\[
  \mathsf H(\Gamma,h)\leT\mathsf H(\Gamma).
\]
\end{lemma}

\begin{proof}
By the designated-carrier composition in item~2, iterating the gadget
produces \(g\otimes h^{\otimes k}\) with one \(g\)-vertex, size \(O(k)\),
and a known nonzero scalar.  Hence an instance containing one \(g\) and \(k\) designated
\(h\)-vertices reduces to an instance containing exactly one \(g\); divide
out the recorded scalar.  For an instance \(\mathcal I\) over \(\Gamma\cup\{h\}\), take the
disjoint union \(\mathcal I\sqcup J\).  Its value is
\(Z_{\mathsf H}(\mathcal I)Z_{\mathsf H}(J)\), so division by
the fixed nonzero number \(Z_{\mathsf H}(J)\) inserts the one carrier required by the
first reduction.  The context and \(J\) are fixed, so the construction has
size linear in the number of \(h\)-occurrences.  Composing with item~1
proves the claim.
\end{proof}

\begin{lemma}[Common unary exit]
\label{lem:common-unary-exit}
Let \(\cF\) be an ordinary-coordinate signature set, let \(\Gamma\) be a
native \(K\)-coordinate signature set, and let \(u\ne0\) be a unary signature.
If \(\Holant(\cF,u)\leT\Holant(\cF)\), then
\(\Holant(\cF)\) is \(\SharpP\)-hard unless \(\Tract(\cF)\) holds.  Likewise,
if \(\KHolant(\Gamma,u)\leT\KHolant(\Gamma)\), then
\(\KHolant(\Gamma)\) is \(\SharpP\)-hard unless \(\Tract(K\Gamma)\) holds.
\end{lemma}

\begin{proof}
The first assertion is \cref{thm:external-odd} applied to the augmented
signature set.  A tractable presentation of that set restricts to one for
\(\cF\).  For the second, transform to ordinary coordinates:
\[
  \KHolant(\Gamma,u)
  \equivT
  \Holant(K\Gamma,Ku).
\]
The signature \(Ku\) is a nonzero unary, so the same argument applies.
\end{proof}

For unary columns \(g,u\), their contraction is the bilinear scalar
\[
  g^{\mathsf T}u:=\sum_{x\in\{0,1\}}g(x)u(x),
\]
with no complex conjugation.
For equality wires we need a carrier whose two unary factors have nonzero
contraction.  The following proof is included because the carrier need not
itself be realizable by a gadget.

\begin{lemma}[Unary companion]
\label{lem:unary-companion}
Let \(\cF\) be a fixed finite algebraic signature set and let \(g\ne0\) be
unary.  There is an algebraic unary \(u\) such that
\[
  g^{\mathsf T}u\ne0,
  \qquad
  \Holant(\cF,(g\otimes u)^{=1})\leT\Holant(\cF).
\]
\end{lemma}

\begin{proof}
We first prove the claim for \(g=\DeltaZero\).  Let \(\mathcal B_\cF\) be
the set of ordered binary tensors realized by finite \(\cF\)-networks with
ordinary equality wires.  Deleting one occurrence of
\(\DeltaZero\otimes u\), where \(u=[s,t]^{\mathsf T}\), leaves a binary
context \(H=(h_{ij})\), and the desired value is
\(sh_{00}+th_{01}\).  Equality-closing the exposed ports gives the oracle value
\[
  S:=h_{00}+h_{11}.
\]
We choose \(u\) by the following exhaustive alternatives.

If every \(H\in\mathcal B_\cF\) satisfies \(h_{00}=h_{11}\), take
\(u=[2,0]^{\mathsf T}\); then the desired value is \(S\).  Otherwise fix
\[
  G_1=\begin{pmatrix}a_1&b_1\\c_1&d_1\end{pmatrix}
  \in\mathcal B_\cF,
  \qquad a_1\ne d_1.
\]
If every realized binary has equal off-diagonal entries, then
\(b_1=c_1=b\).  For \(b=0\), take \(u=[1,0]^{\mathsf T}\).  Inserting
\(G_1\) gives the oracle value \(T:=a_1h_{00}+d_1h_{11}\), whence
\[
  h_{00}=\frac{T-d_1S}{a_1-d_1}.
\]
For \(b\ne0\), take \(u=[a_1-d_1,2b]^{\mathsf T}\).  The same insertion,
now using \(h_{01}=h_{10}\), gives
\[
  T-d_1S=(a_1-d_1)h_{00}+2bh_{01}.
\]

It remains that some
\[
  G_2=\begin{pmatrix}a_2&b_2\\c_2&d_2\end{pmatrix}
  \in\mathcal B_\cF
  \quad\text{satisfies}\quad b_2\ne c_2.
\]
Take \(u=[a_1-d_1,b_1+c_1]^{\mathsf T}\).  Let \(T_1\) be the value obtained
by inserting \(G_1\), and let \(T_2,T_3\) use the two orientations of
\(G_2\).  Then
\[
  \frac{T_2-T_3}{b_2-c_2}=h_{01}-h_{10},
\]
and therefore
\[
  T_1-d_1S+
  \frac{c_1}{b_2-c_2}(T_2-T_3)
  =(a_1-d_1)h_{00}+(b_1+c_1)h_{01}.
\]
In every case the first entry of \(u\) is nonzero, so
\(\DeltaZero^{\mathsf T}u\ne0\).

Now write the unary column vector \(g=[a,b]^{\mathsf T}\).  If
\(a^2+b^2=0\), then projectively
\(g\doteq[1,i]^{\mathsf T}\) or \(g\doteq[1,-i]^{\mathsf T}\).  In the first
case take \(u=[1,-i]^{\mathsf T}\), and in the second take
\(u=[1,i]^{\mathsf T}\).  Under the ordinary-to-native
\(K^{-1}\)-transformation the pair is, up to
a known nonzero scalar and port order,
\(\DeltaZero\otimes\DeltaOne\).  The reduction follows from
\cref{lem:universal-directed-pair}, and \(g^{\mathsf T}u\ne0\).

If \(a^2+b^2\ne0\), choose the scalar \(r\in\AlgNums\) with
\(r^2=a^2+b^2\) and define the matrix
\[
  Q=\frac1r\begin{pmatrix}a&b\\b&-a\end{pmatrix}.
\]
Then \(Q^{\mathsf T}Q=I\) and \(Qg=r\DeltaZero\).  Apply the point-mass
case to \(Q\cF\), obtaining \(u'=[s,t]^{\mathsf T}\) with \(s\ne0\), and
define the unary vector \(u=Q^{-1}u'\).  Because \(Q^{\mathsf T}Q=I\), the
holographic identity leaves every ordinary equality edge unchanged and pulls
the reduction back to \(\cF\), while
\[
  g^{\mathsf T}u=(Qg)^{\mathsf T}u'=rs\ne0.
\]
\end{proof}

\begin{corollary}[Higher equality exposes unaries]
\label{cor:higher-equality-unary}
Let \(d\ge4\).  If \(\Eq_d\) is accessible over a finite algebraic signature
set in either ordinary Holant or \(\KHolant\), then both unary point
masses are accessible in that setting.  In particular, access to one nonzero
unary is enough to invoke \cref{lem:common-unary-exit}.
\end{corollary}

\begin{proof}
First work with ordinary equality wires and adjoin the accessible \(\Eq_d\).
For \(j\in\{0,1\}\), apply \cref{lem:unary-companion} to
\(e_j=\delta_j\), obtaining \(u_j\) with \(u_j(j)\ne0\) and the
one-occurrence carrier \(v_j=e_j\otimes u_j\).  Attach the \(e_j\)-port of
that carrier to one port of \(\Eq_d\).  Up to port order the output is
\[
  v_j\otimes e_j^{\otimes(d-2)}.
\]
Closing the carrier through equality has value \(u_j(j)\ne0\), so
\cref{lem:single-occurrence-replication} gives access to
\(e_j^{\otimes(d-2)}\).  Passing temporarily to \(K\)-coordinates, the unary
\(K^{-1}e_j\) has two nonzero entries; hence this tensor power is not in the
common-one-sided exception of \cref{thm:external-factor}.  Factor extraction
therefore exposes \(e_j\) after transforming back.

Now work in \(\KHolant\) and use the carrier \(\Delta\).  Attach its two factors,
through native \(X\)-edges and in the two orientations, to two copies of
\(\Eq_d\).  The output is
\[
  \Delta^{\otimes(d-1)}
  =\Delta\otimes\Delta^{\otimes(d-2)}
\]
up to port order.  The native \(X\)-closure of \(\Delta\) has value one, so
\cref{lem:universal-directed-pair,lem:single-occurrence-replication} give
access to \(\Delta^{\otimes(d-2)}\).  The available \(\Eq_d\) contains both
endpoint words \(\zero^d,\one^d\), of doubled charges \(-d,+d\), and hence
excludes the common-one-sided exception; applying
\cref{thm:external-factor} exposes \(\DeltaZero\) and \(\DeltaOne\).  Finally
compose with the assumed reduction that removes \(\Eq_d\).
\end{proof}

\subsection{Equality and generalized-equality anchors}
\label{proofsubsec:equality-anchor}

An equality anchor is a quaternary gadget that becomes \(\Eq_4\) under
the same basis change that makes the native edge an equality.  This gives
a whole-problem reduction to the unary result above.

\begin{definition}[Equality anchor]
\label{def:equality-anchor}
A signature set \(\cG\) has an \emph{equality anchor} if it directly
realizes an ordered quaternary signature \(g\) for which there exist a matrix
\(U\in\GL_2(\AlgNums)\) and nonzero algebraic scalars
\(\kappa,\lambda\in\AlgNums^\times\) satisfying
\[
  (U^{-1})^{\mathsf T}XU^{-1}=\kappa\,\Eq_2,
  \qquad
  U^{\otimes4}g=\lambda\,\Eq_4.
\]
\end{definition}

\begin{lemma}[Equality-anchor lifting]
\label{lem:equality-anchor}
Let \(\cG\) be a finite algebraic signature set with an equality anchor
\(g\), and fix witnesses \(U,\kappa,\lambda\) as in
\cref{def:equality-anchor}.  Then either \(\KHolant(\cG)\) is
\(\SharpP\)-hard or \(\Tract(K\cG)\) holds.
\end{lemma}

\begin{proof}
Apply the native-edge holographic identity with \(M=U^{-1}\) to the whole
problem.  It sends
each right signature \(f\) to \(Uf\) and the native edge to
\((U^{-1})^{\mathsf T}XU^{-1}=\kappa\Eq_2\); for a queried instance
\(\Omega\), the edge factor \(\kappa^{|E_{\rm int}(\Omega)|}\) is known.
Thus the direct anchor becomes
\(Ug=\lambda\Eq_4\), so \(\Eq_4\) is accessible in the ordinary-Holant
problem over \(U\cG\).  By
\cref{cor:higher-equality-unary}, a nonzero unary \(u\) is accessible there.
Pulling the reduction back exposes the nonzero unary \(U^{-1}u\) over
\(\cG\).  The conclusion is now \cref{lem:common-unary-exit}.
\end{proof}

\begin{lemma}[Pure generalized-equality lifting]
\label{lem:pure-ge}
Suppose a retained finite algebraic signature set \(\cG\) directly realizes
the arity-\(q\) pure generalized-equality signature
\[
  G_q:=c_0\delta_{\zero^q}+c_1\delta_{\one^q},
  \qquad c_0,c_1\in\AlgNums^\times,\qquad q\ge4.
\]
Then either \(\KHolant(\cG)\) is \(\SharpP\)-hard or
\(\Tract(K\cG)\) holds.
\end{lemma}

\begin{proof}
Define the retained signature set \(\Lambda=\cG\cup\{G_q\}\).  Direct
realizability gives
\(\KHolant(\Lambda)\leT\KHolant(\cG)\); work over \(\Lambda\).
Use the ordered carrier \(\Delta\) from
\cref{subsec:single-occurrence-carriers}.  Attach its
two factors through native \(X\)-edges, in opposite orientations, to one port
of each of two copies of \(G_q\).  The two copies are forced onto opposite
endpoint words, and the external tensor is
\[
  c_0c_1\,\Delta^{\otimes(q-1)}
  =
  c_0c_1\,\Delta\otimes\Delta^{\otimes(q-2)}
\]
up to port order.  The native \(X\)-closure of the carrier has value one.
Thus \cref{lem:universal-directed-pair,lem:single-occurrence-replication}
give access to \(\Delta^{\otimes(q-2)}\).  Since \(G_q\in\Lambda\) contains
both endpoint words, of doubled charges \(-q,+q\),
\cref{thm:external-factor} has no common-one-sided
exception and exposes \(\DeltaZero\) and \(\DeltaOne\).  Apply
\cref{lem:common-unary-exit} over \(\Lambda\).  A hard conclusion composes
back to \(\cG\), while
\(\Tract(K\Lambda)\) implies \(\Tract(K\cG)\) by
\cref{lem:downward-inheritance}.
\end{proof}
\section{Eight-Vertex Signatures with Nonzero Endpoints}
\label{sec:p1-eight-vertex}

The quaternary dichotomy in \cref{thm:external-eight-vertex} leaves
product, affine, and local-affine exceptions.  We analyze these exceptions
over the entire retained signature set, using the transformability
convention of \cref{def:common-transformability} and the unary/equality
reductions of \cref{sec:p1-shared-terminals}.  The continuing outcome will
be a weighted-equality factor of the original quaternary.

\begin{definition}[Endpoint-nondegenerate eight-vertex signature]
\label{def:p1-endpoint-eight-vertex}
An \emph{endpoint-nondegenerate eight-vertex signature} is a quaternary
signature \(q\) satisfying
\[
 q(x)=0\quad\text{when \(\wt(x)\) is odd},
 \qquad q(0000)q(1111)\ne0.
\]
Thus its support is contained in the eight even-parity words, while both
endpoint values are nonzero.
\end{definition}
We use the two entries \(q(0000),q(1111)\), the six weight-two entries, and
the three flattenings with row--column splits
\(12\mid34\), \(13\mid24\), and \(14\mid23\), respectively; for example,
\(12\mid34\) uses \((x_1,x_2)\) as row index and \((x_3,x_4)\) as column
index.  We list their ranks in this order.
A complementary pair of two-port blocks is \(J,J^c\subseteq[4]\) with
\(|J|=2\).

We classify endpoint-nondegenerate eight-vertex signatures first in the
product and local-affine cases and then in the sole finite affine base
case.  The resulting quaternary interface is used by the stable
equality-accessible matching-deck proof and yields the independent
quaternary dichotomy in \cref{thm:p1-quaternary-kholant-dichotomy}.
The two exact finite statements used in this section are given in
Appendix~\ref{appsec:cert-eight-vertex}.

Throughout this section, weighted equality and nonsingularity have the
meaning fixed in \cref{subsec:holant-definitions}.  For a signature class
\(\cC\), the statement that one signature \(M\) is \(\cC\)-transformable
uses the singleton convention: the set \(\{M\}\) and the transformed native
edge have one common transformability witness, as in
\cref{def:common-transformability}.

\subsection{Product and local-affine exceptions}

\begin{lemma}[Product-transformable eight-vertex signatures]
\label{lem:p1-eight-vertex-product-transform}
Let \(M\) be an endpoint-nondegenerate eight-vertex signature.  If \(M\) is
\(\cP\)-transformable in the sense of
\cref{thm:external-eight-vertex}, then exactly one of the following
structural outcomes applies (up to a port permutation):
\begin{enumerate}
\item \(M\) is a pure generalized equality;
\item \(M\) is a tensor product of two nonsingular weighted equalities;
\item \(M\) has full even support, has the dense form
      \cref{eq:p1-eight-vertex-product-dense-form}, and every retained signature set
      directly realizing \(M\) also directly realizes an equality anchor.
\end{enumerate}
In particular, every
\(\cP\)-transformable endpoint-nondegenerate eight-vertex signature has support size \(2\), \(4\), or
\(8\).
\end{lemma}

\begin{proof}
Choose an algebraic witness matrix
\(T=(\begin{smallmatrix}a&b\\c&d\end{smallmatrix})\) by
\cref{lem:algebraic-transform-witness}, and define the transformed signature
\(H=(T^{-1})^{\otimes4}M\in\cP\).  The transformed edge is
\[
 T^{\mathsf T}XT=
 \begin{pmatrix}2ac&ad+bc\\ad+bc&2bd\end{pmatrix}.
\]
Because this transformed edge is nonsingular and belongs to \(\cP\), its
support is either equality or disequality.  Disequality support makes
\(T\) monomial.  Equality support gives \(ad+bc=0\); define the nonzero
scalar ratio \(\lambda:=b/a\).  Then
\[
 T=\diag(a,c)\begin{pmatrix}1&\lambda\\1&-\lambda\end{pmatrix},
 \qquad ac\lambda\ne0.
\]

If \(H\) is tensor-prime, product type makes it a generalized equality
for some \(s\in\Ftwo^4\), with nonzero coefficients
\(A,B\in\AlgNums^\times\):
\[
 H=A\delta_s+B\delta_{\bar s},\qquad AB\ne0.
\]
In the monomial case the two nonzero endpoint values force
\(\{s,\bar s\}=\{\zero^4,\one^4\}\), giving outcome~1.  Otherwise, for
\(k=\wt(s)\) and \(w=\wt(x)\),
\[
 M(x)=a^{4-w}c^w(-1)^{s\cdot x}
 \bigl(A\lambda^k+B\lambda^{4-k}(-1)^w\bigr).
\]
The odd-weight zeros are equivalent to \(B/A=\lambda^{2k-4}\).  Hence all
even entries are nonzero.  Define the nonzero overall scale
\(\rho:=2A\lambda^k a^4\) and weight-ratio scalar \(r:=c/a\); the entries
then have the form
\begin{equation}
 M(x)=\rho r^{\wt(x)}(-1)^{s\cdot x}
 \one[\wt(x)\text{ is even}],\qquad \rho r\ne0.
 \label{eq:p1-eight-vertex-product-dense-form}
\end{equation}

For this dense case choose nonzero algebraic scalars
\(q_0,q_1\in\AlgNums^\times\) satisfying the first two relations below,
define the diagonal matrix \(Q:=\diag(q_0,q_1)\), and define the
normalization matrix \(R:=QT^{-1}\):
\[
 q_0^2=2ac,\qquad q_1/q_0=i\lambda,\qquad R=QT^{-1}.
\]
Then, since
\(\lambda^{2k-4}(i\lambda)^{4-2k}=(-1)^k\), there is a nonzero scalar
\(\eta\in\AlgNums^\times\) such that
\begin{equation}
 (R^{-1})^{\mathsf T}XR^{-1}=I,\qquad
 R^{\otimes4}M=\eta\bigl(\delta_s+(-1)^k\delta_{\bar s}\bigr).
 \label{eq:p1-eight-vertex-product-anchor-normal-form}
\end{equation}
Define the signature
\(F_s=\delta_s+(-1)^k\delta_{\bar s}\).  In these coordinates the
native edge is equality by \cref{eq:p1-eight-vertex-product-anchor-normal-form}.  A nonzero
multiple of \(\Eq_4\) uses at most two copies of \(F_s\): one for
\(k=0,4\); for \(k=2\), expose the two
\(s=0\) ports of the first and the two \(s=1\) ports of the second, pairing
the remaining ports; for \(k=1\), expose the three majority-zero ports of
the first and the minority-one port of the second, join the first's
minority port to one majority port of the second, and equality-close the
second's last two majority ports.  The \(k=3\) wiring is complementary.
For \(k=0,4\), the single copy is \(\Eq_4\).  In each two-copy wiring,
an all-equal external word has exactly one compatible internal assignment,
while every other external word has zero contribution.  The resulting
signature is \(\Eq_4\) for \(k=2\) and \(-\Eq_4\) for \(k=1,3\).
Pulling it back through \(R\) gives a
direct gadget \(g\) over \(M\) with \(R^{\otimes4}g=\lambda'\Eq_4\),
\(\lambda'\ne0\); together with \cref{eq:p1-eight-vertex-product-anchor-normal-form}
this is outcome~3.

If \(H\) is not tensor-prime, neither is \(M\).  Its nonzero endpoints make
every factor nonzero at both endpoints.  Even parity excludes an odd-sized
block, so the only proper tensor-factor partition is
\([4]=S\sqcup S^c\) with \(|S|=|S^c|=2\), where
\(S^c=[4]\setminus S\); the odd-weight zeros then force both binary factors
to be diagonal.  Thus they are nonsingular weighted equalities and
\[
 M(\zero^4)M(\one^4)=M(\one_S)M(\one_{S^c}).
\]
This is outcome~2 and completes the classification.
\end{proof}

\begin{corollary}[Reductions from product-transformable eight-vertex signatures]
\label{cor:p1-eight-vertex-product-operational}
Let \(\Gamma\) be a retained, factor-saturated finite algebraic signature
set that directly realizes a product-transformable endpoint-nondegenerate eight-vertex signature
\(M\).  Then at least one of the following holds:
\begin{enumerate}[label=\textnormal{(\roman*)}]
\item \(\KHolant(\Gamma)\) is \(\SharpP\)-hard;
\item \(\Tract(K\Gamma)\) holds;
\item there is a finite retained, factor-saturated augmentation
      \(\Lambda\supseteq\Gamma\) such that
      \[
        \KHolant(\Lambda)\leT\KHolant(\Gamma),
      \]
      and \(\Lambda\) retains both factors in the two-binary factorization
      of \(M\), including a nonsingular weighted equality exposed as a
      proper factor.  The factors are Turing-exposed by retained factor
      saturation, not asserted to be directly realizable over \(\Gamma\).
\end{enumerate}
\end{corollary}

\begin{proof}
Apply \cref{lem:p1-eight-vertex-product-transform}.  Its pure-generalized-equality
form is closed by \cref{lem:pure-ge}, and its dense form is closed by
\cref{lem:equality-anchor}.  These give alternatives~(i)--(ii), with
hardness composing to \(\Gamma\) and tractability descending to
\(\Gamma\) by \cref{lem:downward-inheritance}.

In the remaining form, adjoin the directly realized \(M\) and retain the
complete factor batch from its decomposition across two two-port blocks.  The
endpoint-nondegenerate eight-vertex signature is nonzero on both endpoint
words \(\zero^4\) and \(\one^4\), whose doubled charges are \(-4\) and
\(+4\), respectively, so the common-one-sided exception in
\cref{lem:factor-saturation} is impossible.  The resulting finite
factor-saturated augmentation \(\Lambda\) reduces to \(\Gamma\) and
contains both nonsingular weighted-equality factors.  This is
alternative~(iii).
\end{proof}

\begin{lemma}[Local-affine reduction to affine]
\label{lem:p1-eight-vertex-local-affine-collapse}
Let \(M\) be a nonzero quaternary even-parity signature.  If one matrix
\(T\in\GL_2(\mathbb C)\) satisfies
\[
  T^{\mathsf T}XT\in\cL,
  \qquad (T^{-1})^{\otimes4}M\in\cL,
\]
then there is a matrix \(T_0\in\GL_2(\mathbb C)\) with
\[
  T_0^{\mathsf T}XT_0\in\cA,
  \qquad (T_0^{-1})^{\otimes4}M\in\cA.
\]
If \(T\in\GL_2(\AlgNums)\), then \(T_0\) is algebraic.
\end{lemma}

\begin{proof}
For any nonempty affine set \(S\subseteq\Ftwo^n\), choose a base point
\(s\in S\) and write \(S=s+C\), where \(C\leq\Ftwo^n\) is a linear subspace.  The subspace

\[
  C=\{x\oplus y:x,y\in S\},
\]

is independent of \(s\).  It is self-orthogonal if
\(u\mathbin{\cdot}v=0\) over \(\Ftwo\) for all \(u,v\in C\).

Define the phase scalar \(\alpha=e^{\pi i/4}\).  For an arity-\(n\) signature
\(F\) and \(z,y\in\Ftwo^n\), define the local twist selected by \(z\) as
\[
  \Theta_z:=\bigotimes_{j=1}^n\diag(1,\alpha)^{z_j},
  \qquad
  (\Theta_zF)(y)
  =\alpha^{\langle z,y\rangle_{\mathbb Z}}F(y),
  \qquad
  \langle z,y\rangle_{\mathbb Z}:=\sum_{j=1}^n z_jy_j.
\]
By the local-affine criterion
\cite[Definition~3.1 and Theorem~3.2]{CaiLuXia2018RealHolantc}, the
translation space of every support is self-orthogonal, and
\(\Theta_zF\in\cA\) for every \(z\in\supp(F)\) when \(F\in\cL\).  Apply
this first to
the symmetric nonsingular binary \(E=T^{\mathsf T}XT\).  Self-orthogonality
rules
out full support, and nonsingularity rules out singleton support; hence
\(\supp(E)\) is equality or disequality.  Disequality support would give
\[
  E=pX\notin\cL,
\]
since \(01\in\supp(E)\) and the twist required by this support word has
relative phase \(\alpha\notin\mu_4\), so it is not affine.  Thus
\(E=\diag(p,q)\), where \(p,q\in\AlgNums^\times\) are its nonzero diagonal
entries (their algebraicity follows from \(E\in\cL\)); the \(00\)-twist is
the identity, so
\(q/p\in\mu_4\).  Writing
\(T=(\begin{smallmatrix}a&b\\c&d\end{smallmatrix})\), the off-diagonal
equation yields \(q/p=-(b/a)^2\).  Hence, with the scalar parameters
\(\rho:=b/a\in\mu_8\) and \(t:=c/a\in\mathbb C^\times\), and, for
\(z\in\mathbb C^\times\), the diagonal scaling matrix
\(D_z:=\diag(1,z)\),
\[
  T=aD_tPD_\rho,
  \quad at\ne0,\quad \rho^8=1.
\]
Define the matrix \(T_0:=aD_tP\), removing the right factor \(D_\rho\), and
define the function
\[
  W_t(y)=
  \sum_{\substack{x\in\Ftwo^4\\\wt(x)\equiv0\pmod 2}}
  t^{-\wt(x)}M(x)(-1)^{x\cdot y}.
\]
Let \(N:=(T^{-1})^{\otimes4}M\) denote the transformed local-affine
signature.  Since \(P^{-1}=P/2\), define
\[
  \kappa_0:=(2a)^{-4}\ne0,
  \qquad
  \widehat W_t:=\kappa_0W_t.
\]
\begin{equation}
  (T^{-1})^{\otimes4}M(y)=\rho^{-\wt(y)}\widehat W_t(y),
  \qquad
  (T_0^{-1})^{\otimes4}M(y)=\widehat W_t(y),
  \label{eq:p1-eight-vertex-local-affine-delete}
\end{equation}
Since \(N\in\cL\) and \(\rho\in\mu_8\),
\(\widehat W_t(y)=\rho^{\wt(y)}N(y)\) is algebraic for every \(y\).
Moreover,
\[
  T_0^{\mathsf T}XT_0
  =2a^2t\diag(1,-1)
  =p\diag(1,-1)\in\cA,
\]
and parity gives \(\widehat W_t(\bar y)=\widehat W_t(y)\).

The support of \(N\), which is also the support of
\(\widehat W_t\), is an affine coset \(s+C\).  Self-orthogonality and the
identity \(\widehat W_t(\bar y)=\widehat W_t(y)\), which makes this support
closed under bitwise complement, give
\(\one^4\in C\) and \(1\le\dim C\le2\).  For dimension one, the support is
a complementary pair on which \(\widehat W_t\) has equal values, so
\(\widehat W_t\in\cA\).
For dimension two, after permuting ports,
\[
  C=\Span_{\Ftwo}\{\one^4,w\},
  \qquad \wt(w)=2.
\]
For \(u,v\in\Ftwo\), index the four support words by
\(x_{uv}:=s\oplus u\one^4\oplus vw\).  For a table \(F\) that is nonzero on
these four words, define its cross-ratio functional by
\(\operatorname{cr}(F):=F(x_{00})F(x_{11})/
(F(x_{10})F(x_{01}))\).  Complement symmetry gives
\(\operatorname{cr}(\widehat W_t)=1\), while the nowhere-zero restriction of
any member of \(\cA\) with this support has cross ratio \(\pm1\).  The
quotient \(\Ftwo^4/C\) has exactly three nonzero classes.  Up to a port
permutation and \(w\leftrightarrow\bar w\), each of the two odd cosets has
representative \((s,w)=(0001,0011)\), and the unique even coset has
representative \((s,w)=(0101,0011)\).  Thus the displayed cases exhaust
all three nonzero quotient classes.  In both representatives
the \(\rho^{-\wt(y)}\) factor has cross ratio one and the direct four-point
check gives
\[
 \langle s,x_{00}\rangle_{\mathbb Z}
 +\langle s,x_{11}\rangle_{\mathbb Z}
 -\langle s,x_{10}\rangle_{\mathbb Z}
 -\langle s,x_{01}\rangle_{\mathbb Z}=2.
\]
The local twist selected by \(s\) is \(N_s:=\Theta_sN\), so
\(N_s(y)=\alpha^{\langle s,y\rangle_{\mathbb Z}}N(y)\).  It would therefore
have cross ratio \(\alpha^2=i\), contradicting \(N_s\in\cA\).  Hence
\(s+C=C\).

If the complement-symmetric values of \(\widehat W_t\) are \(A\) on the endpoints
and \(B\) on the central pair, then \(N\) has values
\((A,\rho^{-4}A,\rho^{-2}B,\rho^{-2}B)\).  Since \(0\in\supp(N)\), its
selected local twist is the identity, so \(N\in\cA\) and
\(\rho^{-2}B/A\in\mu_4\).  As \(\rho^2\in\mu_4\), also \(B/A\in\mu_4\),
whence \(\widehat W_t\in\cA\).  Now
\cref{eq:p1-eight-vertex-local-affine-delete} proves the
claim.  If the original witness is algebraic, so are \(t\) and \(T_0\).
\end{proof}

\subsection{The finite affine quaternary certificate}
\label{proofsubsec:affine-certificate}

The affine branch uses the finite classification in
Appendix~\ref{appsec:cert-eight-vertex}, checked by
\path{certificates/verify_affine_q4.py}.  The next proposition derives the
complete list of possible basis matrices algebraically and invokes that
classification for the resulting quaternary tables.

\begin{proposition}[Affine quaternary certificate]
\label{prop:p1-affine-eight-vertex-certificate}
Let \(M\) be an algebraic affine-transformable
endpoint-nondegenerate eight-vertex signature.  Then \(M\) admits an
algebraic invertible diagonal matrix
\(D_{\mathrm{aff}}=\diag(d_0,d_1)\), with
\(d_0,d_1\in\AlgNums^\times\), such that
\[
  (D_{\mathrm{aff}}^{-1})^{\otimes4}M\in\cA,
  \qquad D_{\mathrm{aff}}^{\mathsf T}XD_{\mathrm{aff}}\doteq X.
\]
For some \(J\subseteq[4]\) with \(|J|=2\), write
\(J^c:=[4]\setminus J\).  In these coordinates its support is exactly one
of
\[
  \{\zero^4,\one^4\},
  \qquad
  \{\zero^4,\one^4,\one_J,\one_{J^c}\},
  \qquad
  \{x:\wt(x)\text{ is even}\}.
\]
The first support is a pure generalized equality.
\end{proposition}

\begin{proof}
By \cref{lem:algebraic-transform-witness}, choose an algebraic affine witness
matrix
\(U=(\begin{smallmatrix}a&b\\c&d\end{smallmatrix})\), and define the
transformed binary edge \(E:=U^{\mathsf T}XU\).  A nonsingular symmetric affine binary has
disequality support, equality support, or full support.  In the first
case,
\[
  U=D_0K_0,\qquad K_0\in\{I,X\},
\]
for an invertible diagonal \(D_0\).
In the second, \(ad+bc=0\), so
\[
 U=D_0\begin{pmatrix}1&r\\1&-r\end{pmatrix},
 \qquad r\in\mu_8,
\]
for an invertible diagonal \(D_0\).  In the full-support case, after
normalizing its \(00\)-entry, \(E\) has the affine form
\[
 \begin{pmatrix}1&u\\u&-u^2\end{pmatrix},
 \qquad u\in\mu_4.
\]
Writing \(U=D_0(\begin{smallmatrix}1&r\\1&s\end{smallmatrix})\) gives
\(r+s=2u\) and \(rs=-u^2\); hence
\(\{r,s\}=\{u(1+\sqrt2),u(1-\sqrt2)\}\).
Denote the resulting eighteen-element set of cores by
\(\mathscr K_{\mathrm{aff}}\); its explicit list appears in
\cref{thm:cert-interface-affine-q4-normal-form}.  These three calculations give precisely
\[
 U=D_0K_0,\qquad K_0\in\mathscr K_{\mathrm{aff}},
\]
so completeness of the list is analytic rather than computer-assumed.

Now define the affine signature \(f=(U^{-1})^{\otimes4}M\in\cA\).  Since
\(M\) is an endpoint-nondegenerate eight-vertex signature and \(D_0\)
is diagonal, \(K_0^{\otimes4}f\) satisfies the hypotheses of
\cref{thm:cert-interface-affine-q4-normal-form}.  The certificate is projective, so the
normalization factor for \(f\) can be absorbed into a global scalar.
The same interface theorem therefore yields
\[
  \gamma,t\in\AlgNums^\times
\]
and a normalized affine signature \(h\in\cA\), with
\(h(\zero^4)=1\), such that
\[
 K_0^{\otimes4}f
 =\gamma\diag(1,t)^{\otimes4}h.
\]
Moreover, \(\supp(h)\) is one of the three supports displayed above.
Consequently the
single algebraic diagonal matrix \(D_\star=D_0\diag(1,t)\) satisfies
\((D_\star^{-1})^{\otimes4}M\doteq h\), and
\(D_\star^{\mathsf T}XD_\star\doteq X\).  Taking
\(D_{\mathrm{aff}}=D_\star\) proves the asserted diagonal
normalization.

\end{proof}

\begin{corollary}[Reductions from affine quaternaries]
\label{cor:p1-affine-eight-vertex-operational}
Let \(\Gamma\) be a retained, factor-saturated finite algebraic signature
set that directly realizes an affine-transformable endpoint-nondegenerate eight-vertex signature
\(M\).
Then at least one of the following holds:
\begin{enumerate}[label=\textnormal{(\roman*)}]
\item \(\KHolant(\Gamma)\) is \(\SharpP\)-hard;
\item \(\Tract(K\Gamma)\) holds;
\item there is a finite retained, factor-saturated augmentation
      \(\Lambda\supseteq\Gamma\) such that
      \[
        \KHolant(\Lambda)\leT\KHolant(\Gamma),
      \]
      and \(\Lambda\) retains the complete factor batch from the physical
      construction.  The original endpoint-nondegenerate eight-vertex signature \(M\) itself has, up to a port
      permutation and a nonzero scalar, a genuine factorization across two
      two-port blocks,
      and \(\Lambda\) contains a nonsingular weighted-equality matrix
      \(E:=\diag(p,q)\), where \(p,q\in\AlgNums^\times\), exposed as one of its proper binary
      factors.  This factor is Turing-exposed by retained factor saturation,
      not asserted to be directly realizable over \(\Gamma\).
\end{enumerate}
\end{corollary}

\begin{proof}
Choose \(D_\star=D_{\mathrm{aff}}\) from
\cref{prop:p1-affine-eight-vertex-certificate}.  Apply the global change
\(D_\star^{-1}\), and normalize the resulting quaternary to an affine
signature \(h\) with \(h(\zero^4)=1\).  Thus
\[
 M\doteq D_\star^{\otimes4}h.
\]
Keep every direct output over the current retained set, and adjoin proper
tensor factors only through retained factor saturation.  Record these
constructions and the chosen normal form in the normalization data of
\cref{lem:normalization-ledger}.  The transported copy of \(M\) remains
nonzero at both endpoint words, so every augmentation is two-sided and
\cref{lem:factor-saturation} applies.

The normal form has three possible supports.  Clause~\textnormal{(ii)} of
\cref{thm:cert-interface-affine-q4-gadget-exits} gives the
support-four construction.  The support-two form is a pure generalized
equality.  On support four, let \(J\sqcup J^c=[4]\) be the two two-port
blocks.  For \(u,v\in\Ftwo\), put
\(x_{uv}:=u\one_J\oplus v\one_{J^c}\), and define the logical table by
\[
 F:=\bigl(h(x_{uv})\bigr)_{u,v\in\Ftwo}
   =\begin{pmatrix}A&B\\ C&D\end{pmatrix},\qquad ABCD\ne0.
\]
Here \(u\) is the row bit and \(v\) is the column bit.
For an ordered binary tensor \(E\), write \(E^{(J)}\) for its placement on
the two ports of \(J\) in their inherited order, and define
\(E^{(J^c)}\) analogously.
If \(AD=-BC\), the two-copy \(X\)-edge gadget in
\cref{thm:cert-interface-affine-q4-gadget-exits}\textnormal{(ii)} directly realizes a nonzero pure
generalized equality.  If \(AD=BC\), then, on these two blocks,
\[
  h=E_1^{(J)}\otimes E_2^{(J^c)}
\]
for nonsingular weighted equalities \(E_1,E_2\).  Since \(D_\star\) is
diagonal,
\[
  M\doteq
  (D_\star E_1D_\star^{\mathsf T})^{(J)}\otimes
  (D_\star E_2D_\star^{\mathsf T})^{(J^c)}.
\]
Thus the original endpoint-nondegenerate eight-vertex signature itself has a
genuine factorization across two two-port blocks into nonsingular weighted
equalities.  Adjoin \(M\) and apply retained factor saturation to
Turing-expose both factors.

On full even support,
\cref{thm:cert-interface-affine-q4-gadget-exits}\textnormal{(iii)}, together with
\cref{lem:p1-eight-vertex-product-transform}, gives, over the current retained set,
either a directly realized equality anchor, a directly realized nonzero
rank-one binary, or a directly realized product
\[
  cE_r\otimes E_r,
  \qquad c\in\AlgNums^\times,
  \qquad E_r:=\diag(1,r),\quad r\in\mu_4,
\]
where \(E_r\) is a nonsingular weighted-equality matrix.
In the last case, \cref{lem:factor-saturation} Turing-exposes \(E_r\).  Choose
the scalar \(s\in\mu_8\) with \(s^2=r^{-1}\), and define the diagonal
equality-normalization matrix \(D_r^{\mathrm{eq}}:=\diag(1,s)\).  Then
\[
  (D_r^{\mathrm{eq}})^{\mathsf T}XD_r^{\mathrm{eq}}=sX,
  \qquad
  D_r^{\mathrm{eq}}E_r(D_r^{\mathrm{eq}})^{\mathsf T}=I.
\]
Apply this global normalization and record its matrix and scalar.  Define
the transformed signature \(h_r=(D_r^{\mathrm{eq}})^{\otimes4}h\).  Attach one port of
\(h_r\) by a native \(X\)-edge
to one port of the retained equality \(I\), leaving the other equality port
external.  This is an actual gadget over the current retained augmentation,
and it complements that input of \(h_r\), producing, up to a known nonzero
scalar, an odd-parity quaternary \(p\) that is nonzero on all eight
odd-weight words and is used in
\cref{thm:cert-interface-affine-q4-odd}.  Because \(s\in\mu_8\), its normalized
nonzero entries are fourth roots of unity.

The twelve native and equality-mediated binary gadgets of
\cref{thm:cert-interface-affine-q4-odd} are therefore actual at this node.  They
expose either a nonzero rank-one binary, a nonsingular binary of infinite
projective order, or a table for which neither of these two outcomes occurs
in the twelve-loop test of \cref{thm:cert-interface-affine-q4-odd}.
This includes the simultaneous-zero case of the theorem's loop tests.
The exceptional remaining family directly
realizes the certified endpoint-nondegenerate eight-vertex signature with
support of cardinality six.  For every other remaining table, choose a certified
\(\delta\in\mathbb Z_4\), choose \(\tau_\delta\in\mu_8\) with
\(\tau_\delta^2=i^\delta\), and apply the further legal global-normalization
matrix \(T_\delta:=\diag(1,\tau_\delta)\).  Let \(p_\delta\) denote the certified
phase-shifted odd table from
\cref{thm:cert-interface-affine-q4-odd}.  Since \(p\) is supported only at
Hamming weights one and three,
\[
  T_\delta^{\otimes4}p\doteq p_\delta.
\]
Record this global change in the normalization data.
\Cref{thm:cert-interface-affine-q4-odd} now gives, through its literal native-loop
and two-copy constructions, either a nonzero rank-one binary or a
product-transformable endpoint-nondegenerate eight-vertex quaternary in the
dense form \cref{eq:p1-eight-vertex-product-dense-form}; the latter directly
yields an equality anchor by
\cref{lem:p1-eight-vertex-product-transform}.

The infinite-order case is handled by \cref{lem:binary-interpolation}.  The
eight-vertex signature with support of cardinality six is outside the
\(\cA\)-, \(\cP\)-, and \(\cL\)-transformable classes by
\cref{thm:cert-interface-affine-q4-odd,lem:p1-eight-vertex-product-transform,%
lem:p1-eight-vertex-local-affine-collapse}, and is hard by
\cref{thm:external-eight-vertex}.  Thus every direct gadget is formed over
the current retained set, every proper factor is only Turing-exposed, and
all global diagonal changes belong to one normalization-record history.

The support-two case and the support-four case \(AD=-BC\) are resolved by
\cref{lem:pure-ge}; equality anchors are resolved by
\cref{lem:equality-anchor}.  A nonzero rank-one binary exposes a unary by
\cref{lem:binary-interpolation}, and an infinite-order binary first
interpolates such a binary.  In both cases
\cref{lem:common-unary-exit} applies.  An exposed odd-arity signature is
resolved by \cref{thm:external-odd} after the \(K\)-coordinate bridge.
Every full-support case above therefore gives hardness or tractability.

The sole continuing case is the support-four form \(AD=BC\).  Its
displayed factorization gives two nonsingular weighted-equality factors of
the original \(M\), after undoing the diagonal normalization.  Retain both
factors and the entire finite factor-saturated augmentation as \(\Lambda\).
By \cref{lem:normalization-ledger}, all hardness reductions compose back
to \(\Gamma\), all tractable conclusions descend to \(K\Gamma\), and the
continuing state returns to the caller's coordinates.  This proves the
three alternatives.
\end{proof}

\subsection{Consequences for the Full Signature Set}

The preceding reductions give the following interface for the entire
signature set.  The factor outcome is retained for the later
equality-accessible analysis.

\begin{corollary}[Reductions from eight-vertex signatures]
\label{cor:p1-eight-vertex-reduced-interface}
Let \(\cG\) be a retained, factor-saturated finite algebraic
signature set containing a directly realizable endpoint-nondegenerate eight-vertex signature \(M\).
Then at least one of the following holds:
\begin{enumerate}[label=\textnormal{(\roman*)}]
\item \(\KHolant(\cG)\) is \(\SharpP\)-hard;
\item \(\Tract(K\cG)\) holds;
\item for some finite retained, factor-saturated augmentation
      \(\Lambda\supseteq\cG\) satisfying
      \(\KHolant(\Lambda)\leT\KHolant(\cG)\), the signature set \(\Lambda\)
      retains the complete factor batch from a genuine factorization of the
      original endpoint-nondegenerate eight-vertex signature \(M\) across two
      two-port blocks and contains the
      nonsingular weighted-equality matrix
      \[
        E:=\diag(p,q),\qquad p,q\in\AlgNums^\times,
      \]
      Turing-exposed as a proper tensor factor.
\end{enumerate}
In item~(iii), the factor \(E\) is Turing-exposed from the original \(M\);
secondary two-copy equality outputs are resolved within the preceding
reductions.  The construction uses only this quaternary analysis.
\end{corollary}

\begin{proof}
If all six weight-two entries vanish, the signature is a pure generalized
equality and is closed by
\cref{lem:pure-ge}.  Otherwise apply the Cai--Fu quaternary boundary
\cref{thm:external-eight-vertex}.  Its hard branch is item~(i).  Its
product-transformable branch is routed through
\cref{cor:p1-eight-vertex-product-operational}.  A local-affine branch first
collapses to affine by \cref{lem:p1-eight-vertex-local-affine-collapse} and then
uses \cref{cor:p1-affine-eight-vertex-operational}; an affine branch uses that operational
corollary directly.

The two operational wrappers internally close every pure generalized
equality, equality anchor, odd-arity signature, and singular binary factor.
In the last case, retained factor saturation exposes a unary by
\cref{lem:factor-saturation}; the corresponding transported diagonal dressing
of the original endpoint-nondegenerate eight-vertex signature remains present,
excludes the common-one-sided exception, and the
common unary exit \cref{lem:common-unary-exit} applies.  Consequently all routes merge into exactly the
three displayed outcomes, and every hardness reduction composes toward
\(\cG\), while every tractable presentation descends to \(K\cG\).  Apply
\cref{lem:normalization-ledger} once to pull back the legal normalizations.
The sole continuing output is the retained nonsingular weighted-equality
factor in item~(iii).
\end{proof}

\begin{proof}[Proof of \cref{thm:p1-quaternary-kholant-dichotomy}]
Choose an endpoint-nondegenerate tensor-prime quaternary
\(q\in\cF\).  Preprocess as in \cref{subsec:zero-nullary} and write
\(\cF^\circ\) for the resulting signature set.  The signature \(q\) survives,
and the preprocessing convention identifies the displayed criterion with
\(\Tract(K(\cF^\circ))\).

If \(\Tract(K(\cF^\circ))\) holds, then
\cref{thm:predicate-algorithms,eq:K-equivalence} gives
\(\KHolant(\cF^\circ)\in\FP\), and the root evaluator restores the deleted
nullary scalars.  Now suppose that \(\Tract(K(\cF^\circ))\) fails.  The two
nonzero endpoint values of \(q\) occur at doubled charges \(-4\) and \(+4\),
so \(\cF^\circ\) is not common one-sided.  Apply
\cref{lem:factor-saturation} and retain
\[
  \Gamma:=\cF^\circ\cup\Pi(\cF^\circ).
\]
Then \(\Gamma\) is finite and factor-saturated,
\(\KHolant(\Gamma)\equivT\KHolant(\cF^\circ)\), and it still contains the
original signature \(q\).  By the contrapositive of
\cref{lem:downward-inheritance}, \(\Tract(K(\Gamma))\) also fails.

Apply \cref{cor:p1-eight-vertex-reduced-interface} to \(q\) over
\(\Gamma\).  Its tractable outcome is excluded by the preceding paragraph.
Its continuing outcome gives a genuine \(2+2\) factorization of the original
\(q\), up to the recorded invertible diagonal dressing, port permutation,
and nonzero scalar.  An invertible portwise dressing preserves
tensor-primality: undoing it would give a positive-arity tensor factorization
of \(q\).  Hence the continuing outcome contradicts the hypothesis, and the
only remaining outcome is that \(\KHolant(\Gamma)\) is \(\SharpP\)-hard.
The retained equivalence and preprocessing transfer this hardness to
\(\KHolant(\cF)\).  Finally, decidability follows from
\cref{thm:predicate-algorithms}.
\end{proof}
\section{Equality-Accessible Matching Decks}
\label{sec:weighted-equality}
\label{sec:p1-binary-phase}

Here equality access means that the literal binary \(I=\Eq_2\) is
available.  We first control all realizable binary transfers, then contract
a minimum-arity nonbinary prime and retain the factors of its outputs.
Stabilization makes these outputs products on perfect matchings.  A finite
group classification then routes the remaining tensor to one of
\cref{sec:normalized-dihedral,sec:proper-deck,sec:v4-deck,sec:platonic-decks}.

The edge and transfer conventions are those of
\cref{subsec:holographic,subsec:reduction-conventions}; the termination
argument uses the retained-factor potential from
\cref{lem:factor-saturation}.  The shared finite quaternary statement is
in Appendix~\ref{appsec:cert-equality-accessible}.
Throughout, the supplied \(I\) remains in the caller's fixed coordinates.
A call to \cref{sec:p1-eight-vertex} returns in those coordinates; any
additional weighted equality is processed as a binary signature.

A \emph{resolved leaf} is hardness, \(\Tract\) for the entire retained set,
or a common tractable \(K\)-presentation in the sense of
\cref{def:common-K-presentation}.

\begin{lemma}[Equality-wire equivalence]
\label{lem:p1-equality-wire-equivalence}
For every signature set \(\cG\),
\begin{equation*}
  \KHolant(\cG,I)
  \equivT
  \Holant(\cG,X),
\end{equation*}
where the right side uses ordinary equality wires and has \(X\) as an
available binary signature.
\end{lemma}

\begin{proof}
Forward, subdivide each native \(X\)-wire by an \(X\)-vertex and suppress
degree-two \(I\)-vertices as equality wires.  Conversely, splice the two
incident ports of the ordinary \(X\)-vertices after first preprocessing the
subgraph induced by those vertices.  Every such component is a path or a
cycle, because \(X\) has arity two.  A path with two incidences at
non-\(X\) vertices collapses to \(I\) or \(X\) according to the parity of
the number of \(X\)-vertices; a cycle with no external incidence contributes
zero for odd length and the scalar \(2\) for even length.  (A self-loop is
the length-one case.)  These are exact identities obtained by summing the
internal equality-wire bits, and they also cover adjacent \(X\)-vertices and
parallel edges.  For each resulting \(X\) kernel use one structural
\(\mathsf E_X\) edge (no additional tensor vertex); for each resulting
\(I\) kernel, and for each ordinary equality wire whose endpoints are not
in an \(X\)-component, use
\[
  \text{port}-X-I-X-\text{port}.
\]
Indeed, \(\sum_{u,v}X(a,u)I(u,v)X(v,b)=I(a,b)\).  Multiply the final
answer by the known product of the even-cycle scalars (or return zero for
an odd cycle).  The preprocessing and substitutions are linear-size and
preserve ordered incidences.
\end{proof}

Consequently, it suffices to work with a retained factor-saturated signature
set containing the distinguished literal equality \(I\).  We first establish
the binary-transfer and finite-group control required by the matching-deck
construction.

\subsection{Binary transfers, complete groups, and local quaternary dispatch}
\label{subsec:binary-transfer-control}
Binary chains turn composition into matrix multiplication.  Bounding the
group they generate gives a finite set of contraction kernels for the
matching deck.

For an ordered binary signature \(B\), recall its transfer
\[
  T_B=BX.
\]
Joining \(B\) to another ordered binary signature \(C\) in series through
one \(X\)-edge realizes \(BXC\),
and therefore
\[
  T_{BXC}=T_BT_C.
\]
For a transfer matrix \(T\), define its reversal by the anti-involution
\begin{equation}
  T^\#:=XT^{\mathsf T}X;
  \qquad
  T^\#=(\det T)ZT^{-1}Z\quad\text{if }T\text{ is nonsingular}.
  \label{eq:transfer-reversal}
\end{equation}
Here ``anti-involution'' means explicitly that
\[
  (ST)^\#=T^\#S^\#,
  \qquad
  (T^\#)^\#=T
\]
for all transfer matrices \(S,T\).

\begin{lemma}[Bounded torsion]
\label{lem:bounded-torsion}
Fix a number field \(L\), and let \(G\) be the set of projective
transfers of all directly realizable nonzero nonsingular ordered
binary gadgets over a fixed signature set with entries in \(L\).
If \(G\ne\varnothing\) and every element of \(G\) has finite
projective order, then \(G\) is a finite subgroup of
\(\PGL_2(L)\).
\end{lemma}

\begin{proof}
Define the field degree \(D=[L:\mathbb Q]\).  Let \([T]\in G\) have
projective order
\(r\).  By definition,
\[
  T^r=\eta I
\]
for some \(\eta\in L^\times\), and \(r\) is the least positive
integer with this property.

The minimal polynomial \(m_T\) divides \(x^r-\eta\), which is square-free
because \(\gcd(x^r-\eta,rx^{r-1})=1\) in characteristic zero.
Thus \(T\) is diagonalizable over the algebraic closure.  Choose
an eigenbasis matrix \(V\in\GL_2(\AlgNums)\) and eigenvalues
\(\lambda_1,\lambda_2\in\AlgNums^\times\) such that
\[
  V^{-1}TV=\diag(\lambda_1,\lambda_2).
\]
Define the eigenvalue ratio
\[
  \zeta:=\frac{\lambda_1}{\lambda_2}.
\]
The equality \(T^r=\eta I\) gives \(\zeta^r=1\), and the minimality of
\(r\) shows that \(\zeta\) is a primitive \(r\)-th root of unity.
Moreover,
\[
  \zeta+\zeta^{-1}
  =
  \frac{(\operatorname{tr}T)^2}{\det T}-2
  \in L.
\]
For \(r>2\),
\[
  [\mathbb Q(\zeta+\zeta^{-1}):\mathbb Q]
  =\frac{\varphi(r)}2,
\]
so \(\varphi(r)\le2D\), also valid for \(r=1,2\).  The integers with this
bound form an effectively finite set: each prime power \(p^a\mid r\)
satisfies \(p^{a-1}(p-1)\le2D\), bounding both \(p\) and \(a\).
Consequently all projective orders divide the computable exponent
\[
  e_D=\operatorname{lcm}\{r:\varphi(r)\le2D\}.
\]

Series composition gives product closure.  A positive power of every
element is the identity, so the identity and inverses also belong to
\(G\).  Hence \(G\) is a subgroup.  Put
\[
 \mathfrak{sl}_2(L)
 :=\{A\in\Mat_2(L):\operatorname{tr}A=0\}.
\]
The adjoint representation is
\[
  \operatorname{Ad}:\PGL_2(L)\hookrightarrow
  \GL(\mathfrak{sl}_2(L))\cong\GL_3(L),
  \qquad
  \operatorname{Ad}([T])(A)=TAT^{-1}.
\]
Its kernel is trivial: a matrix commuting with all trace-zero matrices
is scalar.  Burnside's finite-exponent theorem states that a subgroup
\(H\le\GL_n(\mathbb C)\) satisfying \(h^e=I\) for one fixed \(e\) and
every \(h\in H\) is finite \cite{Burnside1905}.  Apply it to
\(\operatorname{Ad}(G)\le\GL_3(\mathbb C)\) with \(e=e_D\).
Thus all binary gadgets, despite their unbounded sizes, have only finitely
many projective transfer classes.
\end{proof}

For a finite algebraic retained, factor-saturated, two-sided signature set
\(\cG\), call a nonzero directly realizable ordered binary \emph{safe} if it
is nonsingular and its transfer has finite projective order.  Call it
\emph{unsafe} if it is singular or its transfer has infinite projective
order.  Call the set \emph{binary-residual} when every nonzero actual ordered
binary is safe.  We work in this case; otherwise
\cref{lem:binary-interpolation,lem:common-unary-exit} resolves the branch.  Series composition and
\cref{lem:bounded-torsion} then give a finite group: take \(L\) to be the
number field generated by the entries of \(\cG\), which also contains the
entries of every direct gadget transfer.
The group is
\begin{equation}
  G(\cG)
  =
  \{[BX]:B\text{ is a safe ordered binary directly realizable over }\cG\}.
  \label{eq:complete-safe-group}
\end{equation}
The definition uses all directly realizable binaries.  Taking or recomputing
this complete group below is a structural case distinction for a fixed
signature set.  The decision algorithm for the resulting tractability
alternatives is given by \cref{thm:predicate-algorithms}; it does not
enumerate gadgets of unbounded size.  Since the group is finite,
fix an oriented actual representative \(B_g\) satisfying
\([B_gX]=g\) for every \(g\in G(\cG)\).

\begin{lemma}[Relative completeness test for actual binaries]
\label{lem:actual-binary-completeness}
Let \(\Gamma\subseteq\Lambda\) be finite algebraic retained,
factor-saturated signature sets with \(I\in\Gamma\), and suppose that
every nonzero binary directly realizable over \(\Gamma\) is safe.  Put
\(H=G(\Gamma)\).  The following exhaustive priority dispatch applies:
\begin{enumerate}
\item if \(\Lambda\) directly realizes a nonzero singular ordered binary,
      take the factor/unary exit;
\item otherwise, if \(\Lambda\) directly realizes a nonsingular ordered
      binary of infinite projective transfer order, take the interpolation
      exit;
\item otherwise every nonzero binary directly realizable over \(\Lambda\) is safe, \(H\le G(\Lambda)\), and
      every nonsingular actual binary \(B\) over \(\Lambda\) with
      \([BX]\notin H\) witnesses the strict inclusion \(H<G(\Lambda)\).
\end{enumerate}
In the third case, if \(G(\Lambda)=H\), every nonzero actual binary over
\(\Lambda\) is projectively one of the fixed representatives \(B_h\),
\(h\in H\).
\end{lemma}

\begin{proof}
Retained \(I\) has both charge signs, so both signature sets are two-sided.
A singular binary is handled by
\cref{lem:binary-interpolation,lem:common-unary-exit}, and an infinite-order
binary by \cref{lem:binary-interpolation}.  If neither exists, \(\Lambda\)
is binary-residual, so \cref{lem:bounded-torsion} defines \(G(\Lambda)\).
Every gadget over \(\Gamma\) remains a gadget over \(\Lambda\), hence
\(H\le G(\Lambda)\).  The remaining assertions follow directly from the
definitions and the fixed representative choice.
\end{proof}

\begin{remark}[Relation to binary-signature classifications]
The singular, nonsingular infinite-projective-order, and nonsingular
finite-projective-order split is the \(\KHolant\) form of the binary split
used by Cai--Lu--Xia \cite{CaiLuXia2009,CaiLuXia2018RealHolantc}.  The four
ordered Bell signatures
\(\mathcal B=\{B_{0,0},B_{0,1},B_{1,0},B_{1,1}\}\) represent the projective
Bell classes of Shao--Cai \cite{ShaoCai2020}.  Their holographic matrix,
denoted \(Z\) in that paper, is our \(K\); throughout this paper,
\(Z=\left(\begin{smallmatrix}1&0\\0&-1\end{smallmatrix}\right)\), as in
\cref{eq:transfer-reversal}.  Beyond the binary classification, we need the
complete transfer group and uniform control of its higher-arity cards.
\end{remark}

For an even finite port set \(R\), a \emph{perfect matching} is a partition
of \(R\) into two-element subsets.  An \emph{oriented perfect matching}
also chooses an order within each pair, and a matching product places one
ordered binary factor on every oriented pair.

For a quaternary tensor \(q\) with ordered ports \(1,2,3,4\), let
\(F_{ij\mid k\ell}(q)\) denote the \(4\times4\) flattening whose row index is
\((x_i,x_j)\) and whose column index is \((x_k,x_\ell)\).  For a finite
complete transfer group \(G\) and its fixed actual representatives
\(B_g\), \(g\in G\), define
\[
 \mathcal M_4(G)=
 \left\{
 c\,B_g(x_i,x_j)B_h(x_k,x_\ell):
 \begin{array}{l}
 c\in\AlgNums^\times,\ ((i,j),(k,\ell))\text{ is an oriented}\\[-2pt]
 \text{perfect matching of }[4],\quad g,h\in G
 \end{array}
 \right\}.
\]
Thus \(\mathcal M_4(G)\) is the arity-four \(G\)-matching notion used
below.  Define the
complete oriented one-copy binary batch as follows.  For an ordered binary
card \(C\), set \(C^{\mathrm{id}}=C\) and
\(C^{\mathrm{rev}}=C^{\mathsf T}\), and put
\[
 \mathscr D_G^{(1)}(q)
 =
 \{0\}\cup
 \left\{
   \bigl(\partial_{(i,j)}^{XB_gX}q\bigr)^\epsilon:
   i,j\in[4],\ i\ne j,\ g\in G,\
   \epsilon\in\{\mathrm{id},\mathrm{rev}\}
 \right\}.
\]
The order \((i,j)\) records the orientation of the deleted pair, and the
choice of \(\epsilon\) records either orientation of the residual pair.
Every displayed member is a direct gadget: \(XB_gX\) is the physical kernel
obtained by attaching
\(B_g\) through its two native \(X\)-wires, and residual reversal is a port
permutation.  In this local case analysis, a
\emph{nonbinary prime} means a tensor-prime signature of arity at least
three.

\begin{lemma}[Binary contractions of a quaternary signature]
\label{lem:finite-q4-dispatch}
Let \(\Lambda\) be a finite algebraic retained, factor-saturated,
  binary-residual signature set containing the distinguished literal equality
\(I\), suppose every retained nonbinary prime has even arity, and put
\(G=G(\Lambda)\).  Let \(q\) be an actual signature over \(\Lambda\), with
its ordered ports, realization, and nonzero scalar recorded.  Assume that
\(q\) is endpoint-nondegenerate, tensor-prime, and
\[
 q\notin\mathcal M_4(G).
\]
The three physical flattenings satisfy
\[
 \rank F_{12\mid34}(q)\ge2,\qquad
 \rank F_{13\mid24}(q)\ge2,\qquad
 \rank F_{14\mid23}(q)\ge2.
\]
These inequalities follow from tensor primality and specify the rank
checks used by later callers.

Let \(\Lambda_q\) be the preprocessed literal augmentation obtained by
adjoining \(q\) and all nonzero tensor cards in
\(\mathscr D_G^{(1)}(q)\), followed by retained factor saturation of this
finite batch.  Then
\[
 \KHolant(\Lambda_q)\equivT\KHolant(\Lambda),
\]
and one of the following holds:
\begin{enumerate}[label=\textnormal{(\roman*)}]
\item an established resolved leaf or terminal applies, including an
      odd-arity/unary exit, a nonzero singular-binary exit, or an actual
      infinite-order-binary interpolation exit; alternatively, a retained
      lower-arity nonbinary prime gives the corresponding outer deck
      successor;
\item \(\Lambda_q\) is binary-residual and
      \[
        G<G_q:=G(\Lambda_q),
      \]
      with the strict inclusion witnessed by an ordered binary \(B\) actual
      over \(\Lambda_q\) such that \([BX]\in G_q\setminus G\);
\item \(\Lambda_q\) is binary-residual, \(G_q=G\), \(q\) is a
      minimum-arity nonbinary tensor-prime in \(\Lambda_q\), and every member
      of \(\mathscr D_G^{(1)}(q)\) is zero or projectively equal to one of the
      fixed \(B_g\), \(g\in G\).  All nonzero card factors are retained; on
      this continuing branch every nonzero card has a trivial factorization tree.
\end{enumerate}
Only item~(iii) is passed to the equality-accessible quaternary boundary.

Moreover, on branch~(iii), let \(r\) be an auxiliary actual quaternary
signature over \(\Lambda_q\), assume \(r\notin\mathcal M_4(G)\), and suppose
that \(r\) has a genuine \(2+2\) factorization.  Let
\(\Lambda_{q,r}\) be the preprocessed retained, factor-saturated augmentation
obtained by adjoining \(r\) and its nonzero factor batch.  Then either an
established terminal applies or
\[
  G<G(\Lambda_{q,r}).
\]
Thus there is no same-group continuing branch.  This finite local case analysis
does not invoke the global matching-synthesis measure.
\end{lemma}

\begin{proof}
The batch \(\mathscr D_G^{(1)}(q)\) is finite: there are six unordered
deleted pairs, two orders for the closing kernel, two residual orders, and
finitely many closing-representative choices \(B_g\).  Its tensor cards are
actual over \(\Lambda\), since \(q\) and every \(B_g\) have recorded direct
realizations.  Adjoining them therefore gives one Turing reduction by
\cref{eq:direct-gadget-reduction}, while inclusion gives the reverse one.
Because literal \(I\) makes every retained augmentation two-sided,
\cref{lem:factor-saturation} may retain the finite factorization tree without
changing the problem up to Turing equivalence.

Dispatch every odd-arity or unary factor and every singular binary factor.
If \(\Lambda_q\) directly realizes any nonsingular binary of infinite
projective transfer order, apply \cref{lem:binary-interpolation}; this is
item~(i).  Otherwise \(\Lambda_q\) is binary-residual, so
\cref{lem:bounded-torsion} defines the finite group
\(G_q=G(\Lambda_q)\).  Retained inclusion gives \(G\le G_q\).  If the
inclusion is strict, the definition of \(G_q\) supplies an actual binary
representative of a class in \(G_q\setminus G\), proving item~(ii).
The witness is an arbitrary actual gadget over \(\Lambda_q\), possibly
using several copies of \(q\).

It remains that \(G_q=G\).  Every nonzero card in
\(\mathscr D_G^{(1)}(q)\) is then nonsingular and has transfer class in
\(G_q=G\), hence is projectively one of the fixed \(B_g\).  Such a binary
has no proper positive-arity tensor factorization.  The tensor \(q\) remains
prime intrinsically.  Any nonbinary prime of arity below four would have
arity three and would already have triggered the odd-arity terminal, so
\(q\) has minimum nonbinary-prime arity.  This proves item~(iii).

Finally, write \(r\doteq C_1\otimes C_2\) for the auxiliary factorization.
Factor saturation makes \(C_1,C_2\) literal resources of
\(\Lambda_{q,r}\); they are not asserted to have been actual over
\(\Lambda_q\).  If either is singular, or if the augmentation realizes an
infinite-order binary, an established terminal applies.  Otherwise
\(\Lambda_{q,r}\) is binary-residual and
\(G\le G(\Lambda_{q,r})\).  If equality held, both factors would be
nonsingular and would represent classes in \(G\), whence
\(r\in\mathcal M_4(G)\), contradicting the stated hypothesis.  The inclusion
is therefore strict.
\end{proof}

We now choose a minimum-arity nonbinary prime and retain the factors
of all its dressed contractions.

\subsection{Stable matching decks and restart control}
\label{subsec:safe-matching-cards}

Fix a retained, factor-saturated signature set \(\cG\ni I\) and its complete transfer group \(G=G(\cG)\).  Deck operations use the actual binaries \(B_g\),
transfers \(R_g\), and multiplication/reversal scalars fixed in
\cref{subsec:reduction-conventions} before passing to projective classes.

\begin{definition}[Matching deck]
\label{def:matching-deck}
Let \(\cG\ni I\) be factor-saturated and let \(G=G(\cG)\).  On a fixed
ordered port set, a \emph{\(G\)-matching signature} of arity \(2m\) is
\begin{equation*}
  c\bigotimes_{e\in M}B_{g_e}^{(e)},
  \qquad
  c\ne0,\quad g_e\in G,
\end{equation*}
where \(M\) is an oriented perfect matching in the preceding sense and
\(B_{g_e}^{(e)}\) places the chosen ordered representative \(B_{g_e}\) on
the oriented edge \(e\).  We allow a simultaneous port permutation; the
scalar \(c\) absorbs all representative scalars, so the notion depends only
on \(G\).

For \(\omega\in\{\mathrm{id},\mathrm{rev}\}\), put
\(B_g^{\mathrm{id}}=B_g\) and
\(B_g^{\mathrm{rev}}=B_g^{\mathsf T}\).  Let \(\delta\) specify a subset of the ports and, at each selected port, a representative
\(B_g\) and one of these two orientations.  Attach its recorded binary through
one native \(X\)-wire and retain the other binary port in the same position;
write \(f^\delta\) for the resulting actual dressed signature.

For each \(\delta\), choose distinct ordered ports \((i,j)\), a closing
representative \(B_h\) with orientation \(\omega\), and a residual port
order \(\rho\), leaving at least two ports.  The resulting card is
\[
 \rho\!\left(\partial_{(i,j)}^{XB_h^\omega X}f^\delta\right).
\]
The \emph{matching deck} \(\mathfrak D_G(f)\) contains zero and all such
one-copy proper cards.  Thus one copy of \(f\) is used and exactly one
ordered port pair is closed; the binary dressings and closing chain use the
fixed actual representatives.  Every nonzero card is a direct gadget output,
never an analytic linear combination.
\end{definition}

Adjoining a Turing-exposed family preserves the problem up to Turing
equivalence: the exposure reduction removes the new signatures, and
inclusion gives the reverse reduction.  Direct outputs use
\cref{eq:direct-gadget-reduction}; their factors use
\cref{lem:factor-saturation}.

For termination, fix \((A,\Lambda,f)\) and
\(H:=\overline G(A,\Lambda)\).  While \(H\) and \(\nu(\Lambda)\) remain
unchanged, keep one fixed actual representative of each group class.  Let
\(\mathscr U(A,f,H)\) be the set of canonical keys, normalized after
pulling a card back by \(A\), of zero and of every one-copy proper card.
Concretely, its finite datum records a choice of binary dressings using the frozen
representatives indexed by \(H\), the ordered deleted pair, the frozen
closing representative \(B_h\), the resulting kernel \(XB_h^\omega X\) and
its orientation \(\omega\), and the residual port order.
The key retains the ordered
external port list and the first-nonzero-entry normalization; the discarded
nonzero scalar and the gadget provenance are stored separately.  Attach to
each nonzero key its finite factorization tree from the worklist in
\cref{lem:factor-saturation}, and write
\(\mathscr U^{\mathrm{fac}}(A,f,H)\) for the resulting finite set of card and
factor tokens.  A fixed-\((H,\nu)\) queue contains each token at most once.  Write
\[
 Q_{\rm done}\subseteq\mathscr U^{\mathrm{fac}}(A,f,H),
 \qquad
 Q=\mathscr U^{\mathrm{fac}}(A,f,H)\setminus Q_{\rm done}
\]
for the completed and unfinished indices, respectively; \(|Q|\) is finite.  During the local deck-stabilization construction, a card token
contributes one of the already enumerated members of
\(\mathscr U(A,f,H)\), and a factor token is processed entirely inside its
attached finite factorization tree from \cref{lem:factor-saturation}.  Thus this
  construction cannot create a fresh fixed-\((H,\nu)\) token outside the fixed finite set.  A lower \(\nu\), a strict enlargement of the pulled-back group,
  or a terminal dispatch stops the local construction.  The literal equality
  \(I\) and the caller's coordinates remain fixed throughout this protocol.

\begin{definition}[Deck-stable state]
\label{def:deck-stable-state}
A retained state \((\Lambda,G,f)\) is \emph{deck-stable} if \(\Lambda\) is
finite, factor-saturated, and contains \(I\); \(f\in\Lambda\) is a
minimum-arity tensor-prime nonbinary signature; \(G=G(\Lambda)\) is the
recomputed complete finite safe group, with fixed actual representatives;
and the regenerated deck, with all card factors retained, consists of zero
and \(G\)-matching signatures.  Actuality is always relative to \(\Lambda\);
retained factors are only Turing-exposed over their predecessor.
\end{definition}

The next lemma either produces a deck-stable state, resolves the complexity
of the current set, or exposes a smaller nonbinary prime.

\begin{lemma}[Deck stabilization under factor access]
\label{lem:minimum-core-deck}
Let \(\Lambda_0\) be a retained, factor-saturated finite algebraic
signature set containing the distinguished literal equality \(I\), and let
\(f\in\Lambda_0\) be an even-arity tensor-prime nonbinary signature of
arity \(n=\nu(\Lambda_0)\).

Then either an established terminal or the signature-set-wide conclusion
\(\Tract(K\Lambda_0)\) applies, or the construction returns one of the
following retained successors:
\begin{enumerate}
\item a finite retained factor-saturated augmentation
      \(\Lambda_{<}\supseteq\Lambda_0\) such that
      \[
        \KHolant(\Lambda_{<})\equivT\KHolant(\Lambda_0),
        \qquad \nu(\Lambda_{<})<n;
      \]
      every nonbinary prime retained in this successor has even arity;
\item a finite retained factor-saturated augmentation
      \(\Lambda_\star\supseteq\Lambda_0\) such that
\begin{equation*}
  \KHolant(\Lambda_\star)
  \equivT
  \KHolant(\Lambda_0),
\end{equation*}
and, for \(G_\star=G(\Lambda_\star)\), every member of
\(\mathfrak D_{G_\star}(f)\) is zero or a
\(G_\star\)-matching signature.
\end{enumerate}

The stable successor retains the actual representatives, the regenerated
complete deck, and every card factor.  Such factors are Turing-exposed over,
but need not be directly realizable from, their predecessor.
\end{lemma}

\begin{proof}
Construct states \((A_i,\Lambda_i,H_i,Q_i)\), where
\(H_i=\overline G(A_i,\Lambda_i)\le\PGL_2(L_0)\) is in the fixed
pulled-back coordinates and \(Q_i\) is the finite canonical card/factor queue
from the preceding convention.  At stage \(i\), dispatch unsafe binaries,
recompute the complete group \(G_i=G(\Lambda_i)\), retain one actual
representative of every class, and generate each key in
\(\mathscr U(A_i,f,H_i)\) at most once.  Preprocess and jointly
factor-saturate the resulting finite batch to obtain \(\Lambda_{i+1}\).
Retained \(I\) excludes the one-sided exception.  Inclusion supplies the
reverse reduction at every exposure step, so every \(\Lambda_i\) is
Turing-equivalent to \(\Lambda_0\).

Every nonbinary card factor has arity at most \(n-2\).  Odd factors are
terminal by \cref{thm:external-odd}; an even one gives successor~1.
Otherwise only binaries remain.
  No global normalization is made in this fixed-\(I\) protocol.  Completeness
  therefore gives \(H_i\le H_{i+1}\); only a strict inclusion
  \(H_i\subsetneq H_{i+1}\) is an outer group successor, and it alone triggers
  regeneration of the complete deck.

Use the frozen pulled-back field \(L_0\) from
\cref{lem:normalization-ledger}.  Gadget evaluation and pivot-normalized
rank-one splitting are performed in pulled-back coordinates, so
\(H_i\le\PGL_2(L_0)\).  Let \(d_{L_0}:=[L_0:\mathbb Q]\) be the field
degree, and define the torsion exponent and group-order bound by
\begin{equation}
  e_{L_0}=\operatorname{lcm}\{r\ge1:\varphi(r)\le2d_{L_0}\},
  \qquad
  M_{L_0}=\max\{2e_{L_0},60\},
  \label{eq:deck-group-bound}
\end{equation}
where \(\varphi\) denotes Euler's totient function.
\Cref{prop:finite-pgl2-classification,lem:bounded-torsion} give
\(|H_i|\le M_{L_0}\).  Hence the strict subgroup chain has at most
\(M_{L_0}-1\) changes.  Between two such changes the finite queue is exhausted:
each card/factor token is removed once, and the inner factor potential
\(\Phi\) of \cref{eq:factor-potential} decreases while its factorization tree is
processed.  This stabilization procedure invokes only the one-copy proper
cards in \(\mathscr U(A_i,f,H_i)\) and the factorization trees attached to those
cards; it does not yet invoke any branch-specific multi-copy reduction.
Therefore every task that can remain in the same pair \((H_i,n)\) is a
token of this locally frozen queue by construction.  A lower-arity factor, a
new group class, or a terminal stops the current stage instead.  Consequently
there is a first stage at which both
\(H_{i+1}=H_i\) and the queue is empty; call this the stabilized stage.
At that stage define the stable retained signature set and its complete group by
\[
  \Lambda_\star=\Lambda_{i+1},
  \qquad G_\star=G_{i+1}.
\]
Every \(0\ne h\in\mathfrak D_{G_\star}(f)\) then has retained binary prime
factors \(B_1,\ldots,B_m\) satisfying
\[
  [B_jX]\in G(\Lambda_\star)=G_\star.
\]
Writing \(B_j=\beta_jB_{g_j}\), their blocks form a perfect matching, so
\[
  h=c\bigotimes_{j=1}^m B_{g_j}
\]
for \(c\ne0\), as required.
\end{proof}

Later constructions can lower \(\nu\) or enlarge the complete group.
We call either change an \emph{outer deck successor}; processing a fixed
family while both remain unchanged is a \emph{within-node closure}.
The following conditions specify when such processing terminates.

\begin{definition}[Processing a finite family of signatures]
\label{def:relative-stratum-protocol}
For deck-stable \((\Lambda,G,f)\), take a finite family of signatures
retained in or directly realized over \(\Lambda\), with their port orders
and realization scalars.  Adjoin them, preprocess, jointly factor-saturate,
apply the reductions for odd and unsafe factors, and recompute the group.
After a strict enlargement, regenerate the complete deck.  We refer to these
steps as processing the family.  Analytic combinations and projective tables
alone do not supply signatures to this family.  The closure carries the partition
\(\mathscr U^{\mathrm{fac}}=Q_{\rm done}\sqcup Q\) from the finite-set construction, where \(Q\) is the unfinished queue.  A token moves from \(Q\)
to \(Q_{\rm done}\) only after its direct card and its entire factorization tree
have been processed; a canonical key is never processed twice.
Whether a later branch-specific batch may return to the same case is
governed by the conditions in \cref{def:frozen-token-coverage} below rather than imposed
by this protocol definition.

A call to the stable equality-accessible quaternary boundary stated below is legal
only after the current
physical quaternary card and its complete oriented one-copy deck have been
regenerated, factor-saturated, and dispatched as in
\cref{lem:finite-q4-dispatch}.  Auxiliary cards used inside the finite
certificate are local to that proof and are not appended to \(Q\).  Only a
retained factor batch returned by the reduced interface can be passed to an
outer closure.

The current case is complete once finitely many such families resolve every
branch and \(Q=\varnothing\).  Under the conditions of \cref{def:frozen-token-coverage}, a transition leaving the pulled-back
group and \(\nu\) unchanged completes a queue token, decreases its stored
factor potential, or certifies that none remains, and hence is a within-node
closure; otherwise, \(\nu\) decreases or the pulled-back group enlarges.
\end{definition}

\paragraph{Standing predicates for relative cases.}
To keep the group-specific statements short, write
\begin{align}
 \Stable(\Lambda,G,f)
 &\Longleftrightarrow
   (\Lambda,G,f)\text{ is deck-stable in the sense of
   \cref{def:deck-stable-state}},
   \label{eq:stable-predicate}\\
 \Even(\Lambda)
 &\Longleftrightarrow
   \forall g\in\Pi(\Lambda),\quad
   \arity(g)\ge3\Longrightarrow \arity(g)\in2\mathbb Z.
   \label{eq:even-prime-predicate}
\end{align}
For a family \(\mathfrak C\) of tractable signature classes,
\begin{equation}
 \Close_{\mathfrak C}(\Lambda,G,f)
 \label{eq:relative-close-predicate}
\end{equation}
means that finitely many applications of \cref{def:relative-stratum-protocol}
resolve all branches for \((\Lambda,G,f)\), with \(Q=\varnothing\), and every surviving
common-presentation rigidity leaf satisfies
\(\PresK(K\Lambda;\cC)\) for some \(\cC\in\mathfrak C\).
Established terminals, strict outer successors, and already resolved
\(\Tract\)-leaves are part of closure but are not rigidity leaves;
\(\mathfrak C=\varnothing\) means that no common-presentation rigidity leaf
remains.  These symbols abbreviate conclusions only and do not replace any
protocol hypothesis or transition.

When \(H\) and \(\nu\) do not change, a continuing step must process an
index already in \(Q\) and decrease \(|Q|\) or its stored factor potential.
The precise conditions are as follows.

\begin{definition}[Conditions for descent]
\label{def:frozen-token-coverage}
Fix a deck-stable tuple \((\Lambda,G,f)\), a coordinate matrix \(A\), the
pulled-back group \(H=\overline G(A,\Lambda)\), and a frozen partition
\[
 \mathscr U^{\mathrm{fac}}(A,f,H)=Q_{\rm done}\sqcup Q,
\]
where \(Q\) is the unfinished queue.  Let \(\mathfrak B\) be a collection of admissible families.  Processing any
member must give, on each branch, one of:
\begin{enumerate}
\item a resolved terminal or \(\Tract\)-leaf;
\item a successor with smaller \(\nu\);
\item a strict enlargement of the pulled-back complete group;
\item unchanged \((H,\nu)\), in which case every retained continuing
      resource is indexed by a token of \(Q\), and its processing is
      confined to that token's stored card, provenance, and attached factorization tree.  The closure either moves at least one such token to
      \(Q_{\rm done}\), or, while processing one fixed unfinished token,
      strictly decreases the factor potential of its stored worklist.
\end{enumerate}
 A token is marked done only after its full factorization tree is exhausted.
\end{definition}

The verification is necessarily branch-by-branch and uses the reductions
proved in
\cref{sec:normalized-dihedral,sec:proper-deck,sec:v4-deck,sec:platonic-decks}.
It is therefore deferred to \cref{sec:p1-synthesis}, after those four
arguments are available.

The following lexicographic measure strictly decreases on every continuing
step satisfying these conditions.

\begin{lemma}[Deck-transition dichotomy and termination]
\label{lem:matching-synthesis-measure}
Fix the initial number field \(L_0\).  At a continuing all-even,
binary-residual deck-stable node, carry the state \((A,\Lambda,f)\), a
finite token queue \(Q\), and the distinguished literal equality \(I\).  Let
\(\mathfrak B\) satisfy \cref{def:frozen-token-coverage}, let
\(\Lambda^+\) be the closure of one batch in \(\mathfrak B\) over
\(\Lambda\), and define
the pulled-back group by
\[
  \overline G(A,\Lambda)=[A]^{-1}G(\Lambda)[A].
\]
Then
\begin{equation}
  \KHolant(\Lambda^+)\equivT\KHolant(\Lambda).
  \label{eq:deck-transition-equivalence}
\end{equation}
Batch outputs are direct over \(\Lambda\), while their factors need only be
Turing-exposed there.

After signature-set-wide odd/unsafe and unary/binary dispatch, the only
continuing transitions are:
\begin{enumerate}
\item \(\overline G(A,\Lambda^+)=\overline G(A,\Lambda)\) and
      \(\nu(\Lambda^+)=\nu(\Lambda)\): process a previously unfinished
      token of \(Q\), complete it or strictly decrease its factor potential,
      or certify that \(Q\) is empty.
\item \(\overline G(A,\Lambda^+)=\overline G(A,\Lambda)\) and
      \(\nu(\Lambda^+)<\nu(\Lambda)\): choose a new minimum core and reapply
      \cref{lem:minimum-core-deck}.
\item \(\overline G(A,\Lambda)\subsetneq\overline G(A,\Lambda^+)\): choose
      a minimum augmented core, regenerate its deck, and restart marked
      routing.
\end{enumerate}

At every continuing node define the pair by
\begin{equation*}
 \Xi(A,\Lambda)
 =\bigl(M_{L_0}-|\overline G(A,\Lambda)|,\nu(\Lambda)\bigr).
\end{equation*}
Refine it to the tuple
\begin{equation}
  \Omega(A,\Lambda,f,Q)
  =\bigl(M_{L_0}-|\overline G(A,\Lambda)|,\nu(\Lambda),
        |Q|,\,
        \Phi_Q\bigr).
  \label{eq:deck-closure-measure}
\end{equation}
For each \(u\in Q\), let \(\mathcal W_u\) be the multiset of still-open
occurrences in its attached factor worklist, and define
\[
 \Phi_Q:=\sum_{u\in Q}\Phi(\mathcal W_u),
\]
using the potential in \cref{eq:factor-potential}.
The displayed tuple is ordered lexicographically.  Items~2--3 strictly
decrease its first two coordinates.  In item~1, completing a token strictly
decreases the third coordinate until the queue is empty.  If a token remains
unfinished during a genuine factor split, the fourth coordinate \(\Phi_Q\)
strictly decreases.  If an exposed factor changes the pulled-back group, the
operation is instead item~3 and is charged to the strict subgroup chain.
\end{lemma}

\begin{proof}
Substitution and factor saturation prove
\cref{eq:deck-transition-equivalence}; \(I\) excludes the one-sided
exception, and inclusion reverses the reduction.  Retained augmentation
cannot increase \(\nu\).  The continuing state remains in fixed-\(I\)
coordinates.  New classes can arise only from an exposed factor; such a
class is precisely a strict successor in the pulled-back group.  By
  \cref{lem:normalization-ledger,eq:deck-group-bound}, every continuing group
lies in \(\PGL_2(L_0)\) and has order at most \(M_{L_0}\).  Therefore a lower
  \(\nu\) decreases the second coordinate of \(\Xi\), strict subgroup
  inclusion decreases the first.  A fixed-\((H,\nu)\) step either decreases the
  finite queue counter in \(\Omega\) or, within its fixed token, decreases
  \(\Phi\), so no factor loop is possible.  By
  \cref{def:frozen-token-coverage,def:relative-stratum-protocol}, item~1
  never opens a fresh
  fixed-\((H,\nu)\) branch: it can only process one previously frozen token of
  \(Q\).
\end{proof}

\begin{definition}[The conclusion required of each construction]
\label{def:protocol-discharged}
Write
\[
 \Sigma=(A,\Lambda,H,f,Q),
 \qquad H=\overline G(A,\Lambda),
\]
with the record, deck-stability, frozen queue, and fixed-\(I\) data of
\cref{def:relative-stratum-protocol,def:frozen-token-coverage}.  Write
\[
 \Omega(\Sigma)
 :=\Omega(A,\Lambda,f,Q)
 =\bigl(M_{L_0}-|H|,\nu(\Lambda),|Q|,\Phi_Q\bigr)
\]
for the lexicographic measure in \cref{eq:deck-closure-measure}.

For each concrete construction at \(\Sigma\), we prove that processing its
outputs by \cref{def:relative-stratum-protocol} gives either
\begin{enumerate}
\item a resolved complexity leaf: hardness, a \(\Tract\)-leaf, or an
      accepted common tractable presentation; or
\item a deck-stable tuple \(\Sigma'\), with the same required realization and
      induction data, satisfying
      \[
        \Omega(\Sigma')<_{\rm lex}\Omega(\Sigma).
      \]
\end{enumerate}
The second outcome includes a decrease of \(\nu\), a strict enlargement of
the pulled-back complete group, completion of a frozen token, or a strict
decrease of that token's factor potential.  Every
same-\((H,\nu)\) continuation must be indexed by the already frozen queue
and may create no fresh unresolved fixed-\((H,\nu)\) resource.  An arbitrary
changed signature set is therefore not a successor in this definition.

The first outcome resolves the complexity question; the second continues
the induction at \(\Sigma'\).  Below, saying that a construction is resolved
allows either outcome.  This conclusion requires a proof of realizability,
factor access, and descent for that construction.
The assertion \(\Close_{\mathfrak C}(\Lambda,G,f)\) additionally covers
every case for the current \((H,\nu)\) and restricts the surviving common
presentations to \(\mathfrak C\), as in \cref{eq:relative-close-predicate}.
\end{definition}

We have therefore reached a deck-stable tuple.  For the active stable tuple
\(\Sigma\), every continuing branch of a resolved batch has
strictly smaller \(\Omega\).  In a signature-set-wide completion, however,
\(\Omega\) controls only that active tuple.  A separate outer worklist
induction over unresolved tensor-prime occurrences completes the current
prime and advances to the next; replacing \(f\) by another prime of the same
arity is charged to that outer induction and is not asserted to decrease
\(\Omega\).

\subsection{The main stable-state theorem and four-way routing}
\label{sec:matching-decks}

The preceding reductions leave a deck-stable tuple
\((\Lambda,G,f)\).  The next theorem is proved in
\cref{sec:p1-synthesis} from the four branches routed below.  Together with
stabilization, it gives the equality-accessible main theorem
\cref{thm:p1-eq2-kholant-dichotomy}.

\begin{theorem}[Relative matching-deck classification]
\label{thm:matching-deck-classification}
Assume
\[
 \Stable(\cG,G,f),\qquad \arity(f)=\nu(\cG).
\]
In particular, the deck-stable state contains the distinguished literal
equality \(I\).  If \(\Tract(K(\cG))\) holds, then
\(\KHolant(\cG)\in\FP\); otherwise \(\KHolant(\cG)\) is
\(\SharpP\)-hard.  Equivalently, by \cref{lem:tract-X-anchor}, the
tractable condition is \(\TractX(K(\cG))\).
\end{theorem}

The proof retains the stronger structural conclusion that every nonhard
rigidity survivor has one common product, affine, or flat-Lagrangian
\(K\)-presentation; by \cref{eq:common-K-presentation-predicate}, the same
basis also places the transformed edge in that class.

To route the stable tuple, we now record the marked finite-group notation in
the fixed coordinates containing the distinguished literal equality.
\label{sec:marked-finite-groups}

Define the diagonal matrix family
\[
  R(t)=\diag(1,t).
\]
For \(m\ge2\), define the matrix family \(C(t)\) and the first two marked
dihedral families
\begin{align*}
 \mathcal M_m
 &=\{[R(t)],[XR(t)]:t\in\mu_m\},
 \\
 C(t)&=(1+t)I+i(1-t)Y,\\
 \mathcal B_m
 &=\{[C(t)],[XC(t)]:t\in\mu_m\}.
\end{align*}
For even \(m\), fix
\(\eta\in\mu_{2m}\setminus\mu_m\), define the matrix family
\[
 A(t)=(1+t)I+(1-t)X,
\]
and define
\begin{equation*}
 \mathcal A_m
 =
 \{[A(t)]:t\in\mu_m\}
 \cup
 \{[ZA(\eta t)]:t\in\mu_m\}.
\end{equation*}
In the same even parameter range, define the unshifted marked form
\begin{equation*}
 \mathcal A_m^0
 =
 \{[A(t)],[ZA(t)]:t\in\mu_m\}.
\end{equation*}
These are projective displays.  Any suppressed factor of two belongs to
the coordinate model and does not normalize an actual representative.

Throughout this routing discussion, retain the globally fixed marks
\((x,j)=([X],[Z])\) and put \(H:=\langle G,j\rangle\).
For routing, the Platonic symbols record the following abstract marked
types.  In \(\mathcal T\), \(G=H\cong A_4\) and \(x,j,xj\) are its three
double transpositions.  In \(\mathcal T_{\rm ext}\),
\(G\cong A_4\), \(G\triangleleft H\), and \(H\cong S_4\); the mark
\(x\in G\) is a double
transposition, and \(j,xj\in H\setminus G\) are disjoint transpositions.
In \(\mathcal O\), \(G=H\cong S_4\) and all three marks are double
transpositions.  In \(\mathcal O'\), the marks \(x,j\) are disjoint
transpositions and \(xj\) is a double transposition.  In
\(\mathcal O_{dt}\), \(x\) is a double transposition while \(j,xj\) are
transpositions; in \(\mathcal O_{td}\), \(j\) is a double transposition
while \(x,xj\) are transpositions.  Finally,
\(\mathcal I\) and \(\mathcal I_-\) denote the two fixed-\(I\)
coordinate displays of the ordered marked \(V_4=\langle x,j\rangle\) in
\(G=H\cong A_5\); all three nonidentity marks lie in the unique involution
class.  \Cref{sec:platonic-decks} supplies explicit fixed-equality representatives and the
direct coefficient dictionaries used for the alternative displays.
The router below uses only this membership and conjugacy-class data; the
notation never authorizes an implicit conjugacy normalization.

The literal \(I\) supplies \(x=[X]\), while binary reversal singles out
\(j=[Z]\).  Keeping their ordered positions refines the classification of
\cref{prop:finite-pgl2-classification} into the following four branches.

\begin{theorem}[Marked finite-group routing]
\label{thm:marked-finite-group-routing}
Let \(\Lambda\) be a retained factor-saturated signature set containing the
distinguished literal equality \(I\), and let \(G=G(\Lambda)\) be its complete finite safe
transfer group.  Recall the ordered physical marked pair and its generated
group:
\[
 (x,j):=([X],[Z]),\qquad x\in G,\qquad H:=\langle G,j\rangle.
\]
In the fixed coordinates containing the distinguished literal \(I\), exactly
one of the following alternatives holds:
\begin{enumerate}
\item \(G\cong D_{2m}\), \(m\ge3\), and the marked configuration
      \((G;x,j)\) has one of the forms denoted by
      \(\mathcal M_m,\mathcal B_m,\mathcal A_m,\mathcal A_m^0\).  The last
      two families occur only for even \(m\).
\item \(G=C_N\) is cyclic, where \(N\) is even and \(j\notin G\).  This is
      the cyclic branch.
\item \(G\cong V_4\), and \(G\) is either the standard Pauli group
      \(\mathcal M_2=\mathcal B_2\) or the exotic group
      \(\mathcal A_2=\mathcal V_{\rm ex}\).
\item \(G\) is one of the two marked tetrahedral forms
      \(\mathcal T,\mathcal T_{\rm ext}\), the four fixed-equality marked
      octahedral forms
      \(\mathcal O,\mathcal O',\mathcal O_{dt},\mathcal O_{td}\), or one of the two icosahedral displays
      \(\mathcal I,\mathcal I_-\).
\end{enumerate}

\smallskip
\noindent\emph{Positions of the marks.}
All conjugacy statements in this paragraph are taken in \(H\).  In every
form, \(x\) and \(j\) are distinct commuting involutions and
\(\langle x,j\rangle=\{1,x,j,xj\}\cong V_4\).  In the four dihedral forms the
marked positions are
\begin{align*}
 \mathcal M_m:&\quad
 x=[XR(1)],\quad j=[R(-1)],\quad xj=[XR(-1)],\qquad
 H=\begin{cases}\mathcal M_m,&2\mid m,\\ \mathcal M_{2m},&2\nmid m;\end{cases}\\
 \mathcal B_m:&\quad
 x=[XC(1)],\quad j=[XC(-1)],\quad xj=[C(-1)],\qquad
 H=\begin{cases}\mathcal B_m,&2\mid m,\\ \mathcal B_{2m},&2\nmid m;\end{cases}\\
 \mathcal A_m:&\quad
 x=[A(-1)]\in G,\quad j=[ZA(1)]\in H\setminus G,\quad
 xj=[ZA(-1)]\in H\setminus G;\\[-2pt]
 &\quad
 H=\{[A(t)]:t\in\mu_{2m}\}\cup
   \{[ZA(t)]:t\in\mu_{2m}\}\cong D_{4m};\\
 \mathcal A_m^0:&\quad
 x=[A(-1)],\quad j=[ZA(1)],\quad xj=[ZA(-1)],\qquad
 H=G=\mathcal A_m^0.
\end{align*}
In \(\mathcal M_m\), \(j\) is the central half-turn and \(x,xj\) are
reflections; in \(\mathcal B_m\), \(x,j\) are reflections and \(xj\) is the
central half-turn.  In both cases \(j\in G\) exactly when \(m\) is even, and
\(x,xj\) in \(\mathcal M_m\), respectively \(x,j\) in \(\mathcal B_m\), are
conjugate exactly when \(4\mid m\).  In
\(\mathcal A_m\), where \(m\) is even, \(x\) is the central half-turn and
\(j,xj\) are conjugate reflections in \(H\setminus G\).
In \(\mathcal A_m^0\), \(x\) is again the central half-turn, while
\(j,xj\in G\) are reflections; they are conjugate in \(G\) exactly when
\(4\mid m\).
When \(m\) is odd in either \(\mathcal M_m\) or \(\mathcal B_m\), one has
\(j,xj\in H\setminus G\); when \(m\) is even, all three marks lie in \(G\).

In the cyclic branch \(G=C_N\), the order \(N\) is even,
\(j,xj\in H\setminus G\), \(H\cong D_{2N}\), \(x\) is the central half-turn, and
\(j,xj\) are reflections, conjugate exactly when \(4\mid N\).

At \(m=2\), the standard Pauli form has \(H=G\), with \(x,j,xj\) the three
nonidentity elements (three singleton conjugacy classes) of \(V_4\); the
exotic form has \(j,xj\in H\setminus G\), \(H\cong D_8\), central half-turn
\(x\), and conjugate outside reflections \(j,xj\).  In \(\mathcal T\),
\(H=G\cong A_4\) and \(x=[Q_1],j=[Q_3],xj=[Q_2]\) are the three double
transpositions, hence one conjugacy class.  In \(\mathcal T_{\rm ext}\),
\(G\cong A_4\), \(G\triangleleft H\), and \(H\cong S_4\);
\(x=[Q_1]\in G\) is a double transposition,
and \(j=[Q_3],xj=[Q_2]\in H\setminus G\) are disjoint transpositions in one
  class.  In \(\mathcal O\), \(H=G\cong S_4\) and \(x,j,xj\) are the three
  double transpositions.  In \(\mathcal O'\), \(x,j\) are disjoint
  transpositions and \(xj\) is a double transposition.  In
  \(\mathcal O_{dt}\), \(x\) is a double transposition and \(j,xj\) are
  transpositions; in \(\mathcal O_{td}\), \(j\) is a double transposition
  and \(x,xj\) are transpositions.  In all four cases \(H=G\cong S_4\).
  Finally, in
each icosahedral coordinate display \(H=G\cong A_5\), and
\(x=[Q_1],j=[Q_3],xj=[Q_2]\) lie in its unique involution class.
\end{theorem}

Each of the four group-specific sections,
\cref{sec:normalized-dihedral,sec:proper-deck,sec:v4-deck,sec:platonic-decks},
proves a relative theorem for its corresponding stable tuple.  It either
determines the complexity at that tuple, or
produces a physical transition whose closure yields a strictly smaller
active tuple under \(\Omega\).  In the second case, the deck is stabilized
again and the new tuple is routed anew.  Movement between distinct
same-arity primes in a signature-set-wide proof remains governed by the
separate outer worklist induction of \cref{def:protocol-discharged}.

The group-specific proofs may use explicitly stated coefficient or
open-network dictionaries, but these do not mutate the retained fixed-\(I\)
state.  If retaining an exposed factor enlarges \(G\), the current route
is invalidated; one must first reapply \cref{lem:minimum-core-deck} and then
reroute before invoking a group-specific case theorem.

\begin{proof}
The reversal identity
\([T]\mapsto j[T]^{-1}j=[ZT^{-1}Z]\) (with \(j=[Z]\)) and completeness give
\[
 jGj=G,\qquad G\triangleleft H,\qquad [H:G]\le2.
\]
The two distinct projective involutions \(x,j\) commute, so
\(\langle x,j\rangle\cong V_4\le H\).  By
\cref{prop:finite-pgl2-classification}, this leaves an even dihedral group
or one of \(A_4,S_4,A_5\).

Suppose first that \(G\cong D_{2m}\) with \(m\ge3\).  Then \(H\) is also
dihedral.  Indeed, if \(H\) were Platonic and \([H:G]=2\), then \(G\)
would be an index-two normal subgroup.  By
\cref{lem:finite-pgl2-facts}, the only such subgroup among the Platonic
groups is \(A_4\triangleleft S_4\), which is not dihedral.  The case
\(H=G\) is immediate.  Choose a rotation generator \(r\) and a reflection
generator \(s\) so that
\(H=\langle r,s:r^n=s^2=1,\ srs=r^{-1}\rangle\).
Since \(H\) contains \(\langle x,j\rangle\cong V_4\), \(n\) is even.
Exactly one of \(x,j,xj\) is the central half-turn \(r^{n/2}\), and
the other two are reflections.  If the central mark is \(j\) or \(xj\),
we obtain \(\mathcal M_m\) or \(\mathcal B_m\), respectively.
If it is \(x\), membership \(j\in G\) gives \(\mathcal A_m^0\),
while \(j\notin G\) gives \(\mathcal A_m\).
For \(j\in G\), one has \(H=G\), so \(n=m\) is even.
For \(j\notin G\), one has \(n=2m\) and
\(G=\langle r^2,r^ks\rangle\) for some \(k\in\{0,1\}\).
The central half-turn \(r^m\) of \(H\) belongs to \(G\) exactly when \(m\) is even.
Thus \(m\) is odd in the first two outside-mark cases and even in
\(\mathcal A_m\), where \(x\in G\).  This exhausts the four forms
and their parameter ranges.

In \(\mathcal A_m^0\), the mark \(j=[Z]\) belongs to the complete group.
Its fixed actual representative \(B_j\) satisfies
\([B_jX]=j\), hence \(B_j\doteq Y\).  This actual binary and the already
available literal \(I\) give the signature-preserving edge switch used in
\cref{thm:marked-A0-reduction}; the distinguished equality is never
replaced or re-anchored.

The marked eigenline equations give the displayed matrices, with
\[
 ZC(t)Z=tC(t^{-1}),\qquad
 x=[XC(1)],
\]
while \(x=[A(-1)]\) is central in both \(\mathcal A_m\) and
\(\mathcal A_m^0\); hence the four marked configurations are legally distinct.

Next suppose that \(G\) is cyclic.  By
\cref{lem:finite-pgl2-facts}, its unique involution is \(x\).  A class
commuting with \([X]\) lies in \(\Span\{I,X\}\) or
\(\Span\{Z,Y\}\), and every nonsingular member of the latter is another
involution.  Thus all transfer classes have representatives in
\(\Span\{I,X\}\).  Writing \(G=C_N\), it follows that \(N\) is even and
\(j\notin G\).  Conjugation by \(j\) inverts the eigenvalue ratio of every
class in this algebra, so
\[
 H=G\rtimes\langle j\rangle\cong D_{2N}.
\]
Consequently, \(x\) is the central half-turn and \(j,xj\) are the two
marked reflections.

It remains to consider \(G\cong V_4\), corresponding to \(m=2\) in the
dihedral notation.  Here
\[
 \mathcal M_2=\mathcal B_2=\{[I],[X],[Z],[Y]\}.
\]
For \(\eta=i\),
\[
 \frac{ZA(i)}{1-i}=Y+iZ,\qquad
 \frac{ZA(-i)}{1+i}=Y-iZ,
\]
which is \(\mathcal V_{\rm ex}\).  Conversely, every \(V_4\ni x\) has the form
\[
 \{[I],x,[aZ+bY],[bZ+aY]\},\qquad a^2\ne b^2,
\]
and \(jGj=G\) gives either \(ab=0\) (the Pauli form) or
\(a^2=-b^2\) (the exotic form \(\mathcal V_{\rm ex}\)).

Thus the Pauli form has \(H=G\), with \(x,j,xj\) its three nonidentity
elements.  In the exotic form, \(j,xj\notin G\) and
\(H=\langle G,j\rangle\cong D_8\); relative to this \(H\), \(x\) is the central
half-turn and \(j,xj\) are conjugate reflections in the outside coset.

It remains to route the Platonic groups while retaining the ordered marked
pair.  Let \(G\cong A_4\).  If \(j\in G\), then \(H=G\), and the three
nonidentity elements \(x,j,xj\) of their marked \(V_4\) are precisely the
double transpositions.  This is \(\mathcal T\).  If \(j\notin G\), then
\(|H|=24\).  By
\cref{prop:finite-pgl2-classification,lem:finite-pgl2-facts}, the
finite-\(\PGL_2\) classification and the subgroup structure of cyclic and
dihedral groups force
\(H\cong S_4\) with \(G=A_4\triangleleft H\).  Thus \(x\in G\) is a
double transposition, whereas the odd involutions \(j,xj\in H\setminus G\)
are transpositions.  Since they are distinct and commute, they are disjoint.
This is \(\mathcal T_{\rm ext}\).

Now let \(G\cong S_4\).  First note that \(j\in G\):
otherwise \(H=\langle G,j\rangle\) would be a finite group of order \(48\).
By \cref{prop:finite-pgl2-classification,lem:finite-pgl2-facts}, a finite
projective group of order \(48\) is
cyclic or dihedral (here \(D_{48}\) denotes the order-\(48\) dihedral
group).  Thus \(H\cong C_{48}\) or \(D_{48}\); every subgroup
of a cyclic or dihedral group is cyclic or dihedral, whereas \(S_4\) is
neither.  Hence \(j\in G\).

  The two commuting involutions may independently lie in the transposition
  or double-transposition class, subject to the class of their product.  The
  four resulting ordered class patterns are
  \[
    (dd;d),\qquad (tt;d),\qquad (dt;t),\qquad (td;t),
  \]
  where the entries record the classes of \((x,j;xj)\).  Conjugation in
  \(S_4\) is transitive on ordered pairs in each pattern by
  \cref{lem:finite-pgl2-facts}.  They are therefore exactly the four
  fixed-equality marked types
  \[
    \mathcal O,\qquad \mathcal O',\qquad
    \mathcal O_{dt},\qquad \mathcal O_{td}.
  \]
  Their ordered marks distinguish the four types in the fixed coordinates.
  Concrete representatives and direct finite exits are recorded in
  \cref{sec:platonic-decks} and
  Appendix~\ref{appsec:cert-platonic}.

For \(G\cong A_5\), necessarily \(j\in G\).  Otherwise
\(G\triangleleft H\) and \([H:G]=2\) would give a finite subgroup
\(H\leq\PGL_2(\mathbb C)\) of order \(120\) containing \(A_5\).  By
\cref{prop:finite-pgl2-classification,lem:finite-pgl2-facts}, such an \(H\)
would be cyclic or dihedral; every subgroup of a cyclic or dihedral group is
cyclic or dihedral, a contradiction.  Hence
\(H=G\).  Finally, \cref{lem:finite-pgl2-facts} shows that \(A_5\) has one
involution class, with \(x,j,xj\) in that class.  Once an ordered pair
\((x,j)\) is fixed, two conjugators differ by
\[
 C_{\PGL_2(\mathbb C)}(\langle x,j\rangle)=\langle x,j\rangle.
\]
Each of these four projective classes has a representative in \(GO(X)\).
  Hence the abstract marked type is unique.  Fixed-equality coordinates give
  the two displays \(\mathcal I,\mathcal I_-\); \cref{sec:platonic-decks} verifies their
  explicit representatives and a direct coefficient-and-card dictionary
  that leaves the supplied literal \(I\) untouched.  This deferred coordinate calculation is not used
to classify the abstract group, so the four-way list is exhaustive.
\end{proof}

\begin{remark}[Relation with Xia's binary-group taxonomy]
\label{rem:xia-binary-taxonomy}
Version~1 of Xia's framework
\cite[Sections~1 and~4]{Xia2026FrameworkV1} organizes the finite projective
binary residue into nine \emph{unmarked} rotation-group types.  Here the
group consists of transfer classes \([T_B]=[BX]\), rather than projective
classes \([B]\) of the underlying binary matrices, and the supplied literal equality fixes the
physical marks
\[
  x=[X]\in G,
  \qquad
  j=[Z],
  \qquad
  jGj=G.
\]
Consequently, our routing refines Xia's taxonomy by the ordered marked pair
\((x,j)\).  Forgetting these marks gives the following dictionary.
\begin{center}
\footnotesize
\setlength{\tabcolsep}{4pt}
\renewcommand{\arraystretch}{0.92}
\begin{tabularx}{0.96\linewidth}{@{}L{0.21\linewidth}X@{}}
\toprule
Xia's unmarked type & Corresponding case in this proof\\
\midrule
\(C_1\) &
absent at a continuing node, where a nonidentity transfer is already retained\\
\(C_2\) &
the proper-cyclic \(C_2\) quaternary boundary\\
\(C_n\), \(n\ge3\) &
  the odd no-equality cyclic residue, or fixed-\(I\) proper cyclic
\(C_N\ni[X]\)\\
\(D_{2m}\), \(m\) odd &
the odd-dihedral residue, or \(\mathcal M_m,\mathcal B_m\)\\
\(V_4=D_4\) &
the standard Pauli and exotic Klein forms\\
\(D_{2m}\), even \(m\ge4\) &
\(\mathcal M_m,\mathcal B_m,\mathcal A_m^0,\mathcal A_m\)\\
\(A_4\) & \(\mathcal T,\mathcal T_{\rm ext}\)\\
\(S_4\) & \(\mathcal O,\mathcal O',\mathcal O_{dt},\mathcal O_{td}\)\\
\(A_5\) & the two fixed-equality coordinate displays\\
\bottomrule
\end{tabularx}
\end{center}
Thus the dihedral, Klein, and Platonic displays above are physical marked
refinements, not additional abstract finite-group types.  The table is only
a notation dictionary: no completeness or hardness result from the cited
framework is used.  This paper independently proves every marked case,
actual-representative statement, and signature-set-wide lift.
\end{remark}

Work in the fixed coordinates containing the distinguished literal \(I\).  Let
\(\Lambda\subseteq\Lambda^+\) be a retained factor-saturated augmentation,
put
\[
 G=G(\Lambda),\qquad G^+=G(\Lambda^+),
\]
and suppose that the unsafe-binary dispatch has not already applied.  If a
newly exposed ordered binary \(B\) is safe and \([BX]\notin G\), then
\cref{lem:actual-binary-completeness} gives the strict inclusion
\(G<G^+\).  If \([BX]\in G\), the binary itself is old, but the complete
group is still recomputed: the augmented set may realize additional
multi-copy binaries not enumerated by the batch that exposed \(B\).
Accordingly, \(G^+\) is the complete semantic closure, and need not equal
the subgroup generated by \(G\) and the first new class.  Whenever
\(G<G^+\), regenerate the deck and reroute before invoking a
group-specific theorem.

\subsection{Shared methods for the four branch proofs}
\label{subsec:shared-matching-methods}
\label{sec:equality-q4-boundary}
\label{proofsec:finite-certificates}

We use the two-port slices and bilinear Frobenius pairing of
\cref{subsubsec:two-port-fibres} to determine what the available
contractions reveal about a tensor.  The quaternary result at the end of
this subsection supplies a common base case for the later branches.

\begin{definition}[The span \(\mathscr K_G\)]
\label{def:effective-kernel-algebra}
For a retained normalized set with complete group \(G\), let
\(\mathscr K_G\) be the complex span of the transfer matrices
\(T_B=BX\) of all finite ordered two-port chains \(B\) made from the fixed
actual representatives \(B_g\), \(g\in G\), the explicit \(I\)- or
native-\(X\) links, and their port reversals.  This is equivalently the span
of independently dressed \(I/X\) kernels after the fixed right
multiplication by \(X\), which is invertible and therefore only relabels
each observable fibre.
\end{definition}

The information determined by binary contractions depends on
\(\mathscr K_G\).  The following theorem gives its two possible values.

\begin{theorem}[The span of realizable binary kernels]
\label{thm:dressed-kernel-observability}
Let \(\cG\) be a retained signature set with normalized \(I,X\), and let
\(G=G(\cG)\).  With \(\mathscr K_G\) as in
\cref{def:effective-kernel-algebra},
\[
  \mathscr K_G\in
  \bigl\{\Span\{I,X\},\Mat_2(\mathbb C)\bigr\}.
\]
Moreover,
\[
 \Span\{I,X\}^{\perp_F}
 =
 \left\{
  \begin{pmatrix}u&v\\-v&-u\end{pmatrix}:u,v\in\mathbb C
 \right\}.
\]
\end{theorem}

If \(\mathscr K_G=\Mat_2(\mathbb C)\), four actual kernel evaluations
determine each two-port slice.  If \(\mathscr K_G=\Span\{I,X\}\), they
determine it modulo the displayed orthogonal complement.  This reconstruction
is an algebraic identity and does not itself realize a gadget.

\begin{proof}
Because \(B_1=X\), the span contains \(I\), and the explicit \(I\)-link
gives \(X\).  Series composition, with the second chain taken in its
inward (reversed) port order, has physical tensor \(BXC\) and satisfies
\(T_{BXC}=T_BT_C\), so the span is an algebra.  Reversal gives
\(T^\#=XT^{\mathsf T}X\); since \(X\) is in the span,
\(T^{\mathsf T}=XT^\#X\) is also in it.  Thus \(\mathscr K_G\) is unital
and transpose-stable.

Conjugate by the Walsh matrix
\(S:=2^{-1/2}\left(\begin{smallmatrix}1&1\\1&-1\end{smallmatrix}\right)\),
so \(SXS=Z\).  For \(A\) in the conjugated algebra, define its diagonal
and off-diagonal parts by \(D:=(A+ZAZ)/2\) and \(O:=(A-ZAZ)/2\).  If some
\(O\ne0\), write
\(O=\left(\begin{smallmatrix}0&b\\c&0\end{smallmatrix}\right)\).  Since \(Z\)
and \(O\) lie in the algebra, so does
\(ZO=\left(\begin{smallmatrix}0&b\\-c&0\end{smallmatrix}\right)\).  If
\(bc\ne0\), \(O+ZO=2bE_{12}\) and \(O-ZO=2cE_{21}\), where \(E_{12}\)
and \(E_{21}\) are the standard \(2\times2\) matrix units; if one entry is zero,
transpose-stability supplies the missing unit.  Products supply the diagonal
units, hence the algebra is \(\Mat_2(\mathbb C)\).  If every \(O=0\), it is diagonal and contains \(I,Z\), so it is
\(\Span\{I,Z\}\), which conjugates back to \(\Span\{I,X\}\).

In the full case, four actual chain matrices can be greedily chosen as a
basis; their Frobenius evaluations reconstruct every two-port slice
algebraically, not as a gadget.  In the proper case, orthogonality to
\(I,X\) is \(a+d=b+c=0\), giving exactly
\(\left(\begin{smallmatrix}u&v\\-v&-u\end{smallmatrix}\right)\).
\end{proof}

For the finite groups routed above, observability identifies the routing split
intrinsically:
\[
 G\text{ is cyclic}
 \quad\Longleftrightarrow\quad
 \mathscr K_G=\Span\{I,X\},
\]
\[
 G\text{ is dihedral, Klein four, or Platonic}
 \quad\Longleftrightarrow\quad
 \mathscr K_G=\Mat_2(\mathbb C).
\]
Indeed, the cyclic transfer classes lie in \(\Span\{I,X\}\), whereas a
finite projective subgroup of the units of this proper algebra is cyclic.
Thus \cref{sec:proper-deck} treats \(\mathscr K_G=\Span\{I,X\}\);
the other three branches have \(\mathscr K_G=\Mat_2(\mathbb C)\).

At arity four, the complete stabilized deck permits the following
gadget construction.  Its output is an actual nonmatching eight-vertex
signature to which \cref{cor:p1-eight-vertex-reduced-interface} applies.

\begin{theorem}[Quaternary signatures with stable matching decks]
\label{thm:equality-q4-boundary}
Assume
\[
 \Stable(\Lambda,G,q),\qquad \arity(q)=\nu(\Lambda)=4,
\]
and that \(G\) is one of
\begin{enumerate}
\item \(G=\{1,[X]\}\cong C_2\);
\item \(G=\{[\diag(1,t)],[X\diag(1,t)]:t\in\mu_m\}\), \(m\ge3\);
\item the standard Pauli \(V_4\);
\item the exotic Klein group
  \begin{equation*}
   \mathcal V_{\rm ex}
   =
   \left\{[I],[X],
   \left[\begin{pmatrix}i&1\\-1&-i\end{pmatrix}\right],
   \left[\begin{pmatrix}-i&1\\-1&i\end{pmatrix}\right]\right\}.
  \end{equation*}
\end{enumerate}
Then an actual gadget using at most two copies of \(q\) produces an
endpoint-nondegenerate eight-vertex signature that is not \(G\)-matching.
\end{theorem}

The quantifiers in this statement include zero cards, both port orientations
for all six contracted pairs, every current-group representative assignment
on the remaining matched edges, and the fixed representative scalars.  By
\cref{cor:p1-eight-vertex-reduced-interface}, the output gives a resolved
hard/tractable leaf or a retained factor batch.  The latter is dispatched by
\cref{lem:actual-binary-completeness}; only a safe factor outside \(G\) passes
to the global measure as a strict outer-group successor.

\begin{proof}
By \cref{thm:cert-interface-equality-q4}, the asserted directly realized signature
\(r\) satisfies
\[
  \KHolant(\Lambda,r)
  \leT\KHolant(\Lambda).
\]
Apply \cref{cor:p1-eight-vertex-reduced-interface}.  In items~2--4 and the one-copy
case of item~1, invertible dressings preserve tensor primality, so no factor
outcome occurs.

In the exceptional \(C_2\) subcase, write
\(a_S=\widetilde q(\one_S)\).  Call it \emph{mixed-sign} when
\(a_\varnothing=a_{[4]}=0\) and two nonzero complementary weight-two
blocks \(\{S,\bar S\}\) have different ratios
\(a_{\bar S}/a_S\in\{1,-1\}\), as in
\cref{thm:cert-interface-equality-q4}\textnormal{(i)}.  Only this mixed-sign \(C_2\)
subcase can reach the continuing factor outcome after the two-copy
construction.  Keep the complete retained factor batch supplied by
the reduced interface and dispatch each binary prime separately using
\cref{lem:actual-binary-completeness}.  A singular or infinite-order prime
is a terminal; a safe prime outside the current \(G\) gives a strict
complete-group successor.  If all safe primes were in \(G\), the genuine
\(2+2\) product would be \(G\)-matching, contrary to the physical
nonmatching certificate.  Thus no same-\((G,\nu)\) successor is created;
every retained factor remains Turing-exposed over the predecessor, and only
the resulting outer successor is passed back to stabilization.
\end{proof}

Every bounded enumeration appears as an exact proposition in
Appendix~\ref{sec:exact-finite-certificates}; the main text supplies its reduction
direction and signature-set-wide lift.  The coverage map
\path{certificates/theorem-to-script-map.md} records the manuscript anchors
and finite claims checked by each entry point, and
\path{certificates/README.md} gives the replay instructions and the
package-wide computational scope.  The proper- and full-kernel branches
meet only at the final algebraic split.

With the stable-state program, four-way routing, and common methods in
place, we begin with the marked dihedral branch.
\section{Marked Dihedral Matching Decks}
\label{sec:normalized-dihedral}

We prove rigidity for the four marked dihedral groups
\(\mathcal M_m,\mathcal B_m,\mathcal A_m^0,\mathcal A_m\) from
\cref{thm:marked-finite-group-routing}.  The input is a stable state: its
proper contractions are already zero or products of available binaries.
The task is to show that a tensor-prime core of arity at least six cannot
have this property.  At arity four, the sector-specific constructions below
produce an actual eight-vertex signature and invoke its dichotomy.

For \(\mathcal M_m\), monomial binaries turn contractions into weighted
parity constraints.  Comparing two such contractions reduces a disagreement
to six ports; their common factors can then be recovered by linear
inversion.  We use the projective tensor geometry of
\cref{subsubsec:projective-tensor-geometry} to identify these common
factors.  The four-port support classification in
\cref{thm:cert-interface-equality-q4} uses the exact certificate described
there; the six-port classification in
\cref{prop:cert-nd-q6-support-lowering} has a complete 48-case table and
a normalization proof.
Whole-network equivalences handle \(\mathcal B_m\) and \(\mathcal A_m^0\), while
explicit weighted kernels handle \(\mathcal A_m\).

For \(m\ge3\), define the diagonal and anti-diagonal matrix families
\[
  \mathsf R(t)=\diag(1,t),
  \qquad
  \mathsf S(t)=X\mathsf R(t),
  \qquad
  t\in\mu_m.
\]
The odd-arity case is dispatched by \cref{thm:external-odd}.
We also establish support and coefficient reconstruction for later
occurrences: four slices determine the support, a minor proportional to
\(1-t^2\) excludes the Reed--Muller alternative, and two-copy contractions
test whether the remaining coefficients factor.

\begin{theorem}[Marked-dihedral deck rigidity]
\label{thm:marked-dihedral}
Assume
\[
 \Stable(\Lambda,G,f),\qquad \Even(\Lambda),
\]
and suppose that the complete marked group is one of
\[
 G\in
 \{\mathcal M_m:m\ge3\}
 \cup
 \{\mathcal B_m:m\ge3\}
 \cup
 \{\mathcal A_m^0:m\ge4,\ 2\mid m\}
 \cup
 \{\mathcal A_m:m\ge4,\ 2\mid m\}.
\]
Then \(\Close_{\{\cP\}}(\Lambda,G,f)\).
\end{theorem}

The four cases are proved in
\cref{thm:normalized-dihedral,thm:marked-B-reduction,%
thm:marked-A0-reduction,thm:marked-A-rigidity}; the final proof combines
them with exactly the parameter ranges displayed above.

\paragraph{Actual-card and sector-pencil setup.}
Let the core arity be \(n:=\arity(f)\).  We use the conventions of
\cref{subsec:reduction-conventions,def:relative-stratum-protocol,%
def:protocol-discharged} throughout.

If \(n=4\), first regenerate and factor-saturate the complete physical deck
of \(f\), apply \cref{thm:equality-q4-boundary}, and route its actual
endpoint-nondegenerate output through
\cref{cor:p1-eight-vertex-reduced-interface}.  The retained factor batch is
closed by \cref{lem:actual-binary-completeness,lem:factor-saturation}; it
gives a resolved leaf or a strict stable successor.  Thus the arity-four
base is resolved, and throughout the remaining standard-sector
proof we may assume
\[
 n=2k+2\ge6.
\]

For every \(t\in\mu_m\), the representative convention of
\cref{subsec:reduction-conventions} fixes actual ordered binary outputs
\begin{equation*}
  B_t^{\mathsf R}=\sigma_t^{\mathsf R}\mathsf R(t)X,
  \qquad
  B_t^{\mathsf S}=\sigma_t^{\mathsf S}\mathsf S(t)X,
  \qquad
  \sigma_t^{\mathsf R},\sigma_t^{\mathsf S}\ne0.
\end{equation*}
For a deleted set \(p=\{i,j\}\) equipped with the port order \((i,j)\),
define the four analytic slices
\begin{equation*}
  A:=f(00,\mathord\cdot),\quad
  B:=f(01,\mathord\cdot),\quad
  C:=f(10,\mathord\cdot),\quad
  D:=f(11,\mathord\cdot),
\end{equation*}
and define the associated \((n-2)\)-ary \emph{sector pencils} by
\begin{equation}
 E(t):=A+tD,
 \qquad
 O(t):=C+tB.
 \label{eq:nd-sector-pencils}
\end{equation}
Here \(E(t)\) is the equal/even sector associated with the deleted-pair
values \(00,11\), while \(O(t)\) is the unequal/odd sector associated with
\(10,01\).  A pencil means a one-parameter linear family \(U+tV\).
The corresponding kernel spaces
\[
 \Span\{E_{00},E_{11}\}
 \quad\text{and}\quad
 \Span\{E_{01},E_{10}\}
\]
are called the even and odd \emph{sector planes}; their displayed matrix-unit
axes are the \emph{monomial parity axes}.
For \(t\in\mu_m\), the native \(X\)-wires give the exact actual outputs
\(\sigma_t^{\mathsf R}O(t)\) and \(\sigma_t^{\mathsf S}E(t)\), since
\[
  XB_t^{\mathsf R}X
    =\sigma_t^{\mathsf R}X\mathsf R(t),
  \qquad
  XB_t^{\mathsf S}X
    =\sigma_t^{\mathsf S}\mathsf R(t).
\]
Thus \(t\in\mu_m\) gives an actual card with a known nonzero scalar.
We use \(A,B,C,D\) and evaluations outside \(\mu_m\) only to reconstruct
tensor identities.  Later \(K_E(t),K_O(t)\in\Mat_2(\mathbb C)\) denote
the single varying binary factor of \(E(t),O(t)\), respectively.

Write a nonsingular monomial binary, up to scalar, as
\begin{equation*}
  M_\alpha(a)(x,y)
  =\one[x\oplus y=\alpha]a^x,
  \qquad
  \alpha\in\Ftwo,\quad a\ne0.
\end{equation*}
Here \(\alpha\) is its support parity; reversal changes only the weight and
a nonzero scalar.

\subsection{Trace covers and alternating-cycle localization}
\label{proofsubsec:nd-trace-cover}

Consider two nonzero sampled cards of the same core.  Their matchings and
binary factor lines can disagree.  We will contract common parts while
preserving a disagreement, leaving only four or six ports to analyze.

For two perfect matchings \(M_0,M_1\), an \emph{alternating
\(M_0/M_1\)-cycle} is an even cycle in \(M_0\mathbin\triangle M_1\) whose
edges alternate between the two matchings.  We call a sequence of the
contractions below a \emph{trace cover}; it closes common edges and
shortens alternating cycles while retaining a chosen discrepancy.  Its
\emph{skeleton} records the contraction multigraph, edge orientations,
and external-port order.

\begin{definition}[Incompatible matching products]
\label{def:nd-live-card-disagreement}
Normalize the actual sampled outputs by their recorded nonzero scalars.
Two nonzero cards
\emph{disagree} if their residual matchings differ, or if their ordered
projective binary factor lines differ on at least two edges of the same
matching.  A parity change counts as a factor-line change.  Variation on
one designated edge, including its parity, is allowed; a single changed
factor is not by itself an obstruction for the ancestor.

Here a \emph{one-factor Segre ruling} is the projective family
\[
 [B_1\otimes\cdots\otimes B_{r-1}\otimes K\otimes
 B_{r+1}\otimes\cdots\otimes B_s]
\]
in which the matching and all factors except the displayed binary factor
\(K\) are fixed.  Thus two simple matching tensors lie on this ruling
exactly when they differ projectively in at most that one factor.
We use the simple/decomposable-tensor and Segre-line conventions recorded in
\cref{subsubsec:projective-tensor-geometry}.
\end{definition}

Consider two actual nonzero sector cards
\begin{equation}
  P_0=\bigotimes_{e\in M_0}B_{0,e},
  \qquad
  P_1=\bigotimes_{e\in M_1}B_{1,e},
  \label{eq:app-nd-two-products}
\end{equation}
where every \(B_{\epsilon,e}\) is a nonsingular monomial binary from the
normalized group.

  An aligned common edge
is \emph{untagged} when its two ordered factors agree projectively.  To
\emph{cross-close} two residual edges means to connect one endpoint of the
first to one endpoint of the second by an actual ordered group binary, so
that the two edges fuse in each product.  A determinant or parity
discrepancy is \emph{protected} when it is designated to survive these
closures; the proof below records that it is multiplied only by nonzero
trace, orientation, and fusion scalars.

\begin{definition}[Parity and coefficient discrepancies]
\label{def:app-nd-tc-tag}
After aligning the common edges of two sector products, a residual matched
factor is a \emph{parity discrepancy} if its support parity changes between the two products.  It is a \emph{coefficient discrepancy} if the parities agree but the two
ordered nonzero-entry vectors are not projectively proportional; in that
case the associated nonzero \(2\times2\) determinant is part of the tag.

\end{definition}

\begin{lemma}[Compatible closure and cycle shortening]
\label{lem:app-nd-compatible-closure}
Common untagged factors in \cref{eq:app-nd-two-products} can be closed by
actual group binaries while keeping both products nonzero.  An alternating
\(M_0/M_1\)-cycle can be shortened by two ports without cancellation, and
a protected coefficient determinant is only multiplied by a nonzero scalar.
Two discrepancies can be cross-closed by an actual connection while preserving
any distinguished remaining discrepancy.
\end{lemma}

\begin{proof}
For ordered monomial factors one has, up to the fixed orientation scalar,
\begin{align*}
 \sum_{x,y}M_\alpha(a)(x,y)M_\alpha(t)(x,y)&=1+at,\\
 \sum_zM_\alpha(a)(u,z)M_\beta(b)(z,v)
  &=b^\alpha M_{\alpha+\beta}
       \bigl(ab^{1-2\alpha}\bigr)(u,v).
\end{align*}
Thus on a common edge the two product traces are nonzero affine functions
of \(t\in\mu_m\); together they exclude at most two phases.  Since
\(m\ge3\), one actual phase keeps both nonzero.  Closing an edge then
splices the two incident edges in the other matching.  The displayed fusion
has a unique internal lift, so there is no cancellation, and every protected
minor is multiplied by nonzero trace and fusion factors.

For a tagged ordered edge \(q=(u_q,v_q)\), let
\(\alpha_{\xi,q}\in\Ftwo\) denote its support parity in product
\(\xi\in\Ftwo\).  Put
\[
 \delta_q:=\alpha_{0,q}+\alpha_{1,q}\in\Ftwo
 \qquad(q\in\{r,s\}),\qquad
 \alpha_{\xi,q}=\alpha_{0,q}+\xi\delta_q.
\]
One ordered connection \(M_\gamma(t)\) fuses the two factors in each
product.  Applying the displayed fusion identity twice gives
\[
 M_\alpha(a)M_\gamma(t)M_\beta(b)
 =t^\alpha b^{\alpha\oplus\gamma}
 M_{\alpha\oplus\gamma\oplus\beta}
 \!\left(a\,t^{1-2\alpha}b^{1-2(\alpha\oplus\gamma)}\right).
\]
For each supported exposed word the internal lift is unique, and the
fusion scalar is nonzero.  The fused parity difference is
\(\delta_r+\delta_s\).  A mixed pair therefore leaves a parity
discrepancy; equal types leave at most one coefficient discrepancy, which
may also cancel.  A distinguished discrepancy outside the two spliced
edges is preserved: its determinant is multiplied only by the nonzero
orientation and fusion scalars of the two products.  Each connection
removes two ports, whether or not the two spliced discrepancies cancel.
All operations use only the native \(X\)-wire, the explicit normalized
\(I\)-link,
and fixed actual group representatives.
\end{proof}

\begin{lemma}[Coefficient-line quaternary or genuine factor]
\label{lem:app-nd-coefficient-line-factor}
For distinct \(t_0,t_1\in\mu_m\), suppose the four nonzero actual cards
represented by \(E(t_0),E(t_1),O(t_0),O(t_1)\) at a deleted pair \(p\)
have a common residual matching edge \(e\), with the same support parity.
If all four ordered factors on \(e\) lie on one projective line, the
ancestor has the exact tensor factorization
\[
  f=B_e\otimes h
\]
with a nonsingular monomial binary \(B_e\).  Otherwise choose two cards
whose ordered factors on \(e\) are nonproportional.  If their other
matching edges and support parities agree, an actual closure retaining
\(p\cup e\) has a nonzero \(2\times2\) flattening minor.  Without this
parity compatibility no simultaneous trace is asserted.  The factor
identity is analytic; its factors are exposed only by factor saturation.
\end{lemma}

\begin{proof}
For the nonproportional pair with the stated compatibility, close every
other common matching edge successively.  At each step
the two traces exclude at most two roots, so
\cref{lem:app-nd-compatible-closure} keeps both products nonzero.  The two
surviving \(e\)-vectors remain nonproportional.  They are contractions of
the resulting actual quaternary by two kernels on \(p\); hence that quaternary has a
nonzero flattening minor.  The two kernels are linearly independent: two
distinct phases are independent inside one sector plane, while an even and
an odd kernel have different support parities.  Thus nonproportional
retained \(e\)-vectors really give flattening rank at least two.

For the factor alternative all four ordered \(e\)-factors lie on one line.  Choose its
representative
\[
 B_e(x_u,x_v)=\one[x_u\oplus x_v=\alpha]b_{x_u},
 \qquad b_0b_1\ne0,
\]
and absorb each card's nonzero scalar into its complementary tensor.  The
two-card inversions
\[
 A=\frac{t_1E(t_0)-t_0E(t_1)}{t_1-t_0},\quad
 D=\frac{E(t_1)-E(t_0)}{t_1-t_0},\qquad
 C=\frac{t_1O(t_0)-t_0O(t_1)}{t_1-t_0},\quad
 B=\frac{O(t_1)-O(t_0)}{t_1-t_0}
\]
put every slice in \(\mathbb C B_e\otimes V\), where \(V\) is the tensor
space on the residual ports outside \(e\).  Reassembling the four
slices gives the displayed tensor identity.
\end{proof}

\begin{lemma}[Two-card localization to six ports]
\label{lem:app-nd-localization}
Fix the deleted pair \(p\), and let two nonzero actual sampled
\(p\)-cards of the same core disagree.  There is an actual
one-copy contraction \(h\), retaining \(p\) and four further ports, whose
two inherited \(p\)-cards remain nonzero and either have different
residual matchings or differ on both edges of their common residual
matching.  The contraction closes only ports outside \(p\) and uses
\(O(\arity(f))\) fixed actual group representatives.  In particular,
\(h\) is not a three-factor current-group matching product.
\end{lemma}

\begin{proof}
Only the selected two cards are preserved in this construction.  The
contraction may depend on that pair of cards; no common phase choice for
four or more cards is required.

First suppose the two residual matchings differ.  Choose one alternating
cycle as distinguished.  On every other alternating cycle, close an edge
of the first matching by a kernel of the same support parity.  The trace
in the first product excludes at most one phase, whereas the other
product has a unique monomial lift.  Thus both products remain nonzero.
Repeat until that nondistinguished cycle becomes one common edge.

Every common edge outside the distinguished cycle can now be absorbed
into the cycle without any parity compatibility assumption.  Write that
edge as \(\{u,w\}\), select a cycle vertex \(v\), and cross-close
\(u,v\) by any actual monomial group kernel.  If the two cycle neighbours
of \(v\) are \(v_0,v_1\), the new matching edges are respectively
\(\{w,v_0\}\) and \(\{w,v_1\}\).  Thus the alternating cycle is
unchanged except that \(w\) replaces \(v\).  Both internal lifts are
unique and have nonzero weights, including when the old common edge had
different support parities in the two products.  Absorb every remaining
common edge in this way.  Finally shorten the distinguished cycle, using
the same trace-and-fusion operation, until it is a four-cycle.  Stop
before closing an edge of that four-cycle.  The two residual four-port
matchings are therefore still different.

Now suppose the residual matching agrees and at least two factor lines
differ.  Protect two such edges.  Every unprotected edge whose two support
parities agree can be traced directly: the two traces exclude at most
two phases, so \(m\ge3\) leaves an actual phase keeping both products
nonzero.  The two protected edges are untouched by these traces.

An unprotected edge with different support parities is instead
cross-closed with one protected edge.  The other protected edge remains
untouched.  If the chosen protected edge has equal parities but different
coefficient lines, the fused edge has different support parities and
remains a tag.  If both fused edges have different support parities, the
new support parities agree, but the ratio of their normalized fused
weights is
\[
  \kappa t^{\pm2},\qquad \kappa\ne0.
\]
This follows directly from the three-factor fusion identity in
\cref{lem:app-nd-compatible-closure}.  At most two phases make this ratio
equal to one.  Choose another actual phase; the fused edge is then a
coefficient tag.  In both cases the operation reduces the number of edges
by one and leaves two protected, distinct factor lines.  Continue until
exactly these two edges remain.

Perform the selected contractions on the original core, leaving the
ports of \(p\) open.  Tensor contraction commutes with either of the two
original \(p\)-kernels, so its inherited cards are precisely the two
nonzero products just constructed, with their recorded nonzero scalars.
Each operation removes two ports, giving the stated size bound.

If \(h\) were a product on three matching edges, deleting \(p\) would
either trace out one of its matching edges, leaving the other two factors
fixed, or fuse its two incident matching edges, leaving the third factor
fixed.  Whenever the two resulting cards are nonzero, their residual
matching is therefore the same and at most the fused factor can vary.
This contradicts either of the two constructed configurations.
\end{proof}

\begin{lemma}[Reduction to four ports or a tensor factor]
\label{lem:app-nd-bounded-consumer}
Every six-port residue in \cref{lem:app-nd-localization} has a complete
reduction through its proper factor and physical deck batches.  It cannot
close as a current-group matching product.  A quaternary child is processed
by its own complete factor and deck batch; a safe crossed factorization
completes that child only.
\end{lemma}

\begin{proof}
Let \(h\) be the actual six-port output, with its inherited nonzero cards,
port order, realization scalar, and originating frozen token retained.
By \cref{lem:app-nd-localization}, \(h\) is not a three-factor matching
product.  Adjoin \(h\), jointly factor-saturate, and close every proper
factor and every proper oriented deck card.  Odd or unary factors and
singular, unsafe, infinite-order, or outside-group binaries follow their
established exits.  Any proper nonbinary prime has arity at most four and
is processed by the lower-arity induction, including the arity-four
physical-deck boundary stated in the setup.  Thus this processing does
not use the six-port case being proved.

If no resolved leaf or strict successor occurs, all surviving proper
factors are current-group matching products.  A proper factorization of
\(h\) would therefore make \(h\) itself a matching product, a
contradiction.  Otherwise \(h\) is tensor-prime and every nonzero proper
pair card is a current-group matching product.  The independent six-port
argument in \cref{lem:nd-six-prime-base} excludes precisely this case.
That argument uses only the independent pencil lemmas and the Boolean
support atlas, not this consumer or a six-port recursive invocation.

For an actual quaternary child, first close its complete factor batch.
If it factors into safe current-group binaries, this completes the child
key.  Otherwise use the established binary exits or retain its prime
quaternary occurrence, regenerate the complete oriented deck, and apply
\cref{thm:equality-q4-boundary,cor:p1-eight-vertex-reduced-interface}.
Every returned factor batch is closed by
\cref{lem:actual-binary-completeness,lem:factor-saturation}.
An exact factorization of this child is never used to infer a
factorization of its ancestor.

In particular, for an active stable core of arity at least eight there
is an immediate contradiction: the first proper pair closure outside
\(p\) is zero or a matching product, and further monomial contractions
preserve that property.  Yet localization keeps the two disagreeing
inherited cards nonzero.  The same argument applies after the proper
lower-arity batches of a later occurrence have been closed.
\end{proof}

Apply this argument separately to each pair of nonzero samples.  The
closures need preserve only that pair, and may differ between comparisons.
On the continuing branch every two samples at a fixed deleted pair have
one residual matching and differ on at most one factor.  This pairwise
conclusion gives a common varying edge for the entire sample family:
if \(S,T\) differ on edge \(e\), a third sample differing from \(S\)
on another edge would differ from \(T\) on two edges.  If every sample is
proportional, designate any edge.  Hence all other factor lines agree
across both sectors.  Two nonzero phase samples in each nonzero sector
(which exist because \(m\ge3\)) then reconstruct its endpoint slices
with these common factors.  An identically zero sector has zero slices.
Thus any ancestor factor used below follows from an explicit four-slice
identity, independently of a safe factorization of a single child.

\begin{lemma}[Sector-ruling rigidity]
\label{lem:app-nd-sector-ruling}
Let \(U+tV\), \(t\in\mu_m\), \(m\ge3\), be a tensor pencil on an
even number of ports whose samples are zero or current
\(\mathcal M_m\)-matching products.  Its nonzero samples have one fixed
matching and all but at most one binary factor line fixed projectively.
In particular this holds for each of \(E(t)\) and \(O(t)\).
\end{lemma}

\begin{proof}
We use only the sampled matching hypothesis.  A proportional pencil has
one projective value.  Otherwise
its \(m\) samples are nonzero and projectively distinct.  Write
\(N=2^s\) for the support size of an \(s\)-edge matching product.
Two different matching supports \(S_0,S_1\) are affine cosets of the
same dimension, so \(|S_0\cap S_1|\le N/2\).  Every nontrivial linear
combination of the corresponding tensors has at least
\(2N-2|S_0\cap S_1|\ge N\) supported words outside the intersection.
For a third sample to have exactly \(N\) words requires equality and
cancellation on the whole intersection.  A single common word fixes that
linear combination projectively.  There is therefore at most one such
third matching point, contradicting \(m\ge4\).

For \(m=3\), every ratio of two nonzero coefficients within a matching
sample lies in \(\mu_3\).  Use a word of \(S_0\cap S_1\) to express
the required cancellation, and choose one word in each of
\(S_0\setminus S_1\) and \(S_1\setminus S_0\).  Their coefficient
ratio in the third sample lies in \(-\mu_3\), disjoint from \(\mu_3\).
This is again impossible.  Thus all sample supports coincide.

The weight-two words of the common support direction determine its
matching uniquely, including the parity of each support coset.  In that
fixed tensor-product presentation every \(2\times2\) Segre flattening
minor is a polynomial in \(t\) of degree at most two.  It vanishes at
all \(m\ge3\) sampled phases and hence identically.  The whole
projective line consists of simple tensors.  The Segre-line criterion
in \cref{subsubsec:projective-tensor-geometry} therefore permits variation
in only one factor, proving the assertion.
\end{proof}

\subsection{Endpoint directions and the parallelogram test}
\label{proofsubsec:nd-endpoints}

By the independent \cref{lem:app-nd-sector-ruling}, each sector has one
fixed residual matching and at most one varying factor line.  This
subsection classifies those binary factor pencils, discharges every zero or
proportional endpoint case, and isolates the complementary-axis
parallelogram survivor.

The M\"obius geometry of \cref{subsubsec:mobius-geometry} controls the
ratio of the two nonzero entries in a binary factor.  That ratio is a
fractional-linear function of the contraction phase.  Requiring it to
remain in the same finite phase group forces the two endpoint matrices
onto complementary coordinate axes.

\begin{lemma}[M\"obius stabilizer of a root polygon]
\label{lem:nd-root-mobius}
Let \(m\ge3\), and let the M\"obius map
\[
 \phi:\widehat{\mathbb C}\longrightarrow\widehat{\mathbb C},
 \qquad
 \phi:z\longmapsto\frac{c+dz}{a+bz},
 \qquad ad-bc\ne0,
\]
have no pole on \(\mu_m\) and permute \(\mu_m\).  Then
\[
 \phi(z)=\alpha z
 \qquad\text{or}\qquad
 \phi(z)=\alpha/z,
 \qquad \alpha\in\mu_m.
\]
\end{lemma}

\begin{proof}
A M\"obius map sends generalized circles to generalized circles.  Three
distinct roots and their images force \(\phi\) to preserve the unit circle.
Its action on the regular \(m\)-gon preserves or reverses cyclic order, so it
is a rotation or a rotation followed by reflection.  The corresponding
M\"obius maps agree with \(\phi\) on at least three points and hence agree
identically.
\end{proof}

The even and odd sector planes introduced above are precisely the two
monomial parity planes in \(\Mat_2(\mathbb C)\):
\[
 \Span\{E_{00},E_{11}\}
 \quad\text{and}\quad
 \Span\{E_{01},E_{10}\}.
\]
Their matrix-unit axes fix the following convention.  For a nonsingular
monomial matrix, its \emph{ordered nonzero-entry ratio} is the coefficient
of the second displayed axis divided by that of the first, in the fixed port
orientation.

\begin{lemma}[Classification of binary pencils]
\label{lem:nd-root-pencil}
Let \(A,B\in\Mat_2(\mathbb C)\) be binary signatures, and define the binary
matrix pencil
\[
 K(t):=A+tB,
 \qquad t\in\mu_m,\qquad m\ge3.
\]
Suppose every sampled binary signature \(K(t)\) is zero or a nonsingular
monomial binary whose ordered nonzero-entry ratio lies in \(\mu_m\).  Then
either \(A,B\) are proportional, or
there are a fixed parity plane with coordinate axes \(E_0,E_1\),
\(\lambda\ne0\), and \(\alpha\in\mu_m\) such that
\[
 K(t)=\lambda(E_0+\alpha tE_1)
 \qquad\text{or}\qquad
 K(t)=\lambda(tE_0+\alpha E_1).
\]
Here \(K(1)=A+B\) is a sampled value, whereas \(A=K(0)\) is an analytic
endpoint because \(0\notin\mu_m\).  In the application below, \(K(t)\) is
the one varying binary factor of a higher-arity actual matching card; it is
not thereby an independently realized binary gadget.
\end{lemma}

\begin{proof}
In the nonproportional case no sample is zero and the projective samples are
distinct.  Two of the at least three samples lie in one of the two monomial
parity planes; they span the pencil, so all samples lie in that plane.  In
these axes, write the matrix pencil with coefficient scalars \(a,b,c,d\)
and projective coordinate ratio \(\phi\) as
\[
 K(t)=(a+tb)E_0+(c+td)E_1,
 \qquad \phi(t)=\frac{c+td}{a+tb}.
\]
Nonproportionality gives \(ad-bc\ne0\), and nonsingularity excludes zeros
and poles on \(\mu_m\).  Thus \(\phi\) permutes \(\mu_m\).  By
\cref{lem:nd-root-mobius}, it equals \(\alpha t\) or \(\alpha/t\).
Cross-multiplication gives a degree-two identity valid on at least three
points.  Its coefficients yield respectively
\(b=c=0,d=\alpha a\), or \(a=d=0,c=\alpha b\), which are the two displayed
forms.  A zero sample already makes the endpoints proportional.
\end{proof}

Use the following names for the three binary-pencil outcomes of
\cref{lem:nd-root-pencil}:
\[
 \begin{array}{ccl}
 Z&:&\text{identically zero},\\
 P&:&\text{nonzero proportional},\\
 C&:&\text{nonproportional complementary-axis}.
 \end{array}
\]

For a deleted pair \(p\), factor the even-sector card into the common matching
factor \(P_E\) and endpoint matrix pencil \(K_E(t)\), after ordering the
\(2k=\arity(f)-2\) residual ports of a nonconstant even ruling:
\begin{equation*}
  E(t)=P_E\otimes K_E(t),
  \qquad
  K_E(t):=K_{E,0}+tK_{E,1},
  \qquad K_E(t),K_{E,0},K_{E,1}\in\Mat_2(\mathbb C),
\end{equation*}
with \(k-1\) fixed factors; define \(O(t)\) and its binary factor pencil
\(K_O(t)\) analogously.  To obtain one tensor identity with a fixed scalar, take
two distinct nonzero samples, use their linear inversion to reconstruct the
two analytic endpoint tensors, normalize the common factor \(P_E\) by a
fixed linear functional \(\eta_E(P_E)=1\), and define
\[
 K_E(t)=(\eta_E\otimes\operatorname{id})(E(t)).
\]
The normalization fixes the common factor once for the whole pencil.
The whole sector inherits the type of this exact residual binary pencil;
an identically zero whole pencil is type \(Z\), and a nonzero proportional
whole pencil is type \(P\) before any varying factor is selected.  In a
\(C/C\) branch the variable endpoints are complementary rank-one matrices.

\begin{lemma}[Direct six-port base case]
\label{lem:nd-six-prime-base}
Let \(h\ne0\) be a retained six-port signature in the standard
\(\mathcal M_m\) sector, where \(m\ge3\).  Suppose every actual oriented
pair card of \(h\), in both kernel parities and at every phase in
\(\mu_m\), is zero or a product of two nonsingular current-group
monomial binaries.  Then, after a port permutation, \(h\) has an exact
nontrivial factorization of type \(1+5\) or \(2+4\).  In particular,
\(h\) is not tensor-prime.
\end{lemma}

\begin{proof}
Apply \cref{lem:app-nd-sector-ruling,lem:nd-root-pencil} at each
deleted pair \(p\).  A type-\(P\)
sector is a scalar pencil times one fixed nonsingular matching product
on the four residual ports.  A type-\(C\) sector is one fixed
nonsingular binary factor tensored with a complementary-axis pencil;
its two endpoint supports are complementary singletons on the varying
edge.  A type-\(Z\) sector vanishes identically.

If a pair has type \(Z/P\), then \(h=U_p\otimes P\) for a nonzero
binary \(U_p\) and a nonzero residual matching product \(P\).  If it
has type \(Z/C\), the fixed binary factor of the \(C\) ruling factors
from \(h\).  The reversed cases are identical, and \(Z/Z\) would give
\(h=0\).  Each remaining case containing \(Z\) therefore proves the
assertion.

Suppose a pair has type \(P/P\).  Write
\[
 h=U_p\otimes P_0+V_p\otimes P_1,
\]
where \(U_p,V_p\ne0\) are supported on the two opposite parities of
\(p\), and \(P_0,P_1\) are nonsingular matching products on the four
residual ports.  Choose a residual pair \(e\) belonging to neither
matching.  Such a pair exists because the two matchings contain at most
four of the six residual pairs.  Contract \(e\) by either actual
parity kernel at phase \(1\).  This crosses the two factors of each
\(P_j\).  For every supported \(p\)-word, the output has precisely
two complementary residual words, each with one internal lift and a
nonzero coefficient.  Different \(p\)-words cannot cancel.  The
resulting actual quaternary therefore has
\[
 2\bigl(\lvert\supp U_p\rvert+\lvert\supp V_p\rvert\bigr)
\]
nonzero entries.  Six or eight entries contradict the card hypothesis.
Otherwise both supports are singletons.  Their opposite-parity
\(p\)-words differ in exactly one bit, so the other physical port of
\(p\) is fixed.  Its point mass is a unary tensor factor of \(h\).

For type \(P/C\), choose a residual pair \(e\) belonging to neither
the \(P\) matching nor the fixed/varying matching of the \(C\)
ruling.  Apply the same actual parity contraction.  Each supported
\(P\)-sector \(p\)-word leaves two residual words.  Each of the two
\(C\)-sector \(p\)-words leaves exactly one: the singleton varying
endpoint and the kernel parity force the bit on the fixed factor, and
then force both remaining residual bits.  These lifts are unique and
all coefficients are nonzero.  If the \(P\) sector has two
\(p\)-words, the card has six nonzero entries.  If it has one, the
card has four nonzero entries but its projection onto \(p\) has
three words.  A product of two nonsingular monomial binaries projects
onto any two ports in exactly two words when those ports are matched,
and in four words otherwise.  Both alternatives contradict the card
hypothesis.  Type \(C/P\) is identical.

Unless one of the stated factorizations has already occurred, every
deleted pair must consequently have type \(C/C\).  Fix one such pair.
Its fixed binary factors and complementary-axis endpoints give exactly
the eight-word four-flat support of
\cref{eq:cert-nd-q6-four-flat-support}.  The independent support atlas
\cref{prop:cert-nd-q6-support-lowering} either gives an injective
\(I/X\)-contraction with six or eight projected words, or identifies
a three-edge matching coset.  In the first case, injectivity gives one
nonzero summand per output word, so this actual quaternary contradicts
the card hypothesis.  In the second case, one parity sector at any
edge of that matching is identically zero, contradicting the conclusion
that every deleted pair is \(C/C\).  This completes the proof.

All contractions in this proof are applied directly to \(h\).
The factors are exact tensor factors, and are exposed only through
\cref{thm:external-factor,lem:factor-saturation} when the reduction
protocol uses them.
\end{proof}

Together with two-card localization, \cref{lem:nd-six-prime-base}
proves the higher-arity part of \cref{thm:normalized-dihedral}, as shown in
its proof.  The reconstruction results below describe the support and
coefficient constraints separately.

For a fixed deleted pair \(p=\{i,j\}\), a \emph{residual matching} of a
nonzero card
\[
 P=\bigotimes_{\{u,v\}\in\mathcal M_p}B_{uv}
\]
is the perfect matching \(\mathcal M_p\) of the residual ports
\([n]\setminus p\).  Its matching-indicator direction is
\[
 M(\mathcal M_p)=
 \Span\{e_u+e_v:\{u,v\}\in\mathcal M_p\}.
\]
The even and odd families are \emph{compatible at \(p\)} when their
residual matchings agree as unordered edge sets, equivalently when their
matching-indicator directions agree.  This concerns port positions only;
support-parity and ordered coefficient-line disagreements remain separate
discrepancies.  For example, on residual ports \(3,4,5,6\), the pairings
\(\{\{3,4\},\{5,6\}\}\) in both sectors are compatible, whereas that
pairing and \(\{\{3,5\},\{4,6\}\}\) form a matching disagreement whose
union contains an alternating four-cycle.  Matchings belonging to different
deleted pairs use different coordinate sets and are not asserted to be
literally equal; their global consistency is proved only in the support
integration below.

For the following support comparison suppose that both sector
pencils have type \(C\); all ordered pairs containing \(Z\) or \(P\) are
handled separately in \cref{lem:app-nd-endpoint-type} below.

Compare one actual nonzero even card with one actual nonzero odd card at the fixed
deleted pair \(p\).  If their residual matchings differ, or their ordered factor lines differ
on at least two edges (including parity changes), this is a cross-sector
disagreement and is consumed
by \cref{lem:app-nd-localization,lem:app-nd-bounded-consumer}.  Hence the
continuing branch has one common residual matching with all nondesignated
factor lines aligned.

If the two sector rulings have the same residual matching but designate
different varying edges, close only their common fixed factors.  Up to affine
translations, the remaining support, in the distinguished bits \(x,y\) and
the two varying-edge bits \(u,v\), is
\begin{equation*}
  \begin{aligned}
    x=y&\ \Longrightarrow\ u=x+\alpha,\quad v\text{ free},\\
    x\ne y&\ \Longrightarrow\ v=x+\beta,\quad u\text{ free}.
  \end{aligned}
\end{equation*}
An actual parity kernel imposing \(u=x+\alpha\) has a unique lift and leaves
a six-word physical quaternary, hence not a two-monomial product.  Its
support contains a complementary pair; actual logical \(X\)-dressings
may move that pair to the two endpoints.  Adjoin and
factor-saturate this physical quaternary.  A retained nonbinary factor has arity at
most four: for the active minimum core it is a strict lower-\(\nu\)
successor, and for a later occurrence it is assigned to the lower-arity
occurrence induction.  Stabilize that quaternary token, regenerate its complete
oriented deck, and only then invoke the equality-accessible quaternary
boundary and reduced interface.  Thus the different-edge branch is closed.
We may therefore assume below that both sectors use one common varying
edge.

After aligning their other common factors, the four analytic endpoint
supports have the form
\begin{equation*}
  \supp h_\xi=a_\xi+L,
  \qquad
  h_{00}=A,\ h_{01}=B,\ h_{10}=C,\ h_{11}=D,
\end{equation*}
where
\begin{equation*}
  L=\Span\{e_{u_r}+e_{v_r}:1\le r<k\},
  \qquad
  \dim L=k-1.
\end{equation*}
Define the endpoint displacement vectors
\begin{equation*}
  d_E=a_{00}+a_{11},
  \qquad
  d_O=a_{01}+a_{10}.
\end{equation*}
Then the actual even/odd card directions are
\begin{equation*}
  L_E=L+\Span_{\Ftwo}\{d_E\},
  \qquad
  L_O=L+\Span_{\Ftwo}\{d_O\},
  \qquad d_E,d_O\notin L.
\end{equation*}
The condition we need is
\begin{equation}
 a_{00}+a_{01}+a_{10}+a_{11}\in L
 \label{eq:app-nd-parallelogram}
\end{equation}
Equivalently,
\[
 d_E+d_O\in L,
 \qquad
 [a_{00}]+[a_{11}]=[a_{01}]+[a_{10}]
       \quad\text{in }\Ftwo^{n-2}/L,
 \qquad
 L+\Span_{\Ftwo}\{d_E\}=L+\Span_{\Ftwo}\{d_O\}.
\]
This condition is independent of the representatives
\(a_\xi\), since changing one representative changes the four-term sum by
an element of \(L\).

\begin{lemma}[Actual-card endpoint test]
\label{lem:app-nd-parallelogram}
In the complementary-axis \(C/C\) branch with one common varying edge, the
residual matchings of
the actual nonzero cards represented by \(E(t)\) and \(O(s)\) agree if
and only if
\cref{eq:app-nd-parallelogram} holds.
Failure is localized using the actual cards, not the endpoint slices.
\end{lemma}

\begin{proof}
A monomial matching is determined by its weight-two words, so the two
matchings agree exactly when the directions agree.  Since
\(d_E,d_O\notin L\),
\[
  L_E=L_O
  \quad\Longleftrightarrow\quad
  d_E+d_O\in L
  \quad\Longleftrightarrow\quad
  a_{00}+a_{01}+a_{10}+a_{11}\in L.
\]
Failure therefore invokes
\cref{lem:app-nd-localization,lem:app-nd-bounded-consumer} on the actual
cards and is resolved.
\end{proof}

\begin{lemma}[Zero and proportional endpoint pencils]
\label{lem:app-nd-exceptional-endpoints}
Every endpoint pair in which at least one sector has type \(Z\) or \(P\)
yields at least one of the following concrete outcomes:
\begin{enumerate}
\item the branch is impossible because it would give \(f=0\);
\item all four slices have one common residual tensor, giving an exact
      binary--residual factorization of \(f\);
\item all but one residual matching factor are common, giving an exact
      factorization \(f=P\otimes q\), where \(q\) is a nonzero quaternary
      factor on the deleted pair and the exceptional residual edge;
\item actual nonzero cards have a typed disagreement, which is localized and
      consumed by \cref{lem:app-nd-localization,lem:app-nd-bounded-consumer}.
\end{enumerate}
In item~3 the quaternary has matching-flattening rank two whenever one of
the two pencils has type \(C\).  The factor \(q\) is Turing-exposed through
the factor theorem and is a strict lower-\(\nu\) successor; it is not
silently treated as a directly realized gadget.
\end{lemma}

\begin{proof}
For a type-\(Z\) even pencil, two-card inversion gives \(A=D=0\), hence the
support relation \(x_i\oplus x_j=1\).  If the odd pencil also has type
\(Z\), then \(f=0\), contrary to the active nonzero prime.  If it has type
\(P\), the same inversion gives
\begin{equation*}
 f=G_p\otimes H_O,
  \qquad
  G_p=\begin{pmatrix}0&b\\ c&0\end{pmatrix},
\end{equation*}
where \(G_p\) is the displayed binary signature on \(p\) and \(H_O\) is the
complementary \((n-2)\)-ary signature.  This is an exact ancestor identity, and its factors
are exposed only by \cref{thm:external-factor}.  If the odd pencil has type
\(C\), write its exact ruling as \(O(t)=P_O\otimes K_O(t)\), where \(P_O\)
is the product of the \(k-1\) fixed residual factors and \(K_O(t)\) is the
nonproportional binary pencil on its varying edge \(e\).  Two-card inversion
puts both \(B\) and \(C\) in \(P_O\otimes\Mat_2(\mathbb C)\); the zero
slices \(A,D\) lie there as well.  Hence
\[
 f=P_O\otimes q_{p,e}.
\]
The complementary-axis endpoints of \(K_O\) are nonproportional, so the
\(p\mid e\) logical flattening of \(q_{p,e}\) has rank two.  Apply
    \cref{thm:external-factor,lem:factor-saturation}; the Turing-exposed quaternary is a
deck-stable lower-\(\nu\) successor, while every binary factor is sent through
\cref{lem:actual-binary-completeness}.  Exchanging the sectors covers every
type-\(Z\) case.  A single zero sample in a type-\(P\) pencil is included in
the exact proportional inversion and is never used by itself to infer a
factor.

For a proportional even pencil, inversion gives
\begin{equation*}
  A=aH_E,
  \qquad
  D=dH_E.
\end{equation*}
Choose a nonzero actual even card as the representative of the line
\(\mathbb C H_E\), and compare it with two nonzero actual odd cards.  Any
matching change or variation on at least two residual factor lines,
including parity changes, is a disagreement and is consumed by
\cref{lem:app-nd-localization,lem:app-nd-bounded-consumer}.  If no bounded
output occurs, the two sectors have one common residual matching and differ
on at most one designated edge \(e\).  Let \(P\) be the product of all fixed
factors outside \(e\).  Four-card inversion then puts
\(A,B,C,D\in P\otimes\Mat_2(\mathbb C)\), and therefore proves the exact
identity
\[
 f=P\otimes q_{p,e}.
\]
If no edge differs, absorb the last common factor into \(P\); this is the
stronger identity \(f=G_p\otimes H\).  If one edge differs, the remaining
quaternary is nonzero; in a \(P/C\) or \(C/P\) case its \(p\mid e\) flattening has
rank two because the \(C\)-endpoints are nonproportional.  In the \(P/P\)
case unequal projective edge lines give the same rank-two conclusion, while
equal lines give the stronger binary--residual factorization.  Only after
this exact identity has been proved are \(P\) and \(q_{p,e}\) exposed by
\cref{thm:external-factor,lem:factor-saturation}.  This covers \(P/P\),
\(P/C\), \(C/P\), and their possible zero samples.
\end{proof}

By \cref{lem:app-nd-exceptional-endpoints}, it remains to analyze two
nonproportional sector pencils.

Put the distinguished pair at
\(p=\{1,2\}\) and the residual ports at \(R=\{3,4,5,6\}\).  For each sector
\(\epsilon=x_1\oplus x_2\), let
\(M_\epsilon=\{e_\epsilon,e'_\epsilon\}\) be its residual matching, where
\(e_\epsilon=\{u_\epsilon,v_\epsilon\}\) is the fixed edge and
\(e'_\epsilon=(r_\epsilon,s_\epsilon)\) is the ordered varying edge.  There
are parameters \(a_\epsilon,b_\epsilon,o_\epsilon\in\Ftwo\) such that its
six-port support is parametrized by
\begin{equation}
 \begin{aligned}
 (x_1,x_2)&=(z,z+\epsilon),\\
 x_{u_\epsilon}\oplus x_{v_\epsilon}&=a_\epsilon,\\
 (x_{r_\epsilon},x_{s_\epsilon})
   &=(o_\epsilon+z,o_\epsilon+z+b_\epsilon),
 \end{aligned}
 \qquad \epsilon,z\in\Ftwo.
 \label{eq:app-nd-four-flat-support}
\end{equation}
For fixed \((\epsilon,z)\), the parity equation on the fixed edge has two
solutions and the varying edge is determined.  Hence each of the four
endpoint slices indexed by \((\epsilon,z)\) is a two-word affine flat, each
sector has four words, and the two sectors have the eight-word support used
in \cref{eq:app-nd-four-flat-support}.  The varying-edge assignments for
\(z=0,1\) are bitwise complements.  The coefficients on these slices need not be equal.

For a perfect matching \(M\) on \(2r\) ports and parities
\((c_e)_{e\in M}\), the associated \emph{matching coset} is
\[
 \{x\in\Ftwo^{2r}:x_u+x_v=c_{\{u,v\}}
       \text{ for every }\{u,v\}\in M\}.
\]

\begin{lemma}[A six-to-four contraction]
\label{lem:app-nd-six-to-four}
Suppose a normalized-dihedral localization leaves exactly the distinguished
pair \(p\) and two residual monomial pairs, with support as in \cref{eq:app-nd-four-flat-support}.  Then either the six-port support is a single three-edge
matching coset, or an actual \(I/X\)-contraction realizes a nonzero
quaternary of support size six or eight.  In the latter case, the output is
not a product of current-group monomial binaries and the protocol does not
re-enter a six-port state.
\end{lemma}

\begin{proof}
The explicit support in \cref{eq:app-nd-four-flat-support} is the Boolean
family exhaustively covered by
\cref{thm:cert-interface-nd-q6}.  That interface either identifies
the global matching coset or supplies a pair and parity \(\gamma\) with
injective projection.  Contracting it by the actual parity-\(\gamma\)
kernel---\(I\) for \(\gamma=0\) and \(X\) for \(\gamma=1\)---gives one
nonzero internal extension per word,
so there is no cancellation and the physical output has the certified
six- or eight-word support.  Adjoin this direct output and close its complete
factor batch.  A singular/unsafe factor, unary or odd factor, or a smaller
prime is dispatched by the stabilized protocol; a safe binary outside the
current group is a strict complete-group successor.  If the output itself
is tensor-prime, retain it as a lower-arity quaternary token, stabilize and
regenerate its complete oriented deck, and only then invoke the
equality-accessible quaternary boundary and reduced interface.  Since a product of two
nonsingular current-group monomial binaries has support size four, the
support-six/eight output cannot close as an unrecorded matching product.
This completes the reduction of the six-port tensor through its
four-port output and factor batch.
\end{proof}

\begin{lemma}[Reduction of the endpoint cases]
\label{lem:app-nd-endpoint-type}
For every deleted pair \(p\), all endpoint combinations containing type
\(Z\) or \(P\) are resolved.  On the remaining branch, the
endpoint pair is \(C/C\), \cref{eq:app-nd-parallelogram} holds, and there is one
compatible residual matching.
\end{lemma}

\begin{proof}
The \(Z/Z\) entry gives \(f=0\).  Every \(Z/P\), \(Z/C\), \(P/P\),
\(P/C\), and transposed entry is covered by
\cref{lem:app-nd-exceptional-endpoints}; the proof compares actual sample
cards and uses the same four-card inversion as
\cref{lem:app-nd-coefficient-line-factor}, so projective coincidence is
promoted to a factor only through an exact tensor identity.

It remains to take both types equal to \(C\).  A matching change or two
different ordered factor lines invokes the typed trace cover
\cref{lem:app-nd-localization,lem:app-nd-bounded-consumer}.  After closing
common untagged factors, two different variable edges give the physical
six-word quaternary constructed above; apply the lower-arity stabilized-quaternary route
given there.  For one common variable edge, failure of
\cref{eq:app-nd-parallelogram} is again an actual-card disagreement.  If the
parallelogram holds, then, because \(n\ge6\), retain the varying edge and one
fixed residual edge together with \(p\).  This is exactly the six-port
tensor of \cref{eq:app-nd-four-flat-support}.  Apply
\cref{lem:app-nd-six-to-four}.  Its nonmatching alternative is a physical
quaternary and follows the factor-saturate, stabilize, deck-regenerate, and
quaternary-boundary route in the proof of
\cref{lem:app-nd-six-to-four}; its sole
support-level survivor is a three-edge matching coset.  These alternatives
exhaust \(C/C\).
\end{proof}

\subsection{Global support integration and the Reed--Muller alternative}
\label{proofsubsec:nd-global-support}

The \(Z/P/C\) endpoint classification has resolved every case
containing \(Z\) or \(P\) and every
incompatible \(C/C\) card family.  The remaining survivor has \cref{eq:app-nd-parallelogram} and compatible residual matchings for every deleted
pair.  This subsection integrates those local supports into a global code,
identifies a global matching support, and excludes the sole Reed--Muller
alternative by the actual \(1-t^2\) obstruction.

\begin{lemma}[Endpoint-support saturation]
\label{lem:app-nd-endpoint-support-saturation}
Assume the normalized-dihedral analysis is on the remaining endpoint-
type branch.  Then, for each deleted pair
\(p=\{i,j\}\), there is a fixed-edge matching direction \(L_p\) on the
residual coordinates and endpoint differences \(d_{E,p},d_{O,p}\notin L_p\)
such that the four endpoint slices have full support on affine cosets
\(a_{ab}+L_p\), with every coefficient on each such coset nonzero.  Moreover
\[
 M_p:=L_p+\Span_{\Ftwo}\{d_{E,p}\}
     =L_p+\Span_{\Ftwo}\{d_{O,p}\}
\]
is the direction of the equal- and unequal-\(p\) actual cards; it is the
indicator code of a perfect matching on the residual coordinates and hence
\(M_p=M_p^\perp\).  This assertion is within the fixed deleted pair \(p\);
no literal equality between matchings on two different residual coordinate
sets is asserted.
\end{lemma}

\begin{proof}
On the remaining \(C/C\) branch,
\cref{lem:app-nd-sector-ruling} gives one fixed
nonzero monomial factor on every residual matching edge except the designated
variable edge.  Two samples with nonzero output and \cref{eq:app-nd-parallelogram} reconstruct
each slice by the four-card inversion used in
\cref{lem:app-nd-coefficient-line-factor}; hence every word in the resulting
affine matching coset has a nonzero coefficient.  More explicitly, the two
complementary-axis endpoint matrices have disjoint monomial supports, so each
residual word is supplied by exactly one endpoint; its coefficient is a
product of nonzero fixed monomial entries and a nonzero endpoint entry, with
no cancellation between endpoints.  If a word were absent, the
corresponding complementary-axis pencil would have a zero sample or exactly
one vanishing endpoint coordinate, producing a singular binary; this case
is resolved by
\cref{lem:app-nd-exceptional-endpoints,lem:app-nd-localization,%
lem:app-nd-bounded-consumer}.  A change of
matching or of a fixed factor likewise gives the typed trace-cover output.
Thus the only unconsumed possibility is the asserted common matching
direction.  For \(d_{E,p}=a_{00}+a_{11}\) and
\(d_{O,p}=a_{01}+a_{10}\), the parallelogram equation says
\(d_{E,p}+d_{O,p}\in L_p\), which identifies the two displayed descriptions
of \(M_p\) and the four offsets.  Matchings belonging to different deleted
pairs are integrated into one common code by
\cref{lem:app-nd-support-integration}, rather than compared directly here.
\end{proof}

We now regard the support as a set of binary words, using the projections
and affine operations of \cref{subsubsec:boolean-affine-tools}.  The
local matching equations will determine a self-dual binary code.  A
weight-two codeword identifies a global matched pair; if no such word
exists, the weight-four words identify the eight-coordinate Reed--Muller
code.  This reconstructs support only; coefficients are treated afterward.

\begin{lemma}[Self-dual Steiner reconstruction]
\label{lem:self-dual-steiner-reconstruction}
Let \(D\le\Ftwo^n\) be self-dual, where \(n>4\), and suppose that \(D\)
has no word of weight two.  If every three distinct coordinates lie in the
support of a unique weight-four word of \(D\), then \(n=8\); after
identifying the coordinates with \(\Ftwo^3\), one has \(D=RM(1,3)\).
\end{lemma}

\begin{proof}
The weight-four supports form a \emph{Steiner quadruple system}: a family of
four-element blocks in which every three coordinates belong to exactly one
block.  Self-orthogonality
of \(D\) makes two blocks meet evenly; uniqueness through triples then makes
the intersection of two distinct blocks either empty or a pair.

Fix a coordinate \(0\).  Define a binary operation \(+\) on the coordinates
by \(x+x=0\), \(x+0=x\), and, for distinct nonzero \(x,y\), let \(x+y\) be
the fourth point in the block through
\(0,x,y\).  This operation is commutative and every element has order two.
For the associativity check, define the auxiliary points \(u:=x+y\),
\(v:=y+z\), and \(w:=u+z\).  The symmetric
difference of the blocks \(\{0,x,y,u\}\) and \(\{0,y,z,v\}\) is the
weight-four codeword \(\{x,u,z,v\}\).  Symmetric difference with
\(\{0,u,z,w\}\) gives the block \(\{0,x,v,w\}\).  Uniqueness of the
block through \(0,x,v\) gives
\[
 (x+y)+z=w=x+v=x+(y+z).
\]
If an auxiliary point coincides---for example \(z=x+y\)---the same identity
reduces immediately to \(x+x=0\); thus the displayed distinct-block
calculation and these collision cases cover all inputs.  The coordinates form an
elementary abelian group \(V\cong\Ftwo^r\), and \(n=2^r\).  If
\(\{a,b,c,d\}\) is any block, symmetric difference with the block through
\(0,a,b\) shows \(d=a+b+c\); hence the blocks are exactly the affine
planes of \(V\).

Let \(B\) be the binary span of the block indicator vectors.  An affine plane is a translate of a two-dimensional
subspace of \(V\).  Orthogonality to every such plane is equivalent to
vanishing of every second additive derivative
\(\Delta_a\Delta_bF(x)\); hence \(B^\perp\) is exactly the space of truth
tables of affine-linear functions \(F:V\to\Ftwo\).  Thus
\(\dim B=2^r-r-1\).  Since \(B\subseteq D\) and
\(\dim D=2^{r-1}\), the inequality \(2^r-r-1\le2^{r-1}\), with \(r>2\),
forces \(r=3\); equality gives \(D=B=RM(1,3)\).
\end{proof}

\begin{lemma}[Support integration]
\label{lem:app-nd-support-integration}
Assume the remaining endpoint-type survivor, so that
\cref{eq:app-nd-parallelogram} holds for every deleted pair.  Then
\(\supp(f)=a+D\) for a self-dual binary linear code \(D\le\Ftwo^n\).
Exactly one of the following support alternatives remains:
\begin{enumerate}
\item \(D\) is the matching-indicator direction of one perfect matching on
      all \(n\) ports, so \(a+D\) is one \emph{global matching coset};
\item \(n=8\) and, after relabelling the ports, \(D=RM(1,3)\).
\end{enumerate}
The second alternative is incompatible with the normalized
\(\mu_m\)-dihedral deck for \(m\ge3\).  Consequently the sole
remaining support is one global matching coset.
\end{lemma}

\begin{proof}
Apply \cref{lem:app-nd-endpoint-support-saturation}.  Translation by the
\(00\)-slice gives, for any fixed deleted pair \(p\),
\begin{equation*}
  a+D,\qquad
  D=(\{00\}\times L_p)
    +\Span\{(01,a_{00}+a_{01}),
             (10,a_{00}+a_{10})\}.
\end{equation*}
The two displayed generators are independent modulo
\(\{00\}\times L_p\); hence \(\dim D=n/2\).  For an arbitrary deleted
pair \(p=\{i,j\}\), define the equal-coordinate slice
\(H_p:=\{x\in D:x_i=x_j\}\).  The endpoint saturation and
the parallelogram give
\[
  \pi_{\{i,j\}}:D\longrightarrow\Ftwo^2\ \text{surjective},\qquad
  \pi_{[n]\setminus\{i,j\}}(H_{\{i,j\}})=M_{\{i,j\}},
\]
where \(M_p=M_p^\perp\) is the residual perfect-matching direction from
\cref{lem:app-nd-endpoint-support-saturation}.  To see
self-orthogonality, take arbitrary \(x,y\in D\).  Among the \(n>4\)
coordinate columns
\((x_\ell,y_\ell)\in\Ftwo^2\), two are equal; relabel them \(i,j\).  Then
\(x,y\in H_p\), so their residual projections lie in \(M_p\) and have zero
mutual dot product.  Their contribution on the deleted pair is
\(x_i y_i+x_j y_j=0\).  Hence \(x\cdot y=0\), proving
\(D\subseteq D^\perp\).  Since \(\dim D=n/2\), equality follows:
\(D=D^\perp\).

Suppose first that \(e_u+e_v\in D\).  Self-duality implies
\(x_u=x_v\) for every \(x\in D\), so \(H_{\{u,v\}}=D\).  Put
\(E=[n]\setminus\{u,v\}\).  Endpoint saturation gives
\[
 \pi_E(D)=M_{\{u,v\}},
\]
the indicator direction of the residual perfect matching.  For
\(y\in M_{\{u,v\}}\), choose a lift \(x\in D\).  Subtracting
\(x_u(e_u+e_v)\) gives the zero extension \(\iota_E(y)\in D\).  Hence
\[
 D=\Span_{\Ftwo}\{e_u+e_v\}
   \oplus\iota_E(M_{\{u,v\}}),
\]
which is exactly the indicator direction of the global perfect matching
\(\{\{u,v\}\}\cup\mathcal M_{\{u,v\}}\).  This is a support identity
only: it does not yet assert a tensor factorization of the coefficients.

It remains that \(D\) has no weight-two word.  Fix \(i\ne j\) and put
\(E=[n]\setminus\{i,j\}\).  The projection
\(H_{\{i,j\}}\to\Ftwo^E\) is injective, because its only possible nonzero kernel
is \(e_i+e_j\).  The functional \(x\mapsto x_i+x_j\) is nonzero on \(D\),
since otherwise \(e_i+e_j\in D^\perp=D\).  Thus
\(\dim H_{\{i,j\}}=n/2-1\), and endpoint saturation identifies its image with
\(M_{\{i,j\}}\).  For a third coordinate \(\ell\), the residual matching
pairs \(\ell\) with a unique \(q\).  The unique lift of
\(e_\ell+e_q\in M_{\{i,j\}}\) cannot have \(00\) on \(i,j\), since that would
be a weight-two word.  It is therefore
\(e_i+e_j+e_\ell+e_q\in D\).  Hence every three coordinates lie in a unique
weight-four support, and
\cref{lem:self-dual-steiner-reconstruction} gives \(n=8\) and
\(D=RM(1,3)\).

To exclude this residue, identify the eight coordinate positions with
\(\Ftwo^3\) and, after an affine relabelling, take the deleted pair to be
\(\{000,001\}\).  On its equal-coordinate slice the residual matching is
\[
 \{010,011\},\qquad \{100,101\},\qquad \{110,111\}.
\]
If \(y_1,y_2,y_3\) are the three logical matching bits, the common bit on
the deleted pair is
\[
 b=y_1+y_2+y_3.
\]
Contract the deleted pair by the actual sector phase \(t\).  The projection
of the equal slice is injective, so every residual support word has one
internal lift and its coefficient is multiplied by \(t^b\).  Fixing
\(y_3\) and inspecting the \(2\times2\) face indexed by \(y_1,y_2\), the
phase exponents are \((0,1,1,0)\), or their complement.  Thus the face has
one of the raw forms
\[
 \begin{pmatrix}A&tB\\tC&D\end{pmatrix},
 \qquad
 \begin{pmatrix}tA&B\\C&tD\end{pmatrix},
 \qquad A,B,C,D\ne0.
\]
The actual \(t=1\) card lies in the fixed Segre ruling, so \(AD=BC\).
Nonzero row and column scalings therefore give, for some
\(\varepsilon\in\Ftwo\), the normalized physical block
\begin{equation}
 \begin{aligned}
  H_t&=\kappa\diag(r_0,r_1)C_\varepsilon(t)\diag(s_0,s_1),\\
  C_0(t)&=\begin{pmatrix}1&t\\t&1\end{pmatrix},
  &C_1(t)&=\begin{pmatrix}t&1\\1&t\end{pmatrix},\\
  \det H_t&=(-1)^\varepsilon
      \kappa^2r_0r_1s_0s_1(1-t^2).
 \end{aligned}
 \label{eq:app-nd-RM-minor}
\end{equation}
with all scalar factors nonzero.  If an orientation
produces \(t^{-1}\), clear its nonzero monomial factor.  Since \(m\ge3\),
some actual \(t\in\mu_m\setminus\{\pm1\}\) exists, and the displayed
nonzero minor contradicts the fixed Segre ruling.
\end{proof}

Hence the support survivor is one global matching coset.  A weight-two word
above records a support pair only; no tensor factor is inferred until the
coefficients have also separated.  The following two stages first close
every degenerate six-port coefficient table and then prove that the
coefficients on the global matching are multiplicative.

\subsection{The six-port coefficient calculation}
\label{proofsubsec:nd-six-port}

We now return to the matching-supported six-port coefficient cases retained
by the support analysis.  Trace-cover six-port disagreements have already been
reduced to quaternary and resolved; the only survivor is the fixed
three-edge-matching coefficient pattern.  This subsection first dispatches
all degenerate coefficient cases and then analyzes the full-support table.

For a quaternary contraction whose support is contained in two fixed residual
matching edges, write its logical coefficient table as
\[
 Q=\begin{pmatrix}a&b\\c&d\end{pmatrix}.
\]
We include the cases \(0<|\supp Q|<4\).  A \emph{support-matching disagreement} means that its physical
support uses a residual pairing different from the fixed pairing; this is
the combinatorial disagreement of the trace-cover reduction.  By contrast,
\emph{not \(G\)-matching} means that the signature has no tensor
presentation by current-group binary factors.  These two notions are not
interchangeable.

For each directly realized quaternary, first adjoin it with its port order, scalar, and
provenance; jointly factor-saturate it; test all three \(2+2\)
flattenings and dispatch every exposed factor; make an actual endpoint
dressing; and invoke \cref{lem:finite-q4-dispatch} only after tensor
primality, endpoint nondegeneracy, non-\(G\)-matching status, and both
crossed rank hypotheses have been verified.

\begin{lemma}[Six-port coefficient dispatch]
\label{lem:nd-six-coefficient-dispatch}
Let \(h\) be a direct six-port card whose support is exactly the
eight-word coset of a fixed three-edge matching.  Before a quotient of its logical coefficients is
formed, close every already constructed four- or six-port batch as follows:
\begin{enumerate}
\item a zero quaternary is converted into an exact factorization of its parent six-port tensor by
      the proportional-slice identity below;
\item a support-matching disagreement is sent to
      \cref{lem:app-nd-localization,lem:app-nd-bounded-consumer};
\item for a quaternary on the two retained matching edges, test \(Q\): if
      \(\det Q=0\), factor it exactly; if \(\det Q\ne0\), use the
      quaternary reduction just described;
\item jointly factor-saturate every output and dispatch every unary,
      singular, unsafe, infinite-order, or new safe binary factor;
\item if no preceding concrete output applies, retain \(h\) only when all
      eight logical coefficients are nonzero.
\end{enumerate}
A quaternary with \(0<|\supp Q|<4\) enters item~3.
Every noncurrent factor or nonmatching quaternary outcome gives a resolved leaf or
strict successor.  If a nonzero child quaternary factors into two current-group
binaries, close that child, record its rank-one equation, and continue the
finite analysis of the same originating six-port tensor.
\end{lemma}

\begin{proof}
Let \(u_0,u_1\) be the two standard-basis binary signatures supported on
the two allowed states of a selected physical matching edge.  Write
\[
 h=u_0\otimes H_0+u_1\otimes H_1.
\]
An actual phase kernel has two nonzero coefficients
\(\alpha_t,\beta_t\), and its quaternary contraction is
\[
 q_t=\alpha_tH_0+\beta_tH_1,
 \qquad \alpha_t\beta_t\ne0.
\]
If \(q_t=0\), then
\[
 \alpha_tH_0+\beta_tH_1=0.
\]
If \(H_0=0\), then \(H_1=0\), so \(h=0\) and its zero-parent token is
completed directly.  Otherwise \(H_0\ne0\) and
\[
 H_1=-\frac{\alpha_t}{\beta_t}H_0,
 \qquad
 h=\left(u_0-\frac{\alpha_t}{\beta_t}u_1\right)\otimes H_0.
\]
Thus the zero child is not itself declared a terminal: it certifies an
exact binary--quaternary factorization of its parent.  Close the complete
factor batch by \cref{thm:external-factor,lem:factor-saturation}.  A
nonbinary factor gives a quaternary or lower-\(\nu\) successor; unary,
singular, unsafe, or infinite-order factors enter their established
reductions; and a new safe binary strictly enlarges the complete group.  If
all prime factors are already safe current-group binaries, the parent is a
current-group matching product and its frozen token is completed.
Completing a zero child alone never completes an unresolved multi-copy
parent without this factor-batch argument.

Now let \(q_t\ne0\) have logical table \(Q\) on the two retained matching
edges.  If \(\det Q=0\), then \(Q\) has rank one and the physical quaternary has an
exact \(2+2\) factorization.  When \(0<|\supp Q|<4\), at least one
logical factor has a zero coordinate, so one exposed binary is singular.
If both exposed factors are safe current-group binaries, close this child
quaternary, retain the equation \(\det Q=0\), and resume the coefficient analysis
of the same parent six-port tensor; the child factorization alone does not complete the
parent token.
If \(\det Q\ne0\), its fixed-matching flattening has rank two.  Adjoin the
physical card with its port order, scalar, and provenance.  If \(Q\) is
supported on exactly two entries, nonzero determinant forces
them to be opposite logical entries.  Actual logical \(X\)-translations
then make the quaternary a nonzero pure generalized equality, and
\cref{lem:pure-ge} applies.  If exactly three entries are nonzero, actual
logical flips put a nonzero opposite pair at the two endpoints, so the card is
endpoint-nondegenerate.  After factor saturation and removal of every
crossed rank-one presentation, it is tensor-prime, non-\(G\)-matching, and
has both crossed ranks at least two.  The full-support case is
endpoint-nondegenerate after the same endpoint choice.  Thus the
quaternary reduction just described is licensed; invoke
\cref{lem:finite-q4-dispatch,thm:equality-q4-boundary,%
cor:p1-eight-vertex-reduced-interface} only after their hypotheses have been
verified.  A support-matching disagreement instead enters the typed
trace-cover reduction.  These routes also process every exposed factor by
\cref{lem:actual-binary-completeness}.

For the matching-supported six-port tensor retained from
\cref{lem:app-nd-support-integration}, the support is already the complete
eight-word global-matching coset, so its eight logical coefficients are
nonzero.  The case analysis does not pre-enumerate every possible later phase
contraction of this full-support table; it supplies the concrete reduction
whenever the calculation below, or a later mixed two-copy batch, produces
one.  We may therefore write its coefficients as \(c_{abc}\),
\(a,b,c\in\Ftwo\), and normalize \(c_{000}=1\).
\end{proof}

\begin{lemma}[One-coordinate six-port minor localization]
\label{lem:nd-q6-minor-localization}
Let \(q=(q_{ab\xi})_{a,b,\xi\in\Ftwo}\) be the logical coefficient table of
a physical six-port tensor whose support is contained in a fixed three-edge matching.
Close every quaternary and factor batch produced below by the
quaternary reduction just described and the binary case analysis.  On any branch
with no resolved leaf or strict stable successor, fully close every factor
batch.  Each resulting physical six-port tensor is then either zero, a current-group
matching product, or has rank at most one in each of
\[
 a\mid(b,\xi),\qquad b\mid(a,\xi),\qquad \xi\mid(a,b)
\]
The first two alternatives also have all three ranks at most one.  Their
rank equations are recorded before their child tokens are completed, so
the conclusion holds at every physical phase assignment in any later
two-copy batch to which this lemma is applied.
\end{lemma}

\begin{proof}
Consider \(a\mid(b,\xi)\).  Close the \(\xi\)-edge by every actual phase,
obtaining the logical matrix pencil
\[
 A+tB
\]
in the \(a,b\) coordinates.  If \(\det(A+tB)\ne0\) at an actual phase, the
resulting physical quaternary enters the preceding nonzero-minor reduction.  If the
determinant vanishes for every \(t\in\mu_m\), it is a degree-two polynomial
with at least three roots and is therefore identically zero.

The two-dimensional rank-one-line classification is elementary.  Write
nonzero rank-one endpoints as \(A=uv^{\mathsf T}\) and
\(B=xy^{\mathsf T}\).  The mixed coefficient of
\(\det(A+tB)\) is, up to sign,
\(\det[u,x]\det[v,y]\).  Hence either \(u,x\) are proportional or
\(v,y\) are proportional; zero or proportional endpoints satisfy the same
conclusion directly.  In the first case \(A,B\) have one common
\(a\)-factor, so the \(a\mid(b,\xi)\) flattening of \(q\) has rank at most
one.  In the second case they have one common \(b\)-factor, giving an exact
binary factor of \(q\); factor-saturate and dispatch that factorization.
If closing this factorization gives a resolved leaf or strict stable
successor, the parent branch ends.  Otherwise all exposed factors are
current-group matching factors, so \(q\) itself is a current-group matching
product and all three logical flattenings have rank at most one; record
these equations before completing its child token.  Thus every branch
continuing the parent analysis has the stated \(a\)-rank equation.
Repeating the same argument after permuting \(a,b,\xi\) proves the other two
rank bounds.
\end{proof}

After \cref{lem:nd-six-coefficient-dispatch}, all eight logical
coefficients on the continuing branch are nonzero.  Retain their notation and
normalize the remaining global-matching six-port support to
\(X^{(1,2)}\otimes X^{(3,4)}\otimes X^{(5,6)}\), where the superscript records
the ordered pair of ports on which the binary disequality tensor \(X\) acts.
For each selected pair of matching edges, use both actual parity-zero and
parity-one kernels.  The parity-zero contraction on the first two edges
gives
\[
 c_{000}c_{111}=c_{001}c_{110}.
\]
The parity-one contractions on the three choices of edge pair give the
remaining equalities among
\[
 c_{001}c_{110},\qquad
 c_{010}c_{101},\qquad
 c_{011}c_{100}.
\]
Every nonzero determinant is routed through the quaternary reduction just described.  Hence, on the continuing branch, all four products are equal:
\begin{equation}
  c_{000}c_{111}
  =c_{001}c_{110}
  =c_{010}c_{101}
  =c_{011}c_{100}.
  \label{eq:app-nd-six-relations}
\end{equation}
Define
\[
  \alpha=\frac{c_{100}}{c_{000}},\qquad
  \beta=\frac{c_{010}}{c_{000}},\qquad
  \gamma=\frac{c_{001}}{c_{000}}.
\]
These three quantities are exposed by explicit physical factor cards.
For \(\alpha\), join ports \(3\) and \(5\), the first ports of the second
and third matching edges, by the actual parity-zero kernel
\(M_0(1)=I\).  The remaining logical table, with the first matching edge
as its row edge, is
\[
 \begin{pmatrix}
 c_{000}&c_{011}\\
 c_{100}&c_{111}
 \end{pmatrix}.
\]
Its determinant vanishes by \cref{eq:app-nd-six-relations}, so exact factor
extraction exposes on the first matching edge a binary with ordered ratio
\(\alpha^{\pm1}\).  Cycling the three matching edges exposes
\(\beta^{\pm1}\) and \(\gamma^{\pm1}\) in the same way.  Every
noncurrent or unsafe factor is a binary exit; hence on the continuing
branch
\[
 \alpha,\beta,\gamma\in\mu_m.
\]

Fuse the fixed actual representatives with transfers
\(\mathsf R(\alpha^{-1})\),
\(\mathsf R(\beta^{-1})\), and
\(\mathsf R(\gamma^{-1})\) onto the three oriented matching edges, and
divide by the recorded common scalar \(c_{000}\).  Equation
\cref{eq:app-nd-six-relations} then gives
\[
 N_r(z)=
 \begin{cases}
  1,&\wt(z)\le1,\\
  r,&\wt(z)\ge2,
 \end{cases}
 \qquad
 r:=\frac{c_{000}c_{110}}{c_{100}c_{010}}.
\]
It remains to prove that \(r\) is also a group phase.  On this physically
normalized six-port card, contract ports \(1\) and \(3\), the first ports
of the first and second matching edges, by the actual parity-one card
\(M_1(u)=M_{\alpha=1}(a=u)\), \(u\in\mu_m\); in the displayed port
orientation this is
\[
 M_1(u)=\begin{pmatrix}0&1\\ u&0\end{pmatrix}.
\]
Write the physical bits on the six ports as
\(x_1,\ldots,x_6\); the support matching means
\(x_2=1-x_1\), \(x_4=1-x_3\), and \(x_6=1-x_5\).
The card joining ports \(1\) and \(3\) therefore imposes
\(x_1\oplus x_3=1\) and contributes the factor \(u^{x_1}\).
After this contraction, use \(x_2\) (with \(x_4=1-x_2\)) as the row bit
and \(x_5\) (with \(x_6=1-x_5\)) as the column bit.  Ordering the rows and
columns changes only the recorded nonzero representative scalar or reverses
a port.
The resulting quaternary table on ports
\(2,4,5,6\), up to reversal and a nonzero representative scalar, is
\[
 \begin{pmatrix}1&r\\ u&ur\end{pmatrix}
 =\begin{pmatrix}1\\u\end{pmatrix}
  \begin{pmatrix}1&r\end{pmatrix}.
\]
This is an exact tensor factorization, and its two binary factors have
ordered ratios \(u^{\pm1}\) and \(r^{\pm1}\).  Expose them through
\cref{thm:external-factor,lem:factor-saturation} before testing the complete
binary group.  A noncurrent factor closes through the active protocol or
the selected occurrence's frozen local queue.  Thus the coefficient
calculation continues only when both factors are safe in the current
\(\mu_m\)-dihedral group.  On that branch,
\(r^{\pm1}\in\mu_m\), and hence \(r\in\mu_m\).  Thus the normalized
logical coefficient function above is
\begin{equation}
  N_r(z)=
  \begin{cases}
    1,&\wt(z)\le1,\\
    r,&\wt(z)\ge2,
  \end{cases}
  \qquad
  z\in\Ftwo^3,
  \quad
  r\in\mu_m.
  \label{eq:app-nd-Nr}
\end{equation}
Let \(Q_r\) denote the \(2\times2\) matrix of the quaternary obtained by
closing the third pair, displayed below.  For \(r\ne-1\), its full-support
cross-ratio is
\[
 \operatorname{cr}(Q_r)=\frac{(Q_r)_{00}(Q_r)_{11}}
                 {(Q_r)_{01}(Q_r)_{10}}
          =\frac{4r}{(1+r)^2}.
\]
If this full-support quaternary enters an affine or local-affine tractable branch,
\cref{lem:p1-eight-vertex-local-affine-collapse} reduces the latter to the
affine four-point plane, namely a coset of a two-dimensional subspace of
\(\Ftwo^4\).  Clause~\textnormal{(ii)} of
\cref{thm:cert-interface-affine-q4-gadget-exits} gives
cross-ratio \(\pm1\) on this
plane.  The value \(+1\) is equivalent here to
\(\det Q_r=0\), hence to \(r=1\), which is the product case removed below.
Therefore a nonproduct affine residue has \(\operatorname{cr}(Q_r)=-1\), giving
\(r+r^{-1}=-6\).  Every other outcome is closed by the physical
quaternary reduction.  Closing the third
pair gives
\begin{equation}
  Q_r=
  \begin{pmatrix}
    2&1+r\\
    1+r&2r
  \end{pmatrix},
  \qquad
  \det Q_r=-(r-1)^2.
  \label{eq:app-nd-Qr}
\end{equation}
This is a direct physical quaternary with its inherited port order and nonzero
representative scalar.  Adjoin it and factor-saturate its complete batch
before applying the quaternary reduction.  If \(r=1\), then
\(N_1\equiv1\), so the parent six-port tensor itself is the product of its
three matching binaries; complete the originating frozen token after the
binary dispatch.  If \(r\ne1\), the displayed fixed-matching minor is
nonzero, and the quaternary reduction of
\cref{lem:app-nd-bounded-consumer,lem:nd-six-coefficient-dispatch} applies.
Two actual \(I\)-dressings change its paired-disequality support to
\(\{0000,0011,1100,1111\}\), so every nonproduct retained quaternary is
endpoint-nondegenerate.  For the exceptional value \(r=-1\), the equality exit is the
following explicit two-copy wiring.  Define the six-port signature
\[
 F_r(z_1,\ldots,z_6)=
 \prod_{j=1}^3\one[z_{2j-1}\oplus z_{2j}=1]\,N_r(z_1,z_3,z_5),
\]
where the three displayed logical bits are the first ports of the three
paired-disequality edges.  Take two copies, leave ports \(1,2\) of each copy
external, and join corresponding ports \(3,4,5,6\) by four actual equality
links.  The resulting quaternary tensor is
\begin{equation}
 W_r(a,b,c,d)=\sum_{u,v,w,z\in\Ftwo}
 F_r(a,b,u,v,w,z)F_r(c,d,u,v,w,z).
 \label{eq:app-nd-r-minus-one-wiring}
\end{equation}
Direct substitution into \(N_{-1}\) gives
\[
 W_{-1}(a,b,c,d)=4\,\one[(a,b,c,d)\in\{0101,1010\}].
\]
Applying an actual \(X\)-translation on the second external port of each
copy sends these two words to \(0000,1111\), so the output is
\(4\Eq_4\) in the displayed normalization (any previously recorded gadget
scalar merely changes this to a known nonzero multiple).  This terminates by
\cref{cor:higher-equality-unary,lem:common-unary-exit}.  Otherwise the
affine cross-ratio test would force
\(r+r^{-1}=-6\), impossible for a root of unity.  Hence every nonproduct
case not already terminated enters
\cref{cor:p1-eight-vertex-reduced-interface}.

\subsection{Polarized coefficient integration}
\label{proofsubsec:nd-polarization}

The remaining tensor has a fixed matching support and nonzero logical
coefficients.  We test their factorization by the following two-copy
construction, called \emph{polarization}.  Each phase assignment gives an
actual six-port tensor.  Its vanishing minors are polynomials in those
phases; interpolation turns the finite family of physical tests into
identities among the original coefficients.

Let \(w:\Ftwo^s\to\mathbb C^\times\), \(s\ge3\), be the logical coefficient
table on the global matching, with every matching edge in its fixed
orientation.  Fix distinct logical coordinates \(i,k\) and take two equally
oriented copies.  Leave the entire \(i\)-edge of each copy external.  On the
two \(k\)-edges, join one pair of corresponding physical ports by an actual
equality link and leave the other two ports external; after one recorded
\(X\)-dressing they form the third external matching edge and carry one
common logical bit \(\xi\).  For every \(\ell\notin\{i,k\}\), close the two
\(\ell\)-edges by one equality link and one phase-dressed equality link.
Unique monomial lifts give, after division by the recorded nonzero
representative scalar, the direct six-port coefficient
\begin{equation}
  q^{\,i,k}_{ab\xi}(\lambda)
  =\sum_z\lambda^z w_i(a,\xi,z)w_i(b,\xi,z),
  \qquad
  \lambda^z=\prod_{\ell\notin\{i,k\}}\lambda_\ell^{z_\ell}.
  \label{eq:app-nd-polarized-card}
\end{equation}
Here \(w_i(a,\xi,z)\) means that coordinates \(i,k\) of \(w\) are
\(a,\xi\), respectively, and the other coordinates are ordered as \(z\).
The fixed realization scalars may replace \(\mu_m\) in each variable by a
nonzero coset \(c_\ell\mu_m\); the interpolation lemma below is stated in
that form.  Their Cartesian product is the \emph{polarization grid}.  All
Laurent monomials and exponent arithmetic use the conventions of
\cref{subsubsec:laurent-polynomials}.

\begin{lemma}[Grid interpolation for bounded Laurent degree]
\label{lem:nd-grid-interpolation}
Let \(m\ge3\), let \(c_1,\ldots,c_r\in\mathbb C^\times\), and let
\(P(\lambda_1,\ldots,\lambda_r)\) be a Laurent polynomial such that,
after multiplying by a monomial, its degree in each variable is at most
two.  If \(P\) vanishes on the product of cosets
\(\prod_{j=1}^r c_j\mu_m\), then \(P\) is identically zero.
\end{lemma}

\begin{proof}
After the monomial shift, substitute \(\lambda_j=c_jt_j\).  Fix all
variables except the last one.  The result is a univariate polynomial of
degree at most two with at least three distinct roots \(t_j\in\mu_m\),
hence is zero.  Induction on the number of variables proves that every
coefficient polynomial vanishes.  The nonzero monomial used to clear the
Laurent exponents never affects a value on the stated cosets.
\end{proof}

\begin{lemma}[Polarized determinant]
\label{lem:app-nd-polarized-determinant}
Assume the global matching support has been fixed by
\cref{lem:app-nd-support-integration}, and use
\cref{lem:nd-q6-minor-localization} to retain every one-coordinate rank
equation across completed current-group factor batches.  Form every physical card
\(q^{\,i,k}(\lambda)\) in \cref{eq:app-nd-polarized-card}.  Suppose that,
after the concrete six-port case analysis has removed every violating batch,
each retained card has rank at most one in each of its three
one-logical-coordinate flattenings
\[
 a\mid(b,\xi),\qquad b\mid(a,\xi),\qquad \xi\mid(a,b).
\]
Any concrete card violating this hypothesis is resolved before
this lemma is invoked
by \cref{lem:app-nd-bounded-consumer,lem:nd-six-coefficient-dispatch,%
lem:actual-binary-completeness}.  Then
\begin{equation}
  w(z)=w(0)\prod_{i=1}^s\rho_i^{z_i}
  \label{eq:app-nd-coefficient-product}
\end{equation}
for nonzero scalars \(\rho_i\).  Consequently the core is a tensor product
of binary signatures on the global perfect matching.  This is an exact
tensor identity; the individual factors are made available only through
\cref{thm:external-factor,lem:factor-saturation} and are then sent through
the binary case analysis.
\end{lemma}

\begin{proof}
Fix \(i\ne k\) and \(\xi\in\Ftwo\), and abbreviate
\[
 u_z:=w_i(0,\xi,z),\qquad v_z:=w_i(1,\xi,z).
\]
Rank at most one of the \(a\mid b\) face, including a zero card, gives the
determinant polynomial
\begin{equation*}
  \Delta_\xi(\lambda)
  :=
  \left(\sum_z\lambda^zu_z^2\right)
  \left(\sum_z\lambda^zv_z^2\right)
  -
  \left(\sum_z\lambda^zu_zv_z\right)^2.
\end{equation*}
This polynomial vanishes on the full product of phase cosets.  Its coordinate
degree is at most two, so \cref{lem:nd-grid-interpolation} makes it the zero
polynomial.  If \(z,z'\) are adjacent in the remaining \((s-2)\)-cube, the
coefficient of \(\lambda^{z+z'}\) is
\begin{equation*}
  (u_zv_{z'}-u_{z'}v_z)^2.
\end{equation*}
All values of \(w\) are nonzero; hence
\(w_i(1,\xi,z)/w_i(0,\xi,z)\) is independent of \(z\) for fixed \(\xi\).

The missing cross-slice relation comes from a different \(a\mid\xi\) face
of the same physical six-port tensor, not from \(\Delta_\xi\).  Rank at most one of
\(a\mid(b,\xi)\) gives the genuine face minor
\begin{equation*}
  \Gamma(\lambda)
  =q^{\,i,k}_{000}(\lambda)q^{\,i,k}_{101}(\lambda)
   -q^{\,i,k}_{001}(\lambda)q^{\,i,k}_{100}(\lambda)=0.
\end{equation*}
The polynomial \(\Gamma\) vanishes on the full product of phase cosets and
has coordinate degree at most two.  Hence
\cref{lem:nd-grid-interpolation} makes it identically zero.
At \(\lambda^{2z}\), only the ordered pair \((z,z)\) contributes, and its
coefficient is
\begin{equation*}
  w_i(0,0,z)w_i(0,1,z)
  \bigl(w_i(0,0,z)w_i(1,1,z)
        -w_i(1,0,z)w_i(0,1,z)\bigr).
\end{equation*}
It vanishes, so the edge ratio in coordinate \(i\) is the same for
\(\xi=0\) and \(\xi=1\).  Combining this equality with the preceding
constancy over the remaining \((s-2)\)-coordinate cube shows that
\[
 \rho_i:=\frac{w(1,x_{-i})}{w(0,x_{-i})}
\]
is independent of every coordinate of \(x_{-i}\), where \(x_{-i}\) is the
ordered tuple of all coordinates except \(i\), and \(w(b,x_{-i})\) inserts
\(b\) in coordinate \(i\).  Repeat for each \(i\).
Starting at \(0^s\) and flipping the nonzero coordinates of \(z\) in any order
then multiplies by the corresponding constant edge ratios, giving
\cref{eq:app-nd-coefficient-product}.  The vanishing two-face minors make
the result independent of the order of the flips.
\end{proof}

The cross-slice minor \(\Gamma\) compares the ratios on the two
\(\xi\)-layers, which \(\Delta_\xi\) alone cannot do.  The phase
grid uses every element of \(\mu_m\), including all four phases when
\(m=4\).

\subsection{Reduction semantics and completion}
\label{proofsubsec:nd-completion}

Every local disagreement and exceptional endpoint batch has now been
resolved.  For the active prime, the remaining survivor has a
global matching support and the polarized coefficient product.  This final
subsection checks reduction direction and provenance, then records the relative-deck hypotheses used by local constructions.
The final active-core closure below resolves the entire signature set directly.

Every contraction above uses fixed actual representatives; the only
analytic objects are slices, pencil endpoints, inverse evaluation formulas,
and coefficient comparisons.  The two-root choices are legal for
\(m\ge3\).  Unique-lift monomial splices are cancellation-free; phase
traces and two-copy sums may cancel and are handled by their explicit zero
dispatch or by grid interpolation.
Actual four- or six-port outputs therefore have the direct-gadget reduction direction.  In
contrast, \cref{lem:app-nd-coefficient-line-factor,%
lem:app-nd-support-integration,lem:app-nd-polarized-determinant} use analytic
reconstruction only to prove tensor identities; their proper factors are
adjoined through \cref{thm:external-factor,lem:factor-saturation} with the Turing reduction supplied by the factor theorem recorded.  All bounded
outputs have size \(O(\arity(f))\).  Their
complete concrete batches are resolved by
\cref{lem:app-nd-bounded-consumer,lem:nd-six-coefficient-dispatch,%
def:relative-stratum-protocol,lem:matching-synthesis-measure}; in particular,
no nonmatching six-port tensor is recursively re-entered.

For relative worklist invocations, let \(\mathcal W\) be the finite
multiset of unresolved tensor-prime occurrences in the retained factorization trees; occurrences, not only canonical signature keys, are counted.  For
each unresolved occurrence \(u\), with underlying tensor-prime
signature \(h_u\), let \(A_{\mathrm{led}}\) denote the current coordinate matrix and put
\(H_{\mathrm{pb}}:=\overline G(A_{\mathrm{led}},\Lambda)\).  Freeze the
occurrence's own finite card/factor universe
\[
 \mathscr U^{\mathrm{fac}}(A_{\mathrm{led}},h_u,H_{\mathrm{pb}})
 =Q_{u,\mathrm{done}}\sqcup Q_u
\]
when that occurrence first enters the worklist.  Let \(\Phi_u\) be its
stored factor potential.  Define
\[
 \Theta_{\mathrm{MD}}(\mathcal W)
 =
 \{\!\{
   (\arity(h_u),|Q_u|,\Phi_u):u\in\mathcal W
 \}\!\},
\]
using strict lexicographic order on each triple and its well-founded multiset
extension, with the conventions of
\cref{subsubsec:worklist-termination}.

Certifying an occurrence removes its triple.  If the selected ancestor
itself factors, its triple is replaced by triples having strictly smaller
first coordinate.  If processing one child-card key exposes lower-arity
occurrences while the parent remains unresolved, that key is first consumed
or its factor potential strictly decreases; hence the old parent triple is
replaced by a smaller parent triple together with child triples of smaller
arity.  This also decreases the multiset measure.  The bounded four- or six-port reductions
create no fresh unresolved same-arity occurrence: their complete factorization trees are processed inside the originating \(Q_u\).

\begin{definition}[Occurrence-level relative deck stability]
\label{def:occurrence-relative-deck-stability}
For a retained occurrence \(h\) at a current state \((\Lambda,G)\), write
\[
 \operatorname{RelStable}(\Lambda,G;h)
\]
when its complete physical oriented deck has been regenerated; every
already exposed proper card and its complete factorization tree have been
closed; the complete group remains \(G\); and, on the continuing branch,
every surviving nonzero proper card is a \(G\)-matching product, while every
unresolved lower-arity nonbinary occurrence has been placed earlier in the
occurrence induction.  An active deck-stable minimum occurrence is
automatically relatively deck-stable.
\end{definition}

\begin{lemma}[Relative re-entry]
\label{lem:nd-relative-reentry}
Let \(h\) be a tensor-prime occurrence of even arity in the current
normalized state \((\Lambda,\mathcal M_m)\), and suppose
\(\operatorname{RelStable}(\Lambda,\mathcal M_m;h)\).  Then every invocation of deck stability
in the preceding active-prime argument remains valid
with \(h\) in place of the minimum core.
\end{lemma}

\begin{proof}
The sector-ruling, endpoint, support, six-port, and polarization arguments
use minimum arity only to exclude a surviving proper nonmatching card or
nonbinary factor.  Under the present hypothesis each such object has
already produced a resolved leaf or strict stable successor, or has been
certified as a current-group matching product.  All other inputs---the
fixed actual representatives, normalized group, even arity, and tensor
primality of \(h\)---are unchanged.  Thus every local construction and
reduction used in the active-prime proof remains licensed.
\end{proof}

\begin{theorem}[Standard \(\mathcal M_m\)-sector deck rigidity]
\label{thm:normalized-dihedral}
Assume
\[
 \Stable(\Lambda,\mathcal M_m,f),\qquad \Even(\Lambda),\qquad
 \arity(f)=\nu(\Lambda)=2k+2,\qquad m\ge3,
\]
where
\[
 \mathcal M_m
 =\{[\mathsf R(t)],[\mathsf S(t)]:t\in\mu_m\}.
\]
Then \(\Close_{\{\cP\}}(\Lambda,\mathcal M_m,f)\).
\end{theorem}

\begin{proof}[Completion of \cref{thm:normalized-dihedral}]
If \(\arity(f)=4\), regenerate and close its complete physical deck.
In the standard \(\mathcal M_m\) case the output of
\cref{thm:equality-q4-boundary} is an invertible dressing of the prime
quaternary.  It is therefore an actual endpoint-nondegenerate tensor-prime
eight-vertex signature.  Apply
\cref{thm:p1-quaternary-kholant-dichotomy} to its retained augmentation.
Direct realization and downward inheritance give a hard or
\(\Tract(K\Lambda)\) leaf for the entire retained set \(\Lambda\).

For \(\arity(f)=6\), \cref{lem:nd-six-prime-base} contradicts the
tensor primality of the stable core.  Now let \(n:=\arity(f)\ge8\).
At any deleted pair \(p\), two disagreeing nonzero sampled cards would
give the six-port output in \cref{lem:app-nd-localization}.  The first
proper closure of the stable core is zero or a matching product, and
further monomial contractions preserve that property.  This contradicts
the two inherited disagreeing nonzero cards.  Hence all samples at \(p\)
share one matching and all factor lines outside one designated edge \(e\).
The pairwise argument following \cref{lem:app-nd-bounded-consumer} gives
one common nonzero product \(P\) of those fixed factors.  In each nonzero
sector choose two distinct nonzero phase samples and invert its pencil;
an identically zero sector has zero endpoint slices.  Thus all four
physical slices lie in \(\mathbb C P\otimes V_e\), proving the exact
identity \(f=P\otimes q_{p,e}\).  Since \(P\) contains at least one
binary factor, this again contradicts tensor primality.

Consequently no higher-arity stable-core branch remains.  The quaternary
case already resolves the whole signature set, so no local product-class
certificate has to be promoted to a matching-deck certificate for another
occurrence.  These are precisely the allowed outcomes of
\(\Close_{\{\cP\}}(\Lambda,\mathcal M_m,f)\).
\end{proof}

The restriction \(m\ge3\) is essential here.  For \(m=2\), two forbidden
trace roots can exhaust \(\mu_2\), and evaluation on \(\mu_2\) does not
determine a polynomial of coordinate degree two.  That case belongs to the
separate normalized \(V_4\) analysis.

\subsection{The two nonstandard marked dihedral sectors}
\label{sec:marked-dihedral-sectors}

For the nonstandard marked groups, the same binary transfers acquire a
different edge matrix.  We first compute that matrix, then determine
which parts of the standard proof survive.  The matrices \(U_\tau\)
below change analytic coordinates.  The \(\mathcal B_m\) and
\(\mathcal A_m^0\) proofs additionally give whole-network
equivalences; the \(\mathcal A_m\) proof instead keeps the original
physical network and tracks the resulting weights.

\begin{lemma}[Marked-sector dictionary]
\label{lem:marked-sector-dictionary}
For \(\tau=\mathcal M,\mathcal B\), let \(m\ge3\) and set
\begin{equation*}
 U_{\mathcal M}=I,\qquad
 U_{\mathcal B}=I-iX.
\end{equation*}
For \(\tau=\mathcal A\), let \(m\ge4\) be even, fix
\(\eta\in\mu_{2m}\setminus\mu_m\), choose \(s_0\in\AlgNums\) with
\(s_0^2=\eta^{-1}\), and set
\begin{equation*}
 U_{\mathcal A}=
 \begin{pmatrix}1&1\\1&-1\end{pmatrix}\diag(1,s_0),
 \quad s_0^2=\eta^{-1}.
\end{equation*}
Conjugation by \(U_\tau\) sends \(\mathcal M_m\) projectively onto the
corresponding marked group.  Define the transformed edge matrix \(S_\tau\) by
\[
 S_\tau=U_\tau^{\mathsf T}XU_\tau.
\]
Then, projectively,
\begin{equation*}
 S_{\mathcal M}=X,\qquad
 S_{\mathcal B}=I,\qquad
 S_{\mathcal A}=\diag(1,-\eta^{-1}).
\end{equation*}
For \(T\in\mathcal M_m\), choose a fixed actual marked-sector binary
\(B_T\) whose sector-coordinate transfer is \([T]\), and set
\(\widehat B_T:=U_\tau^{-1}B_TU_\tau^{-\mathsf T}\).  Then, for
\(T,U\in\mathcal M_m\),
\begin{equation}
 \widehat B_T\doteq T S_\tau^{-1},
 \qquad
 (TS_\tau^{-1})S_\tau(US_\tau^{-1})
 =TU S_\tau^{-1}.
 \label{eq:marked-sector-fusion}
\end{equation}
Thus multiplication and reversal are preserved, and every monomial fusion
has the same unique internal lift; consequently such a fusion introduces no
cancellation.
\end{lemma}

\begin{proof}
\[
 U_{\mathcal B}R(t)U_{\mathcal B}^{-1}\doteq C(t),\qquad
 U_{\mathcal A}R(t)U_{\mathcal A}^{-1}\doteq A(t),
 \qquad
 U_{\mathcal A}XR(t)U_{\mathcal A}^{-1}
 \doteq ZA(\eta t).
\]
\[
 U_{\mathcal B}^{\mathsf T}XU_{\mathcal B}=-2iI,
 \qquad
 U_{\mathcal A}^{\mathsf T}XU_{\mathcal A}
 =2\diag(1,-\eta^{-1}).
\]
Congruence gives \(\widehat B_T\doteq TS_\tau^{-1}\), and the fusion
identity is matrix multiplication; all omitted scalars are fixed nonzero.
\end{proof}

\paragraph{Sector coordinates and available kernels.}
For an \(r\)-ary current-coordinate signature \(h\), write
\[
 \widehat h_\tau=(U_\tau^{-1})^{\otimes r}h
\]
for its analytic sector-coordinate table.  A \emph{sector-\(X\) dressing} is physical
port dressing by the fixed actual binary \(B_X\); by
\cref{eq:marked-sector-fusion}, its transported transfer fuses as \(X\).
Put
\[
 \beta_{\mathcal B}:=1,
 \qquad
 \delta:=-\eta^{-1},\quad
 \beta_{\mathcal A}:=\delta^{-1},
 \qquad
 \mathscr R_\tau:=\beta_\tau\mu_m,
\]
where the quantities involving \(\eta\) and \(\delta\) apply only in the
\(\mathcal A\)-sector.  For a quotient in \(\mathscr R_\tau\), division
by \(\beta_\tau\) gives a point of \(\mu_m\).  This analytic
normalization is not a gadget.  In one monomial parity plane, a
\emph{complementary-axis pencil} has its two coefficient matrices on
the two coordinate axes.  A sample with exactly one vanishing coordinate
is a singular binary.

\begin{lemma}[Supports of marked quaternaries]
\label{lem:marked-q4-support-atlas}
Let \(q\ne0\) be a retained quaternary occurrence in a current
\(\mathcal B_m\)-state (\(m\ge3\)) or \(\mathcal A_m\)-state (even
\(m\ge4\)).  Suppose its complete physical oriented deck has been
regenerated and every member of that deck is zero or a safe current-group
binary.  Let
\(\tau=\mathcal B\) or \(\mathcal A\), respectively, and define the
sector-coordinate signature
\[
  \widehat q:=(U_\tau^{-1})^{\otimes4}q.
\]
After actual sector-\(X\) dressings, for some
\(K\subseteq\{1,2,3\}\), nonzero \(a,h\), and nonzero
\(b_\ell,d_\ell\) (\(\ell\in K\)),
\begin{equation}
 \widehat q
 =a\delta_{0000}+h\delta_{1111}
  +\sum_{\ell\in K}
   \bigl(b_\ell\delta_{p_\ell}
        +d_\ell\delta_{\bar p_\ell}\bigr),
 \label{eq:marked-q4-complement-form}
\end{equation}
where \(p_1,p_2,p_3\) represent the three complementary pairs of
weight-two words.  The assertion includes all six contracted pairs, both
orientations, all zero-card patterns, and the fixed representative scalars.
\end{lemma}

\begin{proof}
For an ordered port pair with residual slices \((A,B,C,D)\), the complete
card pencils, up to their recorded nonzero representative scalars, are
\begin{equation*}
\begin{array}{ccc}
\toprule
 &E(t)&O(t)\\ \midrule
 \mathcal B_m&A+tD&tB+C\\
 \mathcal A_m&A+\delta tD&tB+\delta C,
 \quad \delta=-\eta^{-1},\\
\bottomrule
\end{array}
\qquad t\in\mu_m.
\end{equation*}
Their effective kernels are \(R(t),XR(t)\) for \(\mathcal B_m\), and
\[
 S R(t)=\diag(1,\delta t),
 \qquad
 SXR(t)=\begin{pmatrix}0&t\\ \delta&0\end{pmatrix},
 \quad S=\diag(1,\delta),
\]
for \(\mathcal A_m\).  The diagonal kernels are symmetric.  Define the
off-diagonal kernel
\(K^{\rm off}(t):=\left(\begin{smallmatrix}0&t\\\delta&0\end{smallmatrix}\right)\);
it satisfies
\begin{equation}
 K^{\rm off}(t)^{\mathsf T}
 =\frac{t}{\delta}K^{\rm off}(\delta^2/t),
 \qquad \delta^2/t\in\mu_m;
 \label{eq:marked-A-card-reversal}
\end{equation}
so reversal permutes the labels and restores a known nonzero scalar; the
\(\mathcal B_m\) formula is \(\delta=1\).

With \(\beta_\tau\) fixed above, three
phases imply that
each two-coordinate pencil is zero, lies on a fixed line with ordered
ratio in \(\mathscr R_\tau\), or has its endpoints on the two
complementary coordinate axes.  Indeed, a zero makes the endpoints
proportional.  Otherwise the quotient \(\phi(t)\) of the two independent
affine coordinates is injective on \(\mu_m\) and takes its values in the
fixed coset \(\beta_\tau\mu_m\).  Hence
\(\psi(t)=\beta_\tau^{-1}\phi(t)\) permutes \(\mu_m\), and
\cref{lem:nd-root-mobius} gives
\(\psi(t)=\lambda t\) or \(\lambda/t\).  Cross-multiplication leaves the
opposite coordinate axes.  Multiplication by \(\beta_\tau^{-1}\) is only an
analytic scaling of one endpoint coordinate; it is not an available binary
or a gadget.  Formula \cref{eq:marked-A-card-reversal} covers the opposite
orientation with its recorded nonzero scalar.

Apply this three-type pencil classification to both card columns for
each of the six deleted pairs.  After analytic quotient normalization, the
Boolean support constraints are exactly those enumerated in
\cref{thm:cert-interface-equality-q4}\textnormal{(ii)}: every nonempty solution consists of
one total parity together with a nonempty subset of its four complementary
word pairs.  Quotient normalization changes no zero pattern, so the same
exact enumeration applies in marked coordinates.  Choose a nonzero
complementary pair and use actual sector-\(X\) dressings, justified by
\cref{eq:marked-sector-fusion}, to send it to \(0000,1111\).  The other
pairs of nonzero entries are then among the three complementary weight-two blocks, giving
\cref{eq:marked-q4-complement-form}.  This includes both orientations and
every zero-card pattern; no analytic sector basis is installed as a gadget.
\end{proof}

\begin{lemma}[Symmetric binary pullback]
\label{lem:symmetric-binary-pullback}
Let \(\cC\in\{\cA,\cP,\cL\}\).  If
\(E,F\in\cC\) are nonsingular symmetric binaries, then
\(EFE\in\cC\).
\end{lemma}

\begin{proof}
Closure under contraction proves the affine case, and the
equality/disequality normal forms prove the product case.  For the
local-affine case, write
\(E=\left(\begin{smallmatrix}a&b\\b&c\end{smallmatrix}\right)\).
The diagonal twists in the definition of \(\cL\) preserve support, which
must be affine.  Nonsingularity excludes support size one, and affine
support excludes size three.  Full support is also impossible: the
\(\alpha=00\) condition requires \(b/a\in\mu_4\), whereas the
\(\alpha=01\) twist requires \(\zeta_8 b/a\in\mu_4\).
On anti-diagonal support symmetry gives \(E=bX\), and the same
\(\alpha=01\) twist again has a nonzero-entry ratio outside \(\mu_4\).
Consequently \(E=a\diag(1,i^k)\), with \(a\ne0\) and
\(k\in\mathbb Z/4\mathbb Z\).  Conversely these four diagonal forms
satisfy both supported twists, \(\alpha=00,11\), and belong to \(\cL\).
Writing \(F=d\diag(1,i^\ell)\) in the same way gives
\(EFE=a^2d\diag(1,i^{2k+\ell})\in\cL\), as required.
\end{proof}

\begin{theorem}[Signature-set-wide reduction of the \(\mathcal B_m\) sector]
\label{thm:marked-B-reduction}
If
\[
 \Stable(\Lambda,\mathcal B_m,f),\qquad
 \Even(\Lambda),\qquad m\ge3,
\]
then \(\Close_{\{\cP\}}(\Lambda,\mathcal B_m,f)\).
\end{theorem}

The resolved \(\Tract\)-alternatives here include
\(\PresK(K\Lambda;\cA)\) and \(\PresK(K\Lambda;\cL)\); the product case is
the only continuing rigidity leaf.

\begin{proof}
Define the transformed signature set
\(\Gamma:=U_{\mathcal B}^{-1}\Lambda\) and the transformed active prime
\[
 \widehat f:=(U_{\mathcal B}^{-1})^{\otimes\arity(f)}f\in\Gamma.
\]
This whole-problem change gives
\begin{equation*}
 U_{\mathcal B}^{\mathsf T}XU_{\mathcal B}=-2iI,
 \qquad
 U_{\mathcal B}^{-1}IU_{\mathcal B}^{-\mathsf T}=\frac{i}{2}X,
 \qquad
 U_{\mathcal B}^{-1}XU_{\mathcal B}^{-\mathsf T}=\frac{i}{2}I.
\end{equation*}
Thus edge and equality exchange, and \cref{lem:p1-equality-wire-equivalence}
gives
\begin{equation*}
 \KHolant(\Lambda)
 \equivT
 \Holant(\Gamma,X)
 \equivT
 \KHolant(\Gamma,I).
\end{equation*}
Define the augmented signature set
\(\Lambda^\sharp:=\Gamma\cup\{I\}\).  Here the \emph{open-network
bijection} is the vertexwise change by \(U_{\mathcal B}^{-1}\), together
with the displayed edge/equality exchange, applied also to the ordered open
boundary legs; the inverse change restores the original network.  It
preserves arity, factorization, minimum nonbinary-prime arity, and parity of
every prime arity; hence \(\Even(\Lambda^\sharp)\).

For a realizable nonsingular \(b\) over \(\Lambda^\sharp\), its pullback
\(B\) satisfies
\begin{equation*}
 B\doteq U_{\mathcal B}bU_{\mathcal B}^{\mathsf T},
 \qquad
 [BX]=[U_{\mathcal B}bU_{\mathcal B}^{-1}],
\end{equation*}
using \(U_{\mathcal B}^{\mathsf T}X=-2iU_{\mathcal B}^{-1}\).  Infinite
order pulls back to the interpolation terminal of
\cref{lem:binary-interpolation}; otherwise
\cref{lem:actual-binary-completeness} applies to the pullback transfer
\([BX]\).  Conjugating by \(U_{\mathcal B}^{-1}\) first puts
\([b]\) in \(\mathcal M_m\).  Since \(X\in\mathcal M_m\), closure under
right multiplication then puts the actual transformed transfer
\([bX]\) in \(\mathcal M_m\).

Conversely, for the chosen actual representative \(B_T\), define its
recorded nonzero representative scalar \(\beta_T\) by
\(B_TX=\beta_TU_{\mathcal B}TU_{\mathcal B}^{-1}\).  Then
\begin{equation*}
 \widehat B_T
 =U_{\mathcal B}^{-1}B_TU_{\mathcal B}^{-\mathsf T}
 =\frac{i\beta_T}{2}T.
\end{equation*}
Since \([R(t)X]=[XR(t^{-1})]\) and
\([XR(t)X]=[R(t^{-1})]\), every class occurs and
\begin{equation*}
 G(\Lambda^\sharp)=\mathcal M_m.
\end{equation*}
The open-network bijection also transports the complete deck: absorbing an
equality wire replaces a label by \(XT,TX\), or \(XTX\), which permute
\(\mathcal M_m\).  Hence zeros, matching products, factor saturation, and
minimum prime arity correspond, so the transformed state is deck-stable.
In particular,
\[
 \Stable(\Lambda^\sharp,\mathcal M_m,\widehat f),
 \qquad
 \arity(\widehat f)=\nu(\Lambda^\sharp)=\arity(f).
\]

Apply \cref{thm:normalized-dihedral}.  Complete closure of each transported
concrete batch is resolved.  On the remaining branch there are
\(\cC\in\{\cA,\cP,\cL\}\) and one \(V\in\GL_2(\AlgNums)\) such that
\[
 V^{-1}(\Gamma\cup\{I\})\subseteq\cC,
 \qquad V^{\mathsf T}XV\in\cC.
\]

Define the common witness matrix \(W:=U_{\mathcal B}V\); then
\(W^{-1}\Lambda\subseteq\cC\).  With
\[
 E=V^{\mathsf T}XV,
 \qquad F=V^{-1}IV^{-\mathsf T},
\]
both \(E,F\) are nonsingular symmetric members of \(\cC\), and
\cref{lem:symmetric-binary-pullback} gives
\[
 W^{\mathsf T}XW
 =-2iV^{\mathsf T}IV
 =-2iEFE\in\cC.
\]
Thus \(W\) is common.  The \(\cA,\cL\) cases are resolved
\(\Tract(K\Lambda)\) leaves; \(\cP\) is the continuing rigidity leaf.
The reductions above use \(U_{\mathcal B}\) for the whole network;
\(W=U_{\mathcal B}V\) records its contribution to the final witness.
\end{proof}

\begin{theorem}[Literal-equality reduction of the
  \texorpdfstring{\(\mathcal A_m^0\)}{A0m} sector]
\label{thm:marked-A0-reduction}
Assume
\[
 \Stable(\Lambda,\mathcal A_m^0,f),\qquad
 \Even(\Lambda),\qquad m\ge4,\quad 2\mid m.
\]
Then \(\Close_{\{\cP\}}(\Lambda,\mathcal A_m^0,f)\).
\end{theorem}

\begin{proof}
Let
\[
 U_0=\begin{pmatrix}1&1\\1&-1\end{pmatrix},
 \qquad U_0^{-1}=\frac12U_0.
\]
Direct multiplication gives, projectively,
\begin{equation}
 U_0R(t)U_0^{-1}\doteq A(t),
 \qquad
 U_0XR(t)U_0^{-1}\doteq ZA(t),
 \label{eq:marked-A0-conjugacy}
\end{equation}
so \(U_0^{-1}\mathcal A_m^0U_0=\mathcal M_m\).  More importantly,
\begin{equation}
 U_0^{\mathsf T}XU_0=2Z,
 \qquad
 U_0^{-1}IU_0^{-\mathsf T}=\frac12I.
 \label{eq:marked-A0-fixed-I}
\end{equation}
Thus this whole-problem change preserves the available literal equality.

Turing-expose the fixed actual representative \(B_j\) of the mark
\(j=[Z]\), retaining its realization scalar and port orientation.  Since
\([B_jX]=j\), one has \(B_j\doteq Y\).  Put
\[
 \Lambda^+=\Lambda\cup\{B_j\},
 \qquad
 \Gamma=U_0^{-1}\Lambda^+,
 \qquad
 \widehat f=(U_0^{-1})^{\otimes\arity(f)}f.
\]
After recording the two nonzero representative scalars, use literal
\(I,Y\in\Gamma\); indeed
\[
 U_0^{-1}IU_0^{-\mathsf T}=\frac12I,
 \qquad
 U_0^{-1}YU_0^{-\mathsf T}=-\frac12Y.
\]
Write \(\Holant^Z(\Gamma)\) for the homogeneous \(Z\)-edge half-edge
problem.  The whole-problem change in
\cref{eq:marked-A0-fixed-I} gives
\[
 \KHolant(\Lambda^+)\equivT\Holant^Z(\Gamma),
\]
where the factor \(2\) per edge is known.  The exact local identities
\begin{equation}
 Z=XIXYX,
 \qquad
 X=-ZIZYZ
 \label{eq:marked-A0-edge-switch}
\end{equation}
give a signature-preserving open-network equivalence
\begin{equation}
 \Holant^Z(\Gamma)\equivT\KHolant(\Gamma).
 \label{eq:marked-A0-open-equivalence}
\end{equation}
For the forward substitution, replace a \(Z\)-edge by three native
\(X\)-edges with the available ordered vertices \(I,Y\); for the reverse
substitution, replace an \(X\)-edge by three \(Z\)-edges with the same two
vertices.  Reversing the ordered \(Y\)-vertex changes only its recorded
nonzero sign.  Define
\[
 \Gamma^\sharp:=\Gamma\cup\{Z\}.
\]
The first identity in \cref{eq:marked-A0-edge-switch} directly realizes
\(Z\) in the native-\(X\) presentation, so adjoining it changes neither the
problem nor the open-gadget closure.

The two edge presentations therefore have the same projective set
\(\mathscr B\) of directly realizable nonsingular ordered binary tensors.
In the \(Z\)-edge presentation, open-network transport and
\cref{eq:marked-A0-conjugacy} give
\[
 \{[bZ]:b\in\mathscr B\}
 =U_0^{-1}\mathcal A_m^0U_0
 =\mathcal M_m.
\]
Since \(m\) is even,
\([Y]=[XR(-1)]\in\mathcal M_m\).  Hence, in the native-\(X\)
presentation,
\begin{equation}
 G(\Gamma^\sharp)
 =\{[bX]:b\in\mathscr B\}
 =\{[bZ][ZX]:b\in\mathscr B\}
 =\mathcal M_m[Y]
 =\mathcal M_m.
 \label{eq:marked-A0-complete-group}
\end{equation}

The complete deck transports as well.  Under \(U_0^{-1}\), the original
native-\(X\) deck is the corresponding \(Z\)-edge deck.  Translating any
target dressing or closing edge by the second identity in
\cref{eq:marked-A0-edge-switch}, and then absorbing the inserted \(I,Y\)
binary chain into the adjacent fixed actual representative, merely permutes
the \(\mathcal M_m\) labels.  The same argument with the first identity gives
the inverse correspondence.  Thus zeros, matching products, factorization trees,
minimum nonbinary-prime arity, and prime-arity parity correspond.  Therefore
\[
 \Stable(\Gamma^\sharp,\mathcal M_m,\widehat f),
 \qquad
 \arity(\widehat f)=\nu(\Gamma^\sharp)=\arity(f),
 \qquad
 \Even(\Gamma^\sharp).
\]

Apply \cref{thm:normalized-dihedral}.  Every terminal, concrete factor
batch, and strict successor transports back through the two open-network
equivalences.  On a surviving common-presentation leaf, there are
\(\cC\in\{\cA,\cP,\cL\}\) and \(V\in\GL_2(\AlgNums)\) such that
\[
 V^{-1}\Gamma^\sharp\subseteq\cC,
 \qquad
 E:=V^{\mathsf T}XV\in\cC.
\]
Set \(W:=U_0V\).  Then \(W^{-1}\Lambda\subseteq\cC\).  Because
\(Z\in\Gamma^\sharp\), also
\[
 F:=V^{-1}ZV^{-\mathsf T}\in\cC.
\]
The binaries \(E,F\) are nonsingular and symmetric, so
\cref{lem:symmetric-binary-pullback} and \(XZX=-Z\) give
\[
 EFE=V^{\mathsf T}XZXV=-V^{\mathsf T}ZV,
 \qquad
 W^{\mathsf T}XW
 =2V^{\mathsf T}ZV
 =-2EFE\in\cC.
\]
Thus \(W\) is a common witness for the original problem.  The
\(\cA,\cL\) cases are resolved \(\Tract(K\Lambda)\) leaves and \(\cP\) is
the continuing rigidity leaf.  The witness \(W=U_0V\) combines the
whole-network equivalence with the transformed common presentation;
literal \(I\) remains available throughout.
\end{proof}

In \cref{eq:marked-q4-complement-form}, write \(g_\ell:=d_\ell\) for
each complementary weight-two block in the sum.  Such a block has
\(b_\ell g_\ell\ne0\); all other blocks vanish.  Ratios below use the
displayed order of the two nonzero binary entries.

\begin{lemma}[The \(\mathcal A_m\) quaternary coset]
\label{lem:marked-A-q4}
Let \(m\ge4\) be even, let \(G(\Lambda)=\mathcal A_m\), and let
\(0\ne q\) be a retained quaternary occurrence satisfying
\[
 \operatorname{RelStable}(\Lambda,\mathcal A_m;q).
\]
Suppose every member of the complete oriented dressed deck of \(q\) is
zero or a safe \(\mathcal A_m\)-binary.  Then the quaternary analysis yields an established
terminal or strict complete-group successor, a genuine factorization of
\(q\) into two binary signatures whose complete factor batch is retained,
or an actual endpoint-nondegenerate eight-vertex signature in the current
physical coordinates; in this last alternative the original
sector-coordinate signature also satisfies
\[
 (U_{\mathcal A}^{-1})^{\otimes4}q\in\cP.
\]
\end{lemma}

\begin{proof}
Apply \cref{lem:marked-q4-support-atlas} and fix the resulting actual
sector-\(X\) dressings.  In sector coordinates those dressings are
invertible monomial transformations, so they preserve membership in
\(\cP\), decomposability, all flattening ranks, and exact factorization.  We may therefore work with
\cref{eq:marked-q4-complement-form} and undo the dressings afterward.

Define the phase scalar \(\delta\) and safe-ratio coset \(\mathscr R_A\) by
\[
 \delta=-\eta^{-1},\qquad
 \mathscr R_A=\delta^{-1}\mu_m=\eta\mu_m.
\]
Since \(\eta\notin\mu_m\) and \(m\) is even,
\begin{equation*}
 \mathscr R_A\cap\mu_m=\varnothing,
 \qquad \mathscr R_A^{-1}=\mathscr R_A.
\end{equation*}
By \cref{lem:marked-sector-dictionary}, every safe sector binary has ratio
in \(\mathscr R_A\).  In \cref{eq:marked-q4-complement-form}, denote the endpoint
pair by \((a,h)\) and an active central pair by \((b_i,g_i)\).

For one active block, the corresponding diagonal pencil is
\begin{equation*}
 \diag(a+\delta t g_i,\ b_i+\delta t h),
 \qquad t\in\mu_m.
\end{equation*}
The root-pencil type classification in
\cref{lem:marked-q4-support-atlas} excludes a
sample with exactly one vanishing coordinate and forces
\begin{equation*}
 (b_i,g_i)=(a\rho_i,h/\rho_i),
 \qquad \rho_i\in\mathscr R_A,
 \qquad b_i g_i=ah,
\end{equation*}
whereas an inactive block gives the diagonal card pencil
\[
 \diag(a,\delta t h),\qquad t\in\mu_m.
\]
Its ordered ratio \(\delta t h/a\) belongs to
\(\mathscr R_A=\delta^{-1}\mu_m\).  Since
\(\delta^2=\eta^{-2}\in\mu_m\), this is equivalent to
\[
 h/a\in\delta^{-2}\mu_m=\mu_m.
\]

If two blocks \(j,k\) are active, the crossing deleted pair and the actual
reflections
\[
 \begin{pmatrix}0&t\\ \delta&0\end{pmatrix},
 \qquad t\in\mu_m
\]
give, up to common endpoint factors,
\[
 t a\rho_j+\delta h/\rho_k,
 \qquad
 t a\rho_k+\delta h/\rho_j.
\]
Writing
\[
 L(t)=ta+\frac{\delta h}{\rho_j\rho_k},
\]
these entries are \(\rho_jL(t)\) and \(\rho_kL(t)\).  Every nonzero sample
therefore has ratio
\(\rho_j/\rho_k\in\mu_m\), because both \(\rho_j,\rho_k\) lie in the same
coset \(\mathscr R_A\).  This ratio is outside the safe coset
\(\mathscr R_A\).  The nonzero affine function \(L(t)\) vanishes for at most
one phase, so another actual phase gives a nonsingular outside-group binary.
Complete binary closure by \cref{lem:actual-binary-completeness} gives a
terminal or a strict group successor, contradicting a remaining
deck-stable branch.

Thus at most one block is active.  If exactly one block \(i\) is active,
then \(ah-b_ig_i=0\), giving a genuine \(2+2\) factorization exposed by
\cref{thm:external-factor,lem:factor-saturation}.  Its two binaries are
then dispatched individually: a new safe class is a strict group successor,
and singular, unsafe, or infinite-order factors reach their established
terminal.

With no active block, let \(\lambda=h/a\in\mu_m\).  Before any further
rotation, the sector table is already
\[
 a\delta_{0000}+h\delta_{1111},
\]
a generalized equality and hence a member of \(\cP\).  Undoing the fixed
sector-\(X\) dressings preserves \(\cP\), proving
\[
 (U_{\mathcal A}^{-1})^{\otimes4}q\in\cP.
\]
Moreover \(\eta^2/\lambda\in\mu_m\), so it is the product of four available
phase parameters.  Actual rotations with this product of port phases give the
sector tensor
\[
 a\bigl(\delta_{0000}+\eta^2\delta_{1111}\bigr).
\]
Since \(s_0^4=\eta^{-2}\),
\begin{align*}
 U_{\mathcal A}^{\otimes4}
 a\bigl(\delta_{0000}+\eta^2\delta_{1111}\bigr)
 &=
 a\bigl((1,1)^{\otimes4}+(1,-1)^{\otimes4}\bigr)\\
 &=2a\,\one[\wt(x)\text{ is even}].
\end{align*}
This is a literal current-coordinate endpoint-nondegenerate eight-vertex signature.  The rotations are
actual by \cref{eq:marked-sector-fusion}; \(U_{\mathcal A}\) only verifies
the identity.
\end{proof}

An actual quaternary, whether it is the selected occurrence or is produced by a
higher-arity marked construction, is not sent directly to
\cref{lem:marked-A-q4}.  First adjoin it with its port order,
scalar, and provenance, jointly factor-saturate it, recompute the complete
binary group, and regenerate its complete oriented one-copy deck.  Test all
three physical \(2+2\) flattenings.  Every rank-one flattening is retained as
an exact factor identity.  Unary, singular, unsafe, or infinite-order factors
reach their established reductions, and a safe binary outside
\(\mathcal A_m\) is a strict group successor.  If all exposed binary factors
are safe current-group binaries, the child quaternary is a current-group matching
product and its originating frozen token is completed; if the quaternary itself is
the selected tensor-prime occurrence, such a factorization contradicts its
primality.

On the remaining branch the quaternary is tensor-prime, its complete physical deck
is zero or safe-group-valued, and the closure just performed gives
\[
 \operatorname{RelStable}(\Lambda,\mathcal A_m;q)
\]
in the sense of \cref{def:occurrence-relative-deck-stability}.  For a later
occurrence, closure of its frozen local queue and prior processing of every
lower-arity occurrence give this certificate directly.  Invoke
\cref{lem:marked-A-q4} directly.  We do not invoke
\cref{lem:finite-q4-dispatch} here: the endpoint support supplied by
\cref{lem:marked-q4-support-atlas} is an analytic sector-coordinate
statement, not a physical endpoint-nondegenerate quaternary.

In the no-active-block branch the physical output is
\(2a\,\one[\wt(x)\text{ is even}]\), with \(a\ne0\).
It is tensor-prime: the weighted generalized equality used to construct
it has rank two across every nontrivial partition, and the invertible
portwise transformations preserve these ranks.  The output is actual,
has nonzero endpoints, and is an eight-vertex signature.  Apply the already
proved \cref{thm:p1-quaternary-kholant-dichotomy} to the retained
augmentation.  This gives a hard or \(\Tract(K\Lambda)\) leaf for the
entire current signature set, with the reduction supplied by actual
realization.  This branch does not return merely a local
\(\cP\)-certificate for a child.  This complete procedure is the
\emph{marked-coordinate reduction}.

For the remainder of the \(\mathcal A_m\)-sector, fix
\[
 \delta:=-\eta^{-1},\qquad D_\delta:=\diag(1,\delta),
\]
and use the portwise-gauge convention of
\cref{subsubsec:gauge-product-weights}; in particular,
\((D_\delta^{(i)}h)(x)=\delta^{x_i}h(x)\).

A finite family of marked-sector cards is a \emph{gauge-coherent
comparison batch} if its members use the same ordered contraction
multigraph, the same orientations on all contracted edges, and the same
external-port order, and differ only in the chosen phase or product at
designated kernels.  For a set \(\mathscr S\) of parent ports write
\[
 D_\delta^{(\mathscr S)}:=\prod_{i\in\mathscr S}D_\delta^{(i)},
 \qquad
 (D_\delta^{(\mathscr S)}h)(x)
   =\delta^{\sum_{i\in\mathscr S}x_i}h(x).
\]
The embedded one-port operators commute.  The set \(\mathscr S\), and
hence the gauged parent \(D_\delta^{(\mathscr S)}h\), is fixed throughout
one comparison batch.  Different contraction skeletons may yield
different sets.

\begin{lemma}[Explicit weighted-gauge transport]
\label{lem:marked-A-gauge-transport}
Let \(m\ge4\) be even, \(H=\mu_m\), and use the fixed
\(\delta=-\eta^{-1}\), where \(\eta\in\mu_{2m}\setminus\mu_m\), and the
diagonal gauge matrix \(D_\delta\) above.  In the analytic \(\mathcal A_m\)-sector the
effective diagonal and reflection kernels are
\begin{equation*}
 K^0_t=D_\delta R(t)=\diag(1,\delta t),
 \qquad
 K^1_t=D_\delta XR(t)=\begin{pmatrix}0&t\\ \delta&0\end{pmatrix},
 \qquad t\in H.
\end{equation*}
For every binary kernel \(K\) and every two-port slice \(g\) on ordered
ports \((i,j)\), pointwise in all remaining variables, every occurrence of
\(D_\delta\) can be moved from the kernel to the corresponding endpoint of
the parent tensor:
\begin{equation}
 \langle D_\delta K,g\rangle_F
   =\langle K,D_\delta^{(i)}g\rangle_F,
 \qquad
 \langle K D_\delta,g\rangle_F
   =\langle K,D_\delta^{(j)}g\rangle_F,
 \label{eq:marked-A-gauge-push}
\end{equation}
where \(i\) and \(j\) are the first and second ordered ports and
\(D_\delta^{(i)}\) means multiplication by \(\delta^{x_i}\).  Consequently
for every gauge-coherent comparison batch there is one fixed set
\(\mathscr S\),
independent of the phase and product choices, such that the entire batch
is exactly the corresponding standard \(\mathcal M_m\) batch on the one
locally gauged parent \(D_\delta^{(\mathscr S)}h\), followed by invertible
row/column gauges and the substitutions
\begin{equation}
 \widehat t=c\,t^{\varepsilon},\qquad
 c\in\{1,\delta,\delta^{-1}\},\quad \varepsilon\in\{1,-1\}.
 \label{eq:marked-A-gauge-phase-substitution}
\end{equation}
In particular:
\begin{enumerate}
\item support, zero or incomplete-support status, matching directions, all flattening ranks,
      and exact factorization are unchanged by the gauges;
\item whenever both product coefficients on a common edge are nonzero and their
      trace is not identically zero, its trace has the form
      \(\kappa(1+\theta t^{\varepsilon})\) with
      \(\kappa\theta\ne0\) and \(\varepsilon\in\{-1,1\}\); a one-term
      trace has no forbidden phase, while an identically-zero product is a
      zero or incomplete-support dispatch according to the other product.  Whenever
      the corresponding protected minor is
      defined, its transported value is a nonzero monomial in \(\delta,t\)
      times the standard protected determinant;
\item reversal sends \(t\) to \(\delta^2/t\in H\), so both orientations
      use the same finite phase set, and all safe ordered ratios lie in the
      single coset \(\delta^{-1}H\).
\item the four samples
      \(E(t_0),E(t_1),O(t_0),O(t_1)\) at one fixed ordered deleted pair,
      the two products of one fixed trace-cover skeleton, the phase family
      producing a quaternary from a six-port tensor, and the complete two-copy polarization grid
      are gauge-coherent.  Hence every joint multi-card comparison in the
      transported proof is made on one common locally gauged parent.
\end{enumerate}
\end{lemma}

\begin{proof}
The two displayed kernel identities are direct multiplication.  For the
first identity in \cref{eq:marked-A-gauge-push}, expand the Frobenius pairing:
\(\sum_{a,b}\delta^aK(a,b)g(a,b)\); the second is identical with the
factor \(\delta^b\).  Transposing or reversing an ordered kernel merely
switches \(i\) and \(j\) and contributes the nonzero scalar in
  \cref{eq:marked-A-card-reversal}.  Induction over the contraction edges
  therefore pushes all \(D_\delta\)'s to vertex ports.  For a fixed ordered
  contraction skeleton, the endpoint at which each \(D_\delta\) lands is
  determined only by the ordered edge, not by its phase or by the choice of
  diagonal versus reflection product.  Hence every member of a
  gauge-coherent batch has the same parent-port set \(\mathscr S\).  Reversed cards
  are first put in the common orientation using
  \cref{eq:marked-A-card-reversal}, which only changes the recorded scalar
  and replaces the phase by \(\delta^2/t\).  On a matching product the
  remaining external gauges are row and column gauges of its binary
  factors, so they preserve support, nonsingularity, ranks, and exact
  tensor-factor identities.  Unrelated contraction skeletons may induce
  different sets \(\mathscr S\), but their raw coefficients are never compared:
  support assertions are gauge-invariant, and coefficient identities pull
  back with their explicit nonzero product-weight factors.

Pushing the entire gauge to the parent leaves the phase unchanged, so
\(c=1\).  Alternatively, keeping the raw parent uses
\(K^0_t=R(\delta t)\) and \(K^1_t=\delta XR(t/\delta)\), giving
\(c=\delta\) and \(c=\delta^{-1}\), respectively.  These are the
substitutions in \cref{eq:marked-A-gauge-phase-substitution}; a reversed reflection has
\(\delta^2/t\), and \(\delta^2=\eta^{-2}\in H\).  A trace of two monomial
factors is a sum of two nonzero monomials.  If the two terms have the same
exponent and cancel, that product is identically zero (the full card may
consequently be zero or partially supported); otherwise, after
extracting one term it is \(\kappa(1+\theta t)\) with
\(\kappa\theta\ne0\) (or a one-term trace).  A unique-lift
fusion multiplies a protected minor by the nonzero entry and orientation
scalars; a row/column gauge multiplies a determinant by the product of its
diagonal entries.  This proves the second item and also shows that a zero
or partially supported card is zero or partially supported before and after transport.
Finally, \(\delta^{-1}H=\eta H\) is disjoint from \(H\) and is invariant
under inversion because \(\eta^{-2}\in H\).  This proves the ratio and
orientation assertions.
\end{proof}

\paragraph{The six-port argument with weighted kernels.}
\label{par:marked-A-six-base}
The six-port conclusion of \cref{lem:nd-six-prime-base} also holds in the marked \(\mathcal A_m\) sector,
for even \(m\ge4\), without a simultaneous normalization of its
safe-ratio cosets.  Work analytically with
\(\widehat h=(U_{\mathcal A}^{-1})^{\otimes6}h\).
The exact sector-ruling argument uses only monomial supports and
Segre factorization.  In its individual binary pencils, the safe ratios
lie in \(\delta^{-1}\mu_m\); dividing that one ratio by
\(\delta^{-1}\) gives the same \(Z/P/C\) support classification.
This is a separate analytic calculation for each pencil, not a common
port gauge or an available binary.  In the proof above replace the
parity kernels \(I,X\) by the actual effective kernels
\[
 D_\delta R(1)=\begin{pmatrix}1&0\\0&\delta\end{pmatrix},
 \qquad
 D_\delta XR(1)=\begin{pmatrix}0&1\\\delta&0\end{pmatrix}.
\]
They have the same respective supports and nonzero entries.  Every
cross-pair lift in the \(P/P\) and \(P/C\) arguments and every
injective lift in the support atlas therefore acquires only a nonzero
weight; all support counts and projection obstructions are unchanged.
A matching coset still has an identically zero opposite-parity sector.
The resulting exact factorization of \(\widehat h\) pulls back by
the portwise tensor powers of \(U_{\mathcal A}\) to an exact
factorization of the physical \(h\), with its recorded nonzero
scalar.  No analytic gauge or factor is used as a gadget.

\begin{lemma}[Primitive transport in the \(\mathcal A_m\) sector]
\label{lem:marked-A-primitive-transport}
Let
\[
 \Stable(\Lambda,\mathcal A_m,f_0),\qquad \Even(\Lambda),
 \qquad m\ge4\text{ even},
\]
and define the phase group \(H\), phase scalar \(\delta\), diagonal gauge
matrix \(D_\delta\), and safe-ratio coset \(\mathscr R_A\) by
\[
 H=\mu_m,\qquad \delta=-\eta^{-1},\qquad
 D_\delta=\diag(1,\delta),\qquad \mathscr R_A=\delta^{-1}H.
\]
Let \(h\in\Pi(\Lambda)\) be a selected unresolved nonbinary prime
occurrence of even arity at least six.  Regenerate and close its complete
actual oriented deck over the current retained set.  Suppose that this
occurrence is relatively deck-stable and has produced neither a resolved
complexity leaf nor a strict stable-tuple successor; in particular, all lower-arity
prime occurrences exposed by its proper cards or factors have already been
processed, and every surviving binary factor belongs to the current
complete group.  Then every concrete local comparison batch in the
normalized-dihedral proof has a gauge-coherent physical counterpart.  Its
required invariants are the following:
\begin{enumerate}
\item The actual kernels are monomial, their traces remain affine in
      \(t\in H\), reversal permutes \(H\), and every safe ordered ratio lies
      in the fixed coset \(\mathscr R_A\).  Analytically dividing that ratio by
      \(\delta^{-1}\) gives the pencil normal form.  Hence all support,
      ruling, trace-cover, coefficient-minor, endpoint-type, and
      bounded-reduction arguments transport with their zero cases;
      unique-lift monomial fusions remain cancellation-free.
\item For a logical coefficient table
      \(c=(c_z)_{z\in\Ftwo^d}\) on an oriented matching and orientation
      signs \(\epsilon_j\in\{1,-1\}\), the product-weight change
      \(\widetilde c_z=\delta^{\sum_j\epsilon_jz_j}c_z\) sends safe
      logical ratios from \(\mathscr R_A\) to \(H\) while preserving
      \cref{eq:app-nd-six-relations} and the cross ratio \(r\).  The final
      physical quaternary is routed through the marked-coordinate reduction.
\item Fixed \(\delta\)-weights do not alter support integration.  They
      multiply the Reed--Muller obstruction by a nonzero scalar.  In the
      polarized determinant, they merely reparameterize the full grid and
      multiply each tested coefficient by a nonzero product-weight value.
\end{enumerate}
Consequently every branch either
\begin{enumerate}
\item reaches a resolved complexity leaf or a strict stable-tuple successor;
\item returns an exact proper-factor or child-card factor batch to the
      lower-arity occurrence induction; or
\item certifies
      \[
       (U_{\mathcal A}^{-1})^{\otimes\arity(h)}h\in\cP.
      \]
\end{enumerate}
When \(h=f_0\), the first two alternatives are resolved in the
sense of \cref{def:protocol-discharged}.  For a later occurrence they are
charged to the occurrence-level induction described in
\cref{proofsubsec:nd-completion}; no decrease of the active tuple's
\(\Omega\) is asserted.
The bounded outputs are actual quaternary or six-port tensors in the
current coordinates.  Factor alternatives are exact tensor identities
pulled back to those coordinates and exposed through
\cref{thm:external-factor,lem:factor-saturation}.  Their Turing reductions
are retained in the worklist.  The conclusion concerns the selected
occurrence \(h\).
\end{lemma}

\begin{proof}
By \cref{lem:marked-sector-dictionary}, an actual sector binary is a
nonzero multiple of \(TD_\delta^{-1}\), and its effective kernel is
\(D_\delta T\).
The gauge identity \cref{lem:marked-A-gauge-transport} moves each kernel
weight to a specified parent port.  It gives the effective kernels
\begin{equation}
 D_\delta R(t)=\diag(1,\delta t),
 \qquad
 D_\delta XR(t)=\begin{pmatrix}0&t\\ \delta&0\end{pmatrix},
 \label{eq:marked-A-effective-kernels}
\end{equation}
so a fiber has pencils \(A+\delta tD\) and \(tB+\delta C\).  For a
nonproportional pencil, multiply the safe ratio by \(\delta\) and use the
  coset-normalized pencil argument; the resulting ratio is in
\(H\), so \cref{lem:nd-root-mobius,lem:nd-root-pencil} applies.  This is
an analytic rescaling only.  Reversal sends \(t\) to
\(\delta^2/t\in H\), and the gauge lemma shows that a common-edge trace
has at most one forbidden phase in each product.  Hence the two-product
trace cover still has an admissible phase because \(m\ge4\), while every
zero or partially supported card is preserved exactly by the invertible gauges.
Every comparison used here belongs to one of the gauge-coherent batches
listed in item~4 of \cref{lem:marked-A-gauge-transport}.  Thus its two
samples, two products, or full phase grid are evaluated on one common
locally gauged parent.  Different batches may carry different portwise
gauges, but their conclusions are either support/rank/factor assertions,
  which are gauge-invariant, or coefficient identities whose product-weight
factors are displayed below.  In particular, no equation compares raw
coefficients from differently gauged parents.  The determinant statement
in the gauge lemma gives, with an explicitly
nonzero scalar \(\kappa\),
\[
 \det_{\mathcal A}(t)
   =\kappa\,\delta^e t^u
      \det_{\mathcal M}(c\,t^{\varepsilon}),
\qquad e,u\in\mathbb Z,\quad c\ne0,\quad\varepsilon\in\{-1,1\},
\]
where \(\det_{\mathcal M}\) denotes the same protected minor for the
locally gauged tensor in the standard sector,
so protected tags, the six-to-four Laurent coefficient, and the
endpoint minor are nonzero exactly in the corresponding standard
calculation.  Therefore
\cref{lem:app-nd-sector-ruling,lem:app-nd-compatible-closure,%
lem:app-nd-coefficient-line-factor,lem:app-nd-localization,%
lem:app-nd-bounded-consumer,lem:app-nd-parallelogram,%
lem:app-nd-exceptional-endpoints,lem:app-nd-endpoint-type,%
lem:app-nd-six-to-four}
transport with the same actual closures.  Every bounded quaternary output
is routed through the marked-coordinate reduction,
and every exact factor identity is pulled
back by the corresponding portwise tensor power of \(U_{\mathcal A}\) and
then handled with a recorded Turing reduction.  For the active minimum occurrence these
transported concrete batches are resolved.  For a later
occurrence, each batch is completely closed inside that occurrence's frozen
local queue and charged to the outer occurrence induction.  On the branch
surviving the corresponding closure, the standard support, ruling, and
coefficient invariants remain available.

Use the product-weight convention of
\cref{subsubsec:gauge-product-weights}; for this logical
coefficient cube put
\[
 \chi(z):=\prod_j\gamma_j^{z_j}\qquad(\gamma_j\in\mathbb C^\times).
\]
Multiplication by \(\chi\) preserves supports and all two-face
parallelogram identities.  It also preserves every complementary-pair
product up to one common nonzero scalar.  For oriented matched pairs, a safe logical ratio is
\(\delta^{-\epsilon_j}h_j\), \(h_j\in H\),
\(\epsilon_j\in\{1,-1\}\).  The product-weight change
\[
  \widetilde c_z
  =\delta^{\sum_j\epsilon_jz_j}c_z
\]
puts all ratios in \(H\), multiplies both sides of every relation in
\cref{eq:app-nd-six-relations} equally, and leaves
\[
  r=\frac{c_{000}c_{110}}{c_{100}c_{010}}
\]
unchanged.  Here the product-weight change is analytic.  Since
\(\delta^2\in H\), absorb orientation changes into the phases and write
the three initial safe ratios as \(\delta^{-1}h_j\), \(h_j\in H\).
The crossing factor cards used in the standard six-port calculation
expose these raw ratios: the extra \(\delta t\) multiplies a column of
the card and does not change the ratio on its unclosed matching edge.
Actual rotations \(R(h_j^{-1})\) therefore leave the physical table
\[
 c_z=\delta^{-|z|}N_r(z),
 \qquad N_r(z)=\begin{cases}1,&|z|\le1,\\r,&|z|\ge2,\end{cases}
\]
after its nonzero realization scalar is recorded.  To check \(r\in H\)
without realizing the analytic weight change, use the standard
\(X_{12}X_{34}X_{56}\) support and close ports \((1,3)\) by \(K^1_t\).
The resulting physical quaternary, with logical row \(x_2\) and column
\(x_5\), has coefficient matrix
\[
 \begin{pmatrix}1&r/\delta\\t/\delta&tr/\delta^2\end{pmatrix}
 =\binom{1}{t/\delta}(1,r/\delta).
\]
Expose its exact binary factors through the factor theorem.  Unless a
binary exit occurs, \(r/\delta\in\delta^{-1}H\), proving \(r\in H\).
Closing the third matching pair of the physical table by \(K^1_t\) gives
\[
 Q_t=\diag(1,\delta^{-1})
 \begin{pmatrix}1+t&r+t\\r+t&r(1+t)\end{pmatrix}
 \diag(1,\delta^{-1}),\qquad
 \det Q_t=\delta^{-2}(r-1)(t^2-r).
\]
If \(r=1\), the displayed six-port table already factors.  Otherwise
at most two phases satisfy \(t^2=r\); choose another phase in \(H\).
The resulting nonzero physical quaternary enters the marked-coordinate
reduction with the displayed nonzero minor.

Support integration is Boolean and unchanged.  On its sole Reed--Muller
residue the raw face, in one orientation, is
\[
 \begin{pmatrix}A&\delta tB\\\delta tC&D\end{pmatrix},
 \qquad A,B,C,D\ne0.
\]
The actual \(t=1\) matching card gives \(AD=\delta^2BC\), so its
determinant is \(\delta^2BC(1-t^2)\).  The complementary exponent pattern
gives the same obstruction up to sign and a nonzero scalar.  It is nonzero
for some \(t\in H\setminus\{\pm1\}\).  In polarization,
each grid variable is merely reparameterized as
\(\lambda_j=c_jt_j\), \(c_j\in\mathbb C^\times\), and the remaining weights form
a nonzero product weight \(\chi\).  The grid is a product of cosets
\(c_j\mu_m\), so \cref{lem:nd-grid-interpolation} applies without requiring
\(c_j\in H\).  Hence the adjacent coefficient becomes
\[
 [\lambda^{z+z'}]\Delta_\xi
 =\chi(z)\chi(z')(u_zv_{z'}-u_{z'}v_z)^2
 \qquad\text{for adjacent }z,z'.
\]
The cross-slice coefficient is similarly multiplied by \(\chi(z)^2\).
Thus the full grid and all vanishing tests in
\cref{lem:app-nd-support-integration,lem:nd-q6-minor-localization,%
lem:app-nd-polarized-determinant}
remain valid.

The transported support integration therefore gives one global matching
support for the selected occurrence, and the transported polarization
identities integrate its nonzero coefficients into a product of binary
factors on that matching.  This is precisely
\((U_{\mathcal A}^{-1})^{\otimes\arity(h)}h\in\cP\).
If an intermediate identity factors the selected ancestor \(h\) itself,
replace its occurrence by the tensor-prime occurrences of its proper
factors.  If instead only a constructed child card factors, process that
child's complete factorization tree inside the selected occurrence's frozen
local queue; the token for \(h\) remains unresolved until its analysis
resumes.  Together with the
product certificate above, this proves the occurrence-level alternatives.
\end{proof}

\begin{theorem}[\(\mathcal A_m\)-sector deck rigidity]
\label{thm:marked-A-rigidity}
Assume
\[
 \Stable(\Lambda,\mathcal A_m,f),\qquad \Even(\Lambda),
 \qquad m\ge4,\quad 2\mid m.
\]
Then \(\Close_{\{\cP\}}(\Lambda,\mathcal A_m,f)\).
\end{theorem}

\begin{proof}
Fix \(U_{\mathcal A}\) from \cref{lem:marked-sector-dictionary} and
let \(n=\arity(f)\) be the active minimum arity.  If \(n=4\), the
marked-coordinate reduction just proved gives a whole-signature-set leaf
or an established strict successor.  For \(n=6\), the independent
weighted six-port support argument above gives a proper
factorization, contradicting tensor primality.

Suppose \(n\ge8\), and work analytically with
\(\widehat f=(U_{\mathcal A}^{-1})^{\otimes n}f\).
The two-card localization proof transports directly to the actual
weighted kernels: common-edge traces exclude at most two phases; joining
two parity discrepancies gives a fused weight ratio
\(\kappa t^{\pm2}\), \(\kappa\ne0\); and cycle shortening or vertex
replacement has the same unique-lift matching support.  These are the
explicit kernel and fusion identities in
\cref{lem:marked-A-gauge-transport,lem:app-nd-localization}.
As \(m\ge4\), each required actual phase exists.  A disagreeing pair
of sampled cards would therefore localize to a nonmatching six-port
tensor.  But the first proper physical closure of the stable core is
zero or a product of safe current-group binaries, and subsequent
contractions preserve that property by the group fusion law.  The
inherited nonzero disagreeing cards give a contradiction.

Thus at a fixed deleted pair all nonzero samples have one matching and
one common varying edge.  The elementary pairwise argument used in the
standard completion fixes the other factor lines across both sectors.
The weighted pencils \(A+\delta tD\) and \(tB+\delta C\) are invertible
from two distinct nonzero samples, since \(\delta\ne0\).  Their four
slices have one common nonzero binary product factor, proving an exact
proper factorization of \(\widehat f\).  Pulling it back by the
invertible portwise matrices \(U_{\mathcal A}\) gives a proper
factorization of the physical core \(f\), again impossible.

No higher-arity stable-core branch survives, and the quaternary route
already resolves the entire retained set.  This proves
\(\Close_{\{\cP\}}(\Lambda,\mathcal A_m,f)\) without promoting a
local product-class certificate to relative matching-deck stability.
\end{proof}

\begin{proof}[Proof of \cref{thm:marked-dihedral}]
If \(G=\mathcal M_m\), then deck stability and \(\Even(\Lambda)\) give
\(\arity(f)=\nu(\Lambda)=2k+2\) for some \(k\), so apply
\cref{thm:normalized-dihedral}.  If \(G=\mathcal B_m\), apply
\cref{thm:marked-B-reduction}.  If \(G=\mathcal A_m^0\), apply
\cref{thm:marked-A0-reduction}.  If \(G=\mathcal A_m\), apply
\cref{thm:marked-A-rigidity}.  These are exactly the four cases in the
statement.
\end{proof}
\section{Proper Cyclic Matching Decks and Reed--Muller Rigidity}
\label{sec:proper-deck}

Throughout this section, an unadorned \(0\) or \(\one\) used in a vector
context denotes the all-zero or all-one vector in the ambient binary space,
whose dimension is clear from context.  For \(S\subseteq\Ftwo^r\), write
\(\one[S]\) for the arity-\(r\) indicator signature
\(x\mapsto\one[x\in S]\).

Here the effective kernels span only \(\Span\{I,X\}\).  In the Walsh
coordinates of \cref{eq:Walsh-coordinates}, these kernels impose equality
of two bits, with a possible phase on their common value.  This restricted
access leaves one higher-arity core, \(\RMcore\).  We first reconstruct
that core from its pair contractions, then use networks of it to test every
retained signature.  The code and self-duality conventions needed for the
second step are in \cref{subsec:binary-affine-conventions}.

\subsection{Cyclic phase pencils and physical dressings}
\label{subsec:proper-cyclic-phase-pencils}

The cyclic group supplies a finite family of phase contractions.  The next
three lemmas describe which binary pencils they permit and which changes
to a quaternary table their actual port dressings can realize.

Here a tensor or matrix \emph{pencil} is a one-parameter linear family
\(U+tV\) over the stated field.  A \emph{sampled phase pencil} restricts
\(t\) to a displayed finite root-of-unity group, here usually \(\mu_N\),
and records zero members as well as the projective classes of nonzero
members.  Two nonzero members are proportional when one is a nonzero scalar
multiple of the other.  A monomial matrix has exactly one nonzero entry in
each row and column; in dimension two it is diagonal or anti-diagonal.

For each integer \(N\ge3\) and \(s\in\mu_N\), define the diagonal phase
matrix
\[
  D_s=\diag(1,s).
\]
We use \(P=\left(\begin{smallmatrix}1&1\\1&-1\end{smallmatrix}\right)\)
for the analytic Walsh change of coordinates.  Actual networks continue to
use the native \(X\)-edge; \(P^{\mathsf T}XP=2Z\) gives their edge
matrix in the tables below.

For \(N\ge3\) and \(a,b,c,d\in\AlgNums\), define the relation
\(\mathcal R_N(a,b,c,d)\) by declaring that, for some \(\alpha\in\mu_N\), at least
one of the following systems holds:
\begin{align}
  c&=-\alpha a,&d&=-\alpha b,
  &&\text{(proportional)},
  \label{eq:odd-dih-R-proportional}\\
  b&=c=0,&d&=\alpha a,
  &&\text{(diagonal monomial)},
  \label{eq:odd-dih-R-diagonal}\\
  a&=d=0,&c&=\alpha b,
  &&\text{(anti-diagonal monomial)}.
  \label{eq:odd-dih-R-antidiagonal}
\end{align}
The alternatives may overlap and include the all-zero quadruple.

For a quaternary signature \(q\), define its analytic Walsh-transform
signature by
\(\widetilde q:=(P^{-1})^{\otimes4}q\).  When this analytic tensor is
supported on even Hamming weight, define its endpoint and central Walsh
coordinates by
\begin{equation}
  p=\widetilde q_{0000},\qquad
  t=\widetilde q_{1111},\qquad
  x_{ab}=\widetilde q_{\one_{\{a,b\}}},\qquad 1\le a<b\le4.
  \label{eq:odd-dih-q4-coordinates}
\end{equation}
Here \(\one_{\{a,b\}}\in\Ftwo^4\) has ones precisely at ports
\(a,b\), so \(x_{ab}\) indexes one of the six weight-two words.

\begin{lemma}[Cyclic phase-pencil rigidity, including zeros]
\label{lem:odd-dih-RN}
For every integer \(N\ge3\), the relation
\(\mathcal R_N(a,b,c,d)\) is equivalent to the following
condition: for every \(s\in\mu_N\),
\[
  \diag(a-sb,c-sd)Z
  \quad\text{is zero or is projectively }D_\eta
  \text{ for some }\eta\in\mu_N.
\]
\end{lemma}

\begin{proof}
If \(a=b=c=d=0\), both conditions hold, so assume otherwise.  Set
\[
  u(s)=-c+sd,\qquad v(s)=a-sb.
\]
The phase-card condition says that \(u(s)\) and \(v(s)\) are either both zero
or both nonzero, and in the latter case
\(u(s)/v(s)\in\mu_N\).  If the two affine linear forms are dependent, the
possibility that exactly one is identically zero is excluded by this
condition.  Hence \(u=\alpha v\) for some \(\alpha\in\mu_N\), giving
\cref{eq:odd-dih-R-proportional}.  If they are independent, neither form
vanishes on \(\mu_N\), and
\(\psi(s):=u(s)/v(s)\) is a nonconstant M\"obius
transformation.  The card condition gives
\(\psi(\mu_N)\subseteq\mu_N\); injectivity and equality of the finite
cardinalities give \(\psi(\mu_N)=\mu_N\).  Thus
\cref{lem:nd-root-mobius} applies and gives
\(\psi(s)=\alpha s\) or \(\alpha/s\), with
\(\alpha\in\mu_N\).  Cross multiplication yields
\cref{eq:odd-dih-R-diagonal,eq:odd-dih-R-antidiagonal}.  Direct substitution
proves the converse, including the all-zero case.
\end{proof}

For the two physical cyclic lemmas below, fix a finite algebraic retained
signature set \(\Lambda\) in \(K\)-coordinates, with native edge \(X\), and
fix an integer \(N\ge3\).  For every \(s\in\mu_N\), fix an actual ordered binary gadget
\(C_s\) over \(\Lambda\), with a specified orientation, whose analytic Walsh
matrix is
\[
  \widetilde C_s=\kappa_sD_sZ,
  \qquad \kappa_s\in\AlgNums^\times.
\]
The nonzero scalar \(\kappa_s\) is recorded with the gadget.  Throughout
these lemmas, ``actual'' and ``directly realizable'' are relative to this same
\(\Lambda\), and every displayed phase-card matrix is divided by its recorded
nonzero scalar.  In the proper-cyclic application, analytic Walsh placement
puts \(G(\Lambda)\) in the standard cyclic form
\(C_N=\{[D_s]:s\in\mu_N\}\), and the \(C_s\) are fixed actual representatives
of all its classes.

\begin{lemma}[Cyclic quaternary phase family]
\label{lem:odd-cyc-q4-family}
Let \(q\) be an actual quaternary signature over \(\Lambda\), and suppose
\(\widetilde q=(P^{-1})^{\otimes4}q\) is supported on the even-weight Walsh
sector.  Define \(p,t,x_{ab}\) by \cref{eq:odd-dih-q4-coordinates}.  Assume
that, for every complementary pair, both orders of its two residual ports,
and every \(s\in\mu_N\), contraction with the fixed actual representative
\(C_s\) gives either the zero card or a nonsingular binary card whose Walsh
transfer class belongs to
\(C_N=\{[D_\eta]:\eta\in\mu_N\}\).  Equivalently, define
\[
  \mathfrak I=\bigl\{\{1,2\},\{1,3\},\{1,4\}\bigr\},
  \qquad I^c=[4]\setminus I,
\]
and write \(x_I:=x_{ab}\) for \(I=\{a,b\}\).  For every
\(I\in\mathfrak I\), \(s\in\mu_N\), and
\(\varepsilon\in\{+,-\}\), the corresponding normalized actual Walsh card
has matrix
\begin{align*}
 B_{I,+}(s)&=\diag(p-sx_I,x_{I^c}-st),\\
 B_{I,-}(s)&=\diag(p-sx_{I^c},x_I-st).
\end{align*}
Here \(+\) records deletion of \(I\) and \(-\) deletion of \(I^c\); both
orders of the remaining two ports are part of the assumption.  Each displayed
card is zero or satisfies
\[
  [B_{I,\varepsilon}(s)Z]\in C_N.
\]
The right factor \(Z\) is only the Walsh transfer convention and is not an
additional gadget.  Organize the eight coordinates into the endpoint block
\(\mathbf v_0=(p,t)\) and the three central complement blocks
\(\mathbf v_I=(x_I,x_{I^c})\), \(I\in\mathfrak I\).  Then exactly one of the
following holds:
\begin{enumerate}
\item If \(p,t\ne0\), then \(\lambda=t/p\in\mu_N\), and each central
      block \(\mathbf v_I\) is zero or
      \begin{equation}
        \mathbf v_I=(-p\lambda\alpha_I^{-1},-p\alpha_I),
        \qquad \alpha_I\in\mu_N.
        \label{eq:odd-cyc-q4-nonzero-endpoints}
      \end{equation}
\item If \(p=t=0\), each central block \(\mathbf v_I\) is zero or
      \[
        \mathbf v_I=\rho_I(1,\alpha_I),
        \qquad \rho_I\in\AlgNums^\times,
        \quad \alpha_I\in\mu_N.
      \]
\end{enumerate}
In particular, exactly one of \(p,t\) cannot be nonzero.
Conversely, every even Walsh table in these alternatives has this
zero-or-cyclic property under every phase contraction.
\end{lemma}

\begin{proof}
For a complementary block \(I,I^c\), the phase card is
\[
  \diag(p-sx_I,x_{I^c}-st),
  \qquad s\in\mu_N.
\]
Apply \cref{lem:odd-dih-RN} to
\((a,b,c,d)=(p,x_I,x_{I^c},t)\), and then interchange \(I,I^c\).
If \(p,t\ne0\), the anti-diagonal system is impossible.  If one of the
two complementary-deletion systems is diagonal, then \(x_I=x_{I^c}=0\) and
\(t/p\in\mu_N\); the other system is then diagonal as well.
In the remaining case both systems are proportional.  Taking their
parameters to be \(\alpha,\beta\in\mu_N\) gives
\[
 x_{I^c}=-\alpha p,\qquad x_I=-\beta p,\qquad
 t=\alpha\beta p.
\]
Thus \(t/p=\alpha\beta\in\mu_N\), and choosing
\(\alpha_I=\alpha\) gives the displayed parametrization.  If exactly one
endpoint is nonzero, neither monomial system is possible and the two
proportional systems force respectively \(x_I=0,x_{I^c}\ne0\) and
\(x_{I^c}=0,x_I\ne0\), a contradiction.  If \(p=t=0\), proportional and
diagonal systems give only the zero block; every nonzero block is therefore
anti-diagonal, \(\mathbf v_I=\rho_I(1,\alpha_I)\).  Direct substitution checks
both complementary deletions and every card arising from a zero complementary
block, proving the converse.
\end{proof}
\begin{lemma}[Actual quaternary cyclic dressing]
\label{lem:odd-cyc-four-port-dressing}
Let \(q\) be an actual quaternary signature over \(\Lambda\).  For
\(r_1,r_2,r_3,r_4\in\mu_N\), attach the fixed representative \(C_{r_i}\), in
its specified orientation, to port \(i\) of \(q\) through one native
\(X\)-edge, leaving the other port of \(C_{r_i}\) exposed.  Denote the
normalized output signature by \(h\).  After division of the direct-gadget
boundary tensor by the recorded nonzero scalar
\(2^4\prod_i\kappa_{r_i}\), its Walsh coordinates are
\begin{equation*}
  \widetilde h_y
  =\widetilde q_y\prod_{i=1}^4r_i^{y_i}.
\end{equation*}
Define the map \(F_N:\mu_N^4\to\mu_N^4\) by
\begin{equation*}
  F_N(r_1,r_2,r_3,r_4)
  =(R,r_{12},r_{13},r_{14}),
  \qquad
  R=\prod_{i=1}^4r_i,
  \quad r_I=\prod_{i\in I}r_i,
\end{equation*}
and put \((\mu_N)^2:=\{z^2:z\in\mu_N\}\).  A target
\((R,u,v,w)\in\mu_N^4\) lies in the image of \(F_N\) exactly when
\begin{equation}
  \frac{uvw}{R}\in(\mu_N)^2.
  \label{eq:cyclic-dressing-image}
\end{equation}
For odd \(N\), \(F_N\) is bijective.
\end{lemma}

\begin{proof}
At port \(i\), the fixed Walsh matrix
\(\kappa_{r_i}D_{r_i}Z\), followed by the analytically transformed native
edge \(P^{\mathsf T}XP=2Z\), acts as
\(2\kappa_{r_i}D_{r_i}\).  The four attachments therefore give the stated
actual dressing and recorded scalar.  Necessity is
\(uvw/R=r_1^2\); for sufficiency choose
\[
  r_1^2=uvw/R,
  \qquad
  r_2=u/r_1,
  \qquad
  r_3=v/r_1,
  \qquad
  r_4=w/r_1.
\]
These choices satisfy
\(F_N(r_1,r_2,r_3,r_4)=(R,u,v,w)\); for odd \(N\), squaring on
\(\mu_N\) is bijective.
\end{proof}

\subsection{Quaternary entry and theorem statement}

This branch, \(\mathscr K_G=\Span\{I,X\}\), is the marked proper-cyclic row
of Xia's taxonomy; see \cref{rem:xia-binary-taxonomy}.  Its arity-four
boundary is independent.

\begin{lemma}[Proper-cyclic quaternary boundary]
\label{lem:proper-cyclic-q4}
Let \((\Lambda,G,q)\) be a deck-stable retained state with
the distinguished literal equality \(I\) available,
\(\arity(q)=\nu(\Lambda)=4\), kernel span \(\mathscr K_G=\Span\{I,X\}\), and finite cyclic \(G\ni[X]\).  Then either an
unsafe actual binary is exposed, factor access strictly enlarges the
recomputed complete group, or an actual gadget using at most two copies of
\(q\) realizes a tensor-prime endpoint-nondegenerate eight-vertex signature.  The boundary thus
returns a resolved leaf, an outer deck successor under
\cref{lem:matching-synthesis-measure}, or \(\Tract(K\Lambda)\).
\end{lemma}

\begin{proof}
Conjugate analytically to Walsh coordinates and continue to write \(G\) for
the conjugated projective group.  The proper algebra is diagonal:
\begin{equation*}
  G=C_N=\{[D_s]:D_s=\diag(1,s),\ s\in\mu_N\},
  \qquad 2\mid N.
\end{equation*}
Here \([X]\) becomes \([Z]=[D_{-1}]\), and actual representatives are
proportional to \(D_sZ\).  To justify the parity reduction, fix an ordered
deleted pair and an odd-weight word \(y\) on the two residual ports.  Up to
a fixed nonzero representative scalar and a possible sign, the corresponding
off-diagonal phase-card entry is \(u+sv\), where \(u\) and \(v\) are the
contributions of the two Walsh coefficients whose deleted-pair assignments
are \(00\) and \(11\), respectively, and whose residual word is \(y\).  A zero or
cyclic card is diagonal, so this entry vanishes.  At the two actual phases
\(s=1,-1\) we obtain \(u+v=u-v=0\), hence \(u=v=0\).  Running through
all six deleted pairs and both port orders (equivalently, the complete
ordered deck) forces every odd-weight Walsh coefficient to vanish.  Thus
\(\widetilde q\) is supported on even Walsh parity, as required by
\cref{lem:odd-cyc-q4-family}.  Item~1 of
\cref{thm:equality-q4-boundary} handles \(N=2\); hence
assume \(N\ge4\) and use the endpoint and complement-block coordinates of
\cref{eq:odd-dih-q4-coordinates}, indexed by
\[
  I_1=\{1,2\},\qquad I_2=\{1,3\},\qquad I_3=\{1,4\}.
\]
The cyclic quaternary-family result in \cref{lem:odd-cyc-q4-family} leaves two
endpoint regimes.  We treat the zero-endpoint regime first; in the nonzero
regime, the dressing-image criterion further separates successful dressing
from the residual square obstruction.
Define the block phase parameters \(\alpha_i:=\alpha_{I_i}\) and, for
nonzero endpoints, the endpoint ratio \(\lambda:=t/p\).
Call the complementary block indexed by \(I_i\) \emph{active} when
\((x_{I_i},x_{I_i^c})\ne(0,0)\), and \emph{inactive} otherwise.

For each actual two-copy output below, let \(\Lambda^+\) be the retained
augmentation after preprocessing, complete factor saturation, and group
recomputation.  Zero, nullary, unary, odd, singular, and infinite-order
factors take their established exits.  The remaining calculations concern
the all-even finite-group branch.  Proper factors become available over
\(\Lambda^+\) through the factor theorem, with its Turing reduction and
nonzero normalization scalars recorded.

\emph{Zero endpoints.}
Write each complementary block with scale \(\rho_i\) as
\((x_{I_i},x_{I_i^c})=\rho_i(1,\alpha_i)\), taking \(\rho_i=0\)
for an inactive block and then choosing any \(\alpha_i\in\mu_N\), and define the auxiliary scalar
\(X_i:=\alpha_i\rho_i^2\).  For each \(i\), choose the remaining
indices \(j<k\), so \(\{i,j,k\}=\{1,2,3\}\), and order
\(I_i=(1,i+1)\).  Connect the first port
of one copy to the first port of the other by the actual representative
\(D_{\alpha_j\alpha_k}Z\), and connect the second ports by \(I\) (with the
same orientation on both copies).  If \(Q_i\) is the \(I_i\mid I_i^c\)
flattening, let \(\widetilde g_i\) denote the normalized analytic output
\[
 \widetilde g_i=Q_i^{\mathsf T}K_iQ_i,\qquad
 K_i=(D_{\alpha_j\alpha_k}Z)\otimes I
     =\diag(1,1,-\alpha_j\alpha_k,-\alpha_j\alpha_k).
\]
Order the residual ports as the two ports of \(I_i^c\) in the first copy
followed by those in the second copy, each in inherited order.  The only
potentially nonzero complementary pairs of \(\widetilde g_i\) are
\[
 E=(\widetilde g_i(0000),\widetilde g_i(1111)),
 \qquad
 D=(\widetilde g_i(0101),\widetilde g_i(1010)).
\]
Their entries are
\begin{equation*}
 \begin{aligned}
  (E_0,E_1)
    &=(-\alpha_j\alpha_k\rho_i^2,\ \alpha_i^2\rho_i^2),\\
  (D_0,D_1)
    &=(\alpha_j(X_j-X_k),\ -\alpha_k(X_j-X_k)),
 \end{aligned}
\end{equation*}
and, whenever both displayed pairs \(E\) and \(D\) are nonzero,
\[
 \frac{D_1/D_0}{E_1/E_0}=\left(\frac{\alpha_k}{\alpha_i}\right)^2.
\]
Set
\[
 \Delta_i:=(X_j-X_k)^2-X_i^2.
\]
We now verify endpoint nondegeneracy, including the case of an inactive
\(i\)-block.  If \(\rho_i\ne0\), the physical dressing criterion
\cref{lem:odd-cyc-four-port-dressing,eq:cyclic-dressing-image} is applied
in the displayed residual-port order.  Put
\((u,v,w)=(r_{12},r_{13},r_{14})\) for these four output ports.
The entries at \(0101\) and \(1010\) are multiplied by \(R/v\)
and \(v\), respectively.  Take
\[
 R=(E_1/E_0)^{-1},
 \qquad u=1,
 \qquad v=-\alpha_j/\alpha_i,
 \qquad w=R/v.
\]
All four parameters lie in \(\mu_N\), since \(N\) is even, and
\(uvw/R=1\).  Moreover
\[
 E_1R=E_0,\qquad
 D_0R/v=D_1v=\frac{\alpha_j\alpha_k}{\alpha_i}(X_j-X_k).
\]
Thus the two entries in each displayed block are equal.  Since the Walsh
support is even, the two physical endpoints coincide; up to one common
nonzero scalar their common value is
\[
 (X_j-X_k)-X_i.
\]
It is nonzero whenever \(\Delta_i\ne0\).  If instead \(\rho_i=0\),
then \(E_0=E_1=0\), while \(\Delta_i\ne0\) implies
\(D_0D_1\ne0\).  Set
\[
 u=v=1,
 \qquad R=D_1/D_0=-\alpha_k/\alpha_j,
 \qquad w=R.
\]
Here \(uvw/R=1\), so the same criterion applies and makes the two dressed
entries of the \(D\)-block equal and nonzero.  Inverse Walsh transformation again gives two
nonzero physical endpoints.  Thus every index with \(\Delta_i\ne0\),
active or inactive, yields an endpoint-nondegenerate eight-vertex
  signature.  In the displayed four-word matching-code support, the two
  crossed \(2+2\) flattenings have rank at least two; the remaining logical
  flattening has determinant, up to a nonzero scalar,
\begin{equation*}
  \alpha_j\alpha_k\Delta_i.
\end{equation*}
If some \(\Delta_i\ne0\), the corresponding actual output has no binary
factor.  The dressing is a product of nonsingular local maps, so it preserves
tensor factorizations.  Even parity and its complementary endpoints exclude a unary factor,
so the dressed endpoint-nondegenerate eight-vertex signature is tensor-prime.

It remains to exclude simultaneous vanishing of the three minors.  If
\(\Delta_1=\Delta_2=\Delta_3=0\), choose signs
\(\epsilon_i\in\{1,-1\}\) such that
\begin{equation}
 \begin{aligned}
  X_2-X_3&=\epsilon_1X_1,\\
  X_3-X_1&=\epsilon_2X_2,\\
  X_1-X_2&=\epsilon_3X_3.
 \end{aligned}
 \label{eq:proper-cyclic-simultaneous-zero}
\end{equation}
The coefficient matrix of this system is
\[
 M_\epsilon=
 \begin{pmatrix}
  -\epsilon_1&1&-1\\
  -1&-\epsilon_2&1\\
  1&-1&-\epsilon_3
 \end{pmatrix},
 \qquad
 \det M_\epsilon
 =-(\epsilon_1\epsilon_2\epsilon_3+
     \epsilon_1+\epsilon_2+\epsilon_3).
\]
For the two constant sign triples this determinant is nonzero.  For each of
the other six, direct substitution gives, up to permutation, precisely the
kernel lines
\(\Span_{\mathbb C}\{(1,1,0)\}\),
\(\Span_{\mathbb C}\{(1,0,1)\}\), and
\(\Span_{\mathbb C}\{(0,1,1)\}\).  Hence every nonzero solution is a permutation of
  \((\xi,\xi,0)\), with \(\xi\ne0\).  Since \(q\ne0\), not all three
  \(X_i\) vanish.

Relabel so that \((X_1,X_2,X_3)=(\xi,\xi,0)\).  Then
\(\rho_1\rho_2\ne0\) and \(\rho_3=0\).  In the
\(I_3\mid I_3^c\) flattening of \(\widetilde q\), order the odd assignments
on either side as \(10,01\).  Its only nonzero block is
\begin{equation*}
 \begin{pmatrix}
  \rho_1&\rho_2\\
  \alpha_2\rho_2&\alpha_1\rho_1
 \end{pmatrix},
 \qquad
 \det=\alpha_1\rho_1^2-\alpha_2\rho_2^2=X_1-X_2=0.
\end{equation*}
All four entries are nonzero, so the block has rank one and
\(\widetilde q\) has a genuine binary-by-binary tensor factorization across
\(I_3\mid I_3^c\).  The Walsh change is a product of invertible local maps,
so \(q\) has the same genuine factorization.  Its factors would be exposed
only through factor saturation; under the present deck-stable hypotheses
the identity already contradicts that \(q\) is the retained tensor-prime
core.  Therefore some \(\Delta_i\ne0\), as required.

\emph{Nonzero endpoints and successful dressing.}
Define the ratio scalar \(R:=\lambda^{-1}\) and the dressing scalar
\(r_I:=\epsilon_I\alpha_I\lambda^{-1}\) on each active block, where
\(\epsilon_I\in\{\pm1\}\); abbreviate
\(\epsilon_i:=\epsilon_{I_i}\).  An inactive block removes the square
constraint: its otherwise irrelevant dressing parameter may be chosen so
that the criterion's product equation \(uvw/R=1\) holds;
with all three active,
\cref{lem:odd-cyc-four-port-dressing,eq:cyclic-dressing-image} applies
exactly when
\(\epsilon_1\epsilon_2\epsilon_3
 \alpha_1\alpha_2\alpha_3\) is a square.  If
\(N\equiv2\pmod4\), signs remove the obstruction; when \(4\mid N\), this
works precisely if \(\alpha_1\alpha_2\alpha_3\) is a square.  The dressed
endpoint block is \((p,p)\); every active complementary block becomes
\((-\epsilon_Ip,-\epsilon_Ip)\), while every inactive block remains
\((0,0)\).  Let \(h\) be
the resulting physical signature.  Equality of complementary Walsh blocks
makes \(h\) even-parity, while its even Walsh support makes it
  complement-invariant, meaning
  \(h(y)=h(y+\one^4)\) for every \(y\in\Ftwo^4\); it is nonzero because
  \(p\ne0\).  Choose an even
support word \(x\) with \(h(x)\ne0\), and attach actual \(X\)-dressings on
exactly the coordinates where \(x\) is one.  This even translation preserves
parity, and the translated signature \(h'\) satisfies
\[
 h'(0000)=h(x)=h(x+1111)=h'(1111)\ne0.
\]
Thus \(h'\) is an endpoint-nondegenerate eight-vertex signature.  Every
dressing used here is an invertible one-port map, so it also preserves the
tensor-primality of \(q\).

\emph{Square obstruction.}
The only remaining case is
\begin{equation*}
 4\mid N,\qquad
 p,t\ne0,\qquad
 \text{all three blocks active},\qquad
 \alpha_1\alpha_2\alpha_3\notin(\mu_N)^2.
\end{equation*}
  Normalize \(p=1\).  For a permutation \((i,j,k)\) of \((1,2,3)\),
  abbreviate
  \[
   (a,b,c):=(\alpha_i,\alpha_j,\alpha_k),
  \]
  and order
the \(I_i\mid I_i^c\) flattening as \(00,01,10,11\):
\begin{equation*}
 Q_i=
 \begin{pmatrix}
  1&0&0&-a\\
  0&-b&-c&0\\
  0&-\lambda c^{-1}&-\lambda b^{-1}&0\\
  -\lambda a^{-1}&0&0&\lambda
 \end{pmatrix}.
\end{equation*}
Let \(D:=Z\otimes Z\) be the diagonal four-state matrix, and let \(S\) be
the permutation matrix that swaps \(01,10\).  The actual joins are
\begin{align*}
 \widetilde G_i^{\parallel}&=4Q_iDQ_i^{\mathsf T},
 \\
 \widetilde G_i^{\times}&=4Q_iDSQ_i^{\mathsf T}.
\end{align*}
The crossed one is
\begin{equation*}
 \frac14\widetilde G_i^{\times}
 =
 \begin{pmatrix}
  1+a^2&0&0&-\lambda(a+a^{-1})\\
  0&-2bc&-2\lambda&0\\
  0&-2\lambda&-2\lambda^2/(bc)&0\\
  -\lambda(a+a^{-1})&0&0&\lambda^2(1+a^{-2})
 \end{pmatrix}.
\end{equation*}
By \cref{lem:actual-binary-completeness}, an unsafe phase card is already a
binary exit.  Otherwise the cyclic pencil
of the displayed central block forces, when \(1+a^2\ne0\),
\[
 \frac{2\lambda}{1+a^2}\in\mu_N.
\]
Since \(a\) is a root of unity, \(|1+a^2|=2\) gives \(a^2=1\), whereas
\(1+a^2=0\) gives \(a^2=-1\).  Hence
\begin{equation*}
 \alpha_i^2\in\{1,-1\}\qquad(i=1,2,3).
\end{equation*}

If \(8\mid N\), all fourth roots are squares, contradicting the
  obstruction.  Thus \(N\equiv4\pmod8\), and an odd number of the three
  squares equals \(-1\).  Choose the permutation \((i,j,k)\) so that
  \(\alpha_i^2=a^2=-1\), define the ratio scalar
  \(r:=\lambda/(bc)\), and evaluate the straight join.  Its sole nonzero
block has rank one and factors through the nonsingular anti-diagonal matrix
\(A_r\), giving
\begin{equation*}
  \frac14\widetilde G_i^\parallel
  =\pm2 A_r\otimes A_r,
  \qquad
  A_r:=\begin{pmatrix}0&1\\ r&0\end{pmatrix},
  \qquad r\ne0,
\end{equation*}
on the two residual copies of \(I_i\).  \Cref{lem:factor-saturation} exposes
the nonsingular anti-diagonal \(A_r\); its transfer obeys
\[
  (A_rZ)^2=-rI,
  \qquad
  (A_rZ)D_s(A_rZ)^{-1}=sD_{s^{-1}}\quad(s\in\mu_N).
\]
Thus \([A_rZ]\) is a safe involution outside \(C_N\) and generates its finite
dihedral extension.  Concretely, close the factor batch first and recompute
\(G(\Lambda^+)\): nonsingularity follows from \(r\ne0\), the displayed
relations prove finite dihedral closure, and the anti-diagonal projective
class is not in the diagonal group \(C_N\).  Hence this is a strict pulled-back
group successor, not an unprocessed fixed-\((H,\nu)\) factor.  The first two cases
pass an actual endpoint-nondegenerate eight-vertex signature tensor to
\cref{cor:p1-eight-vertex-reduced-interface}, while the third strictly enlarges the
complete group.  This proves the lemma.
\end{proof}

The preceding lemma supplies the quaternary entry to the branch: arity four
is either resolved or passed to a strict outer successor.  With that boundary
in place, the main statement can isolate the only common-presentation
rigidity leaf that remains.

\begin{theorem}[Reed--Muller deck rigidity]
\label{thm:reed-muller-rigidity}
Assume
\[
 \Stable(\cG,G,f),\qquad
 \arity(f)=\nu(\cG),\qquad
 \mathscr K_G=\Span\{I,X\}.
\]
Then \(\Close_{\{\FLag\}}(\cG,G,f)\).  On its sole rigidity branch,
\[
 G=\{[I],[X]\},\qquad
 f\sim\RMcore=\one[RM(1,3)],
\]
where \(\sim\) means equality up to a nonzero scalar, a permutation of
ports, and actual one-port \(X\)-dressings.
\end{theorem}

\subsection{Reconstruction of the minimum Reed--Muller core}

Let \(f\) be a minimum-arity nonbinary prime of the retained signature set,
and let its arity be \(n:=\arity(f)\).  If \(n\) is odd, apply
\cref{thm:external-odd} to the ordinary-coordinate set \(K\cG\).  Its
hardness conclusion pulls back through the holographic \(K\)-coordinate
equivalence, while its tractable conclusion is exactly \(\Tract(K\cG)\).
Thus, on the continuing branch, \(n\) is even.  Every
proper card is zero or a nonsingular \(G\)-matching signature; in Walsh
coordinates the latter is a weighted-equality product.  Define its analytic
Walsh coefficients by
\[
  a_S:=\widetilde f(\one_S),
  \qquad S\subseteq[n].
\]
For distinct \(i,j\), label and normalize the two actual cards obtained
from the \(I/X\)-closures as \(C_{ij}^{+}\) and \(C_{ij}^{-}\) so that,
for every \(T\subseteq[n]\setminus\{i,j\}\),
\begin{equation}
  \widetilde C_{ij}^{+}(\one_T)=2(a_T+a_{Tij}),
  \qquad
  \widetilde C_{ij}^{-}(\one_T)=2(a_T-a_{Tij}),
  \label{eq:rm-card-pencil}
\end{equation}
where \(Tij=T\cup\{i,j\}\).  A nonzero card has the form
\begin{equation}
  c\prod_{\{p,q\}\in M}
  \one[z_p=z_q]r_{pq}^{z_p},
  \qquad c\prod_{\{p,q\}\in M}r_{pq}\ne0,
  \label{eq:rm-weighted-matching}
\end{equation}
for a perfect matching \(M\); it is nonzero at zero, and its supported
weight-two words are exactly the edges of \(M\).  Arity four is
\cref{lem:proper-cyclic-q4}, so assume \(n\ge6\).

Define the zero-card graph \(Z_f\) on \([n]\), writing \(ij\) for the
unordered pair \(\{i,j\}\), by
\[
  ij\in E(Z_f)
  \quad\Longleftrightarrow\quad
  C_{ij}^{+}=C_{ij}^{-}=0.
\]
By \cref{eq:rm-card-pencil},
\begin{equation*}
  ij\in E(Z_f)
  \quad\Longleftrightarrow\quad
  \supp\widetilde f\subseteq\{z:z_i+z_j=1\}.
\end{equation*}

\begin{lemma}[Elimination of the zero-card graph]
\label{lem:rm-zero-graph}
\label{lem:rm-no-zero-edge}
Let \(f\ne0\) be the fixed minimum core of arity \(n\).  Every connected
component of \(Z_f\) is complete bipartite.  If \(n>4\), then
\(E(Z_f)=\varnothing\).
\end{lemma}

\begin{proof}
Zero-edge equations alternate along paths, so every nontrivial connected
component is complete bipartite.  Suppose that \(uv\in E(Z_f)\).  Since
\(n>4\), the set \([n]\setminus\{u,v\}\) has at least three vertices.  The
graph induced there is a disjoint union of bipartite graphs and therefore
is not complete; choose distinct \(i,j\notin\{u,v\}\) with
\(ij\notin E(Z_f)\).  At least one of \(C_{ij}^{+},C_{ij}^{-}\) is then a
nonzero actual card.  Because \(u,v\) remain external and
\(\supp\widetilde f\subseteq\{z:z_u+z_v=1\}\), that card vanishes at the
all-zero boundary word.  This contradicts \cref{eq:rm-weighted-matching},
under which every nonzero card is nonzero at zero.  Hence
\(E(Z_f)=\varnothing\).
\end{proof}

Thus, for the remaining arities, no coordinate pair is invisible to both
physical contractions: every pair supports at least one nonzero matching card.
We may therefore compare the local card matchings through the common pencil
\cref{eq:rm-card-pencil}.  The first global datum extracted from that
comparison is the following weight-two graph.

\paragraph{The weight-two graph.}

Assume henceforth that \(n\ge6\) and \(E(Z_f)=\varnothing\).  Define the
weight-two edge set
\begin{equation*}
  W_f:=\{ij:1\le i<j\le n,\ a_{ij}\ne0\}.
\end{equation*}
Let \(\mathcal W_f:=([n],W_f)\) be the corresponding simple graph.  For
\(S\subseteq[n]\), write
\[
 W_f[S]:=\{pq\in W_f:p,q\in S\},
 \qquad
 W_f[S]-k:=W_f[S\setminus\{k\}]
\]
for the restricted edge sets.

\begin{lemma}[Matching consistency]
\label{lem:rm-matching-consistency}
Let \(f\) be the fixed minimum core of arity \(n\ge6\), and assume
\(E(Z_f)=\varnothing\).  Then the following statements hold.
\begin{enumerate}
\item \(a_\varnothing\ne0\).
\item \(W_f\) is a matching.
\item If both \(C_{ij}^{+}\) and \(C_{ij}^{-}\) are nonzero, their
      support matchings are equal.
\item Every edge \(pq\in W_f\) is a global Walsh equality:
\[
  \supp\widetilde f\subseteq\{z:z_p=z_q\}.
\]
\end{enumerate}
\end{lemma}

\begin{proof}
If \(a_\varnothing=0\), then
\(\widetilde C_{ij}^{\pm}(0)=\pm2a_{ij}\).  Since
\(E(Z_f)=\varnothing\)
and matching cards are nonzero at zero, every \(a_{ij}\ne0\).  The two
\(ij\)-cards would then cover every edge on the remaining at least four
ports, impossible for two matchings.  Hence \(a_\varnothing\ne0\).

For a deleted pair \(i,j\), let \(P,Q\) be the support matchings when both
cards survive, and write \(P\mathbin\triangle Q\) for their symmetric
difference.  Their weight-two entries imply
\begin{equation*}
  W_f[[n]\setminus\{i,j\}]\subseteq P\cup Q,
  \qquad
  P\mathbin\triangle Q
  \subseteq W_f[[n]\setminus\{i,j\}].
\end{equation*}
Indeed, for \(pq\in W_f[[n]\setminus\{i,j\}]\) the two pencil values cannot
both vanish; if one whole card vanishes, the other has value
\(\pm4a_{pq}\) at \(pq\).  Thus, with one nonzero card,
\(W_f[[n]\setminus\{i,j\}]\) lies in its perfect
matching.  We use the following elementary graph argument.  If a vertex of
\(\mathcal W_f\) had degree at least three, delete two vertices outside it and its
three neighbours (possible since \(n\ge6\)); the induced graph has degree at
least three, whereas the union of two perfect matchings has maximum degree
two (and one matching has maximum degree one), a contradiction.  Hence every
component of \(\mathcal W_f\) is a path or a cycle.

Suppose a path component has at least two edges.  If two vertices lie outside
that component, delete them; otherwise (then the component has at least
\(n-1\ge5\) vertices) delete two consecutive vertices at one end, leaving a
path with at least two edges.  For this deleted pair, one nonzero card is
impossible.  If both cards are nonzero, the two edges at the middle of the
retained path must lie one in
\(P\) and one in \(Q\), hence in \(P\triangle Q\).  Alternating along the
path forces an edge of \(P\triangle Q\) at an endpoint, but
\(P\triangle Q\) is a disjoint union of cycles and
\(P\triangle Q\subseteq W_f[[n]\setminus\{i,j\}]\), whose endpoint has no
second edge.  This is impossible.  A cycle is handled similarly: a triangle
is retained by deleting vertices outside it and cannot alternate; for a
4-cycle delete one cycle vertex and one outside vertex, and for a cycle of
length at least five delete two adjacent cycle vertices, leaving a path with
at least two edges.  Thus every component has at most one edge, so \(W_f\)
is a matching.  Finally, for any pair with both cards nonzero,
\(P\triangle Q\) is a union of cycles contained in the matching \(W_f\);
therefore \(P\triangle Q=\varnothing\) and \(P=Q\).

For \(pq\in W_f\), suppose \(z\in\supp\widetilde f\) and
\(z_p\ne z_q\).  Choose distinct \(i,j\in[n]\setminus\{p,q\}\) with
\(z_i=z_j\).  Both \(ij\)-cards vanish on this restriction because every
nonzero one contains \(pq\); \cref{eq:rm-card-pencil} then kills both
lifts, contradicting \(z\in\supp\widetilde f\).  Hence \(z_p=z_q\) on the
support.
\end{proof}

When \(W_f\ne\varnothing\), \cref{lem:rm-matching-consistency} supplies
global equality constraints.  The remaining issue is whether the
coefficients separate across these pairs.  The following lemma recovers
that product structure from diagonal restrictions.

\begin{lemma}[Diagonal-slice integration]
\label{lem:diagonal-slice-integration}
Let \(r\ge3\), and let \(h:\Ftwo^r\to\mathbb C^\times\) satisfy
\(h(0)=1\).  Suppose that, for every \(i\ne j\), the restriction
\(h|_{t_i=t_j}\), viewed as a function of the common bit and the other
\(r-2\) bits, is a product of unary functions (equivalently, is a simple
tensor).  If \(r\ge4\), then
\[
  h(t)=\prod_{i=1}^r\rho_i^{t_i}
  \qquad(\rho_i\in\mathbb C^\times).
\]
If \(r=3\), there are
\(\rho_1,\rho_2,\rho_3,\kappa\in\mathbb C^\times\) such that, for distinct
\(i,j\),
\begin{equation}
  h(e_i)=\rho_i,\qquad
  h(e_i+e_j)=\kappa\rho_i\rho_j,\qquad
  h(111)=\kappa\rho_1\rho_2\rho_3.
  \label{eq:diagonal-slice-q3-table}
\end{equation}
\end{lemma}

\begin{proof}
Use the multiplicative subset expansion
\begin{equation}
  h(t)=\prod_{\varnothing\ne S\subseteq[r]}
       \theta_S^{\prod_{i\in S}t_i},
   \qquad \theta_S\in\mathbb C^\times.
  \label{eq:multiplicative-moebius-expansion}
\end{equation}
Multiplicative M\"obius inversion determines its parameters:
\[
  \theta_S
  =\prod_{T\subseteq S}h(\one_T)^{(-1)^{|S|-|T|}}.
\]
The parameter \(\theta_S\) measures the interaction among the coordinates
in \(S\); unary factors account exactly for the parameters with
\(|S|=1\).  Thus the table is simple if and only if
\(\theta_S=1\) for every \(|S|\ge2\).  On setting \(t_i=t_j\), a parameter indexed by a
nonempty \(U\subseteq[r]\setminus\{i,j\}\) and the common bit is
\begin{equation}
  \theta_{U\cup\{i\}}
  \theta_{U\cup\{j\}}
  \theta_{U\cup\{i,j\}}=1.
  \label{eq:diagonal-slice-moebius-relation}
\end{equation}
Parameters \(\theta_S\) with \(S\cap\{i,j\}=\varnothing\) and
\(|S|\ge2\) must also equal one.

Assume \(r\ge4\).  Choosing \(i,j\) outside \(S\) first gives
\(\theta_S=1\) for \(2\le|S|\le r-2\).  For \(|S|=r-1\), choose distinct
\(i,j\in S\) and define the index set \(U:=S\setminus\{i,j\}\) in
\cref{eq:diagonal-slice-moebius-relation}; the other two factors have
index-set cardinality \(r-2\), so \(\theta_S=1\).  Taking
\(U=[r]\setminus\{i,j\}\) then gives \(\theta_{[r]}=1\).  Only the unary
parameters remain in \cref{eq:multiplicative-moebius-expansion}, proving
the product formula.

Let \(r=3\).  The three instances of
\cref{eq:diagonal-slice-moebius-relation} imply
\[
  \theta_{12}=\theta_{13}=\theta_{23}=:\kappa,
  \qquad \theta_{123}=\kappa^{-2}.
\]
Substitution in \cref{eq:multiplicative-moebius-expansion}, with
\(\rho_i=\theta_{\{i\}}\), gives
\cref{eq:diagonal-slice-q3-table}.
\end{proof}

\begin{lemma}[Integration of a nonempty weight-two graph]
\label{lem:rm-weight-two-integration}
Let \(f\) be the fixed minimum core of arity \(n\ge6\), with
\(E(Z_f)=\varnothing\).  If \(W_f\ne\varnothing\), then either \(n=6\) and a
constant-size two-copy gadget realizes the even-parity signature of arity
four, or \(f\) admits a nontrivial tensor factorization.  In particular, a
minimum core of arity
at least eight has \(W_f=\varnothing\).
\end{lemma}

\begin{proof}
Fix \(E_1=\{u,v\}\in W_f\) and define the residual port set
\(R:=[n]\setminus E_1\).  For \(k\in R\), let
\(A_k:=C_{uk}^{+}\).  Since
\(\widetilde A_k(0)=2a_\varnothing\ne0\), this is a nonzero matching card.
The global equality on \(E_1\) gives the unique lift
\begin{align*}
  \widetilde A_k(z_v,z_{R\setminus\{k\}})
  &=2\widetilde f(z_u=z_v,z_k=z_v,z_{R\setminus\{k\}}),\\
  \widetilde A_k(0)&=2a_\varnothing\ne0,
  \\
  \widetilde A_k(e_p+e_q)&=2a_{pq}
  \quad(p,q\in R\setminus\{k\}),\\
  \widetilde A_k(e_v+e_q)&=2a_{\{u,v,k,q\}}
  \quad(q\in R\setminus\{k\}).
\end{align*}
Its matching therefore consists of \(W_f[R]-k\) and at most one edge from
\(v\).  To justify the last assertion, let \(U\) be the set of vertices not
covered by \(W_f\); it excludes \(u,v\).  For every \(k\), all vertices of
\(U\setminus\{k\}\) remain on the card boundary and have no available
weight-two edge, except that at most one can be paired through \(v\).
If \(W_f=\{E_1\}\), then \(|U|=n-2\ge4\), and the preceding restriction
already gives a contradiction.  Otherwise a covered vertex of \(W_f\) lies
in \(R\), so choose \(k\in R\setminus U\); hence \(|U|\le1\).
Since \(|U|=n-2|W_f|\) is even, \(U=\varnothing\).  Thus \(W_f\) is
perfect; write
its edges \(E_1,\ldots,E_s\), where \(s=n/2\).

By the global equality conclusion of
\cref{lem:rm-matching-consistency}, for every word \(z\) constant on each
\(E_r\), let \(t_r(z)\) denote that common edge bit.  There is a logical
table \(h:\Ftwo^s\to\mathbb C\) such that
\begin{equation*}
  \widetilde f(z)=
  \begin{cases}
    h(t_1(z),\ldots,t_s(z)),
      &z|_{E_r}\text{ is constant for every }r,\\
    0,&\text{otherwise}.
  \end{cases}
\end{equation*}
For two distinct \(W_f\)-edges \(E_i,E_j\), close one endpoint of \(E_i\)
to one endpoint of \(E_j\), forming separately the two cards with the
actual \(+\) and \(-\) kernels.  Since
cross coefficients between different \(W_f\)-edges vanish and
\(a_\varnothing\ne0\), both cards are nonzero.  Both matching supports
impose \(t_i=t_j\); the \(+/-\) choice changes only the sign character in
that common logical bit.  Since \(s=n/2\ge3\),
every \(t\in\Ftwo^s\) has a pair \(i\ne j\) with \(t_i=t_j\); use the
plus card for that pair.  The boundary word obtained from \(t\)
by retaining the mates of \(E_i,E_j\) and the two ports of every other
\(E_r\) lies in this card's matching support.  The closure equation has a
unique compatible lift in the global equality code, namely the word with
edge bits \(t\), so the plus-card value is a fixed nonzero scalar times \(h(t)\)
(there is no second lift with which it could cancel).  A nonzero matching
product is nonzero at every supported boundary word; hence \(h(t)\ne0\) for
every \(t\), and \(h\) has full support.
Rescale it by \(a_\varnothing^{-1}\) (absorbing
this nonzero scalar into \(f\)), so that \(h(0)=1\).  Every cross-pair
diagonal slice \(t_i=t_j\) is nonzero and rank one.  By
\cref{lem:diagonal-slice-integration},
\(s\ge4\) would make \(f\) a product of \(s\) nonsingular binaries, which
is the factorization alternative of the lemma.  Hence, on the remaining
branch, \(s=3\).  Rename the parameters
\(\rho_1,\rho_2,\rho_3,\kappa\ne0\) from
\cref{eq:diagonal-slice-q3-table} as
\(r_1,r_2,r_3,\kappa\).  Let \(H_i^\pm\) be the \(2\times2\)
logical coefficient matrix obtained by closing \(E_i\) with the \(\pm\)
kernel.  For \(\{i,j,k\}=\{1,2,3\}\), one has
\begin{equation*}
  \det H_i^{\pm}
  =r_j r_k(\kappa-1)(1-\kappa r_i^2).
\end{equation*}
If \(\kappa=1\), \(h\), and hence \(f\), has the asserted product
factorization.  On the remaining branch \(\kappa\ne1\), and rank one of the
nonzero cards forces
\begin{equation*}
  \kappa r_1^2=\kappa r_2^2=\kappa r_3^2=1.
\end{equation*}

In this remaining case, write each matching edge as its two port labels, e.g.
\(E_i=\{i^0,i^1\}\), in the Walsh-index matching (this notation does not
assert that the physical tensor itself is an equality signature).  An
\emph{actual \(I/X\)-link} is the constant-size binary chain allowed by the
kernel-span convention: in Walsh coordinates its kernel is a nonzero
scalar times \(\epsilon^u\one[u=v]\), with \(\epsilon=1\) for \(I\) and
\(\epsilon=-1\) for \(X\).  On a paired join with common bit \(u=v\), the
two link signs \(\epsilon_0,\epsilon_1\) therefore contribute
\((\epsilon_0\epsilon_1)^u\); choosing their parity realizes either desired
character without treating \(I\) as a structural edge.

Let \(E'_1,E'_2,E'_3\) denote the corresponding matching edges in the second
copy.  Use two actual \(I/X\)-links for each displayed edge-to-edge join:
join the two ports of \(E_3\) in the first copy to the two ports of \(E'_1\)
in the second copy, and join the two ports of \(E'_2\) to those of \(E'_3\)
inside the second copy; leave both ports of \(E_1\) and \(E_2\) in the first
copy external.  Evaluating these actual contractions by
\cref{eq:rm-card-pencil} gives the logical characters below, and every
internal matching assignment has a unique lift.  Choose the actual
\(I/X\)-signs
\[
  \theta=\kappa r_2 r_3,
  \qquad
  \eta=-\kappa r_1 r_3,
\]
which lie in \(\{\pm1\}\).  Writing \(x=t_1\) and \(y=t_2\) for the two
external logical bits, the external table is
\begin{equation*}
  Q(x,y)=
  \sum_{a,c\in\Ftwo}
  \eta^a\theta^c h(x,y,a)h(a,c,c).
\end{equation*}
Using \(\sum_c\theta^c h(a,c,c)=2r_1^a\), this becomes
\begin{equation*}
  Q\doteq
  \frac{\kappa-1}{\kappa}
  \begin{pmatrix}1&0\\0&\kappa r_1 r_2\end{pmatrix}.
\end{equation*}
The calculation is in the analytic Walsh coordinates of
\cref{eq:Walsh-coordinates}; it is not yet the physical boundary tensor.
In the logical Walsh variables the quaternary tensor therefore is
\[
 \widetilde F
 =c\bigl(\delta_{0000}+\sigma\delta_{1111}\bigr),
 \qquad
 c\ne0,
 \qquad \sigma=\kappa r_1r_2\in\{\pm1\},
\]
where the two nonzero coefficients follow from \(\kappa\ne1\).  Applying
the inverse analytic change of coordinates, namely
\(F=P^{\otimes4}\widetilde F\) with
\(P=\left(\begin{smallmatrix}1&1\\1&-1\end{smallmatrix}\right)\), gives
\[
 F(x)=c\bigl(1+\sigma(-1)^{x_1+x_2+x_3+x_4}\bigr).
\]
Define the even-parity signature by
\[
 \mathrm{Even}_4(x):=\one[x_1+x_2+x_3+x_4=0].
\]
Thus the actual tensor is a nonzero scalar multiple of
\(\mathrm{Even}_4\) when \(\sigma=1\), and of its odd-parity coset when
\(\sigma=-1\); an actual \(X\)-dressing on one exposed port toggles the
coset.  The Walsh matrix is used only for this analytic reconstruction and
is not being installed as a physical gadget.  For completeness, the
equality-anchor witness for \(\mathrm{Even}_4\) is explicit: with the
diagonal matrix \(D\), Walsh matrix \(P\), and transform matrix \(U\) defined by
\[
 D:=\operatorname{diag}(1,i),\qquad
 P:=\begin{pmatrix}1&1\\1&-1\end{pmatrix},\qquad
 U:=DP^{-1},
\]
one has
\[
 U^{\otimes4}\mathrm{Even}_4\doteq\mathrm{Eq}_4,\qquad
 (U^{-1})^{\mathsf T}XU^{-1}
 =D^{-1}P^{\mathsf T}XP D^{-1}=2I=2\,\mathrm{Eq}_2.
\]
Consequently an actual \(\mathrm{Even}_4\) equality anchor is obtained, and
\cref{lem:equality-anchor} applies.
\end{proof}

A nonempty \(W_f\) has therefore already produced an equality anchor or a
proper factorization.  After those alternatives are dispatched, the
rigidity branch continues only with
\[
  W_f=\varnothing.
\]
Scale \(a_\varnothing=1\).  With all weight-two coefficients absent, the
common card matchings determine the entire Walsh support.

\begin{lemma}[Reed--Muller support reconstruction]
\label{lem:rm-local-slices}
\label{lem:rm-support-classification}
Let \(f\) be the fixed minimum core of arity \(n\ge6\), normalized by
\(a_\varnothing=1\), and assume \(E(Z_f)=\varnothing\) and
\(W_f=\varnothing\).  For \(i\ne j\), let \(M_{ij}\) be the common support
matching of the two \(ij\)-cards.  Fix an arbitrary ordering of its
\(n/2-1\) edges and, for \(t\in\Ftwo^{n/2-1}\), write \(t_e\) for the
coordinate indexed by \(e\in M_{ij}\).  Let
\(D:=\supp\widetilde f\) be the support set.  Its equal-\(ij\) slice is
the \((n/2-1)\)-dimensional code
\begin{equation}
  D_{ij}=
  \left\{
    \one_{T(t)}+(\wt(t)\bmod2)(e_i+e_j):
    t\in\Ftwo^{n/2-1}
  \right\},
  \label{eq:rm-local-code}
\end{equation}
where
\(
  T(t)=\bigcup_{e\in M_{ij}:\,t_e=1}e.
\)
Every word of \(D_{ij}\) has weight divisible by four.  Moreover,
\[
  n=8,
  \qquad
  D=RM(1,3),
\]
after a bijective relabelling of the ports by \(\Ftwo^3\).
\end{lemma}

\begin{proof}
By \cref{lem:rm-matching-consistency}, the two \(ij\)-cards survive with a
common support matching \(M_{ij}\).  Normalize them at the zero word
\(t=0\) in matching coordinates, and write \(r_{ij,e}^{+}\) and
\(r_{ij,e}^{-}\) for the corresponding normalized edge ratios (the
\(t_e=1\) factor entry divided by its \(t_e=0\) entry).  For each
\(e\in M_{ij}\), set \(\rho_{ij,e}:=r_{ij,e}^{-}/r_{ij,e}^{+}\).  Their
quotient on the matching code---the repetition code consisting of assignments
constant on every edge of \(M_{ij}\), parametrized by \(t\)---is
\[
  \frac{\widetilde C_{ij}^{-}(t)}
       {\widetilde C_{ij}^{+}(t)}
  =\prod_{e\in M_{ij}}\rho_{ij,e}^{t_e}.
\]
Each single-edge quotient is \(-1\) by \(a_{pq}=0\) and
\cref{eq:rm-card-pencil}; hence the quotient is \((-1)^{\wt(t)}\).  Inverting
the pencil selects the \(ij=00\) lift for even \(\wt(t)\) and \(ij=11\) for
odd \(\wt(t)\), proving \cref{eq:rm-local-code}; the selected word has
weight \(2\wt(t)\) in the even case and \(2\wt(t)+2\) in the odd case, so
its weight is divisible by four in either case.

For \(x,y\in D\), some two coordinates have the same pair
\((x_i,y_i)\); the linear slice \(D_{ij}\) then contains \(x,y,x+y\).
Thus \(D\) is linear.  For the linear functional
\(\ell_{ij}:\Ftwo^n\to\Ftwo\) defined by
\(\ell_{ij}(z):=z_i+z_j\), one has
\[
  D\cap\ker\ell_{ij}=D_{ij}.
\]
This subspace has dimension \(n/2-1\).  If
\(\dim D=n/2-1\), then \(D=D_{ij}\subseteq\ker\ell_{ij}\) for every
\(i\ne j\).  Hence every word of \(D\) has all coordinates equal, so
\(\dim D\le1\), contradicting \(n/2-1\ge2\).  Therefore
\(\dim D=n/2\).
Every word of \(D\) lies in one of the displayed slices and hence has weight
divisible by four.  For \(x,y\in D\),
\[
 0\equiv \wt(x)+\wt(y)-\wt(x+y)=2x\cdot y\pmod 4.
\]
Thus \(D\subseteq D^\perp\), and equal dimensions give
\begin{equation*}
  D=D^\perp.
\end{equation*}

For distinct \(i,j,p\), the mate \(q\) of \(p\) in \(M_{ij}\) gives the
unique weight-four word \(\{i,j,p,q\}\).  Now
\cref{lem:self-dual-steiner-reconstruction} yields
\(n=8\) and \(D=RM(1,3)\).
\end{proof}

\paragraph{Coefficient rigidity and \texorpdfstring{\(C_2\)}{C2} collapse.}

The combinatorial reconstruction is complete: the surviving core has arity
eight and Reed--Muller support.  It remains to control its nonzero
coefficients.  The next lemma shows that they carry only a linear sign
character, so an available \(X\)-dressing removes the resulting affine
shift.

\begin{lemma}[Coefficients on \(\mathrm{RM}(1,3)\)]
\label{lem:rm-coefficient-atlas}
Let \(f\) be an eight-port signature with
\(\supp\widetilde f=RM(1,3)\) whose proper \(I/X\)-cards are zero or have
the weighted-equality matching form of
\cref{eq:rm-weighted-matching}.  There are \(c\ne0\) and
\(w\in\Ftwo^8\) such
that
\begin{equation}
  \widetilde f(z)
  =c(-1)^{w\cdot z}\one[z\in RM(1,3)].
  \label{eq:rm-character-form}
\end{equation}
Consequently, back in the physical \(K\)-coordinates,
\begin{equation}
  f(x)\doteq\one[x\in w+RM(1,3)].
  \label{eq:rm-coset-form}
\end{equation}
\end{lemma}

\begin{proof}
Let the support code be \(D:=RM(1,3)\), and define
\(b(z):=\widetilde f(z)/\widetilde f(0)\).  For an affine plane
\(A=a+L\subseteq\Ftwo^3\), write
\(\operatorname{inc}(A)\in\Ftwo^8\) for its incidence vector and identify
\(A\) with \(\operatorname{inc}(A)\in D\).

Fix distinct points \(p,q\in\Ftwo^3\), and let
\(A_1,A_2,A_3\) be the three affine planes containing them.  Under the
identification of the equal-\(pq\) plus card with its three logical
matching coordinates, the words \(A_1,A_2,A_3\) are the three singleton
logical words.  Its zero word lifts only to \(0\in D\), because \(D\)
contains no word of weight two.  Normalizing the card at that word, the
weighted-product form therefore gives
\begin{equation}
 b\left(\sum_{r\in S}A_r\right)
   =\prod_{r\in S}b(A_r)
 \qquad\left(S\subseteq\{1,2,3\}\right).
 \label{eq:rm-plane-slice-product}
\end{equation}
In particular, \(A_1+A_2+A_3=\one^8\).  The same argument shows that
whenever two affine planes \(A,B\) meet in a two-point affine line,
\begin{equation}
 b(A+B)=b(A)b(B).
 \label{eq:rm-intersecting-plane-product}
\end{equation}

For a plane \(A\), write \(A^c:=\Ftwo^3\setminus A\).  Now
\(A_1+A_2=A_3^c\) is itself an affine plane.  Since
\(A_2\cap A_3=\{p,q\}\), it meets \(A_2\) in the other two points of
\(A_2\), hence in a two-point line.  Applying
\cref{eq:rm-intersecting-plane-product} first to \(A_1,A_2\) and then to
\(A_1+A_2,A_2\) gives
\[
 b(A_1+A_2)=b(A_1)b(A_2),
 \qquad
 b(A_1)=b(A_1+A_2)b(A_2).
\]
Thus \(b(A_2)^2=1\).  Every affine plane can play the role of \(A_2\), so
\begin{equation}
 b(A)\in\{\pm1\}
 \qquad\text{for every affine plane }A.
 \label{eq:rm-plane-signs}
\end{equation}
For any affine plane \(A\), choose distinct \(p,q\in A\) and call the other
two planes through them \(B,C\).  Then \(A^c=B+C\), and
\cref{eq:rm-plane-slice-product,eq:rm-plane-signs} give
\begin{equation}
 b(A^c)=b(B)b(C)=b(\one^8)b(A).
 \label{eq:rm-complement-plane-product}
\end{equation}

The code \(D\) consists of \(0,\one^8\), and the fourteen affine-plane
incidence words.  We can now check multiplicativity on all of \(D\).
The zero case is immediate.  Equal planes use
\cref{eq:rm-plane-signs}; distinct noncomplementary planes meet in a
two-point line and use \cref{eq:rm-intersecting-plane-product};
complementary planes and sums of \(\one^8\) with a plane use
\cref{eq:rm-complement-plane-product}.  Finally,
\cref{eq:rm-plane-slice-product,eq:rm-plane-signs} imply
\(b(\one^8)^2=1\), which handles \(\one^8+\one^8\).  Hence
\[
 b(x+y)=b(x)b(y)
 \qquad(x,y\in D).
\]
Thus \(b\) is a sign character of \(D\).  Its binary exponent extends to a
linear functional on \(\Ftwo^8\), so there is \(w\in\Ftwo^8\) with
\(b(z)=(-1)^{w\cdot z}\) for \(z\in D\).  This proves
\cref{eq:rm-character-form}.  Since \(D=D^\perp\), Fourier inversion gives
\[
  f(x)\doteq\sum_{z\in D}(-1)^{(x+w)\cdot z}
  =|D|\one[x+w\in D].
\]
This is \cref{eq:rm-coset-form}.
\end{proof}

Available \(X\)-dressings translate the coset, so the standard core
\(\RMcore\) of \cref{eq:RM-core} is available.  In lexicographic order it is the
Shao--Cai signature \(f_8=\widehat f_8\); see \cref{rem:sc-f8-core}.

\begin{remark}[Limited appearance and internal \(C_2\) rigidity]
\label{rem:rm-binary-enlargement-shortcut}
The limited-appearance construction of Shao and Cai
\cite[full version, Section~8.2 and Lemma~8.10]{ShaoCai2020} provides a
related reduction in their real-valued setting.  For the complex-valued
signature sets considered here, the present branch is discharged by the
independent \(C_2\) rigidity argument below.  On the remaining
continuation, it first proves that the complete group is exactly
\(\{[I],[X]\}\cong C_2\), and then determines its internal structure.  The exact
Reed--Muller network calculus, actual Radon access, localization, affine
patching, and recognition prove that the entire retained signature set,
together with its edge, has one common flat-Lagrangian presentation.

The final matching-deck synthesis invokes
\cref{thm:reed-muller-rigidity} to resolve the proper-\(C_2\) case.
The no-inverse result \cref{cor:no-rm-to-H6} also records why the later
\(\Hcore\)-to-\(\RMcore\) circuit cannot be reversed using only pure
\(\RMcore\)/Pauli resources.  We retain the internal network calculus,
Radon access, patching, and recognition arguments below.
\end{remark}

Thus the core itself is now fixed, up to the recorded equivalences.  The last
initial-core issue is the complete transfer group.  On the continuing
matching-deck branch, the coefficient classification also rules out every safe class
beyond \([I]\) and \([X]\).

\begin{lemma}[Safe-group collapse]
\label{lem:rm-safe-group-collapse}
Let \(\cG\) be a retained normalized signature set, let
\(G:=G(\cG)\) be its complete transfer group, and assume that its
kernel span \(\mathscr K_G\) is
\(\mathscr K_G=\Span\{I,X\}\).  Let \(f\in\cG\) be an irreducible
Reed--Muller core such that every safe one-port dressing of \(f\) remains
in the standing matching-deck branch.  Then the complete safe projective
transfer group is
\[
  \{[I],[X]\}\cong C_2.
\]
\end{lemma}

\begin{proof}
Put \(D:=RM(1,3)\).  By \cref{lem:rm-coefficient-atlas}, choose
\(c\ne0\) and \(w\in\Ftwo^8\) such that
\(\widetilde f(z)=c(-1)^{w\cdot z}\) for \(z\in D\), and define
\(b(z):=\widetilde f(z)/c\).  Let \(T\) be a nonsingular safe physical
transfer, and write
\[
 T^{\mathrm W}:=P^{-1}TP
\]
for its Walsh-coordinate matrix.  In the proper algebra,
\(T^{\mathrm W}\doteq\diag(1,r)\) for some \(r\ne0\).
Dress one port \(p\) by \(T\).  This actual constant-size gadget preserves
irreducibility, support, and the matching-deck hypotheses; its coefficient
table, after division by its value at \(z=0\), is
\[
 b_T(z)=(-1)^{w\cdot z}r^{z_p}
 \qquad(z\in D).
\]
By \cref{lem:rm-coefficient-atlas}, \(b_T\), hence also the quotient
\(r^{z_p}=b_T(z)/b(z)\), is a character of \(D\) and takes values
\(\pm1\).  Evaluating at \(\one^8\in D\) gives \(r=\pm1\), precisely the
physical transfer classes \([I]\) and \([X]\).
\end{proof}

\subsection{Flat-Lagrangian networks and Radon access}

Retain the literal \(\RMcore\) to build contexts around an arbitrary
signature.  A flat-Lagrangian signature is constant on an affine coset
\(a+L\) with \(L=L^\perp\); see
\cref{subsubsec:symmetric-binary-spaces,subsec:standard-classes}.  The next lemma
shows both what these networks preserve and what they can construct.
Self-duality tracks the boundary after internal variables are summed out;
weight-four transvections provide the reverse construction from a basis
of \(L\).

\begin{lemma}[Exact flat-Lagrangian network calculus]
\label{lem:rm-pure-boundary}
\label{lem:rm-flat-contexts}
\begin{enumerate}
\item Let \(\Omega\) be a tensor network with ordered boundary \(B\).
      Suppose that the signature at each tensor vertex \(u\) is
      proportional to \(\one[x_u\in a_u+L_u]\), where either
      \(L_u=RM(1,3)\le\Ftwo^8\) or
      \(L_u=\Span_{\Ftwo}\{(1,1)\}\le\Ftwo^2\), and \(a_u\)
      is a shift in the corresponding ambient space.  The binary case
      includes ordinary \(I/X\)-vertices.  Suppose every internal
      connection is labelled by \(I\) or \(X\).  Then the boundary signature
      is either zero or
\begin{equation}
  x_B\longmapsto c_B\one[x_B\in a_B+L],
  \qquad c_B\ne0,\quad a_B\in\Ftwo^B,\quad L=L^\perp\le\Ftwo^B.
  \label{eq:rm-pure-flat}
\end{equation}
\item Conversely, for every \(a_B\in\Ftwo^{2k}\) and every
      \(L=L^\perp\le\Ftwo^{2k}\), a polynomial-size
      \(\RMcore/I/X\)-network directly realizes a nonzero scalar multiple of
      \(x_B\longmapsto\one[x_B\in a_B+L]\).
\end{enumerate}
In particular every nonzero boundary produced by a network in item~1 has
even arity.
\end{lemma}

\begin{proof}
For item~1, first apply the boundary normalization following
\cref{eq:half-edge-boundary-tensor}: remove one-free-endpoint tails and
separate isolated \(I/X\)-edge boxes.  Each isolated box is itself a
flat-Lagrangian binary and is restored at the end by a direct sum of
support directions.  The remaining ordered boundary \(B\) therefore
consists of unpaired ordinary vertex ports, with no free stub or already
paired boundary ports.  Let \(E_{\rm int}\) be the set of internal
half-edge pairs and put
\[
 H_{\rm int}:=\bigcup E_{\rm int},
 \qquad
 V_{\rm he}:=\Ftwo^B\oplus\Ftwo^{H_{\rm int}}.
\]
For each \(e=\{h,h'\}\in E_{\rm int}\), let \(\eta_e=0\) when \(e\) is
labelled by \(I\) and \(\eta_e=1\) when it is labelled by \(X\).  Define
\[
 E_\eta:=
 \{y\in\Ftwo^{H_{\rm int}}:
   y_h+y_{h'}=\eta_e
   \text{ for every }e=\{h,h'\}\in E_{\rm int}\},
\]
and let
\[
 E_{\rm eq}:=
 \{y\in\Ftwo^{H_{\rm int}}:
   y_h=y_{h'}\text{ for every }\{h,h'\}\in E_{\rm int}\},
 \qquad
 H_{\rm wire}:=\Ftwo^B\oplus E_{\rm eq}.
\]
The direction \(E_{\rm eq}\) is self-dual in
\(\Ftwo^{H_{\rm int}}\), and hence
\[
 H_{\rm wire}^\perp
 =\{\mathbf 0_B\}\oplus E_{\rm eq},
\]
where, for any finite index set \(R\), \(\mathbf 0_R\) denotes the zero
vector in \(\Ftwo^R\).

Let \(V(\Omega)\) denote the tensor-vertex set of \(\Omega\), and put
\[
 C_0:=
 \bigoplus_{u\in V(\Omega)}L_u
 =C_0^\perp,
 \qquad
 \mathcal K:=C_0\cap H_{\rm wire}.
\]
Let \(a_{\rm vert}:=(a_u)_{u\in V(\Omega)}\) be the concatenated vertex
shift.  The feasible assignments form
\[
 (a_{\rm vert}+C_0)
 \cap\bigl(\Ftwo^B\oplus E_\eta\bigr).
\]
This intersection is either empty or an affine coset
\(y_0+\mathcal K\).  Define
\[
 \pi_B:V_{\rm he}\to\Ftwo^B,
 \qquad
 L:=\pi_B(\mathcal K),
 \qquad
 a_B:=\pi_B(y_0).
\]
Every nonempty boundary fibre is a coset of
\(\ker(\pi_B|_{\mathcal K})\), so all fibres have size
\(2^{\dim\ker(\pi_B|_{\mathcal K})}\).  Multiplying this number by
the recorded nonzero vertex scalars gives \(c_B\ne0\) in
\cref{eq:rm-pure-flat}.  The support and scalar are therefore determined
by affine Gaussian elimination.

To prove \(L\subseteq L^\perp\), take \(k,k'\in\mathcal K\).  Since
\(k,k'\in C_0=C_0^\perp\), we have \(k\cdot k'=0\).  For an internal pair
\(e=\{h,h'\}\), put
\[
 t:=k_h=k_{h'},
 \qquad
 s:=k'_h=k'_{h'}.
\]
That pair contributes \(ts+ts=0\) to \(k\cdot k'\).  Consequently
\(\pi_B(k)\cdot\pi_B(k')=0\), proving \(L\subseteq L^\perp\).

Moreover,
\[
 \mathcal K^\perp=C_0+H_{\rm wire}^\perp.
\]
If \(b\in L^\perp\), then
\((b,\mathbf 0_{H_{\rm int}})\in\mathcal K^\perp\), so write
\[
 (b,\mathbf 0_{H_{\rm int}})=c_0+d,
 \qquad
 c_0\in C_0,\quad d\in H_{\rm wire}^\perp.
\]
Both \((b,\mathbf 0_{H_{\rm int}})\) and \(d\) belong to
\(H_{\rm wire}\), hence \(c_0\in C_0\cap H_{\rm wire}=\mathcal K\).
Therefore \(b=\pi_B(c_0)\in L\), proving \(L=L^\perp\).

For the converse, order one RM vertex's input ports as
\((000,100,010,001)\) and output ports as \((111,011,101,110)\).  Direct
evaluation gives
\[
 y=(I_4+J_4)x=x+(x_1+x_2+x_3+x_4)\one^4.
\]
Here \(x,y\in\Ftwo^4\) are respectively the ordered input and output words,
\(I_4\) is the \(4\times4\) identity matrix, and
\(J_4:=\one^4(\one^4)^{\mathsf T}\) is the \(4\times4\) all-ones matrix over
\(\Ftwo\).
Thus actual RM boxes and port permutations realize every weight-four
transvection \(\tau_s(x)=x+(s\cdot x)s\); retained \(I\)-vertices (or
\(X-I-X\) chains) compose them.  We claim by induction on \(k\) that every
\(L=L^\perp\le\Ftwo^{2k}\) is obtained from
 \(L_0=\Span_{\Ftwo}\{e_1+e_2,e_3+e_4,\ldots\}\).  The cases \(k=0,1\) are
immediate.  For \(k\ge2\), self-orthogonality gives \(x\cdot x=0\) for
every \(x\in L\), so all words are even and \(\one\in L\).  Choose
\(z\in L\setminus\{0,\one\}\); one of \(z,z+\one\) has even weight in
\([2,k]\).  If its weight is larger than two, three support coordinates and
one zero coordinate give a weight-four \(s\) with \(z\cdot s=1\) and
\(\wt(z+s)=\wt(z)-2\).  Apply the corresponding transvection to the whole
current code and repeat, then permute coordinates, to obtain
\(p=e_1+e_2\in L\).  Every codeword has equal first two coordinates.
 Projection deleting them has kernel \(\Span_{\Ftwo}\{p\}\), image dimension
\(k-1\), and preserves dot products (the two deleted contributions cancel),
so its image is self-dual.  Recurse and invert the accumulated actual
transvections and permutations.  The resulting circuit has size \(O(k^2)\).
Closing the standard input pairs by \(I\) gives one internal lift per
boundary word.  For each \((a_B)_j=1\), attach a retained ordinary
\(I\)-vertex to the corresponding circuit port through one native
\(X\)-edge and expose the other ordinary port of that vertex.  Since
\(XI=X\), this attachment complements exactly that boundary bit.
Leave the other circuit ports unchanged.  All exposed ports in this
construction are ordinary, so no free-endpoint boundary action cancels
the required flip.  Applying these dressings independently
realizes every affine coset \(a_B+L\), rather than only the zero coset.
\end{proof}

An orthogonal map \(A\) has a self-dual graph relation
\(\{(x,Ax):x\in\Ftwo^n\}\).  The converse of
\cref{lem:rm-flat-contexts} realizes that relation as a network.
Attaching a signature to its input and closing one output pair then sums
the signature over prescribed two-point orbits.  This supplies the
following Radon cards for the high-arity argument.

\paragraph{Actual Radon access.}

We use the symmetric-binary-space terminology and ambient-relative
orthogonal notation of
\cref{subsubsec:symmetric-binary-spaces}.

Let \(g\) have even arity \(n\).  Call
\(v\in(\one^n)^\perp\setminus\{0,\one^n\}\) \emph{accessible}.  Such a
\(v\) is isotropic.  For \(\epsilon\in\Ftwo\), define
\[
 H_{v,\epsilon}:=\{x\in\Ftwo^n:x\cdot v=\epsilon\}.
\]
Translation by \(v\) preserves \(H_{v,\epsilon}\).  For
\(x\in H_{v,\epsilon}\), write
\[
 [x]_v:=x+\Span_{\Ftwo}\{v\}
\]
for its translation orbit, and let
\(H_{v,\epsilon}/\Span_{\Ftwo}\{v\}\) denote the set of these orbits.  Define
\[
 [u]_v\cdot[w]_v:=u\cdot w
 \qquad(u,w\in v^\perp).
\]
This is a well-defined nondegenerate symmetric form on the translation
space \(v^\perp/\Span_{\Ftwo}\{v\}\); each
\(H_{v,\epsilon}/\Span_{\Ftwo}\{v\}\) is an affine space over that quotient.
Define
\begin{equation}
  \mathcal R_{v,\epsilon}g([x]_v)
  =g(x)+g(x+v),
  \qquad [x]_v\in H_{v,\epsilon}/\Span_{\Ftwo}\{v\}.
  \label{eq:rm-radon-card}
\end{equation}
The right side is invariant under replacing \(x\) by \(x+v\), so the
quotient signature is well defined.  The construction below identifies
the quotient with \(n-2\) ordinary boundary ports.  Thus the quotient
sum is an actual arity-reducing operation in the original network model.
Write
\[
 O(n,2):=\{A\in\operatorname{GL}_n(\Ftwo):A^{\mathsf T}A=I_n\}.
\]

\begin{lemma}[Accessible Radon directions]
\label{lem:rm-radon-accessible}
Let \(g\) be a retained signature of even arity \(n\ge4\).  For each
\(v\in(\one^n)^\perp\setminus\{0,\one^n\}\), one can construct and
record distinct ports \(i_v,j_v\) and an orthogonal matrix
\(A_v\in O(n,2)\) such that
\[
 A_v^{\mathsf T}(e_{i_v}+e_{j_v})=v.
\]
Define
\[
 r_v:=A_v^{\mathsf T}e_{i_v},\qquad
 \Pi_v:=\Span\{r_v,v\},\qquad W_v:=\Pi_v^\perp,
\]
and, for \(\epsilon\in\Ftwo\),
\[
 \iota_{v,\epsilon}:W_v\longrightarrow
 H_{v,\epsilon}/\Span\{v\},\qquad
 w\longmapsto[\epsilon r_v+w]_v.
\]
Then \(\Ftwo^n=\Pi_v\perp W_v\), and \(\iota_{v,\epsilon}\) is
an affine isometry from \(W_v\) to
\(H_{v,\epsilon}/\Span\{v\}\).  An \(O(n^2)\)-size network using
one copy of \(g\), \(\RMcore\), and \(I/X\) realizes
\(\mathcal R_{v,\epsilon}g\) in the boundary coordinates supplied by
\(A_v|_{W_v}\), up to a recorded nonzero scalar.  The available
directions are exactly the displayed vectors \(v\).
\end{lemma}

\begin{proof}
First take \(A_v,i_v,j_v\) as in the statement; their construction is
given at the end of the proof.  The definitions imply
\(\Pi_v=A_v^{\mathsf T}\Span\{e_{i_v},e_{j_v}\}\).
Then
\[
 r_v\cdot r_v=1,\qquad
 r_v\cdot v=1,\qquad
 v\cdot v=0.
\]
Hence \(\Pi_v\) is a nondegenerate two-plane and
\[
 \Ftwo^n=\Pi_v\perp W_v,
 \qquad
 v^\perp=\Span_{\Ftwo}\{v\}\oplus W_v.
\]
The isometry \(A_v\) carries \(\Pi_v\) to the physical output-coordinate
plane \(\Span\{e_{i_v},e_{j_v}\}\) and \(W_v\) to its ordinary-coordinate
complement.  We use \(A_v\) to identify these pairs of spaces and call the
corresponding output ports the \(\Pi_v\)-ports and the \(W_v\)-ports.
Consequently \(\iota_{v,\epsilon}\) has the asserted domain, range,
and isometry property.

For \(A\in O(n,2)\), define
\[
 \operatorname{Graph}(A):=\{(x,Ax):x\in\Ftwo^n\}.
\]
Its direction is self-dual in \(\Ftwo^n\oplus\Ftwo^n\): it is isotropic
because
\(x\cdot y+(Ax)\cdot(Ay)=0\), and it has half the ambient dimension.
Hence \cref{lem:rm-flat-contexts} realizes a nonzero scalar multiple of
\[
 (x,y)\longmapsto\one[y=A_vx]
\]
by an actual \(\RMcore/I/X\)-network.  Attach one copy of \(g\) to the
graph input using an effective \(I\)-kernel on each input connection:
insert one retained ordinary \(I\)-vertex between two native \(X\)-edges,
so that \(XIX=I\).  The output before closing a pair is therefore
\(g(A_v^{\mathsf T}y)\), with the graph's recorded nonzero scalar.
Close the ordered output pair \((i_v,j_v)\) with either actual kernel
\(I\) or actual kernel \(X\).  The two internal lifts differ by
\(v\); the choice of \(I\) or \(X\) fixes \(x\cdot v=\epsilon\).  Under the
isometry \(A_v|_{W_v}\), the remaining output coordinates are precisely
the \(W_v\)-coordinates in the inverse of \(\iota_{v,\epsilon}\).  Thus the
resulting boundary is \(\mathcal R_{v,\epsilon}g\), up to its recorded
nonzero scalar.

The weight-four transvections constructed in
\cref{lem:rm-flat-contexts}, together with coordinate permutations, send
every nonzero even \(v\ne\one^n\) to a weight-two vector.  Write the
resulting orthogonal map as \(B\), with \(Bv=e_{i_v}+e_{j_v}\), and
record \(A_v=B\).  Then
\(A_v^{\mathsf T}(e_{i_v}+e_{j_v})=B^{-1}Bv=v\), as required.  Conversely, orthogonal maps preserve parity and fix
\(\one^n\), so neither \(0\) nor \(\one^n\) lies in the orbit of a
weight-two vector.  This proves necessity as well.  The transvection
construction and graph network use \(O(n^2)\) vertices; the associated
binary linear systems are solved by Gaussian elimination.
\end{proof}

\subsection{Relative flat-Lagrangian closure and completion of
\texorpdfstring{\cref{thm:reed-muller-rigidity}}{Reed--Muller deck rigidity}}

We now consider an arbitrary unresolved tensor-prime occurrence \(g\).
For arity at most eight, its smaller flat cards are matching products, so
the reconstruction above applies.  At higher arity, Radon cards compare
two quotient layers.  A failure to combine them into an affine signature
can be retained on at most eight ports; after that obstruction is excluded,
affine patching and dimension counts recover a flat-Lagrangian \(g\).

For an \(n\)-ary signature \(h\) and a shift vector \(s\in\Ftwo^n\), define
the translation operator \(T_s\) by \((T_sh)(x):=h(x+s)\).  Write
\(\operatorname{Id}\) for the identity operator on the current signature
space and define the orbit-sum operator
\[
 \mathsf P_s^+:=\operatorname{Id}+T_s.
\]
Thus
\((\mathsf P_s^+h)(x)=h(x)+h(x+s)\) is the sum over the two-point
translation orbit \(\{x,x+s\}\).
For accessible \(s\), the two quotient signatures
\(\mathcal R_{s,0}h\) and \(\mathcal R_{s,1}h\) are called the
\(0\)- and \(1\)-layer \emph{plus cards} of \(h\) in direction \(s\);
together they are the two quotient layers of \(\mathsf P_s^+h\).

\begin{lemma}[Low-arity flat signatures are matching products]
\label{lem:rm-low-flat-matching}
Let \(L=L^\perp\le\Ftwo^{2r}\), where \(r\le3\).  After a port
permutation, \(L\) is a direct sum of \(r\) two-coordinate repetition
lines \(\Span_{\Ftwo}\{e_{2j-1}+e_{2j}\}\).  Consequently
every flat-Lagrangian signature of arity at most six is a product of
projective \(I/X\) binaries on a perfect matching.
\end{lemma}

\begin{proof}
Every vector of \(L\) has even weight, so \(\one\in L^\perp=L\).
If \(r>1\), choose \(z\in L\setminus\{0,\one\}\).  One of \(z,z+\one\)
has weight at most \(r\), hence weight two when \(2r\le6\).  Thus
\(e_i+e_j\in L\) for some pair \(i,j\).  Self-duality now gives
\(x_i=x_j\) for every \(x\in L\), so this repetition line splits
orthogonally from \(L\).  Induction proves the assertion.  Translating the
affine support changes an equality factor only to its disequality coset.
\end{proof}

Consequently every flat-Lagrangian proper card of an occurrence of arity at
most eight is already a matching product.  A nonmatching flat-Lagrangian
proper card can first occur at arity ten; its appearance is not by itself a
protocol exit.  This is precisely the remaining case for the Radon
localization argument.

\begin{lemma}[Two flat quotient charts localize nonaffine gluing]
\label{lem:rm-two-chart-localization}
Let \(g\) have even arity \(n\ge10\), let \(v\) be accessible, and suppose
that, for both \(\epsilon\in\Ftwo\), the card
\(\mathcal R_{v,\epsilon}g\) is zero or flat-Lagrangian.  Define, only for
the analysis, the orbit-sum signature
\begin{equation*}
  F_v(x)=g(x)+g(x+v).
\end{equation*}
If \(F_v\notin\cA\), an \(O(n^2)\)-size direct
\(\RMcore/I/X\)-context using one copy of \(g\) produces a nonzero
\(H\notin\cA\) of arity at most eight.
\end{lemma}

\begin{proof}
Fix the recorded data
\((A_v,i_v,j_v,r_v,\Pi_v,W_v)\) from
\cref{lem:rm-radon-accessible}, and write \(W:=W_v\).  Before the final
output pair is closed, the graph realization gives the orthogonal
decomposition
\[
 \Ftwo^n=\Pi_v\perp W.
\]
By the recorded identification, the two physical output ports
\(i_v,j_v\) are the \(\Pi_v\)-ports, while the remaining output ports are
the \(W\)-ports.

Under the isometry \(W\simeq v^\perp/\Span_{\Ftwo}\{v\}\), the
characteristic vector of \(W\) corresponds to the class of \(\one^n\).
This class is nonzero because accessibility excludes \(v=\one^n\).
Thus \(W\) is even-dimensional and nonalternating.

For \(\epsilon\in\Ftwo\), define the analytic quotient chart
\[
 h_\epsilon(w):=F_v(\epsilon r_v+w),
 \qquad w\in W.
\]
This is well defined because \(F_v\) is invariant under translation by
\(v\).  Each nonzero chart has the form
\[
 h_\epsilon(w)
 =c_\epsilon\one[w\in a_\epsilon+L_\epsilon],
 \qquad
 c_\epsilon\ne0,\quad a_\epsilon\in W,\quad
 L_\epsilon=L_\epsilon^{\perp_W}\le W.
\]
If \(h_{1-\epsilon}=0\) and \(h_\epsilon\ne0\), then, in the decomposition
\(\Ftwo^n=\Pi_v\perp W\),
\[
 F_v(x)
 =c_\epsilon\one[
   x\in\epsilon r_v+a_\epsilon+
        (\Span_{\Ftwo}\{v\}\oplus L_\epsilon)].
\]
Thus \(F_v\) is affine whenever one chart is zero, and the assumed
nonaffinity ensures that both displayed charts are nonzero.
Here ``the charts glue affinely'' means that their union, with the layer
coordinate restored, is the support coset of one affine signature.  This
happens exactly when \(L_0=L_1\); once their supports glue, their constant
phases glue exactly when \(c_1/c_0\in\mu_4\).

If \(L_0\ne L_1\), take \(s\in L_0\setminus L_1\).  Since
\(s\notin L_1^{\perp_W}=L_1\), choose \(t\in L_1\) with
\(s\cdot t=1\).  Elementary completion in a
nondegenerate binary symmetric space gives an even-dimensional,
nonalternating, nondegenerate subspace
\[
 U\supseteq\Span\{s,t,a_1-a_0\},
 \qquad \dim U\le6.
\]
To obtain this dimension bound, use the symmetric-space conventions in
\cref{subsubsec:symmetric-binary-spaces}.
Put \(S:=\Span\{s,t,a_1-a_0\}\).  Since \(s\cdot t=1\), the plane
\(\Span\{s,t\}\) is nondegenerate, so
\(\dim\operatorname{rad}(S)\le1\).  If
\(\operatorname{rad}(S)=0\), put \(N:=S\).  Otherwise write
\(\operatorname{rad}(S)=\Span_{\Ftwo}\{r\}\), choose a nondegenerate
complement \(S_0\) of \(\Span_{\Ftwo}\{r\}\) in \(S\), and choose
\(d\in W\) satisfying
\[
 r\cdot d=1,
 \qquad
 d\in S_0^\perp.
\]
Put \(N:=S+\Span_{\Ftwo}\{d\}\).  In either case \(N\) is nondegenerate
and \(\dim N\le4\).
If \(\dim N\) is odd, its nondegenerate orthogonal complement has odd
dimension and therefore contains an anisotropic vector; adjoining that
orthogonal line makes \(N\) even-dimensional without exceeding dimension
four.  If the resulting even-dimensional \(N\) is alternating, then
\(N^\perp\) is nonalternating because \(W\) is; adjoin a nonalternating
nondegenerate plane from \(N^\perp\).  The resulting \(U\) has all the
asserted properties and dimension at most six.

If instead \(L_0=L_1\) and \(c_1/c_0\notin\mu_4\), every vector in an
even-dimensional nondegenerate nonalternating symmetric binary space lies in a
nonalternating nondegenerate plane.  For a nonzero isotropic vector \(z\),
choose \(d\) with \(z\cdot d=1\) and \(d\cdot d=1\); for an anisotropic
\(z\), choose \(d\) with \(z\cdot d=1\) and \(d\cdot d=0\).  For \(z=0\),
choose any nonalternating nondegenerate plane.  We may therefore
choose an even-dimensional, nonalternating, nondegenerate
\[
 U\supseteq\Span\{a_1-a_0\},
 \qquad \dim U\le2.
\]

Put
\[
 P:=U^{\perp_W},
 \qquad
 W=U\perp P,
\]
and let \(\pi_U,\pi_P\) be the corresponding orthogonal projections.  The
space \(P\) is even-dimensional and nondegenerate.  By the
symmetric-binary Witt decomposition recalled in
\cref{subsubsec:symmetric-binary-spaces}, \(P\) contains a
half-dimensional totally isotropic subspace \(J\); equivalently,
\[
 J=J^{\perp_P}\le P.
\]
Because \(a_1-a_0\in U\), the two shifts have the same \(P\)-projection;
call it \(p\).  Choose an isometry
\(\phi:U\to\Ftwo^{\dim U}\), and define the affine relation
\[
 \Gamma_p
 :=\{(u+p+j,\phi(u)):u\in U,\ j\in J\}
 \subseteq W\oplus\Ftwo^{\dim U}.
\]
Its direction is self-dual: it has half the ambient dimension, and
\[
 (u+j)\cdot(u'+j')+\phi(u)\cdot\phi(u')
 =j\cdot j'=0
\]
for \(u,u'\in U\) and \(j,j'\in J\).  Thus
after transporting the \(W\)-coordinate through \(A_v|_{W_v}\),
\cref{lem:rm-flat-contexts} realizes the corresponding signature
\[
 (w,z)\longmapsto\one[(w,z)\in\Gamma_p]
\]
by an actual \(O(n^2)\)-size \(\RMcore/I/X\)-network.  Apply this relation
to the \(W\)-ports and leave the two \(\Pi_v\)-ports external.

In the different-direction case, the boundary points induced by
\(a_0,a_0+s,a_1\) survive but the fourth parallelogram point induced by
\(a_1+s\) does not.  Otherwise there would be \(z\in L_1\) with
\(\pi_Uz=s\), and
\[
  0=z\cdot t=(\pi_Uz)\cdot t=s\cdot t=1.
\]
The output therefore retains a nonaffine parallelogram on at most
\(\dim U+2\le8\) ports.  In the equal-direction case, the two fibers have
equal size because their support directions agree.  Thus neither
cancellation nor projection changes the ratio \(c_1/c_0\) outside
\(\mu_4\), and the output is again nonaffine.

Let \(H\) be the contraction of this actual context with \(g\), and let
\(H^+\) be the contraction with \(F_v=\mathsf P_v^+g\).  Order the boundary
with the two retained \(\Pi_v\)-ports first and put
\[
 w_\partial:=(1,1,\mathbf 0_{\dim U}).
\]
Indeed, translation by \(v\) changes the graph output by
\(e_{i_v}+e_{j_v}\), so it flips exactly the two \(\Pi_v\)-ports and acts
trivially on the remaining boundary ports.  Consequently, with the same
recorded scalar on both terms,
\[
 H^+=\mathsf P_{w_\partial}^+H.
\]
The missing parallelogram or forbidden layer ratio proves
\(H^+\notin\cA\).  By \cref{lem:affine-two-translate-closure},
\(H\in\cA\) would imply \(H^+=\mathsf P_{w_\partial}^+H\in\cA\), a
contradiction.  Therefore \(H\ne0\), \(H\notin\cA\), and its arity is at
most eight.
\end{proof}

In the relative induction, the localized nonaffine signature supplied by the
preceding lemma is excluded by the already closed lower-arity cases.  Hence
every accessible orbit sum \(F_v\) is affine.  The next lemma patches these
hyperplane-wise affine conclusions into a global affine conclusion for
\(g\).

\begin{lemma}[Affine hyperplane patching]
\label{lem:rm-affine-hyperplane-patching}
Let \(n\ge10\) be even, and let \(g:\Ftwo^n\to\AlgNums\).  For every
accessible \(v\), define the orbit-sum signature
\(F_v:=\mathsf P_v^+g\).  If \(F_v\in\cA\) for every such \(v\),
then \(g\in\cA\).
\end{lemma}

\begin{proof}
We first record the Walsh-invariance of \(\cA\) used below.  The Walsh
kernel \((-1)^{\xi\cdot x}\) has an allowed quadratic modulo-four phase,
so its product with an affine signature is affine in the joint variables
\((\xi,x)\).  Gaussian-eliminate the affine support constraints and sum
the remaining internal variables one at a time.  If the current phase
depends on a free bit \(y\) as
\[
 i^{\alpha y}(-1)^{y\ell},
 \qquad \alpha\in\mathbb Z_4,
\]
where \(\ell\) is an affine linear form in the remaining variables and is
independent of \(y\),
then summing over \(y\) either imposes the affine equation
\(\ell=\alpha/2\pmod2\) when \(\alpha\) is even, or contributes a nonzero
scalar times a fourth-root phase of the affine linear form \(\ell\) when
\(\alpha\) is odd.  Thus every elimination remains in \(\cA\), proving
that the Walsh transform preserves \(\cA\).  The converse follows from
Walsh inversion.

Let \(\mathcal F_2 g\) denote the (unnormalized) Walsh transform
\[
  (\mathcal F_2g)(\xi)=\sum_{x\in\Ftwo^n}(-1)^{\xi\cdot x}g(x),
\]
and write \(g^{\mathrm W}:=\mathcal F_2g\).  This superscript distinguishes
the Walsh transform here from the manuscript's global hat notation for
\(K\)-coordinates.  Since
\begin{equation*}
  (\mathcal F_2F_v)(\xi)
  =(1+(-1)^{\xi\cdot v})g^{\mathrm W}(\xi),
\end{equation*}
each restriction \(g^{\mathrm W}|_{v^\perp}\) is affine or zero.  The claim
is immediate when \(g^{\mathrm W}=0\), so assume otherwise and choose
\(\zeta\in\supp g^{\mathrm W}\).

Given \(x,y\in\supp g^{\mathrm W}\), the space
\[
  E\cap\Span\{\zeta,x,y\}^\perp,
  \qquad E=\one^\perp,
\]
has dimension at least \(n-4\), hence contains accessible \(v\).  Affinity
on \(v^\perp\) puts \(x+y+\zeta\) in the support, so
\(\supp g^{\mathrm W}\) is a
coset; write it as \(\zeta+U\).

Choose an isomorphism \(\beta:\Ftwo^{\dim U}\to U\) and define the
nowhere-zero coefficient function
\[
 \psi(u):=g^{\mathrm W}(\zeta+\beta u).
\]
Fix arbitrary
\[
 u_0,u_1,u_2,u_3\in\Ftwo^{\dim U},
\]
where \(u_0\) is the base point and \(u_1,u_2,u_3\) are the three
directions of a multiplicative third-difference cube of \(\psi\).  The four
vectors
\[
 \zeta+\beta u_0,\quad
 \beta u_1,\quad
 \beta u_2,\quad
 \beta u_3
\]
span at most four dimensions.  Hence
\[
 E\cap
 \Span\{\zeta+\beta u_0,\beta u_1,\beta u_2,\beta u_3\}^\perp
\]
has dimension at least \(n-5\ge5\) and contains an accessible direction
\(s\).  The entire cube lies in \(s^\perp\), where the nonzero restriction
of \(g^{\mathrm W}\) is affine; hence
\[
 \mathcal T_\psi(u_0,u_1,u_2,u_3)=1.
\]
By \cref{lem:affine-phase-criterion}, \(\psi\) has an allowed quadratic
modulo-four phase.  Consequently \(g^{\mathrm W}\in\cA\), and the
Walsh-invariance proved above gives \(g\in\cA\).
\end{proof}

Affine patching gives \(g\in\cA\), but does not yet give the required
flat-Lagrangian form.  The final recognition lemma uses the dimensions and
phases of all actual Radon cards to force a self-dual support direction and
constant coefficient on that support.

\begin{lemma}[Affine-to-flat Radon recognition]
\label{lem:rm-affine-flat-radon}
Let \(0\ne h\in\cA\) have even arity \(2r\ge8\).  Suppose that for every
accessible direction \(v\) and both \(\epsilon\in\Ftwo\), the actual card
\(\mathcal R_{v,\epsilon}h\) is zero or flat-Lagrangian.  Then \(h\) is
flat-Lagrangian.
\end{lemma}

\begin{proof}
Write the affine signature \(h\) in standard normal form as
\[
  h(x)=c\one[x\in a+U]i^{Q(x)},
  \qquad d:=\dim U.
\]
Here \(c\ne0\) is the scale, \(a+U\) is the affine support, \(Q\) is its
quadratic phase polynomial, and \(d\) is the support dimension.  Define the even hyperplane
\(E:=\one^\perp\).  For \(v\notin U\) there is no cancellation:
the nonempty card has support dimension \(d\) if \(v\in U^\perp\), and
both have dimension \(d-1\) otherwise.  A flat card has dimension \(r-1\).

Neither \(E\subseteq U\) nor \(E\subseteq U^\perp\).  Suppose first that
\(E\subseteq U\).  Choose \(p,q\in E\) so that
\(p,q,p+q\) are accessible.  They exist because \(\dim E=2r-1\ge7\).
If
\[
 \mathsf P_p^+h=\mathsf P_q^+h=0,
\]
then \(T_ph=T_qh=-h\), so \(T_{p+q}h=h\) and
\(\mathsf P_{p+q}^+h=2h\ne0\).  Thus at least one layer plus card arising
from the three directions \(p,q,p+q\) is nonzero.  Let
\(s\in\{p,q,p+q\}\) be its direction.  Passing from \(h\) to such a card
can lower support dimension by at most one when restricting to
its layer hyperplane and by one when quotienting by
\(\Span_{\Ftwo}\{s\}\).
On an overlapping support orbit, factoring out \(h(x)\) leaves
\[
 1+\zeta(-1)^{\ell(x)},
 \qquad \zeta\in\mu_4.
\]
Here \(\ell\) is a binary affine-linear form on the current support.
For \(\zeta\in\{\pm i\}\) this factor is nowhere zero; for
\(\zeta\in\{\pm1\}\) its nonzero set is empty, the whole current support,
or one affine hyperplane.  Since the selected card is nonzero, cancellation
therefore lowers support dimension by at most one more.
Hence its support dimension is at least
\[
 d-3\ge2r-4>r-1,
\]
contradicting flatness.

If instead \(E\subseteq U^\perp\), then \(d\le1\).  Choose accessible
\(v\in E\setminus U\).  Here \(v\in U^\perp\setminus U\), so there is no
cancellation and the nonzero card has support dimension \(d\), whereas a
flat-Lagrangian card has dimension \(r-1\ge3\), again a contradiction.

Thus both \(E\cap U\) and \(E\cap U^\perp\) are proper subspaces of \(E\).
Two proper subspaces cannot cover \(E\); since \(\dim E\ge7\), we may also
avoid \(0,\one\) and choose accessible
\(v\notin U\cup U^\perp\).  Its two cards have dimension
\(d-1=r-1\), so \(d=r\).

No accessible vector lies in \(U^\perp\setminus U\), since it would give a
dimension-\(r\) card.  Put \(A_0:=E\cap U^\perp\).  Then
\[
 A_0\setminus\{0,\one\}\subseteq U.
\]
If \(\one\notin A_0\), or if \(\one\in A_0\cap U\), this gives
\(A_0\subseteq U\).  If instead \(\one\in A_0\setminus U\), choose
\(z\in A_0\setminus\{0,\one\}\), which is possible because
\(\dim A_0\ge r-1\ge3\); then both \(z\) and \(z+\one\) are
accessible and lie in \(U\) by the preceding display, forcing
\(\one\in U\), a contradiction.  Hence \(A_0\subseteq U\) in every case.
If
\(U^\perp\subseteq E\), dimensions give
\(U=U^\perp\).  Otherwise \(\dim A_0=r-1\).  Since
\(A_0\subseteq U\) and \(\dim U=r\), choose \(u\in U\setminus A_0\).
Then \(U=A_0\oplus\Span_{\Ftwo}\{u\}\), and
\(A_0\subseteq U^\perp\) makes \(u\) orthogonal to \(A_0\).  If \(u\) were
even, then \(u\cdot u=0\), so \(u\in U^\perp\) and
\(U\subseteq U^\perp\); equal dimensions would give \(U=U^\perp\),
contrary to the present case.  Thus \(u\) is odd and
\begin{equation}
 U=A_0\oplus\Span_{\Ftwo}\{u\},
  \qquad u\cdot u=1.
  \label{eq:rm-near-self-dual}
\end{equation}

For accessible \(s\in A_0\), put
\[
 \epsilon_s:=a\cdot s.
\]
Because \(s\in U^\perp\), the support \(a+U\) lies wholly in
\(H_{s,\epsilon_s}\); the other layer plus card is identically zero.  The
direction of the supported card
\(\mathcal R_{s,\epsilon_s}h\) is
\(U/\Span_{\Ftwo}\{s\}\), whose orthogonal is
\(U^\perp/\Span_{\Ftwo}\{s\}\); these differ because \(U\ne U^\perp\).
Write its phase difference as
\[
 Q(x+s)-Q(x)=\alpha_s+2\ell_s(x)\pmod4,
 \qquad \alpha_s\in\mathbb Z_4,
\]
where \(\ell_s\) is affine linear on \(a+U\).  Thus
\[
 \frac{h(x+s)}{h(x)}
 =i^{\alpha_s}(-1)^{\ell_s(x)}.
\]
If \(\alpha_s\) is odd, the supported plus card is nonzero on the full
quotient and has the non-self-dual direction
\(U/\Span_{\Ftwo}\{s\}\).
If \(\alpha_s\) is even and \(\ell_s\) is nonconstant, its nonzero support
is a codimension-one hyperplane of dimension \(r-2\).  If \(\ell_s\) is
constant, the card is either zero or again has the full non-self-dual
direction.  The hypothesis therefore forces
\[
 \mathcal R_{s,\epsilon_s}h=0,
 \qquad\text{equivalently}\qquad
 \mathsf P_s^+h=0.
\]

Because \(\dim A_0=r-1\ge3\), choose a two-dimensional subspace
\(B_0\le A_0\) with \(\one\notin B_0\), and let \(p,q\) be a basis of
\(B_0\).  Then \(p,q,p+q\) are all accessible, and
\(\mathsf P_p^+h=\mathsf P_q^+h=0\), whereas
\(T_{p+q}=T_pT_q\) gives
\(\mathsf P_{p+q}^+h=2h\ne0\), a contradiction.
Therefore \(U=U^\perp\).

For accessible \(s\in U=U^\perp\), put
\(\epsilon_s:=a\cdot s\).  The support lies wholly in
\(H_{s,\epsilon_s}\), so \(\mathcal R_{s,\epsilon_s}h\) is the unique
potentially nonzero layer plus card.  Use the same phase-difference form
as above.
If \(\alpha_s\) is even and \(\ell_s\) is nonconstant, the nonzero part of
the plus card has quotient-support dimension \(r-2\), so it cannot be
flat-Lagrangian.  In every remaining case the plus card is either zero or
has the full self-dual quotient support \(U/\Span_{\Ftwo}\{s\}\); by the
flat-Lagrangian hypothesis it is then constant on that support.  Taking
\(c_s=0\) in the zero case, we obtain
\begin{equation*}
  h(x)+h(x+s)=c_s
  \qquad(x\in a+U).
\end{equation*}
Since \(\dim U=r\ge4\) and \(U=U^\perp\), choose a two-dimensional
subspace \(B_1\le U\) with \(\one\notin B_1\), and let \(p,q\) be a basis.
Then \(p,q,p+q\) are all accessible.  Using \(T_{p+q}=T_pT_q\) to compare
their three equations gives
\[
  2h(x)=c_{p+q}-c_q+c_p,
\]
so \(h\) is constant on \(a+U\), hence flat-Lagrangian.
\end{proof}

\begin{proof}[Completion of \cref{thm:reed-muller-rigidity}]
We first close the initial minimum core \(f\).  If it has odd arity, apply
\cref{thm:external-odd} to \(K\cG\).  Its hardness conclusion pulls back
through the holographic \(K\)-coordinate equivalence, while its tractable
conclusion is exactly \(\Tract(K\cG)\).  If
\(\arity(f)=4\), \cref{lem:proper-cyclic-q4} and the quaternary interface
give either a resolved leaf or a strict stable successor, and the complete
batch is resolved by
\cref{thm:equality-q4-boundary,cor:p1-eight-vertex-reduced-interface,%
lem:actual-binary-completeness,def:relative-stratum-protocol,%
lem:matching-synthesis-measure}.  Hence, on the sole remaining initial
branch, \(f\) has even arity \(n\ge6\), and every proper card is zero or a
nonsingular current-group matching signature.

Apply \cref{lem:rm-zero-graph,lem:rm-matching-consistency}.  If
\(W_f\ne\varnothing\), then \cref{lem:rm-weight-two-integration} gives
either the actual \(\mathrm{Even}_4\) equality anchor constructed in its
proof or a proper tensor factorization.  After factor saturation and group
recomputation, both alternatives are resolved by
\cref{lem:equality-anchor,lem:factor-saturation,%
def:relative-stratum-protocol,lem:matching-synthesis-measure}.
Thus the remaining branch has \(W_f=\varnothing\).  Now
\cref{lem:rm-support-classification,lem:rm-coefficient-atlas} give
\[
 n=8,
 \qquad
 f(x)=c\one[x\in a+RM(1,3)]
\]
after a port permutation, for a recorded \(c\ne0\).  Available actual
one-port \(X\)-dressings remove \(a\), so the retained set contains the
literal standard core \(\RMcore\).  Regenerate the complete deck of every
safe one-port dressing of this core.  If one leaves the matching branch,
its complete batch is resolved by the same case analysis;
otherwise \cref{lem:rm-safe-group-collapse} gives
\[
 G=\{[I],[X]\}\cong C_2.
\]
This proves the asserted description of the sole initial rigidity branch.

Retain this literal \(\RMcore\) and normalized \(I/X\) throughout the
remaining signature-set-wide closure.  At every continuing augmentation,
recompute the complete group and process the Reed--Muller safe-dressing
deck before re-entering the argument.  A decrease of \(\nu\) or a strict
group enlargement is charged to the stable-state measure \(\Omega\) of
\cref{lem:matching-synthesis-measure}; on the unchanged remaining branch
the preceding safe-dressing argument reapplies and the complete group
remains \(C_2\).

We next make the outer induction explicit.  Let \(\mathcal W\) be the
finite multiset of unresolved tensor-prime occurrences in the
retained factorization trees.  When a token \(u\), carrying \(h_u\), first
enters the worklist, freeze the following finite Reed--Muller reduction
universe \(\mathscr U_{\rm RM}(u)\):
\begin{enumerate}
\item the complete physical \(C_2\) matching-deck universe of \(h_u\), with
      every oriented proper contraction and safe \(I/X\)-dressing, together
      with every ordered two-copy \(I/X\)-integration network displayed in
      \cref{lem:rm-weight-two-integration}, for all port, orientation, and
      link-sign choices;
\item for every accessible \(v\) and \(\epsilon\), the fixed actual Radon
      network of \cref{lem:rm-radon-accessible}, including its recorded
      data \((A_v,i_v,j_v,r_v,\Pi_v,W_v)\), and, for every finite
      parameter tuple \((v,U,J,p,\phi)\) allowed in
      \cref{lem:rm-two-chart-localization}, its actual localization network
      applied to \(h_u\) using that same recorded Radon data;
\item the complete factorization tree of every output in the preceding two
      items.
\end{enumerate}
Each network-output key records the ordered physical network, its boundary
order, actual scalar, and direct provenance; each attached factor token
records its canonical signature and Turing-exposure provenance.  The
universe is finite because the port set and
all of its binary subspaces, vectors, and isometries are finite; every
individual listed network has the size bounds proved in the cited lemmas.
These choices depend on the fixed retained signature set, not on the
input Holant graph.  Once selected, each network replaces a vertex by a
fixed-size gadget, as required for the stated reductions.
Partition it as
\[
 \mathscr U_{\rm RM}(u)=Q_{u,\mathrm{done}}\sqcup Q_u.
\]
For \(q\in Q_u\), let \(\mathcal W_q\) be the multiset of open occurrences in its attached factorization tree, and set
\[
 \Phi_u:=\sum_{q\in Q_u}\Phi(\mathcal W_q),
\]
using \cref{eq:factor-potential}.  Define
\[
 \Theta_{\rm RM}(\mathcal W)
 =
 \{\!\{(\arity(h_u),|Q_u|,\Phi_u):u\in\mathcal W\}\!\},
\]
using lexicographic order on each triple and the well-founded multiset
extension of that order; explicitly, replacing one triple by any finite
multiset of strictly smaller triples is a strict decrease.

Certifying an occurrence removes its triple.  If the selected occurrence
factors, its triple is replaced by triples of strictly smaller arity.  Every
batch used below starts from a frozen key: the matching-deck and integration
steps use item~1, the Radon steps use the first family in item~2, and the
two-chart step uses the second family in item~2.  Completing such a key moves
it to \(Q_{u,\mathrm{done}}\), while an intermediate genuine factor split
strictly decreases its summand \(\Phi(\mathcal W_q)\).  Any newly opened
child occurrence has smaller arity: proper and Radon cards have arity at
most \(\arity(h_u)-2\), localization outputs have arity at most eight when
they are invoked from arity at least ten, and every proper factor is smaller
than its parent.  Thus the old parent triple is replaced by a strictly
smaller parent triple together with child triples whose first coordinate is
smaller.  This strictly decreases \(\Theta_{\rm RM}\).  No fixed-\((H,\nu)\)
batch creates a fresh unresolved key, because all of its possible outputs
and complete factorization trees were included in \(\mathscr U_{\rm RM}(u)\) at
the freeze.  Consequently this outer induction, rather than a multiset of
arities alone, also controls passage to another occurrence of the same
arity.

We use the following relative re-entry observation.  Suppose a
tensor-prime occurrence \(g\) has had its complete physical deck regenerated
and every proper card and proper factor of smaller arity has already been
closed.  On the continuing branch, assume every surviving proper card is
zero or a current-\(C_2\) matching signature.  Then the proofs of
\cref{lem:rm-zero-graph,lem:rm-matching-consistency,%
lem:rm-weight-two-integration,lem:rm-support-classification,%
lem:rm-coefficient-atlas}
apply verbatim with \(g\) in place of the initial minimum core.  Indeed,
minimum arity is used there only to exclude a surviving proper nonmatching
card or nonbinary factor; regenerated-deck and prior-closure supply exactly
that exclusion.  All displayed local networks and coefficient identities
are unchanged.

Run induction on \(\Theta_{\rm RM}\).  Before processing an occurrence,
recursively close all its lower-arity proper cards and factors.  Unary,
odd-arity, and quaternary outputs are dispatched by
\cref{thm:external-odd,thm:equality-q4-boundary}.  Hence consider a
relatively re-entered nonflat tensor-prime \(g\) of even arity at least six.

There are now two possible continuations.  First suppose that every surviving
proper card of \(g\) is zero or a current-\(C_2\) matching signature.  The
relative re-entry observation and
\cref{lem:rm-zero-graph,lem:rm-matching-consistency,%
lem:rm-weight-two-integration,lem:rm-support-classification,%
lem:rm-coefficient-atlas}
then give a resolved \(\mathrm{Even}_4\) equality anchor, a
proper factorization, or the flat eight-ary \(RM(1,3)\) core.  Thus no
nonflat occurrence with an all-matching proper deck survives.

Otherwise some nonzero proper card is nonmatching.  Factor saturation and
prior closure nevertheless make every proper card zero or a tensor product
of already certified flat-Lagrangian factors, hence zero or
flat-Lagrangian.  If \(6\le\arity(g)\le8\), every such card has arity at most
six, and \cref{lem:rm-low-flat-matching} makes it an \(I/X\)-matching
signature after a port permutation, contradicting the present case.  The
appearance of a nonmatching flat-Lagrangian card is not itself a protocol
exit; it identifies the remaining Radon branch, where
\(\arity(g)=n\ge10\).

By \cref{lem:rm-radon-accessible}, every
Radon card is actual, is indexed by its frozen item~2 key, and has arity
\(n-2\).  Processing that key includes its complete factor saturation.
Every surviving factor has already been certified
flat-Lagrangian; tensor closure therefore makes each card zero or
flat-Lagrangian.  Fix an accessible \(v\).  If the analytic orbit sum
 \[
 F_v=\mathsf P_v^+g
 \]
were nonaffine, \cref{lem:rm-two-chart-localization} would produce an
actual nonaffine signature \(H\) of arity at most eight.  Its parameter
tuple and actual context have a frozen item~2 key, whose attached factorization tree is now processed.  Factor saturation of \(H\) must expose a
nonaffine tensor-prime factor, since a tensor product
of affine factors is affine.  That factor has smaller arity and is processed
earlier in the induction, contradicting the remaining branch.  Hence
\(F_v\in\cA\) for every accessible \(v\).

It follows from \cref{lem:rm-affine-hyperplane-patching} that
\(g\in\cA\).  All its actual Radon cards are zero or flat-Lagrangian, so
\cref{lem:rm-affine-flat-radon} gives \(g\in\FLag\), contradicting the
choice of a surviving nonflat occurrence.  Therefore every retained
tensor-prime occurrence is eventually certified flat-Lagrangian.

Every concrete terminal, anchor, factor, or group-change batch above is
closed before the next frozen key is selected.  A stable-tuple transition
strictly decreases \(\Omega\); an unchanged-tuple worklist step strictly
decreases \(\Theta_{\rm RM}\).  The two measures therefore have disjoint
and well-founded responsibilities, so no re-entry or same-arity cycle is
possible.

All certifications were made in the same retained normalized coordinates.
The native \(X\)-edge is itself flat-Lagrangian, and tensor products of
flat-Lagrangian signatures remain flat-Lagrangian by taking direct sums of
their self-dual support directions.  Retained factor saturation reconstructs
every member of the signature set from its certified prime factors.  Hence
the entire retained set and its edge have one common flat-Lagrangian
presentation, which is exactly
\(\Close_{\{\FLag\}}(\cG,G,f)\).
\end{proof}

\subsection{Character-flat extensions of the Reed--Muller calculus}
\label{proofsubsec:rm-pauli-tools}
\label{proofsec:reed-muller}

Allowing all four Bell binaries adds linear signs to the Reed--Muller
network calculus.  We describe the resulting class and the corresponding
signed Radon cards.  These are additional structural results; the
standard-\(V_4\) closure in \cref{thm:standard-v4-relative} uses the
tensor-line argument given there.  Orthogonal complements and
Witt decomposition retain the conventions of
\cref{subsubsec:symmetric-binary-spaces}.

Recall from \cref{eq:character-flat-class} that
\(\mathsf{CFlat}_{2k}\) consists of zero together with signatures
\[
 c(-1)^{\ell(x)}\one[x\in a+L],
 \qquad c\ne0,\quad L=L^\perp.
\]
Thus \(\mathsf{CFlat}\) differs from \(\FLag\) by allowing a linear sign
character on a self-dual affine support.  Here a
\emph{pure-Pauli \(\RMcore\)-network} means a tensor network whose
nonbinary vertices are retained actual copies of \(\RMcore\), whose binary
vertices are actual ordered Bell binaries, and whose internal connections
are native \(X\)-edges.  The next lemma gives the exact Pauli network
calculus for this class.  Its corollary distinguishes the resources
\(\Hcore\) and \(\RMcore\): the four-copy circuit in
\cref{proofsubsec:v4-H6-to-R8} realizes the latter from the former, while
pure Reed--Muller/Pauli networks cannot realize the reverse direction.
The remaining lemmas extend Radon localization and recognition from
\(\FLag\) to \(\mathsf{CFlat}\).

\begin{lemma}[Character-flat Pauli network calculus]
\label{lem:v4-character-flat-calculus}
\begin{enumerate}
\item Every network whose vertices are \(\RMcore\) and actual ordered Bell
      binaries \(B_{d,\ell}\), and whose internal edges are native
      \(X\)-edges, has zero
      boundary or a boundary in \(\mathsf{CFlat}\).
\item If \(\RMcore\) and actual nonzero representatives of all four ordered
      Bell binaries \(B_{d,\ell}\) are retained, then, for every
      \(0\ne f\in\mathsf{CFlat}_{2k}\), an \(O(k^2)\)-size actual
      \(\RMcore\)/ordered-Bell network has boundary equal to a nonzero
      scalar multiple of \(f\).
\item Every nonzero character-flat signature of arity at most six is, after
      a port permutation, a nonzero scalar multiple of a product of ordered
      Bell binaries on a perfect matching.
\item \(\mathsf{CFlat}\) is closed under nonzero scalars, tensor products,
      port permutations, one-port ordered-Bell attachments through native
      \(X\)-edges, and nonzero tensor factors.
\end{enumerate}
\end{lemma}

\begin{proof}
Let \(B\) be the displayed boundary-port set and let
\(V_{\rm vert}\) be the tensor-vertex set.  As in
\cref{lem:rm-flat-contexts}, all displayed boundary ports in this
calculus are unpaired ordinary vertex ports; the pairs below are the internal
wire pairs.  Let \(H_{\rm int}\) be the set of internal half-edges.  For every
vertex let \(L_v\) be the linear direction of its affine support.  The
Reed--Muller direction and the two-coordinate repetition direction of every
Bell binary are self-dual.  Define the internal equality subspace
\(E_{\rm eq}\), the ambient
constraint subspace \(H_{\rm con}\), and the vertex-code sum \(C_0\) by
\[
 \begin{aligned}
 E_{\rm eq}
   &=\{y\in\Ftwo^{H_{\rm int}}:
       y_h=y_{h'}\text{ on every internal pair }\{h,h'\}\},\\
 E_{\rm eq}&=E_{\rm eq}^\perp,
 &H_{\rm con}&=\Ftwo^B\oplus E_{\rm eq},\\
 C_0&=\bigoplus_{v\in V_{\rm vert}}L_v=C_0^\perp.
 \end{aligned}
\]
Exactly as in the proof of \cref{lem:rm-pure-boundary}, the feasible boundary
direction is
\(L=\pi_B(C_0\cap H_{\rm con})=L^\perp\).

Define the feasible half-edge code
\(\mathcal K:=C_0\cap H_{\rm con}\) and its
internal-fiber subspace
\[
 J:=\mathcal K\cap
   (\{\mathbf 0_B\}\oplus\Ftwo^{H_{\rm int}}),
\]
where \(\mathbf 0_B\) is the zero boundary word.
The product of the ordered-Bell signs is a linear character
\((-1)^{\lambda(z)}\) on the half-edge variables.  For a fixed feasible
boundary word the internal fiber is \(z_0+J\), and
\[
 \sum_{j\in J}(-1)^{\lambda(z_0+j)}
 =(-1)^{\lambda(z_0)}\sum_{j\in J}(-1)^{\lambda(j)}.
\]
If \(\lambda|_J\ne0\), every nonempty boundary fiber cancels; otherwise
the same nonzero scalar remains on every fiber and the character descends
to the boundary.  This proves item~1 and shows that partial cancellation
cannot create a non-Lagrangian support.

Conversely, the actual network in \cref{lem:rm-flat-contexts}
realizes the indicator of \(a+L\).  Attaching a Bell binary through \(X\)
induces, up to a nonzero sign, one of \(I,X,Z,Y\), so actual ordered-Bell
attachments implement
  the shift and linear character, proving item~2.  More explicitly,
  \(B_{d,\ell}X\) is respectively \(X,Y,I,Z\) for
  \((d,\ell)=(0,0),(0,1),(1,0),(1,1)\), up to the fixed orientation signs.
Arity zero is the empty matching.  At positive arity at most six,
\cref{lem:rm-low-flat-matching} splits the self-dual direction into
two-coordinate repetition lines.  Support shifts and characters turn
their factors into ordered Bell signatures, proving item~3.

For the factor claim in item~4, let \(f=f_1\otimes f_2\ne0\).
The two projections of its affine support are affine; write its Cartesian
support as \((a_1+L_1)\times(a_2+L_2)\), with
\(L=L_1\oplus L_2\).  Since \(L=L^\perp\), each \(L_j\) is
self-dual in its own coordinate space.  Fixing a supported word of the
other factor restricts the character of \(f\) to a character on
\(a_j+L_j\), up to a nonzero scalar.  Hence \(f_j\) is character-flat.
The other closure properties follow by the same direct-sum and character
identities.
\end{proof}

\begin{corollary}[No pure-Pauli inverse from \(\RMcore\) to \(\Hcore\)]
\label{cor:no-rm-to-H6}
No network using only \(\RMcore\), actual ordered Bell binaries, and native
\(X\)-edges has boundary proportional to \(\Hcore\).
\end{corollary}

\begin{proof}
Every nonzero such boundary lies in \(\mathsf{CFlat}\) by
\cref{lem:v4-character-flat-calculus}.  A character-flat six-port support is
a coset of a self-dual subspace of \(\Ftwo^6\), hence has dimension three.
By direct expansion of \cref{eq:H6}, the support of \(\Hcore\) is the
six-bit even-parity space, of dimension five.  Thus
\(\Hcore\notin\mathsf{CFlat}\).  Consequently the direct
\(\Hcore\)-to-\(\RMcore\) circuit constructed below has no inverse using
only \(\RMcore\), ordered Bell binaries, and native \(X\)-edges.
\end{proof}

The support-coset and character data in
\cref{eq:character-flat-class} can be evaluated by affine Gaussian
elimination.

For every even \(n\), define the accessible-direction set
\[
\operatorname{Acc}_n
 =(\one^n)^\perp\setminus\{0,\one^n\}.
\]
Equivalently, \(\operatorname{Acc}_n\) consists of the nonzero even-weight
vectors other than \(\one^n\).  In the two smallest even arities,
\[
 \operatorname{Acc}_2=\varnothing,
 \qquad
 \operatorname{Acc}_4
   =\{e_i+e_j:1\le i<j\le4\}.
\]
Recall
\[
 O(n,2):=\{A\in\operatorname{GL}_n(\Ftwo):A^{\mathsf T}A=I_n\}.
\]
For \(A\in O(n,2)\), write
\[
 \operatorname{Graph}(A)=\{(x,Ax):x\in\Ftwo^n\}.
\]
As proved in \cref{lem:rm-radon-accessible}, this direction is self-dual;
let \(\Gamma_A\) denote a fixed actual \(\RMcore/I/X\)-network supplied by
\cref{lem:rm-flat-contexts} whose boundary is a nonzero scalar multiple of
\(\one[\operatorname{Graph}(A)]\).

\begin{lemma}[Actual signed Radon cards]
\label{lem:v4-signed-radon}
Assume \(\RMcore\) and actual representatives of all four ordered Bell
binaries are retained, and let \(g\) have even arity \(n\).  For each
\(v\in\operatorname{Acc}_n\), use the recorded data
\((A_v,i_v,j_v,r_v,\Pi_v,W_v)\) fixed in
\cref{lem:rm-radon-accessible}, and fix the actual graph network
\(\Gamma_{A_v}\).  For \(d,\ell\in\Ftwo\), using one copy of \(g\), this
graph network, native \(X\)-edges, and the actual ordered Bell binary
\(B_{d,\ell}\), one directly realizes, up to a known nonzero scalar, the
signed card
\begin{equation}
 \mathcal S_{v,d,\ell}g([x]_v)
 =(-1)^{\ell(A_vx)_{i_v}}
  \bigl(g(x)+(-1)^\ell g(x+v)\bigr),
 \qquad [x]_v\in H_{v,d}/\Span_{\Ftwo}\{v\}.
 \label{eq:v4-signed-radon}
\end{equation}
The cases \(\ell=0\) and \(\ell=1\) are called respectively the
\emph{plus} and \emph{minus} signed cards on layer \(d\).  The prefactor
\((-1)^{\ell(A_vx)_{i_v}}\) is the exterior Bell sign in the chosen
quotient coordinates; after stripping this sign and the recorded nonzero
scalar, the two orbit numerators are
\(g(x)+g(x+v)\) and \(g(x)-g(x+v)\).
The remaining \(n-2\) graph-output coordinates give an explicit affine
isometry from this quotient to the boundary coordinate space.  Replacing
\((A_v,i_v,j_v)\) by another such choice changes the displayed analytic
card only by a quotient affine isometry, a nonzero scalar, and a linear sign
character.
\end{lemma}

\begin{proof}
Connect \(g\) to the input side of \(\Gamma_{A_v}\) through native
\(X\)-edges, and let \(x\) denote the variables at \(g\).  The graph input
is therefore \(x+\one\).  Every \(A\in O(n,2)\) fixes \(\one\): invariance
of \(z\cdot z=z\cdot\one\) gives \(A^{\mathsf T}\one=\one\), and hence
\(A\one=\one\).  Thus the graph output is
\(z=A_v(x+\one)=A_vx+\one\).  Connect output ports \(i_v,j_v\) to
\(B_{d,\ell}\) through their two native \(X\)-edges.  The ordered Bell
arguments are then
\[
 u=z_{i_v}+1=(A_vx)_{i_v},
 \qquad
 w=z_{j_v}+1=(A_vx)_{j_v}.
\]
Consequently its support equation is
\[
 u+w
 =x\cdot A_v^{\mathsf T}(e_{i_v}+e_{j_v})
 =x\cdot v=d.
\]
For fixed remaining graph-output coordinates there are two lifts, differing
by \(e_{i_v}+e_{j_v}\) in \(z\) and hence by \(v\) in \(x\).  The Bell sign
is \((-1)^{\ell u}=(-1)^{\ell(A_vx)_{i_v}}\) and changes between those
lifts by \((-1)^\ell\), giving
\cref{eq:v4-signed-radon}.  The graph relation is single-valued before the
two-port contraction, so the realization scalar is nonzero.

Each choice identifies the quotient affinely and isometrically with its
remaining graph-output coordinates.  The transition between two choices is
therefore an affine isometry.  Their selected-coordinate functionals differ
on the quotient by a linear form and a constant, which gives the last
assertion.
\end{proof}

\begin{lemma}[Character-aware two-chart localization]
\label{lem:v4-character-two-chart}
Assume the resources of \cref{lem:v4-signed-radon} are retained.  Let
\(g\) have even arity \(n\ge10\), fix
\(v\in\operatorname{Acc}_n\), and suppose that the two plus cards
\(\mathcal S_{v,0,0}g\) and \(\mathcal S_{v,1,0}g\) are zero or
character-flat.  Define, only for the analysis, the orbit-sum signature
\[
 F_v=\mathsf P_v^+g.
\]
If \(F_v\notin\cA\), an actual \(\RMcore\)/ordered-Bell context using one
copy of \(g\) realizes a nonzero \(H\notin\cA\) of arity at most eight.
\end{lemma}

\begin{proof}
Use the recorded data
\((A_v,i_v,j_v,r_v,\Pi_v,W_v)\) from
\cref{lem:rm-radon-accessible}, and put \(W:=W_v\).  The plane \(\Pi_v\)
is nondegenerate, and the map \(\iota_{v,\epsilon}\) identifies \(W\)
isometrically with
\(H_{v,\epsilon}/\Span_{\Ftwo}\{v\}\).  Since \(F_v\) is
invariant under \(T_v\), its two quotient charts are
\[
 f_\epsilon(x):=F_v(\epsilon r_v+x),
 \qquad x\in W.
\]
For each nonzero chart
choose, without asserting uniqueness, representatives
\begin{equation}
 f_\epsilon(x)
 =c_\epsilon(-1)^{\lambda_\epsilon(x)}
  \one[x\in a_\epsilon+M_\epsilon],
 \qquad M_\epsilon=M_\epsilon^\perp\le W.
 \label{eq:v4-chart-presentation}
\end{equation}
Here \(a_\epsilon\) is chosen modulo \(M_\epsilon\), and
\(\lambda_\epsilon\) is an ambient extension of the character on the
support.  Changing either choice changes \(c_\epsilon\) only by a sign, so
membership of \(c_1/c_0\) in \(\mu_4\) is independent of these choices.

If one chart is zero, the other gains only one affine layer equation and
\(F_v\in\cA\).  Suppose both are nonzero.  Their union has affine support
exactly when \(M_0=M_1\).  When the directions agree, the difference of
the two characters is absorbed by an allowed quadratic mixed term with the
layer bit, and the layer scales extend to an affine fourth-root phase
exactly when \(c_1/c_0\in\mu_4\).  Thus \(F_v\notin\cA\) gives
\begin{equation}
 M_0\ne M_1,
 \qquad\text{or}\qquad
 M_0=M_1\ \text{ and }\ c_1/c_0\notin\mu_4.
 \label{eq:v4-chart-obstruction}
\end{equation}

Represent \(\lambda_0+\lambda_1\) by \(\rho\in W\), so that
\((\lambda_0+\lambda_1)(x)=\rho\cdot x\).  When \(M_0\ne M_1\), choose
\(s\in M_0\setminus M_1\) and \(t\in M_1\) with
\(s\cdot t=1\), and put \(\Delta=a_1-a_0\).  The plane
\(\Span_{\Ftwo}\{s,t\}\) is nondegenerate.  Put
\(S_0=\Span\{s,t,\Delta,\rho\}\).  Its radical has dimension at most two;
pairing a basis of that radical with dual vectors gives a nondegenerate
extension \(N\supseteq S_0\) of dimension at most six.  If \(N\) has odd
dimension, then its nondegenerate orthogonal complement is odd-dimensional
and contains an anisotropic vector; adjoining that orthogonal line preserves
the bound and makes the extension even-dimensional.  Denote the result by
\[
 U\supseteq\Span\{s,t,\Delta,\rho\},
 \qquad \dim U\le6.
\]
In the equal-direction case, start from
\(S_0=\Span\{\Delta,\rho\}\).  Pairing its radical needs at most two dual
vectors and gives a nondegenerate extension of dimension at most four; if
its dimension is odd, adjoin an anisotropic vector from its odd-dimensional
nondegenerate complement.  This gives an even-dimensional nondegenerate
\(U\supseteq\Span\{\Delta,\rho\}\) with \(\dim U\le4\).
Put
\[
 P:=U^{\perp_W},
 \qquad W=U\perp P,
\]
and let \(\pi_U,\pi_P\) be the corresponding orthogonal projections.  A
Witt decomposition of the even-dimensional nondegenerate space \(P\)
supplies
\[
 J\le P,
 \qquad J=J^{\perp_P}.
\]
Because \(\Delta,\rho\in U\), the two shifts have the same
\(P\)-projection
\[
 p_0:=\pi_P(a_0)=\pi_P(a_1),
\]
and the two characters have the same restriction
\(\lambda_P:=\lambda_0|_P=\lambda_1|_P\).

Let \(S:=\Pi_v\perp U\), put \(m:=\dim S\), and choose an isometry
\(\phi:S\to\Ftwo^m\).  Such an isometry exists because \(S\) contains the
nonalternating plane \(\Pi_v\).  We have \(m\le8\) in the
different-direction case and \(m\le6\) in the equal-direction case.  Define
\[
 R_0:=\{(s_0+j,\phi(s_0)):s_0\in S,\ j\in J\}
 \le\Ftwo^n\oplus\Ftwo^m.
\]
For two elements of \(R_0\), their inner product is
\[
 (s_0+j)\cdot(s'_0+j')
 +\phi(s_0)\cdot\phi(s'_0)=j\cdot j'=0.
\]
Moreover \(\dim R_0=m+(n-m)/2=(n+m)/2\), so \(R_0=R_0^\perp\).
On \(R_0\), define the linear sign character
\[
 \chi(s_0+j,\phi(s_0)):=(-1)^{\lambda_P(j)}.
\]
Choose an ambient linear-character extension and denote it again by
\(\chi\).  Hence, for
\((x,y)\in\Ftwo^n\oplus\Ftwo^m\),
\begin{equation}
 C(x,y):=\chi(x-p_0,y)\,
 \one[(x,y)\in(p_0,0)+R_0]
 \label{eq:v4-localizing-context}
\end{equation}
is character-flat.  By \cref{lem:v4-character-flat-calculus}, choose an
actual \(O(n^2)\)-size \(\RMcore\)/ordered-Bell realization whose boundary
is \(\kappa C\), where the realization scalar
\(\kappa\in\AlgNums^\times\) is recorded.  Any fixed
coordinate complements contributed by its native \(X\)-attachments are
absorbed into the displayed affine shift.

Contract this context with \(F_v\).  Its boundary is
\begin{equation}
 H^+(\phi(s_0))
 =\kappa\sum_{j\in J}(-1)^{\lambda_P(j)}
   F_v(s_0+p_0+j),
 \qquad s_0\in S.
 \label{eq:v4-localized-sum}
\end{equation}
In the different-direction case, write
\(a_\epsilon=u_\epsilon+p_0\) with \(u_\epsilon\in U\).  The three boundary
points \(\phi(u_0)\), \(\phi(u_0+s)\), and
\(\phi(r_v+u_1)\) give nonzero sums.  On each support fiber the context
character cancels the common \(P\)-restriction of the chart character, so
there is no cancellation.  The fourth point
\(\phi(r_v+u_1+s)\) has value zero: otherwise some \(j\in J\subseteq P\)
would satisfy \(s+j\in M_1\), and then
\[
 0=(s+j)\cdot t=s\cdot t=1.
\]
Thus the support of \(H^+\) is not affine.  In the equal-direction case, the
two corresponding fibers have the same nonzero cardinality
\(|J\cap M_0|\).  Their value ratio differs from \(c_1/c_0\) only by a
sign, and therefore remains outside \(\mu_4\).  In both cases
\(H^+\ne0\) and \(H^+\notin\cA\).

Finally, contract the same actual context with \(g\), and call its boundary
\(H\).  Put \(w:=\phi(v)\).  Translation by \(v\) on the input of
\cref{eq:v4-localizing-context} is translation by \(w\) on its boundary,
and the context character is trivial on \((v,w)\).  Therefore
\[
 H^+=\mathsf P_w^+H.
\]
If \(H\in\cA\), then \cref{lem:affine-two-translate-closure} would give
\(H^+\in\cA\), a contradiction.  Thus \(H\ne0\), \(H\notin\cA\), and its
arity \(m\) is at most eight.
\end{proof}

\begin{lemma}[Affine-to-character-flat recognition]
\label{lem:v4-affine-character-radon}
Let \(0\ne g\in\cA\) have even arity \(2r\ge10\).  If every signed card
\(\mathcal S_{v,d,\ell}g\), for
\(v\in\operatorname{Acc}_{2r}\) and \(d,\ell\in\Ftwo\), is zero or
character-flat, then \(g\in\mathsf{CFlat}\).
\end{lemma}

\begin{proof}
Input translation permutes the two values of \(d\) and reindexes the
corresponding quotient by an affine isometry.  Under this reindexing, it
changes a signed card only by a nonzero scalar and a linear sign character.
We may therefore
translate the affine support to a linear subspace and write
\begin{equation}
 g(x)=c\,i^{Q(x)}\one[x\in U],
 \qquad c\ne0,
 \qquad d_U:=\dim U.
 \label{eq:v4-recognition-normal-form}
\end{equation}
Put \(E:=\one^\perp\).  Every nonzero character-flat card has arity
\(2r-2\) and hence support dimension \(r-1\).

We first prove \(U=U^\perp\).  If \(E\subseteq U\), then
\(d_U\ge2r-1\).  For an accessible \(v\in E\), at least one signed card on
a nonempty layer is nonzero.  Restriction to the layer and quotienting by
\(\Span_{\Ftwo}\{v\}\) can each remove at most one dimension.  More
explicitly, on an overlapping support orbit, factoring out \(g(x)\) leaves
\[
 1\pm\zeta(-1)^{\theta(x)},
 \qquad \zeta\in\mu_4,
\]
where \(\theta\) is a binary linear form.  If
\(\zeta\in\{\pm i\}\), this factor is nowhere zero.  If
\(\zeta\in\{\pm1\}\), its nonzero set is either empty, the whole current
support, or one affine hyperplane.  Thus a nonzero signed card loses at
most one additional support dimension through cancellation and has support
dimension at least
\[
 d_U-3\ge2r-4>r-1,
\]
a contradiction.  If \(E\subseteq U^\perp\), then \(d_U\le1\); choosing
accessible \(v\in E\setminus U\) gives a nonzero card of support dimension
\(d_U\), again a contradiction.  Hence neither containment holds.

The two proper subspaces \(E\cap U\) and \(E\cap U^\perp\) do not cover
\(E\).  Since \(\dim E=2r-1\ge9\), choose
\[
 v\in\operatorname{Acc}_{2r}\setminus(U\cup U^\perp).
\]
There is no cancellation because \(v\notin U\), and each nonempty card has
support dimension \(d_U-1\).  Thus \(d_U=r\).

No accessible vector can lie in \(U^\perp\setminus U\), because such a
vector would give a nonzero card of support dimension \(r\).  Put
\(A:=E\cap U^\perp\).  It follows that
\begin{equation}
 A\setminus\{0,\one\}\subseteq U.
 \label{eq:v4-A-contained}
\end{equation}
If \(\one\notin A\), or if \(\one\in A\cap U\), this gives
\(A\subseteq U\).  If instead \(\one\in A\setminus U\), then
\(\dim A\ge r-1\ge4\); choose \(z\in A\setminus\{0,\one\}\).  Both
\(z\) and \(z+\one\) are accessible and lie in \(U\) by the preceding
display, forcing \(\one\in U\), a contradiction.  Thus \(A\subseteq U\)
in every case.

If \(U^\perp\subseteq E\), equal dimensions give \(U=U^\perp\).  Otherwise
\(\dim A=r-1\).  Since \(A\subseteq U\) and \(\dim U=r\), choose
\(u\in U\setminus A\), so
\(U=A\oplus\Span_{\Ftwo}\{u\}\).  Because
\(A\subseteq U^\perp\), the vector \(u\) is orthogonal to \(A\).  If \(u\)
were even, then \(u\cdot u=0\), hence \(u\in U^\perp\) and
\(U\subseteq U^\perp\); equal dimensions would again give
\(U=U^\perp\).  Therefore \(u\) is odd and
\begin{equation}
 U=A\oplus\Span_{\Ftwo}\{u\},
 \qquad u\cdot u=1.
 \label{eq:v4-near-self-dual}
\end{equation}
Since \(u\) is odd, \(\one\cdot u=1\), whereas every vector of \(A\) is
orthogonal to \(u\); hence \(\one\notin A\).  Fix \(0\ne s\in A\).
Here \(\dim A=r-1\ge4\), so \(s\) is accessible.  For this \(s\), a nonzero signed
card has either the full quotient direction
\(U/\Span_{\Ftwo}\{s\}\), of
dimension \(r-1\), or a codimension-one subdirection of dimension \(r-2\).
The first is not self-dual because its orthogonal direction is
\(U^\perp/\Span_{\Ftwo}\{s\}\ne
U/\Span_{\Ftwo}\{s\}\); the second has the
wrong dimension.  On the supported layer \(d=0\), at least one of the plus
and minus cards is nonzero: after stripping their respective exterior Bell
signs and recorded nonzero scalars, their orbit brackets are
\[
 g(x)+g(x+s)
 \quad\text{and}\quad
 g(x)-g(x+s),
\]
whose sum is \(2g(x)\).  This
contradicts the hypothesis.  Therefore
\[
 U=U^\perp.
\]

The orthogonal reduction in the proof of
\cref{lem:rm-flat-contexts} sends \(U\) to the code
\[
 \mathsf{Pair}_r
 :=\{x(t):t\in\Ftwo^r\},
 \qquad
 x(t):=(t_1,t_1,\ldots,t_r,t_r).
\]
Orthogonal maps preserve \(\one\) and hence accessible directions.  Their
induced quotient affine isometries, together with nonzero scalar and
linear-character changes, preserve character-flatness, so the signed-card
hypothesis is invariant under this coordinate reduction.
Put
\[
 g(x(t))=c\,h(t),
\]
where
\begin{equation}
 h(t)=i^{q_0+\sum_j\alpha_jt_j+
             2\sum_{j<k}\beta_{jk}t_jt_k},
 \qquad q_0,\alpha_j\in\mathbb Z_4,
 \quad \beta_{jk}\in\Ftwo.
 \label{eq:v4-pair-phase}
\end{equation}
Extend the quadratic coefficients symmetrically by setting
\(\beta_{jk}:=\beta_{kj}\) whenever \(k<j\), and put \(\beta_{jj}:=0\).
Let \(\bar e_j\) be the \(j\)-th basis vector of \(\Ftwo^r\), and define
the \(j\)-th logical pair and its physical accessible direction by
\[
 P_j:=\{2j-1,2j\},
 \qquad
 v_j:=e_{2j-1}+e_{2j}.
\]
Let
\[
 \mathcal T_j:=\{t\in\Ftwo^r:t_j=0\}.
\]
Using the representatives \(t\in\mathcal T_j\), and stripping the exterior
Bell signs and recorded nonzero scalars, the \(d=0\) plus and minus cards
in the \(v_j\)-direction have coefficient functions
\begin{equation}
 A_j(t)=h(t)+h(t+\bar e_j),
 \qquad B_j(t)=h(t)-h(t+\bar e_j),
 \qquad t\in\mathcal T_j.
 \label{eq:v4-pair-cards}
\end{equation}
Whenever one of these cards has support \(\mathcal T_j\),
character-flatness makes its coefficient function a nonzero scalar times
a sign character on \(\mathcal T_j\).

Define
\[
 L_j(t):=\sum_{k\ne j}\beta_{jk}t_k\in\Ftwo,
 \qquad
 \frac{h(t+\bar e_j)}{h(t)}=i^{\alpha_j}(-1)^{L_j(t)}.
\]
Suppose first that some \(\alpha_j\) is even.  If \(L_j\ne0\), the two
functions in \cref{eq:v4-pair-cards} are supported on complementary proper
hyperplanes, so neither nonzero card is character-flat.  Thus \(L_j=0\),
which gives \(\beta_{jk}=0\) for every \(k\ne j\).  One of \(A_j,B_j\) is
then zero and the other is a nonzero scalar multiple of
\(h|_{t_j=0}\).  Its being a sign character forces
\[
 \alpha_k\equiv0\pmod2\quad(k\ne j),
 \qquad \beta_{km}=0\quad(k,m\ne j).
\]
Together with the evenness of \(\alpha_j\), this proves that all
\(\alpha_k\) are even and all \(\beta_{km}\) vanish.

It remains to exclude the alternative in which every \(\alpha_j\) is odd.
Write \(i^{\alpha_j}=\delta_j i\), where \(\delta_j\in\{\pm1\}\).  The
plus bracket in \(A_j\) satisfies
\begin{equation}
 1+\delta_j i(-1)^{L_j}
 =(1+\delta_j i)(-\delta_j i)^{L_j}.
 \label{eq:v4-parity-bracket}
\end{equation}
The multilinear modulo-four expansion of the Boolean parity \(L_j\) is
\begin{equation}
 L_j(t)\equiv
 \sum_{k\ne j}\beta_{jk}t_k
 +2\sum_{\substack{k<m\\k,m\ne j}}
    \beta_{jk}\beta_{jm}t_kt_m
 \pmod4.
 \label{eq:v4-parity-expansion}
\end{equation}
Since the nonzero full-support card \(A_j\) is a sign character, all mixed
quadratic coefficients in its phase vanish.  Hence, for all distinct
\(j,k,m\),
\begin{equation}
 \beta_{km}=\beta_{jk}\beta_{jm}.
 \label{eq:v4-beta-relation}
\end{equation}
View \(\beta\) as the adjacency matrix of a simple graph.  These relations
make the graph either empty or complete.  Indeed, if
\(\beta_{ab}=1\), then for every \(c\notin\{a,b\}\),
\[
 1=\beta_{ab}=\beta_{ca}\beta_{cb},
\]
so \(c\) is adjacent to both \(a\) and \(b\); applying
\cref{eq:v4-beta-relation} with \(j=a\) joins every remaining pair.

The empty case is impossible: when \(L_j=0\), either nonzero card in
\cref{eq:v4-pair-cards} is a scalar multiple of \(h|_{t_j=0}\), whose
ratio in any remaining coordinate \(k\) is
\(i^{\alpha_k}=\pm i\), not the \(\pm1\) ratio of a sign character.

In the complete case, choose distinct logical pairs \(a,b,c\) and the
physical weight-two direction
\[
 w:=e_{2a-1}+e_{2b-1}\in\operatorname{Acc}_{2r}.
\]
This vector is not in \(\mathsf{Pair}_r\).  The support of the plus card
\(\mathcal S_{w,0,0}g\) is the image of
\(\{t:t_a+t_b=0\}\), and on the two-point \(w\)-translation orbit
represented by \(x(t)\) its
coefficient is \(c\,h(t)\), because \(x(t)+w\notin U\).  Toggling the
third logical pair \(c\) preserves this support, but its coefficient ratio
is
\[
 \frac{h(t+\bar e_c)}{h(t)}
 =i^{\alpha_c}(-1)^{t_c+\sum_{k\ne c}t_k}\in\{\pm i\}.
\]
A character-flat signature has ratio \(\pm1\) along every direction of its
support.  This contradiction eliminates the complete case.

Thus the all-odd alternative is impossible.  The earlier even-\(\alpha_j\)
case applies, and \cref{eq:v4-pair-phase} reduces to a nonzero scalar times
a linear sign character on the self-dual support \(U\).  This is exactly
\cref{eq:character-flat-class}.
\end{proof}
\section{Klein-Four Matching Decks}
\label{sec:v4-deck}

The two marked Klein-four groups require different arguments.  For the
standard Pauli group, pair contractions detect whether a complex tensor
is a scalar multiple of a real tensor; this reduces the remaining case to
the real-valued Holant dichotomy.  For the exotic group, contractions
compare four matrix labels on matching edges and force a proper tensor
factor.  We use the Bell notation and cores \(\Hcore,\RMcore\) of
\cref{subsec:pauli-notation}.  The six-port exotic classification remains
the finite certificate in Appendix~\ref{appsec:cert-klein-four}.

\begin{theorem}[Klein-four relative rigidity]
\label{thm:v4-relative}
Assume
\[
 \Stable(\Lambda,G,f),\qquad \arity(f)=\nu(\Lambda),
 \qquad G\in\{V_4,\mathcal V_{\rm ex}\}.
\]
In the standard Pauli case \(G=V_4\), assume additionally that
\(\nu(\Lambda)\ge6\).  Then
\(\Close_{\{\cA\}}(\Lambda,G,f)\).
\end{theorem}

The two cases are proved in
\cref{thm:standard-v4-relative,thm:exotic-v4-relative-rigidity}.
As in the definition of \(\Close\), a resolved \(\Tract\)-leaf
completes the branch.

\begin{proofstrategy}
In the standard Pauli form, pair contractions determine a tensor line
above arity four.  A retained signature that is not a scalar times a real
tensor therefore yields an actual quaternary tensor with the same
property.  Factor saturation gives a terminal or a strict outer successor.
When every retained signature is a scalar times a real tensor, the
real-valued Holant dichotomy closes the entire retained set.  The explicit
\(\Hcore\)-to-\(\RMcore\) circuit is retained for the Platonic analysis.

In the exotic Klein form, the physical exotic calculus and the six-port base
first force a common residual ruling, after which fusion constrains the
label maps to permutation fibres.  Two choices of a dual-cross contraction
exclude a nonfactor at arity eight; at arity at least ten, one such
contraction suffices.  These factor conclusions exclude a higher-arity
stable prime, while the quaternary interface gives a resolved leaf or a
strict outer successor.
\end{proofstrategy}

\subsection{The standard \texorpdfstring{\(V_4\)}{V4} form: actual Pauli contractions}
\label{proofsec:v4}
\label{proofsubsec:v4-pauli-calculus}

Let \(\Lambda\) be retained and factor-saturated, with the distinguished
literal equality \(I\) and complete group \(V_4\).  Let the Bell-label space be
\(\mathsf L:=\Ftwo^2\).  For a Bell label
\(\gamma=(d,\ell)\in\mathsf L\), let \(B_\gamma=B_{d,\ell}\) be the ordered
Bell signature defined in \cref{subsec:pauli-notation}.  The chosen actual
representative fixed in \cref{subsec:reduction-conventions} is
\(\beta_\gamma B_\gamma\), where
\(\beta_\gamma\in\AlgNums^\times\) is its known nonzero realization scalar.
The formulas below divide actual outputs by these recorded scalars.

Here and below, a \emph{companion} is a retained tensor-prime occurrence
other than the distinguished minimum core; its occurrence carries its
actual/Turing provenance.  For
\(\alpha=(d,\ell)\in\mathsf L\), put
\begin{equation*}
 P_\alpha:=Z^\ell X^d.
\end{equation*}
All contractions below are bilinear, without complex conjugation.
For an ordered pair \(p\), define
\begin{equation*}
 C_p^\gamma(h):=\left\langle B_\gamma^{(p)},h\right\rangle_p,
 \qquad \gamma\in\mathsf L,
\end{equation*}
and omit the argument only when the parent tensor is unambiguous.  For a
list \(\mathbf p=(p_1,\ldots,p_s)\) of disjoint ordered pairs and
\(\boldsymbol\gamma=(\gamma_1,\ldots,\gamma_s)\), put
\begin{equation*}
 C_{\mathbf p}^{\boldsymbol\gamma}(h)
 :=C_{p_s}^{\gamma_s}\cdots C_{p_1}^{\gamma_1}(h).
\end{equation*}
The order is immaterial for disjoint pairs.  A \emph{proper Pauli card}
means such an actual iterated contraction leaving at least two ports; a
\emph{Pauli pair card} has \(s=1\).  Dividing the physical output by its
recorded nonzero realization scalar gives the displayed canonical card.
  The phrase
\emph{local-Pauli
orbit} has the precise projective meaning
\begin{equation*}
 \operatorname{Orb}_{V_4}(h)
 :=\left\{c\left(\bigotimes_{j=1}^{\arity(h)}P_{\alpha_j}\right)(\pi h):
 c\in\mathbb C^\times,\ \alpha_j\in\mathsf L,\ \pi\in S_{\arity(h)}\right\};
\end{equation*}
port reversal signs and realization scalars are included in the recorded
nonzero scalar \(c\).

A \emph{complete Bell block} is an ordered tensor \(B_\gamma^{(p)}\) with
both ports of \(p\) available for attachment.  An \emph{\(I\)-path} between
two equally oriented complete blocks is one copy of the retained equality
\(I\) joining corresponding endpoints; connecting the blocks uses the two
corresponding \(I\)-paths, with their recorded realization scalars.

\subsection{The Pauli cores}
\label{proofsubsec:v4-hcore}
\label{proofsubsec:v4-core-atlas}

For the fixed ordered pairs \(A=(1,2),B=(3,4),C=(5,6)\), and with the
order-three Bell-label map \(\Theta\) defined in
\cref{subsec:pauli-notation}, define the six-port Pauli core signature
\begin{equation*}
 \Hcore=\sum_{\gamma\in\mathsf L}
 B_\gamma^A B_{\Theta\gamma}^B B_{\Theta^2\gamma}^C.
\end{equation*}
Direct expansion gives
\begin{align*}
 \supp(\Hcore)
   &=\left\{x\in\Ftwo^6:\sum_{j=1}^6x_j=0\right\},\\
 \Hcore(x)&=(-1)^{x_1+x_1x_2+x_1x_6+x_2x_5+x_3x_5+x_3x_6}
   \quad(x\in\supp(\Hcore)).
\end{align*}
Hence \(\Hcore\in\cA\), and its even-parity support is tensor-prime.
By \cref{rem:sc-f8-core}, \(\Hcore\) and \(\RMcore\) are
the Shao--Cai \(\widehat f_6\) and \(f_8=\widehat f_8\), up to the stated
port permutation.  Consequently every Pauli pair card of \(\Hcore\) is a
nonzero product of two Pauli binaries, while every Pauli pair card of
\(\RMcore\) is, up to a nonzero scalar, a product of three copies of the
contracted Pauli binary; see
\cite[full version, Definition~6.12 and Lemma~6.13]{ShaoCai2020} and
\cite[Theorems~13--14]{CaiFuShao2020Entanglement}.

\subsection{The direct \texorpdfstring{\(\Hcore\)-to-\(\RMcore\)}{H6-to-RM8} circuit}
\label{proofsubsec:v4-H6-to-R8}

In four copies \(\Hcore^{(v)}\), write the ordered blocks
\((\mathsf A_v,\mathsf B_v,\mathsf C_v)
=(B_{\gamma_v},B_{\Theta\gamma_v},B_{\Theta^2\gamma_v})\)
and leave the \(\mathsf A_v\) external.  Connect equally oriented complete
blocks by
\begin{equation*}
  \mathsf B_1-\mathsf B_2,\qquad
  \mathsf C_2-\mathsf C_3,\qquad
  \mathsf B_3-\mathsf B_4,\qquad
  \mathsf C_4-\mathsf C_1.
\end{equation*}

\begin{figure}[htbp]
\centering
\begin{tikzpicture}[
  block/.style={
    draw,rounded corners=1pt,
    minimum width=10mm,minimum height=5.5mm,
    inner sep=1pt,font=\scriptsize\bfseries
  },
  Ablock/.style={
    block,draw=blue!65!black,fill=blue!16,
    text=blue!55!black
  },
  Bblock/.style={
    block,draw=orange!75!black,fill=orange!20,
    text=orange!65!black
  },
  Cblock/.style={
    block,draw=teal!70!black,fill=teal!18,
    text=teal!60!black
  },
  corebox/.style={
    draw=black!45,dashed,rounded corners=2pt,
    line width=.45pt
  },
  corelabel/.style={font=\scriptsize},
  apath/.style={draw=blue!65!black,line width=.55pt},
  bpath/.style={draw=orange!75!black,line width=.5pt},
  cpath/.style={draw=teal!70!black,line width=.5pt},
  outdot/.style={
    circle,fill=blue!65!black,
    minimum size=1.6pt,inner sep=0pt
  }
]
  \draw[corebox] (-4.25,.55) rectangle (-1.85,2.95);
  \draw[corebox] ( 1.85,.55) rectangle ( 4.25,2.95);
  \draw[corebox] ( 1.85,-2.95) rectangle ( 4.25,-.55);
  \draw[corebox] (-4.25,-2.95) rectangle (-1.85,-.55);

  \node[corelabel] at (-3.55, 1.75) {$\Hcore^{(1)}$};
  \node[corelabel] at ( 3.55, 1.75) {$\Hcore^{(2)}$};
  \node[corelabel] at ( 3.55,-1.75) {$\Hcore^{(3)}$};
  \node[corelabel] at (-3.55,-1.75) {$\Hcore^{(4)}$};

  \node[Ablock] (A1) at (-3.55, 2.55) {$\mathsf A_1$};
  \node[Bblock] (B1) at (-2.40, 1.75) {$\mathsf B_1$};
  \node[Cblock] (C1) at (-3.55,  .95) {$\mathsf C_1$};

  \node[Ablock] (A2) at ( 3.55, 2.55) {$\mathsf A_2$};
  \node[Bblock] (B2) at ( 2.40, 1.75) {$\mathsf B_2$};
  \node[Cblock] (C2) at ( 3.55,  .95) {$\mathsf C_2$};

  \node[Ablock] (A3) at ( 3.55,-2.55) {$\mathsf A_3$};
  \node[Bblock] (B3) at ( 2.40,-1.75) {$\mathsf B_3$};
  \node[Cblock] (C3) at ( 3.55,-.95) {$\mathsf C_3$};

  \node[Ablock] (A4) at (-3.55,-2.55) {$\mathsf A_4$};
  \node[Bblock] (B4) at (-2.40,-1.75) {$\mathsf B_4$};
  \node[Cblock] (C4) at (-3.55,-.95) {$\mathsf C_4$};

  \draw[bpath]
    ([yshift=1.1mm]B1.east) -- ([yshift=1.1mm]B2.west);
  \draw[bpath]
    ([yshift=-1.1mm]B1.east) -- ([yshift=-1.1mm]B2.west);

  \draw[cpath]
    ([xshift=-1.6mm]C2.south) -- ([xshift=-1.6mm]C3.north);
  \draw[cpath]
    ([xshift=1.6mm]C2.south) -- ([xshift=1.6mm]C3.north);

  \draw[bpath]
    ([yshift=1.1mm]B4.east) -- ([yshift=1.1mm]B3.west);
  \draw[bpath]
    ([yshift=-1.1mm]B4.east) -- ([yshift=-1.1mm]B3.west);

  \draw[cpath]
    ([xshift=-1.6mm]C4.north) -- ([xshift=-1.6mm]C1.south);
  \draw[cpath]
    ([xshift=1.6mm]C4.north) -- ([xshift=1.6mm]C1.south);

  \draw[apath] ([xshift=-1.6mm]A1.north) -- ++(0,.60)
    node[outdot] {};
  \draw[apath] ([xshift= 1.6mm]A1.north) -- ++(0,.60)
    node[outdot] {};
  \draw[apath] ([xshift=-1.6mm]A2.north) -- ++(0,.60)
    node[outdot] {};
  \draw[apath] ([xshift= 1.6mm]A2.north) -- ++(0,.60)
    node[outdot] {};

  \draw[apath] ([xshift=-1.6mm]A3.south) -- ++(0,-.60)
    node[outdot] {};
  \draw[apath] ([xshift= 1.6mm]A3.south) -- ++(0,-.60)
    node[outdot] {};
  \draw[apath] ([xshift=-1.6mm]A4.south) -- ++(0,-.60)
    node[outdot] {};
  \draw[apath] ([xshift= 1.6mm]A4.south) -- ++(0,-.60)
    node[outdot] {};
\end{tikzpicture}
\caption{The four-copy \(\Hcore\)-to-\(\RMcore\) circuit.  Each doubled
internal connection represents the two corresponding \(I\)-paths between
complete blocks.  The blue external legs are the eight ports of the
resulting \(\RMcore\).}
\label{fig:v4-H6-R8-gadget}
\end{figure}

Each connection uses two legal \(I\)-paths and contributes
\(2\delta_{\gamma_u,\gamma_v}\), so all labels agree without cancellation:
\begin{equation}
  2^4\sum_{\gamma\in\mathsf L}B_\gamma^{\otimes4}=32\RMcore.
  \label{eq:app-v4-H6-R8-output}
\end{equation}
Restoring the known nonzero realization and path scalars gives
\begin{equation*}
  \KHolant(\Lambda,\RMcore)\leT
  \KHolant(\Lambda,\Hcore).
\end{equation*}
If \(\Hcore\) is already actual or Turing-exposed over \(\Lambda\), then
\(\KHolant(\Lambda,\Hcore)\leT\KHolant(\Lambda)\) as well.
In the actual case, adjoining the circuit output cannot enlarge the
complete actual binary group, which remains \(V_4\).
The circuit itself only requires \(\Hcore\) and actual \(I\)-paths, so it
also applies in an ambient language with a larger complete group.
A reverse gadget construction using only \(\RMcore\) and pure Pauli
resources is excluded by \cref{cor:no-rm-to-H6}.


\subsection{The relative implication}
\label{proofsubsec:v4-relative-implication}

The next identity extracts information from all pair contractions without
reconstructing a tensor from an enumerated list.  It compares two tensors
by their antisymmetrized product \(u\otimes v-v\otimes u\): this
product vanishes exactly when they lie on the same line.  Pair contractions
will force it to vanish once there are more than four ports.

\begin{lemma}[Pauli contractions determine a tensor line]
\label{lem:v4-pauli-projective-determination}
Let \(n\ge5\), and let \(u,v\in(\mathbb C^2)^{\otimes n}\).
If \(C_{rs}^\gamma(u)\) and \(C_{rs}^\gamma(v)\) are linearly dependent
for every pair \(r<s\) and every \(\gamma\in\mathsf L\), with zero
cards allowed, then \(u,v\) are linearly dependent.
Consequently, if every Pauli pair card of \(u\) is zero or a nonzero
complex scalar times a real tensor, the same holds for \(u\).
\end{lemma}

\begin{proof}
Put \(W=u\otimes v-v\otimes u\) and group its two copies by physical
port, so that \(W\in(\mathbb C^4)^{\otimes n}\).  The hypothesis says
\[
 (C_{rs}^\gamma\otimes C_{rs}^\gamma)W=0
 \qquad(r<s,\ \gamma\in\mathsf L).
\]
On the regrouped four-dimensional spaces at \(r,s\), the four
contraction matrices are
\[
 I\otimes I,\quad Z\otimes Z,\quad X\otimes X,
 \quad (ZX)\otimes(ZX).
\]
Use the same orthonormal Bell basis
\(2^{-1/2}(\operatorname{vec}I,\operatorname{vec}Z,
\operatorname{vec}X,\operatorname{vec}(ZX))\) on each regrouped port.
These matrices are diagonal in that basis, with eigenvalue rows
\[
 \begin{pmatrix}
  1&1&1&1\\
  1&1&-1&-1\\
  1&-1&1&-1\\
  1&-1&-1&1
 \end{pmatrix}.
\]
The rows are independent.  Thus the four contraction equations, for any
fixed values on the other ports, force each coefficient with equal Bell
basis indices at \(r,s\) to vanish.  Applying this to every pair leaves
only words whose \(n\) indices are all distinct.  There are just four
indices, so \(n\ge5\) gives \(W=0\), proving the first assertion.

For the last assertion take \(v=\overline u\).  The Pauli kernels are
real, and a vector and its conjugate are dependent precisely when it is
zero or a scalar times a real vector: divide by its first nonzero
coordinate to see the reverse implication.  Conjugation and the Bell
basis above are used only to prove this identity; they are not gadgets.
\end{proof}

\begin{theorem}[Standard Pauli \(V_4\) relative closure]
\label{thm:standard-v4-relative}
If
\[
 \Stable(\cG,V_4,f),\qquad \nu(\cG)=\arity(f)\ge6,
\]
then \(\Close_{\{\cA\}}(\cG,V_4,f)\).
\end{theorem}

\begin{proof}
Write \(\Lambda=\cG\).  Dispatch odd, unary, and unsafe binary factors
by the established reductions in
\cref{thm:external-odd,lem:common-unary-exit,lem:binary-interpolation,%
lem:factor-saturation,lem:actual-binary-completeness}.
On the continuing branch every nonbinary prime is even, and all binary
prime factors are nonsingular Pauli binaries up to nonzero scalars.

Suppose some retained signature is not a scalar times a real tensor.
Factor saturation supplies such a tensor-prime signature \(g\), necessarily
of even arity \(n\ge\nu(\Lambda)\ge6\).  The contrapositive of
\cref{lem:v4-pauli-projective-determination} selects a Pauli pair card
with the same property.  Repeat while its arity is at least six, reaching
a nonzero quaternary tensor \(q\) that is not a scalar times a real tensor.
All selected contractions are actual: complete \(V_4\) supplies the four
Pauli kernels, including their reversal and native-edge scalars.  Their
composition uses one \(g\)-vertex and \(O(n)\) binary vertices.  The
intermediate coefficient tables only select this one finite network.

Adjoin \(q\) and process its complete factor batch by
\cref{def:relative-stratum-protocol}.  An odd or unsafe factor is resolved;
a new finite binary class gives a strict group successor, with the deck
regenerated.  Otherwise, if all factors were binaries, they would all be
Pauli and their product would be a scalar times a real tensor, contrary to
the choice of \(q\).  Hence an even nonbinary prime remains.  It has
arity four, so the retained augmentation strictly lowers
\(\nu(\Lambda)\ge6\).  Retained \(I\) excludes the one-sided exception
in factor saturation.  Thus this entire batch gives an established leaf
or a strict outer successor; it opens no new unresolved branch with
unchanged group and \(\nu\).

It remains to treat the case in which every retained signature is a scalar
times a real tensor.  Divide each by its first nonzero coefficient and
call the resulting real algebraic set \(\Lambda_0\).  These known
scalars are accounted for per vertex, and \(I\in\Lambda_0\).  The
equality-wire equivalence \cref{lem:p1-equality-wire-equivalence} gives
\[
 \KHolant(\Lambda)\equivT\KHolant(\Lambda_0)
 \equivT\Holant(\mathcal F),
 \qquad \mathcal F:=\Lambda_0\cup\{X\}.
\]
Apply the real-valued dichotomy of Shao and Cai
\cite[full version, Theorems~1.2 and~2.22]{ShaoCai2020} to this whole
finite real set.  If its condition \textnormal{(T)} fails, hardness
composes along the displayed reductions.  If it holds, the following
argument supplies the required \(\Tract(K\Lambda)\) leaf in our coordinates.

The alternative \(\mathcal F\subseteq\angles{\cT}\) immediately implies
\(K\Lambda\subseteq\angles{\cT}\), since invertible one-port maps and
nonzero scalars preserve tensor products of unary and binary signatures.
In each other alternative there is one algebraic matrix
\(M\in\GL_2(\AlgNums)\), by
\cref{lem:algebraic-transform-witness}, and one
\(\cC\in\{\cP,\cA,\cL\}\) with
\[
 E:=M^{\mathsf T}M\in\cC,
 \qquad M^{-1}\mathcal F\subseteq\cC.
\]
In particular \(B:=M^{-1}XM^{-\mathsf T}\in\cC\), and
\[
 M^{\mathsf T}XM=EBE.
\]
Both \(E\) and \(B\) are symmetric nonsingular members of \(\cC\).
Therefore \cref{lem:symmetric-binary-pullback} gives
\(EBE\in\cC\) for each of the three possible classes.
In all cases the same \(U=M\) satisfies
\[
 U^{\mathsf T}XU\in\cC,
 \qquad U^{-1}\Lambda\subseteq\cC.
\]
This is a common native presentation, hence
\(\Tract(K\Lambda)\), by \cref{def:common-K-presentation}.

The real case is therefore a resolved complexity leaf.  The other case
uses one finite actual batch and strictly decreases the outer parameters
unless already resolved.  The starting deck is stable; all factors of the
new batch are processed before return, and strict successors reapply
\cref{lem:minimum-core-deck}.  Thus the queue requirements of
\cref{def:relative-stratum-protocol,def:frozen-token-coverage} hold.
There is no remaining common-presentation rigidity leaf to restrict, so
the definition of \(\Close\) gives the stated conclusion.
\end{proof}

\subsection{The exotic Klein deck}
\label{sec:exotic-v4-deck}

The exotic Klein form is not \(GO(X)\)-conjugate to the Pauli form, so we
work in its physical basis
\begin{equation}
\begin{aligned}
 E_0&=I,&
 E_1&=X,\\
 E_2&=\begin{pmatrix}i&1\\-1&-i\end{pmatrix},&
 E_3&=\begin{pmatrix}-i&1\\-1&i\end{pmatrix}.
\end{aligned}
\label{eq:exotic-v4-basis}
\end{equation}
Identify \(0,1,2,3\) with \(00,01,10,11\in V=\Ftwo^2\), define the exotic
projective Klein set \(\mathcal V_{\rm ex}:=\{[E_a]:a\in V\}\), write
addition as \(\oplus\), and define the label permutation
\(\vartheta:=(0)(1)(2\ 3)\).  A superscript that is an ordered port pair,
as in \(E_a^{13}\), specifies the tensor ports; an upper label, as in
\(E^a\), denotes the Frobenius-dual basis element defined below.

For an oriented perfect matching \(M\), an \emph{exotic
\(M\)-matching product} is a nonzero scalar multiple of
\[
 \bigotimes_{e\in M}E_{\tau_e}^e,
 \qquad \tau\in V^M.
\]
The products with these representatives form a basis of
\((\mathbb C^2)^{\otimes 2|M|}\), called the \emph{\(M\)-basis};
one product spans a \emph{matching line}.  Recoupling below means
expressing a product on one matching in the basis of another matching.
Its coefficient support records how many matching labels must vary.

\begin{lemma}[Exotic Klein calculus and recoupling]
\label{lem:exotic-v4-calculus}
For the bilinear Frobenius form,
\begin{equation}
 \langle E_a,E_b\rangle_F
 =d_a\one[b=\vartheta(a)],
 \qquad (d_0,d_1,d_2,d_3)=(2,2,4,4).
 \label{eq:exotic-v4-Gram}
\end{equation}
Thus define
\[
 E^a:=d_a^{-1}E_{\vartheta(a)},\qquad
 \langle E^a,E_b\rangle_F=\one[a=b].
\]
This is the Frobenius-dual basis, with an actual contraction up to fixed
scalar, and
\begin{equation}
 E_a^{\mathsf T}\doteq E_{\vartheta(a)},
 \qquad E_aE_b\doteq E_{a\oplus b},
 \label{eq:exotic-v4-transpose-fusion}
\end{equation}
with nonzero omitted scalars.  Relative to the matching \(12|34\),
\begin{align*}
 \delta^\times_{a,b,\gamma}
   &:=\vartheta(\gamma)\oplus b\oplus\vartheta(a),\\
 \delta^\parallel_{a,b,\gamma}
   &:=\gamma\oplus\vartheta\bigl(b\oplus\vartheta(a)\bigr),
\end{align*}
and define the recoupling coefficients by
\begin{align*}
 \rho^\times_{a,b}(\gamma)
   &:=
   \left\langle
    (E^\gamma)^{12}
    (E^{\delta^\times_{a,b,\gamma}})^{34},
    E_a^{13}E_b^{24}
   \right\rangle_F,\\
 \rho^\parallel_{a,b}(\gamma)
   &:=
   \left\langle
    (E^\gamma)^{12}
    (E^{\delta^\parallel_{a,b,\gamma}})^{34},
    E_a^{14}E_b^{23}
   \right\rangle_F .
\end{align*}
Here the pairing is the bilinear Frobenius pairing on quaternary tensors.
The recoupling identities are
\begin{align*}
 E_a^{13}E_b^{24}
 &=\sum_{\gamma\in V}\rho^{\times}_{a,b}(\gamma)
   E_\gamma^{12}E_{\delta^\times_{a,b,\gamma}}^{34},
 \\
 E_a^{14}E_b^{23}
 &=\sum_{\gamma\in V}\rho^{\parallel}_{a,b}(\gamma)
   E_\gamma^{12}E_{\delta^\parallel_{a,b,\gamma}}^{34},
\end{align*}
Every displayed coefficient lies in \(\mathbb Q(i)^\times\).
For matchings \(M,N\) of \(2r\)
ports, let \(c(M,N)\) count alternating cycles, including common edges.
An \(N\)-matching product then has exactly
\begin{equation}
 4^{r-c(M,N)}
 \label{eq:exotic-v4-overlap-size}
\end{equation}
nonzero coefficients in the \(M\)-basis.  Thus support four means one
alternating four-cycle and exactly two varying \(M\)-edges.
\end{lemma}

\begin{proof}
The displayed matrices satisfy
\[
 \begin{gathered}
 E_1^2=I,\qquad E_2^2=E_3^2=-2I,\qquad
 E_1E_2=-E_2E_1=-iE_3,\\
 E_0^{\mathsf T}=E_0,\qquad E_1^{\mathsf T}=E_1,\qquad
 E_2^{\mathsf T}=-E_3,\qquad E_3^{\mathsf T}=-E_2,\\
 \operatorname{tr}(E_a)=2\one[a=0].
 \end{gathered}
\]
These identities give the fusion and Gram formulas, including all
nonzero scalar factors.  For example, the coefficient of
\(E_\gamma^{12}E_\delta^{34}\) in the first recoupling is
\[
 \operatorname{tr}\!\left((E^\delta)^{\mathsf T}
                  E_a^{\mathsf T}E^\gamma E_b\right).
\]
Fusion makes its matrix product a nonzero multiple of
\(E_{\delta\oplus\vartheta(a)\oplus\vartheta(\gamma)\oplus b}\).
The trace therefore vanishes except at
\(\delta=\delta^\times_{a,b,\gamma}\), and is nonzero there.
For the second recoupling the coefficient is
\[
 \operatorname{tr}\!\left((E^\delta)^{\mathsf T}
                  E_b^{\mathsf T}(E^\gamma)^{\mathsf T}E_a\right),
\]
which is nonzero exactly at
\(\delta=\gamma\oplus\vartheta(b)\oplus a
=\delta^\parallel_{a,b,\gamma}\).  All factors have entries in
\(\mathbb Q(i)\), proving the coefficient assertion without enumerating
the labels.

For general \(M,N\), contracting an \(N\)-product against a dual
\(M\)-basis element separates over the alternating cycles of
\(M\cup N\).  On a cycle of length \(2\ell\), the coefficient is
one trace of a product of the \(\ell\) fixed \(N\)-matrices and
\(\ell\) dual \(M\)-matrices, transposed according to their
orientations.  Each \(M\)-label enters its fused label once, through
either the identity or the bijection \(\vartheta\).  After
\(\ell-1\) labels are chosen, the zero-label trace condition determines
the last label uniquely.  The resulting trace is a nonzero scalar times
\(\operatorname{tr}(I)=2\).  Thus the cycle contributes exactly
\(4^{\ell-1}\) nonzero coefficients, including one for a common edge
when \(\ell=1\).  Multiplying over cycles gives
\(4^{\sum(\ell-1)}=4^{r-c(M,N)}\).  If this equals four, exactly one
cycle has length four and the others are common edges, proving the last
assertion.
\end{proof}

\paragraph{Exotic-card convention.}
For an ordered pair \(p\), define the exotic pair card
\[
 C_p^a(f):=\partial_p^{E^a}f.
\]
For disjoint ordered pairs \(p_1,\ldots,p_s\), an \emph{exotic card} is
the corresponding iterated contraction by fixed actual representatives of
the duals \(E^{a_1},\ldots,E^{a_s}\); it is \emph{proper} when at least two
ports remain.  Disjoint contractions commute, and recorded nonzero dual
scalars are divided out in the canonical notation.
With the fixed representatives, write
\[
 E_aE_b=\omega(a,b)E_{a\oplus b},
 \qquad \omega(a,b)\in\mathbb Q(i)^*;
\]
the nonzero factors \(\omega(a,b)\) and the corresponding transpose
factors are recorded.  To
\emph{dual-cross} two residual matching edges \(e=\{u,u'\}\) and
\(e'=\{v,v'\}\), choose one endpoint \(u\) of \(e\), one endpoint \(v\) of
\(e'\), and one fixed \(a\in V\); contract the single pair \(u,v\) by the
fixed actual representative of \(E^a\), and leave \(u',v'\) as the endpoints
of the fused residual edge.  The endpoint orientations are part of the
choice, and all \(\omega\)- and transpose scalars are recorded.  A
\emph{scale or cancellation locus} is a set
of nonzero-card scalar parameters on which a specified card becomes
proportional to a matching line or vanishes.

The duals are actual up to the recorded scalars.  Tomography
\begin{equation}
 f=\sum_{a\in V}E_a^{p}\otimes C_p^a,
 \qquad C_p^a=\langle E^a,f\rangle_p,
 \label{eq:exotic-v4-tomography}
\end{equation}
is coordinatewise only; neither its sum nor a matching basis is a gadget.

\begin{lemma}[The exotic-Klein six-port base]
\label{lem:exotic-v4-q6}
If every exotic pair card of a nonzero six-port tensor \(h\) is zero or a
product of two exotic binaries, then \(h\) is a product of three exotic
binaries.  Thus there is no exotic analogue of \(\Hcore\).
\end{lemma}

\begin{proof}
Normalize a nonzero card to \(C_{12}^0=E_0^{34}E_0^{56}\) by actual dressings
and a permutation.  The six-port interface
\cref{thm:cert-interface-exotic-v4-q6} covers every nonzero/zero,
scale, and cancellation locus; all survivors are three-binary products.
Undo the normalization.
\end{proof}

\begin{lemma}[Fixed ruling for exotic cards]
\label{lem:exotic-v4-fixed-ruling}
Let \(f\ne0\) have even arity at least eight, and suppose every proper exotic
card of \(f\) is zero or an exotic matching product.  For every deleted
pair \(p\), all nonzero cards \(C_p^a\) in
\cref{eq:exotic-v4-tomography} use one residual perfect matching.
\end{lemma}

\begin{proof}
Fix a deleted pair \(p\), and suppose that two nonzero cards
\(C_p^a\) and \(C_p^b\) use distinct residual matchings \(M\ne N\).  We
construct a sequence of actual exotic contractions that preserves both
cards as nonzero matching products and preserves one residual disagreement.

The only local operation needed is the following.  If two contracted ports
belong to distinct matching edges in a given distinguished card, then, up to the fixed
transpose convention, the new edge factor is
\[
 E_\alpha E_\kappa E_\beta\doteq
 E_{\alpha\oplus\kappa\oplus\beta}.
\]
It is nonzero because every \(E_\gamma\) is nonsingular.  Choose two
vertices at alternating distance two on an alternating cycle of
\(M\cup N\).  They belong to distinct edges in both distinguished cards, so contracting
them shortens that cycle by two vertices while keeping both distinguished
cards nonzero.  Repeating reduces a marked nontrivial cycle to length four.

Reduce every other nontrivial alternating cycle and splice it into the
marked cycle by one cross-contraction; the same three-matrix identity keeps
both distinguished cards nonzero and merges the two cycles.  A common edge
carrying the same label in both cards may be closed with its actual
Frobenius dual.  If
a common edge carries different labels, splice it into the marked cycle
instead of closing it, and reduce the enlarged cycle again.  After all
common edges and all other cycles have been removed, the resulting actual
six-port tensor \(h\) has two nonzero \(p\)-cards with different residual
two-edge matchings.

Every exotic pair card of \(h\) is zero or a matching product.  Indeed,
bilinear contractions commute, so any card of \(h\) may be reordered so
that its final pair contraction is applied first to \(f\).  By hypothesis
that first card is zero or a matching product, and every subsequent
contraction preserves this property by
\(E_\alpha E_\kappa E_\beta\doteq
E_{\alpha\oplus\kappa\oplus\beta}\).  Thus
\cref{lem:exotic-v4-q6} applies and makes \(h\) a product of three exotic
binaries.

For a product of three binaries, all nonzero cards at a fixed deleted pair
have one residual matching: if the deleted pair is a factor edge, only its
contraction in the dual direction is nonzero; otherwise fusion of the two incident factor edges
leaves the same residual matching for every kernel giving a nonzero card.  This contradicts
the two preserved rulings of \(h\).  Hence all nonzero \(p\)-cards of \(f\)
use one residual perfect matching.
\end{proof}

\begin{lemma}[Fixed-ruling fiber dichotomy]
\label{lem:exotic-v4-fiber-dichotomy}
Let \(f\) satisfy \cref{lem:exotic-v4-fixed-ruling}, fix \(p\) and its
residual matching \(M\), and suppose \(f\) has no nontrivial
positive-arity tensor factor.  Then
all four cards \(C_p^a\) are nonzero and have the form
\begin{equation*}
 C_p^a=\lambda_a\bigotimes_{e\in M}E_{\tau_e(a)}^e,
\end{equation*}
where \(\lambda_a\in\mathbb C^\times\) are nonzero scalar coefficients and
\(\tau_e:V\to V\) are label maps; every \(\tau_e\) is a permutation.
\end{lemma}

\begin{proof}
Let
\[
 \Gamma=\{a\in V:C_p^a\ne0\}.
\]
Tomography and the fixed ruling give
\[
 f=\sum_{a\in\Gamma}\lambda_a E_a^p
      \bigotimes_{e\in M}E_{\tau_e(a)}^e,
 \qquad \lambda_a\ne0.
\]
For an edge \(e\in M\), contraction with the actual Frobenius dual of
\(E_b\) selects the fibre
\[
 \Gamma_{e,b}=\{a\in\Gamma:\tau_e(a)=b\}.
\]
In the displayed \(p\cup(M\setminus\{e\})\)-basis, the resulting card has
exactly \(|\Gamma_{e,b}|\) nonzero coefficients, since distinct \(a\)'s
remain distinct on the \(p\)-coordinate.  By
\cref{eq:exotic-v4-overlap-size}, every nonempty fibre therefore has size
one or four.

If \(|\Gamma|=1\), the factor \(E_a^p\) is common to the whole tensor,
contrary to the no-factor hypothesis.  If \(|\Gamma|\in\{2,3\}\), every
\(\tau_e\) is injective on \(\Gamma\).  Choose ports from two distinct
edges of \(M\) and cross-contract them by any actual exotic kernel.
Projective fusion is nonzero in every tomographic summand, and the \(p\)-labels remain
distinct, so the resulting proper card has exactly \(|\Gamma|\)
coefficients in the induced matching basis.  This is impossible because
\cref{eq:exotic-v4-overlap-size} permits only powers of four.  Consequently
\(\Gamma=V\).

Now every nonempty \(\tau_e\)-fibre has size one or four.  A four-point
fibre makes \(\tau_e\) constant and exposes \(E_b^e\) as a genuine edge
factor.  Since no factor exists, all fibres are singletons.  Hence every
\(\tau_e:V\to V\) is a permutation, proving the stated form.
\end{proof}

\begin{lemma}[The exotic-Klein eight-port base]
\label{lem:exotic-v4-q8}
Let \(h\ne0\) have arity eight and suppose every exotic pair card is zero
or a product of three exotic binaries.  Then \(h\) has a genuine tensor
factor.
\end{lemma}

\begin{proof}
Contraction closure under \cref{eq:exotic-v4-transpose-fusion} makes every
iterated proper exotic card zero or a matching product.  For a nonfactor,
\cref{lem:exotic-v4-fixed-ruling,%
lem:exotic-v4-fiber-dichotomy} gives four cards on one matching.  Normalize
one actually to \(C_{12}^0=E_0^{34}E_0^{56}E_0^{78}\).  Each residual
 label map is affine because the convention in
 \cref{subsubsec:boolean-affine-tools} identifies
 \(\operatorname{AGL}(2,2)\) with \(\operatorname{Sym}(V)\).  Actual
edge dressings remove translations and absorb fixed nonzero cocycles:
\begin{equation}
 h=\sum_{\gamma\in V}\lambda_\gamma
 E_\gamma^{12}E_{A_1\gamma}^{34}
 E_{A_2\gamma}^{56}E_{A_3\gamma}^{78},
 \qquad
 A_j\in\GL_2(\Ftwo),\quad \lambda_\gamma\ne0.
 \label{eq:exotic-v4-q8-linear-normal-form}
\end{equation}
Contract \((3,5)\) or \((4,5)\) by \(E^0=I/2\).  On the two remaining
endpoints of the affected edges, the respective matrices are
\[
 \tfrac12 E_{A_1\gamma}^{\mathsf T}E_{A_2\gamma},
 \qquad \tfrac12 E_{A_1\gamma}E_{A_2\gamma}.
\]
By \cref{eq:exotic-v4-transpose-fusion}, their labels are
\[
 D_1\gamma=(\vartheta A_1+A_2)\gamma,
 \qquad D_2\gamma=(A_1+A_2)\gamma,
\]
with all sums over \(\Ftwo\) and all omitted scalars nonzero.
The two maps cannot both vanish, since that would imply
\(\vartheta A_1=A_1\), contrary to the invertibility of \(A_1\) and
\(\vartheta\ne I\).  Choose a nonzero one, denoted by \(D\).
In the resulting six-port matching basis the four labels are
\( (\gamma,D\gamma,A_3\gamma)\).
They are distinct because their first coordinates are distinct, and every
coefficient is nonzero.  All three coordinates vary: the first and third
maps are bijective, and \(D\ne0\).  This contradicts
\cref{eq:exotic-v4-overlap-size}, which allows only two varying edges in
a four-term matching product.  Thus \(h\) has a genuine proper tensor factor.
\end{proof}

\begin{proposition}[High-arity exotic factor]
\label{prop:exotic-v4-high-arity-factor}
Let \(f\ne0\) have arity \(2r\ge10\), and suppose every proper exotic card is
zero or an exotic matching product.  Then \(f\) has a genuine proper
tensor factor.
\end{proposition}

\begin{proof}
Assume no factor.  At a fixed pair, \cref{lem:exotic-v4-fixed-ruling,%
lem:exotic-v4-fiber-dichotomy} gives four nonzero cards on one matching with
permutation label maps.  Choose two residual edges and perform one
dual-cross contraction as defined above.  Every tomographic summand remains nonzero
because the three matrices in its endpoint-to-endpoint fusion are
nonsingular.  The four resulting terms still vary at the fixed pair and at
all \(r-3\) untouched residual edges, hence in at least \(r-2\ge3\)
coordinates.  But
\cref{eq:exotic-v4-overlap-size} allows only two varying coordinates in a
four-term matching product.  This contradiction proves that a genuine
proper factor exists.
\end{proof}

\begin{theorem}[Exotic-Klein relative rigidity]
\label{thm:exotic-v4-relative-rigidity}
Assume
\[
 \Stable(\Lambda,\mathcal V_{\rm ex},f),\qquad
 \arity(f)=\nu(\Lambda).
\]
Then
\[
 \Close_{\{\cA\}}(\Lambda,\mathcal V_{\rm ex},f).
\]
More explicitly, every continuing branch reaches an established resolved
leaf or a strict outer deck successor after processing the actual batch
from the quaternary interface by \cref{lem:matching-synthesis-measure}.
No rigidity leaf remains at the current node.
\end{theorem}

\begin{proof}
First dispatch odd-arity factors, including unaries, by
\cref{thm:external-odd}, and unsafe binaries by
\cref{lem:binary-interpolation,lem:actual-binary-completeness}.
On a continuing branch all retained primes have even arity.
Deck stability makes every exotic pair card zero or an exotic matching
product.  Repeated transpose/fusion using
\cref{eq:exotic-v4-transpose-fusion} therefore makes every iterated proper
exotic card zero or an exotic matching product.
At arity four, item~4 of \cref{thm:equality-q4-boundary} invokes the endpoint-nondegenerate eight-vertex signature
interface.  By \cref{cor:p1-eight-vertex-reduced-interface}, this
arity-four branch is a resolved leaf or a strict outer successor, so it
contributes no surviving product or local-affine rigidity leaf.  At arity six \cref{lem:exotic-v4-q6} contradicts primality;
at arity eight and at least ten, the factor conclusions of
\cref{lem:exotic-v4-q8,prop:exotic-v4-high-arity-factor} contradict
primality as well.  These cases exhaust the possible arities of the
distinguished nonbinary prime \(f\).

Process any concrete batch from the quaternary interface, including its
full factorization, by \cref{def:relative-stratum-protocol,%
def:protocol-discharged,lem:matching-synthesis-measure}.
The interface supplies a resolved leaf or a strict outer successor, while
the higher-arity cases leave no stable tuple.  Consequently there is no
surviving rigidity leaf at this node.  This gives
\(\Close_{\{\cA\}}(\Lambda,\mathcal V_{\rm ex},f)\).
\end{proof}

\begin{proof}[Proof of \cref{thm:v4-relative}]
If \(G=V_4\), apply \cref{thm:standard-v4-relative}.  If
\(G=\mathcal V_{\rm ex}\), apply
\cref{thm:exotic-v4-relative-rigidity}.  These are exactly the two
Klein-four forms in the statement.
\end{proof}
\section{Platonic Matching Decks}
\label{sec:platonic-decks}

The finite groups in \cref{thm:marked-finite-group-routing} leave three
Platonic cases: tetrahedral, octahedral, and icosahedral.  We prove
\cref{thm:platonic-decks} by reducing an arbitrary-arity core to a six-port
tensor.  A pair contraction is a linear map from the four-dimensional
matrix space to residual tensors.  Incidence among the group matrices
then forces either a common tensor factor or three collinear cards with
different matchings.  In the latter case four residual ports contain all
the variation, and the remaining factors can be closed by actual kernels.
The resulting finite classifications are proved in
Appendix~\ref{appsec:cert-platonic}.

We use the quaternion coordinates of \cref{eq:quaternion-basis}, the
marked-group classification of \cref{thm:marked-finite-group-routing}, and
the projective and tensor conventions of
\cref{subsubsec:group-projective,subsubsec:projective-tensor-geometry}.
Thus \([F]\) is the line spanned by any nonzero matrix or tensor \(F\),
including a singular matrix, and \(\mathbb P(0)=\varnothing\).

The three nonstandard octahedral displays are as follows.  For the
transposition--transposition form put
\begin{equation}
 H_{\mathcal O}:=
 \begin{pmatrix}\sqrt2&-i\sqrt2\\1-i&1+i\end{pmatrix},
 \qquad
 \mathcal O':=H_{\mathcal O}^{-1}\mathcal O H_{\mathcal O}.
 \label{eq:octahedral-second-form}
\end{equation}
This realizes the marked type \(\mathcal O'\) of
\cref{thm:marked-finite-group-routing} in coefficient coordinates.

For the two mixed fixed-equality forms, put
\begin{equation}
 \rho:=\sqrt2-1,\qquad
 H_{dt}:=Q_0-\rho Q_1,\qquad
 H_{td}:=Q_0+\rho Q_3,
 \label{eq:octahedral-mixed-conjugators}
\end{equation}
and, for \(\alpha\in\{dt,td\}\), define
\begin{equation}
 \theta_\alpha(U):=H_\alpha^{-1}UH_\alpha,\qquad
 \mathcal O_\alpha:=\theta_\alpha(\mathcal O).
 \label{eq:octahedral-mixed-forms}
\end{equation}
Direct quaternion multiplication gives
\begin{align}
 H_{dt}H_{dt}^{\mathsf T}&=2\rho(Q_0-Q_1),&
 \theta_{dt}(Q_1)&=Q_1,&
 \theta_{dt}(Q_2-Q_3)&=-\sqrt2Q_3,\notag\\
 H_{td}H_{td}^{\mathsf T}&=2\rho(Q_0+Q_3),&
 \theta_{td}(Q_3)&=Q_3,&
 \theta_{td}(Q_1-Q_2)&=\sqrt2Q_1.
 \label{eq:octahedral-mixed-mark-check}
\end{align}
Thus \(\mathcal O_{dt}\) has marked class pattern \((dt;t)\), whereas
\(\mathcal O_{td}\) has pattern \((td;t)\).  All three matrices \(H_{\mathcal O},H_{dt},H_{td}\) are used as coefficient
similarities.  Actual contractions and dressings use representatives of
the current group; the retained signature set and its literal equality
remain in the fixed coordinates.

Throughout this section, for a finite projective group \(G\) put
\[
 G^\#:=\{[T^\#]:[T]\in G\},
 \qquad T^\#:=XT^{\mathsf T}X,
\]
 and call \(G\) \emph{reversal-stable} when \(G^\#=G\).  A \emph{marked
 Platonic state} means a retained, factor-saturated state containing the
 distinguished literal \(I\), whose
recomputed complete group is one of
\[
 \mathcal T,\ \mathcal T_{\rm ext},\ \mathcal O,\ \mathcal O',\
 \mathcal O_{dt},\ \mathcal O_{td},\
 \mathcal I,\ \mathcal I_-,
\]
with one fixed actual representative and realization record for every
projective group class.  All statements about group matrices distinguish a
matrix representative \(K\) from its projective class \([K]\).

\begin{lemma}[The two marked tetrahedral forms]
\label{lem:marked-a4-forms}
Let \(G\leq\PGL_2(\mathbb C)\) be a finite reversal-stable complete group in the fixed coordinates containing literal \(I\), and define the physical transfer
mark \(x\), the reversal mark \(j\), and the extension group \(H\) by
\[
  x=[X]\in G,\qquad j=[Z],\qquad H=\langle G,j\rangle.
\]
If \(G\cong A_4\), then, up to one legal global fixed-\(I\)-preserving
\(GO(X)\)-normalization, exactly one of the following
holds.
\begin{enumerate}
\item \(j\in G\), \(H=G\), and \(G=\mathcal T\), where
      \(\mathcal T\) is the standard group in
      \cref{eq:tetrahedral-domain}.
\item \(j\notin G\), \(H\cong S_4\), and
\begin{equation}
\begin{aligned}
 \mathcal T_{\rm ext}
 ={}&\{[Q_0],[Q_1],[Q_2+Q_3],[Q_2-Q_3]\}\\
 &{}\cup
 \{[Q_0+\epsilon Q_1+\delta\sqrt2Q_2],
    [Q_0+\epsilon Q_1+\delta\sqrt2Q_3]:
       \epsilon,\delta\in\{1,-1\}\}.
\end{aligned}
\label{eq:external-a4-domain}
\end{equation}
Here \(x=[Q_1]\in\mathcal T_{\rm ext}\), whereas
\(j=[Q_3]\notin\mathcal T_{\rm ext}\) and
\(xj=[Q_2]\notin\mathcal T_{\rm ext}\).
\end{enumerate}
\end{lemma}

\begin{proof}
Reversal gives \(G\triangleleft H\) and \([H:G]\le2\).  If \(j\notin G\),
then \(H\cong S_4\), with the commuting pair \((x,j)\) marked as a double
transposition and a transposition; if \(j\in G\), the three involutions form
the unique normal Klein four of \(A_4\) by
\cref{lem:finite-pgl2-facts}.  After the corresponding
finite-subgroup embedding the remaining ambiguity is
\[
 C_{\PGL_2(\mathbb C)}(\langle x,j\rangle)=\langle x,j\rangle,
\]
represented inside \(GO(X)\).  Direct multiplication gives the two stated
domains.
\end{proof}

For algebraic comparison only, define the scalar \(s_{\rm ext}\), the matrix
\(R_{\rm ext}\),
and the coefficient-similarity map \(\theta_{\rm ext}\) by
\begin{equation*}
 s_{\rm ext}:=\sqrt2-1,\qquad R_{\rm ext}:=Q_0+s_{\rm ext}Q_1,
 \qquad \theta_{\rm ext}(U):=R_{\rm ext}^{-1}UR_{\rm ext}.
\end{equation*}
On projective sets, apply \(\theta_{\rm ext}\) representative-wise.  Then
\begin{equation*}
 \theta_{\rm ext}(\mathcal T)=\mathcal T_{\rm ext},
 \qquad
 [R_{\rm ext}^{\mathsf T}XR_{\rm ext}]=[Q_0-Q_1]\ne[X].
\end{equation*}
Thus \(R_{\rm ext}\notin GO(X)\): \(\theta_{\rm ext}\) transports coefficients only;
all gadgets and normalizations use actual external representatives.

Write \(\mathcal I_-\) for the opposite icosahedral coordinate display.  The
two displays have the same abstract marked classification, but both are
retained because their physical representatives are different.  Put
\begin{equation}
  D_{\mathcal I}:=Q_0+Q_3,
  \qquad
  D_{\mathcal I}^{\mathsf T}XD_{\mathcal I}=2X,
  \qquad
  \mathcal I_-=D_{\mathcal I}\mathcal I D_{\mathcal I}^{-1}.
  \label{eq:icosahedral-chirality}
\end{equation}
The signed coordinate permutations defining \(\mathcal T\) also give
\[
 D_{\mathcal I}\mathcal T D_{\mathcal I}^{-1}=\mathcal T,
 \qquad
 \mathcal T\subseteq\mathcal I\cap\mathcal I_-.
\]
Indeed, conjugation by \(D_{\mathcal I}\) sends
\((Q_1,Q_2,Q_3)\) to \((-Q_2,Q_1,Q_3)\).  We use this identity only as the
coefficient dictionary below; no retained signature, and in particular no
equality signature, is transformed.

\begin{lemma}[Fixed-equality dictionary for the opposite icosahedral display]
\label{lem:icosahedral-fixed-I-dictionary}
Set \(D=D_{\mathcal I}\) and
\[
 \kappa_D(K):=D^{-\mathsf T}KD^{-1}.
\]
Let \((\Lambda,\mathcal I_-,f)\) be a marked Platonic state in the fixed
coordinates containing the supplied literal \(I\), and introduce only for
coefficient analysis the tensor
\(\widehat f:=D^{-\otimes\arity(f)}f\).  For every ordered pair \(p\),
\begin{equation}
 \Phi_p^f\!\left(\kappa_D(K)\right)
 =D^{\otimes(\arity(f)-2)}\Phi_p^{\widehat f}(K).
 \label{eq:icosahedral-card-covariance}
\end{equation}
The map \([K]\mapsto[\kappa_D(K)]\) bijects the standard
\(\mathcal I\)-kernel classes with the actual \(\mathcal I_-\)-kernel
classes.  Consequently every finite card or network identity proved in the
standard display has a direct identity in the opposite display, obtained by
replacing each kernel by its \(\kappa_D\)-image and recording the nonzero
representative and edge scalars.  This dictionary preserves zero cards,
tensor factors, all flattening ranks, matching labels, endpoint support,
and the property that an output binary transfer lies outside the complete
group.  It does not change \(\Lambda\), the active occurrence, or the
distinguished literal \(I\).
\end{lemma}

\begin{proof}
Direct contraction gives \cref{eq:icosahedral-card-covariance}.  Moreover
\(D^{-\mathsf T}=\tfrac12XDX\), and hence
\[
 [\kappa_D(K)]
 =\bigl[X\,D(XK)D^{-1}\bigr]\in\mathcal I_-
 \qquad([K]\in\mathcal I).
\]
Here \([X],[XK]\in\mathcal I\), conjugation sends the latter class into
\(\mathcal I_-\), and \([X]\in\mathcal T\subseteq\mathcal I_-\).  The
inverse formula proves bijectivity.  Since \(D^{\mathsf T}XD=2X\), applying
the covariance successively at the vertices of a finite network cancels the
internal \(D,D^{-1}\) factors across every native edge, leaving only a known
nonzero scalar.  A residual binary \(B\) is carried to
\(DBD^{\mathsf T}\), whose transfer satisfies
\[
 [DBD^{\mathsf T}X]=[D(BX)D^{-1}].
\]
Thus nonsingularity and exclusion from the current group are preserved.
Finally, \(D\) is diagonal and invertible, so it preserves support (hence
endpoint nondegeneracy), tensor products, matching labels, and all
flattening ranks.  All resulting kernels are replaced by the already fixed
actual \(\mathcal I_-\) representatives in their projective classes, with
their nonzero realization scalars recorded.  At no point is the analytical
tensor \(\widehat f\) adjoined or realized.
\end{proof}

\begin{theorem}[Platonic deck rigidity]
\label{thm:platonic-decks}
Assume
\[
 \Stable(\Lambda,G,f),\qquad
 G\in\{\mathcal T,\mathcal T_{\rm ext},\mathcal O,\mathcal O',
         \mathcal O_{dt},\mathcal O_{td},
         \mathcal I,\mathcal I_-\}.
\]
Then
\[
 \begin{cases}
  \Close_{\{\cA\}}(\Lambda,G,f),&G=\mathcal T,\\
  \Close_{\varnothing}(\Lambda,G,f),
    &G\in\{\mathcal T_{\rm ext},\mathcal O,\mathcal O',
            \mathcal O_{dt},\mathcal O_{td},
            \mathcal I,\mathcal I_-\}.
 \end{cases}
\]
\end{theorem}

\subsection{Common card-map interface and quaternary boundary}

At a continuing fixed-\(I\) deck-stable node, fix for each \(g\in G\)
a matrix \(K_g\) with \([K_g]=g\).  Since \([X]\in G\), choose a fixed
actual binary representative \(B_g\) of transfer class \([XK_g]\).
The convention \(T_{B_g}=B_gX\) gives \(B_g\doteq XK_gX\).
Attach its two ordered ports to the selected pair through two native
\(X\)-edges; the effective bilinear kernel is then
\(XB_gX\doteq K_g\).  Store its nonzero realization
scalar and ordered-port record, and define the actual representative set
\begin{equation}
  \mathcal K_G^{\rm act}:=\{K_g:g\in G\}.
  \label{eq:platonic-actual-kernel-lift}
\end{equation}
Every contraction below uses these fixed representatives.  A resulting
binary \(B\) is tested through its transfer \(BX\).  A zero card contributes
zero; a nonzero singular \(B\) is processed by factor saturation.  If
\(B\) is nonsingular and \([BX]\notin G\), then
\cref{lem:binary-interpolation,lem:actual-binary-completeness} gives an
infinite-order terminal or a strict enlargement of the complete group.
The ordered ports and nonzero realization scalar are retained in each
case.  These are bounded local operations on the selected card.

For a signature \(u\) and an ordered pair \(p=(a,b)\) of its ports, define
the physical card map
\begin{equation}
  \Phi_p^u:\Mat_2(\mathbb C)\longrightarrow
  {(\mathbb C^2)}^{\otimes(\arity(u)-2)},
  \qquad
  \Phi_p^u(K)(y)=\sum_{s,t\in\{0,1\}}u(s,t,y)K(s,t).
  \label{eq:platonic-card-map}
\end{equation}
When the signature is the active \(f\), write \(\Phi_p=\Phi_p^f\), and write
\(\langle K^{(p)},u\rangle:=\Phi_p^u(K)\) for this physical bilinear
(not Hermitian) contraction; equivalently
\(\Phi_p^u(K)(y)=\operatorname{tr}(A_y^{\mathsf T}K)\), where
\(A_y:=(u(s,t,y))_{s,t}\in\Mat_2(\mathbb C)\) is the \(p\)-slice matrix.
Port permutations
supply the other orderings.

For every two-element port set \(\{a,b\}\), use the canonical orientation
\(p=(a,b)\) with \(a<b\), and order complementary ports increasingly unless
both orientations are explicitly requested.  Thus a quaternary has six
canonical pairs and a six-port tensor has fifteen; all other orientations
are obtained by actual port permutations.

The rank of \(\Phi_p^h\) is the flattening rank across the selected pair
and the remaining ports.  Since its domain is \(\Mat_2(\mathbb C)\), it is
at most four.  We treat ranks below four and rank four separately.  All signatures used as gadgets are directly realized
over \(\Lambda\), with ordered ports, nonzero scalar, and predecessor
reduction recorded.  The required condition is
\begin{equation}
 \Phi_p^h(K_g)=0\quad\text{or}\quad\Phi_p^h(K_g)\text{ is \(G\)-matching},
 \qquad \text{for every canonical \(p\) and \(K_g\in\mathcal K_G^{\rm act}\)}.
 \label{eq:platonic-all-card-condition}
\end{equation}

Put
\(\mathcal D_G:=\{[K]:K\in\mathcal K_G^{\rm act}\}\).  A \emph{rich
domain line} is a projective line containing at least three points of
\(\mathcal D_G\).  Let \(\mathscr G_G^{\rm dom}\) be the graph on
\(\mathcal D_G\) joining two points when they lie on a common rich domain
line.  Such a line is \emph{kernel-avoiding} for a linear map \(\Phi\) when
it is disjoint from \(\mathbb P(\ker\Phi)\).  

For an ancestor of arity at least six, localization retains \(p\) and four
residual ports.  The rank of the ancestor map is used only to choose a
kernel-avoiding rich domain line and to bound the resulting six-port rank.  After
that signature is directly realized, its own card-map rank is
recomputed before either the rank-deficient or full-rank finite reduction is
invoked.

\paragraph{Quaternion coordinates and quaternary matching products.}

Recall the quaternion basis from \cref{eq:quaternion-basis}.  Its
Frobenius Gram matrix is
\begin{equation*}
  {\bigl(\operatorname{tr}(Q_\mu^{\mathsf T}Q_\nu)\bigr)}_{\mu,\nu}
  =2\diag(1,-1,1,-1),
\end{equation*}
so the bilinear dual basis is
\begin{equation}
  Q^0=Q_0/2,\qquad Q^1=-Q_1/2,\qquad
  Q^2=Q_2/2,\qquad Q^3=-Q_3/2.
  \label{eq:platonic-dual-basis}
\end{equation}
These dual matrices recover the coefficients of the original tensor:
their role is linear reconstruction.  For \(K\in\Mat_2(\mathbb C)\), let
\[
 c_Q(K)=(c_0,c_1,c_2,c_3)^{\mathsf T}
 \quad\Longleftrightarrow\quad
 K=c_Q(K)\mathbin{\cdot}Q:=\sum_{\mu=0}^3c_\mu Q_\mu.
\]
Write \(U_{ab}\) for the binary tensor with entries \(U(x_a,x_b)\) on
the ordered pair \((a,b)\).

Frobenius-dual reconstruction gives, for every signature \(f\) and
canonical pair \(p\),
\begin{equation}
  f=\sum_{\mu=0}^3 Q_p^\mu\otimes\Phi_p^f(Q_\mu).
  \label{eq:platonic-frobenius-reconstruction}
\end{equation}
More generally, for an ordered pair \(p\), a finite-dimensional tensor
space \(W\), and a linear map \(\Phi:\Mat_2\to W\), write
\(f_{\Phi,p}:=\sum_{\mu=0}^3Q_p^\mu\otimes\Phi(Q_\mu)\).  Every assertion below
that an abstract card map factors refers to this reconstructed tensor.

On four ordered residual ports define the three matching products
\begin{align*}
 P_0(U,V)&=U(x_1,x_2)V(x_3,x_4),\\
 P_1(U,V)&=U(x_1,x_3)V(x_2,x_4),\\
 P_2(U,V)&=U(x_1,x_4)V(x_2,x_3).
\end{align*}
For a finite projective transfer group \(G\), define its quaternary projective
matching deck by
\begin{equation*}
  \mathfrak D_4(G)=
  \{[P_r(U,V)]:r\in\{0,1,2\},\ [UX],[VX]\in G\}.
\end{equation*}
Define the nonzero scalar cone
\(\operatorname{cone}(\mathfrak D):=\{\lambda F:\lambda\ne0,
[F]\in\mathfrak D\}\).

A \emph{rich deck line} is a projective line containing at least three
points of \(\mathfrak D_4(G)\).  It is an \(s\)-line when it contains
exactly \(s\) deck points.  Its \emph{matching-label type} is the partition
of \(s\) formed by the multiplicities of its residual matching labels.
Type \((s)\) is \emph{monochromatic}, type \((1,1,1)\) is a
\emph{rainbow three-line}, and every other nonmonochromatic type is
\emph{mixed}.  A rich deck line is \emph{maximal} under the analogous
deck-point containment order.  Every deck point has rank triple
\((1,4,4)\), up to the fixed
order \((12\mid34,13\mid24,14\mid23)\).  For every matching product \(F\)
with matching \(M\), and every port subset \(S\),
\begin{equation}
  \rank F_{S\mid\bar S}=2^{c_M(S)},
  \label{eq:matching-schmidt-rank}
\end{equation}
where \(c_M(S)\) counts edges crossing the cut.

For \(L\in\mathcal K_G^{\rm act}\), the binary \(B_L=XL^{\mathsf T}\) has
\([B_LX]=[XL^{\mathsf T}X]=[L^\#]\in G\), and one native edge satisfies
\(\sum_uX(x,u)B_L(u,y)=L(y,x)\).  Hence, for
\(\mathbf L=(L_1,L_2,L_3,L_4)\in(\mathcal K_G^{\rm act})^4\),
\[
  (\mathscr D_{\mathbf L}q)(y)
  =\sum_xq(x)\prod_{j=1}^4L_j(y_j,x_j)
\]
is a constant-size actual dressing with recorded nonzero scalar.

\begin{lemma}[Eight-vertex signatures with rank triples \((2,2,2)\) or \((3,3,3)\)]
\label{lem:p1-platonic-eight-vertex-rank-terminal}
Let \(\Lambda\) be a retained, factor-saturated finite algebraic signature
set containing the distinguished literal equality \(I\).  If \(\Lambda\) directly
realizes an endpoint-nondegenerate eight-vertex signature \(M\) of
flattening ranks \((2,2,2)\) or \((3,3,3)\), then
\[
  \KHolant(\Lambda)\text{ is }\SharpP\text{-hard}
  \qquad\text{or}\qquad
  \Tract(K\Lambda)\text{ holds}.
\]
\end{lemma}

\begin{proof}
Apply \cref{cor:p1-eight-vertex-reduced-interface} to \(M\).  Its first two
outcomes are the asserted conclusions.  In its sole continuing outcome,
\(M\) itself has, up to port order and a nonzero scalar, a factorization
\(E\otimes E'\) into two nonsingular binary signatures.  Its matching
flattening ranks are therefore \((1,4,4)\), up to permutation, by
\cref{eq:matching-schmidt-rank}, contradicting either assumed rank triple.
The full-even two-copy equality outputs are already resolved inside
\cref{cor:p1-affine-eight-vertex-operational} and are not continuing leaves.  Hence only
the hard or tractable outcomes remain.
\end{proof}

\begin{lemma}[The Platonic quaternary boundary]
\label{lem:platonic-quaternary-boundary}
Let \((\Lambda,G,q)\) be a deck-stable retained, factor-saturated fixed-\(I\)
state with nonzero tensor-prime quaternary core \(q\).
Let \(G\) be the internal \(A_4\) form \(\mathcal T\), any of the four
marked \(S_4\) forms
\(\mathcal O,\mathcal O',\mathcal O_{dt},\mathcal O_{td}\), or a marked
\(A_5\) display, and
define the transfer-valued pair map
\begin{equation*}
  L_p(K)=\Phi_p^q(K)X\in\Mat_2(\mathbb C).
\end{equation*}
If every fixed actual representative \(K_g\) satisfies
\begin{equation*}
  L_p(K_g)=0\quad\hbox{or}\quad [L_p(K_g)]\in G
  \qquad\text{for all six canonical pairs }p,
\end{equation*}
then internal \(A_4\) and either \(A_5\) display have no tensor-prime
survivor.  In the first \(S_4\) form, there are
\(\lambda\in\mathbb C^\times\) and actual matrices \(U_i\) with
\([U_i]\in\mathcal O\) such that
\begin{equation*}
  q=\lambda\,(U_1\otimes U_2\otimes U_3\otimes U_4)\Eq_4.
\end{equation*}
In the transposition--transposition form, four actual
\(\mathcal O'\)-matrices dress \(q\) to
\begin{equation}
  M_{\mathcal O'}(x)
  =\one[\wt(x)\equiv0\pmod2]
    (-1)^{\sum_{i<j}x_ix_j},
  \label{eq:p1-platonic-q4-Oprime-eight-vertex}
\end{equation}
of ranks \((2,2,2)\).  In the two mixed forms the normalized rank-two
survivor is already an endpoint-nondegenerate eight-vertex signature, with
all odd entries zero and, in the even-word order
\((0000,0011,0101,0110,1001,1010,1100,1111)\), coefficient vectors
\begin{equation}
 \mathbf e(M_{dt})=(1,1,1,1,1,1,1,1),\qquad
 \mathbf e(M_{td})=(1,i,i,i,1,1,1,i).
 \label{eq:p1-platonic-q4-mixed-eight-vertex}
\end{equation}
Both have ranks \((2,2,2)\).  Hence arity four always factors or reaches the
unsafe-binary, pure generalized equality, or endpoint-nondegenerate
eight-vertex signature terminal and never recurses.
\end{lemma}

\begin{proof}
Let \(A_p\in\Mat_4(\mathbb C)\) represent \(L_p\) in the quaternion
coordinates above, so that \(c_Q(L_p(K))=A_pc_Q(K)\).  Up to scale, the
bilinear Frobenius form is \(J_{\rm Fr}=\diag(1,-1,1,-1)\).
Write \(\bar p\) for the complementary pair and define the right-\(X\)
multiplication map \(\mathsf R_X:\Mat_2(\mathbb C)\to\Mat_2(\mathbb C)\) and
its coordinate matrix \(\mathbf R_X\) by
\[
 \mathsf R_X(K):=KX,\qquad
 c_Q(\mathsf R_X(K))=\mathbf R_Xc_Q(K).
\]
Complementary-pair contraction gives
\begin{equation}
  L_{\bar p}=\mathsf R_X\circ L_p^*\circ\mathsf R_X,
  \label{eq:platonic-q4-map-adjoint}
\end{equation}
where \(^*\) is the Frobenius adjoint.  Thus
\begin{equation}
  A_{\bar p}=\mathbf R_XA_p^\sharp\mathbf R_X,
  \qquad A_p^\sharp=J_{\rm Fr}A_p^{\mathsf T}J_{\rm Fr}.
  \label{eq:platonic-q4-adjoint}
\end{equation}
Actual right \(X\) and reversals turn the six pair conditions into
\begin{equation*}
 \begin{aligned}
  &A_pr=0\ \text{or}\ [A_pr]\in\mathcal R_G,\\
  &A_p^\sharp r=0\ \text{or}\ [A_p^\sharp r]\in\mathcal R_G
 \end{aligned}
 \qquad\text{for every }[r]\in\mathcal R_G,
 \qquad
 \mathcal R_G=\{[c_Q(K_g)]:g\in G\}.
\end{equation*}
The exact two-sided classification
\cref{thm:cert-interface-platonic-q4}\textnormal{(i)} has no rank three: rank zero gives
\(q=0\), ranks one and four factor, and rank two occurs only for \(S_4\).
In the first marked form, actual inverse dressings expose \(\Eq_4\), so
\(\KHolant(\Lambda,\Eq_4)\leT\KHolant(\Lambda)\) and
\cref{cor:higher-equality-unary,lem:common-unary-exit} apply.  In the second form, its non-\(GO(X)\)
similarity identifies coefficient tables only; the same proposition checks
the dressing to \cref{eq:p1-platonic-q4-Oprime-eight-vertex}, directly in actual
\(\mathcal O'\)-representatives, with ranks \((2,2,2)\).  For
\(\mathcal O_{dt},\mathcal O_{td}\), the direct reconstructed tensors in
\cref{thm:cert-interface-platonic-q4}\textnormal{(i)} are exactly
\cref{eq:p1-platonic-q4-mixed-eight-vertex}.  These three cases invoke
\cref{lem:p1-platonic-eight-vertex-rank-terminal}.  Finally, the direct
dictionary \cref{lem:icosahedral-fixed-I-dictionary} proves the same
no-survivor conclusion for \(\mathcal I_-\) while leaving the literal
equality untouched.
\end{proof}

\begin{lemma}[Quadratic extension of a fixed matching]
\label{lem:platonic-global-ruling}
Let \(Q_0,Q_1,Q_2,Q_3\) be the quaternion basis in
\cref{eq:quaternion-basis}, and let
\(\mathcal D\subset\mathbb P(\Mat_2)\) contain the twelve tetrahedral
points
\begin{equation}
 [Q_\mu]\ (0\le\mu\le3),
 \qquad
 [Q_0+\eps_1Q_1+\eps_2Q_2+\eps_3Q_3]
 \quad(\eps_j\in\{1,-1\}).
 \label{eq:platonic-quadratic-test-set}
\end{equation}
Let \(M\) be a nonempty finite index set, let \(W_e\) be a nonzero
finite-dimensional complex vector space for every \(e\in M\), and let
\(\Phi:\Mat_2\to\bigotimes_{e\in M}W_e\) be a nonzero linear map.  Define
\[
 \Sigma_M:=
 \left\{\left[\bigotimes_{e\in M}w_e\right]:0\ne w_e\in W_e\right\}
 \subseteq\mathbb P\!\left(\bigotimes_{e\in M}W_e\right),
 \qquad
 \operatorname{Cone}_0(\Sigma_M):=\{0\}\cup
 \left\{\bigotimes_{e\in M}w_e:w_e\in W_e\right\},
\]
so that \(\Sigma_M\) and \(\operatorname{Cone}_0(\Sigma_M)\) are,
respectively, the Segre variety and its affine cone in the terminology of
\cref{subsubsec:projective-tensor-geometry}.
Suppose
\(\Phi(K)\in\operatorname{Cone}_0(\Sigma_M)\) for every
\([K]\in\mathcal D\).
Then there is one index \(e_*\in M\) and fixed nonzero tensors
\(U_e\in W_e\), \(e\ne e_*\), such that
\begin{equation}
  \Phi(K)=
  \left(\bigotimes_{e\ne e_*}U_e\right)\otimes\psi(K)
  \qquad(K\in\Mat_2)
  \label{eq:platonic-global-factorization}
\end{equation}
for a linear map \(\psi:\Mat_2\to W_{e_*}\).  For the reconstruction
conclusion, suppose \(M\) partitions the residual physical ports into
nonempty blocks, \(W_e=(\mathbb C^2)^{\otimes e}\), and
\(\Phi=\Phi_p^f\) under the corresponding port identification.  If
\(|M|\ge2\), then \(f=f_{\Phi,p}\) has a genuine proper tensor factor.
\end{lemma}

\begin{proof}[Proof of \cref{lem:platonic-global-ruling}]
Every defining Segre minor pulls back to a homogeneous quadratic in the
quaternion coordinates.  The four axes kill its square coefficients; at the
eight sign points the remaining expression is
\[
 \sum_{0\le\mu<\nu\le3}c_{\mu\nu}\eps_\mu\eps_\nu=0,
 \qquad \eps_0=1.
\]
The six displayed functions are distinct characters of
\(\{\pm1\}^3\), so every \(c_{\mu\nu}\) vanishes.  All Segre minors
therefore vanish identically, including at zero, and
\(\operatorname{im}\Phi\subseteq\operatorname{Cone}_0(\Sigma_M)\).

If a linear subspace lies in \(\operatorname{Cone}_0(\Sigma_M)\), take two nonzero simple tensors
\(u=\bigotimes u_e\) and \(v=\bigotimes v_e\) in it.  If \(u_i,v_i\) and
\(u_j,v_j\) are both independent for \(i\ne j\), apply dual functionals on
all other factors; a \(2\times2\) flattening minor of
\(u_i\otimes u_j+t\,v_i\otimes v_j\) is then nonzero, contradicting that
\(u+tv\) is simple.  Thus each pair varies in at most one factor.  Fixing
one nonzero \(u\), two vectors that vary from \(u\) in different factors
would violate this pairwise conclusion, so all variation is in one factor.
Hence the image lies in one ruling, giving
\cref{eq:platonic-global-factorization}.  At a test point with
\(\Phi(K)=0\), every minor vanishes automatically; the quadratic
interpolation never divides by a coordinate.  Thus zero values are allowed
throughout and no nonzero-test-point assumption was used.

In the card-map applications \(\Phi=\Phi_p^f\), \cref{eq:platonic-dual-basis}
is only a coefficient identity, and Frobenius reconstruction gives
\[
 f=
 \left(\bigotimes_{e\ne e_*}U_e\right)
 \otimes
 \left(\sum_{\mu=0}^3Q^\mu_p\otimes\psi(Q_\mu)\right).
\]
The fixed factor occupies the nonempty set
\(\bigcup_{e\ne e_*}e\), while the other factor contains \(p\), so the
factorization is proper.  \Cref{lem:factor-saturation} exposes the resulting factors under a
Turing reduction.
\end{proof}

\begin{lemma}[Deficient Platonic card maps]
\label{lem:platonic-deficient-rank}
Let \((\Lambda,G,h)\) be a marked Platonic state, where
\(G\in\{\mathcal T,\mathcal O,\mathcal O',\mathcal O_{dt},
\mathcal O_{td},\mathcal I,\mathcal I_-\}\).
Assume \(h\ne0\) is a directly realized six-port signature over \(\Lambda\), satisfies \cref{eq:platonic-all-card-condition}
for \(G\), and, for a specified canonical ordered pair \(p\), satisfies
\(\rank\Phi_p^h\le3\).  Then \(h\)
factors, or an actual binary gives a singular/infinite-order terminal or
a strict complete-group successor.  No rank-deficient incidence branch remains.
\end{lemma}

\begin{proof}
For \(G=\mathcal I_-\), apply the card covariance and direct network
dictionary of \cref{lem:icosahedral-fixed-I-dictionary}.  It identifies the
coefficient calculation with the standard \(\mathcal I\) calculation while
replacing every physical kernel by its actual \(\kappa_D\)-image in
\(\mathcal I_-\).  Thus its fixed-matching branch uses
\cref{lem:platonic-global-ruling}, and its mixed branch uses
\cref{thm:cert-interface-platonic-q4}\textnormal{(ii)--(iii)}, directly in
the original fixed-\(I\) state.  Hence assume below that
\(G\ne\mathcal I_-\).
Rank zero contradicts \(h\ne0\) by
\cref{eq:platonic-frobenius-reconstruction}, and rank one factors.  If all nonzero cards
use one matching, apply \cref{lem:platonic-global-ruling}; in any
nonstandard \(S_4\) form precompose with its fixed coefficient similarity
and pull the varying factor back by the inverse similarity.  This analytic pullback,
not a gadget, yields
\(\Phi_p^h(K)=(\bigotimes_{e\ne e_*}U_e)\otimes\psi(K)\);
\cref{lem:factor-saturation} exposes the proper factor.  If two nonzero cards use different matchings, the
physical interface \cref{thm:cert-interface-platonic-q4}\textnormal{(ii)} gives an
actual rainbow localization and a quaternary whose three
flattening ranks are \((3,3,3)\), while
the same interface gives its actual outside transfer.
The direct-gadget direction \cref{eq:direct-gadget-reduction} and binary
dispatch \cref{lem:binary-interpolation,lem:actual-binary-completeness} then
apply.
\end{proof}

\begin{lemma}[Finite Platonic quaternary exits]
\label{lem:platonic-finite-q4-terminal}
Let \((\Lambda,G,f)\) be a marked Platonic state, and let \(q\) be directly
realized over \(\Lambda\), with its canonical port
order and nonzero realization scalar recorded.  Suppose that \(q\) is one
of the explicitly enumerated representatives in
\cref{thm:cert-interface-platonic-q4}\textnormal{(iii)}, including every normalized output later
supplied by a marked rank-four separator.  Then
\(\KHolant(\Lambda,q)\leT\KHolant(\Lambda)\), and \(q\) directly realizes
one of the following: a rank-\((2,2,2)\) endpoint-nondegenerate
eight-vertex signature, a quaternary pure generalized equality, or a
nonsingular binary \(B\) with \([BX]\notin G\).
These invoke
\cref{lem:p1-platonic-eight-vertex-rank-terminal,lem:pure-ge,%
lem:actual-binary-completeness}; none recurses.
\end{lemma}

\begin{proof}
This is exactly the representative-wise three-way direct-realization
statement of \cref{thm:cert-interface-platonic-q4}\textnormal{(iii)}.
\end{proof}

Consequently the reductions
\cref{lem:platonic-quaternary-boundary,lem:platonic-deficient-rank,%
lem:platonic-finite-q4-terminal} are bounded local reductions.  After their
displayed physical output is constructed and its finite factorization tree is
processed, every branch is a resolved leaf, a lower-arity factor route, or
a strict complete-group successor.  An intermediate factor split remains
inside the originating card's finite factorization tree and strictly decreases
the factor potential of \cref{eq:factor-potential}.  None of these
conclusions invokes a queue-termination or restart lemma.

\subsection{Rainbow localization and matching-contraction closure}

\Cref{lem:platonic-global-ruling} handles cards with one matching.  We now
show that three collinear cards with different matchings differ on only
four ports.  Closing their common factors then converts this geometric
statement into a six-port gadget.

\begin{lemma}[Localization of a rainbow matching line]
\label{lem:platonic-rainbow-localization}
Let \(F_0,F_1,F_2\) be three distinct nonsingular matching products on one
projective line, with pairwise distinct perfect matchings \(M_0,M_1,M_2\).
Then there is a unique set \(S\) of four ports such
that the three matchings restrict on \(S\) to the three pairwise distinct
perfect matchings and the three tensors have common projective binary
factors on \(S^c\).  Moreover, this projective line contains exactly three
matching-product points.  After compatible projective rescaling one may
write
\[
 F_i=R\otimes E_i\qquad(0\le i\le2),
\]
where \(R\ne0\) is the fixed matching product on \(S^c\), and the
quaternary tensors \(E_0,E_1,E_2\) span a two-dimensional subspace.

In the physical setting, let \((\Lambda,G,f)\) be a deck-stable marked
Platonic state, where \(f\) has even arity at least six,
and suppose these three tensors occur as the images of one rich domain line
under the card map
\cref{eq:platonic-card-map}.  If every common binary factor \(U\) satisfies
\([UX]\in G\), and the actual kernel set
\(\mathcal K_G^{\rm act}\) spans
\(\Mat_2(\mathbb C)\), those factors can be closed by actual
\(G\)-kernels without annihilating the three selected cards.  This realizes
a six-port signature \(h\) with one \(f\)-vertex and exactly
\((\arity(f)-6)/2\) actual binary kernels, and
\begin{equation}
  \KHolant(\Lambda,h)
  \leT
  \KHolant(\Lambda).
  \label{eq:platonic-localization-reduction}
\end{equation}
Every other actual \(G\)-card of \(h\) remains zero or a nonsingular safe
quaternary \(G\)-matching product.
\end{lemma}

\begin{proof}[Proof of \cref{lem:platonic-rainbow-localization}]
For a port set \(S\), let \(c_i(S)\) count the edges of \(M_i\) crossing
\(S\mid\bar S\).  Nonsingularity of every binary factor makes the matching
recoverable from flattening ranks and gives
\[
 \rank (F_i)_{S\mid\bar S}=2^{c_i(S)}.
\]
In particular, proportional nonsingular matching products have the same
matching.  Collinearity also gives, for every permutation \((i,j,k)\),
\[
 2^{c_i(S)}\le2^{c_j(S)}+2^{c_k(S)}.
\]

Decompose \(M_0\cup M_1\) into common edges and alternating cycles.  For
each component \(C\), both \(F_0\) and \(F_1\) have flattening rank one
across \(C\mid C^c\), so \(\rank_C(F_2)\le2\).  The number \(c_2(C)\) is
even; hence \(2^{c_2(C)}\le2\) forces \(c_2(C)=0\).  Regard the components
as tensor factors.  A nontrivial linear combination of two simple tensors
is simple only when the two summands differ projectively in at most one
factor.  Since \(M_2\) is distinct from \(M_0,M_1\), it differs on at least
one alternating cycle (a common-edge component has only one internal
matching).  The simple-sum criterion permits only one such varying
component; hence there is exactly one exceptional alternating cycle, and
the common-edge factors outside it agree projectively in all three tensors.

The matching \(M_2\) cannot repeat \(M_0\) or \(M_1\) on the exceptional
component.  Indeed, flatten across the \(2\mid(2\ell-2)\) cut defined by
the two endpoints of an edge of the repeated matching: the two same-matching
tensors have rank one, so the third
has rank at most two, whereas its different matching has two edges crossing
that cut and hence rank four.  Thus the three restrictions on the
exceptional component are pairwise distinct.

Write the exceptional cycle as \(2\ell\) cyclic ports.  If \(\ell\ge3\),
consider its cyclic consecutive triples.  Each \(M_2\)-edge belongs to at
most two of them.  If every triple contained an \(M_2\)-edge, equality in
this incidence count would force every such edge to be a cycle edge and
hence \(M_2=M_0\) or \(M_1\).  Choose instead a triple \(T\) containing no
\(M_2\)-edge.  Then
\[
 (c_0(T),c_1(T),c_2(T))=(1,1,3),
\]
contradicting \(8\le2+2\).  Hence \(\ell=2\), so the exceptional component
is the asserted quaternary set and carries the three perfect matchings.  It
is unique by unique factorization.  A fourth matching-product point on the
same projective line would repeat one of the three quaternary matchings;
applying the preceding repeated-matching rank argument to it and either of
the other two points gives a contradiction.  Thus the line has exactly
three matching-product points, and the displayed decomposition
\(F_i=R\otimes E_i\) follows.

In the physical card-map setting, for each common factor \(U\), spanning by
\(\mathcal K_G^{\rm act}\) supplies an actual \(K\) with
\(\operatorname{tr}(U^{\mathsf T}K)\ne0\).  Closing the common factors
therefore preserves all three selected cards.  For any other card,
superpose its matching with the closure matching.  A closed common edge
contributes \(\operatorname{tr}(B^{\mathsf T}K)\), while a crossed path
splices two factors to \(B_1^{\mathsf T}KB_2\).  Put
\(T_1:=B_1X\), \(T_K:=KX\), and \(T_2:=B_2X\).  The latter is nonsingular
and, with \(T^\#=XT^{\mathsf T}X\), has transfer
\[
 \bigl[T_{B_1^{\mathsf T}KB_2}\bigr]
 =\bigl[T_1^\#XT_KXT_2\bigr]\in G.
\]
Internal cycles contribute scalars and may only make the card zero.
Iteration over the components proves \cref{eq:platonic-all-card-condition} and
\cref{eq:platonic-localization-reduction}.
\end{proof}

\begin{lemma}[Platonic matching-contraction closure]
\label{lem:platonic-matching-contraction-closure}
Let \(G\le\PGL_2\) contain \([X]\) and be reversal-stable.  Let \(F\) be
zero or a product of nonsingular binary
factors on a perfect matching, with every factor transfer in \(G\).  Contract
any disjoint collection of its ports by actual nonsingular \(G\)-binaries
(including the explicit \(I\)- or native-\(X\) links).  The result is zero or
again a product of nonsingular binary factors whose transfers lie in \(G\).
\end{lemma}

\begin{proof}
Superpose the factor matching with the contraction matching.  Components are
cycles and paths.  A cycle contributes a scalar (possibly zero); a path
splices \(B_1,K,B_2\) to \(B_1^{\mathsf T}KB_2\).  Put
\(T_1:=B_1X\), \(T_K:=KX\), and \(T_2:=B_2X\).  Its transfer class is
\[
 [T_1^\#XT_KXT_2],\qquad T^\#=XT^{\mathsf T}X.
\]
Reversal stability, \([X]\in G\), and closure of \(G\) keep this class in
\(G\); the omitted orientation scalar is nonzero.  Iterating over
components proves the claim.
\end{proof}

For the second octahedral display, with \(H_{\mathcal O}\) already fixed in
\cref{eq:octahedral-second-form}, define the coefficient-similarity map
\begin{equation}
 \theta_{\mathcal O}(U):=H_{\mathcal O}^{-1}UH_{\mathcal O},
 \label{eq:octahedral-coefficient-map}
\end{equation}
the transformation tuple
\[
 \mathbf H_{\mathcal O}:=
 (H_{\mathcal O}^{-1},H_{\mathcal O}^{\mathsf T},
  H_{\mathcal O}^{-1},H_{\mathcal O}^{\mathsf T}),
\]
and its induced quaternary coefficient operator
\[
 \mathcal L_{\mathcal O}:=\mathscr D_{\mathbf H_{\mathcal O}},
 \qquad
 \mathcal L_{\mathcal O}(F)=
 (H_{\mathcal O}^{-1}\otimes H_{\mathcal O}^{\mathsf T}\otimes
  H_{\mathcal O}^{-1}\otimes H_{\mathcal O}^{\mathsf T})F.
\]
These maps compare coefficient tables under the fixed-coordinate
convention at the start of the section.

For \(\alpha\in\{dt,td\}\), put
\[
 S_\alpha:=H_\alpha H_\alpha^{\mathsf T},\qquad
 \mathcal L_\alpha:=
 H_\alpha^{-1}\otimes H_\alpha^{\mathsf T}\otimes
 H_\alpha^{-1}\otimes H_\alpha^{\mathsf T}.
\]
By \cref{eq:octahedral-mixed-mark-check}, \([S_\alpha]\in\mathcal O\).
Ordered contraction gives the three matching formulas
\begin{align}
 \mathcal L_\alpha P_0(U,V)&=
  P_0\bigl(\theta_\alpha(U),\theta_\alpha(V)\bigr),\notag\\
 \mathcal L_\alpha P_1(U,V)&=
  P_1\bigl(\theta_\alpha(US_\alpha^{-1}),
            \theta_\alpha(S_\alpha V)\bigr),\notag\\
 \mathcal L_\alpha P_2(U,V)&=
  P_2\bigl(\theta_\alpha(U),
            \theta_\alpha(S_\alpha VS_\alpha^{-1})\bigr).
 \label{eq:octahedral-mixed-coefficient-transport}
\end{align}
Thus \(\mathcal L_\alpha\) bijects the standard and mixed matching decks,
preserves their matching labels, and intertwines their actual local group
actions after the displayed current-group orientation factors are absorbed.
This is a coefficient identity; the physical operations on the right use
only actual \(\mathcal O_\alpha\)-representatives.

The incidence counts, actual rainbow-line orbits, and deletion
connectivity used below are supplied by
\cref{thm:cert-interface-platonic-incidence}.

\begin{lemma}[High-arity Platonic rank-deficient localization]
\label{lem:platonic-high-arity-localization}
Let \((\Lambda,G,f)\) be a deck-stable retained, factor-saturated fixed-\(I\)
state, where \(f\ne0\) is a tensor-prime minimum core
of even arity \(n\ge6\) and
\[
 G\in\{\mathcal T,\mathcal O,\mathcal O',\mathcal O_{dt},
          \mathcal O_{td},\mathcal I,\mathcal I_-\}.
\]
Fix an ordered pair \(p\), define the ancestor card-map rank
\(r:=\rank\Phi_p^f\in\{2,3,4\}\), and
suppose that, for every ordered pair of ports and every actual
\(G\)-kernel, the corresponding ancestor card is zero or a nonsingular
safe \(G\)-matching product.  Then at least one of the following structural
routes applies; when both are available, choose route~1.
\begin{enumerate}
\item All nonzero \(p\)-cards use one residual matching, and
      \(f\) has a genuine proper tensor factor.
\item One \(f\)-vertex and exactly \((n-6)/2\) actual binary kernels
      directly realize a nonzero six-port signature \(h\), retaining
      \(p\), such that
      \begin{equation}
        2\le\rank\Phi_p^h\le r\le4,
        \label{eq:platonic-localized-rank-bound}
      \end{equation}
      and every actual \(G\)-card of \(h\) is zero or a nonsingular safe
      quaternary \(G\)-matching product.  In particular,
      \(\KHolant(\Lambda,h)\leT\KHolant(\Lambda)\).
\end{enumerate}
The second route recomputes the rank of the directly realized six-port card map;
the ancestor rank is only the upper bound in
\cref{eq:platonic-localized-rank-bound}.
\end{lemma}

\begin{proof}
Use the domain graph \(\mathscr G_G^{\rm dom}\) defined above.  The kernel-deleted
connectivity assertion of
\cref{thm:cert-interface-platonic-incidence}\textnormal{(iii)}
says
that
\[
 \mathcal D_G\setminus\mathbb P(\ker\Phi_p^f)
\]
is nonempty and connected.  If all its nonzero images have one matching,
\cref{lem:platonic-global-ruling} gives a proper factor.  For any
nonstandard octahedral form, precompose the card map with its fixed
invertible coefficient similarity and pull the resulting tensor identity
back.  For \(\mathcal I_-\), use
\cref{lem:icosahedral-fixed-I-dictionary}.  These are analytic
reconstructions, and factor access is supplied by
\cref{lem:factor-saturation}.

Otherwise choose the first matching-changing edge along a path in this
deleted graph, and let \(\ell\) be a rich domain line containing that edge.  If
\(r\in\{2,3\}\) and \(\ell\) met
\(\mathbb P(\ker\Phi_p^f)\), the restriction of \(\Phi_p^f\) to the
two-dimensional vector span of \(\ell\) would have rank at most one.
All nonkernel images on \(\ell\) would then be projectively proportional
and have the same matching, a contradiction.  If \(r=4\), the kernel is
zero.  Thus \(\ell\) is kernel-avoiding.  Let \(V_\ell\) be its
two-dimensional vector span.  The matching-changing edge gives
\(\rank(\Phi_p^f|_{V_\ell})\ge2\), while kernel-avoidance gives
\(\rank(\Phi_p^f|_{V_\ell})\le2\); hence the restriction has rank two and
induces a projective injection on \(\ell\).  Let \(F_0,F_1\) be the
matching-changing images, and choose any third domain point on this rich
line, with image \(F_2\).  Its matching differs from both of theirs.
Indeed, if \(M_2=M_0\), choose an edge of \(M_0\setminus M_1\) and
flatten across its two ports.  The ranks of \(F_0,F_2\) are one, whereas
the rank of \(F_1\) is four, because two nonsingular factors cross this
cut.  This contradicts collinearity; the case \(M_2=M_1\) is identical.
Thus the three images satisfy the general-arity hypotheses of
\cref{lem:platonic-rainbow-localization}.  That lemma says their line
contains exactly three matching-product points, so injectivity also forces
\(\ell\) to have exactly three domain points.  The quaternary incidence
classification is used only after this localization.

The same lemma supplies a unique
quaternary set \(S\) and compatible rescalings
\[
 \Phi_p^f(D_i)=R\otimes E_i,
 \qquad
 R=\bigotimes_{j=1}^{(n-6)/2}U_j,
 \qquad 0\le i\le2,
\]
where the \(E_i\) form the residual rainbow deck line and
\(D_i\in\mathcal K_G^{\rm act}\) are the selected domain representatives.  The actual kernels span
\(\Mat_2\), so for every \(U_j\ne0\) choose an actual \(K_j\) with
\(\operatorname{tr}(U_j^{\mathsf T}K_j)\ne0\).  Close the corresponding
common-factor ports by these \(K_j\), leaving \(p\cup S\).  The resulting
one-vertex signature \(h\) is nonzero because on each selected card the
closing scalar is
 \[
  \kappa=\prod_j\operatorname{tr}(U_j^{\mathsf T}K_j)\ne0.
 \]
Record the retained pair \(p\), selected line \(\ell\), ordered kernels
\((K_j)_j\), output order \(\pi\), and scalar \(\kappa\).
If \(C\) denotes the resulting contraction on the residual ports and
\(P_\pi\) the invertible
output reindexing induced by the recorded port permutation, then
\[
 \Phi_p^h=P_\pi\circ C\circ\Phi_p^f.
\]
This proves the upper rank bound, while the surviving rainbow deck line proves
\(\rank\Phi_p^h\ge2\).

It remains to verify every card, not only the three selected ones.  Each is
obtained from an ancestor \(G\)-matching card by the remaining actual
closures, so \cref{lem:platonic-matching-contraction-closure} gives zero or
a nonsingular safe \(G\)-matching product.  The tensors \(U_j\) are factors
found inside the selected cards; the only gadgets inserted are the actual
kernels \(K_j\).
\end{proof}

Whenever localization supplies a rainbow deck line, use actual dressings
on the domain pair and on the four residual ports to send the two lines to
their specified representatives.  The two operations act on disjoint port
sets and commute; retain their scalars and port orders.
In the standard forms \(G\in\{\mathcal T,\mathcal O,\mathcal I\}\),
let \(\Phi\) be this actual card map after one common nonzero scalar is
removed, and put \(G_0=G\).
For a nonstandard octahedral form, write \(\theta,\mathcal L\) for its
fixed coefficient maps and \(\Phi_{\rm phys}\) for the actual normalized
map.  Its domain and image representatives are the \(\theta\)- and
\(\mathcal L\)-images of the standard ones.  Only for coefficient
calculations put
\begin{equation}
 \Phi=\mathcal L^{-1}\circ\Phi_{\rm phys}\circ\theta,
 \qquad
 \Phi_{\rm phys}(\theta K)=\mathcal L\Phi(K),
 \qquad G_0=\mathcal O.
 \label{eq:platonic-nonstandard-card-dictionary}
\end{equation}
Thus a kernel \(K\) in a standard formula means the actual current-form
kernel \(\theta K\), and its physical output is \(\mathcal L\Phi(K)\).
Neither coefficient map is a gadget.  The opposite icosahedral display uses
its own \(\kappa_D\)-dictionary in
\cref{lem:icosahedral-fixed-I-dictionary}.
In these standard coefficient coordinates define
\begin{equation}
 A=P_0(Q_0,Q_0),\qquad
 B=P_1(Q_0,Q_0),\qquad
 C=P_2(Q_2,Q_2)=A-B.
 \label{eq:platonic-rainbow-normalization}
\end{equation}
The two displayed input matrices sum to \(2Q_0\).  Together with
\(C=A-B\), this fixes the relative image scalars.  Dividing by their one common recorded nonzero
scalar, the normalized map satisfies
\begin{equation}
 \Phi(Q_0+Q_1+Q_2+Q_3)=A,
 \qquad
 \Phi(Q_0-Q_1-Q_2-Q_3)=-B.
 \label{eq:platonic-rainbow-images}
\end{equation}
Thus, defining the image tensors \(W_j:=\Phi(Q_j)\),
\begin{equation}
 \Phi(Q_0)=C/2,
 \qquad
 W_1+W_2+W_3=(A+B)/2,
 \label{eq:platonic-rainbow-sum}
\end{equation}
and each remaining sign point imposes
\begin{equation}
 W_j,\quad A-2W_j,\quad 2W_j-B
 \in\{0\}\cup\operatorname{cone}\bigl(\mathfrak D_4(G_0)\bigr).
 \label{eq:platonic-separated-sign-tests}
\end{equation}
These are coefficient identities of actual cards, read through
\cref{eq:platonic-nonstandard-card-dictionary} in the nonstandard forms.
Every terminal below is checked on its physical current-form output.

\subsection{The four fixed-equality marked \texorpdfstring{\(S_4\)}{S4} forms}

The projective octahedral group has four ordered marked-involution forms in
the coordinates containing the supplied literal \(I\).  They are the
standard \((dd;d)\) form \(\mathcal O\), the \((tt;d)\) form
\(\mathcal O'\), and the mixed forms
\(\mathcal O_{dt},\mathcal O_{td}\) fixed in
\cref{eq:octahedral-second-form,eq:octahedral-mixed-forms}.  They are
exhaustive by \cref{thm:marked-finite-group-routing}.  The three displayed
similarity matrices compare coefficient configurations only; physical
operations always use current-form representatives.
For a normalized rainbow map, the exact bridge has six rank-deficient solutions,
handled on that directly realized six-port map by
\cref{lem:platonic-deficient-rank}, and the twenty-four orderings of the
four frames in
\cref{thm:cert-interface-platonic-bridges}\textnormal{(i)}.

Retain the even-entry order
\begin{equation}
  \mathbf e(M)=
  (M_{0000},M_{0011},M_{0101},M_{0110},
   M_{1001},M_{1010},M_{1100},M_{1111}).
  \label{eq:octahedral-even-order}
\end{equation}

\begin{lemma}[Uniform octahedral full-rank exits]
\label{lem:octahedral-separator}
Let \(\Phi\) be a normalized full-rank map for which
\((\Phi(Q_1),\Phi(Q_2),\Phi(Q_3))\) is one of the twenty-four orderings of
the four frames in
\cref{thm:cert-interface-platonic-bridges}\textnormal{(i)}.  In a nonstandard
form, \(\Phi\) is the coefficient map associated to the actual
\(\Phi_{\rm phys}\) by \cref{eq:platonic-nonstandard-card-dictionary}.
Define
\[
 K_{\mathcal O}:=Q_0+Q_1,\qquad
 K_{\mathcal O'}:=\theta_{\mathcal O}(K_{\mathcal O}),\qquad
 K_{\mathcal O_\alpha}:=\theta_\alpha(K_{\mathcal O})
 \quad(\alpha\in\{dt,td\}).
\]
For each marked form, apply its displayed kernel to the actual card map.
At most four further actual group binaries, used in the indicated port
dressings or pair closure, then directly realize one of:
an endpoint-nondegenerate eight-vertex signature of ranks \((2,2,2)\), a
quaternary pure generalized equality, or a nonsingular binary whose transfer
lies outside the current complete group.  For \(\mathcal O_{dt}\), the
first-frame entries \(R_0,R_1,R_3\) give the eight-vertex exits and the
other six named entries give outside binaries.  For
\(\mathcal O_{td}\), \(R_0,R_3\) give the eight-vertex exits, \(R_1\)
gives the pure generalized equality, and the other six give outside
binaries.  Thus every full-rank map is a resolved local exit.
\end{lemma}

\begin{proof}
This is \cref{thm:cert-interface-platonic-bridges}\textnormal{(i)} together
with the direct current-coordinate tables in
\cref{prop:p1-cert-S4-uniform-eight-vertex-separator}.  Every listed kernel
and dressing is an actual representative in its own marked form; the
similarities only identify the coefficient tables.  Hence
\cref{eq:direct-gadget-reduction} gives the stated directions, followed by
\cref{lem:p1-platonic-eight-vertex-rank-terminal,lem:pure-ge,%
lem:actual-binary-completeness} according to the listed exit.
\end{proof}

For the full-rank six-port statements below, assume
\begin{equation}
 \begin{aligned}
 &\Stable(\Lambda,G,h),\qquad \arity(h)=\nu(\Lambda)=6,\\
 &h\text{ is directly realized over }\Lambda,\qquad \rank\Phi_p^h=4.
 \end{aligned}
 \label{eq:platonic-six-four-predicate}
\end{equation}
The rank in \cref{eq:platonic-six-four-predicate} is always recomputed from
the actual coefficients of \(h\), never inherited from an ancestor.  Under
this condition, put \(\Phi:=\Phi_p^h\).  Varying card matchings give a
rainbow deck line by clauses~\textnormal{(i)--(iii)} of
\cref{thm:cert-interface-platonic-incidence}, together with
\cref{eq:matching-schmidt-rank};
\cref{lem:platonic-rainbow-localization} then supplies the actual
normalization.  Apply \cref{lem:octahedral-separator} to its normalized
coefficient map, using \cref{eq:platonic-nonstandard-card-dictionary} in a
nonstandard form, to obtain a resolved eight-vertex,
pure-generalized-equality, or outside-binary leaf.

If all nonzero cards use one matching \(M\), define the auxiliary card map
\(\Psi:=\Phi\) in the first form,
\(\Psi:=\Phi\circ\theta_{\mathcal O}\) for \(\mathcal O'\), and
\(\Psi:=\Phi\circ\theta_\alpha\) for \(\mathcal O_\alpha\),
\(\alpha\in\{dt,td\}\).  The twelve test points
map to zero or an \(M\)-matching product, so
\cref{lem:platonic-global-ruling} factors
\(\Psi\).  In a nonstandard form, composing the varying factor with the
inverse coefficient similarity pulls this tensor identity back to \(\Phi\);
it does not use a similarity matrix as a gadget.  The reconstructed tensor
\(h\) has a
genuine proper factor.

\begin{theorem}[Octahedral deck rigidity]
\label{thm:octahedral-deck-rigidity}
If \cref{eq:platonic-six-four-predicate} holds and
\[
 G\in\{\mathcal O,\mathcal O',\mathcal O_{dt},\mathcal O_{td}\},
\]
then \(\Close_{\varnothing}(\Lambda,G,h)\).
\end{theorem}

\begin{proof}
If card matchings vary,
\cref{lem:platonic-rainbow-localization,lem:octahedral-separator} gives a
resolved local exit.  If they are fixed,
\cref{lem:platonic-global-ruling,lem:factor-saturation} exposes a proper
factor of the active six-port occurrence, including after the coefficient
pullback for any of the three nonstandard forms.

Close that finite factorization tree directly.  A nonbinary prime factor has
arity below six and therefore lowers \(\nu\).  A singular or unsafe binary
is terminal, and a finite safe binary whose transfer class lies outside
\(G\) strictly enlarges the recomputed complete group.  If every factor
were a safe binary whose transfer class already lies in \(G\), the original
signature would be a \(G\)-matching product, contrary
to the retained tensor-prime branch.  Hence every branch is a resolved leaf
or a strict outer successor, proving
\(\Close_{\varnothing}(\Lambda,G,h)\).  The rank and factorization used here belong to the realized tensor \(h\).
\end{proof}

\subsection{The internal \texorpdfstring{\(A_4\)}{A4} form}

The internal tetrahedral case uses the standard twelve-point model
\(\mathcal T\) in \cref{eq:tetrahedral-domain}.  It is closed under
multiplication, inverse, and transpose.  All normalizations below use actual
left and right \(\mathcal T\)-multiplication and port permutations.

Let \(\Phi:\Mat_2\to{(\mathbb C^2)}^{\otimes4}\) be the full-rank map
localized from a tetrahedral rainbow deck line.  We use the common normalization
\cref{eq:platonic-rainbow-normalization,eq:platonic-rainbow-images,%
eq:platonic-rainbow-sum,eq:platonic-separated-sign-tests} with
\(G=\mathcal T\).

\begin{lemma}[Tetrahedral bridge]
\label{lem:tetrahedral-bridge}
The simultaneous solution set of
\cref{eq:platonic-separated-sign-tests} is
\begin{equation*}
  \mathcal B_{\mathcal T}
  =\left\{0,\frac A2,\frac B2,R_0,R_1,R_3\right\},
\end{equation*}
where
\begin{equation*}
  R_0=\frac12P_2(Q_0,Q_0),\qquad
  R_1=-\frac12P_2(Q_1,Q_1),\qquad
  R_3=-\frac12P_2(Q_3,Q_3).
\end{equation*}
The sum equation leaves exactly twelve ordered maps: the six rank-two
permutations of \((0,A/2,B/2)\) and the six rank-four permutations of
\((R_0,R_1,R_3)\).
\end{lemma}

\begin{proof}
The exact separated-sign and sum calculation, including dependent lines,
is \cref{thm:cert-interface-platonic-bridges}\textnormal{(ii)}.
\end{proof}

For the rank-four maps write \(\sigma=(a,b,c)\) for a permutation of
\((0,1,3)\).  The corresponding six-port ancestor is uniquely reconstructed
as
\begin{equation*}
  f_\sigma=
  \sum_{\mu=0}^3Q^\mu_{12}\otimes F^\sigma_\mu,\qquad
  (F^\sigma_0,F^\sigma_1,F^\sigma_2,F^\sigma_3)
  =(C/2,R_a,R_b,R_c).
\end{equation*}

\begin{lemma}[The two permutation classes]
\label{lem:tetrahedral-orientation}
The six rank-four maps split into two sets:
\begin{align*}
  \Sigma_+&=\{(0,1,3),(1,3,0),(3,0,1)\},\notag\\
  \Sigma_-&=\{(0,3,1),(1,0,3),(3,1,0)\}.
\end{align*}
For \(\sigma\in\Sigma_+\), all fifteen deleted-pair contractions by all
twelve actual tetrahedral kernels are nonzero matching products.  For
\(\sigma\in\Sigma_-\), contracting ports \(1,3\) by the actual kernel
\(Q_0=I\) gives a directly realizable quaternary \(q_\sigma\) whose three
flattening ranks are \((2,2,2)\).  For each such \(\sigma\), four
actual tetrahedral port dressings transform \(q_\sigma\) into an
endpoint-nondegenerate eight-vertex signature of the same ranks.  The
contraction and dressings are direct gadgets from \(f_\sigma\), so
\cref{lem:p1-platonic-eight-vertex-rank-terminal} applies in the required
direction.
\end{lemma}

\begin{proof}
\Cref{thm:cert-interface-platonic-bridges}\textnormal{(ii)} gives the split,
common \(Q_0\)-card, and actual
endpoint-nondegenerate eight-vertex signature dressings.  Thus the two
direct-gadget reductions point back to \(f_\sigma\), and
\cref{lem:p1-platonic-eight-vertex-rank-terminal} applies.
\end{proof}

\begin{lemma}[The orientation-preserving tensors are affine]
\label{lem:tetrahedral-affine}
For every \(\sigma\in\Sigma_+\), the tensor \(f_\sigma\) belongs to the
affine class \(\cA\) in the same computational basis as all twelve
tetrahedral binaries.
\end{lemma}

\begin{proof}
\Cref{thm:cert-interface-platonic-bridges}\textnormal{(ii)} gives affine
support and even mixed quadratic coefficients; \cref{eq:tetrahedral-domain} puts all twelve
binaries in that same affine basis.
\end{proof}

To lift this statement to the signature set, define the tetrahedral matrices
\[
  R_+=Q_0+Q_1+Q_2+Q_3,\qquad
  R_-=Q_0-Q_1-Q_2-Q_3.
\]
Then
\begin{equation*}
 R_+R_-=4I,\qquad
 R_+Q_1=-Q_0+Q_1-Q_2+Q_3,
\end{equation*}
so \([R_+],[R_-],[Q_1],[R_+Q_1]\in\mathcal T\), and all four matrices
have actual inverses up to nonzero scalar.  For a matrix \(M\), let
\(L_j(M)\) denote its physical action \(M(y_j,x_j)\) on port \(j\), and
let \(\Pi_{ab}\) exchange output ports \(a,b\); products of permutations
below act from right to left.  Entrywise expansion in the fixed physical
port convention gives the identities
\begin{align}
 \Pi_{45}\Pi_{23}\Pi_{34}
 \bigl[I\otimes R_+\otimes I^{\otimes3}\otimes R_+\bigr]\Hcore
   &=-16 f_{(0,1,3)},\notag\\
 \Pi_{34}
 \bigl[I\otimes R_-\otimes(R_+Q_1)\otimes R_+
              \otimes I^{\otimes2}\bigr]\Hcore
   &=-32 f_{(1,3,0)},\notag\\
 \Pi_{23}
 \bigl[I\otimes R_+\otimes R_+\otimes I\otimes R_+\otimes Q_1\bigr]\Hcore
   &=32 f_{(3,0,1)}.
  \label{eq:tetrahedral-H6-equivalence}
\end{align}
The brackets are tensor products of the six \(L_j\)-actions, not a
coefficient-side similarity.  Because \(R_+^{-1}=R_-/4\) and
\(Q_1^{-1}=-Q_1\), every local action and permutation is projectively
actual-invertible.  Thus all three equivalences in
\cref{eq:tetrahedral-H6-equivalence} are reversible actual identities.
Every \(f_\sigma\), \(\sigma\in\Sigma_+\), therefore directly gives literal
\(\Hcore\); four copies then directly give \(\RMcore\) by
\cref{eq:app-v4-H6-R8-output}.  The complete group is \(A_4\), so the signature-set conclusion uses the
twelve-card relative theorem \cref{thm:a4-relative}.  For any ambient retained set
\(\Lambda\) over which \(f_\sigma\) is directly realizable, both
direct-gadget reductions point from
\(\KHolant(\Lambda,\Hcore,\RMcore)\) through
\(\KHolant(\Lambda,f_\sigma)\) to \(\KHolant(\Lambda)\).

After the quaternary boundary, let an ancestor \(f\) have even arity at
least six.  Rank zero is impossible and rank one factors.  In ranks two
through four, \cref{lem:platonic-high-arity-localization} either gives the
proper fixed-matching factor or materializes the safe six-port bridge.
The bridge rank is then recomputed from its own coefficients before
\cref{lem:tetrahedral-bridge,lem:tetrahedral-orientation} is invoked; an
ancestor rank is never transferred to that reduction.

\paragraph{From one core to the retained signature set.}

To obtain a common affine basis, every retained occurrence \(g\) must be
classified with all twelve tetrahedral kernels.  A binary factor may have
transfer in \(\mathcal T\setminus V_4\), so the complete family is
\(\{\Phi_p^g(K):K\in\mathcal K_{\mathcal T}^{\rm act}\}\) for each pair
\(p\).  The six-port classification supplies the two sets below; an
eight-port base and a stabilizer-code shortening argument then handle
larger arities.

Define the projective six-port matching set, H-core orbit, and safe union by
\begin{align*}
 \mathscr M_6
   &=\left\{\left[\prod_{\{a,b\}\in M}U_{ab}\right]:
       M\text{ is a perfect matching of }[6],\ [U_{ab}X]\in\mathcal T
     \right\},\\
 \mathscr H_6
   &=\operatorname{Orb}_{\mathcal T^6\rtimes S_6}([\Hcore]),
 \qquad
 \mathscr S_6=\mathscr M_6\mathbin{\dot\cup}\mathscr H_6.
\end{align*}
Define their affine cone by
\(\mathscr C_6:=\{0\}\cup\operatorname{cone}(\mathscr S_6)\).
Here the local action uses the fixed actual representatives of
\(\mathcal T\), so every orbit equivalence is an actual reversible port
dressing.  Their disjointness and exact cardinalities are given by
\cref{thm:cert-interface-A4-q8}.

\begin{proposition}[Hereditary tetrahedral eight-port base]
\label{prop:a4rel-required-eight-port}
Let \(g\ne0\) be an eight-port signature in the normalized tetrahedral
branch whose every physical tetrahedral pair card lies in \(\mathscr C_6\).
Then \(g\) factors across a \(2|6\) cut or \(g\in\cA\) in the normalized
basis.  Thus every tensor-prime eight-port survivor is affine; factors are
exposed only through factor saturation.
\end{proposition}

\begin{proof}
For a pair map \(\Phi=\Phi_p^g\), actual dressings normalize a nonzero card
to a three-tetrahedral-binary product or \(\Hcore\).  On the four domain
lines through that card, opposite endpoints sum to twice it.  If both
endpoint images are dependent on every line, their four sign differences
span the subspace \(\operatorname{span}\{Q_1,Q_2,Q_3\}\), so
\(\rank\Phi=1\); dual-basis
reconstruction gives a genuine \(2|6\) factor.  Otherwise the
scale-sensitive exhaustive audit
\cref{thm:cert-interface-A4-q8} says that every nonaffine parent has
an actual card outside \(\mathscr C_6\), contrary to the hypothesis.
The remaining parents are affine.  Actuality and arbitrary complex scales
are part of the cited audit; factors use only \cref{lem:factor-saturation}.
\end{proof}

We use the stabilizer description of affine signatures from
\cref{subsubsec:pauli-stabilizer,subsubsec:stabilizer-error-correction} to
locate the varying part of the twelve cards.  The next two lemmas show
that this variation occupies at most four logical qubits.  The encoding
is an analytic description of their coefficient space.  The shortening
lemma then selects an actual tetrahedral kernel from that description.

\begin{lemma}[Sharp stabilizer-line bound]
\label{lem:a4rel-sharp-stabilizer-line}
Let \(u,v\) be nonproportional nonzero affine signatures of common
arity \(m\), and let \(c\ne0\).  If \(u+cv\) is also nonzero and affine, then, after applying the same analytic Clifford change to both
tensors, followed by a common nonzero scalar normalization sending
\(u\) to \(\delta_{0^m}\), the support of the transformed
\(v\) is either a singleton disjoint from zero or a linear subspace of
dimension at most two.
\end{lemma}

\begin{proof}
Choose a unitary Clifford \(C\) with \(Cu=\alpha\delta_{0^m}\),
where \(\alpha\ne0\), and apply \(\alpha^{-1}C\) to both tensors.
The scalar normalization is separate from the Clifford operator.
Relabel the resulting second tensor as \(v\).
Write \(v\) in affine normal form as
\(v=\lambda\one[a+L]i^Q\), where \(\lambda\in\mathbb C^\times\) is the scale,
\(a+L\) is the affine support, \(Q\) is the quadratic phase, and
\(d:=\dim L\).  If
\(0\notin a+L\), the support size \(1+2^d\) forces \(d=0\).  If
\(0\in a+L=L\) and the zero coefficient cancels, \(2^d-1\) forces \(d=1\).
Otherwise rescale the phase quotient to one off zero.  For \(d\ge3\), its
multiplicative third difference on a three-cube is one (both phases are
quadratic), so its value at zero is also one, contradicting
nonproportionality.  Hence \(d\le2\).
\end{proof}

\begin{lemma}[The tetrahedral card cube has at most four active qubits]
\label{lem:a4rel-four-active-qubits}
Let
\(
 \Phi:\Mat_2(\mathbb C)\to(\mathbb C^2)^{\otimes m}
\)
be linear.  Suppose that every one of the twelve values
\(\Phi(K)\), \([K]\in\mathcal T\), is zero or affine.
Then there are an analytic Clifford operator \(C\) on the \(m\) output
qubits and a subspace \(R\le\Ftwo^m\), \(\dim R\le4\), such that
\begin{equation}
 \supp(C\Phi(K))\subseteq R
 \qquad(K\in\Mat_2(\mathbb C)).
 \label{eq:a4rel-common-active-support}
\end{equation}
Equivalently, the image of \(\Phi\) is contained in one stabilizer code
with at most four logical qubits.
\end{lemma}

\begin{proof}
If \(\Phi=0\), take \(R=\{0\}\).  Otherwise at least one tetrahedral test
point has nonzero image because those points span \(\Mat_2\).  Precompose
\(\Phi\) by an actual tetrahedral multiplication carrying that point to
\(A:=Q_0-Q_1-Q_2-Q_3\), relabel the resulting map as \(\Phi\), and hence
assume \(\Phi(A)\ne0\).  Define
\[
 B_0:=2Q_0-A,\qquad B_j:=A+2Q_j\quad(1\le j\le3).
\]
These are the four opposite-parity tetrahedral sign points paired with
\(A\); they
form a basis and satisfy
\[
 A+B_0=2Q_0,\qquad B_j-A=2Q_j\ (1\le j\le3),\qquad
 2A=-B_0+B_1+B_2+B_3.
\]
Choose an analytic Clifford \(C\) with
\(C\Phi(A)=\alpha\delta_{0^m}\), where \(\alpha\ne0\).
For the following analysis put \(\Psi(K):=\alpha^{-1}C\Phi(K)\);
this common scalar normalization is separate from the Clifford operator
and leaves every support unchanged.  By the normalized-delta case in the
proof of \cref{lem:a4rel-sharp-stabilizer-line}, each nonzero
\(\Psi(B_j)\) not proportional to \(\delta_{0^m}\) has, off zero, either
one point or the three nonzero points of a two-dimensional
\(\Ftwo\)-subspace.  The last displayed relation makes
every point in their union occur at least twice.  If its span had dimension
at least five, total incidence at most twelve would force sizes
\((3,3,3,1)\) or \((3,3,3,3)\).  In type \((3,3,3,1)\), double incidence
leaves at most five distinct nonzero points, and any three-point support is
the nonzero part \(\{x,y,x+y\}\) of a two-dimensional
\(\Ftwo\)-subspace.  Its line relation is
nontrivial, so those at most five points have span at most four.

In type \((3,3,3,3)\), double incidence leaves at most six points.  The only
remaining extremal case has six points, each occurring in exactly two of
the four supports.  Regard a point shared by supports \(i,j\) as an edge
between vertices \(i,j\).  The four supports form a loopless
three-regular multigraph on four vertices, and the line relation at each
vertex is the sum of its three incident edge labels.  The binary
vertex--edge incidence matrix has rank at least two: its rank is
\(4-c\), where the graph has \(c\le2\) connected components.  Hence the six
edge labels satisfy at least two independent relations and again span at
most four.  Thus their union spans some \(R\) with \(\dim R\le4\).  Since
the \(B_j\)'s form a basis and the common nonzero scalar leaves
supports unchanged, \cref{eq:a4rel-common-active-support} follows.
\end{proof}

For the physical shortening step, let \(\mathcal O\) denote the projective
one-qubit Clifford group in the quaternion coordinates
\cref{eq:quaternion-basis}.  The normal subgroup
\(\mathcal T\triangleleft\mathcal O\) has index two.  Each of its two cosets
contains the following eight full-support representatives:
\begin{align}
 \mathcal C_+
 &=\{Q_0+\epsilon_1Q_1+\epsilon_2Q_2+\epsilon_3Q_3:
       \epsilon_j\in\{\pm1\}\},\notag\\
 \mathcal C_-
 &=\{Q_d+\epsilon Q_o:
       d\in\{0,3\},\ o\in\{1,2\},\ \epsilon\in\{\pm1\}\}.
 \label{eq:a4rel-full-support-cosets}
\end{align}
On the open set \(K_{01}K_{10}\ne0\), define the cross-ratio rational
function \(\chi(K):=K_{00}K_{11}/(K_{01}K_{10})\).  Direct use of
\cref{eq:quaternion-basis} gives
\begin{equation}
 K_{xy}\ne0,\qquad \chi(K)=-1
 \qquad(K\in\mathcal C_+\cup\mathcal C_-).
 \label{eq:a4rel-full-support-cross-ratio}
\end{equation}
Indeed, in \(\mathcal C_+\) the diagonal and off-diagonal products are
\(2\) and \(-2\), while in \(\mathcal C_-\) they are \(1\) and \(-1\).
After division by one entry, every such table and its pointwise inverse
have entries in \(\mu_4\), so both are nowhere-zero affine phase tables.

\begin{lemma}[Stabilizer erasure and column normal form]
\label{lem:stabilizer-erasure-column-normal-form}
Let
\(
 V:(\mathbb C^2)^{\otimes k}\to(\mathbb C^2)^{\otimes N}
\)
be a stabilizer isometry and let \(r\) be a physical output port.  If
erasure of \(r\) is correctable, meaning
\(V^\dagger R_rV\) is scalar for every one-qubit operator \(R\),
exactly one of the following normal forms applies after a fixed one-qubit
Clifford on \(r\):
\begin{enumerate}
\item there are a one-qubit stabilizer state \(\lvert\chi\rangle_r\) and a
      stabilizer isometry \(W\) such that
      \begin{equation}
       V\lvert\psi\rangle
       =\lvert\chi\rangle_r\otimes W\lvert\psi\rangle;
       \label{eq:stabilizer-erasure-pure}
      \end{equation}
\item there are an auxiliary qubit \(a\) and a stabilizer isometry
      \(
       W:\mathbb C^2_a\otimes(\mathbb C^2)^{\otimes k}
       \to(\mathbb C^2)^{\otimes(N-1)}
      \)
      such that
      \begin{equation}
       V\lvert\psi\rangle
       =2^{-1/2}\sum_{c\in\Ftwo}
        \lvert c\rangle_r W(\lvert c\rangle_a\otimes\lvert\psi\rangle).
       \label{eq:stabilizer-erasure-mixed}
      \end{equation}
\end{enumerate}
If no single-site erasure is correctable, then, after postcomposing the
outputs by a tensor product of one-qubit Cliffords and precomposing the input
by one logical Clifford, there are a vector
\(c\in\Ftwo^N\), and a matrix
\(M=(m_1\ \cdots\ m_N)\in\Ftwo^{k\times N}\) such that
\begin{equation}
 V\lvert z\rangle
 =\zeta(z)\lvert c+M^{\mathsf T}z\rangle,
 \qquad \rank M=k,\qquad m_r\ne0\quad(1\le r\le N),
 \label{eq:stabilizer-column-normal-form}
\end{equation}
where \(\zeta(z)=\zeta_0i^{Q(z)}\) is a nowhere-zero stabilizer quadratic
phase in the precise sense of
\cref{subsubsec:pauli-stabilizer}.
\end{lemma}

\begin{proof}
By the Hilbert-space and error-correction conventions of
\cref{subsubsec:hilbert-conventions,subsubsec:stabilizer-error-correction},
correctability of erasure
at \(r\) is the single-site Knill--Laflamme condition
\begin{equation}
 V^\dagger P_rV=\alpha_P I
 \qquad(P\in\{I,X,-iY,Z\}).
 \label{eq:stabilizer-single-erasure-KL}
\end{equation}
These are Hermitian Pauli representatives; recall that \(Y=ZX\).
Conjugate \(P_r\) through a Clifford encoder and then evaluate its fixed
ancillas.  The three nonidentity expectations in
\cref{eq:stabilizer-single-erasure-KL} are either all zero or exactly one is
\(\pm1\).  Thus the input-independent marginal at \(r\) is respectively
maximally mixed or a pure one-qubit stabilizer state.  In the pure case its
Schmidt rank is one, which is \cref{eq:stabilizer-erasure-pure}.

In the maximally mixed case a Schmidt decomposition gives
\[
 V\lvert\psi\rangle
 =2^{-1/2}\bigl(\lvert0\rangle_rW_0\lvert\psi\rangle
                 +\lvert1\rangle_rW_1\lvert\psi\rangle\bigr).
\]
Applying \cref{eq:stabilizer-single-erasure-KL} to the four matrix units on
\(r\) gives \(W_c^\dagger W_d=\delta_{cd}I\).  Hence
\(
 W(\lvert c\rangle_a\otimes\lvert\psi\rangle)=W_c\lvert\psi\rangle
\)
is an isometry.  Symplectic Gaussian elimination of the encoder tableau, in
the sense fixed in \cref{subsubsec:pauli-stabilizer}, isolates
a hyperbolic row pair meeting \(r\); the remaining rows are a
stabilizer tableau for \(W\).  This proves
\cref{eq:stabilizer-erasure-mixed} without inserting the displayed
Clifford coordinate changes as gadgets.

Suppose now that no single-site erasure is correctable.  For each \(r\),
failure of the displayed Knill--Laflamme condition gives a one-site Pauli
whose pullback through the stabilizer encoder is nonscalar.  By the logical
Pauli convention in \cref{subsubsec:pauli-stabilizer}, this
is a nontrivial logical Pauli with a physical representative supported on
\(r\) alone.  Representatives on
different ports commute, so their logical Paulis commute.  Independent
physical one-qubit Cliffords send them to \(Z_r\), and one logical Clifford
sends the generated commuting logical subgroup to \(Z\)-type.  Consequently
each \(V\lvert z\rangle\) is a simultaneous eigenvector of every physical
\(Z_r\), and hence is one computational-basis ket times a phase.  Its
\(r\)-th bit is \(c_r+m_r\cdot z\); nontriviality of the logical Pauli gives
\(m_r\ne0\).  If \(\rank M<k\), some
\(0\ne d\in\ker M^{\mathsf T}\) makes \(V\lvert z\rangle\) and
\(V\lvert z+d\rangle\) proportional, contrary to
isometry.  Thus \(\rank M=k\), and the stabilizer amplitude normal form of
\cref{subsubsec:pauli-stabilizer} makes the remaining
nowhere-zero phase quadratic modulo four, proving
\cref{eq:stabilizer-column-normal-form}.
\end{proof}

\begin{lemma}[An actual tetrahedral card shortens a nonaffine encoded
residue]
\label{lem:a4rel-physical-code-shortening}
Let \(t\ge1\), let
\(V:(\mathbb C^2)^{\otimes k}\to(\mathbb C^2)^{\otimes N}\)
be a stabilizer isometry with \(k\le4\) and \(N\ge k+2\), and let
\(\xi\in(\mathbb C^2)^{\otimes(t+k)}\setminus\cA\).  Put
\(F:=(I^{\otimes t}\otimes V)\xi\).  Then \(F\) factors or one actual
tetrahedral contraction of two physical ports is nonzero and nonaffine.
\end{lemma}

\begin{proof}
Apply \cref{lem:stabilizer-erasure-column-normal-form}.  If erasure of a
code port \(r\) is correctable, its pure-marginal case
\cref{eq:stabilizer-erasure-pure} gives a genuine unary factor.  In the
maximally-mixed case \cref{eq:stabilizer-erasure-mixed} isolates an
auxiliary qubit \(a\).  Contract \(r\) with one of the \(t\) untouched
ports using the actual kernel \(Q_0\), whose class lies in
\(\mathcal T\).  In the analytic local
frame of \(r\), this acts on the untouched-port/\(a\) coefficient table by
an invertible one-qubit Clifford.  The residual encoder is a nonzero
stabilizer isometry, and a Clifford and its inverse preserve \(\cA\).
Consequently the shortened card is nonaffine.

Assume now that no one-site erasure is correctable.  Write the normal
form of \cref{eq:stabilizer-column-normal-form} explicitly as
\[
 \widetilde V=\left(\bigotimes_{j=1}^N C_j^{-1}\right)VL,
 \qquad
 \widetilde\xi=(I^{\otimes t}\otimes L^{-1})\xi,
\]
where \(L\) is the logical Clifford and the \(C_j\)'s map normal-form
output coordinates back to physical coordinates.  Then
\(\widetilde\xi\notin\cA\), since Clifford maps and their inverses
preserve \(\cA\), and
\(F=(I^{\otimes t}\otimes(\bigotimes_jC_j)\widetilde V)
\widetilde\xi\).  Let \(c\) and \(M=(m_1\ \cdots\ m_N)\) be the
data for \(\widetilde V\).  Choose a \(k\)-column basis of
\(M\).  Since \(N\ge k+2\), choose two positions \(r,s\) outside it; deleting
their columns leaves rank \(k\).  Choose matrix representatives \(C_r,C_s\)
of the two analytic local Clifford frames used in the normal form, with
\([C_r],[C_s]\in\mathcal O\).  Normality and
index two imply that
\begin{equation*}
 \{[C_r^{\mathsf T}KC_s]:[K]\in\mathcal T\}
\end{equation*}
is one of the two projective \(\mathcal T\)-cosets in \(\mathcal O\)
(the native \(X\)-twist only relabels that coset).  By
\cref{eq:a4rel-full-support-cosets,eq:a4rel-full-support-cross-ratio}, choose
an actual \(K\) with \([K]\in\mathcal T\) whose effective table
\(K_{\rm eff}=C_r^{\mathsf T}KC_s\) has full support and cross ratio
\(-1\).  The analytic frames select this physical \(K\); they are not
inserted into the network.

For a logical word \(z\), the two contracted bits are
\[
 \ell_r(z)=c_r+m_r\cdot z,\qquad
 \ell_s(z)=c_s+m_s\cdot z.
\]
There is a unique physical codeword for each \(z\), so the contraction
contains no sum of competing lifts: it multiplies its coefficient by
\[
 \omega(z)=K_{\rm eff}(\ell_r(z),\ell_s(z)).
\]
The full-support property makes \(\omega\) nowhere zero.  By
\cref{eq:a4rel-full-support-cross-ratio}, after multiplication by a common
nonzero scalar, \(\omega=i^{q(z)}\) for a modulo-four quadratic \(q\), and
\(\omega^{-1}\) has the same property.  The undeleted columns retain rank
\(k\), so the residual computational-word map is injective.  Multiplication
by \(\omega\), its inverse, and this injective stabilizer encoding preserve
affine membership in both directions, as does the normal-form phase
\(\zeta(z)\).  Since \(\widetilde\xi\notin\cA\), the actual two-port
\(K\)-card is nonzero and nonaffine.
\end{proof}

\begin{proposition}[Physical tetrahedral localization]
\label{prop:a4rel-required-localization}
Let \(g\) be an even-arity tensor-prime signature in a fixed-\(I\)
standard marked \(A_4\) state whose recomputed complete group remains
\(\mathcal T\) after factor saturation.  Assume the established terminal
exits do not apply.  If \(g\notin\cA\) and \(\arity(g)\ge10\), some actual
tetrahedral pair card of \(g\) is nonzero and nonaffine.  Iteration therefore
reaches a realized nonaffine signature on at most eight ports, unless a factor,
terminal, or strict group restart occurs, using one \(g\)-vertex and
\(O(\arity(g))\) actual binaries.
\end{proposition}

\begin{proof}
If every actual pair card were zero or affine, tomography and
\cref{lem:a4rel-four-active-qubits} would place each card-map image in a
stabilizer code with \(k\le4\) logical qubits.  Frobenius reconstruction
would then write \(g=(I^{\otimes2}\otimes V)\xi\), with \(V\) encoding
those \(k\) qubits into \(N=\arity(g)-2\ge8\) sites and
\(\xi\notin\cA\).  Since \(N\ge k+2\),
\cref{lem:a4rel-physical-code-shortening} gives a factor or a nonaffine
actual card, a contradiction.  Iterating this two-port descent reaches
arity at most eight unless a terminal, factor, or recomputed-group restart
occurs.  The direct-gadget and factor-saturation directions are
\cref{eq:direct-gadget-reduction,lem:factor-saturation}.
\end{proof}

\paragraph{Section-local occurrence induction.}
We close the relative tetrahedral argument by a finite induction on
unresolved occurrences and their remaining contractions.  On entering the argument,
dispatch nonzero odd signatures by \cref{thm:external-odd}, unaries by
\cref{lem:common-unary-exit}, nullaries by preprocessing, and newly exposed
binaries by
\cref{lem:binary-interpolation,lem:actual-binary-completeness}.  Thus, on a
continuing branch, every unresolved nonbinary tensor-prime occurrence has
even arity and every binary transfer belongs to the recomputed group
\(\mathcal T\).

Let \(\mathcal P\) be the finite multiset of unresolved nonbinary
tensor-prime occurrences.  When a token \(u\), carrying \(g_u\), is
selected, freeze the finite current-coordinate universe consisting of all
iterated proper physical tetrahedral pair cards of \(g_u\), all actual
dressings and port permutations used by the finite reductions below, and the
complete factorization tree of every output.  This universe is finite: every
proper contraction or factor lowers arity, \(\mathcal T\) and the reduction
list are finite, each arity has finitely many port choices and dressings,
and each tensor has a finite prime-factorization tree.  Write its remaining keys as \(Q_u\)
and let \(\varphi_u\in\mathbb N\) be the sum of their factor potentials from
\cref{eq:factor-potential}.  Define
\begin{equation}
 \Theta_{\rm Pl}(\mathcal P)
 =\{\!\{(\arity(g_u),|Q_u|,\varphi_u):u\in\mathcal P\}\!\},
 \label{eq:platonic-occurrence-rank}
\end{equation}
using lexicographic order on triples and its well-founded multiset
extension.

Before applying the next induction step to \(u\), regenerate its complete
oriented current-group deck.  Process each proper card and factor, either
resolving it or entering it as a lower-arity occurrence with smaller
\(\Theta_{\rm Pl}\).  The queue \(Q_u\) then contains only
unconsumed finite reduction or certificate keys, not unresolved proper
signatures.  Before a proper card is returned, consume its unique
originating key in \(Q_u\) and attach that key, marked completed, to the
child's provenance; it is not copied into the child's unfinished key set.
If the card factors, its finite factorization tree is processed under that same
originating key, and every intermediate split strictly lowers the stored
factor potential.

Processing \(u\) therefore replaces its triple by a finite multiset of
child triples.  Every proper child has smaller first coordinate; if a
residual copy of \(u\) remains, it has the same arity but one fewer frozen
card/factor key or smaller factor potential.  Thus every inserted
triple is strictly smaller than the removed parent triple in lexicographic
order.  Certification removes the triple altogether.  The standard
multiset extension therefore makes every same-coordinate local step strictly decrease
\(\Theta_{\rm Pl}\); a terminal, lower \(\nu\), or strict complete-group
enlargement is instead an outer exit.

\begin{lemma}[The quaternary case for Platonic groups]
\label{lem:platonic-prepared-q4}
Let \(\Lambda\) be the current retained factor-saturated fixed-\(I\)
signature set with recomputed marked group
\[
 G\in\{\mathcal T,\mathcal O,\mathcal O',\mathcal O_{dt},
          \mathcal O_{td},\mathcal I,\mathcal I_-\}.
\]
Let \(q\in\Lambda\) be a nonzero tensor-prime quaternary occurrence whose
complete oriented deck and all proper cards and factors have been processed
as above.
For each canonical pair \(p\), put
\(L_p:=\mathsf R_X\circ\Phi_p^q\).  Suppose
\[
 L_p(K_g)=0\quad\text{or}\quad [L_p(K_g)]\in G
 \qquad
 (K_g\in\mathcal K_G^{\rm act})
\]
for all six canonical pairs.  Then \(q\) has exactly the factor, generalized
equality, endpoint-nondegenerate eight-vertex, or outside-binary outcomes of
\cref{lem:platonic-quaternary-boundary}; in particular, no tensor-prime
quaternary occurrence continues in the internal tetrahedral or either
icosahedral form.
\end{lemma}

\begin{proof}
The proof of \cref{lem:platonic-quaternary-boundary} uses deck stability
only to supply the displayed pair conditions and to exclude an unresolved
proper card or factor.  Here the pair conditions are assumed explicitly,
and the preceding steps have already resolved or entered every such proper object at
a smaller value of \(\Theta_{\rm Pl}\).  The two-sided finite preserver
classification and all of its actual terminal gadgets therefore apply
verbatim.
\end{proof}

\begin{lemma}[Stable tetrahedral core augmentation]
\label{lem:tetrahedral-stable-HR-augmentation}
Assume \cref{eq:platonic-six-four-predicate} with \(G=\mathcal T\), and suppose the full-rank
analysis reaches an orientation-preserving tensor \(f_\sigma\),
\(\sigma\in\Sigma_+\), by reversible actual tetrahedral dressings and a port
permutation.  Then nonzero scalars \(\alpha,\beta\) are recorded such
that the literal augmentation
\[
 \Lambda_{\rm HR}:=
 \Lambda\cup\{\alpha\Hcore,\beta\RMcore\}
\]
satisfies \(\KHolant(\Lambda_{\rm HR})\equivT\KHolant(\Lambda)\).  Close
the finite core batch and recompute its group.  On an unchanged
\(\mathcal T\)-branch, regenerate the complete physical deck of
\(\alpha\Hcore\), retain and factor-saturate every card and its factorization tree, dispatch all outputs, and recompute the group and deck once more.
Every branch is then an outer exit or has a retained residue
\(\Lambda_{\rm HR}^\star\) with
\[
 \Stable(\Lambda_{\rm HR}^\star,\mathcal T,\alpha\Hcore),
 \qquad
 \{\alpha\Hcore,\beta\RMcore\}\subseteq\Lambda_{\rm HR}^\star.
\]
Only this stable residue is passed to \cref{thm:a4-relative}.
\end{lemma}

\begin{proof}
The reversible identities \cref{eq:tetrahedral-H6-equivalence} produce
\(\alpha\Hcore\), and four copies followed by the forward circuit
\cref{eq:app-v4-H6-R8-output} produce \(\beta\RMcore\).  Direct realization
and literal inclusion give the displayed Turing equivalence.  First
preprocess and jointly factor-saturate the two cores, dispatch every
terminal output, and recompute the complete group.  A terminal, lower
minimum arity, or strict complete-group enlargement is an outer exit.
Otherwise the group remains \(\mathcal T\) and the minimum nonbinary arity
remains six.

Now regenerate the complete oriented physical deck of \(\alpha\Hcore\).
Retain every card with its ordered kernel and scalar provenance, process
its entire finite factorization tree, dispatch all resulting binaries and
terminals, and recompute the complete group.  The core \(\Hcore\) is
tensor-prime by \cref{proofsubsec:v4-hcore}; its pair cards are matching
products because this holds for \(f_\sigma\) by
\cref{lem:tetrahedral-orientation} and the displayed equivalence is an
actual reversible \(\mathcal T^6\rtimes S_6\)-dressing.
\Cref{lem:platonic-matching-contraction-closure} then makes every iterated
proper physical card zero or a safe tetrahedral matching product.  Thus
any new singular or outside binary, lower-arity nonbinary factor, or group
change is an outer exit.  On the unchanged branch, regenerate the deck
once more.  It is complete, all of its cards and factors are retained, and
every card is zero or a safe \(\mathcal T\)-matching product; hence
\cref{def:deck-stable-state} gives the asserted stable tuple.  The two
cores are certified affine signatures and are excluded from every subsequent
unresolved-occurrence set; their augmentation does not create a fresh
fixed-\((H,\nu)\) token.
\end{proof}

\begin{theorem}[Tetrahedral relative rigidity]
\label{thm:a4-relative}
Assume
\[
 \Stable(\Lambda,\mathcal T,f),\qquad
 \{\alpha\Hcore,\beta\RMcore\}\subseteq\Lambda,\qquad
 \alpha\beta\ne0.
\]
Then \(\Close_{\{\cA\}}(\Lambda,\mathcal T,f)\).
\end{theorem}

\begin{proof}
Let \(\mathcal P\) contain every unresolved nonbinary tensor-prime
occurrence of \(\Lambda\); the two displayed cores are certified affine
signatures, not unresolved tokens.  Prepare a selected occurrence and induct
on \(\Theta_{\rm Pl}(\mathcal P)\).

Arity four is \cref{lem:platonic-prepared-q4}.  For the next step, assume
the complete oriented deck and all proper cards and factors have been
processed as above, and consider a six-port
occurrence \(g\).  Its quaternary cards have already been closed, so every
continuing one is a product of two safe tetrahedral binaries.  For a
canonical pair \(p\), recompute \(\rho=\rank\Phi_p^g\).  Rank zero gives
\(g=0\), rank one factors, and ranks two or three are closed by
\cref{lem:platonic-deficient-rank}.  At rank four, fixed residual matchings
factor by \cref{lem:platonic-global-ruling}; otherwise clauses~\textnormal{(i)--(iii)}
of \cref{thm:cert-interface-platonic-incidence}, together with
\cref{lem:platonic-rainbow-localization,lem:tetrahedral-bridge,%
lem:tetrahedral-orientation}, give either the resolved
endpoint-nondegenerate eight-vertex leaf or an
orientation-preserving \(f_\sigma\).  The latter lies in \(\mathscr H_6\)
by \cref{eq:tetrahedral-H6-equivalence} and is affine by
\cref{lem:tetrahedral-affine}.  Because \(g\) already has six ports, this
rainbow localization closes no common-factor ports: \(f_\sigma\) differs
from \(g\) only by a nonzero scalar, a port permutation, and reversible
actual \(\mathcal T^6\)-dressings.  Affine membership and the
\(\mathscr H_6\)-orbit conclusion therefore pull back to \(g\).  Thus every continuing six-port occurrence
lies in \(\mathscr M_6\cup\mathscr H_6\).

Let \(g\) have arity eight.  Prepare every physical tetrahedral pair card.
For a factored six-port card, process its factorization tree under the
originating \(g\)-key; a continuing result consists only of safe
tetrahedral binaries and hence lies in \(\mathscr M_6\).  Before a
tensor-prime six-port card is entered as a child occurrence, consume its
originating \(g\)-key.  Replacing the parent triple by the residual
eight-port triple with one fewer key and the six-port child triple is a
strict multiset decrease of \(\Theta_{\rm Pl}\).  The preceding six-port
base then places that child in \(\mathscr H_6\).  Therefore all pair cards
lie in \(\mathscr C_6\), and
\cref{prop:a4rel-required-eight-port} says that \(g\) factors or is affine.

Finally let \(\arity(g)\ge10\).  If \(g\notin\cA\),
\cref{prop:a4rel-required-localization} gives a terminal, factor, strict
group successor, or a nonzero nonaffine actual card of smaller arity.  The
first three are outer exits or decrease the current token.  In the last
case the parent card key is consumed before its factor-saturated child is
entered, so \(\Theta_{\rm Pl}\) strictly decreases.  Induction makes every
child factor affine; their tensor product would then be affine, contradicting
the nonaffine card.  Hence \(g\) is affine.

After every retained augmentation recompute the complete group.  A strict
change is an outer exit; otherwise prepare the next occurrence from its
regenerated current-group deck.  Every continuing step removes or
strictly lowers a triple of \(\Theta_{\rm Pl}\), so every tensor-prime
occurrence is eventually affine.  The native edge and all twelve
tetrahedral binaries are affine in this basis by
\cref{eq:tetrahedral-domain}; \(\Hcore\) is affine by
\cref{proofsubsec:v4-hcore}, and \(\RMcore\) is affine because
\(RM(1,3)\) in \cref{eq:RM-core} is a linear code.  Equivalently their
literal relationship is certified by
\cref{lem:tetrahedral-affine,eq:tetrahedral-H6-equivalence,%
eq:app-v4-H6-R8-output}.  Tensor closure, the normalization data, and
\cref{def:common-K-presentation} give \(\PresK(K\Lambda;\cA)\), proving the
claimed relative closure.
\end{proof}

\begin{theorem}[Tetrahedral deck rigidity]
\label{thm:tetrahedral-deck-rigidity}
If \cref{eq:platonic-six-four-predicate} holds with \(G=\mathcal T\), then
\(\Close_{\{\cA\}}(\Lambda,\mathcal T,h)\).
\end{theorem}

\begin{proof}
Varying card matchings give an actual rainbow localization by
\cref{lem:platonic-rainbow-localization}.  The bridge and orientation split
then give either the resolved endpoint-nondegenerate eight-vertex leaf or an
orientation-preserving tensor \(f_\sigma\).  Fixed matchings expose a proper
factor by \cref{lem:platonic-global-ruling,lem:factor-saturation}; closing
its finite factorization tree gives a resolved leaf, lower \(\nu\), or strict
complete-group successor, as in the octahedral proof.

On the orientation-preserving branch, apply
\cref{lem:tetrahedral-stable-HR-augmentation}.  Every outer-exit branch is
already closed.  On its sole stable \(\mathcal T\)-residue,
\cref{thm:a4-relative} gives the common affine presentation, and the
recorded Turing equivalence transports that conclusion back to \(\Lambda\).
The bridge rank throughout is the rank recomputed on the directly realized
six-port signature, never an inherited ancestor rank.  This proves \(\Close_{\{\cA\}}(\Lambda,\mathcal T,h)\).
\end{proof}

\subsection{The external \texorpdfstring{\(A_4\)}{A4} form}

The external form has the same domain incidence as \(\mathcal T\), but
its matching deck has no rainbow line.  Its four-port case follows from
the two-sided matrix conditions below; at higher arity, connected domain
incidence forces all cards to share one matching and hence a proper factor.

\begin{proposition}[External tetrahedral two-sided preservers]
\label{prop:external-a4-q4-preservers}
For a projective matrix set \(\mathcal S\), define its projective
coefficient set by
\[
\mathcal R_{\mathcal S}
 :=\{[c_Q(K)]:[K]\in\mathcal S\}.
\]
Coordinate matrices act on these projective sets representative-wise.
Let
\(\Theta_{\rm ext}:=c_Q\circ\theta_{\rm ext}\circ c_Q^{-1}\), define the
Frobenius-form matrix
\(J_{\rm Fr}:=\diag(1,-1,1,-1)\), and define the associated adjoint map by
\(A^\sharp:=J_{\rm Fr}A^{\mathsf T}J_{\rm Fr}\).  A coordinate linear map
\(A\) is a \emph{two-sided preserver} of
\(\mathcal T_{\rm ext}\) if, for every
\([r]\in\mathcal R_{\mathcal T_{\rm ext}}\),
\begin{equation}
 \begin{aligned}
  &Ar=0\quad\text{or}\quad[Ar]\in\mathcal R_{\mathcal T_{\rm ext}},\\
  &A^\sharp r=0\quad\text{or}\quad
    [A^\sharp r]\in\mathcal R_{\mathcal T_{\rm ext}}.
 \end{aligned}
 \label{eq:external-a4-two-sided-definition}
\end{equation}
This is representative-independent.  There is one zero map; up to nonzero
scalar, the nonzero maps have rank distribution
\((144,0,0,576)\) in ranks \((1,2,3,4)\), respectively.  Under the
Frobenius-dual reconstruction
\cref{eq:platonic-frobenius-reconstruction}, every rank-one or rank-four
map reconstructs a tensor with a binary-pair factorization.
\end{proposition}

\begin{proof}
The adjoint is taken for the bilinear Frobenius form, without complex
conjugation.  Since \(R_{\rm ext}^{\mathsf T}=R_{\rm ext}\), conjugation
by \(R_{\rm ext}\) is
self-adjoint, so \(\Theta_{\rm ext}^\sharp=\Theta_{\rm ext}\).  Moreover,
\[
 [R_{\rm ext}^2]=[Q_0+Q_1]\in N_{\PGL_2}(\mathcal T),
 \qquad
 \Theta_{\rm ext}^2(\mathcal R_{\mathcal T})
 =\mathcal R_{\mathcal T}.
\]
Put \(A'=\Theta_{\rm ext}^{-1}A\Theta_{\rm ext}\).  If
\([r]\in\mathcal R_{\mathcal T}\), then
\([\Theta_{\rm ext}r]\in\mathcal R_{\mathcal T_{\rm ext}}\), so the first
condition for \(A\) implies the standard first condition for \(A'\).  Also
\[
 (A')^\sharp=\Theta_{\rm ext}A^\sharp\Theta_{\rm ext}^{-1}.
\]
Because
\(\Theta_{\rm ext}^2(\mathcal R_{\mathcal T})
=\mathcal R_{\mathcal T}\), one has
\[
 \Theta_{\rm ext}^{-1}(\mathcal R_{\mathcal T})
 =\Theta_{\rm ext}(\mathcal R_{\mathcal T})
 =\mathcal R_{\mathcal T_{\rm ext}};
\]
the second condition transports in the same way,
and the converse follows by reversing the calculation.  Thus the external
condition is conjugate to the standard tetrahedral two-sided condition,
equivalently the \(A_4\) row of
\cref{thm:cert-interface-platonic-q4}\textnormal{(i)}.  Rank-one reconstruction gives a
binary-pair factorization.  In rank four, for some
\(U,V\in\GL_2(\mathbb C)\), the standard map is \(K\mapsto UKV\) or
\(K\mapsto UK^{\mathsf T}V\).  Since
\(A=\Theta_{\rm ext}A'\Theta_{\rm ext}^{-1}\), transport gives,
respectively,
\[
 K\mapsto\theta_{\rm ext}(U)K\theta_{\rm ext}(V),
 \qquad
 K\mapsto\theta_{\rm ext}(U)R_{\rm ext}^{-2}K^{\mathsf T}
 R_{\rm ext}^2\theta_{\rm ext}(V).
\]
Either display reconstructs two crossing nonsingular binaries.  These are
analytic tensor identities, not actual gadgets.
\end{proof}

\begin{theorem}[External tetrahedral deck rigidity]
\label{thm:external-a4-deck-rigidity}
Let \((\Lambda,\mathcal T_{\rm ext},f)\) be a deck-stable retained,
factor-saturated state containing the distinguished literal \(I\) and
\(\arity(f)=\nu(\Lambda)\).  Then an established terminal, a
group-enlargement restart, or a genuine proper factor of \(f\) occurs.
Thus this case has no tensor-prime survivor or common-affine leaf.
\end{theorem}

\begin{proof}
At arity four, use the transfer-valued pair maps
\(L_p=\mathsf R_X\circ\Phi_p^f\).  The exact
complementary-pair adjoint identity
\cref{eq:platonic-q4-map-adjoint,eq:platonic-q4-adjoint} gives the
representative-wise two-sided condition
\cref{eq:external-a4-two-sided-definition}.  By
\cref{prop:external-a4-q4-preservers}, rank zero makes \(f=0\), ranks one
and four factor, and ranks two and three do not occur.

For \(\arity(f)\ge6\), fix a pair and let \(\Phi\) be its card map.  The
actual kernel set \(\mathcal K_{\mathcal T_{\rm ext}}^{\rm act}\) spans
\(\Mat_2\), so rank zero is impossible and rank one
factors.  In ranks two through four, the kernel-deleted rich-line graph is
connected by
\cref{thm:cert-interface-platonic-incidence}\textnormal{(iii)}.  A change of
matching along an edge would give a kernel-avoiding rich domain line with
three distinct images.  Their matchings are pairwise distinct: if two
repeated a matching, flattening across one of its edges absent from the
third would give ranks \(1,1,4\), contradicting collinearity.
Apply \cref{lem:platonic-rainbow-localization} before using the quaternary
incidence classification.  Each common factor has transfer in
\(\mathcal T_{\rm ext}\), and the actual kernels span \(\Mat_2\), so
actual closures preserve all three images and produce a quaternary
\(\mathcal T_{\rm ext}\)-rainbow line.  This contradicts
\cref{thm:cert-interface-platonic-incidence}\textnormal{(i)}.
Thus all nonzero cards use one matching \(M\).

Now \(\Psi=\Phi\circ\theta_{\rm ext}\) maps the twelve standard test points
to zero or an \(M\)-matching product.
\Cref{lem:platonic-global-ruling} factors \(\Psi\), and composing
its varying factor with \(\theta_{\rm ext}^{-1}\) pulls the tensor identity,
not a gadget, back to \(\Phi\).  Since \(|M|\ge2\), dual reconstruction and
\cref{lem:factor-saturation} expose a proper factor; its binaries give the
stated terminal or restart.
\end{proof}

\subsection{The \texorpdfstring{\(A_5\)}{A5} form}

Recall \(\phi=(1+\sqrt5)/2\), \(\tau=\phi^{-1}=\phi-1\), and the
standard sixty-point model \(\mathcal I\) from
\cref{eq:icosahedral-domain}.  Its axes, tetrahedral sign points, and
additional even-permutation points contribute \(4,8,48\) classes.
\Cref{thm:cert-interface-platonic-bridges}\textnormal{(iii)} proves directly
that these points are closed under
multiplication, inverse, and transpose.

For the opposite display, use the fixed-equality card and network dictionary
\cref{lem:icosahedral-fixed-I-dictionary}.  It replaces every standard
kernel by its actual \(\kappa_D\)-image while leaving the retained state and
the supplied literal \(I\) unchanged.

Use the common rainbow normalization and separated-sign system
\cref{eq:platonic-rainbow-normalization,eq:platonic-rainbow-sum,%
eq:platonic-separated-sign-tests} with \(G=\mathcal I\).

\begin{lemma}[Icosahedral bridge and uniform separator]
\label{lem:icosahedral-bridge}
\label{lem:icosahedral-separator}
The separated-sign system for normalized maps
\(\Phi:\Mat_2(\mathbb C)\to(\mathbb C^2)^{\otimes4}\) with
\(G=\mathcal I\) has exactly six rank-deficient solutions and thirty full-rank
solutions.  For each full-rank solution \(\Phi\), the actual kernel
\begin{equation*}
  K_{\mathcal I}=\phi Q_1-\tau Q_2+Q_3,
  \qquad [K_{\mathcal I}]\in\mathcal I
\end{equation*}
realizes the quaternary \(q_{\mathcal I}:=\Phi(K_{\mathcal I})\), whose
three flattening ranks are \((3,3,3)\), and the actual kernel on
the canonical pair \((1,2)\)
\begin{equation*}
  L_{\mathcal I}=Q_0-\tau^2Q_1+\tau Q_3,
  \qquad [L_{\mathcal I}]\in\mathcal I
\end{equation*}
gives the nonsingular sequential contraction
\(B_{\rm out}:=\Phi_{(1,2)}^{q_{\mathcal I}}(L_{\mathcal I})\), with
\([B_{\rm out}X]\notin\mathcal I\) in every case.
For \(\mathcal I_-\), the corresponding actual kernels are
\[
 K_{\mathcal I_-}:=\kappa_D(K_{\mathcal I}),\qquad
 L_{\mathcal I_-}:=\kappa_D(L_{\mathcal I}).
\]
They give the covariant quaternary and a nonsingular binary whose transfer
lies outside \(\mathcal I_-\), directly in the fixed-\(I\) coordinates.
\end{lemma}

\begin{proof}
The exact line buckets, bridge, five frames, representative-wise thirty-case
coverage, and arbitrary scale and cancellation loci are summarized in
\cref{thm:cert-interface-platonic-bridges}\textnormal{(iii)}.
The same interface supplies the uniform actual outsider.  Thus
\cref{eq:direct-gadget-reduction} gives the required direction;
\cref{lem:binary-interpolation,lem:actual-binary-completeness} dispatches
infinite order and makes finite order contradict completeness.
\Cref{lem:icosahedral-fixed-I-dictionary} gives the stated direct physical
certificate for the other display.
\end{proof}

For a directly realized six-port localization whose recomputed card-map rank is
four, varying matchings and connected incidence supply a mixed line.  It
cannot be a mixed five-line by
\cref{thm:cert-interface-platonic-incidence}\textnormal{(i)}, so its rainbow three-line and
\cref{lem:icosahedral-separator} give the uniform outsider.  If instead all
nonzero cards use one matching \(M\), the twelve test points map to zero or
an \(M\)-matching product;
since \(|M|=2\),
\cref{lem:platonic-global-ruling,lem:factor-saturation} expose a proper
factor.

\begin{theorem}[Icosahedral deck rigidity]
\label{thm:icosahedral-deck-rigidity}
If \cref{eq:platonic-six-four-predicate} holds and
\[
 G\in\{\mathcal I,\mathcal I_-\},
\]
then \(\Close_{\varnothing}(\Lambda,G,h)\).
\end{theorem}

\begin{proof}
For \(G=\mathcal I\), connected incidence and
\cref{lem:platonic-rainbow-localization,lem:icosahedral-separator} give
the uniform outsider whenever card matchings vary.  For
\(G=\mathcal I_-\), the same argument is direct after replacing every
kernel and finite network identity by its image under
\cref{lem:icosahedral-fixed-I-dictionary}.  The literal equality remains
unchanged.

If matchings are fixed,
\cref{lem:platonic-global-ruling,lem:factor-saturation} exposes a proper
factor.  Close its finite factorization tree: a nonbinary factor lowers arity, a
singular or unsafe binary is terminal, and a finite safe binary whose
transfer class lies outside the current group strictly enlarges that group.
All-safe current-group binary
factors would contradict the tensor-prime branch.  Thus every branch is a
resolved leaf or strict outer successor, proving the conclusion without a
global termination lemma.  The separator uses only the rank recomputed on
the actual six-port coefficients.
\end{proof}

\begin{proof}[Proof of \cref{thm:platonic-decks}]
Work at the stable tuple in the theorem.  First dispatch a nonzero odd
signature by \cref{thm:external-odd}, a unary by
\cref{lem:common-unary-exit}, a nullary factor by preprocessing, and every
new binary by
\cref{lem:binary-interpolation,lem:actual-binary-completeness}.  Thus a
continuing branch is all-even and every binary transfer lies in the
recomputed complete group.

If \(G=\mathcal T_{\rm ext}\), apply
\cref{thm:external-a4-deck-rigidity}.  Its factor branch is closed by
retained factor saturation.  For every other Platonic form, including
\(\mathcal I_-\), use the direct coefficient interfaces already proved
above.  Hence assume
\[
 G\in\{\mathcal T,\mathcal O,\mathcal O',\mathcal O_{dt},
          \mathcal O_{td},\mathcal I,\mathcal I_-\}.
\]
At arity four,
\cref{lem:platonic-quaternary-boundary,%
lem:p1-platonic-eight-vertex-rank-terminal,%
cor:higher-equality-unary,lem:common-unary-exit}
and the local-reduction conclusion following
\cref{lem:platonic-finite-q4-terminal} close every output.  Suppose
\(n:=\arity(f)=\nu(\Lambda)\ge6\).

Fix a canonical pair \(p\) and compute \(r=\rank\Phi_p^f\) from the actual
coefficients.  Rank zero gives \(f=0\) by
\cref{eq:platonic-frobenius-reconstruction}.  At rank one there are nonzero
tensors \(U,F_{\rm res}\) with
\[
 \Phi_p^f(K)=\operatorname{tr}(U^{\mathsf T}K)F_{\rm res},
 \qquad f=U_p\otimes F_{\rm res},
\]
so \cref{lem:factor-saturation} gives a proper factor route.  Closing that
finite factorization tree either lowers \(\nu\), reaches a binary terminal or
strict group successor, or expresses \(f\) entirely as safe current-group
binaries, contradicting its tensor-prime choice.

Let \(r\in\{2,3,4\}\).  Deck stability gives \cref{eq:platonic-all-card-condition}, so
\cref{lem:platonic-high-arity-localization} applies.  The
fixed-matching route again gives a proper factor.  Close its finite factorization tree by the same alternatives; an all-safe current-group binary product
again contradicts tensor primality.  The mixed route
materializes a nonzero
six-port \(h\), using one \(f\)-vertex and exactly \((n-6)/2\) actual
kernels, with
\[
 2\le\rank\Phi_p^h\le r
\]
and \cref{eq:platonic-all-card-condition}.  Form the literal augmentation
\(\Lambda_h:=\Lambda\cup\{h\}\).  Direct realization and literal inclusion
give
\[
 \KHolant(\Lambda_h)\leT\KHolant(\Lambda),
 \qquad
 \KHolant(\Lambda)\leT\KHolant(\Lambda_h).
\]

Preprocess and factor-saturate this entire finite batch, dispatch its newly
exposed terminal outputs, and recompute the complete group.  If \(h\) factors, its
selected varying matchings exclude a product entirely of safe current-group
binaries, so a lower-arity nonbinary factor or an earlier terminal occurs.
If \(n>6\), a retained six-port prime lowers \(\nu\); a strict
complete-group enlargement is likewise an outer successor.  If \(n=6\),
the localization has no closing kernels: the contraction \(C\) in the
proof of \cref{lem:platonic-high-arity-localization} is the identity, and
\(h\) is only the recorded port permutation and nonzero scalar multiple of
\(f\).  We identify it with the same occurrence and reuse its frozen
card/factor key; no fresh fixed-\((H,\nu)\) occurrence is inserted.

On the only remaining prime six-port branch with unchanged marked group,
regenerate the complete oriented physical deck of \(h\).  Retain every card
with its ordered kernel, scalar, and predecessor provenance; factor-saturate
its complete finite factorization tree; dispatch every resulting binary and
terminal; and recompute the complete group.  A lower-arity factor or group
change is an outer successor.  On the unchanged branch regenerate the deck
again.  Condition \cref{eq:platonic-all-card-condition} and
\cref{lem:platonic-matching-contraction-closure} make every card zero or a
safe current-group matching product, and all card factors are retained.
Consequently \cref{def:deck-stable-state} gives the stable tuple
\[
 \Stable(\Lambda_h^\star,G,h),
 \qquad \arity(h)=\nu(\Lambda_h^\star)=6.
\]

Only on this stable residue recompute
\(\rho_h:=\rank\Phi_p^h\).  If \(\rho_h\le3\), apply
\cref{lem:platonic-deficient-rank}.  If \(\rho_h=4\), apply, according to
the recomputed marked group,
\cref{thm:tetrahedral-deck-rigidity,thm:octahedral-deck-rigidity,%
thm:icosahedral-deck-rigidity}.  In the internal tetrahedral
orientation-preserving case,
\cref{lem:tetrahedral-stable-HR-augmentation} first closes the literal
\(\Hcore,\RMcore\) augmentation and re-establishes a stable
\(\mathcal T\)-tuple; only then does \cref{thm:a4-relative} supply the one
signature-set-wide affine presentation.  Its common basis is certified by
\cref{eq:tetrahedral-domain,proofsubsec:v4-hcore,eq:RM-core,%
eq:tetrahedral-H6-equivalence,eq:app-v4-H6-R8-output}.

These reductions exhaust the stable six-port residue.  The octahedral and
icosahedral forms leave no rigidity survivor, while the internal
tetrahedral form leaves only the stated common affine class.  No second
prime worklist is needed: the internal relative theorem already classifies
the retained set, and the other marked forms have no common-presentation
leaf.  Each construction records its port order and nonzero scalar, and its
reduction points to the preceding retained problem.  The finite bridges
have constant size; localization uses \((n-6)/2\) kernels.  These bounds
and the local occurrence induction complete the proof.
\end{proof}
\section{Synthesis of the Binary-Disequality Dichotomy}
\label{sec:p1-synthesis}

We now prove the two dichotomies stated in the introduction.  First,
\cref{thm:matching-deck-classification} combines the four group branches
into a terminating classification of an equality-accessible native-\(X\)
problem.  \Cref{cor:equality-accessible-closure} applies that classification
to the entire retained set.  The final two proofs then remove the
factor-saturation hypothesis and use the equality-wire equivalence to
recover the ordinary binary-disequality formulation.

Every continuing signature set \(\Lambda\) in this section contains \(I\),
so \(Z=K^{\otimes2}I\in K(\Lambda)\).  Hence
\cref{lem:tract-X-anchor} gives
\[
  \Tract(K(\Lambda))\quad\Longleftrightarrow\quad
  \TractX(K(\Lambda)).
\]
This equivalence identifies the branchwise criterion \(\Tract\) with the
equality-accessible criterion \(\TractX\).

\subsection{Synthesis of the Equality-Accessible Matching-Deck Classification}
\label{proofsec:matching-deck-synthesis}

\begin{lemma}[Descent in the four group cases]
\label{lem:token-coverage}
For every deck-stable tuple \((\Lambda,G,f)\), work in the fixed coordinate matrix \(A=I\), put
\[
 H=\overline G(I,\Lambda)=G(\Lambda),
 \qquad
 \mathscr U^{\mathrm{fac}}(I,f,H)=Q_{\rm done}\sqcup Q,
\]
and freeze the displayed partition.  Then the unfinished-token batches and
the further reductions invoked by the active-core completion proofs in
\cref{sec:normalized-dihedral,sec:proper-deck,sec:v4-deck,sec:platonic-decks}
satisfy \cref{def:frozen-token-coverage}.
\end{lemma}

\begin{proof}
Fix a deck-stable tuple \((\Lambda,G,f)\) in the fixed-\(I\) coordinates and
use the displayed frozen partition.  By construction, every oriented
one-copy proper card of \(f\), with every assignment of frozen actual
representatives, has a canonical card token in \(\mathscr U(I,f,H)\).
Factor saturation attaches exactly its finite factorization tree.  Consequently,
a same-\((H,\nu)\) continuation produced by a one-copy card is confined to
one token of \(Q\).  Closing its stored card and full factorization tree moves
that token to \(Q_{\rm done}\); during any intermediate genuine factor
split, its stored factor potential strictly decreases.

It remains to check the reductions beyond the frozen one-copy deck.
In the standard marked-dihedral branch, the active-core proof of
\cref{thm:normalized-dihedral} uses an actual tensor-prime eight-vertex
output at arity four, closed for the whole retained set by
\cref{thm:p1-quaternary-kholant-dichotomy}.  At arity six,
\cref{lem:nd-six-prime-base} excludes the stable prime.  At arity at least
eight, disagreeing sampled cards would localize by
\cref{lem:app-nd-localization} to a nonmatching six-port tensor, although
the first proper closure and all later monomial contractions preserve
matching products.  Thus the samples agree; their common factor lines
and the two invertible pencils give a genuine factorization of the active
core, again contradicting primality.  Thus the dihedral construction closes its selected occurrence by a
local factorization or a terminal for the retained set.

For the \(\mathcal B_m\) form, the local identities
\[
 U_{\mathcal B}^{\mathsf T}XU_{\mathcal B}=-2iI,
 \qquad
 U_{\mathcal B}^{-1}IU_{\mathcal B}^{-\mathsf T}=\frac{i}{2}X,
 \qquad
 U_{\mathcal B}^{-1}XU_{\mathcal B}^{-\mathsf T}=\frac{i}{2}I
\]
induce an open-network bijection with the \(\mathcal M_m\) deck.  It
preserves arity, factorization, ordered actual cards, and their factorization trees; hence it transports the frozen partition bijectively, and the preceding local reduction applies to the transported card and factor
tokens.  For \(\mathcal A_m^0\), the
literal-\(I\) identities \cref{eq:marked-A0-fixed-I,eq:marked-A0-edge-switch}
and the deck relabelling in the proof of
\cref{thm:marked-A0-reduction} give the same bijective transport to the
\(\mathcal M_m\) reductions, including their frozen card and factor tokens;
the transport is supplied by the displayed local identities.
For \(\mathcal A_m\), the weighted kernel identities and two-card
localization in \cref{lem:marked-A-gauge-transport} give the same active-core
argument, as detailed in the proof of \cref{thm:marked-A-rigidity}:
the six-port prime is excluded, and at higher arity the common factor
lines and weighted pencils factor the core.  Its arity-four case is the
physical whole-set reduction in \cref{lem:marked-A-q4}.  Neither transport
leaves an unclassified fixed-\((H,\nu)\) resource.

In the proper-cyclic branch, the \(C_2\) boundary is the common quaternary
method \cref{thm:equality-q4-boundary}; its endpoint output is directly realized
and sent to \cref{cor:p1-eight-vertex-reduced-interface}.  For \(C_N\),
\(N\ge4\), the
independent phase and dressing calculations
\cref{lem:odd-cyc-q4-family,lem:odd-cyc-four-port-dressing,%
eq:cyclic-dressing-image} give either an actual tensor-prime
endpoint-nondegenerate eight-vertex signature or a nonsingular
anti-diagonal binary factor of a direct two-copy output, as in
\cref{lem:proper-cyclic-q4}.  The former is directly realized,
factor-saturated, and closed by
\cref{thm:p1-quaternary-kholant-dichotomy}.  The latter is Turing-exposed
by retained factor saturation; only in that augmented set is it available
as a binary signature.  After the full unsafe-binary dispatch, its safe
transfer class lies outside \(C_N\) and strictly enlarges the recomputed
complete group, which contains the dihedral subgroup generated by the old
cyclic group and this involution.  Higher in the Reed--Muller
argument, \cref{lem:rm-weight-two-integration} gives either a proper factor
or an actual \(\mathrm{Even}_4\); the latter is consumed by the explicit
equality anchor in \cref{lem:equality-anchor}.  Thus every cyclic
multi-copy reduction yields a resolved leaf, a lower-arity factor, or a
strict complete-group successor.

For the standard Pauli branch, arity four is the common quaternary method
followed by \cref{cor:p1-eight-vertex-reduced-interface}.  At higher
arity, \cref{thm:standard-v4-relative} either resolves the whole retained
set through the real-valued dichotomy or supplies one actual quaternary
batch.  Its complete factor processing yields a terminal, a strict group
enlargement, or a decrease of \(\nu\); it has no unresolved same-state
return.  In the exotic branch, arity four is the
common quaternary method followed by
\cref{cor:p1-eight-vertex-reduced-interface}, while the remaining finite
reductions are
\cref{lem:exotic-v4-q6,lem:exotic-v4-q8,%
prop:exotic-v4-high-arity-factor}.  Their proofs use the explicit six- and eight-port classifications and
the higher-arity factorization.  The tetrahedral base construction is
closed by the finite arity induction in \cref{thm:a4-relative}.  Their
certified outputs are resolved leaves, proper factors, strict successors,
or named cores whose dedicated local theorem closes the invocation; none is
appended as a fresh same-state key.

Finally, a Platonic localization is directly realized with its actual kernel
tuple, port order, and scalar by
\cref{lem:platonic-high-arity-localization}.  If it has arity below the
current minimum, it lowers \(\nu\).  Otherwise its card-map rank is
recomputed before re-entry.  Rank at most three is closed by
\cref{lem:platonic-deficient-rank}; rank four is closed, according to the
recomputed marked group, by
\cref{thm:tetrahedral-deck-rigidity,thm:octahedral-deck-rigidity,%
thm:icosahedral-deck-rigidity}.  Hence no ancestor rank label creates a new
same-state resource.

We have verified the reductions using actual contractions, finite arity
induction, factorization, and binary dispatch.  They therefore supply the
hypotheses needed for \cref{lem:matching-synthesis-measure}.  Every genuinely multi-copy reduction
has outcomes~1--3 of \cref{def:frozen-token-coverage}.  Any fixed-\((H,\nu)\)
proper-factor continuation arising from a one-copy reduction remains within
its originating unfinished token and satisfies outcome~4: completion
decreases \(|Q|\), while an intermediate split decreases its stored factor
potential.  The retained state remains in the same fixed-\(I\) coordinates,
so no coordinate transport can create or reactivate a token.  Thus outcome~4
has the required strict progress, proving the conditions of \cref{def:frozen-token-coverage}.
\end{proof}

\begin{proof}[Proof of \cref{thm:matching-deck-classification}]
Run the stabilized protocol of
\cref{def:relative-stratum-protocol,lem:matching-synthesis-measure} in the
fixed-\(I\) coordinates, with the existing coordinate matrix instantiated as
\(A=I\).  Any normalization used internally by
\cref{sec:p1-eight-vertex} is pulled back before a continuing branch returns
to the outer protocol.
At every continuing state initialize a
finite canonical occurrence worklist: one token for each current
tensor-prime root, together with the factor-forest and finite card universe
attached to that root.  A token is removed only after its complete deck queue
and factor queue have been exhausted.  At each outer successor, dispatch an
odd-arity member on the augmented signature set, or choose a minimum-arity
nonbinary prime and reapply \cref{lem:minimum-core-deck}; an interrupted token
is copied to the successor rather than silently marked closed.  Thus every
continuing node consists of a deck-stable all-even tuple
\((\Lambda,H,f)\), together with its existing queue \(Q\), where
\(H=\overline G(I,\Lambda)=G(\Lambda)\) is finite by
\cref{lem:bounded-torsion}.  Use the previously defined measures
\(\Xi(I,\Lambda)\) and \(\Omega(I,\Lambda,f,Q)\) from
\cref{lem:matching-synthesis-measure,eq:deck-closure-measure}.
By \cref{lem:token-coverage}, all four routed batch families satisfy
\cref{def:frozen-token-coverage}; hence
\cref{lem:matching-synthesis-measure} applies at every restart.
Strict group enlargement lowers the first coordinate, a lower minimum arity
lowers the second, fixed-\((H,\nu)\) closure decreases \(|Q|\) in the third, and a
genuine factor split lowers \(\Phi_Q\) in the fourth.  A newly exposed safe
transfer class outside \(H\) is a strict-group successor, whereas a class
already in \(H\) is handled within the current node and an unsafe binary is a
terminal.  Reductions compose toward the predecessor, common presentations
restrict, and signature-set-wide \(\Tract\) descends under retained inclusion
by \cref{lem:downward-inheritance}.

First suppose \(\arity(f)=4\).  Apply
\cref{thm:marked-finite-group-routing}; the following list is exhaustive.
\begin{enumerate}
\item If the dressed-kernel algebra is proper and \(G=C_2\), use item~1
      of \cref{thm:equality-q4-boundary}.  If \(G=C_N\), with even
      \(N\ge4\), use
      \cref{lem:proper-cyclic-q4,thm:p1-quaternary-kholant-dichotomy}.
\item If \(G=\mathcal M_m\), \(m\ge3\), use item~2 of
      \cref{thm:equality-q4-boundary}.
\item If \(G=\mathcal B_m\), with \(m\ge3\), use
      \cref{thm:marked-B-reduction}.  If
      \(G=\mathcal A_m^0\), with even \(m\ge4\), use
      \cref{thm:marked-A0-reduction}.  If
      \(G=\mathcal A_m\), with even \(m\ge4\), use
      \cref{lem:marked-A-q4,cor:p1-eight-vertex-reduced-interface}.
\item For the standard Pauli or exotic Klein group, use items~3 and~4,
      respectively, of \cref{thm:equality-q4-boundary}.
\item For \(G=\mathcal T\), any of
      \(\mathcal O,\mathcal O',\mathcal O_{dt},\mathcal O_{td}\), or either
      of the two \(\mathcal I\)-displays, use
      \cref{lem:platonic-quaternary-boundary}; for
      \(G=\mathcal T_{\rm ext}\), use
      \cref{thm:external-a4-deck-rigidity}.
\end{enumerate}
These are physical quaternary interfaces in their stated signature-set-wide
problems: \(U_{\mathcal B},U_0\), and \(U_{\mathcal A}\) do not change the
continuing outer state, and no analytic sector-coordinate eight-vertex
tensor is passed directly to
\cref{cor:p1-eight-vertex-reduced-interface}.  A directly realized
endpoint-nondegenerate tensor-prime quaternary is closed by
\cref{thm:p1-quaternary-kholant-dichotomy}; the reduced interface is used
only while a genuine factorization or its weighted-equality successor still
has to be processed.  Once the finite factor/card queue is exhausted, the
outcome is a terminal, a successor with smaller \(\Omega\), or
\(\TractX(K(\Lambda))\).  Indeed,
\(Z=K^{\otimes2}I\in K(\Lambda)\), so
\cref{lem:tract-X-anchor} identifies \(\TractX(K(\Lambda))\) with
\(\Tract(K(\Lambda))\).  Tensor primeness excludes the arity-two
alternative; hence a tractable quaternary leaf has one common
\(\cA\)-, \(\cP\)-, or \(\cL\)-presentation.  It is retained as a generic
\(\TractX(K(\Lambda))\)-terminal, and no local-affine presentation is
upgraded to affine.

Assume henceforth that \(\arity(f)\ge6\).  Apply
\cref{thm:dressed-kernel-observability}.  If the kernel algebra is
\(\Span\{I,X\}\), then
\cref{thm:reed-muller-rigidity}, whose arity-four base is
\cref{lem:proper-cyclic-q4}, gives an established exit or one common
flat-Lagrangian \(K\)-presentation, which is a structured affine subcase.

It remains to consider the full matrix algebra.  Apply
\cref{thm:marked-finite-group-routing}.  Its marked list is exhaustive
because
\[
 G\triangleleft\langle G,[Z]\rangle,
 \qquad [\langle G,[Z]\rangle:G]\le2,
 \qquad \langle[X],[Z]\rangle\cong V_4.
\]
Dispatch the resulting form by
\begin{center}
\small
\begin{tabular}{@{}ll@{}}
\toprule
complete marked form&relative theorem\\
\midrule
\(\mathcal M_m,\mathcal B_m\) (\(m\ge3\));
\(\mathcal A_m^0,\mathcal A_m\) (even \(m\ge4\))&
\cref{thm:marked-dihedral}\\
standard Pauli \(V_4\) and \(\mathcal V_{\rm ex}\)&
\cref{thm:v4-relative}\\
all marked Platonic forms&
\cref{thm:platonic-decks}\\
\bottomrule
\end{tabular}
\end{center}
Within the marked-dihedral row, the \(\mathcal B_m\) branch is a
whole-problem equality-wire equivalence.  The \(\mathcal A_m^0\) branch uses
the literal-\(I\) whole-problem/open-network equivalence
\cref{eq:marked-A0-fixed-I,eq:marked-A0-open-equivalence}; its two edge
substitutions use only \(I,Y\), and its final common witness is pulled back
inside \cref{thm:marked-A0-reduction}.  For \(\mathcal A_m\),
\cref{thm:marked-A-rigidity} resolves arity four by its physical
whole-set reduction and excludes every higher-arity stable prime.
The sector matrix \(U_{\mathcal A}\) expresses the weighted pencils
and fusion identities used in that proof; the realized contractions
remain in the current fixed-\(I\) coordinates.
Common local-affine conclusions in the \(\mathcal B_m\) and exotic Klein
branches are direct generic \(\TractX(K(\Lambda))\)-terminals; they are not
promoted to affine presentations.  By contrast, the flat-Lagrangian
rigidity survivor is already a structured affine subcase.

Each dispatched theorem closes the stabilized protocol.  Factor batches
create outer successors only by lowering the minimum nonbinary arity or
strictly enlarging the recomputed complete group, so \(\Xi\) decreases.
A safe binary in \(G\) is within-node closure; an unsafe binary is terminal;
and an actual safe binary outside the asserted complete group forces
recomputation by \cref{lem:actual-binary-completeness}.  Thus no unresolved
branch recurs.

Observability exhausts the two kernel algebras and marked routing exhausts
the full-algebra case.  Every branch therefore ends in a terminal, a direct
\(\TractX(K(\Lambda))\)-leaf, or one common product or affine
\(K\)-presentation, with the flat-Lagrangian survivor forming a structured
affine subcase.  The outer synthesis remains in the fixed-\(I\) coordinates;
any normalization used internally by \cref{sec:p1-eight-vertex} is
discharged before a continuing branch re-enters this protocol.  Non-\(GO(X)\)
sector matrices occur only in whole-problem equivalences or as final
common-presentation witnesses.  By
\cref{def:common-K-presentation}, such a witness is one basis that
simultaneously places the transformed edge and every signature in the
retained set in the asserted class.  Common product and affine
presentations, including the flat-Lagrangian subcase, are clauses of
\(\TractX(K(\Lambda))\); common local-affine conclusions are already direct
\(\TractX(K(\Lambda))\)-leaves.  Since
\(Z=K^{\otimes2}I\in K(\Lambda)\),
\cref{lem:tract-X-anchor} identifies this criterion with
\(\Tract(K(\Lambda))\).  Hence the two cases in
\cref{thm:matching-deck-classification} are exhaustive.
\end{proof}

\begin{corollary}[Equality-accessible signature-set-wide closure]
\label{cor:equality-accessible-closure}
Let \(\Lambda_0\) be a retained, factor-saturated finite algebraic
signature set containing the distinguished literal equality \(I\).  If
\(\TractX(K(\Lambda_0))\) holds, then
\(\KHolant(\Lambda_0)\in\FP\); otherwise
\(\KHolant(\Lambda_0)\) is \(\SharpP\)-hard.  Here
\(\TractX(K(\Lambda_0))\) is equivalent to \(\Tract(K(\Lambda_0))\), since
\(Z=K^{\otimes2}I\in K(\Lambda_0)\).  The factor access used in reaching
this conclusion is Turing exposure; every deck is regenerated over the
current augmented signature set and its recomputed complete group.
\end{corollary}

\begin{proof}
Because \(Z=K^{\otimes2}I\in K(\Lambda_0)\),
\cref{lem:tract-X-anchor} identifies
\(\TractX(K(\Lambda_0))\) with \(\Tract(K(\Lambda_0))\).  If this
criterion holds, \cref{thm:predicate-algorithms} and the fixed holographic
equivalence give \(\KHolant(\Lambda_0)\in\FP\).

Suppose the criterion fails.  After the odd-arity and unsafe-binary
terminals, a state with only unary and binary prime factors would satisfy the
arity-two alternative and is therefore impossible on this branch.
Otherwise \cref{lem:minimum-core-deck} supplies a minimum nonbinary core.
Strong induction on its arity handles every lowering successor, while a
stable successor is closed by \cref{thm:matching-deck-classification}; that
theorem's \(\Omega\)-descent also closes all later restarts.  Reductions
compose toward \(\Lambda_0\), tractability descends, and common
presentations restrict.  Thus every nonhard leaf would imply
\(\Tract(K(\Lambda_0))\), and hence \(\TractX(K(\Lambda_0))\), contrary to
the assumed failure.  Every remaining branch is therefore
\(\SharpP\)-hard.
\end{proof}

\subsection{Proofs of the Two Main Dichotomies}

\Cref{cor:equality-accessible-closure} gives polynomial-time evaluation
when \(\TractX\) holds and hardness otherwise.  Its tractable alternatives
are the arity-two tensor closure and the common affine, product, or
local-affine classes; the flat-Lagrangian case lies in the affine class.
We now match this criterion to the input sets of the two main theorems.

\begin{proof}[Proof of \cref{thm:p1-eq2-kholant-dichotomy}]
Preprocess \(\cF\) according to \cref{subsec:zero-nullary} and put
\[
  \Lambda:=\cF^\circ.
\]
Because \(I\in\cF\) is nonzero and has positive arity, \(I\in\Lambda\), and
the exact preprocessing gives
\(\KHolant(\cF)\equivT\KHolant(\Lambda)\).
The preprocessing convention and \cref{lem:tract-X-anchor} give
\[
 \TractX(K(\cF))
 \quad\Longleftrightarrow\quad
 \TractX(K(\Lambda))
 \quad\Longleftrightarrow\quad
 \Tract(K(\Lambda)),
\]
because \(Z=K^{\otimes2}I\in K(\Lambda)\).
If this criterion holds, then \cref{thm:predicate-algorithms}, together
with the fixed holographic equivalence, gives
\(\KHolant(\Lambda)\in\FP\).

Suppose the criterion fails.  The available equality \(I\) has support on
both charge signs, so \(\Lambda\) is not common one-sided.  Apply
\cref{lem:factor-saturation} and put
\[
  \Lambda^+=\Lambda\cup\Pi(\Lambda).
\]
Then \(\KHolant(\Lambda^+)\equivT\KHolant(\Lambda)\), and
\(\Lambda^+\) is a finite retained factor-saturated set still containing
\(I\).  If \(\TractX(K(\Lambda^+))\) held, then
\cref{lem:tract-X-anchor} and downward inheritance would imply
\(\Tract(K(\Lambda))\), contrary to the failure above.  Hence
\cref{cor:equality-accessible-closure} makes
\(\KHolant(\Lambda^+)\), and therefore \(\KHolant(\Lambda)\),
\(\SharpP\)-hard.  Decidability follows from
\cref{thm:predicate-algorithms}.  The coordinate convention of the closure theorem is the fixed-\(I\)
convention used here.
\end{proof}

\begin{proof}[Proof of \cref{thm:p1-deq2-dichotomy}]
Preprocess \(\cF\) according to \cref{subsec:zero-nullary} and define
\[
  \Lambda:=\cF^\circ\cup\{I\}.
\]
Then
\[
  \Holant(\cF,X)
  \equivT
  \Holant(\cF^\circ,X)
  \equivT
  \KHolant(\Lambda).
\]
The first equivalence performs the exact zero/nullary evaluation, and the
second is \cref{lem:p1-equality-wire-equivalence}, with the literal equality
already included in \(\Lambda\).  The set in the tractability criterion is the fixed holographic image
\(K(\Lambda)\).  Apply
\cref{thm:p1-eq2-kholant-dichotomy}; the
preprocessing and wire substitutions are linear-size, and every known
nonzero scalar is restored exactly.  Under the preprocessing convention,
its tractability condition is precisely
\[
 \TractX(K(\Lambda))
 \quad\Longleftrightarrow\quad
 \TractX(K(\cF\cup\{I\})),
\]
which is the criterion stated in \cref{thm:p1-deq2-dichotomy}.
\end{proof}
\section{Conclusion}
\label{sec:p1-conclusion}

The supplied binary disequality relation yields a decidable
\(\FP\)-versus-\(\SharpP\)-hard classification for every fixed finite
set of algebraic complex-valued Boolean signatures.
\Cref{thm:p1-deq2-dichotomy,thm:p1-eq2-kholant-dichotomy} express this
classification using, respectively, ordinary edges with an \(X\)-vertex
and native \(X\)-edges with an \(I\)-vertex.  The binary anchor leaves
four tractable alternatives: the arity-two tensor closure and common
affine, product, or local-affine transformability.

The proof connects contractions of a smallest nonbinary prime to the
finite group of available binary transfers.  Arity descent, strict
group enlargement, and complete factor processing reduce the problem
to the four marked group analyses.
\Cref{thm:p1-quaternary-kholant-dichotomy} supplies a signature-set
boundary without requiring literal equality, extending the
single-signature eight-vertex classification
\cite{CaiFu2023} in the stated endpoint-nondegenerate, tensor-prime
setting.  The finite classifications in the appendix support the
remaining low-arity cases; their precise scope is stated at each use.

\appendix
\section{Finite Classifications and Explicit Contractions}
\label{sec:exact-finite-certificates}

This appendix supplies the finite classifications and explicit contractions
used in \cref{sec:p1-eight-vertex,sec:weighted-equality,%
sec:normalized-dihedral,sec:v4-deck,sec:platonic-decks}.  Each subsection
specifies the tensors being classified, the equations imposed on them,
and the consequence needed in the main proof.  Coefficient identities and
short finite classifications are given explicitly.  Where a completeness
claim still relies on exhaustive exact computation, its proof identifies
the state space and tests, together with the companion verifier in
\path{certificates/}.  The code supplements these stated mathematical
arguments and is included with the source.

In a histogram, \(s^m\) means \(m\) objects of size \(s\).
Projective representatives are normalized by their first nonzero
coordinate.  Rank, span, and equality tests use the exact
characteristic-zero field specified in each calculation.  We use the
projective, tensor, affine, and algebraic conventions of
\cref{subsubsec:group-projective,subsubsec:projective-tensor-geometry,%
subsubsec:symmetric-binary-spaces,subsubsec:boolean-affine-tools,%
subsubsec:effective-algebraic-geometry}.

\subsection{Endpoint-Nondegenerate Eight-Vertex Signatures}
\label{appsec:cert-eight-vertex}

This subsection gives the two exact finite statements used in
\cref{sec:p1-eight-vertex}.

\subsubsection{The fourth-root odd quaternary residue}
\label{proofsubsec:cert-mu4}

Normalize a full-support odd-parity quaternary signature by
\(a_j=f(e_j)\), \(b_j=f(\one-e_j)\), and \(a_0=1\).  Its fourth-root
state is
\begin{equation}
  a_j=i^{\alpha_j},\qquad b_j=i^{\beta_j},
  \qquad
  (\alpha_1,\alpha_2,\alpha_3,\beta_0,\ldots,\beta_3)
  \in\mathbb Z_4^7.
  \label{eq:cert-mu4-state}
\end{equation}
There are \(4^7=16\,384\) states; write \(N(z)=z\bar z\) for the Gaussian
norm.  For \(i<j\), with \(k<l\) the complementary pair, the six native
and six equality-mediated loops are
\begin{equation}
 L^X_{ij}=\diag(a_i+a_j,b_i+b_j),
 \qquad
 L^I_{ij}=
 \begin{pmatrix}0&a_l+b_k\\a_k+b_l&0\end{pmatrix}.
  \label{eq:cert-twelve-loops}
\end{equation}
The latter inserts one equality vertex using two native \(X\)-wires.  A
series equality support-flips the former:
\begin{equation}
 \diag(P,Q)XI=\begin{pmatrix}0&P\\Q&0\end{pmatrix},
 \qquad
 T_{\diag(P,Q)XI}=\diag(P,Q).
  \label{eq:cert-diagonal-loop-support-flip}
\end{equation}
Thus in either case \(P/Q\) is the eigenvalue ratio of an actual transfer.

\begin{proposition}[Deterministic twelve-loop test]
\label{prop:cert-twelve-loop}
For each state in \cref{eq:cert-mu4-state}, inspect the twelve ordered pairs
in \cref{eq:cert-twelve-loops}:
\begin{enumerate}[label=\textnormal{(\roman*)}]
  \item If exactly one coordinate of some pair is zero, declare a rank-one
  exit.
  \item Otherwise, for every pair with two nonzero coordinates \(P,Q\),
  compare their Gaussian norms.  If \(N(P)\neq N(Q)\), declare an
  infinite-projective-order exit.
  \item Otherwise declare the state remaining, retaining zero--zero pairs.
\end{enumerate}
The resulting disjoint partition is
\begin{equation}
 16\,384=15\,792+224+368,
 \label{eq:cert-mu4-histogram}
\end{equation}
where the summands are respectively rank-one, infinite-order, and remaining
states.
\end{proposition}

\begin{proof}
The entries lie in \(\mathbb Z[i]\).  Up to a fourth root of unity, every
nonzero sum of two fourth roots is either \(2\) or \(1+i\).
Consequently two such sums have the same Gaussian norm precisely when their
quotient lies in \(\mu_4\).  If their norms differ, the
quotient has modulus different from one and therefore cannot have finite
multiplicative order.  Thus, for exactly the nonzero sums occurring here,
\(P/Q\) is torsion if and only if \(N(P)=N(Q)\).  Hence
\cref{eq:cert-diagonal-loop-support-flip} makes the three tests exact and
disjoint; exhaustive evaluation gives \cref{eq:cert-mu4-histogram}.
\end{proof}

Exactly 64 remaining states satisfy
\begin{equation}
 b_j=\eps a_j\quad(0\leq j\leq3),\qquad
 \eps\in\{1,-1\},\qquad \prod_{j=0}^3a_j\in\{i,-i\}.
  \label{eq:cert-mu4-exceptional64}
\end{equation}
Indeed, after the normalization \(a_0=1\), exactly half of the \(4^3\)
choices of \((a_1,a_2,a_3)\) have product in \(\{i,-i\}\), and \(\eps\)
has two choices; hence the count is \(32\cdot2=64\).
Join ports \(0,1\) of two copies \(A,B\) by native \(X\)-edges.  Denote the
resulting quaternary signature by \(Q\).  In the order
\((A_2,A_3,B_2,B_3)\), its coefficients are
\[
\begin{array}{ccccccc}
\toprule
x&0000&1111&0011,1100&0101&1010&0110,1001\\ \midrule
Q(x)&2a_0a_1&2b_0b_1&a_0b_0+a_1b_1&2a_3b_2&2a_2b_3&
a_2b_2+a_3b_3.\\
\bottomrule
\end{array}
\]
The contraction of two odd-parity signatures along two \(X\)-edges has even
external parity, and the two endpoint values \(2a_0a_1\) and \(2b_0b_1\)
are nonzero.  Under \cref{eq:cert-mu4-exceptional64}, exactly one of
\(a_0^2+a_1^2,a_2^2+a_3^2\) vanishes: indeed, each \(a_j^2\) is a sign and
\(\prod_j a_j^2=-1\), so exactly one of the two displayed pairs consists of
opposite signs.  Hence \(Q\) is a support-six eight-vertex signature.

For each other remaining state and \(\delta\in\mathbb Z_4\), define the
phase-shifted signature \(f_\delta\) by
\[
 f_\delta(e_j)=i^{\alpha_j},
 \qquad
 f_\delta(\one-e_j)=i^{\beta_j+\delta}
\]
and keep \(f_\delta(x)=0\) on every even-weight input.  Parameterize the odd
coset by
\[
 x(y)=(y_0,y_1,y_2,1\oplus y_0\oplus y_1\oplus y_2),
 \qquad y\in\Ftwo^3.
\]
Define the phase-exponent function
\(\varphi_\delta:\Ftwo^3\to\mathbb Z_4\) by
\(f_\delta(x(y))/f_\delta(x(000))=i^{\varphi_\delta(y)}\), and define its
linear coefficients \(\ell_j\) and mixed-quadratic coefficients \(q_{jk}\)
by
\[
 \ell_j=\varphi_\delta(e_j),\qquad
 q_{jk}=\varphi_\delta(e_j+e_k)-\ell_j-\ell_k.
\]
Then \(f_\delta\) is affine if and only if
\begin{equation}
 q_{jk}\in\{0,2\}\ (j<k),\qquad
 \varphi_\delta(111)
 =\ell_0+\ell_1+\ell_2+q_{01}+q_{02}+q_{12}\pmod 4.
  \label{eq:cert-mu4-affine-test}
\end{equation}
Indeed, in the Boolean M\"obius convention of
\cref{subsubsec:boolean-affine-tools}, the expansion of
\(\varphi_\delta\) over \(\mathbb Z_4\) has the three
\(q_{jk}\) as its quadratic coefficients, while the difference between the
two sides of the second congruence is its cubic coefficient.  A modulo-four
phase is affine exactly when the quadratic coefficients are even
and that cubic coefficient vanishes.
Exactly two \(\delta\)'s pass: \(\{0,2\}\) for 176 states and
\(\{1,3\}\) for 128 states.

The complete 512-state affine odd family is
\begin{equation}
 (a_0,a_1,a_2,a_3)=(1,A,B,C),\qquad
 (b_0,b_1,b_2,b_3)=(xyzABC,zBC,yAC,xAB),
 \label{eq:cert-affine-odd-param}
\end{equation}
where \(A,B,C\in\mu_4\) and \(x,y,z\in\{1,-1\}\); the \(4^3\cdot2^3\)
choices give 512 distinct normalized states.  The locus having no rank-one
native \(X\)-loop is the disjoint union
\begin{align}
 &x=y=z &&(128\text{ states}),\label{eq:cert-affine-odd-rank-a}\\
 &x,y,z\text{ not all equal},\quad
   (A^2,B^2,C^2)=(xy,xz,yz)
 &&(48\text{ states}).\label{eq:cert-affine-odd-rank-b}
\end{align}
The other 336 states have a rank-one loop.  For the remaining 176, let
\(G_{0j}\) join ports \(0,j\) of two copies by corresponding native
\(X\)-edges and retain the other ports in copy order.  Call a resulting
quaternary \emph{dense product-transformable} when it has the full-even form
in \cref{eq:p1-eight-vertex-product-dense-form}, namely
\[
 g(x)=\rho r^{\wt(x)}(-1)^{s\cdot x}
       \one[\wt(x)\text{ is even}]
 \quad\text{for some }\rho,r\ne0,\ s\in\Ftwo^4.
\]
Then
\begin{align}
G_{01}\text{ is dense product-transformable}
  &\Longleftrightarrow x=y\text{ and }B^2=C^2,\notag\\
G_{02}\text{ is dense product-transformable}
  &\Longleftrightarrow x=z\text{ and }A^2=C^2,\notag\\
G_{03}\text{ is dense product-transformable}
  &\Longleftrightarrow y=z\text{ and }A^2=B^2.
\label{eq:cert-affine-odd-cover}
\end{align}
In \cref{eq:cert-affine-odd-rank-a} two squares agree; in
\cref{eq:cert-affine-odd-rank-b} the unique equal pair of \(x,y,z\) has
the corresponding equal squares.  Hence one test in
\cref{eq:cert-affine-odd-cover} always succeeds.  These formulas are a
self-contained mathematical certificate.

\subsubsection{Affine endpoint normal forms}
\label{proofsubsec:cert-affine-q4-endpoint}

The affine classification used in \cref{sec:p1-eight-vertex} reduces to
eighteen coefficient transforms and three endpoint-support types.  For a nonzero affine signature with support \(s+L\),
choose the lexicographically first support word as origin and the
lexicographically first spanning basis of \(L\) as intrinsic coordinates,
and divide by the value at that origin.  Define the eighteen core matrices
\begin{equation}
\begin{split}
\mathscr K_{\mathrm{aff}}
={}&\{I,X\}\\
&{}\cup
\left\{
 \begin{pmatrix}1&r\\1&-r\end{pmatrix}:r\in\mu_8
\right\}\\
&{}\cup
\left\{
 \begin{pmatrix}
 1&u(1+\varepsilon\sqrt2)\\
 1&u(1-\varepsilon\sqrt2)
 \end{pmatrix}:
 u\in\mu_4,\ \varepsilon\in\{1,-1\}
\right\}.
\end{split}
\label{eq:cert-affine-q4-cores}
\end{equation}

\begin{proposition}[Exact affine endpoint-nondegenerate eight-vertex certificate]
\label{prop:cert-affine-q4-endpoint}
For every canonically projectively normalized nonzero arity-four affine signature
\(f\in\cA\) and every \(K\in\mathscr K_{\mathrm{aff}}\), define the
transformed signature \(q\) by
\[
  q:=K^{\otimes4}f.
\]
Retain \((K,f)\) precisely when
\begin{equation}
  q(x)=0\quad\text{for odd }\wt(x),
  \qquad
  q(\zero^4)q(\one^4)\ne0.
  \label{eq:cert-affine-q4-retention}
\end{equation}
Then the following assertions hold.

\begin{enumerate}[label=\textnormal{(\roman*)}]
\item Writing each affine support as \(s+L\), there are \(36\,720\)
normalized projective affine signatures:
\[
\begin{array}{cccccc}
\toprule
\dim L&0&1&2&3&4\\ \midrule
\#\text{ affine supports}&16&120&140&30&1\\
\#\text{ normalized signatures}&16&480&4480&15360&16384.\\
\bottomrule
\end{array}
\]
Among the \(36\,720\cdot18\) pairs \((K,f)\), exactly \(4200\)
satisfy \cref{eq:cert-affine-q4-retention}.  Six cores retain \(612\)
inputs each and the other twelve retain \(44\) inputs each.

\item Every retained \(q\) has the form
\begin{equation}
  q=\gamma\,\diag(1,t)^{\otimes4}h,
  \qquad
  \gamma,t\in\AlgNums^\times,
  \label{eq:cert-affine-q4-diagonal-form}
\end{equation}
where \(h\in\cA\), \(h(\zero^4)=1\), and
\[
 \supp(h)\in
 \left\{
  \{\zero^4,\one^4\},\
  \{\zero^4,\one^4,\one_I,\one_{I^c}\},\
  \{x:\wt(x)\equiv0\pmod2\}
 \right\},
\]
with \(|I|=2\).  Let \(\mathscr R_{\mathrm{aff}}\) be the finite set of all
such normalized affine tensors \(h\).  Their support-size histogram is
\[
  2^4,\qquad4^{96},\qquad8^{512}.
\]
The three support cases form, respectively, \(1,24,256\) diagonal orbits.
Here
\[
 g\sim_{\mathrm{diag}}g'
 \quad\Longleftrightarrow\quad
 g'=\gamma\,\diag(1,t)^{\otimes4}g
 \quad\text{for some }\gamma,t\in\AlgNums^\times.
\]

More explicitly, fix the lexicographic order on the six weight-two words.
If \(v_1,\ldots,v_6\) are the corresponding central coefficients of an
endpoint-nondegenerate eight-vertex signature \(g\), define its
diagonal-orbit invariant vector by
\begin{equation}
 \mathcal J(g):=
 \left(
 \frac{v_rv_s}{g(\zero^4)g(\one^4)}
 \right)_{1\le r\le s\le6}
 \label{eq:cert-affine-q4-invariant}
\end{equation}
Then \(\mathcal J(g)\) is a complete invariant of these diagonal orbits.

\item For a support-four reference, let \(I\sqcup I^c=[4]\) be its two
two-port blocks and write its logical table as
\[
 F=\begin{pmatrix}A&B\\ C&D\end{pmatrix},
 \qquad ABCD\ne0.
\]
Here the row bit is the common value on the \(I\)-block, the column bit is
the common value on the \(I^c\)-block, and
\[
 A=h(\zero^4),\qquad B=h(\one_{I^c}),\qquad
 C=h(\one_I),\qquad D=h(\one^4).
\]
Then \(AD/BC\in\{1,-1\}\).  If \(AD=BC\), the signature is a tensor
product of two nonsingular weighted equalities.  If \(AD=-BC\), take two
labelled copies, join each of the two physical ports in the \(I^c\)-block
of the first copy to the corresponding port of the second copy by a native
\(X\)-edge, and retain the two \(I\)-blocks as the four external ports.
The resulting logical table is
\begin{equation}
 FXF^{\mathsf T}
 =\begin{pmatrix}
 2AB&AD+BC\\
 AD+BC&2CD
 \end{pmatrix}
 =\begin{pmatrix}2AB&0\\0&2CD\end{pmatrix},
 \label{eq:cert-affine-q4-support-four-gadget}
\end{equation}
and hence the physical output is a nonzero pure generalized equality on
its four external ports.

\item The \(512\) full-even references have the disjoint first-witness
partition
\begin{equation}
  512=16+336+96+48+16.
  \label{eq:cert-affine-q4-full-partition}
\end{equation}
The first \(16\) have the dense product-transformable form of
\cref{eq:p1-eight-vertex-product-dense-form}.  The next \(336\) have a
rank-one native \(X\)-self-loop.  For each remaining state, take two
labelled copies and, successively for \(j=1,2,3\), join ports \(0,j\) of
the first copy by crossed native \(X\)-edges to ports \(j,0\) of the
second.  Order the external ports as the uncontracted ports of the first
copy in increasing order, followed by the corresponding ports of the second
copy.  The first support-four product occurs for, respectively,
\(96,48,16\) new states.  It has the form
\[
  cE_r\otimes E_r,
  \qquad c\in\mathbb C^\times,\qquad E_r:=\diag(1,r),
  \qquad r\in\mu_4,
\]
with respect to the induced perfect matching and displayed external-port
order.  Here \(c\) is a nonzero scalar and \(E_r\) is the diagonal binary
matrix.  The histogram in the order \(r=1,-1,i,-i\) is
\[
  (56,56,24,24).
\]

\item Every normalized full-support fourth-root odd table \(p\) of
\cref{eq:cert-mu4-state} is classified as follows.  The twelve actual
loops in \cref{eq:cert-twelve-loops} give the partition
\[
  16\,384=15\,792+224+368
\]
of \cref{prop:cert-twelve-loop}.  Among the \(368\) remaining states,
exactly \(64\) satisfy \cref{eq:cert-mu4-exceptional64}; the displayed
two-copy gadget there is a support-six endpoint-nondegenerate eight-vertex signature.

For each of the other \(304\) states, let \(p_\delta\) multiply the four
weight-three values of \(p\) by \(i^\delta\) while leaving its weight-one
values fixed.  Exactly two \(\delta\in\mathbb Z_4\) make \(p_\delta\)
affine.  The valid sets are
\(\{0,2\}\) for \(176\) states and \(\{1,3\}\) for \(128\) states.
Every such \(p_\delta\) lies in the \(512\)-state family
\cref{eq:cert-affine-odd-param}.  Of that family, \(336\) states have a
rank-one native \(X\)-loop; every one of the remaining \(176\) states has
at least one two-copy output of the dense product-transformable form among
the three gadgets in
\cref{eq:cert-affine-odd-cover}.
\end{enumerate}
\end{proposition}

\begin{proof}
For a support \(s+L\), put \(d=\dim L\).  In the canonical intrinsic
coordinates \(x=(x_1,\ldots,x_d)\), the normalized phases are exactly
\[
 i^{\sum_{j=1}^d\ell_jx_j+
      2\sum_{1\le j<k\le d}q_{jk}x_jx_k},
 \qquad
 \ell_j\in\mathbb Z_4,\qquad q_{jk}\in\Ftwo.
\]
There are
\[
 2^{4-d}\genfrac[]{0pt}{}{4}{d}_2\,4^d2^{\binom d2}
\]
such signatures of support dimension \(d\): the Gaussian binomial counts
the linear spaces, where
\(\genfrac[]{0pt}{}{4}{d}_2=\prod_{j=0}^{d-1}(2^4-2^j)/(2^d-2^j)\) and the empty
product is one.  The factor \(2^{4-d}\) counts their cosets, and the last two factors the
linear and mixed phase coefficients.  This proves the first table in
\emph{(i)}.  Applying the eighteen matrices in
\cref{eq:cert-affine-q4-cores} and testing
\cref{eq:cert-affine-q4-retention} gives the \(4200\) retained pairs.

For completeness of the orbit test, normalize the first endpoint to one,
write the central layer as \(v\), and write the second endpoint as \(e\).
If \(v=0\), two states are diagonally equivalent after choosing
\(t^4=e'/e\).  Otherwise equality of
\cref{eq:cert-affine-q4-invariant}, using one nonzero coordinate of \(v\),
gives
\[
  v'=\lambda v,
  \qquad e'=\lambda^2e.
\]
Choosing \(t^2=\lambda\) proves
\cref{eq:cert-affine-q4-diagonal-form}.  The support and orbit histograms
then follow from the exhaustive retained list.

On a support-four affine plane, choose the logical coordinates \(u,v\) used
in the displayed table.  After projective normalization its phase is
\[
 h(u,v)=i^{\ell_1u+\ell_2v+2quv},
 \qquad \ell_1,\ell_2\in\mathbb Z_4,\quad q\in\Ftwo.
\]
Therefore \(AD/(BC)=(-1)^q\in\{1,-1\}\), proving the asserted cross-ratio;
the literal wiring in \cref{eq:cert-affine-q4-support-four-gadget} is then a
direct coefficient identity.  For full support, the exact replay tests the six literal native
self-loops and the three crossed two-copy contractions in the stated
first-witness order; this gives
\cref{eq:cert-affine-q4-full-partition} and the \(r\)-histogram.

Finally it evaluates all \(4^7\) states of
\cref{eq:cert-mu4-state}, the twelve ordered pairs in
\cref{eq:cert-twelve-loops}, the exceptional two-copy output, all four
tests in \cref{eq:cert-mu4-affine-test}, and the three literal gadgets
\(G_{0j}\).  This gives the last assertion.

The companion exact replay is
\path{certificates/verify_affine_q4.py}.  It ranges over every
affine support \(s+L\), every affine phase in the intrinsic coordinates,
all eighteen cores, all \(4^7\) normalized fourth-root odd states, all
twelve loops, all four phase shifts, and all three two-copy contractions
listed above.  The enumeration regenerates these finite spaces from their definitions
and evaluates their identities over \(\mathbb Z[i,\sqrt2]\) and
\(\mathbb Q(i,\sqrt2)\).
\end{proof}

\begin{remark}[Use in the main proof]
\label{rem:cert-affine-q4-endpoint-scope}
\Cref{prop:cert-affine-q4-endpoint} classifies coefficients for the eighteen
displayed cores.  The main proof supplies the transformability witness,
reduces it to those cores, and implements the required constructions in
the retained Holant problem.  In particular, the support-four gadget is
the literal two-copy wiring in
\cref{eq:cert-affine-q4-support-four-gadget}.  The retained-list and
full-support partitions still use exhaustive exact computation in
\path{certificates/verify_affine_q4.py}; the program emits assertions,
rather than a separate proof object.
\end{remark}

\begin{theorem}[Affine endpoint normal forms]
\label{thm:cert-interface-affine-q4-normal-form}
The core set used by the affine endpoint certificate is exactly the
eighteen-element set \(\mathscr K_{\mathrm{aff}}\) displayed in
\cref{eq:cert-affine-q4-cores}.  Let \(f\in\cA\) be a nonzero canonically
projectively normalized arity-four signature, let
\(K\in\mathscr K_{\mathrm{aff}}\), and put \(q=K^{\otimes4}f\).  If
\[
 q(x)=0\quad(\wt(x)\text{ odd}),\qquad
 q(\zero^4)q(\one^4)\ne0,
\]
then there are \(\gamma,t\in\AlgNums^\times\) and a normalized
\(h\in\cA\), with \(h(\zero^4)=1\), such that
\[
 q=\gamma\,\diag(1,t)^{\otimes4}h.
\]
Moreover, \(\supp(h)\) is exactly one of the following:
\begin{enumerate}[label=\textnormal{(\roman*)}]
\item the complementary endpoint pair \(\{\zero^4,\one^4\}\);
\item a four-point plane
      \(\{\zero^4,\one^4,\one_I,\one_{I^c}\}\) with \(|I|=2\);
\item the full even-parity support.
\end{enumerate}
\end{theorem}

\begin{proof}
This is the direct-use form of
\cref{prop:cert-affine-q4-endpoint}\textnormal{(i)--(ii)}.  The three
supports have distinct cardinalities, so the alternatives are disjoint.
\end{proof}

\begin{theorem}[Gadgets from affine reference states]
\label{thm:cert-interface-affine-q4-gadget-exits}
Let \(h\) be supplied by
\cref{thm:cert-interface-affine-q4-normal-form}.
\begin{enumerate}[label=\textnormal{(\roman*)}]
\item If \(|\supp(h)|=2\), then \(h\) is a pure generalized equality.
\item If \(|\supp(h)|=4\), its logical table
      \(F=\begin{psmallmatrix}A&B\\ C&D\end{psmallmatrix}\) satisfies
      \(AD/BC\in\{1,-1\}\).  If \(AD=BC\), then \(h\) genuinely factors
      across its two two-port blocks into nonsingular weighted equalities;
      if \(AD=-BC\), the literal two-copy \(X\)-edge construction directly
      realizes a nonzero pure generalized equality.
\item If \(h\) has full even support, then either it has the dense
      product-transformable form of
      \cref{eq:p1-eight-vertex-product-dense-form}, a native \(X\)-self-loop
      is nonzero of rank one, or one of the three crossed two-copy
      contractions specified in
      \cref{prop:cert-affine-q4-endpoint}\textnormal{(iv)} directly realizes
      \(cE_r\otimes E_r\), where \(c\ne0\),
      \(E_r=\diag(1,r)\), and \(r\in\mu_4\).
\end{enumerate}
All constructions asserted to be direct use the literal wirings specified
in this appendix.
\end{theorem}

\begin{proof}
The support-two assertion is immediate, and the other two assertions are
exactly \cref{prop:cert-affine-q4-endpoint}\textnormal{(iii)--(iv)}.
\end{proof}

\begin{theorem}[Contractions of fourth-root odd signatures]
\label{thm:cert-interface-affine-q4-odd}
Let \(p\) be a normalized full-support fourth-root odd table.  Its twelve
actual loops give exactly one of the following routes: a nonzero rank-one
binary, a nonsingular binary of infinite projective transfer order, or a
remaining table, with every zero--zero loop retained.  In the remaining case,
either the exceptional two-copy construction directly realizes a support-six
endpoint-nondegenerate eight-vertex signature, or a certified phase shift
\(p_\delta\in\cA\) has a rank-one native loop or a dense
product-transformable two-copy output.  These alternatives exhaust all
fourth-root odd tables and include every scale and cancellation locus.
\end{theorem}

\begin{proof}
Combine the deterministic test in \cref{prop:cert-twelve-loop} with
\cref{prop:cert-affine-q4-endpoint}\textnormal{(v)}.
\end{proof}

\subsection{Equality-Accessible Matching Decks}
\label{appsec:cert-equality-accessible}
\label{proofsubsec:cert-equality-q4-boundary}

This subsection gives the shared finite quaternary certificate statement used
in \cref{sec:weighted-equality} and by the later finite-group branches.

\begin{proposition}[Physical equality-accessible quaternary certificate]
\label{prop:cert-equality-q4-boundary}
Let \((\Lambda,G,q)\) be a deck-stable retained state, with \(\Lambda\)
factor-saturated and containing the distinguished literal equality \(I\), and
\(q\ne0\) a quaternary tensor-prime core.  Suppose its complete physical
deck consists only of zero cards and nonzero scalar multiples of the fixed
actual binary representatives of \(G\), and suppose
\(G\) is \(C_2=\{1,[X]\}\), a
monomial dihedral group
\(\{[\diag(1,t)],[X\diag(1,t)]:t\in\mu_m\}\) with \(m\ge3\), the
standard Pauli \(V_4\), or the exotic group
\[
 \left\{[I],[X],
 \left[\begin{pmatrix}i&1\\-1&-i\end{pmatrix}\right],
 \left[\begin{pmatrix}-i&1\\-1&i\end{pmatrix}\right]\right\}.
\]
Then an actual gadget using at most two copies
of \(q\) produces an endpoint-nondegenerate eight-vertex signature.  The assertion covers all
zero cards, both port orientations, all six unordered contracted pairs,
and the fixed nonzero scalars of the actual group representatives.  In every
case, the endpoint-nondegenerate eight-vertex signature can moreover be chosen not to be a
\(G\)-matching signature.
\end{proposition}

\begin{proof}
Port dressings and reversal only permute the complete group, multiply cards
by invertible factors, and contribute known nonzero scalars.  It therefore
suffices to check every standard group kernel on the six unordered pairs.

\emph{The \(C_2\) case.}
Use Walsh coordinates only for this calculation, and define the Walsh
coefficients
\[
 a_S=\widetilde q(\one_S),\qquad S\subseteq[4].
\]
For any Walsh table indexed by subsets of \([4]\), call it
\emph{Walsh-complement-symmetric} when
\(a_{\bar S}=a_S\) for every \(S\), where
\(\bar S=[4]\setminus S\).
For a deleted pair \(ij\), the two actual \(I/X\)-cards satisfy
\begin{equation*}
 \widetilde C_{ij}^{\pm}(\one_T)
 =2(a_T\pm a_{Tij}).
\end{equation*}
Here \(Tij:=T\cup\{i,j\}\), with \(T\) a subset of the two residual ports.
The safe lines are \(I,Z\).  Singleton residual coordinates force
\(a_S=0\) for odd \(|S|\).  For
\(\{i,j\}\sqcup\{k,l\}=[4]\), define the four Walsh coefficient scalars
\[
 u=a_\varnothing,\quad v=a_{1234},\quad
 x=a_{ij},\quad y=a_{kl}.
\]
Safety for the \(ij\)-cards gives
\[
 (y+v)^2=(u+x)^2,\qquad (y-v)^2=(u-x)^2,
\]
and safety for the complementary \(kl\)-cards gives
\[
 (x+v)^2=(u+y)^2,\qquad (x-v)^2=(u-y)^2.
\]
Adding and subtracting these four equations gives
\begin{equation*}
 v^2=u^2,\qquad y^2=x^2,\qquad
 yv=ux,\qquad xv=uy.
\end{equation*}
If \(u\ne0\), one \(\epsilon\in\{1,-1\}\) satisfies
\[
 v=\epsilon u,\qquad a_{\bar S}=\epsilon a_S.
\]
When \(u=v=0\), call a complementary weight-two block \(S,\bar S\)
\emph{active} if \(a_S\ne0\) (equivalently \(a_{\bar S}\ne0\)), and call
\(a_{\bar S}/a_S\in\{1,-1\}\) its sign.  If all active blocks have one
common sign, the same conclusion holds.  In both cases Fourier inversion
gives one physical parity and complement closure: support only on even-weight
Walsh indices gives \(q(x)=q(\bar x)\), while
\(a_{\bar S}=\epsilon a_S\) makes the inverse transform vanish unless
\((-1)^{\wt(x)}=\epsilon\).  An actual \(X\)-dressing then moves a nonzero
complementary physical pair to \(\zero^4,\one^4\).  The result is an
endpoint-nondegenerate eight-vertex signature and remains tensor-prime,
whereas every quaternary \(C_2\)-matching signature is a two-binary product.

It remains to treat mixed signs, which necessarily require at least two
active complementary blocks.  Port permutations permute the three block signs, and an
actual \(X\)-dressing reverses all of them.  We may therefore choose a
positive active block and a negative active block, make any third active
block positive, and write the Walsh coefficient scalars as \(x,y,z\), with
\(xz\ne0\), such that one \(X\)-dressing and a permutation give
\begin{equation*}
 a_{12}=a_{34}=x,\qquad
 a_{13}=a_{24}=y,\qquad
 a_{14}=z,\qquad a_{23}=-z,
 \qquad xz\ne0,
\end{equation*}
where \(y\) may vanish.  For \(p=12,13,14\), join the \(p\)-ports of two
copies by the actual path \(X-I-X\), and call the output \(r_p\).  Its
physical kernel is \(I\), its Walsh kernel is \(2I\), and direct contraction
gives
\begin{equation*}
\begin{array}{ccc}
\toprule
 &\text{common value on }0000,1111
 &\text{values on the three complementary weight-two blocks}\ \\ \midrule
 \widetilde r_{12}/4&x^2&(y^2+z^2,0,0)\\
 \widetilde r_{13}/4&y^2&(x^2+z^2,0,0)\\
 \widetilde r_{14}/4&z^2&(x^2+y^2,2xy,0).\\
\bottomrule
\end{array}
\end{equation*}
Every nonzero row is supported only on even-weight Walsh indices and is
Walsh-complement-symmetric.  The same inverse-transform identities just used
therefore make it a complement-symmetric, one-parity physical signature; an
actual \(X\)-dressing makes its two endpoints nonzero.  Hence it becomes an
endpoint-nondegenerate eight-vertex signature.

A \(C_2\)-matching tensor whose Walsh coefficients are
complement-symmetric has Walsh support on the endpoints and exactly one
weight-two block, with the two common values differing by a sign.  Indeed,
in Walsh coordinates each actual \(C_2\)
binary factor is projectively \(I\) or \(Z\).  Their tensor product on a
perfect matching has precisely that support, and Walsh-complement symmetry
forces the two factor signs to agree.

If \(y=0\) and \(x^2+z^2\ne0\), \(r_{13}\) has zero Walsh endpoints and
is not matching.  If \(x^2+z^2=0\), the nonzero part of the
\(13\mid24\) flattening is
\[
 \begin{pmatrix}x&-z\\ z&x\end{pmatrix}.
\]
All other rows and columns vanish, so this matrix is the complete nonzero
block and has rank one, contradicting tensor primeness.  The local Walsh
transforms are invertible and therefore preserve flattening rank and tensor
primeness.

If \(y\ne0\) and all three outputs were matching, the first two rows would
give \(\sigma,\tau\in\{1,-1\}\) with
\begin{equation*}
 y^2+z^2=\sigma x^2,
 \qquad
 x^2+z^2=\tau y^2.
\end{equation*}
Since \(2xy\ne0\), the last row forces \(x^2+y^2=0\).  Then the first
equation gives \(\sigma=1,z^2=2x^2\), while the second gives
\(3x^2=-\tau x^2\), impossible.  Thus an actual endpoint-nondegenerate eight-vertex signature output is outside the
\(C_2\)-matching products.

\emph{The monomial dihedral case.}
For one pair, write its four residual slices as \(A,B,C,D\).  Define the
even and odd residual-card pencils produced by the standard actual kernels by
\[
 E(t):=A+tD,\qquad O(t):=C+tB,\qquad t\in\mu_m.
\]
The pencil normal form \cref{lem:nd-root-pencil} applies verbatim to
each of these pencils.  It says that the two endpoints are proportional on
one monomial line, or, in one fixed parity plane with axes \(E_0,E_1\),
\[
 U+tV=\lambda(E_0+\alpha tE_1)
 \quad\text{or}\quad
 U+tV=\lambda(tE_0+\alpha E_1),
 \qquad \alpha\in\mu_m.
\]
Its proof uses the exact M\"obius classification
\cref{lem:nd-root-mobius}; in particular, a zero sampled value belongs only
to the proportional branch.

Apply this fact to both pencils for all six pairs.  The resulting support
conditions have exactly the following solutions: choose one total parity
and a nonempty subset of its four complementary word pairs.  Their size
histogram is
\begin{equation}
\begin{array}{crrrr}
\toprule
 |\supp q|&2&4&6&8\\ \midrule
 \#&8&12&8&2.\\
\bottomrule
\end{array}
\label{eq:equality-q4-monomial-supports}
\end{equation}
Hence the support is already one parity and complement-closed.  An actual
\(X\)-dressing makes the same copy an endpoint-nondegenerate eight-vertex signature.

\emph{The two Klein cases.}
In the Pauli Bell basis, one flattening writes
\[
 q=\sum_{\alpha,\beta\in\Ftwo^2}
 t_{\alpha\beta}B_\alpha\otimes B_\beta.
\]
Its row and column cards first make \(T=(t_{\alpha\beta})\) a partial
monomial matrix.  Intersecting the four allowed Bell-line preimages for
every kernel on the other four pairs leaves exactly the four maximal
four-planes
\begin{equation}
 V_{\pi,\epsilon}
 =\{q:q(x)=0\text{ if }\wt(x)\not\equiv\pi\pmod2,\
          q(\bar x)=\epsilon q(x)\},
 \quad \pi\in\Ftwo,\quad\epsilon\in\{1,-1\}.
 \label{eq:equality-q4-pauli-planes}
\end{equation}
Thus an actual Pauli \(X\)-dressing makes every nonzero member an endpoint-nondegenerate eight-vertex signature.

For \(\mathcal V_{\rm ex}\), define the representative matrices
\[
 E_0=I,\quad E_1=X,\quad
 E_2=\begin{pmatrix}i&1\\-1&-i\end{pmatrix},\quad
 E_3=\begin{pmatrix}-i&1\\-1&i\end{pmatrix}.
\]
Their Frobenius duals are
\(E_0/2,E_1/2,E_3/4,E_2/4\), so every dual contraction is again an
actual group direction.  The simultaneous six-pair safe locus is the union
of twenty-four maximal two-dimensional linear subspaces.  The
support-stratified calculation also produces sixteen one-dimensional
solution subspaces, but every one is contained in one of those twenty-four
planes and hence contributes no additional component to the safe locus.
Actual local exotic dressings and port permutations act transitively on
the twenty-four planes.  A representative is
\[
 q_{a,b}
 =a(E_0\otimes E_0+E_1\otimes E_1)
  +b(E_2\otimes E_2+E_3\otimes E_3),
\]
whose table is
\begin{equation}
q_{a,b}(x)=
\begin{cases}
a-2b,&x\in\{0000,0110,1001,1111\},\\
a+2b,&x\in\{0011,0101,1010,1100\},\\
0,&\wt(x)\text{ odd}.
\end{cases}
\label{eq:equality-q4-exotic-plane}
\end{equation}
If \(a-2b\ne0\), this is an endpoint-nondegenerate eight-vertex signature.  If \(a=2b\ne0\), actual \(X\)-dressings
on the last two ports make both endpoints \(4b\).  This proves the exotic
case.

In each of the latter three cases, the displayed endpoint-nondegenerate eight-vertex signature is obtained
from the original \(q\) only by invertible group dressings and a port
permutation.  It is therefore tensor-prime.  Since every quaternary
\(G(\Lambda)\)-matching signature is a tensor product of two nonsingular
group binaries, these endpoint-nondegenerate eight-vertex signatures are not \(G(\Lambda)\)-matching.  The
strengthened conclusion in the \(C_2\) case was proved separately above.

The support and subspace completeness claims in
\cref{eq:equality-q4-monomial-supports,eq:equality-q4-pauli-planes,%
eq:equality-q4-exotic-plane} use the exact enumeration in
\path{certificates/verify_equality_q4.py}.  It tests all \(2^{16}-1\)
nonempty computational supports in the monomial case and all \(209\)
partial-injection supports in each Bell basis.  For the latter, it
intersects the allowed preimages for four cross-pairs and four kernels,
retaining zero cards and reducing the equations over \(\mathbb Q\) or
\(\mathbb Q(i)\).  The preceding pencil proof supplies the coefficient
constraint before this finite support classification.
\end{proof}

\begin{theorem}[Equality-accessible quaternary interface]
\label{thm:cert-interface-equality-q4}
Let \((\Lambda,G,q)\) be a deck-stable retained, factor-saturated state
containing the distinguished literal equality \(I\), and let \(q\ne0\) be a
tensor-prime quaternary core.  Suppose every physical card is zero or a
nonzero scalar multiple of one of the fixed actual representatives of \(G\), where
\(G\) is \(C_2\), a monomial dihedral group of rotation order at least three,
the standard Pauli Klein group, or the exotic Klein group.  Then an actual
gadget using at most two copies of \(q\) realizes an
endpoint-nondegenerate eight-vertex signature that is not \(G\)-matching.
More precisely:
\begin{enumerate}[label=\textnormal{(\roman*)}]
\item In the \(C_2\) row, put
      \(a_S:=\widetilde q(\one_S)\).  A complementary weight-two block
      \(S,\bar S\) is \emph{active} when \(a_S\ne0\), and its sign is
      \(a_{\bar S}/a_S\in\{1,-1\}\).  If the two Walsh endpoints
      \(a_\varnothing,a_{[4]}\) are nonzero or all active blocks have one
      common sign, one dressed copy suffices.  Otherwise at least two active
      blocks have different signs; this is the \emph{mixed-sign} case.  For
      \(p\in\{12,13,14\}\), let \(r_p\) join the \(p\)-ports of two copies
      by the actual path \(X-I-X\).  One of these three contractions is
      endpoint-nondegenerate and nonmatching.
\item In the monomial-dihedral row, the safe support is exactly one total
      parity together with a nonempty subset of its four complementary word
      pairs.  One actual \(X\)-dressing makes the same copy
      endpoint-nondegenerate.
\item In either Klein row, actual dressings and port permutations give the
      asserted tensor-prime endpoint-nondegenerate nonmatching copy.
\end{enumerate}
All clauses include zero cards, both orientations, all six contracted pairs,
and the fixed nonzero scalars of the actual representatives.
\end{theorem}

\begin{proof}
The overall conclusion is \cref{prop:cert-equality-q4-boundary}.  Its
\(C_2\), monomial-dihedral, and two Klein proof cases give clauses
\textnormal{(i)}, \textnormal{(ii)}, and \textnormal{(iii)}, respectively.
The containment sentence in the exotic-Klein case makes its
representative-plane argument exhaustive.
\end{proof}

\subsection{Marked Dihedral Matching Decks}
\label{appsec:cert-marked-dihedral}
\label{proofsubsec:cert-normalized-dih-q6}

The six-port support classification below is used by the standard
\(\mathcal M_m\) branch.  The sector identities in
\cref{sec:marked-dihedral-sectors} transport it to \(\mathcal B_m\) and
\(\mathcal A_m\), and \cref{thm:marked-A0-reduction} transports it to
\(\mathcal A_m^0\).  We reduce its \(2304\) labelled parameter choices to
the explicit \(48\)-case table below.

Fix the distinguished pair \(p=\{1,2\}\) and residual ports
\(R=\{3,4,5,6\}\).  For each sector
\(\epsilon\in\Ftwo\), choose a perfect matching
\[
 M_\epsilon=\{e_\epsilon,e'_\epsilon\}
 \quad\text{of }R,
\]
distinguish the fixed edge
\(e_\epsilon=\{u_\epsilon,v_\epsilon\}\), and order the other edge as
\(e'_\epsilon=(r_\epsilon,s_\epsilon)\) with
\(r_\epsilon<s_\epsilon\).  This order is canonical and is not an
additional choice.  Choose
\(a_\epsilon,b_\epsilon,o_\epsilon\in\Ftwo\).  Define the eight-word
support \(S\subseteq\Ftwo^6\) by taking, for every
\(\epsilon,z\in\Ftwo\), the two words satisfying
\begin{equation}
 \begin{aligned}
 (x_1,x_2)&=(z,z+\epsilon),\\
 x_{u_\epsilon}+x_{v_\epsilon}&=a_\epsilon,\\
 (x_{r_\epsilon},x_{s_\epsilon})
   &=(o_\epsilon+z,o_\epsilon+z+b_\epsilon),
 \end{aligned}
 \qquad\text{over }\Ftwo.
 \label{eq:cert-nd-q6-four-flat-support}
\end{equation}
This is exactly the support form obtained from four complementary flat
endpoints: the fixed monomial factor supplies the two words in a slice,
and the two endpoints of the variable factor are complementary.

\begin{proposition}[Exact normalized-dihedral six-to-four support lowering]
\label{prop:cert-nd-q6-support-lowering}
Among the
\[
 3^2\cdot2^8=2304
\]
labelled choices in \cref{eq:cert-nd-q6-four-flat-support}, exactly \(96\)
give a matching coset: for some perfect matching \(M\) and constants
\(c_e\in\Ftwo\), the support is
\[
  \{x\in\Ftwo^6:x_u+x_v=c_{\{u,v\}}
      \text{ for every }\{u,v\}\in M\}.
\]
For every other choice there are a port pair \(e=\{e_1,e_2\}\) and
\(\gamma\in\Ftwo\) such that projection away from \(e\) is injective on
\[
 S\cap\{x:x_{e_1}+x_{e_2}=\gamma\},
\]
and the projected quaternary support has size six or eight.  Choosing a
largest such projection gives \(1152\) support-six cases and \(1056\)
support-eight cases.
\end{proposition}

\begin{proof}
Permute the four residual ports so that the fixed edge in sector zero is
\(34\), and its ordered varying edge is \((5,6)\).  Independently flipping
bits on ports \(3,5,6\) removes \(a_0,o_0,b_0\), respectively: use the
flips \(a_0,o_0,o_0+b_0\).  Thus sector zero is
\[
 S_0=\{(z,z,w,w,z,z):z,w\in\Ftwo\}.
\]
These permutations and bit flips preserve matching cosets, projection
injectivity, and projected support sizes.  For each of the
\(3\cdot2\cdot2^3=48\) choices of sector-zero data, fix one such
normalization.  It bijects the \(48\) choices of sector-one data with the
following canonical choices: a fixed edge \(e\subseteq\{3,4,5,6\}\),
the increasingly ordered complementary edge \(e'=(r,s)\), and
\((a,b,o)\in\Ftwo^3\).  Their four words satisfy
\[
 (x_1,x_2)=(z,z+1),\qquad
 x_u=w,\quad x_v=w+a\quad(e=\{u,v\},\ u<v),\qquad
 (x_r,x_s)=(o+z,o+z+b).
\]

The table gives a projection witness for every canonical choice.  An
entry \(ij_\gamma^s\) means: retain words with \(x_i+x_j=\gamma\),
delete ports \(i,j\), and obtain \(s\) distinct projected words from
exactly \(s\) retained words.  Thus the projection is injective.
\(\mathrm M\) denotes a matching coset.
\begin{center}
\small
\setlength{\tabcolsep}{4pt}
\begin{tabular}{c|cccccccc}
\toprule
\(e\backslash(a,b,o)\)&000&001&010&011&100&101&110&111\\
\midrule
\(34\)&\(15_0^8\)&\(25_0^8\)&\(\mathrm M\)&\(\mathrm M\)&\(15_0^8\)&\(25_0^8\)&\(15_0^8\)&\(16_0^8\)\\
\(35\)&\(16_0^8\)&\(26_0^8\)&\(24_1^6\)&\(14_1^6\)&\(14_0^6\)&\(24_0^6\)&\(26_0^8\)&\(16_0^8\)\\
\(36\)&\(15_0^8\)&\(25_0^8\)&\(24_1^6\)&\(14_1^6\)&\(14_0^6\)&\(24_0^6\)&\(25_0^8\)&\(15_0^8\)\\
\(45\)&\(16_0^8\)&\(26_0^8\)&\(23_1^6\)&\(13_1^6\)&\(13_0^6\)&\(23_0^6\)&\(26_0^8\)&\(16_0^8\)\\
\(46\)&\(15_0^8\)&\(25_0^8\)&\(23_1^6\)&\(13_1^6\)&\(13_0^6\)&\(23_0^6\)&\(25_0^8\)&\(15_0^8\)\\
\(56\)&\(13_0^6\)&\(13_1^6\)&\(15_0^6\)&\(15_0^6\)&\(13_0^6\)&\(13_1^6\)&\(13_0^6\)&\(13_1^6\)\\
\bottomrule
\end{tabular}
\end{center}
Each entry follows by substituting \((z,w)\in\Ftwo^2\) in the two
displayed sectors.  For example, at \(e=35\), \((a,b,o)=010\), the
\(24_1^6\) projection leaves
\[
 \{0001,0100,0111,1000,1011,1110\}
\]
on ports \((1,3,5,6)\), with one preimage per word.

For maximality of the entries marked \(6\), an eight-word projection
would require a pair with constant XOR on both sectors.  The pairs with
constant XOR on \(S_0\) are
\(12,15,16,25,26,34,56\).  Substitution of the sector-one equations in
each \(6\)-entry eliminates these pairs, or leaves a projection that
identifies two words.  A seven-word parity slice is impossible: on each
four-word affine sector its size is \(0,2\), or \(4\).  Thus the displayed
six-word images are maximal.  At the two \(\mathrm M\)-entries the
matching equations are, respectively,
\[
 x_3=x_4,\ x_1=x_5,\ x_2=x_6;
 \qquad
 x_3=x_4,\ x_1=x_6,\ x_2=x_5.
\]
A parity contraction of a three-edge matching coset has at most four
projected words, so none of the other entries is a matching coset.
The canonical table therefore has \(2,24,22\) cases of the three stated
types.  Multiplying by the \(48\) sector-zero choices gives
\[
 48(2,24,22)=(96,1152,1056),
\]
as required.
\end{proof}

\begin{remark}[From support projection to contraction]
\label{rem:cert-nd-q6-support-scope}
The selected parity is implemented by an actual normalized \(I\)-kernel
for \(\gamma=0\) and an \(X\)-kernel for \(\gamma=1\).  Injectivity leaves
one nonzero summand for each projected word, so the conclusion holds for
arbitrary nonzero coefficients on the support.  The main proof supplies
the sector-ruling hypothesis and exposes any resulting factors by
\cref{lem:factor-saturation}.
\end{remark}

\begin{theorem}[Six-to-four contractions for normalized dihedral groups]
\label{thm:cert-interface-nd-q6}
Let \(h\) be a six-port signature whose support is a set \(S\) of the form
\cref{eq:cert-nd-q6-four-flat-support}, and suppose every coefficient of
\(h\) on \(S\) is nonzero.  Then either \(S\) is a three-edge matching
coset, or there are a port pair \(e=\{e_1,e_2\}\) and
\(\gamma\in\Ftwo\) for which projection away from \(e\) is injective on
\[
 S\cap\{x:x_{e_1}+x_{e_2}=\gamma\}
\]
and has image of size six or eight.  In the latter case, contraction by the
actual kernel \(I\) for \(\gamma=0\), or \(X\) for \(\gamma=1\), directly
realizes a cancellation-free quaternary of support size six or eight.  That
      output is not a product of two nonsingular monomial binaries.
\end{theorem}

\begin{proof}
\Cref{prop:cert-nd-q6-support-lowering} gives the exhaustive support
alternative and injectivity.  The selected parity kernel leaves exactly one
nonzero summand above each projected word, proving noncancellation.  A product
of two nonsingular monomial binaries has support size four.
\end{proof}

\subsection{Klein-Four Matching Decks}
\label{appsec:cert-klein-four}
\label{proofsubsec:extended-v4-certificates}

The following certificate supplies the exotic Klein finite cases used
in \cref{sec:v4-deck}.

\subsubsection{The exotic Klein form
  \texorpdfstring{\(\mathcal V_{\rm ex}\)}{Vex}}
\label{proofsubsec:cert-exotic-v4}

\begin{proposition}[Exotic-Klein six- and eight-port certificate]
\label{prop:cert-exotic-v4-q6-q8}
Use the four exotic binary directions \(E_a\), their Frobenius duals, and
the label set \(V=\Ftwo^2\) from
\cref{eq:exotic-v4-basis,eq:exotic-v4-Gram}.  The following finite
statements hold.
\begin{enumerate}[label=\textnormal{(\roman*)}]
\item Let \(h\ne0\) have arity six, suppose every exotic pair-card is zero
      or a product of two exotic binaries, and normalize one nonzero card by
      \[
        C_{12}^0=E_0^{34}E_0^{56}.
      \]
      The other three card directions range over the complete normalized
      state space
      \[
        \bigl(\{0\}\cup\mathfrak D_4(\mathcal V_{\rm ex})\bigr)^3,
        \qquad (1+48)^3=117{,}649.
      \]
      Over arbitrary algebraic nonzero-card scales, including every
      cancellation locus in every other card context, exactly seventeen
      projective scale lines survive.  Every one reconstructs a product of
      three exotic binaries; no tensor-prime arity-six line survives.
\item Let \(h\) have arity eight and suppose the analytic argument in the
      main text has already placed it in the displayed normal form \cref{eq:exotic-v4-q8-linear-normal-form}, with one matching and
all four cards nonzero.  If every one of its
      \(28\cdot4=112\) exotic pair-cards is zero or a product of three
      exotic binaries, then \(h\) has a genuine proper tensor factor.
      The four specified cards are required to be nonzero.
\end{enumerate}
\end{proposition}

\begin{proof}
All matching lines and card maps below are generated directly from the four
matrices in \cref{eq:exotic-v4-basis}.  The calculation first verifies the
Gram matrix, the dual basis
\[
 (E_0/2,E_1/2,E_3/4,E_2/4),
\]
and projective fusion closure, so its contractions use the same normalized
directions as the main argument.

For part~\textnormal{(i)}, there are \(3\cdot4^2=48\) projective
quaternary matching lines.  Encode zero by \(-1\), and encode the line with
matching \(m\in\{0,1,2\}\) and labels \(a,b\in V\) by \(16m+4a+b\).
Thus a nonnegative code denotes the canonically normalized literal line
\(P_m(E_a,E_b)\).  On each base residual edge \(34\) and \(56\), the
physical stabilizer contains the paired endpoint dressings with label pairs
\[
 (1,1),\qquad(2,3),\qquad(3,2),
\]
which fix the base \(E_0\)-edge projectively.  Together with the port
permutations
\[
 (3\,4),\qquad(5\,6),\qquad(3\,5)(4\,6),
\]
these generate the physical stabilizer used here.  Its
induced action on the 48 residual matching lines has order \(64\).  It
partitions
the full \(49^3\) card-direction assignments into 4,045 orbits, and the generated
orbit-size histogram is
\[
 1^{27},\quad2^{49},\quad4^{351},\quad8^{521},\quad
 16^{577},\quad32^{1830},\quad64^{690}.
\]
The orbit sizes weighted by these multiplicities sum to \(117\,649\), so
the generated orbits cover every state.  On one representative of every
orbit, exact \(\mathbb Q(i)\)-row reduction computes the union of the
preimages of the forty-eight matching lines within each nonbase pair/kernel
context, then intersects these unions over all 56 contexts.  A zero
output lies in every matching-line preimage, and nonzero-card scales remain
independent linear variables; in particular, no cancellation locus is divided away.

The only surviving direction triples are
\begin{equation*}
 (-1,-1,-1),\quad (1,2,3),\quad(1,3,2),\quad
 (4,8,12),\quad(4,12,8).
\end{equation*}
The first has the single scale line \((1)\).  Each of the other four has
exactly the four scale lines
\begin{equation*}
 (1,\epsilon,\delta/2,\epsilon\delta/2),
 \qquad \epsilon,\delta\in\{1,-1\}.
\end{equation*}
Thus the nonzero-card histogram consists of one line with a single nonzero card
and sixteen lines with every card nonzero.  An independent literal generation of all
\(15\cdot4^3=960\) three-binary product lines reconstructs exactly these
seventeen tensors, proving part~\textnormal{(i)}.

For part~\textnormal{(ii)}, write the stipulated normal form as
\[
 h=\sum_{\gamma\in\Ftwo^2}\lambda_\gamma
 E_\gamma^{12}E_{A_1\gamma}^{34}E_{A_2\gamma}^{56}E_{A_3\gamma}^{78},
 \qquad A_j\in\GL_2(\Ftwo),\quad\lambda_\gamma\ne0.
\]
Contracting \((3,5)\) and \((4,5)\) by \(E^0=I/2\) gives, respectively,
label maps \(D_1=\vartheta A_1+A_2\) and \(D_2=A_1+A_2\), where
\(\vartheta\ne I\) is the transpose permutation from
\cref{eq:exotic-v4-transpose-fusion}.  If both vanished, then
\(\vartheta A_1=A_1\), contradicting invertibility of \(A_1\).
Choose \(D\ne0\) among them.  The resulting six-port card has four
nonzero terms with labels
\[
 (\gamma,D\gamma,A_3\gamma),\qquad\gamma\in\Ftwo^2.
\]
The terms are distinct and all three edge coordinates vary.  By
\cref{eq:exotic-v4-overlap-size}, a four-term matching product can vary
on only two edges.  This contradicts the assumed safe-card condition
and proves part~\textnormal{(ii)}.  This argument applies to arbitrary
nonzero \(\lambda_\gamma\) and uses only two specified actual contractions.

Every scale condition in part~\textnormal{(i)} is a finite union of homogeneous
linear systems over \(\mathbb Q(i)\).  Prescribed liveness is also exact:
over an infinite field, a linear space contains a vector avoiding all
coordinate hyperplanes exactly when no coordinate vanishes identically on
that space.  Base change to \(\overline{\mathbb Q}\), or to
\(\mathbb C\), therefore creates no additional component, and zero target
cards are kept as zero-card subspaces.  The exact calculation for part~\textnormal{(i)} regenerates the states,
orbits, scale intersections, and factor tensors.
The companion exact replay is
\path{certificates/verify_exotic_v4_q6_q8.py}.
\end{proof}

\begin{theorem}[Six-port tensors for the exotic Klein group]
\label{thm:cert-interface-exotic-v4-q6}
Let \(h\ne0\) have arity six.  If every exotic pair-card of \(h\) is zero
or a product of two exotic binaries, then \(h\) is a product of three
exotic binaries.  Consequently no tensor-prime arity-six survivor exists.
This includes arbitrary algebraic scales and every cancellation locus.
\end{theorem}

\begin{proof}
This is \cref{prop:cert-exotic-v4-q6-q8}\textnormal{(i)} after undoing the
actual normalization of one nonzero card.
\end{proof}

\begin{theorem}[Eight-port tensors for the exotic Klein group]
\label{thm:cert-interface-exotic-v4-q8}
Let \(h\) have arity eight and satisfy the displayed normal form \cref{eq:exotic-v4-q8-linear-normal-form}, with one matching and
all four cards nonzero.  If all \(112\) exotic
pair-cards of \(h\) are zero or products of three exotic binaries, then
\(h\) has a genuine proper tensor factor.  
\end{theorem}

\begin{proof}
This is \cref{prop:cert-exotic-v4-q6-q8}\textnormal{(ii)}.
\end{proof}

\subsection{Platonic Matching Decks}
\label{appsec:cert-platonic}
\label{proofsec:platonic-matching-certificates}

All marked Platonic configurations use the quaternion basis fixed in
\cref{eq:quaternion-basis}.
Whenever a mixed three-point incidence is localized, actual dressings put
its rainbow image in the common normalization
\begin{equation}
 A=P_0(Q_0,Q_0),\qquad
 B=P_1(Q_0,Q_0),\qquad
 C=P_2(Q_2,Q_2)=A-B.
 \label{eq:cert-A4-normal-incidence}
\end{equation}
For a rich line in a quaternary matching deck, its
\emph{matching-label type} is the partition of the line size obtained by
counting points on each of the three residual perfect matchings.  Thus
\((1,1,1)\) is a rainbow three-line, while \((3)\), \((4)\), and \((5)\)
are monochromatic lines.  We also fix, before their first use,
\begin{equation}
\begin{gathered}
 R_0=\tfrac12P_2(Q_0,Q_0),\qquad
 R_1=-\tfrac12P_2(Q_1,Q_1),\qquad
 R_3=-\tfrac12P_2(Q_3,Q_3),\\
 G_+=Q_0+Q_1+Q_2+Q_3,\qquad
 G_-=Q_0-Q_1-Q_2-Q_3.
\end{gathered}
\label{eq:cert-platonic-common-vectors}
\end{equation}

\begin{lemma}[The matching-product bridge]
\label{lem:platonic-bridge-formula}
Let \(S\) be the swap matrix, defined by \(S(u\otimes v)=v\otimes u\).
For \(U,V\in\GL_2(\mathbb C)\) and \(\lambda\ne0\),
\begin{equation}
 \rank\bigl((U\otimes V^{\mathsf T})S-\lambda I_4\bigr)=1
 \quad\Longleftrightarrow\quad V=\lambda U^{-\mathsf T}.
 \label{eq:platonic-swap-rank-one}
\end{equation}
For \(G\in\{\mathcal T,\mathcal O,\mathcal I\}\), put
\(\mathscr C_G=\{0\}\cup\operatorname{cone}(\mathfrak D_4(G))\).
Then
\begin{equation}
 \{W:W,A-2W,2W-B\in\mathscr C_G\}
 =\{0,A/2,B/2\}\cup
 \left\{\tfrac12P_2(J,J^{-\mathsf T}):[J]\in G,
                                      \ J^{\mathsf T}=J\right\}.
 \label{eq:platonic-symmetric-bridge}
\end{equation}
Write a symmetric group matrix as \(J_u=aQ_0+bQ_1+dQ_3\), choosing a
nonzero real coefficient vector \(u=(a,b,d)\), and put
\(R_u=\tfrac12P_2(J_u,J_u^{-\mathsf T})\).
A triple of bridge vectors sums to \((A+B)/2\) precisely when it is a
permutation of \((0,A/2,B/2)\), or is \((R_u,R_v,R_w)\) with
\(u,v,w\) pairwise orthogonal in \(\mathbb R^3\).
\end{lemma}

\begin{proof}
Conjugation by \(U^{-1}\otimes I\) reduces the rank condition to
\((I\otimes W^{\mathsf T})S-\lambda I_4\), with
\(W=U^{\mathsf T}V=\begin{psmallmatrix}a&b\\c&d\end{psmallmatrix}\).
Its matrix is
\[
 \begin{pmatrix}
 a-\lambda&0&c&0\\
 b&-\lambda&d&0\\
 0&a&-\lambda&c\\
 0&b&0&d-\lambda
 \end{pmatrix}.
\]
The second row has nonzero second entry.  Rank one forces the first row
to vanish, hence \(a=\lambda,c=0\); comparing the second and third rows
then gives \(b=0,d=\lambda\).  Conversely, \(S-I_4\) has rank one.
This proves \cref{eq:platonic-swap-rank-one}.

The three exceptional bridge vectors \(0,A/2,B/2\) satisfy the conditions,
since \(C=A-B=P_2(Y,Y)\) and \([Y]\in G\).  For any other solution,
all three tested vectors are nonzero.  A line with two points of one
matching and a point of another would give flattening ranks \(1,1,4\),
which is impossible.  If \(W=P_0(U,V)\), the condition on \(2W-B\)
therefore gives matching \(2\).  Its rank-one flattening and
\cref{eq:platonic-swap-rank-one} imply
\(W=\tfrac12P_0(U,U^{-\mathsf T})\).
The condition on \(A-2W\) has matching \(0\), so
\(I\otimes I-U\otimes U^{-\mathsf T}\) is simple.  The simple-tensor
sum criterion forces \(U\) to be scalar, giving \(W=A/2\), a
contradiction.  The matching-\(1\) case similarly gives \(W=B/2\).

Thus \(W=P_2(U,V)\), and \(A-2W\) has matching \(1\).
Its \(13\mid24\) flattening is
\(I_4-2(U\otimes V^{\mathsf T})S\); hence
\(W=\tfrac12P_2(U,U^{-\mathsf T})\).
Apply the invertible local matrices \((U^{-1},I,U^{\mathsf T},I)\),
which fix \(B\).  They send \(2W\) to
\(P_2(I,U^{-\mathsf T}U)\).  Since \(2W-B\) has matching \(0\),
\cref{eq:platonic-swap-rank-one} with \(\lambda=1\) now gives
\(U^{-\mathsf T}U=I\), or \(U^{\mathsf T}=U\).
Conversely, for symmetric \(U\), the two differences are local images
of \(P_1(Y,Y)\) and \(-P_0(Y,Y)\).  All leg matrices belong projectively
to \(G\), so both differences are in its matching cone.  This proves
\cref{eq:platonic-symmetric-bridge}, including arbitrary scales.

For the sum equation, polarization of
\(\tfrac12P_2(J_u,\overline J_u)\), where
\(\overline J_u=aQ_0-bQ_1-dQ_3\), defines an injective linear map
\(T\) from the symmetric \(3\times3\) matrices to four-port tensors with
\[
 T(uu^{\mathsf T})=\tfrac12P_2(J_u,\overline J_u),\qquad
 T(I_3)=(A+B)/2,\qquad
 R_u=T\!\left(\frac{uu^{\mathsf T}}{u^{\mathsf T}u}\right).
\]
The six coefficient tensors are independent and \(C\notin\operatorname{im}T\):
in the \(14\mid23\) binary basis they use only
\(Q_\mu\otimes Q_\nu\) with \(\mu,\nu\in\{0,1,3\}\), whereas
\(C=Q_2\otimes Q_2\).  Also
\[
 A/2=\tfrac12T(I_3)+C/4,\qquad
 B/2=\tfrac12T(I_3)-C/4.
\]
Comparison of the \(C\)-coefficient forces equal numbers of \(A/2\)
and \(B/2\); if they occur, the third summand is zero.  Otherwise the
trace of \(T^{-1}R_u\) is one, so all three summands are nonzero
\(R\)-vectors.  Their sum is \(T(I_3)\) exactly when their normalized
real coefficient columns form a matrix \(E\) with
\(EE^{\mathsf T}=I_3\).  This is equivalent to
\(E^{\mathsf T}E=I_3\), proving the orthogonality criterion.
The three orthogonal rank-one projectors are independent; together with
\(C\notin\operatorname{im}T\), they give card-map rank four.
The exceptional triple spans \(A,B\) and gives rank two.
\end{proof}

\subsubsection{The four fixed-equality octahedral forms
  \texorpdfstring{\(\mathcal O,\mathcal O',\mathcal O_{dt},\mathcal O_{td}\)}%
  {O, O-prime, O-dt, and O-td}}
\label{proofsubsec:cert-S4}

For the octahedral certificate, use the standard \(24\)-point set
\(\mathcal O\) in
\cref{eq:quaternion-basis,eq:octahedral-domain}, ordered as the four axes, the
eight sign points, and the twelve pair points.  It is the quaternion model of \(S_4\) fixed in the preliminaries.  The second
unmixed distinguished-involution normal form is
\begin{equation}
 \mathcal O'=H_{\mathcal O}^{-1}\mathcal O H_{\mathcal O}.
 \label{eq:cert-S4-second-form}
\end{equation}
Use the coefficient-similarity map \(\theta_{\mathcal O}\), transformation
tuple \(\mathbf H_{\mathcal O}\), and induced coefficient operator
\(\mathcal L_{\mathcal O}\) fixed in
\cref{eq:octahedral-second-form,eq:octahedral-coefficient-map}; the
similarity induces the corresponding map on projective group classes.
Here \(P'_m\) denotes the same ordered matching placement as \(P_m\), but
in the \(\mathcal O'\) coefficient coordinates.
Writing \(X=Q_1\) projectively, direct ordered contraction gives
\begin{equation}
\begin{aligned}
\mathcal L_{\mathcal O}P_0(U,V)&\doteq
 P'_0(\theta_{\mathcal O}(U),\theta_{\mathcal O}(V)),\\
\mathcal L_{\mathcal O}P_1(U,V)&\doteq
 P'_1(\theta_{\mathcal O}(UX),\theta_{\mathcal O}(XV)),\\
\mathcal L_{\mathcal O}P_2(U,V)&\doteq P'_2(\theta_{\mathcal O}(U),
                                  \theta_{\mathcal O}(XVX)).
\end{aligned}
\label{eq:cert-S4-coefficient-transport}
\end{equation}
The operator \(\mathcal L_{\mathcal O}\) compares coefficient arrays.
Actual dressings use the corresponding \(\theta_{\mathcal O}(U)\)
representatives with their native-wire orientations and nonzero scalars,
as in the fixed-coordinate convention of \cref{sec:platonic-decks}.

Order the 24 points as the four axes, the eight sign points, and the twelve
pair points \(Q_\mu\pm Q_\nu\), with pairs lexicographic.  The coset classification in \cref{prop:cert-platonic-incidence-atlas} gives
\begin{center}
\small
\begin{tabular}{c@{\qquad}c@{\qquad}c}
\toprule
domain points&three-point lines&four-point lines\\
\midrule
24&32&18\\
\bottomrule
\end{tabular}
\end{center}
Every point lies on seven rich lines.  The three-lines and
four-lines are single actual left/right orbits, represented by
\[
 \{[Q_0],[G_+],[G_-]\}
\]
and
\[
 \{[Q_0],[Q_1],[Q_0+Q_1],[Q_0-Q_1]\},
\]
respectively.

The quaternary deck has \(3\cdot24^2=1728\) projective points.  Its complete
maximal-rich-line histogram is
\begin{center}
\small
\begin{tabular}{ccr}
\toprule
line size&matching-label type&number\\
\midrule
3&\((1,1,1)\)&13\,824\\
3&\((3)\)&4\,608\\
4&\((4)\)&2\,592\\
\midrule
&&21\,024\\
\bottomrule
\end{tabular}
\end{center}
There is no \((2,1)\) line and no mixed four-line.  Representatives of
the three rows are
\begin{align*}
&\{P_0(Q_0,Q_0),P_1(Q_0,Q_0),P_2(Q_2,Q_2)\},\\
&\{P_0(Q_0,Q_0),P_0(Q_0,G_+),P_0(Q_0,G_-)\},\\
&\{P_0(Q_0,Q_0),P_0(Q_0,Q_1),
       P_0(Q_0,Q_0+Q_1),P_0(Q_0,Q_0-Q_1)\},
\end{align*}
where \(G_+\) and \(G_-\) are the two explicitly fixed representatives in
\cref{eq:cert-platonic-common-vectors}.

For a mixed three-point incidence, use the normalization
\eqref{eq:cert-A4-normal-incidence}.  For a selected ordered frame, let
\(\Phi\) be the linear card map determined by \(\Phi(Q_0)=C/2\) and by
\((\Phi(Q_1),\Phi(Q_2),\Phi(Q_3))\) equal to that ordered frame; define the
image tensors \(W_j:=\Phi(Q_j)\), and write
\(\rank\Phi:=\dim\operatorname{im}\Phi\).  Before imposing the twelve pair
points of \eqref{eq:octahedral-domain}, the bridge is
\begin{equation*}
 \{0,A/2,B/2\}\cup\mathcal R,
\end{equation*}
where the subscripts in \(R_{\mu\nu}^{\epsilon\delta}\) specify the
two quaternion axes and the superscripts specify the displayed signs:
\begin{align*}
\mathcal R=\{&R_0,R_1,R_3,\\
&R_{01}^{+-}=\tfrac14P_2(Q_0+Q_1,Q_0-Q_1),
\quad R_{01}^{-+}=\tfrac14P_2(Q_0-Q_1,Q_0+Q_1),\\
&R_{03}^{+-}=\tfrac14P_2(Q_0+Q_3,Q_0-Q_3),
\quad R_{03}^{-+}=\tfrac14P_2(Q_0-Q_3,Q_0+Q_3),\\
&R_{13}^{++}=-\tfrac14P_2(Q_1+Q_3,Q_1+Q_3),
\quad R_{13}^{--}=-\tfrac14P_2(Q_1-Q_3,Q_1-Q_3)\}.
\end{align*}
By \cref{lem:platonic-bridge-formula}, these are all bridge vectors:
the symmetric octahedral directions are the three axes and the six
\(Q_\mu\pm Q_\nu\) directions with \(\mu,\nu\in\{0,1,3\}\).
An orthogonal triple is either the three axes or a pair of opposite
diagonal directions in one coordinate plane together with the remaining
axis.  There are therefore four unordered frames.  The sum equation
\(W_1+W_2+W_3=(A+B)/2\) leaves 30 ordered maps: the six
\(\rank\Phi=2\) maps permute \((0,A/2,B/2)\), and the other twenty-four are
the six orderings of each of the four \(\rank\Phi=4\) frames
\begin{equation}
 \begin{split}
 &\{R_0,R_1,R_3\},\qquad
 \{R_0,R_{13}^{++},R_{13}^{--}\},\\
 &\{R_1,R_{03}^{+-},R_{03}^{-+}\},\qquad
 \{R_3,R_{01}^{+-},R_{01}^{-+}\}.
 \end{split}
 \label{eq:cert-S4-four-frames}
\end{equation}
The 1,727 deck points other than \([A]\) lie on 1,683 projective lines
through \([A]\); the numbers of other deck points on those lines are
\[
 1^{1645},\qquad2^{32},\qquad3^6.
\]
The single actual kernel
\begin{equation*}
 K_*=Q_0+Q_1,
 \qquad [K_*]\in\mathcal O
\end{equation*}
has the following property: for every ordered rank-four frame,
\(q_R=\Phi(K_*)=C/2+R\) is not a matching-deck point.  The 24 supports
have sizes
\[
 2^4,\qquad8^{12},\qquad16^8.
\]
The second normal form is identified analytically by
\cref{eq:cert-S4-coefficient-transport} and handled physically by the actual
\(\mathcal O'\)-kernels and dressings in
\cref{prop:p1-cert-S4-uniform-eight-vertex-separator}; no naive conjugation
of a contraction kernel is used.  Invertible local coefficient maps
preserve the three ranks.

For the two mixed forms, use the coefficient similarities
\(H_{dt},H_{td}\), groups \(\mathcal O_{dt},\mathcal O_{td}\), and
operators \(\mathcal L_{dt},\mathcal L_{td}\) fixed in
\cref{eq:octahedral-mixed-conjugators,eq:octahedral-mixed-forms,%
eq:octahedral-mixed-coefficient-transport}.  The three exact transport
identities there identify every matching label and absorb their displayed
orientation factors into actual matrices of the current mixed group.  They
are used only to compare coefficient tables; neither similarity is applied
to the retained signature set or to the supplied literal \(I\).

\begin{proposition}[Physical octahedral signature-separator certificate]
\label{prop:p1-cert-S4-uniform-eight-vertex-separator}
For every ordering of every frame in
\cref{eq:cert-S4-four-frames}, use the actual kernel
\(K_*=Q_0+Q_1\) in the \(\mathcal O\) form and the actual kernel
\(K'_*:=\theta_{\mathcal O}(K_*)\) in the \(\mathcal O'\) form.  The
resulting direct quaternary can be dressed, using only actual local
matrices in the current marked octahedral group, to an endpoint-nondegenerate eight-vertex signature of
flattening ranks \((2,2,2)\).  Similarity by \(H_{\mathcal O}\)
is used only
to identify coefficient tables; it is never used as a local gadget.

For \(\alpha\in\{dt,td\}\), put
\(K_{*,\alpha}:=\theta_\alpha(K_*)\), so
\([K_{*,\alpha}]\in\mathcal O_\alpha\), and use the fixed actual kernel in
that projective class, with its realization scalar recorded.  The resulting
direct quaternary has, using only actual current-group contractions and
dressings, one of three certified exits: an endpoint-nondegenerate
eight-vertex signature of ranks \((2,2,2)\), a pure generalized equality,
or a nonsingular binary \(B\) with \([BX]\notin\mathcal O_\alpha\).
The same three-way conclusion holds for each of the six rank-deficient maps in
each mixed form.  All representatives, determinants, and outside-transfer
directions are given in the proof below.
\end{proposition}

\begin{proof}
Define the tetrahedral sign matrices by
\[
  G_{\epsilon_1\epsilon_2\epsilon_3}
  =Q_0+\epsilon_1Q_1+\epsilon_2Q_2+\epsilon_3Q_3.
\]
For the first marked form the separator is \(q_R=C/2+R\), where \(R\)
is the first member of the ordered frame.  In the even-word order
\cref{eq:octahedral-even-order}, exact substitution gives the following
actual dressings; every omitted odd-weight coefficient is zero, and the
displayed vectors are projective.
\begin{center}
\small
\setlength{\tabcolsep}{4pt}
\begin{tabular}{ccc}
\toprule
\(R\)&\(\mathbf L\)&\(\mathbf e(\mathscr D_{\mathbf L}q_R)\)\\
\midrule
\(R_0\)&\(I^{\otimes4}\)&\((1,1,-1,1,1,-1,1,1)\)\\
\(R_1\)&\((I,I,Q_1,Q_1)\)&\((1,0,0,0,0,0,0,1)\)\\
\(R_3\)&\(I^{\otimes4}\)&\((1,1,-1,-1,-1,-1,1,1)\)\\
\(R_{01}^{+-}\)&\(G_{+++}^{\otimes4}\)&\((1,-1,1,-i,i,1,-1,1)\)\\
\(R_{01}^{-+}\)&\(G_{+++}^{\otimes4}\)&\((1,-1,1,i,-i,1,-1,1)\)\\
\(R_{03}^{+-}\)&\(I^{\otimes4}\)&\((1,1,-1,i,-i,-1,1,1)\)\\
\(R_{03}^{-+}\)&\(I^{\otimes4}\)&\((1,1,-1,-i,i,-1,1,1)\)\\
\(R_{13}^{++}\)&\(G_{++-}^{\otimes4}\)&\((1,i,-i,-i,-i,-i,i,-1)\)\\
\(R_{13}^{--}\)&\(G_{++-}^{\otimes4}\)&\((1,-i,i,i,i,i,-i,-1)\)\\
\bottomrule
\end{tabular}
\end{center}
Both endpoints are nonzero in every row.  Since the separator depends
only on the first frame member, these nine rows cover all twenty-four
ordered maps.  Direct row reduction of the three flattenings of
every displayed \(q_R\) gives \((2,2,2)\).

For the second marked form, use the coefficient-transport operator
\(\mathcal L_{\mathcal O}\) above, and define the transported card map
\(\Phi'\) and
transported direct card \(q'_R\) by
\[
  \Phi'(\theta_{\mathcal O}K)=\mathcal L_{\mathcal O}\Phi(K),
  \qquad
  q'_R:=\Phi'(K'_*)=\mathcal L_{\mathcal O}q_R.
\]
Here \([K'_*]\in\mathcal O'\), so \(q'_R\) is a direct card in the
second marked form.
With the nonzero comparison scalar \(\lambda:=2\sqrt2(1-i)\), the identities
\begin{equation*}
  H_{\mathcal O}H_{\mathcal O}^{\mathsf T}=\lambda X,
  \qquad
  H_{\mathcal O}^{-\mathsf T}H_{\mathcal O}^{-1}=\lambda^{-1}X
\end{equation*}
show that \(\mathcal L_{\mathcal O}\) bijects the two coefficient decks.
This is only
an analytic correspondence.  The following port matrices are separately
verified to belong projectively to the actual group \(\mathcal O'\); in the
table, \(\theta_{\mathcal O}\) acts componentwise on every displayed tuple:
\begin{center}
\small
\setlength{\tabcolsep}{4pt}
\begin{tabular}{ccc}
\toprule
\(R\)&\(\mathbf L'\)&\(\mathbf e(\mathscr D_{\mathbf L'}q'_R)\)\\
\midrule
\(R_0\)&\(\theta_{\mathcal O}(G_{++-},G_{++-},G_{++-},G_{+-+})\)&
\((1,-1,-1,-1,-1,-1,-1,1)\)\\
\(R_1\)&\(\theta_{\mathcal O}(I,I,I,I)\)&\((1,1,1,-1,-1,1,1,1)\)\\
\(R_3\)&\(\theta_{\mathcal O}(G_{+++})^{\otimes4}\)&\((1,-1,1,1,1,1,-1,1)\)\\
\(R_{01}^{+-}\)&\(\theta_{\mathcal O}(G_{+++},G_{--+},G_{--+},G_{+++})\)&
\((1,0,0,0,0,0,0,1)\)\\
\(R_{01}^{-+}\)&\(\theta_{\mathcal O}(G_{+++},G_{+--},G_{+--},G_{+++})\)&
\((1,-1,-1,1,1,-1,-1,1)\)\\
\(R_{03}^{+-}\)&\(\theta_{\mathcal O}(I,I,Q_1,Q_1)\)&\((1,1,-1,-1,-1,-1,1,1)\)\\
\(R_{03}^{-+}\)&\(\theta_{\mathcal O}(I,I,Q_2,Q_2)\)&\((1,0,0,0,0,0,0,1)\)\\
\(R_{13}^{++}\)&\(\theta_{\mathcal O}(G_{++-},G_{-++},G_{++-},G_{-++})\)&
\((1,-1,-1,1,1,-1,-1,1)\)\\
\(R_{13}^{--}\)&\(\theta_{\mathcal O}(G_{++-},G_{---},G_{+-+},G_{-++})\)&
\((1,0,0,0,0,0,0,1)\)\\
\bottomrule
\end{tabular}
\end{center}
Again both endpoints are nonzero.  All dressings are invertible, so the
three certified ranks remain \((2,2,2)\).

It remains to verify the mixed rows directly in their fixed-
\(I\) coordinates.  For \(\alpha\in\{dt,td\}\), absorb the three
matching-dependent orientation factors in
\cref{eq:octahedral-mixed-coefficient-transport} into the actual current-
group representatives.  The kernel \(K_{*,\alpha}\) first gives the
actual card \(\mathcal L_\alpha q_R\).  On ports \(2,4\), apply the
actual current-group dressing \(\theta_\alpha(S_\alpha^{-1})\), where
\(S_\alpha=H_\alpha H_\alpha^{\mathsf T}\).  Its product with the
\(H_\alpha^{\mathsf T}\) leg of \(\mathcal L_\alpha\) is
\(H_\alpha^{-1}\), so the resulting direct card
\(q_{R,\mathrm{act}}^\alpha\) has the normalized representative
\begin{equation}
 q_R^\alpha:=(H_\alpha^{-1})^{\otimes4}q_R,
 \qquad q_{R,\mathrm{act}}^\alpha\doteq q_R^\alpha.
 \label{eq:cert-S4-mixed-direct-card}
\end{equation}
This equality follows from the displayed actual dressings; it is not a
state transformation.  When another dressing on ports \(2,4\) is listed
below, compose the two group elements into one actual representative on
that port.  The total is at most four dressings, or two dressings and one
pair closure for an outside-binary row.  All exact
determinants below refer to this normalized representative; the corresponding
actual output differs by a known nonzero scalar, so only determinant
nonvanishing is used.  Put
\[
 \delta_\alpha(B):=
 [c_Q(H_\alpha(BX)H_\alpha^{-1})].
\]
Thus \(\delta_\alpha(B)\notin\mathcal R_{\mathcal O}\) is exactly the assertion
\([BX]\notin\mathcal O_\alpha\).

For \(\mathcal O_{dt}\), the three first-frame entries in the following
table give endpoint-nondegenerate eight-vertex signatures.  All odd
coefficients vanish, and the last column is in the even-word order fixed
above.
\begin{center}
\small
\setlength{\tabcolsep}{4pt}
\begin{tabular}{ccc}
\toprule
\(R\)&actual dressing&\(\mathbf e(\mathscr D q_R^{dt})\)\\
\midrule
\(R_0\)&\(\theta_{dt}(G_{--+})^{\otimes4}\)&
 \((1,-1,1,-1,-1,1,-1,1)\)\\
\(R_1\)&\(\theta_{dt}(G_{---})^{\otimes4}\)&
 \((1,1,-1,-1,-1,-1,1,1)\)\\
\(R_3\)&identity&\((1,1,-1,-1,-1,-1,1,1)\)\\
\bottomrule
\end{tabular}
\end{center}
The other six possible first entries have the uniform outside-binary exits
listed below; the kernel acts on ports \(1,2\).
\begin{center}
\small
\setlength{\tabcolsep}{4pt}
\begin{tabular}{cccc}
\toprule
\(R\)&actual kernel&\(\delta_{dt}(B)\)&\(\det B\)\\
\midrule
\(R_{01}^{+-}\)&\(\theta_{dt}(G_{---})\)&\([1,-1,1,0]\)&
 \((51+36\sqrt2)/64\)\\
\(R_{01}^{-+}\)&\(\theta_{dt}(G_{---})\)&\([1,-1,0,-1]\)&
 \((51+36\sqrt2)/64\)\\
\(R_{03}^{+-}\)&\(\theta_{dt}(Q_0)\)&\([1,-3,1,-1]\)&
 \((51+36\sqrt2)/256\)\\
\(R_{03}^{-+}\)&\(\theta_{dt}(Q_0)\)&\([1,-3,-1,1]\)&
 \((51+36\sqrt2)/256\)\\
\(R_{13}^{++}\)&\(\theta_{dt}(Q_0)\)&\([1,-3,1,1]\)&
 \((51+36\sqrt2)/256\)\\
\(R_{13}^{--}\)&\(\theta_{dt}(Q_0)\)&\([1,-3,-1,-1]\)&
 \((51+36\sqrt2)/256\)\\
\bottomrule
\end{tabular}
\end{center}

For \(\mathcal O_{td}\), no dressing is needed for the two endpoint rows:
\begin{equation*}
\begin{aligned}
 \mathbf e(q_{R_0}^{td})&=(1,i,-i,i,i,-i,i,-1),\\
 \mathbf e(q_{R_3}^{td})&=(1,i,-i,-i,-i,-i,i,-1).
\end{aligned}
\end{equation*}
Both have zero odd coefficients and ranks \((2,2,2)\).  The \(R_1\) row is
\[
 q_{R_1}^{td}=\frac{3+2\sqrt2}{8}
  (\one[0011]+\one[1100]),
\]
so actual \(X\)-dressings on ports \(3,4\) give
\(\frac{3+2\sqrt2}{8}\Eq_4\), a pure generalized equality.
Compose the dressing on port \(4\) with its preceding correction; together
with the corrections on ports \(2,4\), this uses only three dressed ports.
The remaining rows have the following outside-binary exits on ports
\(1,2\):
\begin{center}
\small
\setlength{\tabcolsep}{4pt}
\begin{tabular}{cccc}
\toprule
\(R\)&actual kernel&\(\delta_{td}(B)\)&\(\det B\)\\
\midrule
\(R_{01}^{+-}\)&\(\theta_{td}(Q_0)\)&\([0,2,-1,1]\)&
 \((51+36\sqrt2)/256\)\\
\(R_{01}^{-+}\)&\(\theta_{td}(Q_0)\)&\([0,2,-1,-1]\)&
 \((51+36\sqrt2)/256\)\\
\(R_{03}^{+-}\)&\(\theta_{td}(G_{--+})\)&\([1,2,0,1]\)&
 \((51+36\sqrt2)/64\)\\
\(R_{03}^{-+}\)&\(\theta_{td}(G_{--+})\)&\([1,-2,0,-1]\)&
 \((51+36\sqrt2)/64\)\\
\(R_{13}^{++}\)&\(\theta_{td}(Q_0)\)&\([1,2,-1,0]\)&
 \((51+36\sqrt2)/256\)\\
\(R_{13}^{--}\)&\(\theta_{td}(Q_0)\)&\([1,-2,1,0]\)&
 \((51+36\sqrt2)/256\)\\
\bottomrule
\end{tabular}
\end{center}

Finally let \(q_\pi^\alpha\) denote the six normalized mixed reconstructed
rank-deficient outputs in the same port convention as
\cref{eq:cert-S4-mixed-direct-card}.  Put
\[
 K_{dt}^{\mathrm{def}}:=\theta_{dt}(Q_0-Q_2),
 \qquad
 K_{td}^{\mathrm{def}}:=\theta_{td}(Q_2-Q_3),
\]
and let
\(B_{\pi,\alpha}^{\mathrm{def}}
 :=\partial_{12}^{K_\alpha^{\mathrm{def}}}q_\pi^\alpha\).
The fixed actual representative of \([K_\alpha^{\mathrm{def}}]\) on ports
\(1,2\) works uniformly for all six outputs in each row:
\begin{equation*}
\begin{array}{c|c|c|c}
\alpha&K_\alpha^{\mathrm{def}}&
 \{\delta_\alpha(B_{\pi,\alpha}^{\mathrm{def}}):\pi\in S_3\}
 &\det B_{\pi,\alpha}^{\mathrm{def}}\\ \hline
dt&\theta_{dt}(Q_0-Q_2)&
 \{[1,-1,0,-1]^2,[1,-3,1,-1]^2,[1,-3,-1,-1]^2\}&
 (51+36\sqrt2)/256\\
td&\theta_{td}(Q_2-Q_3)&
 \{[0,2,1,1]^2,[1,2,0,1]^2,[1,2,1,0]^2\}&
 (51+36\sqrt2)/256.
\end{array}
\end{equation*}
Here \([a_0,a_1,a_2,a_3]\) denotes the projective quaternion-coordinate
point \([(a_0,a_1,a_2,a_3)^{\mathsf T}]\), equivalently the matrix class
\([a_0Q_0+a_1Q_1+a_2Q_2+a_3Q_3]\) under \(c_Q\), and an exponent records
multiplicity.  Every determinant displayed in the mixed tables is nonzero.
Every pulled-back direction has support three, or support four with unequal
nonzero magnitudes, whereas the quaternion-coordinate points of
\(\mathcal R_{\mathcal O}\) have support one, two, or four with equal nonzero
magnitudes.  Hence all the listed transfer
classes are outside their actual mixed group.  Direct row reduction gives
rank pattern \((2,2,2)\) in every endpoint row.  This proves all full-rank
and rank-deficient mixed exits using only actual current-form kernels and
dressings, while the similarities serve only to write the coefficient
arrays.
\end{proof}

For the fixed-matching residue in the standard form, the four axes and
eight tetrahedral sign points give the quadratic test set of
\cref{lem:platonic-global-ruling}.  In any nonstandard form, apply the same
test after precomposing \(\Phi\) with the corresponding one of
\(\theta_{\mathcal O},\theta_{dt},\theta_{td}\), and then pull the resulting
tensor identity back by its inverse.  This is a coefficient pullback only,
not a use of any similarity as a gadget.  In all four forms, the evaluation
matrix of the ten homogeneous quadratic
 monomials in the four quaternion coordinates has full column rank.  Hence
 every pulled-back Segre minor vanishes identically, giving a genuine
common factor rather than a new octahedral survivor.

\Cref{prop:cert-platonic-incidence-atlas} proves the domain and deck
incidences by cyclic cosets, tensor factorization, and actual group actions.
\Cref{lem:platonic-bridge-formula} proves the bridge intersections and
reduces their sum equations to the displayed orthogonal frames.

\subsubsection{The internal tetrahedral permutation classes}
\label{proofsubsec:cert-A4}

For the tetrahedral certificate, order the standard twelve points
\(\mathcal T\) in \cref{eq:tetrahedral-domain} as the four axes
followed by the eight sign points.  The coset classification in \cref{prop:cert-platonic-incidence-atlas}
gives 16 three-point lines in one actual left/right \(A_4\)-orbit.  The quaternary
deck has \(3\cdot12^2=432\) distinct projective points.  Exactly 1728 of
its three-point lines meet all three matching classes, again in one orbit
under actual local dressings and port permutations.  Thus the common
normalization above applies.

Scale the normalized card map
\(\Phi:\Mat_2(\mathbb C)\to(\mathbb C^2)^{\otimes4}\)
so that
\[
 \Phi(G_+)=A,
 \qquad
 \Phi(G_-)=-B.
\]
Then \(\Phi(Q_0)=C/2\); define the image tensors \(W_j:=\Phi(Q_j)\).  Thus
\(W_1+W_2+W_3=(A+B)/2\).  For each
\(W\in\{W_1,W_2,W_3\}\), require
\[
 W,\ A-2W,\ 2W-B
 \in\{0\}\cup\operatorname{cone}\bigl(\mathfrak D_4(\mathcal T)\bigr).
\]
The resulting exact bridge is
\begin{equation}
 \mathcal B=\{0,A/2,B/2,R_0,R_1,R_3\},
 \label{eq:cert-A4-bridge}
\end{equation}
where \(R_0,R_1,R_3\) are defined in
\cref{eq:cert-platonic-common-vectors}.
The remaining bridge vectors obey
\begin{center}
\small
\begin{tabular}{cL{0.32\linewidth}L{0.32\linewidth}}
\toprule
\(W\)&\(A-2W\)&\(2W-B\)\\
\midrule
\(R_0\)&\(P_1(Q_2,Q_2)\)&\(-P_0(Q_2,Q_2)\)\\
\(R_1\)&\(-P_1(Q_3,Q_3)\)&\(P_0(Q_3,Q_3)\)\\
\(R_3\)&\(-P_1(Q_1,Q_1)\)&\(P_0(Q_1,Q_1)\)\\
\bottomrule
\end{tabular}
\end{center}
\Cref{lem:platonic-bridge-formula} proves completeness of this list.
The symmetric tetrahedral directions are exactly \(Q_0,Q_1,Q_3\), so
their three coordinate axes form the unique orthogonal frame.

The sum equation on \(\mathcal B^3\) leaves precisely twelve maps:
\begin{equation}
 \begin{array}{ll}
 6\text{ maps with }\rank\Phi=2:&
 (W_1,W_2,W_3)\text{ permutes }(0,A/2,B/2),\\[1mm]
 6\text{ maps with }\rank\Phi=4:&
 (W_1,W_2,W_3)\text{ permutes }(R_0,R_1,R_3).
 \end{array}
 \label{eq:cert-A4-twelve-maps}
\end{equation}
For a permutation \(\sigma=(a,b,c)\) of \((0,1,3)\), the Frobenius-dual
basis is
\[
 Q^0=Q_0/2,\quad Q^1=-Q_1/2,\quad
 Q^2=Q_2/2,\quad Q^3=-Q_3/2,
\]
and the uniquely reconstructed six-port ancestor is
\begin{equation*}
 f_\sigma=\sum_{\mu=0}^3Q^\mu_{12}\otimes F^\sigma_\mu,
 \qquad
 (F_0^\sigma,F_1^\sigma,F_2^\sigma,F_3^\sigma)
 =(C/2,R_a,R_b,R_c).
\end{equation*}
The duals in this formula are coefficient-reconstruction devices, not
available gadgets.

Contract all fifteen port pairs by all twelve actual tetrahedral
kernels.  The exact 1080-card histogram is
\begin{center}
\small
\begin{tabular}{c@{\qquad}rr}
\toprule
\(\sigma\)&zero or matching-deck cards&nonmatching cards\\
\midrule
\((0,1,3),(1,3,0),(3,0,1)\), each&180&0\\
\((0,3,1),(1,0,3),(3,1,0)\), each&36&144\\
\bottomrule
\end{tabular}
\end{center}
All \(180\) cards for each of the three orientations in the first data row
are nonzero.
For the three rejected orientations, the single actual contraction of ports
1 and 3 by \(Q_0=I\) gives the transparent certificate
\begin{center}
\small
\begin{tabular}{c@{\qquad}c@{\qquad}c}
\toprule
\(\sigma\)&support size&three flattening ranks\\
\midrule
\((0,3,1)\)&8&\((2,2,2)\)\\
\((1,0,3)\)&16&\((2,2,2)\)\\
\((3,1,0)\)&16&\((2,2,2)\)\\
\bottomrule
\end{tabular}
\end{center}
Normalize each card by its first nonzero entry, as stipulated for the
projective enumerations above.  In the three \(\sigma\)-rows of the
preceding table, respectively, one nonzero \(2\times2\) minor in each of the flattening orders
\((12\mid34,13\mid24,14\mid23)\) is
\[
 (-1,-1,-2),\qquad(-2,2i,2i),\qquad(-2i,2,-2i).
\]
Thus none of the three outputs is a hidden two-binary matching product.

\begin{proposition}[Tetrahedral permutation classes and eight-vertex signatures]
\label{prop:p1-cert-A4-eight-vertex-orientation-terminal}
Let the exceptional orientation set be
\[
  \Sigma_-:=\{(0,3,1),(1,0,3),(3,1,0)\}.
\]
For each \(\sigma\in\Sigma_-\), let \(q_\sigma\) be the quaternary
obtained by bilinearly contracting ports \(1,3\) of \(f_\sigma\) with the
actual kernel \(Q_0=I\).  Then \(q_\sigma\)
admits actual tetrahedral port dressings to an endpoint-nondegenerate eight-vertex signature.  The three
flattening ranks remain \((2,2,2)\).
\end{proposition}

\begin{proof}
Define the tetrahedral representative matrix
\(G_{\epsilon_1\epsilon_2\epsilon_3}
 :=Q_0+\epsilon_1Q_1+\epsilon_2Q_2+\epsilon_3Q_3\), and use the
even-word order in \cref{eq:octahedral-even-order} on the residual ports
\((2,4,5,6)\).  Direct expansion gives
\begin{center}
\small
\setlength{\tabcolsep}{5pt}
\begin{tabular}{ccc}
\toprule
\(\sigma\)&\(\mathbf L\)&\(\mathbf e(\mathscr D_{\mathbf L}q_\sigma)\)\\
\midrule
\((0,3,1)\)&\((Q_0,G_{++-},G_{+++},Q_1)\)&
\((1,0,0,0,0,0,0,-i)\)\\
\((1,0,3)\)&\((Q_0,Q_0,G_{+++},Q_0)\)&
\((1,1,-i,i,1,-1,i,i)\)\\
\((3,1,0)\)&\((Q_0,G_{+++},Q_0,G_{++-})\)&
\((1,-1,1,1,-i,-i,i,-i)\)\\
\bottomrule
\end{tabular}
\end{center}
Every omitted odd entry is zero and both endpoints are nonzero.  Each
displayed dressing is actual and invertible, so it preserves the certified
flattening ranks.
\end{proof}

For the three surviving orientations, the following coefficient formulas
prove membership in the same affine class.  Let
\(f_{013},f_{130},f_{301}\) be the three surviving six-port signatures and
\(q_{013},q_{130},q_{301}\) their quadratic phase polynomials, and put
\[
 x=(x_1,\ldots,x_6)\in\Ftwo^6.
\]
In the
display below, the indicator constraints and their additions are over
\(\Ftwo\), the exponents \(q_\sigma\) are read modulo four, and the displayed
prefactors are complex scalars:
\begin{align*}
f_{013}(x)&=-\frac{i}{4}\,
 \one[x_3+x_4+x_5+x_6=0]\,i^{q_{013}(x)},\\
q_{013}&=2x_3+3x_4+x_5
 +2(x_1x_2+x_1x_3+x_1x_5+x_2x_3+x_2x_4+x_4x_5),\\[1mm]
f_{130}(x)&=-\frac{i}{4}\,
 \one[x_3+x_4+x_5+x_6=0]\,i^{q_{130}(x)},\\
q_{130}&=3x_1+x_2+3x_4+x_5
 +2(x_1x_2+x_1x_3+x_1x_4+x_2x_3+x_2x_5+x_3x_4+x_3x_5),\\[1mm]
f_{301}(x)&=\frac{1-i}{4}\,
 \one[\substack{x_3+x_4+x_5+x_6=0,\\
                     x_1+x_2+x_4+x_5=1}]\,i^{q_{301}(x)},\\
q_{301}&=2+x_1+2x_2+x_5+x_6
 +2(x_1x_2+x_1x_3+x_1x_4+x_3x_4).
\end{align*}
The displayed supports are affine and every mixed phase coefficient is
even, so the three ancestors are affine.  Separately, direct
expansion of the twelve actual binary representatives \(KX\),
\([K]\in\mathcal T\), verifies that they are affine in this same normalized
basis.

\subsubsection{The icosahedral forms
  \texorpdfstring{\(\mathcal I\) and \(\mathcal I_-\)}{I and I-}}
\label{proofsubsec:cert-A5}

For the icosahedral certificate, use the standard set \(\mathcal I\) in
\cref{eq:quaternion-basis,eq:icosahedral-domain}, ordered as the four axes, the
eight tetrahedral sign points, and the 48 distinct even-permutation points.
Exact arithmetic over
\(\mathbb Q(\phi,i)\) gives
\begin{equation}
 |\mathcal I|=60,
 \qquad
 \#\{g:\operatorname{ord}(g)=1,2,3,5\}=(1,15,20,24).
 \label{eq:cert-A5-orders}
\end{equation}
The set is closed under multiplication, inverse, and transpose, and hence
is the projectivization of the binary icosahedral group, hence the
order-\(60\) projective group \(A_5\).
The apparent opposite chirality gives no additional
distinguished-involution case.  If \(\mathcal I_-\) is obtained by using
odd coordinate permutations in the additional part of
\eqref{eq:icosahedral-domain}, then, for \(D=Q_0+Q_3\),
\begin{equation}
 D^{\mathsf T}XD=2X,
 \qquad
 \mathcal I_-=D\mathcal I D^{-1},
 \qquad
 (DQ_1D^{-1},DQ_2D^{-1},DQ_3D^{-1})=(-Q_2,Q_1,Q_3).
 \label{eq:cert-A5-chirality}
\end{equation}
\Cref{lem:icosahedral-fixed-I-dictionary} uses this identity to translate
kernel and card formulas while preserving the retained fixed-\(I\) problem.

The maximal-cyclic-coset classification in
\cref{prop:cert-platonic-incidence-atlas} gives 200 three-point
\(C_3\)-coset lines and 72 five-point \(C_5\)-coset lines.  Every point is
on ten three-lines and six five-lines.  The quaternary matching-product deck
has \(3\cdot60^2=10\,800\) distinct projective points, and its complete
maximal-rich-line histogram is
\begin{center}
\small
\begin{tabular}{ccr}
\toprule
line size&matching-label type&number\\
\midrule
3&\((1,1,1)\)&216\,000\\
3&\((3)\)&72\,000\\
5&\((5)\)&25\,920\\
\bottomrule
\end{tabular}
\end{center}
There is no \((2,1)\) line and no mixed five-line.

Normalize a mixed three-point incidence by
\eqref{eq:cert-A4-normal-incidence}, and scale its card map so that
\(\Phi(G_+)=A\) and \(\Phi(G_-)=-B\).  Define the image tensors
\(W_j:=\Phi(Q_j)\).  Then
\begin{equation}
 \Phi(Q_0)=C/2,
 \qquad
 W_1+W_2+W_3=(A+B)/2.
 \label{eq:cert-A5-sum}
\end{equation}
The other six tetrahedral sign points require each of
\(W_j,A-2W_j,2W_j-B\) to lie in the zero/matching-product cone
\[
 \mathscr C_{\mathcal I}:=
 \{0\}\cup
 \{\lambda P_m(U,V):\lambda\in\mathbb C^\times,\
      m\in\{0,1,2\},\ [U],[V]\in\mathcal I\}.
\]
Define the following \(15\)-point slice in coefficient projective space:
\begin{align*}
 \mathcal H&=
 \{[a:b:d]\in\mathbb P^2:
      [aQ_0+bQ_1+dQ_3]\in\mathcal I\},
 &|\mathcal H|&=15,\notag\\
 J_u&=aQ_0+bQ_1+dQ_3,
 &\bar J_u&=aQ_0-bQ_1-dQ_3,\notag\\
  R_u&=\frac{P_2(J_u,\bar J_u)}{2(a^2+b^2+d^2)}.
\end{align*}
Here \(u=(a,b,d)\) is any nonzero representative of
\([a:b:d]\in\mathcal H\).  The displayed quotient is unchanged when \(u\)
is rescaled, so \(R_u\) is well defined by the projective coefficient
point.  Moreover,
\(a^2+b^2+d^2=\det J_u\ne0\): the equality follows directly from the
quaternion matrices in \cref{eq:quaternion-basis}, and nonvanishing follows
from \([J_u]\in\mathcal I\subset\PGL_2(\mathbb C)\).  Thus the displayed
denominator never vanishes.
The exact bridge intersection is the eighteen-point set
\begin{equation}
 \mathcal B:=\{W:W,A-2W,2W-B\in\mathscr C_{\mathcal I}\}
 =\{0,A/2,B/2\}\cup\{R_u:[u]\in\mathcal H\}.
 \label{eq:cert-A5-bridge}
\end{equation}
\Cref{lem:platonic-bridge-formula} proves this equality.  The symmetric
icosahedral directions are precisely the three axes with no \(Q_2\)
component and the twelve additional points in that coordinate plane.
For such a direction,
\(J_u^{-1}=\overline J_u/(a^2+b^2+d^2)\), which gives the displayed
formula for \(R_u\).  Thus the list includes every complex scale and
every cancellation case.
As a finite completeness check, the lines through \(A\) determined by the
other 10,799 deck points have intersection-size buckets
\begin{equation}
 1^{10\,591},\qquad2^{80},\qquad4^{12}.
 \label{eq:cert-A5-line-buckets}
\end{equation}

The sum equation on \(\mathcal B^3\) has exactly \(36\) ordered solutions.
Six permute \((0,A/2,B/2)\); the other thirty are the six orderings of
the following five unordered frames, where a vector \(u\) represents
the tensor \(R_u\):
\begin{align}
\mathcal F_0={}&\{(0,0,1),(0,1,0),(1,0,0)\},\notag\\
\mathcal F_1={}&\{(1,-\tau^2,\tau),(1,-\phi,-\phi^2),
                         (1,\phi,-\tau)\},\notag\\
\mathcal F_2={}&\{(1,-\tau^2,-\tau),(1,-\phi,\phi^2),
                         (1,\phi,\tau)\},\notag\\
\mathcal F_3={}&\{(1,-\phi,\tau),(1,\phi,\phi^2),
                         (1,\tau^2,-\tau)\},\notag\\
\mathcal F_4={}&\{(1,-\phi,-\tau),(1,\phi,-\phi^2),
                         (1,\tau^2,\tau)\}.
\label{eq:cert-A5-frames}
\end{align}
\Cref{lem:platonic-bridge-formula} reduces the nondegenerate sum
solutions to orthogonal triples in \(\mathcal H\).
The five displayed frames contain the fifteen points of \(\mathcal H\)
once each.  Using \(\phi^2=\phi+1\), their dot products show that each
point is orthogonal to precisely the other two points in its displayed
frame.  These are therefore all orthogonal triples.  This proves the
\(6+5\cdot6=36\) solutions without a search through \(\mathcal B^3\);
the eighteen-point bridge is proved by the same lemma.

One legal group element suffices for all thirty orderings:
\begin{equation*}
 K_*=\phi Q_1-\tau Q_2+Q_3,
 \qquad [K_*]\in\mathcal I.
\end{equation*}
For an ordering \((u_1,u_2,u_3)\) of any frame, its direct card is
\begin{equation}
 q_{u_1,u_2,u_3}
 =\phi R_{u_1}-\tau R_{u_2}+R_{u_3}.
 \label{eq:cert-A5-output}
\end{equation}
Exact row reduction of the thirty explicit tensors gives, uniformly,
\begin{equation}
 (\rank_{12\mid34},\rank_{13\mid24},\rank_{14\mid23})=(3,3,3),
 \label{eq:cert-A5-ranks}
\end{equation}
with support-size histogram \(8^6,12^4,16^{20}\).  Each output is
therefore a direct nonmatching quaternary card.

For the no-mixed-line branch, first discard rank at most one, which is
already rank lowering.  On the remaining branch
\(\dim\ker\Phi\le2\).  Along every rich line avoiding
\(\mathbb P(\ker\Phi)\), the absence of a mixed line forces all images to
use one matching.  The kernel-deleted connectedness assertion
\cref{eq:cert-platonic-kernel-deleted-graph} propagates this matching to all
nonzero card images among the sixty domain points, placing them in the
fixed-matching Segre variety: the projective simple-tensor locus
\(\Sigma_M\) of \cref{lem:platonic-global-ruling}, with the two matching
edges as its factors.
Because the icosahedral
domain contains the test set in \cref{eq:platonic-quadratic-test-set},
\cref{lem:platonic-global-ruling} then gives, after a port permutation,
\begin{equation}
 \Phi(K)=U_*\otimes\psi(K)
 \quad\text{for all }K,
 \qquad\text{or}\qquad
 \Phi(K)=\psi(K)\otimes U_*.
 \label{eq:cert-A5-common-factor}
\end{equation}
Here \(U_*\) is a fixed nonzero binary factor and
\(\psi:\Mat_2(\mathbb C)\to\Mat_2(\mathbb C)\) is a binary-valued linear
map.
Combining the rank-two solutions, \eqref{eq:cert-A5-ranks}, and
\eqref{eq:cert-A5-common-factor} gives the bounded-interface trichotomy:
rank lowering, an actual nonmatching quaternary, or a genuine common
factor.

The group parametrization and \cref{lem:platonic-bridge-formula} prove
completeness of the five frames.  The rejecting kernel and rank assertions
are evaluated on the thirty displayed outputs over \(\mathbb Q(\phi,i)\).
The common-factor alternative follows from
\cref{lem:platonic-global-ruling}.

\begin{lemma}[Actual rainbow-line orbit certificate]
\label{lem:platonic-rainbow-orbit}
For a marked physical group \(G\), let \(\Gamma_G\) be the finite
permutation group of the projective deck induced by the actual local
operations
\[
 \mathscr L_{j,g}:[F]\longmapsto
 [(I^{\otimes(j-1)}\otimes K_g\otimes
 I^{\otimes(4-j)})F]
 \quad(g\in G,\ 1\le j\le4),
 \qquad (12),(23),(34),
\]
where the first operations mean the physical attachment of an actual binary
\(B_g\doteq X K_g^{\mathsf T}\) in the orientation with effective leg map
\(B_g^{\mathsf T}X\doteq K_g\), which induces the displayed projective
action on that leg and records its nonzero representative scalar; the last three are port
permutations.  Let the standard representative rainbow line be
\[
 \ell_0^{\rm std}
 =\{[P_0(Q_0,Q_0)],[P_1(Q_0,Q_0)],[P_2(Q_2,Q_2)]\}.
\]
Write \(\mathscr S_G\) for the displayed finite generating set (all
\(\mathscr L_{j,g}\) together with the three adjacent swaps), so that
\(\Gamma_G=\langle\mathscr S_G\rangle\).  Here \(K_g\) is the fixed actual
representative of \(g\), not a coefficient-space point.  For the second
octahedral form, write \(\theta:=\theta_{\mathcal O}\) and use
\(\mathcal L_{\mathcal O}\) only to name the initial projective line
\[
\ell_0^{\mathcal O'}
 =\{[P'_0(\theta(Q_0),\theta(Q_0))],
     [P'_1(\theta(Q_1),\theta(Q_1))],
     [P'_2(\theta(Q_2),\theta(Q_2))]\},
\]
which is \(\mathcal L_{\mathcal O}(\ell_0^{\rm std})\) projectively by the
explicit matching formulas below.  The operator
\(\mathcal L_{\mathcal O}\) is not a generator.  Define
\(\ell_0^G:=\ell_0^{\rm std}\) for
\(G\in\{\mathcal T,\mathcal O,\mathcal I\}\).  The orbit of \(\ell_0^G\) under these actual operations has the following
size:
\begin{center}
\small
\begin{tabular}{lrrr}
\toprule
marked form & deck points & all rainbow three-lines & orbit size\\
\midrule
\(\mathcal T\) & 432 & 1\,728 & 1\,728\\
\(\mathcal O\) & 1\,728 & 13\,824 & 13\,824\\
\(\mathcal O'\) & 1\,728 & 13\,824 & 13\,824\\
\(\mathcal I\) & 10\,800 & 216\,000 & 216\,000\\
\(\mathcal T_{\rm ext}\) & 432 & 0 & 0\\
\bottomrule
\end{tabular}
\end{center}
For \(\mathcal O'\), the representative in the orbit column is
\(\ell_0^{\mathcal O'}\); for the other rows with a nonzero rainbow count it
is \(\ell_0^{\rm std}\).  Consequently every rainbow line in the first four rows is
sent to that row's representative by a sequence of
actual local dressings and port permutations.
The corresponding scalar and ordered port permutation are part of the
finite witness; no coefficient similarity is used in this assertion.
For \(\alpha\in\{dt,td\}\), define the initial mixed line as the
\(\mathcal L_\alpha\)-image of \(\ell_0^{\rm std}\).  The three identities in
\cref{eq:octahedral-mixed-coefficient-transport} intertwine every standard
orbit generator with the corresponding actual
\(\mathcal O_\alpha\)-representative and preserve its nonzero scalar and port
order.  Hence each mixed form has the same \(1\,728\) deck points and
\(13\,824\) rainbow three-lines, in one actual local-dressing orbit.  These
two rows are omitted from the table only because their three numerical
entries coincide with the two displayed octahedral rows.
\end{lemma}

\begin{proof}
First let \(G\) contain \([X]\) and be reversal-stable, so it is closed
under inverse transpose.  Every binary factor of a \(G\)-matching tensor
has its matrix class in \(G\), because its transfer does and \([X]\in G\).
Consider any rainbow line.  Actual local dressings send its matching-\(0\)
point to \(A=P_0(I,I)\); write its matching-\(1\) point as \(P_1(U,V)\).
The further dressing \((U^{-1},U^{\mathsf T},I,I)\) fixes \(A\) and sends
that point to \(P_1(I,W)\), where \(W=U^{\mathsf T}V\).

The third point has rank one across \(14\mid23\).  Since the three
points are distinct, both coefficients of its linear combination of
\(A\) and \(P_1(I,W)\) are nonzero.
\Cref{eq:platonic-swap-rank-one} therefore forces \(W\) to be scalar,
and the line is \(\{[A],[B],[C]\}\), with \(C=A-B=P_2(Y,Y)\).
Thus every rainbow line has an actual local normalization to this line.
It exists exactly when \([YX]=[Z]\in G\).  In particular, the external
form has none by \cref{lem:marked-a4-forms}.

For \(G\in\{\mathcal T,\mathcal O,\mathcal I\}\), the local action of
\(G^4\) preserves the three matching labels.  A stabilizer of the
normalized line therefore fixes its three points, and the first two give
\[
 T_2=T_1^{-\mathsf T},\qquad T_3=T_1^{-\mathsf T},\qquad T_4=T_1
\]
projectively.  Conversely, \((T,T^{-\mathsf T},T^{-\mathsf T},T)\)
fixes all three, using \(TYT^{\mathsf T}=\det(T)Y\).
The stabilizer has \(|G|\) elements, so the complete rainbow orbit has
\(|G|^3=1728,13824,216000\) elements, respectively.  This establishes both
transitivity and the counts directly.
For \(\mathcal O'\), recall that
\(\mathbf H_{\mathcal O}\) is the transformation tuple and
\(\mathcal L_{\mathcal O}=\mathscr D_{\mathbf H_{\mathcal O}}\) is its
induced coefficient operator.  The identity
\(H_{\mathcal O}H_{\mathcal O}^{\mathsf T}\doteq X\) gives a
matching-dependent but
explicit bijection
\[
 \mathcal L_{\mathcal O} P_M(U,V)
 \doteq P'_M(U'_M,V'_M),
\]
where \(P'_M\) denotes the same ordered matching placement in the
\(\mathcal O'\) coordinates.  More explicitly, writing
\(\theta=\theta_{\mathcal O}\) and \(X=Q_1\) projectively, the three cases are
\[
\begin{aligned}
\mathcal L_{\mathcal O}P_0(U,V)&\doteq P'_0(\theta(U),\theta(V)),\\
\mathcal L_{\mathcal O}P_1(U,V)&\doteq P'_1(\theta(UX),\theta(XV)),\\
\mathcal L_{\mathcal O}P_2(U,V)&\doteq P'_2(\theta(U),\theta(XVX)).
\end{aligned}
\]
These formulas are the ordered
\(H_{\mathcal O}^{-1}/H_{\mathcal O}^{\mathsf T}\) pair contractions;
the displayed \(X\)'s record the fixed orientation factors.  Thus this
analytic map preserves matching labels and rich-line incidence.  It also
conjugates the local group \(\mathcal O^4\) to \((\mathcal O')^4\):
on an \(H^{-1}\) leg this is \(U\mapsto H^{-1}UH\), while on an
\(H^{\mathsf T}\) leg it is
\[
 H^{\mathsf T}UH^{-\mathsf T}
 =H^{-1}(SUS^{-1})H,\qquad S=HH^{\mathsf T},\quad [S]\in\mathcal O.
\]
The same argument applies to \(H_{dt},H_{td}\), with their already
verified \(S\)-matrices.  Hence all three nonstandard octahedral forms
have the same transitive local orbit, starting at their displayed
\(\mathcal L\)-image of the standard line.  The operator \(\mathcal L\)
only proves this correspondence; each conjugated leg matrix belongs to the
actual current group.  For \(\mathcal T_{\rm ext}\), the preceding rank-one argument and
\([Z]\notin\mathcal T_{\rm ext}\) exclude every rainbow line.

To obtain an actual normalization witness, choose coset representatives
as follows.  In the
standard forms, right multiplication by a stabilizing tuple gives each
local coset a unique representative \((I,T_2,T_3,T_4)\).  Enumerate these
\(|G|^3\) tuples and their images; each rainbow line occurs once.  The
nonstandard correspondence gives the same enumeration in its actual
current-group representatives.  Applying the inverse recorded tuple sends
any selected line to the prescribed one.  The port order and all nonzero
representative scalars are obtained by these exact matrix products.
\end{proof}

\subsubsection{The complete Platonic incidence atlas}
\label{proofsubsec:cert-platonic-incidence-atlas}

\begin{proposition}[Exact Platonic incidence atlas]
\label{prop:cert-platonic-incidence-atlas}
For the standard tetrahedral, first octahedral, external tetrahedral, and
icosahedral marked configurations, the domain and quaternary matching decks
have the following complete maximal-rich-line data.
\begin{center}
\small
\setlength{\tabcolsep}{4pt}
\begin{tabular}{cL{0.25\linewidth}cL{0.34\linewidth}}
\toprule
configuration&domain points and lines&deck points&deck lines by matching-label type\\
\midrule
\(\mathcal T\)&\(12;\ 16_{[3]}\)&\(432\)&
  \(1728(1,1,1),\ 1152(3)\)\\
\(\mathcal O\)&\(24;\ 32_{[3]},18_{[4]}\)&\(1728\)&
  \(13824(1,1,1),\ 4608(3),\ 2592(4)\)\\
\(\mathcal T_{\rm ext}\)&\(12;\ 16_{[3]}\)&\(432\)&
  \(1152(3)\)\\
\(\mathcal I\)&\(60;\ 200_{[3]},72_{[5]}\)&\(10800\)&
  \(216000(1,1,1),\ 72000(3),\ 25920(5)\)\\
\bottomrule
\end{tabular}
\end{center}
In the second column the entry before the semicolon is the number of domain
points, and \(n_{[s]}\) denotes \(n\) rich lines of size \(s\).  A
parenthesized partition in the last column is the matching-label type.  Thus the
standard tetrahedral deck has only rainbow or monochromatic three-lines;
the external tetrahedral deck has only monochromatic three-lines; and the
octahedral and icosahedral decks have no type \((2,1)\) line and no mixed
four- or five-line, respectively.
The octahedral row also gives the abstract incidence of
\(\mathcal O'=H_{\mathcal O}^{-1}\mathcal O H_{\mathcal O}\) under the
coefficient transport in \cref{eq:cert-S4-second-form}.  Since
\(H_{\mathcal O}\notin GO(X)\), this transport is
not a gadget assertion; every bridge kernel and port dressing used in the
second marked form is checked separately in actual
\(\mathcal O'\)-representatives.
Explicitly, the coefficient operator
\(\mathcal L_{\mathcal O}=\mathscr D_{\mathbf H_{\mathcal O}}\) satisfies
\(\mathcal L_{\mathcal O}P_M(U,V)\doteq P'_M(U'_M,V'_M)\).  Writing
\(\theta=\theta_{\mathcal O}\) and \(X=Q_1\) projectively, the three explicit
instances are
\[
\begin{aligned}
\mathcal L_{\mathcal O}P_0(U,V)&\doteq P'_0(\theta(U),\theta(V)),\\
\mathcal L_{\mathcal O}P_1(U,V)&\doteq P'_1(\theta(UX),\theta(XV)),\\
\mathcal L_{\mathcal O}P_2(U,V)&\doteq P'_2(\theta(U),\theta(XVX)).
\end{aligned}
\]
The \(X\)'s are the fixed orientation factors from the ordered
\(H_{\mathcal O}^{-1}/H_{\mathcal O}^{\mathsf T}\) pair contractions, and
each displayed matrix lies in the current \(\mathcal O'\)-representatives.
Thus \(\mathcal L_{\mathcal O}\)
preserves matching labels and rich-line
incidence as an analytic bijection only.  The actual \(\mathcal O'\) port
witnesses are supplied by \cref{lem:platonic-rainbow-orbit} and the separate
bridge tables.

For \(\mathcal O_{dt}\) and \(\mathcal O_{td}\), the direct transport
\cref{eq:octahedral-mixed-coefficient-transport} likewise preserves every
matching label, rich-line incidence, and kernel-deleted adjacency, while
intertwining the generators with actual representatives of the current mixed
form.  Thus both mixed rows have the octahedral counts and connectivity in
the table.  As in the preceding paragraph, the coefficient operator is only
the proof dictionary; the physical witnesses are the transported actual
kernels and dressings.

For \(\mathcal T\), each domain point lies on four three-lines.  For
\(\mathcal O\), each lies on seven rich lines.  For \(\mathcal I\), each
lies on ten three-lines and six five-lines.  In each standard marked form,
every displayed domain-line family is an actual left/right group orbit.
The intersection graphs of rich domain lines---vertices are rich lines and adjacency means nonempty intersection---are connected; their edge
counts for \(\mathcal T,\mathcal O,\mathcal I\) are respectively
\begin{equation}
  72,\qquad 504,\qquad 7200.
  \label{eq:cert-platonic-cross-line-edges}
\end{equation}
For the localization interface, define the actual projective domain set
\[
 \mathcal D_G=\{[K_g]:g\in G\}
\]
and join two distinct domain points when they lie on a common rich line.  For every complex subspace \(W\le\Mat_2\) with \(\dim W\le2\), the
induced graph on
\begin{equation}
 \mathcal D_G\setminus\mathbb P(W)
 \label{eq:cert-platonic-kernel-deleted-graph}
\end{equation}
is nonempty and connected.  The distinct deletion sets tested for
\(\mathcal T,\mathcal O,\mathcal I\) number, respectively,
\begin{equation}
 \begin{aligned}
  47&=1+12+(18+16),\\
 147&=1+24+(72+32+18),\\
 783&=1+60+(450+200+72).
 \end{aligned}
 \label{eq:cert-platonic-kernel-deletion-counts}
\end{equation}
Here the parenthesized summands count the distinct two-point spans by their
numbers of domain points: \(18\) of size two and \(16\) of size three for
\(\mathcal T\); \(72,32,18\) of sizes two, three, four for
\(\mathcal O\); and \(450,200,72\) of sizes two, three, five for
\(\mathcal I\).  The same assertion for \(\mathcal O'\) and
\(\mathcal T_{\rm ext}\) is transported by their respective invertible
coefficient similarities.  Each \(\mathcal O'\) rainbow line
used below nevertheless has a separately checked normalization witness
from actual \((\mathcal O')^4\rtimes S_4\) port operations; the similarity
is never installed as a gadget.

Finally, the actual kernel representatives span \(\Mat_2\).  One may take
\begin{equation}
 \begin{aligned}
  \mathcal T,\mathcal O,\mathcal I &: Q_0,Q_1,Q_2,Q_3,\\
  \mathcal O' &: \theta_{\mathcal O}(Q_0),\theta_{\mathcal O}(Q_1),
                  \theta_{\mathcal O}(Q_2),\theta_{\mathcal O}(Q_3),\\
  \mathcal T_{\rm ext} &: Q_0,Q_1,Q_2+Q_3,Q_2-Q_3.
 \end{aligned}
 \label{eq:cert-platonic-spanning-kernels}
\end{equation}
The tetrahedral incidence counts, and the rainbow deck-line normalizations
used in the octahedral and icosahedral bridges, are realized by actual
independent port dressings and port permutations.  For
any finite configuration in \(\mathbb P(\mathscr V)\), and distinct points
\([F_0],[F]\), define the \emph{quotient direction of \([F]\) at
\([F_0]\)} to be
\[
  [F+\Span_{\mathbb C}\{F_0\}]
  \in\mathbb P\bigl(\mathscr V/\Span_{\mathbb C}\{F_0\}\bigr),
\]
using arbitrary nonzero representatives.  This class is independent of
those representatives, and two points have the same quotient direction
exactly when they lie with \([F_0]\) on one projective line.  For
\(\mathcal T_{\rm ext}\), actual external-group port actions are transitive
on the \(432\) deck points and permute the three matching labels; at
\(P_0(Q_0,Q_0)\) the quotient-direction histogram is
\begin{equation}
  1^{415},\qquad 2^8,
  \label{eq:cert-external-A4-line-buckets}
\end{equation}
and all eight rich lines there are monochromatic.  No use of the analytic
twist \(R\) is made in this deck transitivity assertion.

For completeness, at the normalized octahedral and icosahedral deck points
the exact quotient-direction histograms are, respectively,
\begin{equation}
  1^{1645},\ 2^{32},\ 3^6,
  \qquad
  1^{10591},\ 2^{80},\ 4^{12}.
  \label{eq:cert-platonic-line-buckets-summary}
\end{equation}
The first gives all octahedral rich-line types through the base point; the
second gives all icosahedral rich-line types through it.  Hence the global
counts in the table follow by actual transitivity and incidence double
counting, not by assuming an abstract group action.
\end{proposition}

\begin{proof}
We first classify every domain secant.  Let \(\ell\) contain distinct
\([g],[h]\in G\).  Left multiplication by \(g^{-1}\) sends it to
\(\mathbb P(\operatorname{span}\{I,T\})\), where \([T]=[g^{-1}h]\ne1\).
A nonidentity finite-order projectivity has two distinct eigenlines, so
\(T\) is diagonal in their basis.  Every matrix in this two-dimensional
span is diagonal there.  The nonsingular group points on the line form
\(C=G\cap\mathbb T\), where \(\mathbb T\) is that diagonal torus.
A finite subgroup of \(\mathbb C^\times\) is cyclic, so \(C\) is cyclic.
It is maximal cyclic in \(G\): any cyclic extension containing its
nonidentity element \([T]\) has the same eigenlines.  Conversely, every
maximal cyclic subgroup lies on its eigenline torus's projective line.
Thus the domain secants are exactly the left cosets of maximal cyclic
subgroups, with line size equal to the subgroup order.

The permutation cycle types of \(A_4,S_4,A_5\) give the following
numbers of maximal cyclic subgroups:
\[
\begin{array}{c|rrrr}
 &|C|=2&|C|=3&|C|=4&|C|=5\\\hline
 A_4&3&4&0&0\\
 S_4&6&4&3&0\\
 A_5&15&10&0&6.
\end{array}
\]
For example, the three double transpositions of \(S_4\) are squares of
its three cyclic order-four subgroups; only the six transpositions
contribute maximal order-two subgroups.  The order-three and order-five
counts divide their numbers of generators by two and four, respectively.
A subgroup of order \(d\) has \(|G|/d\) left cosets, and distinct maximal
cyclic subgroups give distinct coset sets.  Hence the complete secant
counts are
\[
 A_4:(18,16),\qquad S_4:(72,32,18),\qquad A_5:(450,200,72),
\]
ordered by the nonzero subgroup-order columns above.  Each point belongs
to one coset of each subgroup.  This also proves every stated number of
rich lines through a point.  Conjugation in \(G\) is transitive on the maximal cyclic subgroups of each
displayed order, and left multiplication moves
between their cosets, proving the asserted actual left/right line orbits.

The deck-line classification then follows structurally.  A line containing
three matching products either has one matching or has all three:
repetition of one matching together with another would give flattening
ranks \(1,1,4\), contradicting collinearity.  In the latter case,
\cref{lem:platonic-rainbow-localization,lem:platonic-rainbow-orbit} gives
exactly three points and the stated \(|G|^3\) count in each standard
form.  For a line with one matching, the simple-tensor sum criterion in
\cref{lem:platonic-global-ruling} implies that exactly one binary factor
varies.  For each domain line there are three matchings, two choices of
varying factor, and \(|G|\) fixed factors.  Thus each monochromatic
line count is \(6|G|\) times its domain-line count.  This proves the
complete deck table from the domain table.

We next prove kernel-deleted connectivity.  The domain-point graph has
\[
 (n,\delta,s)=(12,8,3),\ (24,17,4),\ (60,44,5),
\]
where \(n\) is its number of vertices, \(\delta\) its degree, and \(s\)
the maximum line size.  The degree is the sum of \(t-1\) over the rich
lines of size \(t\) through a point: respectively \(4\cdot2\),
\(4\cdot2+3\cdot3\), and \(10\cdot2+6\cdot4\).
A subspace of dimension at most two removes \(k\le s\) points.  The
remaining graph has \(n-k>0\) vertices and minimum degree at least
\(\delta-k\).  If it were disconnected, two components would each have
at least \(\delta-k+1\) vertices.  But every displayed triple satisfies
\[
 2(\delta-k+1)>n-k\qquad(0\le k\le s),
\]
a contradiction.  This proves \cref{eq:cert-platonic-kernel-deleted-graph}
without graph traversal.  The deletion counts follow from the domain
secant counts by adding the empty and singleton deletions, as displayed
in \cref{eq:cert-platonic-kernel-deletion-counts}.  Every complex
subspace is covered: an intersection with at least two domain points
spans its two-dimensional vector space; smaller intersections are empty
or singletons.  The rich-line intersection graphs are also connected:
a path in the connected domain-point graph gives a sequence of rich
lines with consecutive nonempty intersections.  Their edge counts
\(72,504,7200\) follow by summing \(\binom{d_v}{2}\) over domain
points, where \(d_v=4,7,16\), respectively.  Distinct projective lines
meet in at most one point, so this counts every edge once.

The coefficient transports preserve these domain and deck incidences.
\Cref{lem:platonic-rainbow-orbit} supplies actual local normalization
matrices.  The spanning assertions in
\cref{eq:cert-platonic-spanning-kernels} follow from the displayed four
independent matrices in each row.

For the external form, the coefficient similarity
\(\theta_{\rm ext}\) transports the twelve domain points and their
sixteen three-lines.  Its deck has no rainbow line by
\cref{lem:platonic-rainbow-orbit}.  The monochromatic count is therefore
\(6\cdot12\cdot16=1152\).  Actual independent port actions are
transitive on each matching class, and port permutations exchange the
three classes, so the full action is transitive on the \(432\) deck
points.  At each point, the two factor choices and four domain lines
through each factor give eight rich three-lines.  These account for
\(16\) of the other \(431\) points, proving the quotient-direction
histogram \(1^{415},2^8\).

The octahedral and icosahedral quotient-direction histograms follow in
the same way.  Rainbow transitivity gives \(|G|\) rainbow lines through
each deck point, while domain incidences give the monochromatic lines.
The remaining points each determine a separate two-point line.
The programs
\path{certificates/verify_a4_s4_certificates.py},
\path{certificates/verify_a4_external_form.py},
\path{certificates/verify_a5_certificate.py}, and
\path{certificates/verify_platonic_extended.py} also reproduce the
complete finite data.  The coset, rank-one, and degree arguments above prove this proposition
independently of those computations.
\end{proof}

\begin{theorem}[Incidence of Platonic kernel lines]
\label{thm:cert-interface-platonic-incidence}
For the marked Platonic domain and quaternary matching decks, the following
statements hold.
\begin{enumerate}[label=\textnormal{(\roman*)}]
\item The four entries in each row below count, in order, rainbow
      three-lines and monochromatic lines of sizes three, four, and five:
      \[
      \begin{array}{c|c}
      \mathcal T&(1728,1152,0,0)\\
      \mathcal O,\mathcal O',\mathcal O_{dt},\mathcal O_{td}
        &(13824,4608,2592,0)\\
      \mathcal T_{\rm ext}&(0,1152,0,0)\\
      \mathcal I,\mathcal I_-&(216000,72000,0,25920).
      \end{array}
      \]
      In particular, no deck has a line of type \((2,1)\), no
      octahedral deck has a mixed four-line, neither icosahedral deck has a
      mixed five-line, and the external tetrahedral deck has no rainbow
      line.
\item Every rainbow line for
      \(G\in\{\mathcal T,\mathcal O,\mathcal O',\mathcal O_{dt},
      \mathcal O_{td},\mathcal I,\mathcal I_-\}\)
      lies in one orbit under actual local \(G\)-dressings and port
      permutations.  A witness includes the ordered port permutation and
      its nonzero realization scalar; no coefficient similarity is used as
      a gadget.
\item Let \(\mathcal D_G\) be the actual projective domain set for any of the
      marked rows above.  For every complex subspace
      \(W\le\Mat_2(\mathbb C)\) with \(\dim W\le2\), the rich-line domain
      graph induced on \(\mathcal D_G\setminus\mathbb P(W)\) is nonempty and
      connected.
\end{enumerate}
\end{theorem}

\begin{proof}
The complete line counts and kernel-deleted connectivity are
\cref{prop:cert-platonic-incidence-atlas}; actual rainbow transitivity,
including the separate \(\mathcal O'\) witnesses, is
\cref{lem:platonic-rainbow-orbit}.  For the mixed rows,
\cref{eq:octahedral-mixed-coefficient-transport} gives the same counts and
intertwines the standard orbit generator-by-generator with actual
current-form matrices.  The fixed-equality dictionary
\cref{lem:icosahedral-fixed-I-dictionary} gives both icosahedral statements
in \(\mathcal I_-\), with actual \(\kappa_D\)-kernels and nonzero scalars,
without changing the retained literal \(I\).
\end{proof}

\begin{remark}[Computational completeness and physical interpretation]
\label{rem:cert-platonic-scope}
The finite preserver and hereditary eight-port completeness claims
retain an exact-computation component over the stated number fields.
The bridge classifications have the structural proof in
\cref{lem:platonic-bridge-formula}.
The incidence classification and deletion connectivity have the structural
proof in \cref{prop:cert-platonic-incidence-atlas}.
The companion programs regenerate the configurations and test the stated
predicates; they emit assertions rather than independent proof objects.
The main lemmas supply the physical interpretation: actual kernels and
dressings, factor exposure, and reductions of the retained Holant problem.
The mixed-octahedral exits are given by the explicit coefficient and
contraction tables in
\cref{prop:p1-cert-S4-uniform-eight-vertex-separator}.
\end{remark}

\subsubsection{Tetrahedral eight-port tensors with safe pair cards}
\label{proofsubsec:cert-A4-b8-parent-audit}

\begin{proposition}[Tetrahedral eight-port classification]
\label{prop:cert-A4-b8-parent-audit}
Let \(\mathscr M_6,\mathscr H_6,\mathscr S_6\), and \(\mathscr C_6\) be
the six-port sets defined before
\cref{prop:a4rel-required-eight-port}; let the base matching-product signature
be \(M_0\) and the base H-core signature be \(H_0\):
\[
 M_0=(Q_0)_{12}(Q_0)_{34}(Q_0)_{56},
 \qquad H_0=\Hcore.
\]
Then
\begin{equation}
 |\mathscr M_6|=25920,\qquad
 |\mathscr H_6|=31104,\qquad
 |\mathscr S_6|=57024,
 \label{eq:cert-A4-b8-safe-set-sizes}
\end{equation}
and the union defining \(\mathscr S_6\) is disjoint.  Through either
\([M_0]\) or \([H_0]\), the quotient-direction histogram of the other
safe points is
\begin{equation}
  1^{56855},\qquad 2^{84}.
  \label{eq:cert-A4-b8-line-histogram}
\end{equation}
Here two safe points \(F,F'\ne F_0\) have the same quotient direction at
\([F_0]\) when
\([F+\Span_{\mathbb C}\{F_0\}]
 =[F'+\Span_{\mathbb C}\{F_0\}]\) in the projectivized
quotient.  Thus the histogram means that
56,855 directions contain one other safe point and 84 directions contain
two.  An
\(\mathscr X\mathscr Y\)-endpoint line means that its two nonbase safe
points lie in \(\mathscr X\) and \(\mathscr Y\), with repetition allowed.
The \(84\) rich lines through \([M_0]\) split into \(48\)
\(\mathscr M_6\mathscr M_6\)-endpoint lines and \(36\)
\(\mathscr H_6\mathscr H_6\)-endpoint lines, with orbit
sizes \(36,36,12\) under the actual stabilizer subgroup specified below.  Those through \([H_0]\) split into \(60\)
\(\mathscr H_6\mathscr M_6\)-endpoint lines and \(24\)
\(\mathscr H_6\mathscr H_6\)-endpoint lines, with orbit
sizes \(24,12,8,12,12,4,12\) under its specified actual stabilizer subgroup.

For each of the ten resulting line-orbit representatives and either endpoint
orientation, denote the oriented endpoint tensors by \(A,D\) and the base
safe tensor by \(F_0\), scaled so that \(A+D=2F_0\); define the bridge set
\begin{equation}
 \mathcal B(A,D)=
 \{W:W,\ A-2W,\ D+2W\in\mathscr C_6\}.
 \label{eq:cert-A4-b8-bridge}
\end{equation}
Then \(|\mathcal B(A,D)|=6\), and exactly twelve ordered triples in
\(\mathcal B(A,D)^3\) sum to \((A-D)/2\).  The resulting
\(10\cdot2\cdot12=240\) reconstructed eight-port parents are projectively
distinct.  For a reconstructed parent \(g\), let
\[
  \Phi_g:\Mat_2(\mathbb C)\longrightarrow
       (\mathbb C^2)^{\otimes6},\qquad
  \Phi_g(K):=\Phi_{12}^g(K),
\]
be its distinguished-pair card map.  Testing the
\(\binom82\cdot12=336\) actual tetrahedral cards in
a fixed order, stopping only after an outside witness and exhausting all
cards for every parent without one, gives
\begin{equation*}
\begin{array}{clr}
\toprule
\rank\Phi_g&\text{outcome}&\text{number}\\ \midrule
2&\text{an actual card outside }\mathscr C_6&120\\
3&\text{no parent}&0\\
4&\text{an actual card outside }\mathscr C_6&60\\
4&\text{all }336\text{ cards in }\mathscr C_6&60.\\
\bottomrule
\end{array}
\end{equation*}
Every one of the final sixty parents is affine in the normalized
computational basis; twenty-four also have an exact nontrivial rank-one
flattening.  For each of the other \(180\) parents there is a specified physical pair
and actual tetrahedral kernel producing a card outside \(\mathscr C_6\).
The count \(180\) is a count of parents with an outside witness.
\end{proposition}

\begin{proof}
Both safe sets are generated over \(\mathbb Q(i)\): matching products are
canonicalized directly, while the \(\Hcore\)-orbit is generated to closure
under actual local tetrahedral matrices and port permutations.  Exact
projective comparison gives 
\cref{eq:cert-A4-b8-safe-set-sizes,eq:cert-A4-b8-line-histogram}.  The
subgroups used for this audit are generated by the following actual
operations.  At \(M_0\), use the pair stabilizers
\((U,U^{-\mathsf T})\), \(U\in\mathcal T\), and the permutations
preserving its three-pair matching.  At \(H_0\), first enumerate the local
two-port tetrahedral actions that permute the four projective Bell
signatures.  Retain the triples of such pair-block actions that fix
\([H_0]\), and add all residual-port permutations that directly fix
\([H_0]\).  Every selected generator fixes its base tensor projectively.
Close their induced permutations on the complete set of \(84\) rich lines.
These explicitly generated stabilizer subgroups give the three plus
seven orbits and sizes stated above.  In a fixed lexicographic ordering of normalized
coefficient vectors, the first line of each orbit is the deterministic
line-orbit representative used below; no unlisted generator is assumed.
Applying the inverse stabilizer word transports every bridge solution,
reconstructed parent, and recorded outside-card witness from that
representative to every other line in its orbit, with the physical pair,
kernel, port order, and nonzero scalar transported along the same actual
operation.  Thus the ten representative calculations cover every rich line
through either base point, not merely one line from each abstract incidence
class.

Fix representatives over \(\mathbb Q(i)\) on one such line.  Since its
endpoints are independent, \(A+D=2F_0\) uniquely fixes both endpoint
scales.  To solve \cref{eq:cert-A4-b8-bridge}, the case \(W=0\) is tested
separately.  If \(W\) and \(A-2W\) are nonzero and nonproportional, their
directions lie in an exhaustively generated rich-line bucket and the
equation fixes both scales.  In the remaining nonzero case write \(W=tA\)
for a nonzero scalar \(t\);
then \(D+2tA\ne0\), and the third safe direction on
\(\langle A,D\rangle\) fixes \(t\).  Thus the six bridge points and twelve
sum triples are a complete symbolic solve over \(\mathbb Q(i)\), including
all zero and complex-scale cases.

For every ordered bridge triple \((W_1,W_2,W_3)\), put
\(F_j=W_j\) for \(1\le j\le3\), retain the displayed base tensor as
\(F_0\), and use the Frobenius-dual quaternion basis to reconstruct
\[
 g=\sum_{\mu=0}^3Q^\mu_{12}\otimes F_\mu.
\]
Exact projective comparison gives \(240\) distinct parents.  The companion
replay \path{certificates/verify_a4_b8.py}, invoked with \texttt{--fresh} by
the suite runner, then scans physical pairs in lexicographic order and the twelve actual
kernels in the point order of \cref{eq:tetrahedral-domain}.  For each card it
tests zero first and otherwise tests projective membership in the
definition-level regenerated set \(\mathscr S_6\).
For an outside parent it records and rechecks the first witness; for a safe
parent it necessarily checks all \(336\) cards.  Affineness of each survivor
is tested from the definition: affine-coset support, fourth-root ratios,
linear and even mixed quadratic phase coefficients, and the recovered phase
on every support point.  Factor cuts are tested by exact flattening minors.
This exhaustive calculation regenerates the two orbits, their line
buckets, bridges, parents, and outside witnesses over \(\mathbb Q(i)\).
The finite classification of the safe parents remains computational.
\end{proof}

\begin{theorem}[Tetrahedral eight-port contractions]
\label{thm:cert-interface-A4-q8}
Let \(\mathscr M_6\) and \(\mathscr H_6\) be the actual tetrahedral orbits of
the base matching product and \(\Hcore\), and put
\(\mathscr S_6=\mathscr M_6\sqcup\mathscr H_6\) and
\(\mathscr C_6=\{0\}\cup\operatorname{cone}(\mathscr S_6)\).  Then
\[
 |\mathscr M_6|=25920,\qquad |\mathscr H_6|=31104,\qquad
 |\mathscr S_6|=57024.
\]
Every normalized eight-port parent reconstructed from a rich line through
either safe base point is actual-equivalent to one of the \(240\) audited
parents.  No audited parent has distinguished card-map rank three; \(180\)
have a recorded actual tetrahedral card outside \(\mathscr C_6\), and the
remaining \(60\) have all \(336\) actual cards safe and belong to the
affine class \(\cA\).  These assertions include arbitrary nonzero
complex scales and both endpoint orientations.
\end{theorem}

\begin{proof}
This is the finite-level conclusion of
\cref{prop:cert-A4-b8-parent-audit}.  Its actual-stabilizer covariance
transports the ten representative line calculations, their reconstructed
parents, and their recorded card witnesses to every rich line through either
base point.  The actual tetrahedral port dressings and port permutations
preserve \(\cA\), because the corresponding actual tetrahedral binaries are
affine, so the affine classification transports as well.
\end{proof}

\subsubsection{Deficient mixed-card localization}
\label{proofsubsec:cert-platonic-deficient-mixed}

\begin{proposition}[Platonic contractions of rank less than four]
\label{prop:cert-platonic-deficient-mixed}
Let \(\Lambda\) be a signature set, let \(G\) be the standard marked
tetrahedral form, either unmixed octahedral form, or the marked icosahedral
form, and let \(h\) be a localized six-port signature directly realizable from
\(\Lambda\), whose
card map \(\Phi\) has rank two or three.  Suppose every actual
\(G\)-card is zero or a nonsingular \(G\)-matching product and that two nonzero
cards use different matchings.  Then the nonkernel domain points contain a
rich line disjoint from the projectivized kernel of \(\Phi\) whose three
images are nonzero and use the three distinct residual matchings.  For the
standard tetrahedral, first
octahedral, and icosahedral forms, actual contractions produce one of six
support-ten quaternaries \(q_\pi\), indexed by \(\pi\in S_3\), each
satisfying
\begin{equation}
  \bigl(\rank_{12\mid34}q_\pi,
        \rank_{13\mid24}q_\pi,
        \rank_{14\mid23}q_\pi\bigr)=(3,3,3).
  \label{eq:platonic-rank-two-q4-minors}
\end{equation}
For the second octahedral form they instead produce the corresponding
transported quaternaries \(q'_\pi\).  These need not have support size ten,
but invertible local transport preserves the same three ranks.  Writing
\(q^\star_\pi=q_\pi\) in the standard forms and
\(q^\star_\pi=q'_\pi\) in the second octahedral form, one has
\[
  \KHolant(\Lambda,q^\star_\pi)
  \leT\KHolant(\Lambda,h)
  \leT\KHolant(\Lambda).
\]
The statement includes the second octahedral form in its own actual
representatives; similarity by \(H_{\mathcal O}\) is not used as a gadget.
\end{proposition}

\begin{proof}
Apply the kernel-deleted connectivity assertion
\cref{eq:cert-platonic-kernel-deleted-graph} to
\(W=\ker\Phi\), whose dimension is two or one.  Along a path between two
nonzero cards with different matchings, choose the first matching-changing
edge and a rich line containing it.  That line cannot meet
\(\mathbb P(\ker\Phi)\): its two-dimensional vector span would otherwise
map with rank at most one, making all nonkernel images projectively
proportional.  The exact deck-line types in
\cref{prop:cert-platonic-incidence-atlas} exclude type \((2,1)\) and mixed
four- or five-lines.  The selected line is therefore a kernel-avoiding
rainbow three-line.

Apply \cref{lem:platonic-rainbow-localization}.  It gives an actual
six-port map \(\Psi\) with the normalized rainbow line and
\(\rank\Psi\le3\).  The three exact bridge classifications have only six
rank-two maps and their listed rank-four frames, so \(\Psi\) has rank two.
In the common coefficient normalization,
\[
  (F^\pi_0,F^\pi_1,F^\pi_2,F^\pi_3)
  =\bigl(C/2,\pi(0,A/2,B/2)\bigr),
  \qquad \pi\in S_3.
\]
Using the Frobenius-dual quaternion basis, reconstruct the six-port signature
\[
  h_\pi:=\sum_{\mu=0}^3Q^\mu_{12}\otimes F^\pi_\mu.
\]
Contract ports \(1,3\) by the actual kernel \(Q_0=I\), and order the
surviving ports \((2,4,5,6)\) as \((1,2,3,4)\).  Direct expansion gives
support size ten for every \(q_\pi\), and exact row reduction gives
\cref{eq:platonic-rank-two-q4-minors}.  All preceding normalizations and
the final contraction are actual; undoing them gives the displayed
reduction in the three standard forms.  In the second octahedral form the
coefficient transport identifies
\(q'_\pi=(H_{\mathcal O}^{-1})^{\otimes4}q_\pi\).  The corresponding
localizations and
contraction kernels are checked separately in actual
\(\mathcal O'\)-representatives.  Local invertibility preserves the three
flattening ranks, but no support-size claim is transported.
\end{proof}

\subsubsection{Two-sided root geometry at the minimum quaternary core}
\label{proofsubsec:cert-platonic-q4-preservers}

For a four-port tensor, contracting either of two complementary pairs
gives adjoint linear maps on the matrix space, as in
\cref{eq:platonic-q4-map-adjoint}.  Thus the group-card condition becomes a
two-sided preservation condition on a finite set of quaternion-coordinate
lines.  We classify these maps by rank and reconstruct the original
tensor from each case.  This is the finite geometric input to
\cref{lem:platonic-quaternary-boundary}; the root-system terminology is
defined in \cref{subsubsec:root-incidence}.  The fixed actual representatives
from \cref{eq:platonic-actual-kernel-lift}, written in quaternion
coordinates, give the following concrete root-point configurations:
\[
\begin{aligned}
 \mathcal R_{A_4}&:=\{[c_Q(K)]:[K]\in\mathcal T\}
    \cong\mathcal R(D_4),\\
 \mathcal R_{S_4}&:=\{[c_Q(K)]:[K]\in\mathcal O\}
    \cong\mathcal R(F_4),\\
 \mathcal R_{A_5}&:=\{[c_Q(K)]:[K]\in\mathcal I\}
    \cong\mathcal R(H_4).
\end{aligned}
\]
where the displayed coordinate models are those of
\cref{eq:tetrahedral-domain,eq:octahedral-domain,eq:icosahedral-domain}
and the displayed isomorphisms mean projective identification in the sense of
\cref{subsubsec:group-projective,subsubsec:root-incidence}.  Their sizes are
\(12,24,60\).  Here \(D_4\) is the root system, not the order-four
dihedral group \(D_4\cong V_4\) used earlier.  Recall the Frobenius-form
matrix and define the associated
adjoint map by
\[
  J_{\rm Fr}:=\diag(1,-1,1,-1),
  \qquad A^\sharp:=J_{\rm Fr}A^{\mathsf T}J_{\rm Fr}.
\]
All calculations are exact over \(\mathbb Q(i)\) for \(D_4,F_4\) and
over \(\mathbb Q(\phi,i)\) for \(H_4\).

\begin{proposition}[Exact Platonic two-sided preserver classification]
\label{prop:cert-platonic-q4-preservers}
For \(G\in\{A_4,S_4,A_5\}\), let \(\mathcal R_G\) denote the corresponding
root-line configuration and define
\begin{equation}
 \mathscr E_G=
 \{A\in\operatorname{End}(\mathbb C^4):
   \ \forall[r]\in\mathcal R_G,\ 
   (Ar=0\ \text{or}\ [Ar]\in\mathcal R_G),\
   (A^\sharp r=0\ \text{or}\ [A^\sharp r]\in\mathcal R_G)\}.
 \label{eq:cert-q4-two-sided-endomorphisms}
\end{equation}
Here \(r\) denotes any nonzero representative of \([r]\); the two
conditions are independent of that choice.
Up to nonzero scalar, the nonzero members of \(\mathscr E_G\), together
with the single zero map recorded in the rank-zero column, have the exact
rank distribution
\begin{equation}
\begin{array}{crrrrr}
\toprule
G&\rank0&\rank1&\rank2&\rank3&\rank4\\ \midrule
A_4&1&144&0&0&576\\
S_4&1&576&2592&0&1152\\
A_5&1&3600&0&0&7200.\\
\bottomrule
\end{array}
\label{eq:cert-q4-preserver-counts}
\end{equation}
Rank one reconstructs a binary-pair tensor factor.  Rank two occurs only
for \(S_4\); all \(2592\) maps form one orbit under actual local
octahedral actions in each current marked form.  In the first normal form the reconstructed tensors
form the actual local-\(\mathcal O\) orbit of \(\Eq_4\); in the second
normal form they admit actual local-\(\mathcal O'\) dressings to the endpoint-nondegenerate eight-vertex signature
\[
  M_{\mathcal O'}(x)=
  \one[\wt(x)\equiv0\pmod2](-1)^{\sum_{i<j}x_ix_j}.
\]
In the mixed forms the normalized directly reconstructed survivors have
the even-entry vectors in
\cref{eq:p1-platonic-q4-mixed-eight-vertex}, all odd entries zero, and
ranks \((2,2,2)\).
Every rank-four map has one of
the forms
\begin{equation}
  K\longmapsto UKV
  \qquad\text{or}\qquad
  K\longmapsto UK^{\mathsf T}V,
  \label{eq:cert-q4-full-rank-preservers}
\end{equation}
and reconstructs a product of two nonsingular binaries on a crossing
matching.  The rank counts hold for all four fixed-equality octahedral
forms \(\mathcal O,\mathcal O',\mathcal O_{dt},\mathcal O_{td}\), while
their actual rank-two terminals are the ones just stated.
\end{proposition}

\begin{proof}
We first exclude rank three without enumerating arbitrary linear maps.
In the following table, the secants are the distinct projective lines
spanned by pairs of points of \(\mathcal R_G\), while the plane sections
are the distinct intersections with planes spanned by noncollinear triples
of those points; exponents record the number of distinct lines or sections
of each size.
Exact row reduction gives the complete line and plane-section data
\begin{equation}
\begin{array}{ccc}
\toprule
G&\text{secant-line sizes}&\text{plane sections}\\ \midrule
A_4&2^{18},3^{16}&3^{12},6^{12}\\
S_4&2^{72},3^{32},4^{18}&4^{96},9^{24}\\
A_5&2^{450},3^{200},5^{72}&4^{600},6^{660},15^{60}.\\
\bottomrule
\end{array}
\label{eq:cert-q4-plane-sections}
\end{equation}
For a projective point \(k=[\mathbf k]\), let \(d(k)\) be the number of
distinct nonzero projective classes
\([\mathbf r+\Span_{\mathbb C}\{\mathbf k\}]\) in
\(\mathbb P(\mathbb C^4/\Span_{\mathbb C}\{\mathbf k\})\), as
\(r=[\mathbf r]\) ranges over the root points other than \(k\).  If
\(k\) lies on at
most one secant containing two root
points, the maximum line sizes \(3,4,5\) give
\begin{equation*}
 d(k)\ge10,\qquad d(k)\ge21,\qquad d(k)\ge56,
\end{equation*}
respectively.  The centers lying on at least two secants are the finite
set of pairwise secant intersections.  Exact quotient row reduction gives
\begin{equation}
\begin{array}{cc}
\toprule
A_4&7^{12},9^{12}\\
S_4&13^{24},19^{96}\\
A_5&31^{60},49^{300},51^{360},55^{600}.\\
\bottomrule
\end{array}
\label{eq:cert-q4-secant-directions}
\end{equation}
A rank-three map is projection from its kernel point to its image plane,
and the induced map from the quotient is injective.  Hence its quotient
directions must fit into a plane section of the target configuration.
The bounds and every count in \cref{eq:cert-q4-secant-directions} exceed
the plane maxima \(6,9,15\) in \cref{eq:cert-q4-plane-sections}.  This
proves the rank-three column of \cref{eq:cert-q4-preserver-counts}.

We next classify rank two.  For \(A\in\mathscr E_G\), define the projective
image line \(W:=\mathbb P(\operatorname{im}A^\sharp)\).  Because the root points span
\(\mathbb C^4\), \(W\) is a secant line.  Moreover
\(\ker A=(\operatorname{im}A^\sharp)^{\perp_{J_{\rm Fr}}}\).  Projecting the root points modulo this kernel
gives the following complete table:
\begin{equation*}
\begin{array}{cc}
\toprule
G&( |W\cap\mathcal R_G|,
   \#\text{ quotient directions},\#\text{ killed points})\\ \midrule
A_4&(2,4,2)^{18},\ (3,6,0)^{16}\\
S_4&(2,6,2)^{72},\ (3,6,3)^{32},\ (4,4,4)^{18}\\
A_5&(2,12,2)^{450},\ (3,12,3)^{200},\ (5,10,5)^{72}.\\
\bottomrule
\end{array}
\end{equation*}
The image of \(A\) is also a projective line, containing at most
\(3,4,5\) root points.  The middle entry in every row is larger than that
bound except for \((4,4,4)^{18}\) in the octahedral row.  The eighteen
surviving four-lines form one orbit under actual left and right
\(\mathcal O\)-actions.  There are eighteen choices for
\(\mathbb P(\operatorname{im}A^\sharp)\), eighteen for
\(\mathbb P(\operatorname{im}A)\).  Each surviving four-point set is
harmonic in the cross-ratio convention of
\cref{subsubsec:root-incidence}, and exactly eight
projective isomorphisms between any two such lines preserve its harmonic
pairings.  Thus there are exactly eight harmonic projectivities, and
consequently
\begin{equation*}
  18\cdot18\cdot8=2592
\end{equation*}
projective rank-two maps survive.  They form one actual local-octahedral
orbit of the rank-two endomorphism
\begin{equation*}
  A_0:=\diag(1,1,0,0).
\end{equation*}
Bilinear-dual reconstruction, including the fixed output \(X\) in the
transfer convention \(L_p(K)=\Phi_p(K)X\), gives an actual local-octahedral
equivalent of the quaternary signature
\begin{equation*}
 q_0(x):=\frac12\one[x_1+x_2+x_3+x_4\equiv0\pmod2].
\end{equation*}
The actual octahedral matrix
\(W_0=\begin{psmallmatrix}1&1\\1&-1\end{psmallmatrix}\), with
\([W_0]=[Q_1+Q_3]\in\mathcal O\), satisfies
\(W_0^{\otimes4}\Eq_4=
2\one[x_1+x_2+x_3+x_4\equiv0\pmod2]\) and
\(W_0^{\otimes4}q_0=4\Eq_4\).  It follows that every
rank-two survivor has, for some actual matrices \(U_j\) with
\([U_j]\in\mathcal O\),
\begin{equation*}
  q=\lambda(U_1\otimes U_2\otimes U_3\otimes U_4)\Eq_4,
  \qquad \lambda\ne0.
\end{equation*}
Conversely, direct substitution shows that this orbit satisfies the
two-sided condition.  The complete four-line orbit above leaves no
additional projective rank-two degeneration; every surviving tensor has
flattening ranks \((2,2,2)\).

For the second octahedral form, analytic conjugation by
\(H_{\mathcal O}\) transports
the two-sided linear equations and therefore preserves the count 2592 and
the rank pattern, but it does not transport actual local gadgets.  Indeed,
the actual \((\mathcal O')^{\otimes4}\)-orbit of the second-form survivor
has 2592 projective tensors with support histogram
\begin{equation*}
  8^{16},\qquad10^{128},\qquad12^{384},\qquad16^{2064},
\end{equation*}
and hence does not meet the actual \((\mathcal O')^{\otimes4}\)-orbit of
the support-two signature \(\Eq_4\).  For the
normalized second-form rank-two survivor, use the four actual matrices
\begin{equation}
\begin{aligned}
 L_1&=\theta_{\mathcal O}(Q_0+Q_2),\\
 L_2=L_3&=\theta_{\mathcal O}(Q_0-Q_1+Q_2-Q_3),\\
 L_4&=X\theta_{\mathcal O}(Q_2-Q_3).
\end{aligned}
\label{eq:cert-platonic-q4-Oprime-dressing}
\end{equation}
Exact substitution gives the endpoint-nondegenerate eight-vertex signature
even-entry vector
\begin{equation}
  (1,-1,-1,-1,-1,-1,-1,1),
  \label{eq:p1-cert-q4-Oprime-eight-vertex-vector}
\end{equation}
with every odd entry zero and ranks \((2,2,2)\).  Acting before and after
this certificate by the actual group proves the stated terminal for the
entire second-form orbit.  No use of \(H_{\mathcal O}\) as a gadget is made.

For \(\alpha\in\{dt,td\}\), complementary-pair reconstruction of the
same normalized coefficient map uses the invertible coefficient operator
\begin{equation*}
 \mathscr T_\alpha=
 H_\alpha^{\mathsf T}\otimes H_\alpha^{-1}\otimes H_\alpha^{-1}
 \otimes XH_\alpha^{\mathsf T}X.
\end{equation*}
Direct exact substitution in \(q_0\) gives, projectively,
\begin{equation*}
 \mathbf e(\mathscr T_{dt}q_0)=(1,1,1,1,1,1,1,1),
 \qquad
 \mathbf e(\mathscr T_{td}q_0)=(1,i,i,i,1,1,1,i),
\end{equation*}
with all odd entries zero.  These are the direct current-coordinate
survivors, not tensors obtained by installing \(H_\alpha\) as a gadget.
Exact row reduction gives ranks \((2,2,2)\); acting by actual
\(\mathcal O_\alpha\)-matrices covers the entire rank-two orbit.

The rank-zero member is the zero map.  Because the quaternion basis spans
\(\Mat_2(\mathbb C)\), its reconstructed quaternary is zero.
For rank one, choose nonzero column vectors \(u,v\in\mathbb C^4\) and write
\(A=uv^{\mathsf T}\).  The two-sided condition forces
\([u],[J_{\rm Fr}v]\in\mathcal R_G\); conversely every such pair works.  This gives
\(|G|^2=144,576,3600\) projective maps and a rank-one
\(p\mid\bar p\) flattening of the reconstructed tensor.

Finally, a rank-four member of \(\mathscr E_G\) permutes
\(\mathcal R_G\).  Exact root-incidence automorphism classification gives
precisely \cref{eq:cert-q4-full-rank-preservers}, with
\begin{equation*}
 [V]\in N_{\PGL_2}(G),\qquad [UV]\in G.
\end{equation*}
The projective counts of the two displayed families are
\begin{equation*}
  2|G|\,|N_{\PGL_2}(G)|
  =2\cdot12\cdot24,\quad2\cdot24^2,\quad2\cdot60^2,
\end{equation*}
namely \(576,1152,7200\).  If the physical card map is
\(K\mapsto UKV\), undoing the output \(X\) gives
\(q_{uvyz}=U_{yu}(VX)_{vz}\); the transpose case gives
\(q_{uvyz}=U_{yv}(VX)_{uz}\).  In either case the quaternary is a product
of two nonsingular binaries on a crossing matching.

The use of \(A^\sharp\) in
\cref{eq:cert-q4-two-sided-endomorphisms} is indispensable.  The one-sided
\(A_4\) preserver equations have \(1536\) projective rank-two solutions,
but every one fails the complementary \(A^\sharp\) equations.  Thus the
zero in the \(A_4\) rank-two column is a two-sided conclusion, not a
one-sided assertion.

The coefficient similarities identify the four rank tables.
The formulas
\cref{eq:cert-platonic-q4-Oprime-dressing,eq:p1-cert-q4-Oprime-eight-vertex-vector}
and the two direct mixed reconstructions supply the required physical
terminals in the respective current forms.

For all three rows, exact enumeration generates the root configurations and
then recovers every secant, plane section, multiple-secant center, and
quotient direction by characteristic-zero row reduction: over
\(\mathbb Q(i)\) for \(D_4,F_4\), and over
\(\mathbb Q(\phi,i)\) for \(H_4\).  It enumerates all projectivities between
the surviving octahedral four-lines, giving \(18^2\cdot8=2592\).
For the full-rank case it constructs the complete finite incidence
hypergraph of root points, rich lines, and plane sections, enumerates its
automorphisms, solves for the inducing projective linear map, and retains the
map only after checking both \(A\) and \(A^\sharp\) on every root point.
The resulting automorphism counts equal the cardinalities of the explicit
left--right/transpose families for \(D_4,F_4,H_4\), proving completeness of
all three full-rank lists rather than merely verifying their displayed
members.
\end{proof}

\subsubsection{Terminal certificates for the Platonic quaternaries}
\label{proofsubsec:cert-platonic-q4-exits}

This subsection records the finite terminal calculation used in
\cref{lem:platonic-finite-q4-terminal}.  Ports of every quaternary are numbered
\(1,2,3,4\) in their displayed order.  If ports \(r,s\) are contracted by
an actual kernel \(K\), the surviving ordered binary is denoted
\(B=\partial_{rs}^{K}q\).  The pairing is bilinear, with no conjugation.
Every group-membership test below is made on the transfer class
\begin{equation*}
  T_B=BX,
\end{equation*}
and never on \([B]\) alone.

\begin{proposition}[Platonic quaternary exit certificate]
\label{prop:cert-platonic-q4-exits}
For each quaternary in the following list, let \(G\) denote the marked
Platonic group from whose row the quaternary is obtained.  Every one of the
six rank-deficient quaternaries \(q_\pi\) constructed in
\cref{prop:cert-platonic-deficient-mixed} and satisfying
\cref{eq:platonic-rank-two-q4-minors}, one of the three
orientation-reversing tetrahedral outputs \(q_\sigma\) specified in
\cref{prop:p1-cert-A4-eight-vertex-orientation-terminal}, one of the twenty-four
octahedral separator outputs (in any of the four marked forms) specified in
\cref{prop:p1-cert-S4-uniform-eight-vertex-separator}, or one of the thirty
icosahedral frame outputs \(q_{u_1,u_2,u_3}\) in
\cref{eq:cert-A5-output,eq:cert-A5-frames} has one of the
following directly realized signatures:
\begin{enumerate}
\item an endpoint-nondegenerate eight-vertex signature with flattening ranks
      \((2,2,2)\);
\item a pure generalized equality of arity four;
\item a nonsingular ordered binary \(B\) whose transfer class \([BX]\) is
      outside \(G\).
\end{enumerate}
For either mixed octahedral form the list also includes the six direct
current-coordinate counterparts of the rank-deficient outputs.  This holds for
all four octahedral distinguished-involution normal forms and for both
icosahedral coordinate displays.  Every such signature uses one
quaternary vertex and at most four actual group binaries.
\end{proposition}

\begin{proof}
First consider the six support-ten rank-deficient outputs in
\cref{prop:cert-platonic-deficient-mixed}.  For each fixed output, exact
enumeration of all six port pairs and all actual kernels gives the following
table.  ``Inside'' means \([BX]\) is in the displayed group, and
``singular'' counts nonzero singular cards; zero cards are kept separate.
\begin{center}
\small
\setlength{\tabcolsep}{4pt}
\begin{tabular}{crrrrrc}
\toprule
group&cards per \(q\)&zero&inside&outside&singular&one uniform outsider\\
\midrule
\(A_4\)&72&3&9&60&0&\(\partial_{12}^{Q_0}q\)\\
\(S_4\)&144&6&108&30&0&\(\partial_{24}^{Q_0}q\)\\
\(A_5\)&360&3&36&321&0&\(\partial_{12}^{Q_0}q\)\\
\bottomrule
\end{tabular}
\end{center}
Thus the same displayed pair and kernel work for all six outputs in each
row.  The first two rows are computed over \(\mathbb Q(i)\), and the
\(A_5\) row is computed over \(\mathbb Q(\phi,i)\).

For the second octahedral form, define
\(\sigma_{\mathcal O}:=\sqrt2\) and the transport matrix
\(S_{\mathcal O}:=H_{\mathcal O}^{-1}\).  The transported rank-deficient outputs
are \(q'_\pi=S_{\mathcal O}^{\otimes4}q_\pi\).  The single actual
kernel
\begin{equation}
 K'_{\mathrm{def}}
 =H_{\mathcal O}^{-1}(Q_0-Q_1-Q_2-Q_3)H_{\mathcal O},
 \qquad [K'_{\mathrm{def}}]\in\mathcal O'
 \label{eq:cert-Oprime-deficient-kernel}
\end{equation}
on ports \(1,2\) works for all six.  Direct expansion gives
\(\det(\partial_{12}^{K'_{\mathrm{def}}}q'_\pi)=-3/512\).  As
\(\pi\) ranges over \(S_3\), the pulled-back transfer directions
\(H_{\mathcal O}(B'X)H_{\mathcal O}^{-1}\), in quaternion coordinates,
form the multiset
\begin{equation*}
 [3,-1,1,1],\ [1,-1,0,1],\ [3,-1,1,1],\
 [1,-1,0,1],\ [3,-1,-1,1],\ [3,-1,-1,1].
\end{equation*}
None is in \(\mathcal O\): its support has size three, or size four with
unequal nonzero magnitudes.  Hence every corresponding \([B'X]\) is outside
\(\mathcal O'=H_{\mathcal O}^{-1}\mathcal O H_{\mathcal O}\).

For \(\mathcal O_{dt}\) and \(\mathcal O_{td}\), the final exact table in
\cref{prop:p1-cert-S4-uniform-eight-vertex-separator} gives one actual
current-group kernel on ports \(1,2\), a common nonzero determinant, and
all six pulled-back outside-transfer directions.  This proves the rank-deficient
rows directly in the retained fixed-\(I\) coordinates.

The notation here is compatible with the earlier separator transport
\(q'_R=\mathcal L_{\mathcal O}q_R\), but the two formulas use different
displayed port-orientation conventions.  The matching-dependent \(X\)-factors
in \cref{eq:cert-S4-coefficient-transport,eq:cert-Oprime-transport} convert
the alternating
\(H_{\mathcal O}^{-1}/H_{\mathcal O}^{\mathsf T}\) convention of
\(\mathcal L_{\mathcal O}\) into the uniform
\(S_{\mathcal O}^{\otimes4}\) convention after those fixed orientation
factors are absorbed into the current \(\mathcal O'\)-representatives.
Thus both notations parametrize the same marked second-form coefficient
families in their stated port orders; neither equality is being used as a
physical gadget.

For the three tetrahedral outputs \(q_\sigma\) specified in
\cref{prop:p1-cert-A4-eight-vertex-orientation-terminal}, the one actual kernel
\begin{equation*}
 K_{\mathcal T}^{\mathrm{out}}
 =Q_0+Q_1-Q_2+Q_3,
 \qquad [K_{\mathcal T}^{\mathrm{out}}]\in\mathcal T
\end{equation*}
on ports \(1,2\) gives, projectively,
\begin{center}
\small
\setlength{\tabcolsep}{5pt}
\begin{tabular}{ccc}
\toprule
\(\sigma\)&\(B=\partial_{12}^{K_{\mathcal T}^{\mathrm{out}}}q_\sigma\)&\(BX\)\\
\midrule
\((0,3,1)\)&\(\begin{psmallmatrix}1&1\\-1&1\end{psmallmatrix}\)&
 \(\begin{psmallmatrix}1&1\\1&-1\end{psmallmatrix}\)\\
\((1,0,3)\)&\(\begin{psmallmatrix}1&-i\\-i&1\end{psmallmatrix}\)&
 \(\begin{psmallmatrix}-i&1\\1&-i\end{psmallmatrix}\)\\
\((3,1,0)\)&\(\diag(1,i)\)&
 \(\begin{psmallmatrix}0&1\\i&0\end{psmallmatrix}\)\\
\bottomrule
\end{tabular}
\end{center}
All three determinants are nonzero, and exact projective comparison with
the twelve points in \cref{eq:tetrahedral-domain} excludes all three
transfer classes.

We next treat the first octahedral form.  Define the octahedral sign matrices by
\[
 S_{\epsilon_1\epsilon_2\epsilon_3}
 =Q_0+\epsilon_1Q_1+\epsilon_2Q_2+\epsilon_3Q_3.
\]
For an ordered frame in \cref{eq:cert-S4-four-frames}, its separator is
\(q=C/2+W_1\), where \(W_1\) is the first entry.  The twelve ordered
outputs with
\[
 W_1\in\{R_0,R_3,R_{03}^{+-},R_{03}^{-+}\}
\]
are endpoint-nondegenerate eight-vertex signatures on all eight even-parity words and have ranks
\((2,2,2)\).  The four with \(W_1=R_1\) have support
\(\{0011,1100\}\); actual \(X\)-dressings on the last two ports turn them
into pure generalized equalities.  The remaining eight outputs occur twice
for each row of the following table:
\begin{center}
\small
\setlength{\tabcolsep}{5pt}
\begin{tabular}{cccc}
\toprule
\(W_1\)&kernel on \(12\)&\(B\)&\([BX]\)\\
\midrule
\(R_{13}^{++}\)&\(S_{++-}\)&
 \(\begin{psmallmatrix}1+i&-i\\-i&1-i\end{psmallmatrix}\)&
 \([Q_0+Q_1-Q_2]\)\\
\(R_{13}^{--}\)&\(S_{+++}\)&
 \(\begin{psmallmatrix}1-i&-i\\-i&1+i\end{psmallmatrix}\)&
 \([Q_0+Q_1+Q_2]\)\\
\(R_{01}^{+-}\)&\(S_{+++}\)&
 \(\begin{psmallmatrix}1-i&-1\\1&1+i\end{psmallmatrix}\)&
 \([Q_1+Q_2-Q_3]\)\\
\(R_{01}^{-+}\)&\(S_{++-}\)&
 \(\begin{psmallmatrix}1+i&-1\\1&1-i\end{psmallmatrix}\)&
 \([Q_1-Q_2-Q_3]\)\\
\bottomrule
\end{tabular}
\end{center}
Every determinant equals \(3\).  The points of \(\mathcal O\) in
\cref{eq:octahedral-domain} have quaternion-coordinate support size one,
two, or four, never three, so all four displayed transfer directions are
outside \(\mathcal O\).

\begin{sloppypar}
For the second octahedral form, we compute the actual kernels in its
fixed coordinates.  Define the nonzero scalar
\(c_{\mathcal O}:=2\sigma_{\mathcal O}(1-i)\).  Direct multiplication gives
\begin{equation}
 H_{\mathcal O}^{\mathsf T}XH_{\mathcal O}=c_{\mathcal O}I,
 \qquad
 S_{\mathcal O}US_{\mathcal O}^{\mathsf T}
 =c_{\mathcal O}^{-1}H_{\mathcal O}^{-1}(UX)H_{\mathcal O}.
 \label{eq:cert-Oprime-transport}
\end{equation}
Thus \(S_{\mathcal O}^{\otimes4}\) transports every
\(\mathcal O\)-matching factor to
an \(\mathcal O'\)-matching factor.  More precisely, if \(h\) is the
standard-form six-port ancestor with card map \(\Phi\), define, in its
displayed port order,
\[
 h':=(H_{\mathcal O}^{\mathsf T}\otimes H_{\mathcal O}^{-1}
       \otimes S_{\mathcal O}^{\otimes4})h.
\]
The card map \(\Phi'\) of \(h'\) is parametrized by
\begin{equation}
 \Phi'(H_{\mathcal O}^{-1}KH_{\mathcal O})
 =S_{\mathcal O}^{\otimes4}\Phi(K).
\label{eq:cert-Oprime-card-map}
\end{equation}
This invertible coefficient correspondence exhausts the second-form
tensors.  The contractions below use the listed actual current-group
kernels, with their ordered ports and nonzero realization scalars.
\end{sloppypar}
Define the conjugated octahedral sign matrices by
\(S'_{\epsilon_1\epsilon_2\epsilon_3}
:=H_{\mathcal O}^{-1}S_{\epsilon_1\epsilon_2\epsilon_3}H_{\mathcal O}\).
For the twelve ordered
outputs in the next table, direct contraction of
\(q'=S_{\mathcal O}^{\otimes4}q\) on ports \(1,2\) gives
\(\det B'=-3/256\):
\begin{center}
\small
\setlength{\tabcolsep}{5pt}
\begin{tabular}{ccc}
\toprule
\(W_1\)&actual kernel&
\(H_{\mathcal O}(B'X)H_{\mathcal O}^{-1}\)\\
\midrule
\(R_{13}^{++}\)&\(S'_{+++}\)&\([2Q_0+Q_1+Q_2]\)\\
\(R_{13}^{--}\)&\(S'_{+-+}\)&\([Q_1+Q_2+2Q_3]\)\\
\(R_{03}^{+-}\)&\(S'_{++-}\)&\([Q_0+2Q_1+Q_3]\)\\
\(R_{03}^{-+}\)&\(S'_{+++}\)&\([Q_0+2Q_2+Q_3]\)\\
\(R_{01}^{+-}\)&\(S'_{++-}\)&\([2Q_0+Q_1-Q_2]\)\\
\(R_{01}^{-+}\)&\(S'_{+++}\)&\([2Q_0-Q_1+Q_2]\)\\
\bottomrule
\end{tabular}
\end{center}
Each row occurs for two orderings.  Every pulled-back direction has three
nonzero quaternion coordinates of magnitudes \(2,1,1\), so it is outside
\(\mathcal O\), proving \([B'X]\notin\mathcal O'\).

The remaining twelve ordered \(\mathcal O'\)-outputs have
\(W_1=R_0,R_1,R_3\), four of each.  Actual local dressings turn them into
endpoint-nondegenerate eight-vertex signatures.  In the even-word order
\[
  (0000,0011,0101,0110,1001,1010,1100,1111),
\]
the exact outputs are
\begin{center}
\small
\setlength{\tabcolsep}{5pt}
\begin{tabular}{ccc}
\toprule
\(W_1\)&dressing&even coefficients, up to scale\\
\midrule
\(R_0\)&\((S'_{++-})^{\otimes4}\)&\((1,1,-1,1,1,-1,1,1)\)\\
\(R_1\)&identity&\((1,-1,1,1,1,1,-1,1)\)\\
\(R_3\)&\((S'_{+++})^{\otimes4}\)&\((1,-1,1,1,1,1,-1,1)\)\\
\bottomrule
\end{tabular}
\end{center}
All odd-weight coefficients vanish, both endpoints are nonzero, and local
invertible dressings preserve ranks \((2,2,2)\).

The two mixed full-rank rows are the direct tables in
\cref{prop:p1-cert-S4-uniform-eight-vertex-separator}.  For
\(\mathcal O_{dt}\) they partition into endpoint and outside-binary exits;
for \(\mathcal O_{td}\) they partition into endpoint, pure-generalized-
equality, and outside-binary exits.  The tables cover every possible first
frame entry, hence all twenty-four orderings, and use only actual matrices
of the current mixed group.

Finally, for each of the thirty icosahedral outputs in
\cref{eq:cert-A5-output}, use the one actual representative
\begin{equation*}
 L_{\mathcal I}=Q_0-\tau^2Q_1+\tau Q_3
 \qquad\text{with}\qquad [L_{\mathcal I}]\in\mathcal I
\end{equation*}
on ports \(1,2\).  Exact arithmetic over \(\mathbb Q(\phi,i)\) proves that
all thirty resulting binaries are nonsingular, their thirty projective
transfers \([BX]\) are distinct, and none belongs to the complete sixty-point
set \(\mathcal I\).  Enumerating all
\(30\cdot6\cdot60=10800\) actual cards gives
\begin{equation*}
 \text{zero }180,\qquad
 \text{inside }1260,\qquad
 \text{outside }9360,\qquad
 \text{nonzero singular }0.
\end{equation*}
For the opposite display, use the fixed-equality dictionary
\cref{lem:icosahedral-fixed-I-dictionary}: replace the two standard kernels
by their actual \(\kappa_D\)-images.  Card covariance conjugates the complete
transfer group and preserves nonsingularity and group exclusion, while the
retained literal \(I\) is untouched.  This exhausts every family in the
proposition.

The displayed kernels and dressing tables prove the listed bounded
exits by exact matrix multiplication.  The complete terminal histograms
and the completeness claims used earlier are also reproduced by
\path{certificates/verify_a4_s4_certificates.py},
\path{certificates/verify_s4_second_form.py},
\path{certificates/verify_a5_certificate.py}, and
\path{certificates/verify_platonic_extended.py} over the number fields
specified above.  In particular, the full preserver and hereditary eight-port classifications retain
the exhaustive computations described in their own proofs.
\end{proof}

\begin{theorem}[Platonic bridge interface]
\label{thm:cert-interface-platonic-bridges}
Under the normalized rainbow and separated-sign hypotheses of
\cref{eq:platonic-rainbow-normalization,eq:platonic-rainbow-sum,%
eq:platonic-separated-sign-tests}, the following complete bridge statements
hold.
\begin{enumerate}[label=\textnormal{(\roman*)}]
\item For each marked octahedral form, there are exactly six rank-deficient maps
      and \(24\) full-rank maps, the latter being the six orderings of each
      of
      \[
      \{R_0,R_1,R_3\},\quad
      \{R_0,R_{13}^{++},R_{13}^{--}\},\quad
      \{R_1,R_{03}^{+-},R_{03}^{-+}\},\quad
      \{R_3,R_{01}^{+-},R_{01}^{-+}\}.
      \]
      The fixed actual kernel is
      \(K_{\mathcal O}=Q_0+Q_1\) in the first form and
      \(K_{\mathcal O'}=\theta_{\mathcal O}(K_{\mathcal O})\) in the
      second, while
      \(K_{\mathcal O_{dt}}=\theta_{dt}(K_{\mathcal O})\) and
      \(K_{\mathcal O_{td}}=\theta_{td}(K_{\mathcal O})\) are the two
      mixed-form kernels.  In the two unmixed forms, actual current-group
      dressings send every full-rank map to an endpoint-nondegenerate
      eight-vertex signature of ranks \((2,2,2)\).  In either mixed form,
      the fixed kernel and the direct current-coordinate tables give one of
      the three exits: such an eight-vertex signature, a pure generalized
      equality, or a nonsingular outside-group binary.
\item For the internal tetrahedral form, the bridge set is
      \(\{0,A/2,B/2,R_0,R_1,R_3\}\), and the sum equation has exactly six
      rank-deficient and six full-rank solutions.  Write
      \[
      \Sigma_+=\{(0,1,3),(1,3,0),(3,0,1)\},\qquad
      \Sigma_-=\{(0,3,1),(1,0,3),(3,1,0)\}.
      \]
      For every
      \(\sigma\in\Sigma_+\), all \(180\) actual pair cards of the reconstructed
      parent are safe and the parent lies in \(\cA\).  For every
      \(\sigma\in\Sigma_-\), contraction of ports \(1,3\) by the fixed
      actual kernel \(Q_0=I\), followed by four actual dressings, gives an
      endpoint-nondegenerate eight-vertex signature of ranks \((2,2,2)\).
\item The sixty displayed projective classes form a projective group
      isomorphic to \(A_5\) and are closed under induced projective
      multiplication, inverse, and transpose.  Its
      bridge has exactly six rank-deficient and \(30\) full-rank solutions, the
      latter being the six orderings of the five displayed frames
      \(\mathcal F_0,\ldots,\mathcal F_4\).  The fixed actual kernel
      \(K_{\mathcal I}:=\phi Q_1-\tau Q_2+Q_3\) gives a quaternary of ranks
      \((3,3,3)\), and the fixed subsequent actual kernel
      \(L_{\mathcal I}:=Q_0-\tau^2Q_1+\tau Q_3\) gives a nonsingular binary whose transfer lies
      outside \(\mathcal I\), for all \(30\) solutions.  The statement
      includes arbitrary complex scales and cancellations.  For
      \(\mathcal I_-\), the actual kernels
      \(K_{\mathcal I_-}=\kappa_D(K_{\mathcal I})\) and
      \(L_{\mathcal I_-}=\kappa_D(L_{\mathcal I})\) give the corresponding
      direct fixed-\(I\) identities.
\end{enumerate}
\end{theorem}

\begin{proof}
Clause~\textnormal{(i)} is the exact octahedral bridge
\cref{eq:cert-S4-four-frames} together with the separate physical certificate
\cref{prop:p1-cert-S4-uniform-eight-vertex-separator}.  Clause~\textnormal{(ii)}
is \cref{eq:cert-A4-bridge,eq:cert-A4-twelve-maps} and the all-card,
affine-phase, and actual-orientation checks in
Appendix~\ref{proofsubsec:cert-A4}, including
\cref{prop:p1-cert-A4-eight-vertex-orientation-terminal}.  Clause~\textnormal{(iii)}
is the exact group, bridge, frame, and rank calculation in
Appendix~\ref{proofsubsec:cert-A5}, followed by the uniform physical exit in
\cref{prop:cert-platonic-q4-exits};
\cref{lem:icosahedral-fixed-I-dictionary} gives the opposite display with
actual \(\kappa_D\)-kernels and leaves the supplied equality unchanged.
\end{proof}

\begin{theorem}[Classification of Platonic quaternaries]
\label{thm:cert-interface-platonic-q4}
Let \(G\in\{\mathcal T,\mathcal O,\mathcal O',\mathcal O_{dt},
\mathcal O_{td},\mathcal I,\mathcal I_-\}\).
The finite quaternary certificates have the following directly usable form.
\begin{enumerate}[label=\textnormal{(\roman*)}]
\item Let \(q\) be quaternary and define
      \(L_p(K)=\Phi_p^q(K)X\).  If, for every one of the six pairs \(p\)
      and every fixed actual representative \(K_g\),
      \[
      L_p(K_g)=0\qquad\hbox{or}\qquad[L_p(K_g)]\in G,
      \]
      then, for each pair \(p\), put \(r_p:=\rank L_p\).  One has
      \(r_p\ne3\).  If \(r_p=0\), then \(q=0\); if
      \(r_p\in\{1,4\}\), then \(q\) has a binary-pair factorization; and
      \(r_p=2\) can occur only for the octahedral forms: it is actually
      dressable to \(\Eq_4\) for \(\mathcal O\), and to an
      endpoint-nondegenerate eight-vertex signature of ranks \((2,2,2)\)
      for \(\mathcal O'\), \(\mathcal O_{dt}\), and \(\mathcal O_{td}\),
      with the mixed vectors in
      \cref{eq:p1-platonic-q4-mixed-eight-vertex}.
\item Let \(h\ne0\) have arity six and let \(\Phi_p^h\) have rank two or
      three.  Suppose, for every \(g\in G\), the card
      \(\Phi_p^h(K_g)\) formed with its fixed actual representative is zero
      or a nonsingular \(G\)-matching product, with two nonzero cards in
      this \(p\)-family using different residual matchings.  Actual localization
      produces a quaternary of ranks \((3,3,3)\), and an actual contraction
      of that quaternary gives a nonsingular binary \(B\) with
      \([BX]\notin G\).  This includes the separately verified
      \(\mathcal O'\) and mixed representatives and both icosahedral
      displays.
\item For each fixed \(G\), every corresponding enumerated rank-deficient
      quaternary has the following conclusion.  The same conclusion holds
      for the tetrahedral-orientation outputs when \(G=\mathcal T\), the
      octahedral-frame outputs when
      \(G\in\{\mathcal O,\mathcal O',\mathcal O_{dt},\mathcal O_{td}\}\),
      and the icosahedral-frame outputs
      when \(G\in\{\mathcal I,\mathcal I_-\}\).  Each such quaternary
      directly realizes one of: an
      endpoint-nondegenerate eight-vertex signature of ranks \((2,2,2)\), a
      pure generalized equality of arity four, or a nonsingular binary
      \(B\) with \([BX]\notin G\).  Each realization uses the quaternary and
      at most four actual \(G\)-binaries.
\end{enumerate}
\end{theorem}

\begin{proof}
For clause~\textnormal{(i)}, use
\cref{lem:icosahedral-fixed-I-dictionary} for the \(\mathcal I_-\)
coefficient calculation, replacing every physical kernel by its actual
\(\kappa_D\)-image and leaving the literal \(I\) fixed.  Fix a pair \(p\), and let \(A_p\)
be the matrix of \(L_p\) in the corresponding quaternion root coordinates.
The card conditions at \(p\) give the first condition in
\cref{eq:cert-q4-two-sided-endomorphisms}.  For the complementary pair
\(\bar p\), let \(\mathbf R_X\) be the coordinate matrix of right
multiplication \(K\mapsto KX\).  Direct complementary-pair reshaping gives
\(A_{\bar p}=\mathbf R_XA_p^\sharp\mathbf R_X\); the fixed \(X\)-reversals
permute the root configuration, so the card conditions at \(\bar p\) give
the second condition.  Hence \(A_p\) belongs to the corresponding one of
\(\mathscr E_{A_4},\mathscr E_{S_4},\mathscr E_{A_5}\), and
\cref{prop:cert-platonic-q4-preservers} gives all stated rank and
reconstruction alternatives, including the separately verified
\(\mathcal O'\) terminal and the direct mixed reconstructions above.

Clause~\textnormal{(ii)} combines the actual localization and rank
calculation in \cref{prop:cert-platonic-deficient-mixed} with the
corresponding uniform outside-binary row of
\cref{prop:cert-platonic-q4-exits}.  Clause~\textnormal{(iii)} is the
representative-wise terminal statement of the latter proposition.  The
card and network covariance in
\cref{lem:icosahedral-fixed-I-dictionary} supplies all three clauses for
\(\mathcal I_-\) directly with actual \(\kappa_D\)-kernels.
\end{proof}

\bibliographystyle{alpha}
\bibliography{paper1}

\end{document}
